\RequirePackage{fix-cm}
\documentclass[fontsize=13pt,twoside,openany,headings=normal,parskip=false,bookmarkpackage=false,DIV=calc]{scrbook}
\usepackage[T1]{fontenc}
\usepackage[utf8]{inputenc}
\usepackage{palatino}
\usepackage{microtype}
\usepackage{amsmath,amssymb}
\allowdisplaybreaks
% Preserve natural paragraph/display spacing when material moves to the next page.
\raggedbottom
% Discourage isolated paragraph lines while allowing long footnotes to split naturally.
\clubpenalty=1000
\widowpenalty=1000
\displaywidowpenalty=1000
\usepackage{graphicx}
\renewcommand{\thechapter}{\Roman{chapter}}
\renewcommand{\thesection}{\arabic{section}}
\renewcommand*{\chapterformat}{\thechapter.\enskip}
\renewcommand*{\sectionformat}{\thesection.\enskip}
\counterwithout{figure}{chapter}
\renewcommand*{\captionformat}{}
\setcounter{tocdepth}{1}
\setcounter{secnumdepth}{1}
\DeclareTOCStyleEntry[dynnumwidth=true]{tocline}{chapter}
\linespread{1.1}\selectfont
\usepackage[
  b5paper,left=0.7cm,right=0.7cm,
  top=0.91cm,bottom=0.91cm,
  includehead,nofoot,
  headheight=\baselineskip,
  headsep=0.55\baselineskip
]{geometry}
\usepackage{scrlayer-scrpage}
\clearpairofpagestyles
\ohead{\ifnum\value{page}=1\relax\else\pagemark\fi}
\setkomafont{pageheadfoot}{\normalfont\normalsize}
\setkomafont{pagenumber}{\normalfont\normalsize}
\setkomafont{disposition}{\normalfont\bfseries}
\setkomafont{descriptionlabel}{\normalfont\bfseries}
\pagestyle{scrheadings}
\renewcommand*{\chapterpagestyle}{scrheadings}
\usepackage{indentfirst}
\setlength{\parindent}{1.5em}
\deffootnote[2.5em]{2.5em}{1em}{\textsuperscript{\thefootnotemark}\enskip}
\usepackage{xcolor}
\definecolor{LinkColor}{HTML}{255A83}
\usepackage[colorlinks=true,allcolors=LinkColor,unicode=true,hypertexnames=false,hyperfootnotes=false]{hyperref}
\usepackage[all]{hypcap}
\AtBeginDocument{\pagenumbering{arabic}\thispagestyle{empty}}

% Original notes are true page footnotes with their original numbering.
% Translator notes within original notes are explicitly marked inline annotations.
\usepackage{perpage}
\newcounter{translatornote}
\MakeSorted{translatornote}
\newcounter{translatornoteid}
\renewcommand{\thetranslatornote}{T\arabic{translatornote}}
\setlength{\skip\footins}{0.8\baselineskip plus 0.2\baselineskip minus 0.1\baselineskip}
\renewcommand{\footnoterule}{\kern-0.4\baselineskip\hrule width 0.4\columnwidth height 0.03em\kern\dimexpr0.4\baselineskip-0.03em\relax}
\makeatletter
\deffootnote[2.5em]{2.5em}{1em}{\@makefnmark\enskip}
\ExplSyntaxOn
\tl_new:N \l__vn_note_target_tl
\tl_new:N \l__vn_note_display_tl
\bool_new:N \l__vn_in_original_note_bool
\cs_set_protected:Npn \@makefnmark
  {\hbox{\hyperlink{\l__vn_note_target_tl}{\textsuperscript{\normalfont\@thefnmark}}}}
\cs_new_protected:Npn \__vn_page_footnote:nn #1#2
  {\protected@edef\@thefnmark{#1}\@footnotemark\@footnotetext{#2}}
\NewDocumentCommand \originalnote {m +m}
  {\cs_gset_protected:cpn {vn_original_#1}
    {\__vn_page_footnote:nn{#1}
      {\bool_set_true:N \l__vn_in_original_note_bool
       \hypertarget{note-#1}{}
       \protected@edef\@currentlabel{#1}\def\@currentHref{note-#1}
       \label{orig-note-#1}#2}}}
\NewDocumentCommand \origNote {m}
  {\group_begin:
    \tl_set:Nn \l__vn_note_target_tl {note-#1}
    \cs_if_exist:cTF {vn_original_seen_#1}
      {\textsuperscript{\hyperlink{note-#1}{#1}}}
      {\cs_gset:cpn {vn_original_seen_#1}{}
       \hypertarget{note-call-#1}{}\use:c {vn_original_#1}}
    \group_end:}
\cs_new_protected:Npn \__vn_translator_body:n #1
  {\hypertarget{\l__vn_note_target_tl}{}
   \protected@edef\@currentlabel{\l__vn_note_display_tl}
   \protected@edef\@currentHref{\l__vn_note_target_tl}
   \label{translator-note-\arabic{translatornoteid}}#1}
\NewDocumentCommand \tnote {+m}
  {\stepcounter{translatornote}\stepcounter{translatornoteid}
   \group_begin:
    \tl_set:Nx \l__vn_note_target_tl {translatornote-\arabic{translatornoteid}}
    \tl_set:Nx \l__vn_note_display_tl {\thetranslatornote}
    \nobreak\hspace{0.1em}
    \hypertarget{translatornote-call-\arabic{translatornoteid}}{}
    \bool_if:NTF \l__vn_in_original_note_bool
      {\__vn_translator_body:n
        {\textsuperscript{\hyperlink{\l__vn_note_target_tl}{\l__vn_note_display_tl}}\nobreak\hspace{0.1em}
         [\textit{Translator's~note:}~#1]}}
      {\__vn_page_footnote:nn{\l__vn_note_display_tl}{\__vn_translator_body:n{#1}}}
   \group_end:}
\ExplSyntaxOff
\makeatother
\hypersetup{pdftitle={Mathematische Grundlagen der Quantenmechanik / Mathematical Foundations of Quantum Mechanics},pdfauthor={Johann von Neumann},pdfsubject={English translation of the German 1932 text; 1981 reprint}}

% Font-relative emergency stretch; no page-specific line or page breaks.
\setlength{\emergencystretch}{2em}


% Complete original notes: definitions expand only at the first body call.
\originalnote{1}{% Source PDF page 245 / original page 238
\hypertarget{origpage-238}{}%
Among the comprehensive accounts available in German are the following: Sommerfeld, supplementary volume to the 4th edition of \emph{Atombau und Spektrallinien}, Braunschweig, 1928; Weyl, \emph{Gruppentheorie und Quantenmechanik}, Leipzig, 1928, second edition, Leipzig, 1931; Fraenkel, \emph{Einführung in die Quantenmechanik}, Berlin, 1929; Born and Jordan, \emph{Elementare Quantenmechanik}, Berlin, 1930; Dirac, \emph{Prinzipien der Quantenmechanik}, translated by W. Bloch, Leipzig, 1931.}
\originalnote{2}{See \emph{Proc. Roy. Soc.}, London, vol.~109 (1925) and subsequent volumes, especially vol.~113 (1926). Independently of Dirac, P. Jordan, \emph{Z. Physik}, vol.~40 (1926), and F. London, \emph{Z. Physik}, vol.~40 (1926), provided related approaches to the foundations of the theory.}
\originalnote{3}{See \hyperref[chap:IV]{Chapter~IV} and \hyperref[sec:VI.3]{VI.3}.}
\originalnote{4}{See \hyperref[chap:V]{Chapter~V}.}
\originalnote{5}{Its principal stages were: Planck's discovery of the quantum laws in the case of ``black'' cavity radiation (see, for example, Planck's account in his book \emph{Wärmestrahlung}, Leipzig, 1906); Einstein's demonstration of the corpuscular nature of light, the light-quantum theory, \emph{Ann. Physik} [4], vol.~17 (1905), which furnished the first instance of the wave--corpuscle duality that, as we now know, governs all microphysics; and Bohr's application of these two groups of laws to the atomic model, \emph{Fysisk Tidskr.}, vol.~12 (1914), and \emph{Z. Physik}, vol.~6 (1920).}
\originalnote{6}{Through the work of Epstein and Sommerfeld, the quantum laws supplementing the laws of mechanics were known for multiply periodic motions (see, for example, Sommerfeld, \emph{Atombau und Spektrallinien}, Braunschweig, 1924). On the other hand, it was established that a freely moving mass point, or a planet on a hyperbolic orbit, as opposed to one on an elliptic orbit, is ``unquantized.'' A complete account of this stage in the development of quantum theory will be found in Reiche, \emph{Die Quantentheorie, ihr Ursprung und ihre Entwicklung}, Berlin, 1921, and Landé, \emph{Fortschritte der Quantentheorie}, Dresden, 1922.}
\originalnote{7}{This was proved by Schrödinger, \emph{Ann. Physik} [4], vol.~79 (1926).}
\originalnote{8}{\emph{Z. Physik}, vol.~37 (1926).}
\originalnote{9}{See the papers cited in \hyperlink{note-2}{Note~2}. Schrödinger's papers have appeared in book form: \emph{Abhandlungen zur Wellenmechanik}, Leipzig, 1928.}
\originalnote{10}{The present situation may be characterized as follows. For individual electrons and the electron shells of atoms or molecules, the theory is entirely successful, both for electrostatic forces and for the electromagnetic processes involved in the production, propagation, and transformation of light. In treating atomic nuclei, however, and in attempting to construct a general relativistic theory of electromagnetism, it appears, despite remarkable partial successes, to lead to great difficulties that can hardly be overcome without essentially new ideas.}
\originalnote{11}{As is well known, classical-mechanical motion is completely governed by the Hamiltonian function, since it gives the equations of motion
\[
\dot q_l=\frac{\partial H}{\partial p_l},\qquad
\dot p_l=-\frac{\partial H}{\partial q_l}\qquad(l=1,\ldots,k).
\]
% Source PDF page 246 / original page 239
\hypertarget{origpage-239}{}%
Before the discovery of quantum mechanics, attempts were made to account for quantum phenomena by retaining these equations of motion and imposing additional ``quantum conditions'' (see \hyperlink{note-6}{Note~6}). For every set of values of $q_1,\ldots,q_k,p_1,\ldots,p_k$ given at time $t=0$, the equations of motion determined the subsequent evolution in time: the system's ``trajectory'' in the $2k$-dimensional ``phase space'' of $q_1,\ldots,q_k,p_1,\ldots,p_k$. An additional condition therefore restricts the set of possible initial values, or trajectories, to a particular subset. Corresponding to the small number of admissible orbital levels, only a small number of energy levels are then possible. Although quantum mechanics has broken completely with this procedure, it is clear from the outset that the Hamiltonian function must also play a major role in it. For all experience establishes the validity of Bohr's correspondence principle, which states that, in the so-called limiting case of large quantum numbers, quantum theory must yield results agreeing with classical mechanics.}
\originalnote{12}{The last three concepts come from the earlier quantum theory, developed chiefly by N. Bohr before the advent of quantum mechanics. We shall later analyze them in detail from the quantum-mechanical standpoint; see Dirac's theory of light, to be presented in \hyperref[sec:III.6]{III.6}. An account in their historical setting is given in Bohr's papers on atomic structure from the years 1913--1916, published in a German translation by H. Stintzing, Braunschweig, 1921.}
\originalnote{13}{As closer mathematical analysis shows, the matrices must necessarily be infinite. We shall not discuss their properties here, since we shall examine them in full detail later. For the present, we need only understand that formal algebraic calculations with these matrices follow the familiar rules of matrix addition and multiplication. In particular, $0$ and $1$ denote the zero and identity matrices: respectively, the matrix all of whose entries vanish, and the matrix with zeros off the diagonal and ones on the diagonal.}
\originalnote{14}{If $Q_1,P_1$ are Hermitian, neither $Q_1P_1$ nor $P_1Q_1$ need be Hermitian, but $\tfrac12(Q_1P_1+P_1Q_1)$ always is. In the case of $Q_1^2P_1$, however, both $\tfrac12(Q_1^2P_1+P_1Q_1^2)$ and $Q_1P_1Q_1$ are possible. Admittedly, if $P_1Q_1-Q_1P_1=\frac{h}{2\pi i}1$, these two expressions happen to be equal. In the case of $Q_1^2P_1^2$, the possibilities include $\tfrac12(Q_1^2P_1^2+P_1^2Q_1^2)$, $Q_1P_1^2Q_1$, $P_1Q_1^2P_1$, and so forth; here the expressions do not all coincide even in the special case just mentioned. We shall not pursue this further here, since the operator calculus to be developed later will make these relations much more transparent.}
\originalnote{15}{We have
\[
\frac{h}{2\pi i}\frac{\partial}{\partial q_1}(q_1\psi)
=q_1\frac{h}{2\pi i}\frac{\partial\psi}{\partial q_1}
+\frac{h}{2\pi i}\psi.
\]
Thus
\[
\frac{h}{2\pi i}\frac{\partial}{\partial q_1}\cdot q_1
-q_1\cdot\frac{h}{2\pi i}\frac{\partial}{\partial q_1}
=\frac{h}{2\pi i}1,
\]
where $1$ is the identity operation, taking $\psi$ into itself. In other words, $\frac{h}{2\pi i}\frac{\partial}{\partial q_1}$ and $q_1$ satisfy the same commutation relation as the matrices $P_1,Q_1$.}
\originalnote{16}{See his first two papers in the book cited in \hyperlink{note-9}{Note~9}, \emph{Ann. Phys.} (4), vol.~79 (1926).}
\originalnote{17}{See the first of Schrödinger's papers cited in \hyperlink{note-16}{Note~16}. The spectrum and its parts will not be defined precisely until \hyperref[sec:II.6]{II.6}--\hyperref[sec:II.9]{9}.}
\originalnote{18}{In the original formulation of matrix mechanics (see our discussion above), there was no general concept of a state encompassing the
% Source PDF page 247 / original page 240
\hypertarget{origpage-240}{}%
stationary states as special cases. The theory dealt only with the stationary states associated with the energy eigenvalues.}
\originalnote{19}{$H=H(q_1,\ldots,q_k,p_1,\ldots,p_k)$ may even contain the time $t$ explicitly. In that case, of course, there will in general be no stationary states at all.}
\originalnote{20}{If there is only a point spectrum; see \hyperref[sec:II.6]{II.6}.}
\originalnote{21}{This series expansion, like all those that follow, converges ``in the mean.'' We discuss this more closely in \hyperref[sec:II.2]{II.2}.}
\originalnote{22}{The absence of such oscillations in stationary states, and only in these states, was one of Bohr's most important fundamental postulates in 1913. Classical electrodynamics contradicts it.}
\originalnote{23}{See the second of Schrödinger's papers cited in \hyperlink{note-16}{Note~16}.}
\originalnote{24}{See \hyperlink{note-7}{Note~7}.}
\originalnote{25}{See, for example, §§20 and 23 of Born and Jordan's book cited in \hyperlink{note-1}{Note~1}.}
\originalnote{26}{Since
\begin{alignat*}{2}
\mathsf S^{-1}1\mathsf S&=1,&\quad  \mathsf S^{-1}(aA)\mathsf S&=a\mathsf S^{-1}A\mathsf S,\\
\mathsf S^{-1}(A+B)\mathsf S&=\mathsf S^{-1}A\mathsf S+\mathsf S^{-1}B\mathsf S,\\
\mathsf S^{-1}AB\mathsf S&=\mathsf S^{-1}A\mathsf S\cdot \mathsf S^{-1}B\mathsf S,
\end{alignat*}
every matrix polynomial $P(A,B,\ldots)$ satisfies
\[
\mathsf S^{-1}P(A,B,\ldots)\mathsf S=P(\mathsf S^{-1}A\mathsf S,\mathsf S^{-1}B\mathsf S,\ldots).
\]
Choosing the left-hand sides of the commutation relations as $P$ proves that these relations are invariant; choosing $H$ as $P$ gives $\mathsf S^{-1}\overline{\mathsf H}\mathsf S=\mathsf H$.}
\originalnote{27}{$\delta_{\mu\nu}=1$ for $\mu=\nu$ and $0$ for $\mu\ne\nu$ is the familiar Kronecker--Weierstrass symbol.}
\originalnote{28}{The columns $s_{1\varrho},s_{2\varrho},\ldots$ of $\mathsf S$ for those $\varrho$ with $w_\varrho=\lambda$ form, as we have seen, a complete system of solutions; and, as columns of a matrix possessing an inverse, they must be linearly independent.}
\originalnote{29}{Any permutation of the columns of $\mathsf S$, together with the corresponding permutation of the rows of $\mathsf S^{-1}$, permutes the diagonal entries of $\mathsf H$ in the same way. Thus the order of $w_1,w_2,\ldots$ is indeed undefined and cannot be determined.}
\originalnote{30}{The theory of integral equations received its definitive form from Fredholm and Hilbert. An exhaustive account, with bibliographical references, is given in Courant--Hilbert, \emph{Methoden der mathematischen Physik}, Berlin, 1931.}
\originalnote{31}{More precisely, if we use Lebesgue's concept of the integral, $\varphi(q)=0$ must hold for $q\gtrless0$ except on a set of measure zero; that is, apart from such a set, $\varphi(q)\equiv0$.}
\originalnote{32}{The graph of $\delta(q)$ is thus conceived as an infinitely narrow and infinitely high spike at $q=0$, of area $1$. This is approximately the limiting behavior of $\sqrt{a/\pi}\,e^{-aq^2}$ as $a\to+\infty$, but is nevertheless impossible.}
\originalnote{33}{Incidentally, such a unification was sought long before quantum mechanics by E. H. Moore, the creator of so-called ``General Analysis.'' See the article by Hellinger--Toeplitz in the \emph{Math. Enzyklopädie}, vol.~II, 3, 9, Leipzig, 1927.}
\originalnote{34}{It has repeatedly been observed in Schrödinger's theory that, for the wave functions $\varphi$, only the finiteness of
\[
\underbrace{\int\cdots\int}_{\Omega}|\varphi(q_1,\ldots,q_k)|^2\,dq_1\cdots dq_k
\]
is essential. Thus $\varphi$ may, for example, be singular, and even become infinite, provided that this integral remains finite. An instructive example is the hydrogen
% Source PDF page 248 / original page 241
\hypertarget{origpage-241}{}%
atom in Dirac's relativistic theory; see \emph{Proc. Roy. Soc.}, London, vol.~117 (1928), and also W. Gordon, \emph{Z. Physik}, vol.~48 (1928).}
\originalnote{35}{A proof of this theorem will emerge in the course of our discussion of Hilbert space (see \hyperref[sec:II.2]{II.2} and \hyperref[sec:II.3]{3}, especially \hyperref[thm:II.5]{Theorem~5} in \hyperref[sec:II.2]{II.2}). It is worth adding that the easier half of the theorem, sufficient for many purposes, is the isomorphism between $F_\Omega$ and a suitable part of $F_Z$; this is due to Hilbert, \emph{Gött. Nachr.}, 1906. Schrödinger's original proof of equivalence (see \hyperlink{note-7}{Note~7}) rests only on this half of the theorem.}
\originalnote{36}{We have
\[
q_mq_n\varphi(q_1,\ldots,q_k)=q_nq_m\varphi(q_1,\ldots,q_k),
\]
\[
\frac{\partial}{\partial q_m}\frac{\partial}{\partial q_n}\varphi(q_1,\ldots,q_k)
=\frac{\partial}{\partial q_n}\frac{\partial}{\partial q_m}\varphi(q_1,\ldots,q_k),
\]
and
\begin{align*}
&\frac{\partial}{\partial q_m}\bigl(q_n\varphi(q_1,\ldots,q_k)\bigr)
-q_n\frac{\partial}{\partial q_m}\varphi(q_1,\ldots,q_k)\\
&\qquad=\begin{cases}0,&m\ne n,\\\varphi(q_1,\ldots,q_k),&m=n,\end{cases}
\end{align*}
from which the desired operator relations follow immediately.}
\originalnote{37}{The characterization of $\mathfrak R_n$ by \hyperref[ax:II.A]{A}, \hyperref[ax:II.B]{B}, and \hyperref[ax:II.Cn]{$C^{(n)}$} is due to Weyl (see, for example, \emph{Raum, Zeit, Materie}, Berlin, 1921). To obtain $\mathfrak R_\infty$ in place of $\mathfrak R_n$, it is natural to replace \hyperref[ax:II.Cn]{$C^{(n)}$} by \hyperref[ax:II.Cinfty]{$C^{(\infty)}$}, but \hyperref[ax:II.D]{D} and \hyperref[ax:II.E]{E} then become necessary; see the discussion below.}
\originalnote{38}{In addition to the zero point or zero vector $0$ of $\mathfrak R$, there is the number $0$, so that the same symbol is used for two things. The context, however, leaves no room for confusion.}
\originalnote{39}{It would suffice to require that $af$ belong to $\mathfrak M$ whenever $f$ does, and that $f+g$ belong to $\mathfrak M$ whenever $f,g$ do. For then, together with $f_1,\ldots,f_k$, the set $\mathfrak M$ contains $a_1f_1,\ldots,a_kf_k$, and successively also $a_1f_1+a_2f_2$, $a_1f_1+a_2f_2+a_3f_3$, \ldots, $a_1f_1+\cdots+a_kf_k$.}
\originalnote{40}{The reality of $(f,f)$ already follows from Hermitian symmetry: with $f=g$, this states that $(f,f)=\overline{(f,f)}$.}
\originalnote{41}{If $f$ has components $x_1,\ldots,x_n$, the discussion in $\gamma)$ of \hyperref[sec:II.1]{II.1}, specialized to finitely many components, gives
\[
\sqrt{(f,f)}=\sqrt{\sum_{\nu=1}^{n}|x_\nu|^2},
\]
which is the usual Euclidean length.}
\originalnote{42}{Since $(f,f)$ is real and $\ge0$, $\lVert f\rVert$ is real; the square root is chosen nonnegative. The same therefore holds for $\lVert f-g\rVert$.}
\originalnote{43}{By \hyperref[thm:II.2]{Theorem~2}, applied here to $f-g,g-h$, we must have $f-g=0$, that is, $g=f$; or $g-h=0$, that is, $g=h$; or $g-h=c(f-g)$ with $c$ real and $>0$, that is,
\[
g=\frac{c}{c+1}f+\frac1{c+1}h.
\]
Thus $g=af+(1-a)h$, with $a=1,0,c/(c+1)$ respectively. Geometrically, this means that the point $g$ lies on the segment joining $f,h$.}
\originalnote{44}{The following definition of a limit point is also customary: for every $\varepsilon>0$, there is an $f'$ in $\mathfrak A$ with $\lVert f-f'\rVert<\varepsilon$. The equivalence of the two definitions is proved exactly as in ordinary analysis.}
\originalnote{45}{For brevity we use the technical term from topology (see, for example, Hausdorff, \emph{Mengenlehre}, Berlin, 1927), which is explained below in the text.}
\originalnote{46}{Indeed, $\lVert f\rVert=\sqrt{(f,f)}=1$.}
\originalnote{47}{As one sees, the complete orthogonal systems correspond to Cartesian coordinate systems in $\mathfrak R_n$, that is, to the unit vectors in the directions of their axes.}
\originalnote{48}{As a linear manifold, such a set must contain $\{\mathfrak A\}$; and, being closed, it must also contain the limit points of $\{\mathfrak A\}$.}
\originalnote{49}{% Source PDF page 249 / original page 242
\hypertarget{origpage-242}{}%
If $\varphi_1,\varphi_2,\ldots$ is complete, then $\varphi_2,\varphi_3,\ldots$ is no longer complete, but is still infinite!}
\originalnote{50}{Recall that a double sequence $g_{nm}$ ($n,m=1,2,\ldots$) can also be written as a single sequence: $g_{11},g_{12},g_{21},g_{13},g_{22},g_{31},\ldots$.}
\originalnote{51}{This section is not needed to understand the sections that follow.}
\originalnote{52}{For example, Carathéodory, \emph{Vorlesungen über reelle Funktionen}, Leipzig, 1927, especially pp.~237--274; Kamke, \emph{Das Lebesguesche Integral}, Leipzig, 1925.}
\originalnote{53}{In general,
\begin{align*}
|x+y|^2&=(x+y)\overline{(x+y)}\\
&=x\bar x+y\bar y+(x\bar y+\bar xy)\\
&=|x|^2+|y|^2+2\operatorname{Re}(x\bar y).
\end{align*}}
\originalnote{54}{This is standard practice in the theory of the Lebesgue integral.}
\originalnote{55}{Although $\alpha$ varies continuously, this is a sum, not an integral: only countably many values of $\alpha$ occur in the sum!}
\originalnote{56}{Naturally, we define $(x(\alpha),y(\alpha))$ as $\sum_\alpha x(\alpha)\overline{y(\alpha)}$.}
\originalnote{57}{This is the set-theoretical theorem on the ``uncountability of the numerical continuum.'' See, for example, Hausdorff's book cited in \hyperlink{note-45}{Note~45}.}
\originalnote{58}{Let, for example, $\mathfrak R_\infty$ be an $F_\Omega$, where $\Omega$ is the space of all real $x$, $-\infty<x<\infty$. Then $d/dx$ is a function of functions, that is, an operator; but in our sense it is defined only for those $f(x)$ that are differentiable and for which
\[
\int_{-\infty}^{\infty}\left|\frac d{dx}f(x)\right|^2dx
\]
is finite (see \hyperref[sec:II.8]{II.8}, where this is discussed more precisely). Of course, $d^2f(x)/dx^2$ need not then exist in general, nor need
\[
\int_{-\infty}^{\infty}\left|\frac{d^2}{dx^2}f(x)\right|^2dx
\]
be finite. The function $f(x)=|x|^{3/2}e^{-x^2}$, for example, behaves in this way.}
\originalnote{59}{The discussion of condition \hyperref[ax:II.E]{E} in \hyperref[sec:II.3]{II.3} shows that it suffices to approximate arbitrarily closely every linear combination of functions that equal $1$ on a finite union of intervals and $0$ elsewhere. For this, it is enough to approximate each such function individually. In turn, we need only consider functions equal to $1$ on a single interval, since the others are sums of these. Let the interval be $a<x<b$. The function
\[
f(x)=\begin{cases}
0,&x<a-\varepsilon\text{ or }x>b+\varepsilon,\\
\cos^2\!\left(\dfrac\pi2\dfrac{a-x}{\varepsilon}\right),&a-\varepsilon\le x\le a,\\
\cos^2\!\left(\dfrac\pi2\dfrac{x-b}{\varepsilon}\right),&b\le x\le b+\varepsilon,\\
1,&a<x<b
\end{cases}
\]
does satisfy our regularity requirements and gives an arbitrarily close approximation when $\varepsilon>0$ is sufficiently small.}
\originalnote{60}{This argument is not rigorous, since it uses linearity for infinite sums, and so forth. It can, however, be completed as follows. Let $\varphi_1,\varphi_2,\ldots$ be a complete orthonormal system, and let $A,A^*$ be adjoint operators. Put $f=\sum_{\nu=1}^{\infty}x_\nu\varphi_\nu$ and $Af=\sum_{\nu=1}^{\infty}y_\nu\varphi_\nu$. Then
\begin{align*}
y_\mu&=(Af,\varphi_\mu)=(f,A^*\varphi_\mu)\\
&=\sum_{\nu=1}^{\infty}(f,\varphi_\nu)\overline{(A^*\varphi_\mu,\varphi_\nu)}
\qquad\text{[by \hyperref[thm:II.7]{Theorem 7, $\gamma)$}]}\\
&=\sum_{\nu=1}^{\infty}x_\nu\overline{(\varphi_\mu,A\varphi_\nu)}
=\sum_{\nu=1}^{\infty}(A\varphi_\nu,\varphi_\mu)x_\nu.
\end{align*}
% Source PDF page 250 / original page 243
\hypertarget{origpage-243}{}%
Thus, if we set $a_{\mu\nu}=(A\varphi_\nu,\varphi_\mu)$, we obtain the formula $y_\mu=\sum_{\nu=1}^{\infty}a_{\mu\nu}x_\nu$ of the text, with absolute convergence assured.

In the Hilbert space of sequences $x_1,x_2,\ldots$, the sequences $\varphi_1=1,0,0,\ldots$, $\varphi_2=0,1,0,\ldots$, and so forth form, as is easily seen, a complete orthonormal system. For $f=\{x_1,x_2,\ldots\}$ we have $f=\sum_{\nu=1}^{\infty}x_\nu\varphi_\nu$, and for $Af=\{y_1,y_2,\ldots\}$, $Af=\sum_{\nu=1}^{\infty}y_\nu\varphi_\nu$. This recovers precisely the representation used in the text.

Forming $a^*_{\mu\nu}$ for $A^*$, we find
\[
a^*_{\mu\nu}=(A^*\varphi_\nu,\varphi_\mu)
=(\varphi_\nu,A\varphi_\mu)
=\overline{(A\varphi_\mu,\varphi_\nu)}=\bar a_{\nu\mu}.
\]}
\originalnote{61}{$(Af,f)$ is in any event real, since it equals $(A^*f,f)=(f,Af)=\overline{(Af,f)}$.}
\originalnote{62}{Thus $U,U^*$ must be defined everywhere. Moreover, they are inverse to one another, so both assume every value, and each value only once.}
\originalnote{63}{For a prescribed value of $\int_{-\infty}^{+\infty}|\varphi(q)|^2dq$, either of the integrals
\[
\int_{-\infty}^{+\infty}q^2|\varphi(q)|^2dq
\qquad\text{or}\qquad
\int_{-\infty}^{+\infty}\left|\frac d{dq}\varphi(q)\right|^2dq
\]
can be made arbitrarily large. For example, let $\varphi(q)=ae^{-bq^2}$. All three integrals are finite for $b>0$, but are respectively proportional to $a^2b^{-1/2}$, $a^2b^{-3/2}$, and $a^2b^{1/2}$, so the values of two of them can be prescribed arbitrarily.}
\originalnote{64}{\emph{Gött. Nachr.}, 1906.}
\originalnote{65}{The Hermitian character of $R$ is required in the third step of the calculation
\begin{align*}
&\left(R\frac{f+g}{2},\frac{f+g}{2}\right)
-\left(R\frac{f-g}{2},\frac{f-g}{2}\right)\\
&\quad=\frac{(Rf,g)+(Rg,f)}2
=\frac{(Rf,g)+\overline{(f,Rg)}}2\\
&\quad=\frac{(Rf,g)+\overline{(Rf,g)}}2
=\operatorname{Re}(Rf,g).
\end{align*}}
\originalnote{66}{See the fifth paper in the book cited in \hyperlink{note-9}{Note~9}, \emph{Ann. Phys.} [4], vol.~80 (1926).}
\originalnote{67}{We deliberately do not discuss finer questions of convergence. They were not formulated precisely in the original versions of matrix and wave theory either, and we shall settle them later in any case; see, for example, \hyperref[sec:II.9]{II.9}.}
\originalnote{68}{The latter assertion is immediately clear only for an everywhere-defined continuous $\mathsf H$, that is, when $f_n\to f$ implies $\mathsf Hf_n\to\mathsf Hf$. However, as is easily seen, the weaker property that $f_n\to f$, $\mathsf Hf_n\to f^*$ imply $\mathsf Hf=f^*$ also suffices. This is the so-called ``closedness'' of $\mathsf H$; see the author's paper in \emph{Math. Ann.}, vol.~102 (1929). It is always satisfied by the operators of quantum mechanics, even the discontinuous ones. More precisely, a Hermitian operator that is not closed can be made closed, while remaining Hermitian, by extending its domain in a uniquely determined way. The analogous statement does not hold for continuity, for example. See \hyperref[sec:II.9]{II.9}, \hyperlink{origpage-76}{p.~76}.}
\originalnote{69}{See, for example, Schrödinger's treatment of the hydrogen atom in the work cited in \hyperlink{note-16}{Note~16}.}
\originalnote{70}{See the reference in \hyperlink{note-64}{Note~64}.}
\originalnote{71}{See, for example, Courant--Hilbert, cited in \hyperlink{note-30}{Note~30}.}
\originalnote{72}{% Source PDF page 251 / original page 244
\hypertarget{origpage-244}{}%
That is, $E(l)=(e_{\mu\nu}(l))$ and
\[
E(l;\xi,\eta)=\sum_{\mu,\nu=1}^{n}e_{\mu\nu}(l)\xi_\mu\bar\eta_\nu.
\]
Thus $e_{\mu\nu}(l)=\sum_{\lambda_\varrho\le l}x_{\varrho\mu}\bar x_{\varrho\nu}$.}
\originalnote{73}{For the concept of the Stieltjes integral, see Perron, \emph{Die Lehre von den Kettenbrüchen}, Leipzig, 1913; and, with particular attention to the requirements of operator theory, Carleman, \emph{Équations intégrales singulières}, Upsala, 1923. For readers less interested in these matters, its definition will suffice. For every partition $\Lambda_0,\Lambda_1,\ldots,\Lambda_k$ of the interval $a,b$, with
\[
a\le\Lambda_0<\Lambda_1<\cdots<\Lambda_k\le b,
\]
form the sum $\sum_{\tau=1}^{k}f(\Lambda_\tau)\{g(\Lambda_\tau)-g(\Lambda_{\tau-1})\}$. If this converges as the partition $\Lambda_0,\ldots,\Lambda_k$ is made progressively finer, the integral $\int_a^bf(x)\,dg(x)$ is declared to exist and to equal this limit. For $g(x)=x$, it becomes the familiar Riemann integral.

In our case, therefore, the equation derived states that $\int_{-\infty}^{+\infty}x\,dE(x;\xi,\eta)$ exists---we used $\lambda$ instead of $x$ for the variable---and equals $\sum_{\mu,\nu=1}^{n}h_{\mu\nu}\xi_\mu\bar\eta_\nu$.}
\originalnote{74}{$\min(a,b,\ldots,e)$ is the smallest, and $\max(a,b,\ldots,e)$ the largest, of the finitely many real numbers $a,b,\ldots,e$.}
\originalnote{75}{See the reference in \hyperlink{note-64}{Note~64} and Carleman's book cited in \hyperlink{note-73}{Note~73}. We shall return in detail to these ``continuous spectra''; see \hyperref[sec:II.8]{II.8}.}
\originalnote{76}{By $A(\lambda)\to B$, where $A(\lambda),B$ are operators in $\mathfrak R_\infty$ and $\lambda$ is a parameter, we mean that $A(\lambda)f\to Bf$ for every $f$ in $\mathfrak R_\infty$. This is thus an abbreviation for convergence statements in Hilbert space.}
\originalnote{77}{This follows from the definition of the Stieltjes integral given in \hyperlink{note-73}{Note~73}. For a proof, see the accounts cited there.}
\originalnote{78}{See \emph{Math. Ann.}, vol.~102 (1929).}
\originalnote{79}{This was proved only for sequences of $\lambda$'s. The limit must, however, be the same for all such sequences ($\lambda\to\lambda_0$ with $\lambda<\lambda_0$, or with $\lambda>\lambda_0$). For two such sequences can be interlaced to give another of the same kind; since the latter has a limit, both constituent sequences must have that same limit. It follows that convergence to the common limit of all the sequences also holds when $\lambda$ varies continuously.}
\originalnote{80}{This is the exact counterpart of the definition of $E(\lambda;\xi,\eta)$ given in \hyperref[sec:II.7]{II.7}.}
\originalnote{81}{As the construction in the proof of \hyperref[thm:II.10]{Theorem~10} shows, $P_{[\varphi]}f=(f,\varphi)\varphi$ when $\lVert\varphi\rVert=1$; hence $\lVert P_{[\varphi]}f\rVert=|(f,\varphi)|=|(\varphi,f)|$.}
\originalnote{82}{By \hyperref[thm:II.7]{Theorem~7},
\[
f=\sum_\varrho(f,\varphi_\varrho)\varphi_\varrho
=\sum_\varrho P_{[\varphi_\varrho]}f.
\]
This also follows from the concluding discussion in \hyperref[sec:II.4]{II.4}.}
\originalnote{83}{By \hyperlink{note-73}{Note~73},
\[
\int_{-\infty}^{+\infty}\lambda^2\,d\left(\sum_{\lambda_\varrho\le\lambda}|x_\varrho|^2\right)
=\lim\sum_{\tau=1}^{k}\Lambda_\tau^2
\sum_{\Lambda_{\tau-1}<\lambda_\varrho\le\Lambda_\tau}|x_\varrho|^2.
\]
% Source PDF page 252 / original page 245
\hypertarget{origpage-245}{}%
If $\Lambda_\tau^2-\Lambda_{\tau-1}^2<\varepsilon$ throughout, that is, if the mesh $\Lambda_0,\ldots,\Lambda_k$ is sufficiently fine, this changes by less than $\varepsilon\sum_{\varrho=1}^{\infty}|x_\varrho|^2=\varepsilon\lVert f\rVert^2$ when replaced by
\[
\sum_{\tau=1}^{k}\sum_{\Lambda_{\tau-1}<\lambda_\varrho\le\Lambda_\tau}
\lambda_\varrho^2|x_\varrho|^2
=\sum_{\Lambda_0<\lambda_\varrho\le\Lambda_k}\lambda_\varrho^2|x_\varrho|^2.
\]
For $\Lambda_0$ sufficiently small and $\Lambda_k$ sufficiently large, this is arbitrarily close to $\sum_{\tau=1}^{\infty}\lambda_\varrho^2|x_\varrho|^2$.\tnote{In the preceding sum, the source uses the summation index $\tau$ but retains $\varrho$ in the summand; this inconsistency is reproduced.} This sum is therefore the required limit, that is, the value of the integral.

The next integral formula is proved in exactly the same way.}
\originalnote{84}{At this point the mathematically correct method we follow diverges from Dirac's symbolic procedure (see, for example, his book cited in \hyperlink{note-1}{Note~1}). Dirac's procedure nevertheless treats the functions $f$ satisfying $(q-\lambda)f(q)=0$ as solutions; for simplicity, we put $l=j=1$, $q_j=q$. Yet every inner product $(f,g)=\int f(q)\overline{g(q)}\,dq$ is zero, whereas $f$ is to be nonzero. The function $f(q)$ is therefore imagined to be infinite at $q=\lambda$, the only point where it is nonzero, in such a way that the inner products $(f,g)$ become nonzero. Since $f(q)=0$ for $q\ne\lambda$, the integral $\int f(q)\overline{g(q)}\,dq$ can depend only on $g(\lambda)$; and its additivity clearly requires it to be proportional to $\overline{g(\lambda)}$. Thus it is $c\overline{g(\lambda)}$, where $c\ne0$. Replacing $f(q)$ by $f(q)/c$ makes $c=1$. We thus have a fictitious function $f(q)$ satisfying $\int f(q)\overline{g(q)}\,dq=\overline{g(\lambda)}$.

It is enough, of course, to consider $\lambda=0$. Then write $f(q)=\delta(q)$; it is defined by
\[
\text{\hypertarget{note84-Delta}{\(\Delta\).}}\qquad q\delta(q)\equiv0,\qquad
\int\delta(q)f(q)\,dq=f(0).
\]
For arbitrary $\lambda$, the solution is $\delta(q-\lambda)$. Although no function $\delta$ with property \hyperlink{note84-Delta}{$\Delta$} exists, there are sequences of functions tending toward such behavior; naturally, the limiting function does not exist. Examples are\tnote{Here, and in the approximate functions in \hyperlink{note-97}{Notes 97} and \hyperlink{note-98}{98}, the source alternates between $q$ and $x$. Its notation is retained.}
\[
f_\varepsilon(q)=\begin{cases}\dfrac1{2\varepsilon},&|x|<\varepsilon,\\0,&|x|\ge\varepsilon,\end{cases}
\qquad\varepsilon\to+0,
\]
or $f_\varepsilon(q)=\sqrt{a/\pi}\,e^{-ax^2}$ as $a\to+\infty$. See also \hyperref[sec:I.3]{I.3}, especially \hyperlink{note-32}{Note~32}.}
\originalnote{85}{A rigorous development of this idea, which is to be regarded here only as a heuristic approach, is given in papers by Hellinger, \emph{J. f. Math.}, vol.~136 (1909), and Weyl, \emph{Math. Ann.}, vol.~68 (1910).}
\originalnote{86}{Naturally, only success in physical applications can justify this standpoint, or its adoption in quantum mechanics.}
\originalnote{87}{Plancherel, \emph{Circ. Math. di Pal.}, vol.~30 (1910); Titchmarsh, \emph{London Math. Soc. Proc.}, vol.~22 (1924).}
\originalnote{88}{That is, $E'(\lambda)$ belongs not to $A'=\frac{h}{2\pi i}\frac{\partial}{\partial q}$ itself, but to an operator whose domain includes that of $A'$ and which agrees with $A'$ on the latter domain. See the discussion in \hyperref[sec:II.9]{II.9}.}
\originalnote{89}{All Fourier coefficients vanish, and hence so does the function itself; see, for example, Courant--Hilbert, cited in \hyperlink{note-30}{Note~30}.}
\originalnote{90}{% Source PDF page 253 / original page 246
\hypertarget{origpage-246}{}%
Although $M^{-1}Mf(q)=f(q)$---it suffices to define $f(q)=0$ for $q<0$ and apply the earlier theorems---we do not always have $MM^{-1}F(p)=F(p)$. For in general $\lVert M^{-1}F\rVert<\lVert F\rVert$, and therefore $\lVert MM^{-1}F\rVert<\lVert F\rVert$. Thus $M^{-1}M=1$, $MM^{-1}\ne1$: $M^{-1}$ is not the actual inverse of $M$ at all. Nor can there be another inverse, for, if there were one, $M^{-1}M=1$ would force it to equal our $M^{-1}$. Consequently, for $E'(\lambda)=ME(\lambda)M^{-1}$, one may still infer, for example, $E'(\lambda)^2=E'(\lambda)$, since only $M^{-1}M$ enters this calculation; but $E'(\lambda)\to MM^{-1}\ne1$ as $\lambda\to+\infty$.}
\originalnote{91}{Strictly speaking, this must be proved nonsymbolically, using the rigorous equation $(Af,g)=\int\lambda\,d(E(\lambda)f,g)$. The calculation then runs as follows:
\begin{align*}
(AFf,g)&=\int\lambda\,d(E(\lambda)Ff,g)
=\int\lambda\,d(FE(\lambda)f,g)\\
&=\int\lambda\,d(E(\lambda)f,Fg)
=(Af,Fg)=(FAf,g).
\end{align*}
It follows that $AF=FA$.}
\originalnote{92}{This follows from the identity, valid generally for Stieltjes integrals,
\[
\int f(\lambda)\,d\left(\int^\lambda g(\lambda')\,dh(\lambda')\right)
=\int f(\lambda)g(\lambda)\,dh(\lambda).
\]
It is immediately plausible from the reciprocal relation between $d$ and $\int^\lambda$. A rigorous proof was given by the author in \emph{Annals of Mathematics}, vol.~32 (1931).}
\originalnote{93}{That is,
\[
(Bf,g)=\int_{-\infty}^{\infty}r(\lambda)\,d(E(\lambda)f,g),\qquad
(Cf,g)=\int_{-\infty}^{\infty}s(\lambda)\,d(E(\lambda)f,g).
\]}
\originalnote{94}{The precise foundation of this notion of a function of an operator was given by the author in \emph{Annals of Math.}, vol.~32 (1931). F. Riesz was the first to define more general functions of operators by passage to limits from polynomials.}
\originalnote{95}{Alongside Hilbert's theory of bounded operators, the following discussion chiefly draws on the theory of unbounded Hermitian operators. The latter theory was established by the author in the work cited in \hyperlink{note-78}{Note~78}. M. Stone independently obtained similar results, \emph{Proc. Nat. Ac.}, 1929 and 1930.}
\originalnote{96}{\emph{Math. Ann.}, vol.~69 (1911).}
\originalnote{97}{The functions $f(q)$ analytic on $-\infty<q<+\infty$, with $\int_{-\infty}^{\infty}|f(q)|^2dq$, $\int_{-\infty}^{\infty}|f'(q)|^2dq$, and so forth finite, are also everywhere dense in $\mathfrak R_\infty$. By \hyperref[ax:II.D]{D} in \hyperref[sec:II.3]{II.3}, the linear combinations of
\[
f_{a,b}(q)=\begin{cases}1,&a<q<b,\\0,&\text{otherwise}\end{cases}
\]
are everywhere dense, so it suffices to approximate these arbitrarily closely by functions of the foregoing kind. For example,
\[
f^{(\varepsilon)}_{a,b}(q)
=\frac12-\frac12\operatorname{Tg\,hyp}\frac{(x-a)(x-b)}{\varepsilon}
=\frac1{e^{2(x-a)(x-b)/\varepsilon}+1}
\]
is of the required kind and converges to $f_{a,b}(q)$ as $\varepsilon\to+0$.}
\originalnote{98}{% Source PDF page 254 / original page 247
\hypertarget{origpage-247}{}%
It again suffices to approximate the $f_{a,b}(q)$, $0\le a<b\le1$, by functions from $\mathfrak A^0$. Suitable examples are
\[
f^{(\varepsilon)}_{a,b}(q)=\frac12-\frac12\operatorname{Tg\,hyp}\left(
\frac1\varepsilon\frac{(x-a-\varepsilon)(x-b+\varepsilon)}{x(1-x)}\right),
\]
with $\varepsilon\to+0$.}
\originalnote{99}{Since we may think of $f(\lambda)$ as approximated by polynomials, it suffices to consider polynomials, and hence their building blocks, the powers $f(\lambda)=\lambda^s$, $s=0,1,2,\ldots$. As a unitary transformation makes no difference here, we may assume that $A$ is diagonal; its diagonal entries are its eigenvalues $\lambda_1,\ldots,\lambda_n$. Thus we need only show that $A^s$ is also diagonal and has diagonal entries $\lambda_1^s,\ldots,\lambda_n^s$, which is evident.}
\originalnote{100}{To verify that these properties characterize Hermitian and unitary matrices respectively, we again need only consider diagonal matrices. The conjugate transpose of the diagonal matrix $A$ with entries $\lambda_1,\ldots,\lambda_n$ is the diagonal matrix $A^*$ with entries $\bar\lambda_1,\ldots,\bar\lambda_n$. Hence $A=A^*$ means $\lambda_1=\bar\lambda_1,\ldots,\lambda_n=\bar\lambda_n$, that is, all $\lambda_1,\ldots,\lambda_n$ are real; and $AA^*=A^*A=1$ means $\lambda_1\bar\lambda_1=1,\ldots,\lambda_n\bar\lambda_n=1$, that is, $|\lambda_1|=\cdots=|\lambda_n|=1$.}
\originalnote{101}{For proofs of these facts see the author's paper cited in \hyperlink{note-78}{Note~78}; see also A. Wintner, \emph{Math. Z.}, vol.~30 (1929). The absolute convergence of all integrals $\int_0^1 f(\sigma)\,d(E(\sigma)f,g)$ with bounded $f(\sigma)$ can be seen as follows. It suffices to consider $\operatorname{Re}(E(\sigma)f,g)$, since replacing $f,g$ by $if,g$ turns this into $\operatorname{Im}(E(\sigma)f,g)$. Now
\[
\operatorname{Re}(E(\sigma)f,g)
=\left(E(\sigma)\frac{f+g}{2},\frac{f+g}{2}\right)
-\left(E(\sigma)\frac{f-g}{2},\frac{f-g}{2}\right),
\]
so we need only consider $(E(\sigma)f,f)$. In $\int_0^1 f(\sigma)\,d(E(\sigma)f,f)$ the integrand is bounded, while the function of $\sigma$ following the $d$ is monotone. The assertion is therefore clear.}
\originalnote{102}{We apply Stieltjes integrals to elements of $\mathfrak R_\infty$ in place of numbers. The meaning of the resulting relations is this: choose an arbitrary fixed $g\in\mathfrak R_\infty$ and replace each element of $\mathfrak R_\infty$ occurring in them by its inner product with $g$; the resulting relations are to hold for every $g$. Compared with the operator Stieltjes integrals in \hyperref[sec:II.7]{II.7}, this is a semisymbolic procedure: there two elements $f,g$ were chosen arbitrarily, and $(\ldots f,g)$ was formed instead of $(\ldots,g)$, the dots standing for an operator.}
\originalnote{103}{We tacitly assume that an operator of this kind exists for every resolution of the identity $F(\lambda)$. That is, whenever $\int_{-\infty}^{\infty}\lambda^2\,d\|F(\lambda)f\|^2$ is finite, there is an $f^*$ such that, for every $g$,
\[
(f^*,g)=\int_{-\infty}^{\infty}\lambda\,d(F(\lambda)f,g),
\]
and these $f$ are everywhere dense. The Hermitian property follows from \hyperref[cond:II.Sbar3]{$\overline{\mathrm S}_3$} by interchanging $f,g$ and taking the complex conjugate of the last equation. Both assertions are proved in the paper cited in \hyperlink{note-78}{Note~78}.}
\originalnote{104}{For the eigenvalue problem of $A$ to be always soluble, the unitarity of $U$ would have to follow, that is, $\mathfrak E=\mathfrak F=\mathfrak R_n$ or $\mathfrak R_\infty$. In $\mathfrak R_\infty$ this is not the case, as we inferred from the existence of nonmaximal $A$. In $\mathfrak R_n$ it must hold, and this can also be seen directly. All linear manifolds are closed; thus the set of $\varphi-U\varphi$, being closed and everywhere dense, is $\mathfrak R_n$. The set $\mathfrak E$ of the $\varphi$ has at least as many dimensions as its linear image, the set of $\varphi-U\varphi$, and therefore has the maximum possible dimension $n$. As a one-to-one linear image of $\mathfrak E$, $\mathfrak F$ has the same
% Source PDF page 255 / original page 248
\hypertarget{origpage-248}{}%
dimension. Hence $\mathfrak E=\mathfrak F=\mathfrak R_n$.}
\originalnote{105}{In the paper cited in \hyperlink{note-78}{Note~78}, the author showed that the following operator is maximal but not hypermaximal. Let $\mathfrak R_\infty$ consist of all functions $f(q)$ defined for $0\le q<\infty$ for which $\int_0^\infty|f(q)|^2\,dq$ is finite. Let $R$ be the operator $i\frac{d}{dq}$, obtained, for example, by first defining it on all continuously differentiable $f$ for which $\int_{-\infty}^{\infty}|f'(q)|^2\,dq$ is finite and $f(0)=0$, and then taking its closure. Thus $R=-\frac{2\pi}{h}A'$, if the $A'$ of \hyperref[sec:II.8]{II.8} is taken for the interval $0\le q<\infty$; consequently $R$ is Hermitian. That $R$ is maximal but not hypermaximal can be verified by calculating $\mathfrak E$ and $\mathfrak F$.

This example is noteworthy because $A'=-\frac{h}{2\pi}R$ is usually interpreted as the momentum operator in a half-space bounded on one side by a wall at $q=0$.}
\originalnote{106}{This concept originates with Erhard Schmidt; see the paper cited in \hyperlink{note-78}{Note~78}.}
\originalnote{107}{According to \hyperref[sec:II.5]{II.5}, $R\cdot1$ and $1\cdot R$ are defined precisely where $R$ is defined; the same holds, for $a\ne0$, for $R\cdot a1$ and $a1\cdot R$. They are equal there, so $R$ commutes with $a1$. The only exception is $a=0$ when $R$ is not everywhere defined. This is rather inelegant and suggests modifying the definition of commutativity.}
\originalnote{108}{The resolution of the identity of $a1$ is
\[
F(\mu)=\begin{cases}1,&\mu\ge a,\\0,&\mu<a,\end{cases}
\]
as is easily verified.}
\originalnote{109}{For Hermitian $A,B$ belonging to the special class of ``completely continuous'' operators (cf.\ \hyperlink{note-70}{Note~70}), Toeplitz proved a theorem implying the assertion in the text: there exists a complete orthonormal system of common eigenfunctions. For the reference see \hyperlink{note-33}{Note~33}. The general theorem for arbitrary $A,B$, or $A,B,C,\ldots$, was proved by the author in the paper cited in \hyperlink{note-94}{Note~94}.}
\originalnote{110}{The $\varkappa_1,\varkappa_2,\ldots$ must be bounded, for example $\varkappa_m=1/m$, for $R$ to be continuous. For $\|Rf\|\le C\|f\|$ implies $\|R\psi_m\|=\|\varkappa_m\psi_m\|=|\varkappa_m|\le C\|\psi_m\|=C$, hence $|\varkappa_m|\le C$. Conversely, if $|\varkappa_m|\le C$, then
\begin{align*}
\|Rf\|^2&=\left\|R\left(\sum_{m=1}^{\infty}x_m\psi_m\right)\right\|^2
=\left\|\sum_{m=1}^{\infty}x_m\varkappa_m\psi_m\right\|^2
=\sum_{m=1}^{\infty}|x_m|^2|\varkappa_m|^2,\\
\|f\|^2&=\left\|\sum_{m=1}^{\infty}x_m\psi_m\right\|^2
=\sum_{m=1}^{\infty}|x_m|^2.
\end{align*}
Consequently $\|Rf\|^2\le C^2\|f\|^2$, and hence $\|Rf\|\le C\|f\|$, which gives continuity.}
\originalnote{111}{Exceptions on a set of $q_1,q_2$ of Lebesgue measure zero are permissible.}
\originalnote{112}{See, for example, \hyperlink{note-45}{Note~45}.}
\originalnote{113}{The matrix $\{a_{\mu\nu}\}$ represents the transformation (that is, the operator) $\eta_\mu=\sum_{\nu=1}^n a_{\mu\nu}\xi_\nu$ ($\mu=1,\ldots,n$; cf.\ \hyperref[sec:II.7]{II.7}). If we transform
\[
\xi_\mu=\sum_{\nu=1}^n\varkappa_{\nu\mu}\mathfrak x_\nu,
\qquad
\eta_\mu=\sum_{\nu=1}^n\varkappa_{\nu\mu}\mathfrak y_\nu,
\]
% Source PDF page 256 / original page 249
\hypertarget{origpage-249}{}%
then $\mathfrak y_\mu=\sum_{\nu=1}^n\mathfrak a_{\mu\nu}\mathfrak x_\nu$. The transformed matrix has entries
\[
\mathfrak a_{\mu\nu}=\sum_{\varrho,\sigma=1}^n
a_{\varrho\sigma}\overline{\varkappa_{\mu\varrho}}\varkappa_{\nu\sigma},
\]
and therefore
\begin{align*}
\sum_{\mu=1}^n\mathfrak a_{\mu\mu}
&=\sum_{\mu,\varrho,\sigma=1}^n a_{\varrho\sigma}
\overline{\varkappa_{\mu\varrho}}\varkappa_{\mu\sigma}\\
&=\sum_{\varrho,\sigma=1}^n a_{\varrho\sigma}
\left(\sum_{\mu=1}^n\overline{\varkappa_{\mu\varrho}}\varkappa_{\mu\sigma}\right)
=\sum_{\varrho=1}^n a_{\varrho\varrho},
\end{align*}
so the trace is invariant.}
\originalnote{114}{More precisely, the assertion is that, for hypermaximal definite $A$, there is exactly one hypermaximal definite $A'$ with $A'^2=A$. We prove existence.

Let $A=\int_{-\infty}^{\infty}\lambda\,dE(\lambda)$. Positivity implies that $E(\lambda)$ is constant for $\lambda<0$, and hence, by \hyperref[cond:II.Sbar1]{$\overline{\mathrm S}_1$}, is zero there. Otherwise there would be $\lambda_1<\lambda_2<0$ with $E(\lambda_2)-E(\lambda_1)\ne0$. Choose $f\ne0$ such that $(E(\lambda_2)-E(\lambda_1))f=f$. Then $E(\lambda)f=f$ for $\lambda\ge\lambda_2$, and $E(\lambda)f=0$ for $\lambda\le\lambda_1$. Consequently
\begin{align*}
(Af,f)&=\int_{-\infty}^{\infty}\lambda\,d(E(\lambda)f,f)
=\int_{\lambda_1}^{\lambda_2}\lambda\,d(E(\lambda)f,f)\\
&\le\int_{\lambda_1}^{\lambda_2}\lambda_2\,d(E(\lambda)f,f)
=\lambda_2((E(\lambda_2)-E(\lambda_1))f,f)
=\lambda_2(f,f)<0.
\end{align*}
Thus
\[
A=\int_{-\infty}^{\infty}\lambda\,dE(\lambda)
=\int_0^{\infty}\lambda\,dE(\lambda)
=\int_0^{\infty}\mu^2\,dE(\mu^2),
\]
and $A'=\int_1^{\infty}\mu\,dE(\mu^2)$ has the required property.\tnote{The lower limit in this formula for $A'$ is printed as $1$ in the source; the preceding calculation requires $0$. The printed limit is retained.}

Notice also that definiteness implies $E(\lambda)=0$ for $\lambda<0$, while the converse is evident. Definiteness is therefore characterized by the whole spectrum being $\ge0$.}
\originalnote{115}{See \hyperlink{note-64}{Note~64}. A direct proof is as follows. Let $\lambda_0<\lambda_1<\cdots<\lambda_n$ all be $\ge\varepsilon$, or all $\le-\varepsilon$, and let successive $E(\lambda)$ be unequal. Choose $\varphi_\nu\ne0$ such that $(E(\lambda_\nu)-E(\lambda_{\nu-1}))\varphi_\nu=\varphi_\nu$. Thus $E(\lambda)\varphi_\nu=\varphi_\nu$ for $\lambda\ge\lambda_\nu$, and $E(\lambda)\varphi_\nu=0$ for $\lambda\le\lambda_{\nu-1}$; normalize $\|\varphi_\nu\|=1$. These elements are orthogonal and can therefore be extended to a complete orthonormal system $\varphi_1,\ldots,\varphi_n,\varphi_{n+1},\ldots$. For $\nu=1,\ldots,n$,
\begin{align*}
\|A\varphi_\nu\|^2
&=\int_{-\infty}^{\infty}\lambda^2\,d\|E(\lambda)\varphi_\nu\|^2
=\int_{\lambda_{\nu-1}}^{\lambda_\nu}\lambda^2\,d\|E(\lambda)\varphi_\nu\|^2\\
&\ge\int_{\lambda_{\nu-1}}^{\lambda_\nu}\varepsilon^2\,d\|E(\lambda)\varphi_\nu\|^2
=\varepsilon^2\bigl(\|E(\lambda_\nu)\varphi_\nu\|^2-\|E(\lambda_{\nu-1})\varphi_\nu\|^2\bigr)
=\varepsilon^2\|\varphi_\nu\|^2=\varepsilon^2.
\end{align*}
Hence $\sum_{\mu=1}^{\infty}\|A\varphi_\mu\|^2\ge\sum_{\mu=1}^n\|A\varphi_\mu\|^2\ge n\varepsilon^2$, whereas the first sum is $\Sigma(A)=C^2$. Therefore $n\le C^2/\varepsilon^2$. For $|\lambda|\ge\varepsilon$, $E(\lambda)$ can consequently assume at most $2C^2/\varepsilon^2$ different values; it has only finitely many jumps, separated by intervals of constancy.
% Source PDF page 257 / original page 250
\hypertarget{origpage-250}{}%
Thus the spectrum is purely a point spectrum for $|\lambda|\ge\varepsilon$. Since this holds for every $\varepsilon>0$, the entire spectrum is purely a point spectrum.}
\originalnote{116}{Besides the original papers cited in the discussion, we refer chiefly to the encyclopaedia article by Hellinger and Toeplitz mentioned in \hyperlink{note-33}{Note~33}.}
\originalnote{117}{By analogy with geometry, the ball in $\mathfrak R_\infty$ with centre $\varphi_0$ and radius $r$ is the set of all $f$ with $\|f-\varphi_0\|\le r$. Its interior is given by $\|f-\varphi_0\|<r$, and its surface by $\|f-\varphi_0\|=r$. For the unit ball, $\varphi_0=0$ and $r=1$.}
\originalnote{118}{The first statistical assertions concerning the behavior of a system in the state $\varphi$ were made by M. Born; they were developed more fully by Dirac and Jordan. See \hyperlink{note-8}{Note~8} and \hyperlink{note-2}{Note~2}.}
\originalnote{119}{The association between physical quantities and Hermitian operators will be discussed more fully in \hyperref[sec:IV.1]{IV.1}. For the present we use the association in \hyperref[sec:I.2]{I.2}: the coordinates $q_1,\ldots,q_k$ correspond to multiplication by $q_1,\ldots,q_k$, the momenta to $\frac{h}{2\pi i}\frac{\partial}{\partial q_1},\ldots,\frac{h}{2\pi i}\frac{\partial}{\partial q_k}$, and the energy to $\mathsf H$.}
\originalnote{120}{An example from the kinetic theory of gases will illustrate this.

One mole (16 grams) of oxygen contains $6\cdot10^{23}$ oxygen molecules $\mathrm O_2$.\tnote{The source gives 16 grams for one mole of molecular oxygen. This figure is retained.} Each consists of two oxygen atoms which, if their internal structure is disregarded, are mass points with three degrees of freedom each. Thus $k=2\cdot3\cdot6\cdot10^{23}=36\cdot10^{23}$ degrees of freedom must be taken into account. A complete causal description requires $2k$ numerical values. The theory of gases, however, uses only two quantities, pressure and temperature, which are complicated functions of those $2k$ values. It can therefore make only statistical, probabilistic assertions. These assertions often come close to causal predictions because the probabilities are very close to $0$ or $1$, but their fundamentally statistical character remains unchanged.}
\originalnote{121}{If $F_t(\lambda)$ depends on $t$, $\frac{\partial}{\partial t}F_t(\lambda)=G_t(\lambda)$, and $\mathsf H$ is Hermitian, then $\frac{\partial}{\partial t}F_t(\mathsf H)=G_t(\mathsf H)$, since differentiation consists of subtraction, division, and passage to a limit. For $F_t(\lambda)=e^{-\frac{2\pi i}{h}(t-t_0)\lambda}$ this gives
\[
\frac{\partial}{\partial t}e^{-\frac{2\pi i}{h}(t-t_0)\mathsf H}
=-\frac{2\pi i}{h}\mathsf H e^{-\frac{2\pi i}{h}(t-t_0)\mathsf H}.
\]
Applying this to $\varphi$ gives the required differential equation. Moreover, $|F_t(\lambda)|=1$, so $F_t(\lambda)\overline{F_t(\lambda)}=1$ and hence $F_t(\mathsf H)\{F_t(\mathsf H)\}^*=1$: the operator is unitary. Finally, at $t=t_0$ it is the identity, so $\varphi_{t_0}=\varphi$.}
\originalnote{122}{\emph{Z. Physik}, vol.~37 (1926). The whole subsequent development (\hyperlink{note-2}{Note~2}) rests on this interpretation.}
\originalnote{123}{\emph{Physic. Rev.}, vol.~26 (1925). See also W. Bothe's survey in \emph{Handbuch der Physik}, vol.~23 (Quanten), Berlin 1926, Chap.~3, especially \S~73.}
\originalnote{124}{Bohr, Kramers, and Slater based their statistical theory of elementary phenomena on this idea. See
% Source PDF page 258 / original page 251
\hypertarget{origpage-251}{}%
\emph{Z. Physik}, vol.~24 (1924), and \hyperlink{note-123}{Note~123}. The Compton--Simon experiment may be regarded as its refutation.}
\originalnote{125}{Jordan recognized the relation between these jumps and the ``quantum jumps'' of the older Bohr theory: \emph{Z. Physik}, vol.~40 (1924).}
\originalnote{126}{All this presupposes that the structure of the observed system and of the measuring apparatus, all force fields, and so forth are known exactly, and that we seek only the state, that is, the instantaneous values of the coordinates. If this idealizing assumption is not satisfied, further sources of uncertainty arise.

Our description of an imprecise measurement is likewise idealized: we decide with absolute certainty whether the value lies in an interval $I=\{\lambda',\lambda''\}$, $\lambda'<\lambda''$, or not. In reality the boundaries are blurred and the decision can be made only with a certain probability. Nevertheless, the description chosen is mathematically the most useful, at least provisionally.}
\originalnote{127}{The last assertion can be checked using \hyperref[post:III.1.W]{W}. Form the resolutions of the identity of $R,S$ as in \hyperref[sec:II.8]{II.8}.}
\originalnote{128}{That the operator $R+S$ corresponds to the sum has been established for simultaneously measurable quantities. It also holds for arbitrary quantities, but in that case it is an unprovable fundamental postulate. See the end of \hyperref[sec:IV.1]{IV.1} and \hyperref[sec:IV.2]{IV.2}.}
\originalnote{129}{We leave it to the reader to work out the details of simultaneous measurability for states $\varphi$ in $\mathfrak M$, for quantities $\mathfrak R,\mathfrak S$ that cannot be measured with absolute precision (continuous spectra), and so forth, following the analogous discussion in \hyperref[sec:III.3]{III.3}.}
\originalnote{130}{Indeed,
\begin{align*}
\varepsilon^2&=((R-\varrho1)^2\varphi,\varphi)
=(R^2\varphi,\varphi)-2\varrho(R\varphi,\varphi)+\varrho^2\\
&=\|R\varphi\|^2-2(R\varphi,\varphi)^2+(R\varphi,\varphi)^2
=\|R\varphi\|^2-(R\varphi,\varphi)^2,
\end{align*}
and similarly for $\eta^2$.}
\originalnote{131}{\emph{Z. Physik}, vol.~43 (1927). Bohr deepened the analysis in \emph{Naturwiss.}, vol.~16 (1928). The mathematical treatment below was initiated by Kennard, \emph{Z. Physik}, vol.~44 (1926), and given its present form by Robertson.}
\originalnote{132}{Bohr (\hyperlink{note-131}{Note~131}) emphasized the fundamental significance of this point. The following description is not entirely classical in one respect: it uses the existence of light quanta, that is, the fact that light of frequency $\nu$ never occurs in amounts of energy smaller than $h\nu$.}
\originalnote{133}{The following discussion is due to Heisenberg and Bohr; see \hyperlink{note-131}{Note~131}.}
\originalnote{134}{See Einstein's original paper, \emph{Ann. Physik}, vol.~14 (1905), or any modern textbook.}
\originalnote{135}{For the microscope see \emph{Handbuch der Physik}, Berlin 1927, vol.~18, Chap.~2, G. Very accurate measurements require very small $\varepsilon$ and $\lambda$, hence $\gamma$-rays or radiation of still shorter wavelength. Ordinary lenses then fail. We could use only lenses whose molecules are neither shattered nor knocked out by the $\gamma$-rays. The existence of such molecules or particles violates no known law of nature, so their use in a thought experiment is legitimate.}
\originalnote{136}{% Source PDF page 259 / original page 252
\hypertarget{origpage-252}{}%
That $\frac{mc\,\Delta\nu}{\nu}$ is large compared with $\frac{h\,\Delta\nu}{c}$ means that $\nu$ is small compared with $\frac{mc^2}{h}$, that is, $\overline E=h\nu$ is small compared with $mc^2$. In other words, the energy of the light quantum $L$ is small compared with the relativistic rest energy of $T$\tnote{The source says ``relativistic rest mass.'' Since the comparison is with $mc^2$, the term ``rest energy'' is used here.}---an assumption that is unavoidable in a nonrelativistic calculation in any case.}
\originalnote{137}{For example, let $F(x)$ be a monochromatic wave train of frequency $\nu_0$ confined to the interval from $a$ to $\tau$:\tnote{The source prints the initial endpoint as $a$, but uses $0$ in the formula; it then switches from the amplitude $\alpha$ to $a$. These inconsistencies and the inconsistent factors of $\pi$ and $2\pi$ are retained.}
\[
F(x)=\begin{cases}\alpha\sin2\pi\nu_0x,&0\le x\le\tau,\\0,&\text{otherwise}.\end{cases}
\]
(Continuity at the junction requires $\sin2\pi\nu_0\tau=0$, that is, $\nu_0=\frac n{2\tau}$, $n=1,2,3,\ldots$.) By the familiar inversion formulae for the Fourier integral (see the reference in \hyperlink{note-87}{Note~87}), $a_\nu^2=b_\nu^2+c_\nu^2$, where
\begin{align*}
\left.\begin{matrix}b_\nu\\c_\nu\end{matrix}\right\}
&=2\int_{-\infty}^{+\infty}F(x)\begin{matrix}\cos\\\sin\end{matrix}2\pi\nu x\,dx
=2a\int_0^\tau\sin2\pi\nu_0x\cdot\begin{matrix}\cos\\\sin\end{matrix}2\pi\nu x\,dx\\
&=\pm a\int_0^\tau\left(\begin{matrix}\sin\\\cos\end{matrix}\pi(\nu+\nu_0)x
-\begin{matrix}\sin\\\cos\end{matrix}\pi(\nu-\nu_0)x\right)dx\\
&=-a\left[\frac{\begin{matrix}\cos\\\sin\end{matrix}\pi(\nu+\nu_0)x}{\pi(\nu+\nu_0)}
-\frac{\begin{matrix}\cos\\\sin\end{matrix}\pi(\nu-\nu_0)x}{\pi(\nu-\nu_0)}\right]_0^\tau\\
&=\begin{cases}
-a\left[\dfrac{(-1)^n\cos\pi\nu\tau-1}{\pi(\nu+\nu_0)}-\dfrac{(-1)^n\cos\pi\nu\tau-1}{\pi(\nu-\nu_0)}\right],\\[6pt]
-a\left[\dfrac{(-1)^n\sin\pi\nu\tau}{\pi(\nu+\nu_0)}-\dfrac{(-1)^n\sin\pi\nu\tau}{\pi(\nu-\nu_0)}\right]
\end{cases}\\
&=\begin{cases}
-\dfrac{2a\nu_0(1-(-1)^n\cos\pi\nu\tau)}{\pi(\nu^2-\nu_0^2)},\\[6pt]
\dfrac{2a\nu_0(-1)^n\sin\pi\nu\tau}{\pi(\nu^2-\nu_0^2)}.
\end{cases}
\end{align*}
Thus
\begin{align*}
a_\nu&=\frac{2a\nu_0\sqrt{2-2(-1)^n\cos\pi\nu\tau}}{\pi(\nu^2-\nu_0^2)}\\
&=\frac{4a\nu_0\left|\begin{matrix}\sin\\\cos\end{matrix}\frac12\pi\nu\tau\right|}{\pi(\nu^2-\nu_0^2)}
=\frac{4a\nu_0|\sin\pi(\nu-\nu_0)\tau|}{\pi(\nu^2-\nu_0^2)}.
\end{align*}
The frequencies near $\nu=\nu_0$ are therefore the most strongly represented. Most of the wave train's energy lies in the frequency ranges where $\pi(\nu-\nu_0)\tau$ has moderate values. The dispersion of $\nu-\nu_0$ (or, equivalently, of $\nu$) is therefore of the order of $1/\tau$. The exact calculation of the relevant expression gives the same result:
\[
\frac{\int_0^\infty a_\nu^2(\nu-\nu_0)^2\,d\nu}{\int_0^\infty a_\nu^2\,d\nu}.
\]}
\originalnote{138}{\emph{Proc. Roy. Soc., Lond.}, vol.~114 (1927). See also the exposition in Weyl, \emph{Gruppentheorie und Quantenmechanik}, 2nd ed., pp.~91 ff., Leipzig 1931.}
\originalnote{139}{The interested reader will find expositions of the electromagnetic theory of light in any textbook of electrodynamics, for example Abraham--Becker, \emph{Theorie der Elektrizität}, Berlin 1930. The same reference applies to the following discussion insofar as it belongs to Maxwell's theory.}
\originalnote{140}{See R. Courant and D. Hilbert, \emph{Methoden der mathematischen Physik I}, pp.~358--363, Berlin 1924.}
\originalnote{141}{For
\[
\iiint_{\boldsymbol H}[\mathfrak a,\operatorname{rot}\mathfrak b]\,dx\,dy\,dz
=\iiint_{\boldsymbol H}[\operatorname{rot}\mathfrak a,\mathfrak b]\,dx\,dy\,dz,
\]
% Source PDF page 260 / original page 253
\hypertarget{origpage-253}{}%
because
\[
[\mathfrak a,\operatorname{rot}\mathfrak b]-[\operatorname{rot}\mathfrak a,\mathfrak b]
=\operatorname{grad}(\mathfrak a\times\mathfrak b)
\]
($\mathfrak a\times\mathfrak b$ is the so-called external, or vector, product of $\mathfrak a,\mathfrak b$), provided that the normal component of $\mathfrak a\times\mathfrak b$ vanishes on the boundary of $\boldsymbol{H}$. Since $\mathfrak a\times\mathfrak b$ is perpendicular to both $\mathfrak a$ and $\mathfrak b$, this is certainly the case if either $\mathfrak a$ or $\mathfrak b$ is perpendicular to the boundary of $\boldsymbol{H}$. Here $\mathfrak a=\operatorname{rot}\overline{\mathfrak A}_m$ and $\mathfrak b=\overline{\mathfrak A}_n$, so the condition stated first is certainly satisfied.}
\originalnote{142}{See, for example, the reference in \hyperlink{note-138}{Note~138}.}
\originalnote{143}{Since $P_\nu^x$ commutes with $Q_\nu^y,Q_\nu^z$, but not with $Q_\nu^x$, to justify the following transformation we need to prove the relation (omitting superfluous indices and replacing $\mathfrak A$ by $F$)
\[
PF(Q)-F(Q)P=\frac h{2\pi i}F'(Q),
\]
when $P=\frac h{2\pi i}\frac\partial{\partial q}\cdots$ and $Q=q\cdots$. This relation, which is of the greatest importance especially in matrix theory, is easily verified by direct calculation.}
\originalnote{144}{The following argument shows most simply that the index tuples $k,M_1,M_2,\ldots$ described here can indeed be arranged in a sequence. Let $\pi_1,\pi_2,\pi_3,\ldots$ be the sequence of prime numbers $2,3,5,\ldots$. The products $\pi_1^{k-1}\pi_2^{M_1}\pi_3^{M_2}\cdots$ are in fact finite, for all but finitely many $M_n$ are zero, that is, the corresponding factors $\pi_{n+1}^{M_n}$ are $1$. As $k,M_1,M_2,\ldots$ run through all our index systems, $\pi_1^{k-1}\pi_2^{M_1}\pi_3^{M_2}\cdots$ runs through all the numbers $1,2,3,\ldots$, each exactly once. We can therefore use $\pi_1^{k-1}\pi_2^{M_1}\pi_3^{M_2}\cdots$ to obtain a single consecutive indexing of the $\alpha_{kM_1M_2\ldots}$.}
\originalnote{145}{The momenta $p_x,p_y,p_z$, for example, can be used as coordinates of the light quantum, together with a coordinate $\pi$ describing the state of polarization. The $p_x,p_y,p_z$ determine its direction, that is, the direction cosines $\alpha_x,\alpha_y,\alpha_z$ ($\alpha_x^2+\alpha_y^2+\alpha_z^2=1$), as well as its frequency $\nu$, wavelength $\lambda$, and energy. For according to Einstein the magnitude of the momentum vector is $h\nu/c$ (cf.\ \hyperlink{note-134}{Note~134}), so
\[
p_x=\frac{h\nu}{c}\alpha_x,\qquad p_y=\frac{h\nu}{c}\alpha_y,\qquad p_z=\frac{h\nu}{c}\alpha_z,
\]
that is,
\begin{gather*}
\nu=\frac ch\sqrt{p_x^2+p_y^2+p_z^2},\qquad \lambda=\frac c\nu,\qquad \text{energy}=h\nu,\\
\alpha_x=\frac{cp_x}{h\nu},\qquad\alpha_y=\frac{cp_y}{h\nu},\qquad\alpha_z=\frac{cp_z}{h\nu}.
\end{gather*}
A difficulty now arises because our normal oscillations $\overline{\mathfrak A}_n(xyz)\cdot\gamma\cos2\pi\varrho_n(t-\tau)$ are standing light waves, as they must be in the cavity $\boldsymbol{H}$ because of the reflecting walls. Thus $\overline{\mathfrak A}_n$ cannot be associated with any one ``ray direction'' $\alpha_x,\alpha_y,\alpha_z$. It is immediately clear that, besides $\alpha_x,\alpha_y,\alpha_z$, at least the opposite direction $-\alpha_x,-\alpha_y,-\alpha_z$ must also be present, and the corresponding statement holds for the momenta. In $\boldsymbol{H}$ we must therefore use coordinates other than $p_x,p_y,p_z,\pi$.

Some recent treatments of the subject have removed this difficulty by taking $\boldsymbol{H}$ to be the rectangular parallelepiped
\[
-A<x<A,\qquad-B<y<B,\qquad-C<z<C,
\]
whose boundary faces $x=\pm A$, $y=\pm B$, $z=\pm C$ are not, however, to be walls. Instead, $x=A$ is identified with $x=-A$, $y=B$ with $y=-B$, and $z=C$ with $z=-C$. Thus light reaching the wall $x=A$ at $A,y,z$
% Source PDF page 261 / original page 254
\hypertarget{origpage-254}{}%
continues its path at $-A,y,z$, back into $\boldsymbol{H}$, as though nothing had happened, and so forth (see, for example, a paper by L. Landau and R. Peierls, \emph{Z. Physik}, vol.~62 [1930]). Equivalently, space is assumed periodic in the $x,y,z$ directions, with respective periods $2A,2B,2C$.

The analytical treatment remains the same, except that the boundary condition now reads
\begin{gather*}
\mathfrak A(A,y,z)=\mathfrak A(-A,y,z),\qquad
\mathfrak A(x,B,z)=\mathfrak A(x,-B,z),\\
\mathfrak A(x,y,C)=\mathfrak A(x,y,-C)
\end{gather*}
(instead of $\frac\partial{\partial n}\mathfrak A=0$ on the boundary), and the ``elementary solutions'' in terms of which one expands are\tnote{The source prints $c$ in the following phase. With $c$ denoting the speed of light, as in the preceding formulas, the spatial term would contain $1/c$. The printed phase is retained.}
\[
\begin{matrix}\cos\\\sin\end{matrix}[2\pi\nu(t-c(\alpha_xx+\alpha_yy+\alpha_zz))]
\]
instead of $\overline{\mathfrak A}(xyz)\cdot\widetilde\varrho(t)$. The $\nu=\varrho_n$ and $\alpha_x=\alpha_{n,x}$, $\alpha_y=\alpha_{n,y}$, $\alpha_z=\alpha_{n,z}$ ($n=1,2,\ldots$) belonging to the eigensolutions are easily determined, and the further development of the theory coincides with that given in the text.}
\originalnote{146}{This passage to the limit $S\to+\infty$ differs from the passage $N\to+\infty$ carried out in the electromagnetic theory. Even if we interpret the $M_1,M_2,\ldots$ there as numbers of light quanta, $N$ was a bound on the number of incoherent light quanta (that is, quanta not agreeing in frequency, direction---which together constitute the momentum---and polarization; cf.\ \hyperlink{note-145}{Note~145}), whereas $S$ is a bound on the number of light quanta as such.}
\originalnote{147}{Considered independently of quantum mechanics, the consequences of replacing $\mathfrak R^{(S)}_\infty$ by $\overline{\mathfrak R}^{(S)}_\infty$ are equivalent to those of replacing ordinary statistics by so-called Bose--Einstein statistics. See, for example, Dirac, in the reference cited in \hyperlink{note-138}{Note~138}.}
\originalnote{148}{The literature of the time gives fuller discussions of how this ``dual nature'' was conceived and how paradoxical it seemed. See, for example, the works cited in \hyperlink{note-6}{Note~6}.

It has often been said that quantum mechanics attributes the same dual nature to matter, since the discrete particles (electrons, protons) are also described by wave functions and exhibit typical wave properties, such as diffraction by gratings, and so forth. [See the experiments of Davison--Germer, \emph{Physic. Rev.}, vol.~30 (1927); \emph{Proc. nat. Acad. Sci. U.S.A.}, vol.~14 (1928); also C. P. Thompson, \emph{Proc. Roy. Soc., Lond.}, vol.~117 (1928), and Rupp, \emph{Ann. Physik}, vol.~85 (1928).] It must be emphasized, however, that quantum mechanics derives both ``natures'' from a single, unified theory of elementary phenomena. The paradox of the earlier quantum theory consisted in the need to invoke, alternately, two mutually contradictory theories to explain experience: the electromagnetic theory of light of Maxwell--Hertz and Einstein's light-quantum theory.}
\originalnote{149}{Weisskopf and Wigner gave the exact solution of these differential equations [\emph{Z. Physik}, vol.~63 (1930)]; their solution confirms these statements.}
\originalnote{150}{As is well known, N. Bohr established in 1913 (see the reference in \hyperlink{note-5}{Note~5}) that, in passing from a stationary state of energy $W^{(1)}$ to a stationary state of energy $W^{(2)}$, an atom emits light of frequency $\frac{W^{(1)}-W^{(2)}}h$ (of course $W^{(1)}>W^{(2)}$). In our case this corresponds to $\frac{W_{\bar k}-W_k}h$.}
\originalnote{151}{For
\begin{align*}
|e^{ix}-1|^2&=(e^{ix}-1)\overline{(e^{ix}-1)}=(e^{ix}-1)(e^{-ix}-1)\\
&=2-e^{ix}-e^{-ix}=2-2\cos x=2(1-\cos x).
\end{align*}}
\originalnote{152}{% Source PDF page 262 / original page 255
\hypertarget{origpage-255}{}%
We have (cf.\ Courant--Hilbert, p.~49)
\[
\int_{-\infty}^{\infty}\frac{1-\cos x}{x^2}\,dx
=2\int_0^{\infty}\frac{1-\cos x}{x^2}\,dx
=\int_0^{\infty}\frac{1-\cos(2y)}{y^2}\,dy
=2\int_0^{\infty}\frac{\sin^2y}{y^2}\,dy=\pi.
\]}
\originalnote{153}{$P_\nu^x$ commutes with all $Q_\mu^x,Q_\mu^y,Q_\mu^z,P_\mu^x,P_\mu^y,P_\mu^z$ except $Q_\nu^x$. In fact,
\[
P_\nu^xQ_\nu^x-Q_\nu^xP_\nu^x=\frac h{2\pi i}1.
\]
Hence
\[
\mathsf H_0Q_\nu^x-Q_\nu^x\mathsf H_0
=\frac1{2m_\nu}(P_\nu^x)^2Q_\nu^x-Q_\nu^x\frac1{2m_\nu}(P_\nu^x)^2
=\frac h{2\pi i}\frac1{m_\nu}P_\nu^x;
\]
cf.\ \hyperlink{note-143}{Note~143}.}
\originalnote{154}{\emph{Physik. Z.}, vol.~18 (1917).}
\originalnote{155}{It is important to emphasize the conceptual distinction between a system as such and a system in a particular state. A system is, for example, a hydrogen atom, that is, an electron and a proton with the known forces acting between them. Formally it is described by specifying that configuration space has six dimensions, that the coordinates are called $q_1,\ldots,q_6$, the momenta $p_1,\ldots,p_6$, and the Hamiltonian function is\tnote{The source prints a plus sign before the Coulomb term. For the electron--proton attraction the sign would be negative; the printed formula is retained.}
\begin{align*}
H(q_1\cdots q_6p_1\cdots p_6)
&=\frac{p_1^2+p_2^2+p_3^2}{2m_e}+\frac{p_4^2+p_5^2+p_6^2}{2m_p}\\
&\quad+\frac{e^2}{\sqrt{(q_1-q_4)^2+(q_2-q_5)^2+(q_3-q_6)^2}}.
\end{align*}
A state is fixed only by further specifications: in classical mechanics, by assigning numerical values $q_1^0,\ldots,q_6^0,p_1^0,\ldots,p_6^0$ to the coordinates and momenta $q_1,\ldots,q_6,p_1,\ldots,p_6$; in quantum mechanics, by the wave function $\varphi(q_1\cdots q_6)$. No further specifications are ever needed: once system and state are given, the theory supplies unambiguous rules of calculation for answering every question.}
\originalnote{156}{Such ensembles, called collectives, are necessary in general if probability theory is to be founded as a theory of frequencies. They were introduced by R. von Mises, who recognized their significance for probability theory and developed the theory accordingly (see, for example, his book \emph{Wahrscheinlichkeit, Statistik und ihre Wahrheit}, Berlin 1928).}
\originalnote{157}{$w(a')$ is the probability of $a\le a'$, that is, the probability assigned to the interval $-\infty,a'$. As we see, $w(a)$, or $w_{\mathfrak R}(a)$ as we shall write to make its dependence on $\mathfrak R$ explicit, has the following properties: as $a\to-\infty$, $w_{\mathfrak R}(a)\to0$; as $a\to+\infty$, $w_{\mathfrak R}(a)\to1$; for $a\ge a_0$ and $a\to a_0$, $w_{\mathfrak R}(a)\to w_{\mathfrak R}(a_0)$; for $a'\le a''$, $w_{\mathfrak R}(a')\le w_{\mathfrak R}(a'')$. (In quantum mechanics, if $E(\lambda)$ is the resolution of the identity belonging to $R$, then $w_{\mathfrak R}(a)=\|E(a)\varphi\|^2=(E(a)\varphi,\varphi)$.)

If $w_{\mathfrak R}(a)$ is differentiable, the familiar ``probability density'' $\frac d{da}w_{\mathfrak R}(a)$ may be introduced in its place. If it is discontinuous at $a=a_0$ (from the left, of course), the point $a=a_0$ has the ``discrete probability'' $w_{\mathfrak R}(a_0)-w_{\mathfrak R}(a_0-0)$. But the general concept applicable in every case is $w_{\mathfrak R}(a)$; see the reference in \hyperlink{note-156}{Note~156}.}
\originalnote{158}{This follows from the so-called law of large numbers, Bernoulli's theorem.}
\originalnote{159}{% Source PDF page 263 / original page 256
\hypertarget{origpage-256}{}%
For example, it was regarded as a fundamental difficulty in defining the electric field that the electric ``test body'' to be used for this purpose cannot be smaller than an electron.}
\originalnote{160}{The sharp value $a_0$ corresponds to the probability function
\[
w_{\mathfrak R}(a)=\begin{cases}1,&a\ge a_0,\\0,&a<a_0.\end{cases}
\]
In this case, and only in this case, the dispersion $\varepsilon^2$ is zero. Incidentally, the mean $\varrho$ and the dispersion $\varepsilon^2$ are calculated as follows (Stieltjes integrals):
\begin{align*}
\varrho&=\int_{-\infty}^{+\infty}a\,dw_{\mathfrak R}(a),\\
\varepsilon^2&=\int_{-\infty}^{+\infty}(a-\varrho)^2\,dw_{\mathfrak R}(a)\\
&=\int_{-\infty}^{+\infty}a^2\,dw_{\mathfrak R}(a)-2\varrho\int_{-\infty}^{+\infty}a\,dw_{\mathfrak R}(a)+\varrho^2\\
&=\int_{-\infty}^{+\infty}a^2\,dw_{\mathfrak R}(a)-\varrho^2\\
&=\int_{-\infty}^{+\infty}a^2\,dw_{\mathfrak R}(a)-\left(\int_{-\infty}^{+\infty}a\,dw_{\mathfrak R}(a)\right)^2.
\end{align*}
(Compare \hyperref[sec:III.4]{III.4} and \hyperlink{note-130}{Note~130}.)}
\originalnote{161}{These are independent measurements on several systems; successive measurements on the same system would yield the same value again and again (cf.\ \hyperref[sec:III.3]{III.3}).}
\originalnote{162}{Consider, for example, what happens if we substitute for $\mathfrak R,\mathfrak S$ the quantities $q$ (a Cartesian coordinate) and $p$ (momentum), which cannot be measured simultaneously because of the uncertainty relations. If $q$ has very little dispersion in an ensemble, a measurement of $p$ with accuracy (that is, dispersion) $\varepsilon$ will increase the dispersion of $q$ to at least $h/(4\pi\varepsilon)$ (cf.\ \hyperref[sec:III.4]{III.4})---and thus spoil everything.}
\originalnote{163}{$a\mathfrak R$, $\mathfrak S^2$, $\mathfrak R+\mathfrak S+\cdots$ mean that, in the sense of the definition given above, the quantities $\mathfrak R$, or $\mathfrak S$, or $\mathfrak R,\mathfrak S,\ldots$, respectively, are substituted in the functions $f(x)=ax$, $f(x)=x^2$, or $f(x,y,\ldots)=x+y+\cdots$.}
\originalnote{164}{For example, the energy operator in Heisenberg's theory of an electron moving in a potential field $V(xyz)$,
\[
\mathsf H_0=\frac{(P^x)^2+(P^y)^2+(P^z)^2}{2m}+V(Q^x,Q^y,Q^z)
\]
(see, for example, \hyperref[sec:III.6]{III.6}), is the sum of two noncommuting operators:
\[
R=\frac{(P^x)^2+(P^y)^2+(P^z)^2}{2m},\qquad S=V(Q^x,Q^y,Q^z).
\]
Whereas measuring the quantity $\mathfrak R$ belonging to $R$ is a momentum measurement, and measuring the quantity $\mathfrak S$ belonging to $S$ is a coordinate measurement, the quantity $\mathfrak R+\mathfrak S$ belonging to $\mathsf H_0=R+S$ is measured quite differently: for example, by measuring the frequency of the spectral lines emitted by this bound electron. By Bohr's frequency relation these determine the energy levels, that is, the values of $\mathfrak R+\mathfrak S$. Nevertheless, under all circumstances,
\[
\operatorname{Erw}(\mathfrak R+\mathfrak S)=\operatorname{Erw}(\mathfrak R)+\operatorname{Erw}(\mathfrak S).
\]}
\originalnote{165}{For dispersion-free ensembles, however, there is no reason not to introduce properly normalized expectation values.}
\originalnote{166}{That is,
\[
U\varphi_m=\sum_nu_{mn}\varphi_n,
\]
% Source PDF page 264 / original page 257
\hypertarget{origpage-257}{}%
which of course requires the finiteness of $\sum_n|u_{mn}|^2$. This can be justified as follows, for example. If $\sum_n|x_n|^2=1$, then for $\varphi=\sum_nx_n\varphi_n$, $R=P_{[\varphi]}$ has the matrix $\bar x_\mu x_\nu$, so the corresponding physical quantity $\mathfrak R$ has expectation $\sum_{mn}u_{nm}\bar x_mx_n$. Since $P_{[\varphi]}=P_{[\varphi]}^2$ and $1-P_{[\varphi]}=(1-P_{[\varphi]})^2$, this expectation is $\ge0$ and $\le\operatorname{Erw}(1)$, hence at least for normalized $\operatorname{Erw}(\mathfrak R)$ it is $\ge0$ and $\le1$. If $x_{N+1}=x_{N+2}=\cdots=0$, this says that the $N$-dimensional Hermitian form $\sum_{m,n=1}^Nu_{nm}\bar x_mx_n$ takes values $\ge0$, $\le1$ when $\sum_{n=1}^N|x_n|^2=1$, that is, the eigenvalues of the matrix $u_{\mu\nu}$, $\mu,\nu=1,\ldots,N$, are $\ge0$, $\le1$. Thus the length of the vector $y_m=\sum_{n=1}^Nu_{mn}x_n$ is always no greater than that of the vector $x_m$. For
\[
x_m=\begin{cases}1,&m=\bar m,\\0,&\text{otherwise},\end{cases}
\]
we have $y_m=u_{m\bar m}$, and therefore
\[
\sum_{m=1}^N|x_m|^2\ge\sum_{m=1}^N|y_m|^2,\qquad
1\ge\sum_{m=1}^N|u_{m\bar m}|^2.
\]
Since this holds for every $N$, $\sum_n|u_{\bar mn}|^2\le1$.}
\originalnote{167}{Strictly speaking, the entire argument is sound only when $\varphi_1,\varphi_2,\ldots$ all belong to the domain of $R$. For each $R$ we can indeed find such a complete orthonormal system $\varphi_1,\varphi_2,\ldots$ (cf.\ \hyperref[sec:II.11]{II.11}), but if $R$ is not everywhere defined this system depends on $R$. What we actually have, then, is for each complete orthonormal system $\varphi_1,\varphi_2,\ldots$ a $U$ depending on that system, such that
\[
\operatorname{Erw}(\mathfrak R)=\operatorname{Spur}(UR)
\]
need hold only for those $R$ whose domains contain $\varphi_1,\varphi_2,\ldots$.

All these $U$ are nevertheless equal. For if $U',U''$ are two of them, the above formula holds for both when $R$ is everywhere defined, so $\operatorname{Spur}(U'R)=\operatorname{Spur}(U''R)$. For $R=P_{[\varphi]}$ this gives $(U'\varphi,\varphi)=(U''\varphi,\varphi)$, or $((U'-U'')\varphi,\varphi)=0$. Since this holds for all $\varphi$ with $\|\varphi\|=1$, and hence for all elements of Hilbert space, we must have $U'-U''=0$, that is, $U'=U''$.}
\originalnote{168}{We cannot simply put $\mathfrak S=\sqrt{\mathfrak R}$, that is, $\mathfrak S=h(\mathfrak R)$, $h(x)=\sqrt x$, because we consider only real-valued functions defined for all real $x$, and $\sqrt x$ is not of this kind: it is imaginary for negative $x$.}
\originalnote{169}{A simple direct proof is also possible. Let $\varphi_1,\varphi_2,\ldots$ be a complete orthonormal system, $a_{\mu\nu}=(A\varphi_\mu,\varphi_\nu)$, $b_{\mu\nu}=(B\varphi_\mu,\varphi_\nu)$, and $\operatorname{Spur}(AB)=\sum_{\mu\nu}a_{\mu\nu}b_{\nu\mu}$. This is $\ge0$ if all the sums $\sum_{\mu,\nu=1}^Na_{\mu\nu}b_{\nu\mu}$ are $\ge0$. If $f=\sum_{\mu=1}^Nx_\mu\varphi_\mu$, then
\[
(Af,f)=\sum_{\mu,\nu=1}^Na_{\mu\nu}x_\mu\bar x_\nu\ge0,\qquad
(Bf,f)=\sum_{\mu,\nu=1}^Nb_{\mu\nu}x_\mu\bar x_\nu\ge0,
\]
so the finite matrices $a_{\mu\nu},b_{\mu\nu}$ ($\mu,\nu=1,\ldots,N$) are likewise definite. Both definiteness and the value of $\sum_{\mu,\nu=1}^Na_{\mu\nu}b_{\nu\mu}$ are orthogonal invariants in $N$-dimensional space. Since $b_{\mu\nu}$ is Hermitian, it can be brought to diagonal form by an orthogonal transformation in that $N$-dimensional space. We may therefore assume at once that it is diagonal, that is, $b_{\mu\nu}=0$ for $\mu\ne\nu$. Thus
\[
\sum_{\mu,\nu=1}^Na_{\mu\nu}b_{\nu\mu}=\sum_{\mu=1}^Na_{\mu\mu}b_{\mu\mu}.
\]
But definiteness gives $a_{\mu\mu}\ge0$, $b_{\mu\mu}\ge0$ (put $x_\nu=1$ for $\nu=\mu$ and $x_\nu=0$ for $\nu\ne\mu$), so the sum above is indeed $\ge0$.}
\originalnote{170}{% Source PDF page 265 / original page 258
\hypertarget{origpage-258}{}%
For $\varphi'=\varphi''$ this is clear; suppose $\varphi'\ne\varphi''$. ``Orthogonalizing'' $\varphi',\varphi''$ (cf.\ \hyperref[sec:II.2]{II.2}) produces a $\varphi_1$, with $\|\varphi_1\|=1$, orthogonal to $\varphi'$, such that $\varphi''$ is a linear combination of $\varphi',\varphi_1$. Thus $\varphi''=a\varphi'+b\varphi_1$, $\|\varphi''\|^2=|a|^2+|b|^2=1$. Let, for example, $|a|=\cos\theta$, $|b|=\sin\theta$, so $a=e^{i\alpha}\cos\theta$, $b=e^{i\beta}\cos\theta$.\tnote{The source prints $\cos\theta$ in this expression for $b$; the preceding and following expressions require $\sin\theta$.} Then for $a(x)=e^{ix\alpha}\cos x\theta$, $b(x)=e^{ix\beta}\sin x\theta$ we also have $|a(x)|^2+|b(x)|^2=1$, and hence for $\varphi(x)=a(x)\varphi'+b(x)\varphi_1$ we have $\|\varphi(x)\|=1$. Moreover, $\varphi(x)$ varies continuously from $\varphi'$ ($x=0$) to $\varphi''$ ($x=1$).}
\originalnote{171}{Strictly speaking, because of \hyperref[remark:IV.2.2]{2} we should also require $V\ne0$, $W\ne0$. However, the cases $V=0$ or $W=0$ are included by $c'=0$, $c''=1$, or $c'=1$, $c''=0$, respectively.}
\originalnote{172}{The deduction in the last two sections, leading to uniform ensembles, was given by the author in \emph{Gött. Nachr.}, 1927. The existence of uniform ensembles and their relation to general ensembles were discovered independently by H. Weyl, \emph{Z. Physik}, vol.~46 (1927), and by the author in the paper just cited. A special case of more general ensembles, namely that of two coupled systems (cf.\ the discussion in \hyperref[sec:VI.2]{VI.2}), had appeared earlier in J. Landau, \emph{Z. Physik}, vol.~45 (1927).}
\originalnote{173}{If the ``hidden parameters'', which we denote collectively by $\pi$, take only discrete values $\pi_1,\pi_2,\ldots,\pi_n$ ($n>1$), the two ensembles whose mixture is the original ensemble are obtained by placing the systems with $\pi=\pi_1$ in one and those with $\pi\ne\pi_1$ in the other. If $\pi$ varies continuously over a region $\Pi$, let $\Pi'$ be a subregion of $\Pi$; one ensemble then contains the systems with $\pi$ in $\Pi'$, and the other those whose $\pi$ does not belong to $\Pi'$.}
\originalnote{174}{Phase space is six-dimensional: its six coordinates are the three Cartesian coordinates $q_1,q_2,q_3$ of the mass point and the three corresponding momenta $p_1,p_2,p_3$. According to \hyperref[sec:III.4]{III.4}, the respective dispersions $\varepsilon_1,\varepsilon_2,\varepsilon_3,\eta_1,\eta_2,\eta_3$ satisfy
\[
\varepsilon_1\eta_1\ge\frac h{4\pi},\qquad
\varepsilon_2\eta_2\ge\frac h{4\pi},\qquad
\varepsilon_3\eta_3\ge\frac h{4\pi},
\]
that is,
\[
\varepsilon_1\varepsilon_2\varepsilon_3\eta_1\eta_2\eta_3\ge\left(\frac h{4\pi}\right)^3.
\]
The position in the phase space of classical mechanics cannot, in any event, be specified more accurately than to within this volume.}
\originalnote{175}{See Schrödinger's exceptionally illuminating discussion of this question in \emph{Naturwiss.}, vol.~17 (1929), no.~37.}
\originalnote{176}{$u_\mu\to-u_\mu$ (or $u_\nu\to-u_\nu$, or $v_\mu\to-v_\mu$) is a symmetry operation of $K$ under which the first integrands change sign; their integrals are therefore zero. The operations $u_\mu\to v_\mu$, $v_\mu\to u_\mu$, or $u_\mu\to u_\nu$, $u_\nu\to u_\mu$, are symmetry operations of $K_1$ that transform the latter integrals into one another. These integrals are therefore equal, and hence each is $1/(2s)$ times their sum,
\[
\iint\cdots\iint_K(u_1^2+v_1^2+\cdots+u_s^2+v_s^2)\,dv
=\iint\cdots\iint_Kdv=\text{surface area of }K,
\]
which we shall call $C$.}
\originalnote{177}{If $\operatorname{Spur}U$ is infinite, the latter formula, obtained by subtraction, may appear questionable. It can, however, also be justified as follows. That $\mathfrak R$ has the value $\lambda^*$ means $w(a',a'')=0$, that is, $\operatorname{Spur}U(E(a'')-E(a'))=0$ for $a''<\lambda^*$ or $a'\ge\lambda^*$. Since this trace is always $\ge0$, is nondecreasing with $a''$, and nonincreasing with $a'$, it suffices for $a''<\lambda^*$ to consider its limit as $a'\to-\infty$, and for $a'\ge\lambda^*$ its limit as $a''\to+\infty$. Thus $\operatorname{Spur}UE(a'')=0$ for $a''<\lambda^*$ and $\operatorname{Spur}U(1-E(a'))=0$ for $a'\ge\lambda^*$.}
\originalnote{178}{% Source PDF page 266 / original page 259
\hypertarget{origpage-259}{}%
\textbf{That is, if the property $\mathfrak E$ holds, the state is $\varphi$. If it does not hold, ``non-$\mathfrak E$'' holds; in place of $E=P_{[\varphi]}$, $\mathfrak M=[\varphi]$, we then have $1-E=1-P_{[\varphi]}$, $\mathfrak R-\mathfrak M=\mathfrak R-[\varphi]$, which does not determine $U$ uniquely.} ($E$ corresponds precisely to the question ``Is the state $\varphi$?'') \textbf{A measurement that uniquely determines the state whatever its outcome is a measurement of a quantity $\mathfrak R$ whose operator $R$ has a pure point spectrum consisting entirely of simple eigenvalues;} cf.\ \hyperref[sec:III.3]{III.3}. After the measurement, the system is in one of the states $\varphi_1,\varphi_2,\ldots$ (the eigenfunctions of $R$); that is, the measurement generally changes the state of $\boldsymbol S$. Similarly, measurement of $\mathfrak E$ also changes the state: afterwards, a positive result gives $U=P_{[\varphi]}$, while a negative result gives $U(1-P_{[\varphi]})=U$, $UP_{[\varphi]}=0$, that is, $U\varphi=0$. Beforehand, neither condition need have held. This quantum-mechanical ``determination'' of the state thus changes it, as indeed was only to be expected.}
\originalnote{179}{For example, because
\begin{align*}
RP_{[\varphi_n]}f&=R((f,\varphi_n)\cdot\varphi_n)
=(f,\varphi_n)\cdot R\varphi_n=\lambda_n(f,\varphi_n)\cdot\varphi_n,\\
P_{[\varphi_n]}Rf&=(Rf,\varphi_n)\cdot\varphi_n
=(f,R\varphi_n)\cdot\varphi_n=\lambda_n(f,\varphi_n)\cdot\varphi_n.
\end{align*}}
\originalnote{180}{Every textbook of classical (Hamiltonian) mechanics gives these relations.}
\originalnote{181}{The uncertainty relations for the time--energy pair have been discussed on several occasions. See the survey in Heisenberg, \emph{Die physikalischen Prinzipien der Quantentheorie}, \S~II.2.d, Leipzig 1930.}
\originalnote{182}{All other details of the measuring arrangement serve only to couple the quantity $\mathfrak R$ actually of interest, or its operator $R$, to the $M_n$, or $P^x,P^y,P^z$, or $Q^x,Q^y,Q^z$ mentioned here. Admittedly, this is precisely the most important practical aspect of measurement technique.}
\originalnote{183}{See, for example, Nernst, \emph{Theoretische Chemie}, Stuttgart (numerous editions since 1893), Book~IV, the discussion of the thermodynamic proof of the ``law of mass action''.}
\originalnote{184}{The reader will find the phenomenological system of thermodynamics built on this basis in numerous textbooks, for example Planck, \emph{Vorlesungen über Thermodynamik}, Berlin 1930. For what follows, it is chiefly the statistical formulation of the fundamental laws that matters, as expounded, for example, in the following papers: Einstein, \emph{Verh. d. dtsch. physik. Ges.}, vol.~12 (1914); L. Szilárd, \emph{Z. Physik}, vol.~32 (1925).}
\originalnote{185}{This is the Maxwell--Boltzmann method of statistical mechanics (see the survey in the article by P. and T. Ehrenfest in \emph{Enzykl. d. Math. Wiss.}, vol.~II.4.D, Leipzig 1907). In the theory of gases, for example, the ``very complicated'' system is the gas, consisting of many interacting molecules, and the molecules are investigated statistically.}
\originalnote{186}{This is Gibbs's method (see the reference in \hyperlink{note-185}{Note~185}). Here the individual system is the entire gas; many specimens of the same system, that is, the same gas, are followed through their histories, and their properties are evaluated statistically.}
\originalnote{187}{See the reference in \hyperlink{note-184}{Note~184}. This method was developed further by L. Szilárd.}
\originalnote{188}{As is well known, the kinetic theory of gases describes the process by which the walls impart their temperature to the gas enclosed by them. See the references in \hyperlink{note-184}{Note~184} and \hyperlink{note-185}{Note~185}.}
\originalnote{189}{If, in this transformation, the quantities of heat $\mathbf{Q}_1,\ldots,\mathbf{Q}_i$ are consumed at the respective temperatures $\mathbf{T}_1,\ldots,\mathbf{T}_i$, the entropy difference is $\frac{\mathbf{Q}_1}{\mathbf{T}_1}+\cdots+\frac{\mathbf{Q}_i}{\mathbf{T}_i}$. See the reference in \hyperlink{note-184}{Note~184}.}
\originalnote{190}{% Source PDF page 267 / original page 260
\hypertarget{origpage-260}{}%
If $c(\mathbf T)$ is the heat capacity of the given quantity of gas at temperature $\mathbf T$, it absorbs the quantity of heat $c(\mathbf T)\,d\mathbf T$ in the temperature interval $\mathbf T,\mathbf T+d\mathbf T$. According to \hyperlink{note-186}{Note~186}, the entropy difference is therefore $\int_{\mathbf T_1}^{\mathbf T_2}\frac{c(\mathbf T)\,d\mathbf T}{\mathbf T}$.\tnote{The source cites \hyperlink{note-186}{note~186}; the entropy relation used here is stated in \hyperlink{note-189}{note~189}.}}
\originalnote{191}{For an ideal gas $c(\mathbf T)$ is constant; at very small $\mathbf T$ this certainly ceases to hold. See, for example, the reference in \hyperlink{note-6}{Note~6}.}
\originalnote{192}{This requires the volume $V$ of $\overline{\boldsymbol K}$ to be large compared with the total volume of $\boldsymbol K_1,\ldots,\boldsymbol K_N$, and also that the ``energy per degree of freedom'', $k\mathbf T$ ($k$ is Boltzmann's constant), be large compared with $\frac{h^2}{\mu V^{2/3}}$ ($h$ is Planck's constant, $\mu$ the mass of the individual molecule; this quantity has the dimensions of an energy). See, for example, Fermi, \emph{Z. Physik}, vol.~36 (1926).}
\originalnote{193}{See, for example, the reference in \hyperlink{note-184}{Note~184}.}
\originalnote{194}{See the reference in \hyperlink{note-185}{Note~185}. The reader will find a detailed discussion of the difficulties associated with the concept of ``Maxwell's demon'' in L. Szilárd, \emph{Z. Physik}, vol.~53 (1929).}
\originalnote{195}{If an ideal gas consists of $M$ molecules, its pressure is $p=\frac{M\varkappa\mathbf T}{\mathcal V}$. In compression from volume $\mathcal V_1$ to volume $\mathcal V_2$, the mechanical work performed is therefore
\[
\int_{\mathcal V_1}^{\mathcal V_2}p\,d\mathcal V
=M\varkappa\mathbf T\int_{\mathcal V_1}^{\mathcal V_2}\frac{d\mathcal V}{\mathcal V}
=M\varkappa\mathbf T\ln\frac{\mathcal V_2}{\mathcal V_1}.
\]
In our case $M=N/2$, $\mathcal V_1=\mathcal V/2$, $\mathcal V_2=\mathcal V$.}
\originalnote{196}{As is well known, the energy of an ideal gas depends only on its temperature.}
\originalnote{197}{We shall determine the eigenvalues of $aP_{[\varphi]}+bP_{[\psi]}$ directly. We require
\[
(aP_{[\varphi]}+bP_{[\psi]})f=\lambda f.
\]
Since the left-hand side is a linear combination of $\varphi,\psi$, so is the right-hand side, and hence so is $f$ if $\lambda\ne0$. Certainly $\lambda=0$ is an eigenvalue of infinite multiplicity, since every $f$ orthogonal to $\varphi,\psi$ belongs to it. It therefore suffices to consider $\lambda\ne0$ and $f=x\varphi+y\psi$. (Let $\varphi,\psi$ be linearly independent; otherwise $\varphi=c\psi$, $|c|=1$, so the two states are identical.)

The equation above then becomes
\[
a(x+y(\psi,\varphi))\cdot\varphi+b(x(\varphi,\psi)+y)\cdot\psi
=\lambda x\cdot\varphi+\lambda y\cdot\psi,
\]
that is,
\[
ax+a\overline{(\varphi,\psi)}y=\lambda x,\qquad
b(\varphi,\psi)x+by=\lambda y.
\]
The determinant of these equations must vanish:
\begin{gather*}
\begin{vmatrix}a-\lambda&a\overline{(\varphi,\psi)}\\b(\varphi,\psi)&b-\lambda\end{vmatrix}=0,
\qquad(a-\lambda)(b-\lambda)-ab|(\varphi,\psi)|^2=0,\\
\lambda^2-(a+b)\lambda+ab(1-\alpha^2)=0,\\
\lambda=\frac{a+b\pm\sqrt{(a+b)^2-4ab(1-\alpha^2)}}2
=\frac{a+b\pm\sqrt{(a-b)^2+4\alpha^2ab}}2.
\end{gather*}
Putting $a=1,b=0$, or $a=\frac12,b=\frac12$, or $a=\frac13,b=\frac23$, gives the formulae in the text.}
\originalnote{198}{% Source PDF page 268 / original page 261
\hypertarget{origpage-261}{}%
Since $(x\ln x)'=\ln x+1$, and hence
\begin{align*}
\left(\frac{1+y}2\ln\frac{1+y}2+\frac{1-y}2\ln\frac{1-y}2\right)'
&=\frac12\left(\ln\frac{1+y}2+1\right)-\frac12\left(\ln\frac{1-y}2+1\right)\\
&=\frac12\ln\frac{1+y}{1-y},
\end{align*}
the derivative of our expression is
\begin{align*}
&-3\cdot\frac12\ln\frac{3+\sqrt{1+8\alpha^2}}{3-\sqrt{1+8\alpha^2}}
\cdot\frac13\frac{8\alpha}{\sqrt{1+8\alpha^2}}+4\cdot\frac12\ln\frac{1+\alpha}{1-\alpha}\\
&\qquad=2\left(\ln\frac{1+\alpha}{1-\alpha}
-\frac{2\alpha}{\sqrt{1+8\alpha^2}}\ln\frac{3+\sqrt{1+8\alpha^2}}{3-\sqrt{1+8\alpha^2}}\right).
\end{align*}
To say that this is $>0$ means
\[
\ln\frac{1+\alpha}{1-\alpha}>
\frac{2\alpha}{\sqrt{1+8\alpha^2}}\ln\frac{3+\sqrt{1+8\alpha^2}}{3-\sqrt{1+8\alpha^2}},
\]
that is,
\[
\frac1{2\alpha}\ln\frac{1+\alpha}{1-\alpha}>
\frac23\cdot\frac1{2\beta}\ln\frac{1+\beta}{1-\beta},
\qquad\beta=\frac{\sqrt{1+8\alpha^2}}3.
\]
We shall even prove this with $\frac89$ in place of $\frac23$. Since $1-\beta^2=\frac89(1-\alpha^2)$ and $\alpha<\beta$ (the latter follows from the former because $\alpha<1$), this says
\[
\frac{1-\alpha^2}{2\alpha}\ln\frac{1+\alpha}{1-\alpha}>
\frac{1-\beta^2}{2\beta}\ln\frac{1+\beta}{1-\beta}.
\]
This will be established once $\frac{1-x^2}{2x}\ln\frac{1+x}{1-x}$ is recognized as monotonically decreasing for $0<x<1$. That can be seen, for example, from its power series:
\begin{align*}
\frac{1-x^2}{2x}\ln\frac{1+x}{1-x}
&=(1-x^2)\left(1+\frac{x^2}3+\frac{x^4}5+\cdots\right)\\
&=1-\left(1-\frac13\right)x^2-\left(\frac13-\frac15\right)x^4-\cdots.
\end{align*}}
\originalnote{199}{For this device, characteristic of the phenomenological thermodynamic method, see, for example, the reference in \hyperlink{note-184}{Note~184}.}
\originalnote{200}{Of course we could omit the factor $N\varkappa$ here and consider $-\operatorname{Spur}(U\ln U)$. Or, to retain proportionality to the number $N$ of elements, we could take $-N\operatorname{Spur}(U\ln U)$.}
\originalnote{201}{See, for example, Planck, \emph{Theorie der Wärmestrahlung}, Leipzig 1913.}
\originalnote{202}{In the paper cited in \hyperlink{note-194}{Note~194}, L. Szilárd showed that this ``knowledge'' cannot be acquired for less than a compensating entropy increase $\varkappa\ln2$. In general, $\varkappa\ln2$ is the ``thermodynamic value'' of knowing which of two alternatives obtains. All attempts to carry out the process described above without knowing in which half of the container the molecule is situated can be shown to be unsound, even though they may lead to quite complicated automatic mechanisms.}
\originalnote{203}{This characterization of the macroscopic observer is due to E. Wigner.}
\originalnote{204}{If $E_n$ commutes with $\mathsf H$, and hence with $A$, the equality does hold, because
\begin{align*}
\operatorname{Spur}(A\cdot UA^{-1}E_n)
&=\operatorname{Spur}(UA^{-1}E_n\cdot A)\\
&=\operatorname{Spur}(UA^{-1}AE_n)=\operatorname{Spur}(UE_n).
\end{align*}
But the $E_n$, that is, all macroscopically observable quantities, by no means
% Source PDF page 269 / original page 262
\hypertarget{origpage-262}{}%
all commute with $\mathsf H$: many, for example the centre of gravity of a gas during diffusion, change appreciably with $t$, so $\operatorname{Spur}(UE_n)$ is not constant. Since all macroscopic quantities commute, $\mathsf H$ cannot be a macroscopic quantity; that is, energy cannot be measured macroscopically with perfect accuracy. This is plausible in any case.}
\originalnote{205}{For the classical $\mathsf H$-theorem see Boltzmann, \emph{Vorlesungen über Gastheorie}, Leipzig 1896, and the extremely instructive discussion by P. and T. Ehrenfest in the reference cited in \hyperlink{note-185}{Note~185}. The ``assumptions of disorder'' that can take the place of Boltzmann's in quantum mechanics were formulated by W. Pauli (\emph{Sommerfeld-Festschrift}, 1928), where they are used to prove the $\mathsf H$-theorem. Recently the author also succeeded in proving the ergodic theorem of classical mechanics; see \emph{Proc. Nat. Ac.}, January and March 1930, and the strengthening by G. D. Birkhoff, \emph{Proc. Nat. Ac.}, December 1929 and March 1930.}
\originalnote{206}{% The source prints the following note number as 266 (in the sequence 205, 266, 207).
\emph{Z. Physik}, vol.~57 (1929).\tnote{The source prints this note number as 266; its position and the reference in the text identify it as note 206.}}
\originalnote{207}{N. Bohr, \emph{Naturwiss.}, vol.~17 (1929), first pointed out that the duality in the description of nature, formally unavoidable in quantum mechanics, is also justified on substantive grounds. He also indicated its connection with the principle of psychophysical parallelism.}
\originalnote{208}{The author owes essential parts of the following discussion, and of that in \hyperref[sec:VI.3]{VI.3}, to conversations with Mr. L. Szilárd. See also the similar considerations in Heisenberg, in the reference cited in \hyperlink{note-181}{Note~181}.}
\originalnote{209}{It is easily shown that if $\mathsf A$ is Hermitian, or hypermaximal, then so is $\boldsymbol A$.}
\originalnote{210}{This is clear for \hyperref[cond:IV.2.I]{I}, and for \hyperref[cond:IV.2.II]{II} as well so long as only polynomials are involved. For general functions it follows from the fact that our mapping $\mathsf A\to\boldsymbol A$ does not destroy the correspondence between a resolution of the identity and a Hermitian operator.}
\originalnote{211}{Because of the large number and variety of indices, we use this notation for matrices, which differs somewhat from the notation used hitherto.}
\originalnote{212}{The projections of a state of $I+II$ are in general mixtures in $I$ and $II$, respectively; see below. This fact was discovered by J. Landau, \emph{Z. Physik}, vol.~45 (1927).}
\originalnote{213}{The mathematical discussion follows a paper by E. Schmidt, \emph{Math. Ann.}, vol.~63 (1907).}
\originalnote{214}{Several variants of this attempted explanation are also possible, but must be set aside for similar reasons.}
\originalnote{215}{The corresponding calculation for the case of position measurement discussed in \hyperref[sec:III.4]{III.4} is contained in a paper by Weizsäcker, \emph{Z. Physik}, vol.~70 (1931).}

% Keep display spacing at its font-dependent natural size inside page footnotes.
% TeX's insertion split may accept a taller natural box by shrinking its display
% skips, but the page builder then reserves the natural box height. The mismatch
% can move the whole footnote and its call to the next page, leaving a large gap.
% Remove only this shrink component, preserving natural size and finite stretch.
% Ordinary body displays, footnote numbering, and the split penalty are unchanged.
\RequirePackage{etoolbox}
\makeatletter
\newcommand*{\vn@unshrinkdisplayskip}[1]{%
  #1=\dimexpr#1\relax plus\gluestretch#1\relax
}
\newcommand*{\vn@footnotedisplayspacing}{%
  \vn@unshrinkdisplayskip\abovedisplayskip
  \vn@unshrinkdisplayskip\belowdisplayskip
  \vn@unshrinkdisplayskip\abovedisplayshortskip
  \vn@unshrinkdisplayskip\belowdisplayshortskip
}
\AtBeginDocument{%
  \patchcmd{\@footnotetext}
    {\reset@font\footnotesize}
    {\reset@font\footnotesize\vn@footnotedisplayspacing}
    {\typeout{VN footnote display spacing: active}}
    {% KOMA-Script wraps the native insertion in a saved implementation.
     \patchcmd{\scr@saved@footnotetext}
       {\reset@font\footnotesize}
       {\reset@font\footnotesize\vn@footnotedisplayspacing}
       {\typeout{VN footnote display spacing: active (KOMA-Script)}}
       {\PackageError{vn-footnotes}{Could not set footnote display spacing}
          {The native footnote insertion changed; inspect its definition.}}}%
}
\makeatother

\begin{document}

% Consolidated source: frontmatter.tex
% The first sequential PDF page is deliberately omitted at the user's request.
% Source PDF page 2.
\begin{center}
{\large Grundlehren der mathematischen Wissenschaften 38\par}
\medskip
A Series of Comprehensive Studies in Mathematics\par
\vfill
{\Large Johann von Neumann\par}
\bigskip
{\huge\bfseries Mathematische Grundlagen der Quantenmechanik\par}
\bigskip
{\huge\bfseries Mathematical Foundations of Quantum Mechanics\par}
\bigskip
{\large Reprint\par}
\vfill
\includegraphics[height=4\baselineskip]{figures/publisher-emblem-reprint.png}\par
\medskip
Springer-Verlag\par
Berlin Heidelberg New York 1981
\end{center}
\thispagestyle{empty}
\clearpage
% Source PDF page 3.
\section*{Publication Data}
Short-title cataloging-in-publication record prepared by the German Library:

Neumann, Johann von:\par
\emph{Mathematical Foundations of Quantum Mechanics} / Johann von Neumann.---Reprint [of the edition] Berlin, Springer, 1932.---Berlin, Heidelberg, New York: Springer, 1981.\par
(Grundlehren der mathematischen Wissenschaften; vol.~38)

ISBN-13: 978-3-642-96049-9\par
DOI: \href{https://doi.org/10.1007/978-3-642-96048-2}{10.1007/978-3-642-96048-2}\par
e-ISBN-13: 978-3-642-96048-2

AMS Subject Classifications (1970):\par
Primary: 81 A 03, 81 A 09\par
Secondary: 46 N 05, 47 A 05, 47 A 10

This work is protected by copyright. The rights thereby established, in particular those of translation, reprinting, extraction of illustrations, broadcasting, reproduction by photomechanical or similar means, and storage in data-processing systems, are reserved, even where only extracts are used. For reproductions for commercial purposes, remuneration must be paid to the publisher under Section~54 of the Copyright Act; the amount is to be agreed with the publisher.

\copyright\ by Springer-Verlag Berlin/Heidelberg 1971.\par
Softcover reprint of the hardcover 1st edition 1971.

Reprographic reprinting: Beltz Offsetdruck, Hemsbach/Bergstrasse.\par
Binding: Graphischer Betrieb Konrad Triltsch, Würzburg.\par
2140/3014-543

\clearpage
% Source PDF page 4.
\begin{center}
{\Large\bfseries The Fundamental Principles of the Mathematical Sciences\par}
\medskip
A Series of Monographs with Particular Attention to Applications\par
\bigskip
In collaboration with\par
W.~Blaschke (Hamburg), M.~Born (Göttingen), C.~Runge (Göttingen)\par
\medskip
Edited by R.~Courant (Göttingen)\par
\bigskip
Volume XXXVIII\par
\bigskip
{\LARGE\bfseries Mathematical Foundations of Quantum Mechanics\par}
\bigskip
by Johann v.~Neumann\par
\vfill
\includegraphics[height=4\baselineskip]{figures/publisher-emblem-1932.png}\par
\bigskip
Berlin\par
Julius Springer, Publisher\par
1932
\end{center}
\clearpage
% Source PDF page 5.
\begin{center}
{\LARGE\bfseries Mathematische Grundlagen der Quantenmechanik\par}
\bigskip
{\LARGE\bfseries Mathematical Foundations of Quantum Mechanics\par}
\bigskip
by\par
{\Large Johann v.~Neumann\par}
\bigskip
With 4 Illustrations\par
\vfill
\includegraphics[height=4\baselineskip]{figures/publisher-emblem-1932.png}\par
\bigskip
Berlin\par
Julius Springer, Publisher\par
1932
\end{center}
\clearpage
% Source PDF page 6.
\null\vfill
\begin{center}
\large Dedicated to the Memory of My Father
\end{center}
\vfill\null
\clearpage
% Source PDF page 7: the translated contents is generated from the full text.
\tableofcontents


% Consolidated source: introduction-I.tex
% Source PDF page 8 / original page 1.
\chapter*{Introduction}\label{chap:introduction}
\addcontentsline{toc}{chapter}{Introduction}
\hypertarget{origpage-1}{}%
This book presents a unified and, as far as possible and appropriate, mathematically sound account of the new quantum mechanics. Over the last few years, this theory has taken on a form that is presumably definitive in its essentials: the so-called ``transformation theory.'' Our main concern will be the general and fundamental questions that have arisen in connection with it. In particular, we shall examine more closely the difficult questions of interpretation, many of which are still not fully resolved. The relation of quantum mechanics to statistics and to classical statistical mechanics is especially important here. We shall generally leave aside applications of quantum-mechanical methods to particular problems and the more specialized theories that branch off from the general theory, at least insofar as we can do so without impairing an understanding of the general connections. This seems all the more appropriate because several excellent accounts of these matters already exist or are being published.\origNote{1}

We shall, however, develop the mathematical tools needed for the theory: the theory of Hilbert space and its so-called Hermitian operators. This calls for a detailed treatment of unbounded operators as well, extending the theory beyond its classical scope as established by \textsc{Hilbert}, E.~\textsc{Hellinger}, F.~\textsc{Riesz}, E.~\textsc{Schmidt}, and O.~\textsc{Toeplitz}. A word about our method is in order. As a rule, we shall calculate with the operators themselves, which represent physical quantities, rather than with the matrices obtained from them by introducing a particular, arbitrarily chosen coordinate system into Hilbert space. This ``coordinate-free,'' or invariant, approach is strongly geometrical and has considerable formal advantages.

\textsc{Dirac}, in several papers and in his recently published book, has given an account of quantum mechanics that can scarcely be surpassed in brevity and elegance and is likewise invariant in character.\origNote{2}
% Source PDF page 9 / original page 2; preceding sentence crosses the page boundary.
\hypertarget{origpage-2}{}%
It may therefore be appropriate to explain why we adopt a substantially different method.

Because of its clarity and elegance, \textsc{Dirac}'s method has now been adopted in much of the literature on quantum mechanics. Yet it fails to meet the requirements of mathematical rigor, even if these are reasonably relaxed to the standard customary in theoretical physics. For example, it consistently maintains the fiction that every self-adjoint operator can be put into diagonal form. For operators that cannot, this requires the introduction of ``improper'' functions with self-contradictory properties. Such mathematical ``fictions'' may sometimes be unavoidable even when one merely wants to calculate the numerical result of an intuitively specified experiment. There would be no objection to this if these conceptual constructions, which do not fit within the present framework of analysis, were really essential to the new physical theory. Just as Newtonian mechanics gave rise to the infinitesimal calculus, which in its original form was undoubtedly self-contradictory, quantum mechanics would then call for a reconstruction of our ``analysis of infinitely many variables.'' The mathematical apparatus would have to change, rather than the physical theory. But this is by no means the case. We shall show that transformation theory can be given a mathematically sound foundation that is just as clear and unified. It must be emphasized that this is not a matter of making \textsc{Dirac}'s method mathematically precise and explicit. A different approach is required from the outset: one based on \textsc{Hilbert}'s spectral theory of operators.

In examining the fundamental questions, we shall show in particular how the statistical formulas of quantum mechanics can be derived from a few qualitative basic assumptions. We shall also discuss in detail whether its statistical character can be traced to ambiguity---that is, incompleteness---in our description of nature. Such an explanation would accord best with the general principle that every probability statement arises from incomplete knowledge. This explanation ``by hidden parameters'' has indeed been proposed repeatedly, as has a related explanation that assigns the ``hidden parameters'' to the observer rather than to the observed system. It turns out, however, that such an explanation can scarcely be made satisfactory; more precisely, it is incompatible with certain qualitative basic postulates of quantum mechanics.\origNote{3}
% Source PDF page 10 / original page 3; preceding sentence crosses the page boundary.
\hypertarget{origpage-3}{}%

We shall also consider the relation of quantum-mechanical statistics to thermodynamics. Closer investigation shows that we can overcome the familiar difficulties of classical mechanics associated with the ``assumptions of disorder'' needed to establish thermodynamics.\origNote{4}

% Source PDF page 11 / original page 4.
\chapter{Introductory Considerations}\label{chap:I}
\hypertarget{origpage-4}{}%
\section{The Origin of Transformation Theory}\label{sec:I.1}
This is not the place to dwell on the great successes attained by quantum theory during the period from 1900 to 1925, a development dominated by the names of \textsc{Planck}, \textsc{Einstein}, and \textsc{Bohr}.\origNote{5} At the end of this development it was clear and practically beyond doubt that all elementary processes, that is, everything occurring on the atomic and molecular scale, are governed by the ``discontinuous'' laws of the quanta. In almost every area there were also quantitative quantum-theoretical methods, which usually yielded results in good or tolerable agreement with experience. Of fundamental importance was the acceptance, in theoretical physics, of the idea that the principle of continuity (``natura non facit saltus'') prevailing in the perceived macroscopic world is merely simulated by a process of averaging in a world that is discontinuous in its essence: we generally perceive only the combined effect of many times $10^{24}$ elementary processes\tnote{The German long-scale \emph{Quadrillionen} is rendered numerically: one Quadrillion is $10^{24}$.} at once, so that the all-leveling law of large numbers completely conceals the true nature of the individual processes.

Nevertheless, there was still no mathematical and physical framework for quantum theory that brought together everything then known, let alone one with the monumental unity of mechanics, electrodynamics, and relativity theory---a unity that quantum phenomena had shattered. Despite quantum theory's evidently justified claim to universality, it lacked the necessary formal and conceptual apparatus; instead, it was a tangle of fundamentally different and independent approaches, heterogeneous in character and partly contradictory. The most conspicuous points were the correspondence principle, still partly rooted in classical mechanics and electrodynamics (but destined to play a decisive role in the eventual clarification), and the self-contradictory duality in the nature of light (waves and corpuscles; cf.\ \hyperlink{note-5}{Note~5} and \hyperlink{note-148}{Note~148}). Finally, there was the existence of unquantized (aperiodic) and quantized (periodic or multiply periodic) motions.\origNote{6}

The year 1925 brought the resolution. An approach of \textsc{Heisenberg}'s was developed by \textsc{Born}, \textsc{Heisenberg}, and \textsc{Jordan}, and shortly afterward by \textsc{Dirac}, into a new quantum theory---the first complete, coherent system that physics had possessed in this field.
% Source PDF page 12 / original page 5.
\hypertarget{origpage-5}{}%
A little later, starting from an entirely different point, \textsc{Schrödinger} found ``wave mechanics,'' which accomplished the same thing and soon proved equivalent to the \textsc{Heisenberg--\allowbreak Born--\allowbreak Jordan} and \textsc{Dirac} system (at least in the mathematical sense; cf.\ \hyperref[sec:I.3]{I.3}, \hyperref[sec:I.4]{I.4}).\origNote{7} On the basis of \textsc{Born}'s statistical interpretation of the quantum-theoretical description of nature,\origNote{8} \textsc{Dirac} and \textsc{Jordan} were able to fuse the two theories into a single theory, ``transformation theory,''\origNote{9} in which the two complement one another and allow physical questions to be treated with particular mathematical simplicity.

We should also mention, although it lies outside our immediate subject, that after \textsc{Goudsmit} and \textsc{Uhlenbeck} had also discovered the magnetic moment and angular momentum of the electron, almost all the difficulties of the earlier quantum theory disappeared, so that we now possess an almost entirely satisfactory system of mechanics. Admittedly, the great unity with electrodynamics and relativity theory mentioned at the outset has not yet been recovered; but at least there is a mechanics of universal validity, into which the quantum laws fit with self-evident necessity and which satisfactorily explains the greater part of our experience.\origNote{10}

\section{The Original Formulations of Quantum Mechanics}\label{sec:I.2}
As a preliminary guide, we shall briefly set out the fundamental questions addressed by \textsc{Heisenberg--\allowbreak Born--\allowbreak Jordan} ``matrix mechanics'' and \textsc{Schrödinger} ``wave mechanics.''

In both theories we are initially given a problem of classical mechanics characterized by a Hamiltonian function $H(q_1,\ldots,q_k,p_1,\ldots,p_k)$. (As is explained in more detail in textbooks of mechanics, this means the following. The system has $k$ degrees of freedom, that is, its state at any given time is specified by the numerical values of $k$ coordinates $q_1,\ldots,q_k$. The energy is a given function of the coordinates and their time derivatives:
\[
 E=L(q_1,\ldots,q_k,\dot q_1,\ldots,\dot q_k),
\]
and is as a rule quadratic in the derivatives $\dot q_1,\ldots,\dot q_k$. The ``conjugate momenta'' $p_1,\ldots,p_k$ of the coordinates $q_1,\ldots,q_k$ are introduced by
\[
 p_1=\frac{\partial L}{\partial\dot q_1},\ldots,
 p_k=\frac{\partial L}{\partial\dot q_k};
\]
under the above assumption about $L$, they depend linearly on $\dot q_1,\ldots,\dot q_k$.
% Source PDF page 13 / original page 6.
\hypertarget{origpage-6}{}%
In any event, we can eliminate the $\dot q_1,\ldots,\dot q_k$ from $L$ by means of the $p_1,\ldots,p_k$, obtaining
\[
 E=L(q_1,\ldots,q_k,\dot q_1,\ldots,\dot q_k)
  =H(q_1,\ldots,q_k,p_1,\ldots,p_k).
\]
This $H$ is the Hamiltonian function.) In both theories one now wishes to learn from this Hamiltonian function as much as possible about the true, that is, quantum-theoretical, behavior of the system\origNote{11}---first of all to determine its possible energy levels, then to study the corresponding ``stationary states,'' to calculate their ``transition probabilities,'' and so forth.\origNote{12}

The prescription given by matrix theory for solving this problem is as follows. Find a system of $2k$ matrices $Q_1,\ldots,Q_k,P_1,\ldots,P_k$\origNote{13} which, first, satisfy the relations
\leavevmode\phantomsection\label{rel:I.2.commutation}
\begin{align*}
 Q_mQ_n-Q_nQ_m&=0,& P_mP_n-P_nP_m&=0,\\
 P_mQ_n-Q_nP_m&=
 \begin{cases}0,&m\ne n,\\ \dfrac{h}{2\pi i}\,1,&m=n,\end{cases}
 &&(m,n=1,\ldots,k),
\end{align*}
and for which, second, the matrix $W=H(Q_1,\ldots,Q_k,P_1,\ldots,P_k)$ is diagonal. (We shall not enter here into the origin of these equations, especially the first group, the so-called ``commutation relations,'' which govern the entire noncommutative matrix calculus of this theory. The reader will find exhaustive information on this in the works cited in \hyperlink{note-1}{Note~1}. Here $h$ is \textsc{Planck}'s quantum of action.) The diagonal elements of $W$, say $w_1,w_2,\ldots$, are then the various possible energy levels of the system. The elements of the matrices $Q_1,\ldots,Q_k$---$q^{(1)}_{mn},\ldots,q^{(k)}_{mn}$---determine in a certain manner the transition probabilities of the system (from the $m$th state with energy $w_m$ to the $n$th state with energy $w_n$, $w_m>w_n$) and the radiation emitted in the process.

Note also that the matrix
\[
 W=H(Q_1,\ldots,Q_k,P_1,\ldots,P_k)
\]
is not completely determined by $Q_1,\ldots,Q_k,P_1,\ldots,P_k$ and the classical Hamiltonian function $H(q_1,\ldots,q_k,p_1,\ldots,p_k)$, because the $Q_l$ and $P_l$ do not all commute under multiplication. In the classical interpretation of $H(q_1,\ldots,q_k,p_1,\ldots,p_k)$, by contrast, distinguishing a term $p_1q_1$ from a term $q_1p_1$ would be meaningless. In $H$, therefore, one must go beyond the classical meaning of this expression and specify the order of the factors $q_l$ and $p_l$ in its terms. This process has not been carried out in complete generality, but suitable conventions are known for the most important special cases.
% Source PDF page 14 / original page 7.
\hypertarget{origpage-7}{}%
(In the simple case in which the system under investigation consists of $\nu$ particles, and hence has $k=3\nu$ coordinates $q_1,\ldots,q_{3\nu}$---so that, for instance, $q_{3\mu-2},q_{3\mu-1},q_{3\mu}$ are the three Cartesian coordinates of the $\mu$th particle, $\mu=1,\ldots,\nu$---and the interaction of these particles is given by a potential energy $V(q_1,\ldots,q_{3\nu})$, there is, however, no such ambiguity. The classical Hamiltonian function is then
\begin{multline*}
 H(q_1,\ldots,q_{3\nu},p_1,\ldots,p_{3\nu})
 =\sum_{\mu=1}^{\nu}\frac{1}{2m_\mu}
 (p_{3\mu-2}^2+p_{3\mu-1}^2+p_{3\mu}^2)\\
 {}+V(q_1,\ldots,q_{3\nu}),
\end{multline*}
where $m_\mu$ is the mass of the $\mu$th particle and $p_{3\mu-2},p_{3\mu-1},p_{3\mu}$ are its momenta. It is perfectly clear what this means after the matrices $Q_1,\ldots,Q_{3\nu},P_1,\ldots,P_{3\nu}$ have been substituted; in particular, the $V$ term causes no difficulties, since all the $Q_1,\ldots,Q_{3\nu}$ commute with one another.) The important point is that only Hermitian matrices are admitted, that is, matrices $A=\{a_{mn}\}$ for which $a_{mn}=\overline{a_{nm}}$ holds identically (the elements $a_{mn}$ may be complex!). Consequently, $H(Q_1,\ldots,Q_k,P_1,\ldots,P_k)$ must be Hermitian whenever all the $Q_1,\ldots,Q_k,P_1,\ldots,P_k$ are Hermitian. This restricts the possible orderings of the factors discussed above, but not enough to determine $H(Q_1,\ldots,Q_k,P_1,\ldots,P_k)$ uniquely from the classical $H(q_1,\ldots,q_k,p_1,\ldots,p_k)$.\origNote{14}

By contrast, the prescription of wave mechanics is as follows. First form the Hamiltonian function $H(q_1,\ldots,q_k,p_1,\ldots,p_k)$, and then, for an arbitrary function $\psi(q_1,\ldots,q_k)$ in the configuration space of the system (and not in phase space, that is, the $p_1,\ldots,p_k$ are not to occur in $\psi$!), set up the differential equation
\begin{multline*}
 H\left(q_1,\ldots,q_k,\frac{h}{2\pi i}\frac{\partial}{\partial q_1},
 \ldots,\frac{h}{2\pi i}\frac{\partial}{\partial q_k}\right)\psi(q_1,\ldots,q_k)\\
 =\lambda\psi(q_1,\ldots,q_k).
\end{multline*}
Here $H\bigl(q_1,\ldots,q_k,\frac{h}{2\pi i}\frac{\partial}{\partial q_1},\ldots,
\frac{h}{2\pi i}\frac{\partial}{\partial q_k}\bigr)$ is understood in the natural sense as an operation on functions. For example, in the case mentioned above,
\begin{multline*}
 H(q_1,\ldots,q_{3\nu},p_1,\ldots,p_{3\nu})
 =\sum_{\mu=1}^{\nu}\frac{1}{2m_\mu}(p_{3\mu-2}^2+p_{3\mu-1}^2+p_{3\mu}^2)\\
 {}+V(q_1,\ldots,q_{3\nu}),
\end{multline*}
it transforms the function $\psi(q_1,\ldots,q_{3\nu})$ into
\[
 \sum_{\mu=1}^{\nu}\frac{1}{2m_\mu}\left(\frac{h}{2\pi i}\right)^2
 \left(\frac{\partial^2\psi}{\partial q_{3\mu-2}^2}
 +\frac{\partial^2\psi}{\partial q_{3\mu-1}^2}
 +\frac{\partial^2\psi}{\partial q_{3\mu}^2}\right)+V\psi
\]
% Source PDF page 15 / original page 8.
\hypertarget{origpage-8}{}%
(we have omitted the variables $q_1,\ldots,q_{3\nu}$ in $V$ and $\psi$). Since the operation $q_1\frac{h}{2\pi i}\frac{\partial}{\partial q_1}$ differs from the operation $\frac{h}{2\pi i}\frac{\partial}{\partial q_1}q_1$,\origNote{15} there is again an ambiguity concerning the order of the factors $q_m$ and $p_m$ in $H(q_1,\ldots,q_k,p_1,\ldots,p_k)$. But \textsc{Schrödinger} showed how a particular variational principle removes this indeterminacy while ensuring that the resulting differential equation is self-adjoint.\origNote{16}

This differential equation (the ``wave equation'') now has precisely the character of an eigenvalue problem: $\lambda$ is regarded as the eigenvalue parameter, and the eigenfunction $\psi=\psi(q_1,\ldots,q_k)$ is required, for example, to vanish on the boundary of configuration space (the space of the $q_1,\ldots,q_k$), and to be regular and single-valued within it. According to wave theory, the eigenvalues $\lambda$ (in both the point and the continuous spectrum\origNote{17}) are the possible energy levels. The associated (complex!) eigenfunctions $\psi$ are in turn related to the corresponding states of the system, which are stationary in \textsc{Bohr}'s sense. For example, consider a system of $\nu$ electrons ($k=3\nu$; cf.\ above), and let $e$ be the electron's charge. According to \textsc{Schrödinger}, the $\mu$th electron is to be thought of as ``smeared out'' over the whole $x,y,z$ ($=q_{3\mu-2},q_{3\mu-1},q_{3\mu}$) space. Its charge density, measured at the point $x,y,z$, is given by
\begin{multline*}
 e\underbrace{\int\cdots\int}_{3\nu-3\text{ fold}}
 |\psi(q_1\cdots q_{3\mu-3}xyzq_{3\mu+1}\cdots q_{3\nu})|^2\\
 {}\times dq_1\cdots dq_{3\mu-3}\,dq_{3\mu+1}\cdots dq_{3\nu}.
\end{multline*}
(For the total charge to come out as $e$, $\psi$ must be normalized by the condition
\[
 \underbrace{\int\cdots\int}_{3\nu\text{ fold}}
 |\psi(q_1\cdots q_{3\nu})|^2\,dq_1\cdots dq_{3\nu}=1
\]
[integration over all $3\nu$ variables!]. The same equation results for every $\mu=1,\ldots,\nu$.)

Wave mechanics also describes systems that are not in stationary states in \textsc{Bohr}'s sense.\origNote{18} If the state is not stationary, that is, if it changes with time, then the wave function $\psi=\psi(q_1,\ldots,q_k;t)$ contains the time $t$ and changes according to the differential equation\origNote{19}
\begin{multline*}
 -H\left(q_1,\ldots,q_k,\frac{h}{2\pi i}\frac{\partial}{\partial q_1},
 \ldots,\frac{h}{2\pi i}\frac{\partial}{\partial q_k}\right)\psi(q_1,\ldots,q_k;t)\\
 =\frac{h}{2\pi i}\frac{\partial}{\partial t}\psi(q_1,\ldots,q_k;t).
\end{multline*}
That is, $\psi$ can be prescribed arbitrarily at $t=t_0$ and is then uniquely determined for every $t$. Comparing \textsc{Schrödinger}'s two differential equations shows that even the wave functions $\psi$ of stationary states depend on $t$, but only in the form
% Source PDF page 16 / original page 9.
\hypertarget{origpage-9}{}%
\[
 \psi(q_1,\ldots,q_k;t)=e^{-\frac{2\pi i}{h}\lambda t}\psi(q_1,\ldots,q_k;0).
\]
Thus $t$ occurs only in a factor independent of $q_1,\ldots,q_k$ (that is, constant in configuration space) and of absolute value $1$, so that, for example, the charge-density distribution defined above does not change. (More generally, we may expect any factor in $\psi$ that is constant in configuration space and has absolute value $1$ to be unobservable in principle. Our later detailed analysis will confirm this.)

Since the eigenfunctions of the first differential equation form a complete orthogonal system,\origNote{20} we can expand every $\psi=\psi(q_1,\ldots,q_k)$ in terms of them. Let the eigenfunctions be $\psi_1,\psi_2,\ldots$ (all independent of $t$ again!), and their respective eigenvalues $\lambda_1,\lambda_2,\ldots$. The expansion is\origNote{21}
\[
 \psi(q_1,\ldots,q_k)=\sum_{n=1}^{\infty}a_n\psi_n(q_1,\ldots,q_k).
\]
If $\psi$ does depend on $t$, then $t$ will enter the expansion coefficients $a_n$ (whereas the eigenfunctions $\psi_1,\psi_2,\ldots$ are, now as always in what follows, to be independent of $t$). If, therefore, the given $\psi=\psi(q_1,\ldots,q_k)$ is in fact $\psi(q_1,\ldots,q_k;t_0)$, then, in view of
\begin{align*}
 \psi=\psi(q_1,\ldots,q_k;t)&=\sum_{n=1}^{\infty}a_n(t)\psi_n,\\
 H\psi&=\sum_{n=1}^{\infty}a_n(t)H\psi_n
       =\sum_{n=1}^{\infty}\lambda_na_n(t)\psi_n,\\
 \frac{h}{2\pi i}\frac{\partial\psi}{\partial t}
   &=\sum_{n=1}^{\infty}\frac{h}{2\pi i}\dot a_n(t)\psi_n,
\end{align*}
comparison of coefficients in the second differential equation gives
\[
 \frac{h}{2\pi i}\dot a_n(t)=-\lambda_na_n(t),\qquad
 a_n(t)=c_ne^{-\frac{2\pi i}{h}\lambda_nt},
\]
that is,
\begin{align*}
 a_n(t)&=e^{-\frac{2\pi i}{h}\lambda_n(t-t_0)}a_n(t_0)
        =e^{-\frac{2\pi i}{h}\lambda_n(t-t_0)}a_n,\\
 \psi=\psi(q_1,\ldots,q_k;t)
   &=\sum_{n=1}^{\infty}e^{-\frac{2\pi i}{h}\lambda_n(t-t_0)}
     a_n\psi_n(q_1,\ldots,q_k).
\end{align*}
If, then, $\psi$ is not stationary, that is, if more than one of the $a_n$ is nonzero, $\psi$ (as $t$ varies) no longer changes merely by a factor constant in configuration space and of absolute value $1$. Consequently, the charge densities defined above will in general also change; that is, actual electrical oscillations take place in space.\origNote{22}
% Source PDF page 17 / original page 10.
\hypertarget{origpage-10}{}%

As we see, the two theories differ considerably in their concepts and practical prescriptions. Nevertheless, they always yielded the same results from the outset, even in details where both departed from the older versions of quantum theory.\origNote{23} As mentioned in \hyperref[sec:I.1]{I.1}, this remarkable fact was soon explained by \textsc{Schrödinger}'s proof of their mathematical equivalence.\origNote{24} We shall now turn to this proof of equivalence and at the same time set out the general ``transformation theory'' of \textsc{Dirac} and \textsc{Jordan}, which encompasses both theories.

\section{Equivalence of the Two Theories: Transformation Theory}\label{sec:I.3}
The fundamental problem of matrix theory was to determine the matrices $Q_1,\ldots,Q_k,P_1,\ldots,P_k$ so that, first, the commutation relations of \hyperref[sec:I.2]{I.2} (\hyperref[rel:I.2.commutation]{p.~5})\tnote{The source cites p.~5 here, but the commutation relations are printed on original p.~6. The link leads directly to those relations.} are satisfied and, second, a certain function of them, $H(Q_1,\ldots,Q_k,P_1,\ldots,P_k)$, is a diagonal matrix. In their very first publication, \textsc{Born} and \textsc{Jordan} divided the problem into the following two steps.

First they sought matrices $\overline Q_1,\ldots,\overline Q_k,\overline P_1,\ldots,\overline P_k$ which need only satisfy the commutation relations---something easily accomplished;\origNote{25} as a rule,
\[
 \overline{\mathsf{H}}=H(\overline Q_1,\ldots,\overline Q_k,\overline P_1,\ldots,\overline P_k)
\]
was then not a diagonal matrix. Next they sought the correct solutions in the form
\begin{align*}
 Q_1&=\mathsf{S}^{-1}\overline Q_1\mathsf{S},\ldots,Q_k=\mathsf{S}^{-1}\overline Q_k\mathsf{S},\\
 P_1&=\mathsf{S}^{-1}\overline P_1\mathsf{S},\ldots,P_k=\mathsf{S}^{-1}\overline P_k\mathsf{S},
\end{align*}
where $\mathsf{S}$ could be an arbitrary matrix (provided, of course, that it had an inverse $\mathsf{S}^{-1}$ with the properties $\mathsf{S}^{-1}\mathsf{S}=\mathsf{S}\mathsf{S}^{-1}=1$). The commutation relations for $\overline Q_1,\ldots,\overline Q_k,\overline P_1,\ldots,\overline P_k$ ensure that $Q_1,\ldots,Q_k,P_1,\ldots,P_k$ also satisfy them, whatever the choice of $\mathsf{S}$! Moreover, the transformation takes $\overline{\mathsf{H}}=H(\overline Q_1,\ldots,\overline Q_k,\overline P_1,\ldots,\overline P_k)$ into $\mathsf{H}=H(Q_1,\ldots,Q_k,P_1,\ldots,P_k)$ according to $\mathsf{H}=\mathsf{S}^{-1}\overline{\mathsf{H}}\mathsf{S}$.\origNote{26} Thus the only remaining requirement on $\mathsf{S}$ is that, for the given $\overline{\mathsf{H}}$, the matrix $\mathsf{S}^{-1}\overline{\mathsf{H}}\mathsf{S}$ be diagonal. (One would also have to ensure that $\mathsf{S}^{-1}\overline Q_1\mathsf{S}$, etc., are Hermitian, just as $\overline Q_1$, etc., were. Closer inspection shows, however, that this additional requirement on $\mathsf{S}$ can always be satisfied afterward, and therefore we shall disregard it entirely in the present preliminary discussion.)

Thus a given $\overline{\mathsf{H}}$ must be transformed into diagonal form according to the scheme $\mathsf{S}^{-1}\overline{\mathsf{H}}\mathsf{S}$. Let us therefore formulate precisely what this means.
% Source PDF page 18 / original page 11.
\hypertarget{origpage-11}{}%

Let the matrix $\overline{\mathsf{H}}$ have the elements $h_{\mu\nu}$, the unknown matrix $\mathsf{S}$ the elements $s_{\mu\nu}$, and the diagonal matrix $\mathsf{H}$, likewise unknown, the diagonal elements $w_\mu$, that is, the general element $w_\mu\delta_{\mu\nu}$.\origNote{27} The statement $\mathsf{H}=\mathsf{S}^{-1}\overline{\mathsf{H}}\mathsf{S}$ is equivalent to $\mathsf{S}\mathsf{H}=\overline{\mathsf{H}}\mathsf{S}$, and this means (equating the corresponding elements on both sides as determined by the familiar rules of matrix multiplication)
\[
 \sum_\nu s_{\mu\nu}w_\nu\delta_{\nu\varrho}
 =\sum_\nu h_{\mu\nu}s_{\nu\varrho},
\]
that is,
\[
 \sum_\nu h_{\mu\nu}s_{\nu\varrho}=w_\varrho s_{\mu\varrho}.
\]
Thus the individual columns $s_{1\varrho},s_{2\varrho},\ldots$ of the matrix $\mathsf{S}$ ($\varrho=1,2,\ldots$), together with the corresponding diagonal elements $w_\varrho$ of $\mathsf{H}$, solve the so-called eigenvalue problem
\[
 \sum_\nu h_{\mu\nu}x_\nu=\lambda x_\mu\qquad(\mu=1,2,\ldots).
\]
(The trivial solution $x_1=x_2=\cdots=0$ is of course excluded.) Indeed, $x_\nu=s_{\nu\varrho}$, $\lambda=w_\varrho$ is a solution. (The possibility $x_\nu\equiv0$, that is, $s_{\nu\varrho}\equiv0$ [for all $\nu$], is ruled out, since then the $\varrho$th column of $\mathsf{S}$ would vanish identically although $\mathsf{S}$ has an inverse $\mathsf{S}^{-1}$!) It is noteworthy that these are essentially the only solutions.

The equation above says that transforming the vector $x=\{x_1,x_2,\ldots\}$ by the matrix $\overline{\mathsf{H}}$ is equivalent to multiplying it by the number $\lambda$. We transform $x=\{x_1,x_2,\ldots\}$ by $\mathsf{S}^{-1}$, obtaining a vector $y=\{y_1,y_2,\ldots\}$. Transforming $y$ by $\mathsf{H}$ amounts to transforming $x$ by $\mathsf{H}\mathsf{S}^{-1}=\mathsf{S}^{-1}\overline{\mathsf{H}}\mathsf{S}\cdot \mathsf{S}^{-1}=\mathsf{S}^{-1}\overline{\mathsf{H}}$, that is, to transforming $\lambda x$ by $\mathsf{S}^{-1}$; the result is therefore $\lambda y$. Now $\overline{\mathsf{H}}y$ has the components\tnote{The source prints $\overline{\mathsf{H}}y$ here, although the diagonal components used in this sentence are those of $\mathsf{H}y$.} $\sum_\nu w_\mu\delta_{\mu\nu}y_\nu=w_\mu y_\mu$, while $\lambda y$ has the components $\lambda y_\mu$. Thus $w_\mu y_\mu=\lambda y_\mu$ is required for all $\mu=1,2,\ldots$, that is, $y_\mu=0$ whenever $w_\mu\ne\lambda$. Denote by $\eta^\varrho$ the vector whose $\varrho$th component is $1$ and all of whose other components are $0$. The preceding statement then says that $y$ is a linear combination of those $\eta^\varrho$ for which $w_\varrho=\lambda$---in particular, it is zero if this never occurs. The vector $x$ is obtained by applying $\mathsf{S}$ to $y$; it is therefore a linear combination of the images of those $\eta^\varrho$ under $\mathsf{S}$. The $\mu$th component of $\mathsf{S}\eta^\varrho$ (since the $\nu$th component of $\eta^\varrho$ was $\delta_{\nu\varrho}$) is $\sum_\nu s_{\mu\nu}\delta_{\nu\varrho}=s_{\mu\varrho}$. Thus, if we regard the $\varrho$th column of $\mathsf{S}$, $s_{1\varrho},s_{2\varrho},\ldots$, as a vector, $x$ is a linear combination of all columns for which $w_\varrho=\lambda$---in particular, zero if this never occurs. This proves our original assertion: the $w_1,w_2,\ldots$ are the only eigenvalues, and $x_\nu=s_{\nu\varrho}$, $\lambda=w_\varrho$ are essentially the only solutions.

This is very important: knowledge of $\mathsf{S},\mathsf{H}$ not only determines all solutions of the eigenvalue problem, but conversely, once we have solved the eigenvalue problem completely, we can determine $\mathsf{S},\mathsf{H}$ from it. For example, $\mathsf{H}$ is obtained as follows: the $w_\mu$ are simply all the solutions $\lambda$, each such $\lambda$ appearing in the sequence $w_1,w_2,\ldots$ as many times as there are linearly independent solution sequences $x_1,x_2,\ldots$ corresponding to it.\origNote{28}
% Source PDF page 19 / original page 12.
\hypertarget{origpage-12}{}%
This already determines the $w_1,w_2,\ldots$ up to their order.\origNote{29}

The central problem of matrix theory is thus the solution of the eigenvalue equation
\begin{equation}
 \sum_\nu h_{\mu\nu}x_\nu=\lambda x_\mu\qquad(\mu=1,2,\ldots).
 \tag{$E_1$}\label{eq:I.3.E1}
\end{equation}
Let us now turn to wave theory. Its fundamental equation is the ``wave equation''
\begin{equation}
 H\varphi(q_1\cdots q_k)=\lambda\varphi(q_1\cdots q_k),
 \tag{$E_2$}\label{eq:I.3.E2}
\end{equation}
where $H$ is the differential operator already discussed. All solutions $\varphi(q_1\cdots q_k)$ and $\lambda$ are sought, excluding the trivial solution $\varphi(q_1\cdots q_k)\equiv0$, $\lambda$ arbitrary. This is analogous to what was required in \hyperref[eq:I.3.E1]{$E_1$}: the sequence $x_1,x_2,\ldots$, which we can also regard as a function $x_\nu$ of the ``discontinuous'' variable $\nu$ (with values $1,2,\ldots$), corresponds to the function $\varphi(q_1\cdots q_k)$ of the ``continuous'' variables $q_1,\ldots,q_k$; $\lambda$ plays the same role in both cases. But the linear transformation
\[
 x_\mu\longmapsto\sum_\nu h_{\mu\nu}x_\nu
\]
bears little resemblance to the other transformation,
\[
 \varphi(q_1\cdots q_k)\longmapsto H\varphi(q_1\cdots q_k).
\]
How can we establish an analogy between them?

We have treated the index $\nu$ as a variable analogous to the $k$ variables $q_1,\ldots,q_k$. Thus a positive integer corresponds to a general point in the $k$-dimensional configuration space, which we shall henceforth denote by $\Omega$. We should therefore not expect a sum over $\Omega$ to be the counterpart of $\sum_\nu$. Instead, the integral $\int_\Omega\cdots\int\cdots\,dq_1\cdots dq_k$ (or, more briefly, $\int_\Omega\cdots\,dv$, where $dv$ is the volume element $dq_1\cdots dq_k$ in $\Omega$) is the proper analogue. The matrix element $h_{\mu\nu}$, which depends on two integer indices, thus corresponds to a function $h(q_1,\ldots,q_k;q'_1,\ldots,q'_k)$, where $q_1,\ldots,q_k$ and $q'_1,\ldots,q'_k$ range independently over $\Omega$. The transformation
\[
 x_\mu\longmapsto\sum_\nu h_{\mu\nu}x_\nu
 \quad\text{or}\quad
 x_\nu\longmapsto\sum_{\nu'}h_{\nu\nu'}x_{\nu'}
\]
thereby passes into
\[
 \varphi(q_1\cdots q_k)\longmapsto
 \int_\Omega\cdots\int h(q_1\cdots q_k;q'_1\cdots q'_k)
 \varphi(q'_1\cdots q'_k)\,dq'_1\cdots dq'_k,
\]
and the eigenvalue problem \hyperref[eq:I.3.E1]{$E_1$}, which we may also write as
\[
 \sum_{\nu'}h_{\nu\nu'}x_{\nu'}=\lambda x_\nu,
\]
passes into
\begin{multline}
 \int_\Omega\cdots\int h(q_1\cdots q_k;q'_1\cdots q'_k)
 \varphi(q'_1\cdots q'_k)\,dq'_1\cdots dq'_k\\
 =\lambda\varphi(q_1\cdots q_k).
 \tag{$E_3$}\label{eq:I.3.E3}
\end{multline}
% Source PDF page 20 / original page 13.
\hypertarget{origpage-13}{}%
Eigenvalue problems of type \hyperref[eq:I.3.E3]{$E_3$} have been extensively investigated in mathematics and can indeed be treated in close analogy with problems of type \hyperref[eq:I.3.E1]{$E_1$}. They are called ``integral equations.''\origNote{30}

Unfortunately, however, \hyperref[eq:I.3.E2]{$E_2$} does not have this form, or, rather, it can be brought into this form only if, for the differential operator
\[
 H=H\left(q_1,\ldots,q_k,\frac{h}{2\pi i}\frac{\partial}{\partial q_1},
 \ldots,\frac{h}{2\pi i}\frac{\partial}{\partial q_k}\right),
\]
a function $h(q_1\cdots q_k;q'_1\cdots q'_k)$ can be found such that identically (that is, for every $\varphi(q_1\cdots q_k)$)
\begin{multline}
 H\varphi(q_1\cdots q_k)\\
 =\int_\Omega\cdots\int h(q_1\cdots q_k;q'_1\cdots q'_k)
 \varphi(q'_1\cdots q'_k)\,dq'_1\cdots dq'_k.
 \tag{I}\label{eq:I.3.I}
\end{multline}
This $h(q_1\cdots q_k;q'_1\cdots q'_k)$, if it exists, is called the ``integral kernel'' of the functional operation $H$, and $H$ itself is then called an ``integral operator.''

Such a transformation is, however, in general impossible: differential operators $H$ are never integral operators. Even the simplest operator on functions, which maps each $\varphi$ to itself and which we denote by $1$, is not an integral operator. To see this, take $k=1$ for simplicity. What is required is
\begin{equation}
 \varphi(q)=\int_{-\infty}^{\infty}h(q,q')\varphi(q')\,dq'.
 \tag{$\Delta_1$}\label{eq:I.3.D1}
\end{equation}
In this equation replace $\varphi(q)$ by $\varphi(q_0+q)$, put $q=0$, and introduce the integration variable $q''=q'-q_0$.\tnote{The source prints $q''=q'-q_0$. The substitution needed for the formula immediately following it is $q''=q'+q_0$. The printed sign is retained here.} We then obtain
\[
 \varphi(q_0)=\int_{-\infty}^{\infty}h(0,q''-q_0)\varphi(q'')\,dq''.
\]
Replacing $q_0,q''$ by $q,q'$, we see that $h(0,q'-q)$ solves our problem along with $h(q,q')$, and hence that we may assume $h(q,q')$ to depend only on $q'-q$. The requirement then becomes
\begin{equation}
 \varphi(q)=\int_{-\infty}^{\infty}h(q'-q)\varphi(q')\,dq'
 \quad\bigl(h(q,q')=h(q'-q)\bigr).
 \tag{$\Delta_2$}\label{eq:I.3.D2}
\end{equation}
As another replacement of $\varphi(q)$ by $\varphi(q_0+q)$ shows, it is enough to consider $q=0$, that is,
\begin{equation}
 \varphi(0)=\int_{-\infty}^{\infty}h(q)\varphi(q)\,dq.
 \tag{$\Delta_3$}\label{eq:I.3.D3}
\end{equation}
Replacing $\varphi(q)$ by $\varphi(-q)$ shows that $h(-q)$ is a solution along with $h(q)$, and hence so is $h_1(q)=\frac{h(q)+h(-q)}2$; thus $h(q)$ may be assumed to be an even function of $q$.

It is clear that these conditions cannot be fulfilled: if we choose $\varphi(q)>0$ for $q\ne0$, $\varphi(0)=0$, then \hyperref[eq:I.3.D3]{$\Delta_3$} implies $\varphi(q)=0$ for $q\ne0$.\origNote{31}\tnote{The source prints $\varphi(q)=0$ in this conclusion; the argument concerns $h(q)=0$. The printed symbol is retained.}
% Source PDF page 21 / original page 14.
\hypertarget{origpage-14}{}%
But if we choose $\varphi(q)\equiv1$, we obtain $\int_{-\infty}^{\infty}h(q)\,dq=1$, whereas the preceding considerations certainly imply $\int_{-\infty}^{\infty}h(q)\,dq=0$.

Nevertheless, \textsc{Dirac} postulated the existence of such a function:
\begin{equation}
 \delta(q)=0\ \text{for }q\ne0,\qquad
 \delta(q)=\delta(-q),\qquad \int_{-\infty}^{\infty}\delta(q)\,dq=1.
 \tag{$\Delta_4$}\label{eq:I.3.D4}
\end{equation}
It would satisfy \hyperref[eq:I.3.D3]{$\Delta_3$}:
\begin{align*}
 \int_{-\infty}^{\infty}\delta(q)\varphi(q)\,dq
 &=\varphi(0)\int_{-\infty}^{\infty}\delta(q)\,dq
   +\int_{-\infty}^{\infty}\delta(q)\bigl(\varphi(q)-\varphi(0)\bigr)\,dq\\
 &=\varphi(0)\cdot1+\int_{-\infty}^{\infty}0\,dq=\varphi(0),
\end{align*}
and hence also \hyperref[eq:I.3.D1]{$\Delta_1$}, \hyperref[eq:I.3.D2]{$\Delta_2$}. It is to be imagined as vanishing everywhere except at zero, but as being so infinitely large there that the integral of $\delta(q)$ nevertheless comes out as $1$.\origNote{32}

Once this fiction has been accepted, a wide variety of differential operators can be represented as integral operators, provided that, besides $\delta(q)$, its derivatives are also introduced. Thus one finds
\begin{align*}
 \frac{d^n}{dq^n}\varphi(q)
 &=\frac{d^n}{dq^n}\int_{-\infty}^{\infty}\delta(q-q')\varphi(q')\,dq'\\
 &=\int_{-\infty}^{\infty}\frac{\partial^n}{\partial q^n}\delta(q-q')
    \cdot\varphi(q')\,dq'\\
 &=\int_{-\infty}^{\infty}\delta^{(n)}(q-q')\varphi(q')\,dq',\\
 q^n\varphi(q)&=\int_{-\infty}^{\infty}\delta(q-q')q^n\varphi(q')\,dq',
\end{align*}
that is, $\frac{d^n}{dq^n}$ and $q^n\cdots$ have integral kernels $\delta^{(n)}(q-q')$ and $\delta(q-q')q^n$, respectively. The same scheme can be used to find integral kernels of arbitrarily complicated differential operators. With several variables $q_1,\ldots,q_k$, products of $\delta$ functions serve the same purpose; for example,
\begin{align*}
 &\int_\Omega\cdots\int\delta(q_1-q'_1)\delta(q_2-q'_2)\cdots
 \delta(q_k-q'_k)\varphi(q'_1\cdots q'_k)\,dq'_1\cdots dq'_k\\
 &\quad=\int_{-\infty}^{\infty}\left[\cdots
 \left[\int_{-\infty}^{\infty}
 \left[\int_{-\infty}^{\infty}\varphi(q'_1q'_2\cdots q'_k)
 \delta(q_1-q'_1)\,dq'_1\right]
 \delta(q_2-q'_2)\,dq'_2\right]\cdots\right]\\
 &\qquad{}\times\delta(q_k-q'_k)\,dq'_k\\
 &\quad=\int_{-\infty}^{\infty}\left[\cdots
 \left[\int_{-\infty}^{\infty}\varphi(q_1q'_2\cdots q'_k)
 \delta(q_2-q'_2)\,dq'_2\right]\cdots\right]
 \delta(q_k-q'_k)\,dq'_k=\cdots\\
 &\quad=\varphi(q_1q_2\cdots q_k),\\[1ex]
 &\int_\Omega\cdots\int\delta'(q_1-q'_1)\delta(q_2-q'_2)\cdots
 \delta(q_k-q'_k)\varphi(q'_1\cdots q'_k)\,dq'_1\cdots dq'_k\\
 &\quad=\frac{d}{dq_1}\int_\Omega\cdots\int
 \delta(q_1-q'_1)\delta(q_2-q'_2)\cdots\delta(q_k-q'_k)
 \varphi(q'_1\cdots q'_k)\,dq'_1\cdots dq'_k\\
 &\quad=\frac{d}{dq_1}\varphi(q_1\cdots q_k),\qquad\text{etc.}
\end{align*}
% Source PDF page 22 / original page 15.
\hypertarget{origpage-15}{}%

Thus the integral representation \hyperref[eq:I.3.I]{I} can be forced upon practically every operator.

Once this representation is available, the analogy between the problems \hyperref[eq:I.3.E1]{$E_1$} and \hyperref[eq:I.3.E3]{$E_3$} is complete: one need only make the replacements
\[
 \nu\mapsto q_1\cdots q_k,\qquad \nu'\mapsto q'_1\cdots q'_k,\qquad
 \sum_{\nu'}\mapsto\int_\Omega\cdots\int\cdots\,dq'_1\cdots dq'_k,\qquad x\mapsto\varphi.
\]
 Just as the vectors $x_\nu$ correspond to the functions $\varphi(q_1\cdots q_k)$, the matrices $h_{\nu\nu'}$ must be paired with the integral kernels $h(q_1\cdots q_k;q'_1\cdots q'_k)$. It is more convenient still to regard the integral kernels directly as matrices, calling the $q_1\cdots q_k$ their row indices and the $q'_1\cdots q'_k$ their column indices (corresponding to $\nu$ and $\nu'$, respectively). In addition to ordinary matrices $\{h_{\nu\nu'}\}$ with discrete sets of rows and columns (identified by the numbers $1,2,\ldots$), we then have others $\{h(q_1\cdots q_k;q'_1\cdots q'_k)\}$ (the integral kernels), for which the two sets are each identified by $k$ variables ranging continuously over all of $\Omega$.

This analogy may seem merely formal, but it is more than that. The indices $\nu$ and $\nu'$ can themselves be regarded as coordinates in a configuration space if they are interpreted as quantum numbers. In \textsc{Bohr}'s theory, these numbers label the possible orbits in phase space, which form a discrete set because the quantum conditions exclude the others.

We shall not pursue these ideas further here. \textsc{Dirac} and \textsc{Jordan} developed them into a unified theory of quantum processes, but ``improper'' objects such as $\delta(x),\delta'(x),\ldots$ play an essential part in that theory. These objects lie outside the framework of generally accepted mathematical methods, and we wish to describe quantum mechanics using those methods. We therefore turn to \textsc{Schrödinger}'s alternative way of unifying the two theories.

\section{Equivalence of the Two Theories: Hilbert Space}\label{sec:I.4}
The method sketched in \hyperref[sec:I.3]{I.3} amounted to setting up an analogy between the ``discrete'' space of index values, $Z=(1,2,\ldots)$, and the continuous configuration space $\Omega$ of the mechanical system ($\Omega$ is $k$-dimensional if $k$ is the number of classical mechanical degrees of freedom). It is no wonder that this cannot succeed without doing some violence to formalism and mathematics: the spaces $Z$ and $\Omega$ really are very different, and any attempt to relate them must encounter great difficulties.\origNote{33}

What actually matters to us, however, is not a relation between $Z$ and $\Omega$ at all, but only one between their respective functions: that is, between the sequences $x_1,x_2,\ldots$, which are functions on $Z$, and the wave functions $\varphi(q_1\cdots q_k)$, which are functions on $\Omega$.
% Source PDF page 23 / original page 16.
\hypertarget{origpage-16}{}%
For these alone enter into the problems of quantum mechanics.

In \textsc{Schrödinger}'s theory, the integral
\[
 \int_\Omega\cdots\int|\varphi(q_1\cdots q_k)|^2\,dq_1\cdots dq_k
\]
played an important role: it had to equal $1$ for $\varphi$ to be used in making physical statements (cf.\ \hyperref[sec:I.2]{I.2}). In matrix theory, on the other hand (cf.\ problem \hyperref[eq:I.3.E1]{$E_1$} in \hyperref[sec:I.3]{I.3}), the vector $x_1,x_2,\ldots$ plays a decisive role, and, in accordance with \textsc{Hilbert}'s theory of such eigenvalue problems (cf.\ the reference in \hyperlink{note-30}{Note~30}), it is always required to satisfy $\sum_\nu|x_\nu|^2<\infty$. Indeed, it is customary, if only to exclude the trivial solution $x_\nu\equiv0$, to impose the normalization $\sum_\nu|x_\nu|^2=1$. Thus, in $Z$ and $\Omega$, respectively, it is natural to restrict the class of admissible functions to those for which
\[
 \sum_\nu|x_\nu|^2
 \quad\text{and}\quad
 \int_\Omega\cdots\int|\varphi(q_1\cdots q_k)|^2\,dq_1\cdots dq_k
\]
are finite; only for such functions can multiplication by a constant make the indicated sum or integral equal to $1$---that is, normalize a solution in the customary sense.\origNote{34} We denote these classes of functions by $F_Z$ and $F_\Omega$, respectively.

We now have the \textbf{theorem}: $F_Z$ and $F_\Omega$ are isomorphic (\textsc{Fischer} and F.~\textsc{Riesz}\origNote{35}). This means precisely the following. It is possible to establish a one-to-one correspondence between $F_Z$ and $F_\Omega$, that is, to associate with every sequence $x_1,x_2,\ldots$ having finite $\sum_\nu|x_\nu|^2$ a function $\varphi(q_1\cdots q_k)$ for which
\[\int_\Omega\cdots\int|\varphi(q_1\cdots q_k)|^2\,dq_1\cdots dq_k\]
is finite, and conversely, in such a way that the correspondence is linear and preserves length. Here ``linearity'' means that, if $x_1,x_2,\ldots$ and $y_1,y_2,\ldots$ correspond to $\varphi(q_1\cdots q_k)$ and $\psi(q_1\cdots q_k)$, respectively, then $ax_1,ax_2,\ldots$ and $x_1+y_1,x_2+y_2,\ldots$ correspond to $a\varphi(q_1\cdots q_k)$ and $\varphi(q_1\cdots q_k)+\psi(q_1\cdots q_k)$, respectively. ``Preservation of length,'' for corresponding $x_1,x_2,\ldots$ and $\varphi(q_1\cdots q_k)$, means
\[
 \sum_\nu|x_\nu|^2
 =\int_\Omega\cdots\int|\varphi(q_1\cdots q_k)|^2\,dq_1\cdots dq_k.
\]
(The expression ``preservation of length'' comes from the custom of regarding $x_1,x_2,\ldots$ and $\varphi(q_1\cdots q_k)$ as vectors, and
\[
 \sqrt{\sum_\nu|x_\nu|^2}
 \quad\text{and}\quad
 \sqrt{\int_\Omega\cdots\int|\varphi(q_1\cdots q_k)|^2\,dq_1\cdots dq_k}
\]
as their ``lengths.'') Indeed, something more holds: if $x_1,x_2,\ldots$ and $y_1,y_2,\ldots$ correspond to $\varphi(q_1\cdots q_k)$ and $\psi(q_1\cdots q_k)$, respectively, then
\[
 \sum_\nu x_\nu\overline{y_\nu}
 =\int_\Omega\cdots\int\varphi(q_1\cdots q_k)
 \overline{\psi(q_1\cdots q_k)}\,dq_1\cdots dq_k
\]
% Source PDF page 24 / original page 17.
\hypertarget{origpage-17}{}%
(and both sides are absolutely convergent). In this connection, one might actually wish for
\[
 \sum_\nu x_\nu=\int_\Omega\cdots\int\varphi(q_1\cdots q_k)\,dq_1\cdots dq_k,
\]
or something similar, that is, a complete analogy between summation on the one hand and integration on the other. Closer inspection shows, however, that in quantum mechanics summation $\sum_\nu$ and integration $\int_\Omega\cdots\int\cdots\,dq_1\cdots dq_k$ are always applied only to expressions $x_\nu\overline{y_\nu}$ and $\varphi(q_1\cdots q_k)\overline{\psi(q_1\cdots q_k)}$, respectively.

We shall leave the construction of this correspondence to the next chapter, where we shall examine it in detail. But its significance should be emphasized. The spaces $Z$ and $\Omega$ are very different; establishing a relation between them must lead to insoluble mathematical difficulties. By contrast, $F_Z$ and $F_\Omega$ are isomorphic, that is, identical in their internal structure (they realize the same abstract properties in different mathematical objects). Since they, and not $Z$ and $\Omega$ themselves, are the actual analytic basis of matrix and wave theory, respectively, this isomorphism means that the two theories must always yield the same numerical results. More precisely, this holds if the isomorphism associates the matrix
\[
 \overline{\mathsf{H}}=H(\overline Q_1,\ldots,\overline Q_k;\overline P_1,\ldots,\overline P_k)
\]
with the operator
\[
 H=H\left(q_1,\ldots,q_k;\frac{h}{2\pi i}\frac{\partial}{\partial q_1},
 \ldots,\frac{h}{2\pi i}\frac{\partial}{\partial q_k}\right).
\]
Since the two arise by the same algebraic operations from the matrices $\overline Q_l,\overline P_l$ ($l=1,\ldots,k$) and the functional operators $q_l\cdots$, $\frac{h}{2\pi i}\frac{\partial}{\partial q_l}\cdots$ ($l=1,\ldots,k$), respectively, it suffices to show that $q_l\cdots$ corresponds to the matrix $\overline Q_l$ and $\frac{h}{2\pi i}\frac{\partial}{\partial q_l}\cdots$ to the matrix $\overline P_l$. Now nothing further was required of the $\overline Q_l,\overline P_l$ ($l=1,\ldots,k$) than that they satisfy the ``commutation relations'' stated in \hyperref[sec:I.2]{I.2}:\tnote{In the second line of the following display, the source reverses the order of the position and momentum factors relative to the commutation relation on original p.~6, without changing the sign on the right. The display here retains the source's order and sign.}
\begin{align*}
 \overline Q_m\overline Q_n-\overline Q_n\overline Q_m&=0,
 &\overline P_m\overline P_n-\overline P_n\overline P_m&=0,\\
 \overline Q_m\overline P_n-\overline P_n\overline Q_m
 &=\begin{cases}0,&m\ne n,\\\dfrac{h}{2\pi i}\,1,&m=n,\end{cases}
 &&(m,n=1,2,\ldots).
\end{align*}
The matrices corresponding to $q_l\cdots$, $\frac{h}{2\pi i}\frac{\partial}{\partial q_l}\cdots$ will certainly do this, for the functional operators $q_l\cdots$, $\frac{h}{2\pi i}\frac{\partial}{\partial q_l}\cdots$ themselves possess the stated properties,\origNote{36} and these are not lost in the isomorphic transfer to $F_Z$.
% Source PDF page 25 / original page 18; preceding sentence crosses the page boundary.
\hypertarget{origpage-18}{}%

The systems $F_Z$ and $F_\Omega$ are isomorphic, and the quantum theories built on them are mathematically equivalent. We may therefore expect to obtain a unified construction by taking as our starting point the intrinsic properties common to these two classes of functions. Such a construction would be independent of the incidental features of the particular formal framework chosen and would display only the essential features of quantum mechanics.

$F_Z$ is generally called ``Hilbert space.'' Our first concern will therefore be to find the intrinsic properties of Hilbert space, independent of its particular realization as $F_Z$ or $F_\Omega$. The mathematical object described by these properties (which, in a concrete calculation, is to be identified with $F_Z$ or $F_\Omega$, but which is more convenient to handle for general purposes) will be called ``abstract Hilbert space.''

We therefore propose to describe abstract Hilbert space and then to prove the following points with full rigor:
\begin{enumerate}
 \item Abstract Hilbert space (H.S.\ for short) is uniquely characterized by the properties to be stated; that is, it admits no essentially different interpretations.
 \item Both $F_Z$ and $F_\Omega$ have these properties.
\end{enumerate}
(This will rigorously establish the matters explained only qualitatively in \hyperref[sec:I.4]{I.4}.) Once this has been done, we shall use the mathematical apparatus thus obtained to construct quantum mechanics.


% Consolidated source: chapter-II.tex
% Source PDF page 25 / original page 18 (chapter begins mid-page).
\chapter{General Considerations on Abstract Hilbert Space}\label{chap:II}
\section{Characterization of Hilbert Space}\label{sec:II.1}

We must carry out the program formulated at the end of \hyperref[sec:I.4]{I.4}: to characterize Hilbert space, which provides the mathematical basis for treating quantum mechanics, in such a way that we use no concepts other than those subsequently required in quantum mechanics. Accordingly, these concepts must make just as much sense in the ``discrete'' function space $F_Z$ of sequences $x_\nu$ ($\nu=1,2,\ldots$) as in the ``continuous'' space $F_\Omega$ of wave functions $\varphi(q_1\cdots q_k)$, where $q_1,\ldots,q_k$ range over the configuration space $\Omega$. As we have already indicated, these concepts are the following:

% Source PDF page 26 / original page 19.
\hypertarget{origpage-19}{}%
$\alpha)$ ``Scalar multiplication,'' that is, multiplication of an element $f$ of Hilbert space by a (complex) number $a$: $af$. In $F_Z$ this turns $x_\nu$ into $ax_\nu$; in $F_\Omega$, it turns $\varphi(q_1\cdots q_k)$ into $a\varphi(q_1\cdots q_k)$.

$\beta)$ Addition and subtraction of two elements $f,g$ of Hilbert space: $f\pm g$. In $F_Z$, $x_\nu$ and $y_\nu$ thus yield $x_\nu\pm y_\nu$; in $F_\Omega$, $\varphi(q_1\cdots q_k)$ and $\psi(q_1\cdots q_k)$ yield $\varphi(q_1\cdots q_k)\pm\psi(q_1\cdots q_k)$.

$\gamma)$ ``Inner multiplication'' of two elements $f,g$ of Hilbert space, which, unlike $\alpha)$ and $\beta)$, yields not an element of Hilbert space but a (complex) number: $(f,g)$. In $F_Z$, $x_\nu$ and $y_\nu$ thus yield $\sum_\nu x_\nu\overline{y_\nu}$; in $F_\Omega$, $\varphi(q_1\cdots q_k)$ and $\psi(q_1\cdots q_k)$ yield
\[
\int\!\cdots\!\int_\Omega
\varphi(q_1\cdots q_k)\overline{\psi(q_1\cdots q_k)}\,dq_1\cdots dq_k.
\]
The definitions in $F_Z$ and $F_\Omega$ still require the corresponding convergence proofs. We shall supply these in \hyperref[sec:II.3]{II.3}.

In what follows we shall consistently denote points of Hilbert space by $f,g,\ldots,\varphi,\psi,\ldots$, complex numbers by $a,b,\ldots,x,y,\ldots$, and positive integers by $k,l,m,\ldots,\mu,\nu,\ldots$. Where necessary we shall also call Hilbert space $\mathfrak{R}_\infty$, as an abbreviation for ``infinite-dimensional Euclidean space,'' analogous to the usual notation $\mathfrak{R}_n$ for ``$n$-dimensional Euclidean space'' ($n=1,2,\ldots$).

What is noteworthy about the operations $af$, $f\pm g$, and $(f,g)$ is that they are precisely the fundamental operations of vector calculus: for example, those on which calculations of lengths and angles in Euclidean geometry, or of force and work in particle mechanics, are based. The analogy is clearest for $F_Z$ when, instead of the $x_1,x_2,\ldots$ in $\mathfrak{R}_\infty$, one considers the ordinary points $x_1,\ldots,x_n$ of an $\mathfrak{R}_n$, for which the operations $\alpha)$, $\beta)$, and $\gamma)$ can be performed in just the same way. In particular, for $n=3$ we obtain the situation in ordinary space. Under some circumstances it is more convenient to regard $x_1,\ldots,x_n$ not as points but as vectors, say from the point $0,\ldots,0$ to the point $x_1,\ldots,x_n$.

To characterize abstract Hilbert space, we therefore take the fundamental vector operations $af$, $f\pm g$, and $(f,g)$ as our basis. As the following discussion will show, we shall include all the $\mathfrak{R}_n$ along with $\mathfrak{R}_\infty$. Accordingly, where we do not yet wish to distinguish between $\mathfrak{R}_\infty$ and the $\mathfrak{R}_n$, we use the neutral designation $\mathfrak{R}$ for the space.

First we postulate the characteristic vector properties for $\mathfrak{R}$:\origNote{37}

\leavevmode\phantomsection\textbf{A.}\label{ax:II.A} $\mathfrak{R}$ is a linear space.

That is, addition $f+g$ and ``scalar'' multiplication $af$ are defined in $\mathfrak{R}$ ($f,g$ are elements of $\mathfrak{H}$, $a$ a complex number; $f+g,af$ belong to $\mathfrak{H}$), and it has an element $0$.\origNote{38} They obey the familiar rules of vector algebra:

% Source PDF page 27 / original page 20.
\hypertarget{origpage-20}{}%
\begin{alignat*}{2}
f+g&=g+f &\quad &\text{(commutativity of addition)},\\
(f+g)+h&=f+(g+h)&\quad &\text{(associativity of addition)},\\
a(f+g)&=af+ag,\quad (a+b)f=af+bf
&\quad &\text{(distributivity of multiplication)},\\
(ab)f&=a(bf)&\quad &\text{(associativity of multiplication)},\\
0f&=0,\quad 1f=1&\quad &\text{(the roles of zero and unity)}.
\end{alignat*}
The rules of calculation not mentioned here follow readily from these postulates.\tnote{The source prints $1f=1$ in the preceding list. The identity required here, and used in the next calculation, is $1f=f$. The printed formula is retained.} For example, the role of zero in addition:
\[
f+0=1\cdot f+0\cdot f=(1+0)\cdot f=1\cdot f=f.
\]
Subtraction is also uniquely determined: define
\[
-f=(-1)\cdot f,\qquad f-g=f+(-g).
\]
Then
\begin{align*}
(f-g)+g&=(f+(-g))+g=f+((-g)+g),\\
(f+g)-g&=(f+g)+(-g)=f+(g+(-g)),
\end{align*}
and both expressions are
\begin{align*}
f+((-1)\cdot g+1\cdot g)
&=f+((-1)+1)\cdot g\\
&=f+0\cdot g=f+0=f.
\end{align*}
The distributive laws also hold for subtraction:
\begin{align*}
a\cdot(f-g)&=a\cdot f+a\cdot(-g)=af+a\cdot((-1)\cdot g)\\
&=af+(a\cdot(-1))\cdot g=af+((-1)\cdot a)\cdot g\\
&=af+(-ag)=af-ag,\\
(a-b)\cdot f&=a\cdot f+(-b)\cdot f=af+((-1)b)\cdot f\\
&=af+(-bf)=af-bf.
\end{align*}
There is no need to pursue these matters further: all the rules of linear vector calculus plainly remain valid here.

As with vectors, we can therefore define when certain elements $f_1,\ldots,f_k$ of $\mathfrak{H}$ are linearly independent:

\leavevmode\phantomsection\textbf{Definition 1.}\label{def:II.1} $f_1,\ldots,f_k$ are linearly independent if $a_1f_1+\cdots+a_kf_k=0$, with complex numbers $a_1,\ldots,a_k$, implies $a_1=\cdots=a_k=0$.

We next define the analogue of the linear figures of vector calculus (a line, plane, and so forth through the origin): the linear manifold.

\leavevmode\phantomsection\textbf{Definition 2.}\label{def:II.2} A subset $\mathfrak{M}$ of $\mathfrak{R}$ is called a linear manifold if, together with any $k$ ($=1,2,\ldots$) of its elements $f_1,\ldots,f_k$, it also contains their linear combinations $a_1f_1+\cdots+a_kf_k$.\origNote{39} If $\mathfrak{A}$ is any subset of $\mathfrak{R}$, the set of all $a_1f_1+\cdots+a_kf_k$ ($k=1,2,\ldots$; $a_1,\ldots,a_k$ arbitrary complex numbers; $f_1,\ldots,f_k$ arbitrary elements of $\mathfrak{A}$) is a linear manifold that evidently contains $\mathfrak{A}$. It is clearly a subset of every other linear manifold containing $\mathfrak{A}$.
% Source PDF page 28 / original page 21.
\hypertarget{origpage-21}{}%
It shall be called ``the linear manifold spanned by $\mathfrak{A}$,'' and denoted by $\{\mathfrak{A}\}$.

Before developing these concepts further, we formulate the next fundamental principle of vector calculus: the existence of the inner product.

\leavevmode\phantomsection\textbf{B.}\label{ax:II.B} A Hermitian inner product is defined in $\mathfrak{R}$.

That is, $(f,g)$ is defined ($f,g$ belong to $\mathfrak{H}$, and $(f,g)$ is a complex number), with the following properties:\origNote{40}
\begin{alignat*}{2}
(f'+f'',g)&=(f',g)+(f'',g)&\quad &\text{(distributivity in the first factor)},\\
(a\cdot f,g)&=a\cdot(f,g)&\quad &\text{(associativity in the first factor)},\\
(f,g)&=\overline{(g,f)}&\quad &\text{(Hermitian symmetry)},\\
(f,f)&\geq0,\quad (f,f)=0\text{ only for }f=0
&\quad &\text{(definiteness)}.
\end{alignat*}
By Hermitian symmetry, the two properties of the first factor imply corresponding properties of the second: interchange $f$ and $g$ and take the complex conjugates of both sides. Thus
\[
(f,g'+g'')=(f,g')+(f,g''),\qquad
(f,a\cdot g)=\bar a\cdot(f,g).
\]
This inner product is of great importance because it permits the definition of distance. In Euclidean space the magnitude of a vector $f$ is, as is well known, defined by $\|f\|=\sqrt{(f,f)}$,\origNote{41} and the distance between two points $f,g$ by $\|f-g\|$. We shall follow this precedent.

\leavevmode\phantomsection\textbf{Definition 3.}\label{def:II.3} The ``norm'' of an element $f$ of $\mathfrak{R}$ is $\|f\|=\sqrt{(f,f)}$; the ``distance'' between two elements $f,g$ is $\|f-g\|$.\origNote{42}

We shall now show that this definition has all the properties required of a distance. To do so, we first prove:

\leavevmode\phantomsection\textbf{Theorem 1.}\label{thm:II.1} The inequality $|(f,g)|\leq\|f\|\cdot\|g\|$ always holds.

\emph{Proof.} First, if $z=u+iv$ is a complex number, with $u,v$ real, $\operatorname{Re}z$ and $\operatorname{Im}z$ denote its real and imaginary parts, respectively: $\operatorname{Re}z=u$, $\operatorname{Im}z=v$. We have
\begin{align*}
\|f\|^2+\|g\|^2-2\operatorname{Re}(f,g)
&=(f,f)+(g,g)-(f,g)-(g,f)\\
&=(f-g,f-g)\geq0,\\
\operatorname{Re}(f,g)&\leq\tfrac12(\|f\|^2+\|g\|^2).
\end{align*}
Replace $f,g$ by $af,\frac1a g$, where $a$ is real and positive. The left side is readily seen to remain unchanged, whereas the right becomes $\frac12(a^2\|f\|^2+\frac1{a^2}\|g\|^2)$. Since this is at least $\operatorname{Re}(f,g)$, the same is true of its minimum, which is $\|f\|\cdot\|g\|$. For $f,g\ne0$ this minimum is attained at $a=\sqrt{\frac{\|g\|}{\|f\|}}$; for $f=0$ or $g=0$ it is approached as $a\to+\infty$ or $a\to+0$, respectively.
% Source PDF page 29 / original page 22.
\hypertarget{origpage-22}{}%
Thus
\[
\operatorname{Re}(f,g)\leq\|f\|\cdot\|g\|.
\]
Replace $f,g$ here by $e^{i\alpha}f,g$, with $\alpha$ real. The right side remains unchanged, since
\[
(af,af)=a\bar a(f,f)=|a|^2(f,f),\qquad
\|af\|=|a|\cdot\|f\|,
\]
so that $\|af\|=\|f\|$ for $|a|=1$. The left side becomes
\[
\operatorname{Re}(e^{i\alpha}(f,g))
=\cos\alpha\operatorname{Re}(f,g)-\sin\alpha\operatorname{Im}(f,g).
\]
Its maximum is evidently
\[
\sqrt{(\operatorname{Re}(f,g))^2+(\operatorname{Im}(f,g))^2}=|(f,g)|,
\]
which gives the assertion:
\[
|(f,g)|\leq\|f\|\cdot\|g\|.
\]

\textbf{Corollary.} Equality can hold only if $f,g$ differ by a constant complex factor.

\emph{Proof.} Equality in $\operatorname{Re}(f,g)\leq\frac12(\|f\|^2+\|g\|^2)$ requires $(f-g,f-g)=0$, hence $f=g$. In passing from this relation to $|(f,g)|\leq\|f\|\cdot\|g\|$, we replace $f,g$ by $e^{i\alpha}af,\frac1a g$ ($a,\alpha$ real, $a>0$), unless $f$ or $g$ is zero. Equality therefore requires $e^{i\alpha}af=\frac1a g$, or $g=a^2e^{i\alpha}f=cf$, with $c\ne0$. Conversely, equality clearly holds if $f$ or $g$ is zero, or if $g=cf$ ($c\ne0$).

\leavevmode\phantomsection\textbf{Theorem 2.}\label{thm:II.2} We always have $\|f\|\geq0$, with equality only for $f=0$. Also, $\|a\cdot f\|=|a|\cdot\|f\|$. The inequality $\|f+g\|\leq\|f\|+\|g\|$ always holds, with equality only when $f,g$ differ by a constant, real, nonnegative factor.

\emph{Proof.} We have already established the first two assertions above. The inequality in the third is proved as follows:
\begin{align*}
(f+g,f+g)&=(f,f)+(g,g)+(f,g)+(g,f)\\
&=\|f\|^2+\|g\|^2+2R(f,g)\\
&\leq\|f\|^2+\|g\|^2+2\|f\|\cdot\|g\|
=(\|f\|+\|g\|)^2,\\
\|f+g\|&\leq\|f\|+\|g\|.
\end{align*}
Equality requires $\operatorname{Re}(f,g)=\|f\|\cdot\|g\|$. By the considerations in the preceding corollary, this implies that $f$ or $g$ is zero, or that $g=a^2f=cf$, with $c$ real and positive. Equality clearly does hold in this case.

It now follows immediately from \hyperref[thm:II.2]{Theorem 2} that the distance $\|f-g\|$ has these properties: $f,g$ have distance zero when $f=g$, and never otherwise. The distance between $g,f$ is the same as between $f,g$. The distance between $f,h$ is no greater than the sum of the distances between $f,g$ and $g,h$. Equality holds only when $g=af+(1-a)h$, with $a$ real and $0\leq a\leq1$.\origNote{43} The distance between $af,ag$ is $|a|$ times the distance between $f,g$.

% Source PDF page 30 / original page 23.
\hypertarget{origpage-23}{}%
But these are precisely the properties of distance that, in geometry and topology, allow the concepts of continuity, boundedness, limit point, and so forth to be reduced to the concept of distance. We therefore make the following definitions:

A function $F(f)$ defined on $\mathfrak{R}$ takes its values either in $\mathfrak{R}$ or in the complex numbers. It is continuous at $f_0$ in $\mathfrak{R}$ if for every $\varepsilon>0$ there is a $\delta>0$ such that $\|f-f_0\|<\delta$ implies $\|F(f)-F(f_0)\|<\varepsilon$ or $|F(f)-F(f_0)|<\varepsilon$, according as the values of $F$ are points of $\mathfrak{R}$ or numbers. The same function is called bounded in $\mathfrak{R}$, or in a given subset of $\mathfrak{R}$, if throughout that set $\|F(f)\|\leq C$ or $|F(f)|\leq C$, where $C$ is a suitably chosen but fixed constant. Analogous definitions apply to several variables. A sequence $f_1,f_2,\ldots$ converges to $f$, or has limit $f$, if the numbers $\|f_1-f\|,\|f_2-f\|,\ldots$ converge to zero. A point is a limit point of a set $\mathfrak{A}$, a subset of $\mathfrak{R}$, if it is the limit of a sequence from $\mathfrak{A}$.\origNote{44} In particular, $\mathfrak{A}$ is called closed if all its limit points belong to it, and everywhere dense if its limit points comprise all of $\mathfrak{R}$.

We also observe that $af$, $f+g$, and $(f,g)$ are continuous in all their variables. The first two assertions are clear from
\[
\|af-af'\|=|a|\cdot\|f-f'\|,
\]
\[
\|(f+g)-(f'+g')\|
=\|(f-f')+(g-g')\|\leq\|f-f'\|+\|g-g'\|.
\]
Further, if $\|f-f'\|<\varepsilon$ and $\|g-g'\|<\varepsilon$, put $f'-f=\varphi$, $g'-g=\psi$. Then
\begin{align*}
|(f,g)-(f',g')|
&=|(f,g)-(f+\varphi,g+\psi)|\\
&=|(\varphi,g)+(f,\psi)+(\varphi,\psi)|\\
&\leq|(\varphi,g)|+|(f,\psi)|+|(\varphi,\psi)|\\
&\leq\|\varphi\|\cdot\|g\|+\|f\|\cdot\|\psi\|
+\|\varphi\|\cdot\|\psi\|\\
&\leq\varepsilon(\|f\|+\|g\|+\varepsilon).
\end{align*}
As $\varepsilon\to0$, this tends to zero, and can therefore be made smaller than any $\delta>0$.

As we see, properties \hyperref[ax:II.A]{A} and \hyperref[ax:II.B]{B} enable us to say a good deal about $\mathfrak{R}$, but not enough to distinguish the $\mathfrak{R}_n$ from one another and from $\mathfrak{R}$: the number of dimensions has not yet been mentioned. This is, however, known to be related to the maximum number of linearly independent vectors. If such a maximum $n=0,1,2,\ldots$ exists, the following holds for that $n$:

\leavevmode\phantomsection\textbf{$C^{(n)}$.}\label{ax:II.Cn} There are exactly $n$ linearly independent vectors.

% Source PDF page 31 / original page 24.
\hypertarget{origpage-24}{}%
That is, $n$ such vectors can be specified, but $n+1$ do not exist.

If there is no maximum, on the other hand, we have:

\leavevmode\phantomsection\textbf{$C^{(\infty)}$.}\label{ax:II.Cinfty} There are arbitrarily many linearly independent vectors.

That is, for every $k=1,2,\ldots$, one can specify $k$ such vectors.

Thus \hyperref[ax:II.Cn]{C} is not really a new postulate: if \hyperref[ax:II.A]{A} and \hyperref[ax:II.B]{B} hold, some \hyperref[ax:II.Cn]{$C^{(n)}$} or \hyperref[ax:II.Cinfty]{$C^{(\infty)}$} must hold. Depending on which we choose, we obtain a different space $\mathfrak{R}$. We shall see that \hyperref[ax:II.Cn]{$C^{(n)}$} implies that $\mathfrak{R}$ has all the properties of $n$-dimensional complex Euclidean space. In contrast, \hyperref[ax:II.Cinfty]{$C^{(\infty)}$} does not yet suffice to ensure that $\mathfrak{R}$ is essentially identical with Hilbert space $\mathfrak{R}_\infty$; two further postulates \hyperref[ax:II.D]{D} and \hyperref[ax:II.E]{E} are needed. More precisely, we shall show that an $\mathfrak{R}$ satisfying \hyperref[ax:II.A]{A}, \hyperref[ax:II.B]{B}, and \hyperref[ax:II.Cn]{$C^{(n)}$} has all the properties of $\mathfrak{R}_n$, in particular \hyperref[ax:II.D]{D} and \hyperref[ax:II.E]{E}, which we shall state presently and which therefore follow from \hyperref[ax:II.A]{A}, \hyperref[ax:II.B]{B}, and \hyperref[ax:II.Cn]{$C^{(n)}$}. We shall further show that an $\mathfrak{R}$ satisfying \hyperref[ax:II.A]{A}, \hyperref[ax:II.B]{B}, \hyperref[ax:II.Cinfty]{$C^{(\infty)}$}, \hyperref[ax:II.D]{D}, and \hyperref[ax:II.E]{E} has all the properties of $\mathfrak{R}_\infty$, but that \hyperref[ax:II.D]{D} and \hyperref[ax:II.E]{E} are essential: they do not follow from \hyperref[ax:II.A]{A}, \hyperref[ax:II.B]{B}, and \hyperref[ax:II.Cinfty]{$C^{(\infty)}$}. We now formulate \hyperref[ax:II.D]{D} and \hyperref[ax:II.E]{E}; the proof that all $\mathfrak{R}_n$ and $\mathfrak{R}_\infty$ also possess these properties will be given later, in \hyperref[sec:II.3]{II.3}.

\leavevmode\phantomsection\textbf{D.}\label{ax:II.D} $\mathfrak{R}$ is complete.\origNote{45}

That is, if a sequence $f_1,f_2,\ldots$ in $\mathfrak{R}$ satisfies Cauchy's convergence condition---for every $\varepsilon>0$ there is an $N=N(\varepsilon)$ such that $m,n\geq N$ implies $\|f_m-f_n\|<\varepsilon$---then it converges: it has a limit $f$, in the sense defined above.

\leavevmode\phantomsection\textbf{E.}\label{ax:II.E} $\mathfrak{R}$ is separable.\origNote{45}

That is, there is a sequence $f_1,f_2,\ldots$ in $\mathfrak{R}$ that is everywhere dense in $\mathfrak{R}$.

In \hyperref[sec:II.2]{II.2}, as promised, we shall use this foundation to develop the ``geometry'' of $\mathfrak{R}$ and see that it agrees with that of $\mathfrak{R}_n$ or $\mathfrak{R}_\infty$.

\section{Geometry of Hilbert Space}\label{sec:II.2}

We begin with two definitions. The first contains just as much of the geometric theory of angles as our purposes require: the concept of a right angle, or orthogonality.

\leavevmode\phantomsection\textbf{Definition 4.}\label{def:II.4} Two elements $f,g$ of $\mathfrak{R}$ are orthogonal if $(f,g)=0$. Two linear manifolds $\mathfrak{M},\mathfrak{N}$ are orthogonal if every element of $\mathfrak{M}$ is orthogonal to every element of $\mathfrak{N}$. A set $\mathfrak{O}$ is called an orthonormal system if, for all $f,g$ in $\mathfrak{O}$,
\[
(f,g)=\begin{cases}1,&f=g,\\0,&f\ne g.\end{cases}
\]
% Source PDF page 32 / original page 25.
\hypertarget{origpage-25}{}%
That is, any two distinct elements are orthogonal, and each element has norm $1$.\origNote{46} In particular, $\mathfrak{O}$ is called complete if it cannot be a subset of another orthonormal system containing additional elements.\origNote{47}

We also note that completeness of the orthonormal system $\mathfrak{O}$ evidently means that no $f$ of norm $1$ exists that is orthogonal to all of $\mathfrak{O}$; see \hyperlink{note-46}{note 46}. But if merely $f\ne0$ and $f$ were orthogonal to all of $\mathfrak{O}$, then $f'=\frac1{\|f\|}\cdot f$ would meet all these conditions, since $\|f\|>0$: $\|f'\|=\frac1{\|f\|}\cdot\|f\|=1$, and $f'$ is orthogonal to $\mathfrak{O}$. Thus completeness of $\mathfrak{O}$ means that every $f$ orthogonal to all of $\mathfrak{O}$ must vanish.

The second definition plays an essential role only in $\mathfrak{R}_\infty$: in $\mathfrak{R}_n$, every linear manifold already has the property it defines; see the end of \hyperref[sec:II.3]{II.3}. We therefore cannot give an intuitive geometric picture of it.

\leavevmode\phantomsection\textbf{Definition 5.}\label{def:II.5} A linear manifold that is also closed is called a closed linear manifold. If $\mathfrak{A}$ is any set in $\mathfrak{R}$ and we adjoin to $\{\mathfrak{A}\}$, the linear manifold spanned by $\mathfrak{A}$, all its limit points, we obtain a closed linear manifold containing $\mathfrak{A}$. It is a subset of every other closed linear manifold containing $\mathfrak{A}$.\origNote{48} We call it the closed linear manifold spanned by $\mathfrak{A}$, and denote it by $[\mathfrak{A}]$.

We now turn to a closer analysis of $\mathfrak{R}$, especially of its complete orthonormal systems. To theorems that assume, in addition to \hyperref[ax:II.A]{A} and \hyperref[ax:II.B]{B}, either \hyperref[ax:II.Cn]{$C^{(n)}$} or all three conditions \hyperref[ax:II.Cinfty]{$C^{(\infty)}$}, \hyperref[ax:II.D]{D}, and \hyperref[ax:II.E]{E}, we append the superscript $(n)$ or $(\infty)$, respectively. Theorems valid in both cases receive no superscript.

\leavevmode\phantomsection\textbf{Theorem $3^{(n)}$.}\label{thm:II.3n} Every orthonormal system has at most $n$ elements; it is complete if and only if it has $n$ elements.

\emph{Remark.} The first assertion implies that there is a maximum number of elements for orthonormal systems. Systems attaining this number are complete by definition. Thus in case \hyperref[ax:II.Cn]{$C^{(n)}$} complete orthogonal systems exist; according to the theorem they have $n$ elements, $\varphi_1,\varphi_2,\ldots,\varphi_n$.

\emph{Proof.} Every finite orthonormal system is linearly independent. Let it be $\varphi_1,\varphi_2,\ldots,\varphi_m$. From
\[
a_1\varphi_1+\cdots+a_m\varphi_m=0
\]
taking the inner product with $\varphi_\mu$ ($\mu=1,2,\ldots,m$) gives $a_\mu=0$. By \hyperref[ax:II.Cn]{$C^{(n)}$}, therefore, it cannot have $n+1$ elements. Consequently, an arbitrary orthonormal system cannot have a subset with $n+1$ elements. It is therefore finite and has at most $n$ elements.

% Source PDF page 33 / original page 26.
\hypertarget{origpage-26}{}%
One with $n$ elements thus admits no extension and is complete. But one with $m<n$ elements, $\varphi_1,\varphi_2,\ldots,\varphi_m$, is not complete. Among the linear combinations $a_1\varphi_1+\cdots+a_m\varphi_m$, no $n>m$ are linearly independent. By \hyperref[ax:II.Cn]{$C^{(n)}$} there are therefore elements $f$ distinct from all these combinations, so that
\[
\psi=f-a_1\varphi_1-\cdots-a_m\varphi_m
\]
is always nonzero. Now $(\psi,\varphi_\mu)=0$ means $a_\mu=(f,\varphi_\mu)$ ($\mu=1,2,\ldots,m$), and can thus be achieved for all $\mu$. Our system is therefore incomplete.

\leavevmode\phantomsection\textbf{Theorem $3^{(\infty)}$.}\label{thm:II.3infty} Every orthonormal system is finite or a countably infinite sequence; if complete, it is necessarily infinite.

\emph{Remark.} We may therefore write all orthonormal systems as sequences $\varphi_1,\varphi_2,\ldots$, possibly terminating and hence finite, and shall do so. Here, having infinitely many elements is necessary for completeness but is not sufficient,\origNote{49} unlike the corresponding condition under \hyperref[ax:II.Cn]{$C^{(n)}$}.

\emph{Proof.} Let $\mathfrak{O}$ be an orthonormal system and $f,g$ two distinct elements. Then
\[
(f-g,f-g)=(f,f)+(g,g)-(f,g)-(g,f)=2,
\qquad \|f-g\|=\sqrt2.
\]
Let $f_1,f_2,\ldots$ be the everywhere dense sequence in $\mathfrak{R}$ provided by \hyperref[ax:II.E]{E}. For every $f$ in $\mathfrak{O}$ there is an $f_m$ of this sequence with $\|f-f_m\|<\frac12\sqrt2$. The corresponding $f_m,f_n$ for $f,g$ must be distinct, since $f_m=f_n$ would imply
\begin{align*}
\|f-g\|&=\|(f-f_m)-(g-f_m)\|\\
&\leq\|f-f_m\|+\|g-f_m\|
<\tfrac12\sqrt2+\tfrac12\sqrt2=\sqrt2.
\end{align*}
Thus to every $f$ in $\mathfrak{O}$ corresponds an $f_m$ from the sequence $f_1,f_2,\ldots$, with distinct $f$ corresponding to distinct $f_m$. Hence $\mathfrak{O}$ is finite or a sequence.

As in the proof of \hyperref[thm:II.3n]{Theorem $3^{(n)}$}, one shows that if $\mathfrak{R}$ contains more than $m$ linearly independent elements, a system $\varphi_1,\varphi_2,\ldots,\varphi_m$ cannot be complete. Since \hyperref[ax:II.Cinfty]{$C^{(\infty)}$} makes this true for every $m$, a complete system must be infinite.

Insofar as the following theorems concern convergence, they apply only to \hyperref[ax:II.Cinfty]{$C^{(\infty)}$}. Nevertheless, their other assertions make it preferable to formulate them generally.

\leavevmode\phantomsection\textbf{Theorem 4.}\label{thm:II.4} Let $\varphi_1,\varphi_2,\ldots$ be an orthonormal system. Then every series $\sum_\nu(f,\varphi_\nu)\overline{(g,\varphi_\nu)}$ with infinitely many terms converges absolutely. In particular, taking $f=g$, we always have $\sum_\nu|(f,\varphi_\nu)|^2\leq|f|^2$.

\emph{Proof.} Put $a_\nu=(f,\varphi_\nu)$, $\nu=1,2,\ldots$. Then $f-\sum_{\nu=1}^Na_\nu\varphi_\nu=\psi$ is orthogonal to all $\varphi_\nu$, $\nu=1,2,\ldots,N$; see the proof of \hyperref[thm:II.3n]{Theorem $3^{(n)}$}.
% Source PDF page 34 / original page 27.
\hypertarget{origpage-27}{}%
Since $f=\sum_{\nu=1}^Na_\nu\varphi_\nu+\psi$, we have
\begin{align*}
(f,f)&=\sum_{\mu,\nu=1}^Na_\mu\bar a_\nu(\varphi_\mu,\varphi_\nu)
+\sum_{\nu=1}^Na_\nu(\varphi_\nu,\psi)\\
&\quad+\sum_{\nu=1}^N\bar a_\nu(\psi,\varphi_\nu)+(\psi,\psi)\\
&=\sum_{\nu=1}^N|a_\nu|^2+(\psi,\psi)
\geq\sum_{\nu=1}^N|a_\nu|^2.
\end{align*}
Thus $\sum_{\nu=1}^N|a_\nu|^2\leq\|f\|^2$. If the system $\varphi_1,\varphi_2,\ldots$ is finite, this immediately gives $\sum_\nu|a_\nu|^2\leq\|f\|^2$. If infinite, letting $N\to\infty$ establishes the absolute convergence of $\sum_\nu|a_\nu|^2$ and shows that its sum is at most $\|f\|^2$. This proves the second assertion.

Since
\[
|(f,\varphi_\nu)\overline{(g,\varphi_\nu)}|
\leq\tfrac12\bigl\{|(f,\varphi_\nu)|^2+|(g,\varphi_\nu)|^2\bigr\},
\]
our convergence result also gives the more general convergence statement of the first assertion.

\leavevmode\phantomsection\textbf{Theorem 5.}\label{thm:II.5} Let $\varphi_1,\varphi_2,\ldots$ be an infinite orthonormal system. Then $\sum_{\nu=1}^\infty x_\nu\varphi_\nu$ converges if and only if $\sum_{\nu=1}^\infty|x_\nu|^2$ does. The latter series has nonnegative real terms and therefore either converges or diverges to $+\infty$.

\emph{Proof.} The assertion concerns only \hyperref[ax:II.Cinfty]{$C^{(\infty)}$}, so we may use \hyperref[ax:II.D]{D}, Cauchy's convergence criterion. Accordingly, $\sum_{\nu=1}^\infty x_\nu\varphi_\nu$ converges---that is, the sequence of partial sums $\sum_{\nu=1}^Nx_\nu\varphi_\nu$ converges as $N\to\infty$---if and only if, for every $\varepsilon>0$, there is an $N=N(\varepsilon)$ such that
\[
\left\|\sum_{\nu=1}^Lx_\nu\varphi_\nu-
\sum_{\nu=1}^Mx_\nu\varphi_\nu\right\|<\varepsilon
\qquad(L,M\geq N).
\]
Assume $L>M\geq N$. Then
\[
\left\|\sum_{\nu=1}^Lx_\nu\varphi_\nu-
\sum_{\nu=1}^Mx_\nu\varphi_\nu\right\|
=\left\|\sum_{\nu=M+1}^Lx_\nu\varphi_\nu\right\|<\varepsilon,
\]
\begin{align*}
\left\|\sum_{\nu=M+1}^Lx_\nu\varphi_\nu\right\|^2
&=\left(\sum_{\mu=M+1}^Lx_\mu\varphi_\mu,
\sum_{\nu=M+1}^Lx_\nu\varphi_\nu\right)\\
&=\sum_{\mu,\nu=M+1}^Lx_\mu\bar x_\nu(\varphi_\mu,\varphi_\nu)
=\sum_{\nu=M+1}^L|x_\nu|^2\\
&=\sum_{\nu=1}^L|x_\nu|^2-\sum_{\nu=1}^M|x_\nu|^2,
\end{align*}
and therefore
\[
0\leq\sum_{\nu=1}^L|x_\nu|^2-\sum_{\nu=1}^M|x_\nu|^2<\varepsilon^2.
\]
But this is precisely Cauchy's convergence condition for the sequence $\sum_{\nu=1}^N|x_\nu|^2$, $N\to\infty$, hence for the series $\sum_{\nu=1}^\infty|x_\nu|^2$.

\textbf{Corollary.} For $f=\sum_\nu x_\nu\varphi_\nu$, we have $(f,\varphi_\nu)=x_\nu$, whether the orthogonal system is finite or infinite, with convergence of course assumed in the latter case.

% Source PDF page 35 / original page 28.
\hypertarget{origpage-28}{}%
\emph{Proof.} For $N\geq\nu$, the inner product $(\sum_{\mu=1}^Nx_\mu\varphi_\mu,\varphi_\nu)$ is
\[
\sum_{\mu=1}^Nx_\mu(\varphi_\mu,\varphi_\nu)=x_\nu.
\]
If $\varphi_1,\varphi_2,\ldots$ is finite, take $N$ equal to the largest index. If infinite, let $N\to\infty$, using continuity of the inner product. In either case the result is $(f,\varphi_\mu)=x_\mu$.

\leavevmode\phantomsection\textbf{Theorem 6.}\label{thm:II.6} Let $\varphi_1,\varphi_2,\ldots$ be an orthonormal system and $f$ arbitrary. The series defining $f'=\sum_\nu x_\nu\varphi_\nu$, with $x_\nu=(f,\varphi_\nu)$ ($\nu=1,2,\ldots$), always converges when it has infinitely many terms. The element $f-f'$ is orthogonal to $\varphi_1,\varphi_2,\ldots$.

\emph{Proof.} Convergence follows from \hyperref[thm:II.4]{Theorems 4} and \hyperref[thm:II.5]{5}. By the corollary to \hyperref[thm:II.5]{Theorem 5}, $(f',\varphi_\nu)=x_\nu=(f,\varphi_\nu)$, so $(f-f',\varphi_\nu)=0$.

After these preparations, we can give general criteria for completeness of orthonormal systems, valid also under \hyperref[ax:II.Cinfty]{$C^{(\infty)}$}.

\leavevmode\phantomsection\textbf{Theorem 7.}\label{thm:II.7} Let $\varphi_1,\varphi_2,\ldots$ be an orthonormal system. Each of the following conditions is necessary and sufficient for completeness:

$\alpha)$ The closed linear manifold $[\varphi_1,\varphi_2,\ldots]$ spanned by $\varphi_1,\varphi_2,\ldots$ equals $\mathfrak{R}$.

$\beta)$ We always have $f=\sum_\nu x_\nu\varphi_\nu$, where $x_\nu=(f,\varphi_\nu)$ ($\nu=1,2,\ldots$; convergence by \hyperref[thm:II.6]{Theorem 6}).

$\gamma)$ We always have
\[
(f,g)=\sum_\nu(f,\varphi_\nu)\overline{(g,\varphi_\nu)}
\]
(absolute convergence by \hyperref[thm:II.4]{Theorem 4}).

\emph{Proof.} If $\varphi_1,\varphi_2,\ldots$ is complete, then $f-\sum_\nu x_\nu\varphi_\nu$, where $x_\nu=(f,\varphi_\nu)$, $\nu=1,2,\ldots$, is zero, since by \hyperref[thm:II.6]{Theorem 6} it is orthogonal to $\varphi_1,\varphi_2,\ldots$. Thus $\beta)$ holds. If $\beta)$ holds, every $f$ is the limit of its sums $\sum_{\nu=1}^Nx_\nu\varphi_\nu$ as $N\to\infty$, if the system is infinite, and therefore belongs to $[\varphi_1,\varphi_2,\ldots]$. Hence this manifold equals $\mathfrak{R}$, and $\alpha)$ holds. If $\alpha)$ holds, argue as follows: an $f$ orthogonal to all $\varphi_1,\varphi_2,\ldots$ is also orthogonal to their linear combinations and, by continuity, to their limit points, hence to all of $[\varphi_1,\varphi_2,\ldots]$. It is therefore orthogonal to all of $\mathfrak{R}$, and thus to itself: $(f,f)=0$, $f=0$. The system is consequently complete.

We thus have the logical scheme
\[
\text{completeness}\longrightarrow\beta)\longrightarrow\alpha)
\longrightarrow\text{completeness},
\]
so $\alpha)$ and $\beta)$ have been shown necessary and sufficient.

To deduce completeness from $\gamma)$, let $f$ be orthogonal to all $\varphi_1,\varphi_2,\ldots$ and put $f=g$. Then $(f,f)=\sum_\nu0\cdot0=0$, so $f=0$; the system is complete.
% Source PDF page 36 / original page 29.
\hypertarget{origpage-29}{}%
Conversely, $\beta)$, which is equivalent to completeness, gives
\begin{align*}
(f,g)&=\lim_{N\to\infty}
\left(\sum_{\nu=1}^N(f,\varphi_\nu)\varphi_\nu,
\sum_{\nu=1}^N(g,\varphi_\nu)\varphi_\nu\right)\\
&=\lim_{N\to\infty}\sum_{\mu,\nu=1}^N
(f,\varphi_\mu)\overline{(g,\varphi_\nu)}(\varphi_\mu,\varphi_\nu)\\
&=\lim_{N\to\infty}\sum_{\nu=1}^N
(f,\varphi_\nu)\overline{(g,\varphi_\nu)}\\
&=\sum_{\nu=1}^\infty(f,\varphi_\nu)\overline{(g,\varphi_\nu)}.
\end{align*}
If the system is finite the limiting processes are unnecessary. This is $\gamma)$, which is therefore also necessary and sufficient.

\leavevmode\phantomsection\textbf{Theorem 8.}\label{thm:II.8} For every sequence $f_1,f_2,\ldots$ there is an orthonormal system $\varphi_1,\varphi_2,\ldots$ spanning the same linear manifold. Both sequences may terminate after finitely many terms.

\emph{Proof.} First replace $f_1,f_2,\ldots$ by a subsequence $g_1,g_2,\ldots$ spanning the same linear manifold and consisting of linearly independent elements. Proceed as follows: let $g_1$ be the first nonzero $f_n$; $g_2$, the first $f_n$ distinct from every $a_1g_1$; $g_3$, the first $f_n$ distinct from every $a_1g_1+a_2g_2$; and so on. If, for some $p$, no $f_n$ exists that differs from all $a_1g_1+\cdots+a_pg_p$, we stop at $g_p$. These $g_1,g_2,\ldots$ evidently meet the requirements.

We now apply E.~Schmidt's so-called ``orthogonalization procedure'':
\begin{alignat*}{2}
\gamma_1&=g_1,&\quad \varphi_1&=\frac1{\|\gamma_1\|}\gamma_1,\\
\gamma_2&=g_2-(g_2,\varphi_1)\varphi_1,
&\quad \varphi_2&=\frac1{\|\gamma_2\|}\gamma_2,\\
\gamma_3&=g_3-(g_3,\varphi_1)\varphi_1-(g_3,\varphi_2)\varphi_2,
&\quad \varphi_3&=\frac1{\|\gamma_3\|}\gamma_3,\\
&\cdots&\quad &\cdots.
\end{alignat*}
Each $\varphi_p$ is well-defined, since its denominator $\|\gamma_p\|$ is nonzero. Otherwise $\gamma_p=0$, and $g_p$ would be a linear combination of $\varphi_1,\ldots,\varphi_{p-1}$, hence of $g_1,\ldots,g_{p-1}$, contrary to the assumption. Further, $g_p$ is clearly a linear combination of $\varphi_1,\ldots,\varphi_p$, and $\varphi_p$ a linear combination of $g_1,\ldots,g_p$. Thus the two sequences determine the same linear manifold.

Finally, by construction $\|\varphi_p\|=1$, and for $q<p$ we have $(\gamma_p,\varphi_q)=0$, hence $(\varphi_p,\varphi_q)=0$. Since we may interchange $p,q$, the last equality holds whenever $p\ne q$. Thus $\varphi_1,\varphi_2,\ldots$ is an orthonormal system.

\leavevmode\phantomsection\textbf{Theorem 9.}\label{thm:II.9} For every closed linear manifold $\mathfrak{M}$ there is an orthonormal system spanning precisely $\mathfrak{M}$ as a closed linear manifold.

% Source PDF page 37 / original page 30.
\hypertarget{origpage-30}{}%
\emph{Proof.} In case \hyperref[ax:II.Cn]{$C^{(n)}$} the theorem is self-evident. If $\mathfrak{R}$ satisfies \hyperref[ax:II.A]{A}, \hyperref[ax:II.B]{B}, and \hyperref[ax:II.Cn]{$C^{(n)}$}, every linear manifold $\mathfrak{M}$ in $\mathfrak{R}$ satisfies \hyperref[ax:II.A]{A}, \hyperref[ax:II.B]{B}, and \hyperref[ax:II.Cn]{$C^{(m)}$} for some $m\leq n$. The remark to \hyperref[thm:II.3n]{Theorem $3^{(n)}$} therefore applies to $\mathfrak{M}$: there is an orthonormal system $\varphi_1,\ldots,\varphi_m$ complete in $\mathfrak{M}$. By \hyperref[thm:II.7]{Theorem 7}, $\alpha)$, this is precisely the assertion. As we see, the assumption of closedness is unnecessary in this case; indeed, closedness is proved. Compare the discussion preceding \hyperref[def:II.5]{Definition 5}.

In case \hyperref[ax:II.Cinfty]{$C^{(\infty)}$}, recall that $\mathfrak{R}$ is separable by \hyperref[ax:II.E]{E}. We shall show that $\mathfrak{M}$ is also separable; indeed every subset of $\mathfrak{R}$ is separable. Take the everywhere dense sequence $f_1,f_2,\ldots$ in $\mathfrak{R}$, as in \hyperref[ax:II.E]{E} of \hyperref[sec:II.1]{II.1}. For each $f_n$ and each $m=1,2,\ldots$, form the ball $\mathfrak{K}_{n,m}$ consisting of all $f$ with $\|f-f_n\|<1/m$. For each $\mathfrak{K}_{n,m}$ containing points of $\mathfrak{M}$, choose one such point $g_{n,m}$. Thus $g_{n,m}$ may be undefined for some $n,m$, but the defined points form a sequence in $\mathfrak{M}$.\origNote{50} Let $f$ be any point of $\mathfrak{M}$, and $\varepsilon<0$.\tnote{The source prints $\varepsilon<0$; the following approximation argument requires $\varepsilon>0$. The printed inequality is retained.} There is an $m$ with $1/m<\varepsilon/2$, and an $f_n$ with $\|f_n-f\|<1/m$. Thus $\mathfrak{K}_{nm}$ contains a point of $\mathfrak{M}$, namely $f$, so $g_{nm}$ is defined and $\|f_n-g_{nm}\|<1/m$. Hence $\|f-g_{nm}\|<2/m<\varepsilon$. Thus $f$ is a limit point of the defined $g_{nm}$, and this sequence meets the requirement.

Denote this sequence, which lies in $\mathfrak{M}$ and is everywhere dense there, by $f_1,f_2,\ldots$. The closed linear manifold $[f_1,f_2,\ldots]$ it determines includes all its limit points, hence all of $\mathfrak{M}$. But since $\mathfrak{M}$ is a closed linear manifold containing $f_1,f_2,\ldots$, $[f_1,f_2,\ldots]$ is also a subset of $\mathfrak{M}$, and therefore equals it. Choose the orthonormal system $\varphi_1,\varphi_2,\ldots$ according to \hyperref[thm:II.8]{Theorem 8}. Then $\{\varphi_1,\varphi_2,\ldots\}=\{f_1,f_2,\ldots\}$. Adjoining the limit points on both sides gives $[\varphi_1,\varphi_2,\ldots]=[f_1,f_2,\ldots]=\mathfrak{M}$, which was our assertion.

We need only put $\mathfrak{M}=\mathfrak{R}$ in \hyperref[thm:II.9]{Theorem 9} to obtain, by \hyperref[thm:II.7]{Theorem 7}, $\alpha)$, a complete orthonormal system $\varphi_1,\varphi_2,\ldots$. Thus complete orthonormal systems exist. We are now in a position to show that $\mathfrak{R}$ is an $\mathfrak{R}_n$ or $\mathfrak{R}_\infty$, according as \hyperref[ax:II.Cn]{$C^{(n)}$} or \hyperref[ax:II.Cinfty]{$C^{(\infty)}$} holds, and is therefore determined in all its properties.

It remains only to show that $\mathfrak{R}$ can be mapped one-to-one onto the set of all $\{x_1,\ldots,x_n\}$, or of all $\{x_1,x_2,\ldots\}$ with $\sum_{\nu=1}^\infty|x_\nu|^2$ finite, in such a way that:

1. $f\longleftrightarrow\{x_1,x_2,\ldots\}$ implies $af\longleftrightarrow\{ax_1,ax_2,\ldots\}$.

2. $f\longleftrightarrow\{x_1,x_2,\ldots\}$ and $g\longleftrightarrow\{y_1,y_2,\ldots\}$ imply $f+g\longleftrightarrow\{x_1+y_1,x_2+y_2,\ldots\}$.

% Source PDF page 38 / original page 31.
\hypertarget{origpage-31}{}%
3. $f\longleftrightarrow\{x_1,x_2,\ldots\}$ and $g\longleftrightarrow\{y_1,y_2,\ldots\}$ imply
\[
(f,g)=\sum_{\nu=1}^{n\text{ or }\infty}x_\nu\overline{y_\nu}.
\]
(In the infinite case, absolute convergence must also be proved in 3.) We now give this correspondence $f\longleftrightarrow\{x_1,x_2,\ldots\}$.

Let $\varphi_1,\varphi_2,\ldots$ be a complete orthonormal system. In case \hyperref[ax:II.Cn]{$C^{(n)}$} it terminates at $\varphi_n$; in case \hyperref[ax:II.Cinfty]{$C^{(\infty)}$} it is infinite (\hyperref[thm:II.3n]{Theorems $3^{(n)}$}, \hyperref[thm:II.3infty]{$3^{(\infty)}$}). We put
\[
f=\sum_{\nu=1}^{n\text{ or }\infty}x_\nu\varphi_\nu.
\]
By \hyperref[thm:II.5]{Theorem 5} this series converges in the infinite case as well, since $\sum_{\nu=1}^\infty|x_\nu|^2$ is finite; thus precisely all of $\mathfrak{R}_n$ or $\mathfrak{R}_\infty$ is covered. By \hyperref[thm:II.7]{Theorem 7}, $\beta)$, and the finiteness of $\sum_{\nu=1}^{n\text{ or }\infty}|(f,\varphi_\nu)|^2$ (\hyperref[thm:II.4]{Theorem 4}), all of $\mathfrak{R}$ is covered as well: put $x_\nu=(f,\varphi_\nu)$. It is clear that each $\{x_1,x_2,\ldots\}$ corresponds to only one $f$; the converse follows from the corollary to \hyperref[thm:II.5]{Theorem 5}. Conditions 1 and 2 are evidently satisfied, and 3 follows from \hyperref[thm:II.7]{Theorem 7}, $\gamma)$.

\section[Digression on Conditions A--E]{Digression on Conditions A--E\texorpdfstring{\origNote{51}}{}}\label{sec:II.3}
We still have to verify assertion 2 at the end of \hyperref[sec:I.4]{I.4}: that $F_Z$ and $F_\Omega$ actually satisfy conditions \hyperref[ax:II.A]{A}--\hyperref[ax:II.E]{E}. It suffices to consider $F_\Omega$, for in \hyperref[sec:II.2]{II.2} we have already shown that a space $\mathfrak{R}$ satisfying \hyperref[ax:II.A]{A}--\hyperref[ax:II.E]{E} must agree in all its properties with $\mathfrak{R}_\infty$, that is, $F_Z$, so \hyperref[ax:II.A]{A}--\hyperref[ax:II.E]{E} must hold for $F_Z$ as well. We shall also prove the independence of conditions \hyperref[ax:II.D]{D} and \hyperref[ax:II.E]{E} from \hyperref[ax:II.A]{A}--\hyperref[ax:II.Cinfty]{$C^{(\infty)}$}, mentioned in \hyperref[sec:II.2]{II.2}, and the fact that they follow from \hyperref[ax:II.A]{A}--\hyperref[ax:II.Cn]{$C^{(n)}$}, that is, that they hold in $\mathfrak{R}_n$. These three purely mathematical questions form the subject of this section.

We begin by verifying \hyperref[ax:II.A]{A}--\hyperref[ax:II.E]{E} in $F_\Omega$. This requires the Lebesgue integral; for its foundations we refer the reader to the relevant specialist literature.\origNote{52} The Lebesgue integral is used only on this occasion, however, and familiarity with it is not required for the later chapters.

In \hyperref[sec:I.4]{I.4} we introduced $\Omega$ as the $k$-dimensional space of the $q_1,\ldots,q_k$, and $F_\Omega$ as the totality of all functions $f(q_1\cdots q_k)$ for which
\[
\underbrace{\int\cdots\int}_{\Omega}|\varphi(q_1\cdots q_k)|^2\,dq_1\cdots dq_k
\]
is finite; here all of $q_1,\ldots,q_k$ range from $-\infty$ to $+\infty$. All our deductions would remain valid, and their proofs could mostly be carried over word for word, if we restricted the ranges of $q_1,\ldots,q_k$, so that $\Omega$ became, for example, a half-space, the interior of a cube, the interior of a sphere, the exterior of these figures, and so on; indeed, we could even choose $\Omega$ to be a curved surface, such as the surface of a sphere.

% Source PDF page 39 / original page 32.
\hypertarget{origpage-32}{}%
To avoid inessential complications, which the reader can readily work through by following the pattern of our proofs, we confine ourselves to the simplest case just described. We now take \hyperref[ax:II.A]{A}--\hyperref[ax:II.E]{E} in succession.

\emph{Ad \hyperref[ax:II.A]{A}.} We must show that if $f,g$ belong to $F_\Omega$, then $af,f\pm g$ do also; that is, the finiteness of
\[
\int_\Omega|f|^2,\qquad\int_\Omega|g|^2
\]
must imply that of $\int_\Omega|af|^2=|a|^2\int_\Omega|f|^2$ and $\int_\Omega|f\pm g|^2$. Since no misunderstanding is possible, we use these abbreviations for
\[
\underbrace{\int\cdots\int}_{\Omega}|f(q_1\cdots q_k)|^2\,dq_1\cdots dq_k,
\quad
\underbrace{\int\cdots\int}_{\Omega}|g(q_1\cdots q_k)|^2\,dq_1\cdots dq_k.
\]
The first assertion is trivial. In view of $|f\pm g|^2=|f|^2+|g|^2\pm2\operatorname{Re}(f\overline g)$,\origNote{53} the second is established as soon as the finiteness of $\int_\Omega|f\overline g|=\int_\Omega|f||g|$ is assured. But $|f||g|\leq\frac12(|f|^2+|g|^2)$, so this follows from the assumptions.

\emph{Ad \hyperref[ax:II.B]{B}.} We define $(f,g)$ as $\int_\Omega f\overline g$; as we have just seen, this integral converges absolutely. All the properties postulated in \hyperref[ax:II.B]{B} are evident except the last: that $(f,f)=0$ implies $f\equiv0$. The equation $(f,f)=0$ means $\int_\Omega|f|^2=0$, so the set of points at which $|f|^2>0$, that is, $f(q_1\cdots q_k)\ne0$, must have Lebesgue measure zero. Suppose we regard two functions $f,g$ as not essentially different when $f\ne g$, that is, $f(q_1\cdots q_k)\ne g(q_1\cdots q_k)$, holds only on a set of points $q_1\cdots q_k$ of Lebesgue measure zero.\origNote{54} We may then conclude that $f\equiv0$.

\emph{Ad \hyperref[ax:II.Cn]{C}.} Let $O_1,\ldots,O_n$ be $n$ regions in $\Omega$, no two of which have a common point, and each of whose Lebesgue measures is positive but finite. Let $f_l(q_1\cdots q_k)$ equal $1$ in $O_l$ and $0$ elsewhere. Since $\int_\Omega|f_l|^2$ equals the measure of $O_l$, the function $f_l$ belongs to $F_\Omega$ ($l=1,\ldots,k$). The $f_1,\ldots,f_n$ are linearly independent. For $a_1f_1+\cdots+a_nf_n\equiv0$ implies that the function on the left vanishes except on a set of Lebesgue measure zero; hence it has zeros in every $O_l$. But it is constant, equal to $a_l$, in $O_l$, so $a_l=0$, $l=1,\ldots,n$. This construction works for every $n$, and therefore \hyperref[ax:II.Cinfty]{$C^{(\infty)}$} holds.

\emph{Ad \hyperref[ax:II.D]{D}.} Let the sequence $f_1,f_2,\ldots$ satisfy Cauchy's convergence condition: for every $\varepsilon>0$ there is an $N=N(\varepsilon)$ such that $\int_\Omega|f_m-f_n|^2<\varepsilon$ whenever $m,n\geq N$. Choose
\[
n_1=N\left(\frac18\right);\qquad
n_2\geq n_1,N\left(\frac1{8^2}\right);\qquad
n_3\geq n_1,n_2,N\left(\frac1{8^3}\right);\quad\ldots.
\]
Thus $n_1\leq n_2\leq\cdots$, and $n_\nu,n_{\nu+1}\geq N(1/8^\nu)$, so
\[
\int_\Omega|f_{n_{\nu+1}}-f_{n_\nu}|^2<\frac1{8^\nu}.
\]

% Source PDF page 40 / original page 33.
\hypertarget{origpage-33}{}%
Consider the set $P^{(\nu)}$ of all points at which $|f_{n_{\nu+1}}-f_{n_\nu}|>1/2^\nu$. If its Lebesgue measure is $\mu^{(\nu)}$, then
\[
\int_\Omega|f_{n_{\nu+1}}-f_{n_\nu}|^2
\geq\mu^{(\nu)}\left(\frac1{2^\nu}\right)^2
=\frac{\mu^{(\nu)}}{4^\nu},\qquad
\frac{\mu^{(\nu)}}{4^\nu}<\frac1{8^\nu},\qquad
\mu^{(\nu)}<\frac1{2^\nu}.
\]
Now let $Q^{(\nu)}$ be the union of $P^{(\nu)},P^{(\nu+1)},P^{(\nu+2)},\ldots$. Its Lebesgue measure is
\[
\leq\mu^{(\nu)}+\mu^{(\nu+1)}+\mu^{(\nu+2)}+\cdots
<\frac1{2^\nu}+\frac1{2^{\nu+1}}+\frac1{2^{\nu+2}}+\cdots
=\frac1{2^{\nu-1}}.
\]
Outside $Q^{(\nu)}$ we have
\[
|f_{n_{\nu+1}}-f_{n_\nu}|<\frac1{2^\nu},\quad
|f_{n_{\nu+2}}-f_{n_{\nu+1}}|<\frac1{2^{\nu+1}},\quad
|f_{n_{\nu+3}}-f_{n_{\nu+2}}|<\frac1{2^{\nu+2}},\quad\ldots,
\]
and hence, in general, for $\nu\leq\nu'\leq\nu''$,
\begin{align*}
|f_{n_{\nu''}}-f_{n_{\nu'}}|&<|f_{n_{\nu'+1}}-f_{n_{\nu'}}|+|f_{n_{\nu'+2}}-f_{n_{\nu'+1}}|+\cdots+|f_{n_{\nu''}}-f_{n_{\nu''-1}}|\\
&<\frac1{2^{\nu'}}+\frac1{2^{\nu'+1}}+\cdots+\frac1{2^{\nu''-1}}<\frac1{2^{\nu'-1}}.
\end{align*}
As $\nu'\to\infty$, this tends to zero independently of $\nu''$; that is, the sequence $f_{n_1},f_{n_2},\ldots$ satisfies Cauchy's convergence condition provided $q_1\cdots q_k$ does not lie in $Q^{(\nu)}$. Since, for fixed $q_1\cdots q_k$, these are numbers, the sequence converges. Equivalently, if $f_{n_1},f_{n_2},\ldots$ fails to converge at some $q_1\cdots q_k$, that point must belong to $Q^{(\nu)}$. Let $Q$ be the set of all $q_1\cdots q_k$ at which convergence fails. Then $Q$ is a subset of $Q^{(\nu)}$, so its measure is less than that of $Q^{(\nu)}$, and thus less than $1/2^{\nu-1}$. This must hold for every $\nu$, although $Q$ is defined independently of $\nu$; hence $Q$ has Lebesgue measure zero. It therefore makes no difference to set, for example, all $f_n$ equal to zero on $Q$ (see \hyperlink{note-54}{note 54}). Then $f_{n_1},f_{n_2},\ldots$ converges on $Q$ as well, and consequently everywhere.

We have thus given a subsequence $f_{n_1},f_{n_2},\ldots$ of $f_1,f_2,\ldots$ which converges at every point $q_1\cdots q_k$; this need not be true of $f_1,f_2,\ldots$ itself. Denote the limit of $f_{n_1},f_{n_2},\ldots$ by $f=f(q_1\cdots q_k)$. We have still to show: 1. that $f$ belongs to $F_\Omega$, that is, $\int_\Omega|f|^2$ is finite; 2. that $f$ is the limit of $f_{n_1},f_{n_2},\ldots$ not merely in the sense of convergence at every $q_1,\ldots,q_k$, but also in the sense of convergence in norm in Hilbert space, that is, $\|f-f_{n_\nu}\|\to0$, or $\int_\Omega|f-f_{n_\nu}|^2\to0$; 3. that in this sense it is even the limit of the entire sequence $f_1,f_2,\ldots$, that is, $\|f-f_n\|\to0$, or $\int_\Omega|f-f_n|^2\to0$.

Let $\varepsilon>0$, choose $\nu_0$ so that $n_{\nu_0}\geq N(\varepsilon)$, for example $1/8^{\nu_0}\leq\varepsilon$, and let $\nu\geq\nu_0$, $n\geq N(\varepsilon)$. Then $\int_\Omega|f_{n_\nu}-f_n|^2<\varepsilon$. As $\nu\to\infty$ the integrand tends to $|f-f_n|^2$, so a convergence theorem for Lebesgue integrals (see the reference in \hyperlink{note-53}{note 53})\tnote{The source cites Note~53 here; the references for Lebesgue integration are in \hyperlink{note-52}{Note~52}.} gives $\int_\Omega|f-f_n|^2\leq\varepsilon$.

% Source PDF page 41 / original page 34.
\hypertarget{origpage-34}{}%
First, therefore, $\int_\Omega|f-f_1|^2$ is finite, that is, $f-f_1$ belongs to $F_\Omega$; since $f_1$ belongs to $F_\Omega$, so does $f$. This proves 1. Second, the inequality above gives $\int_\Omega|f-f_n|^2\to0$ as $n\to\infty$, proving 2 and 3.

\emph{Ad \hyperref[ax:II.E]{E}.} We must give a sequence of functions $f_1,f_2,\ldots$ everywhere dense in $F_\Omega$.

Let $\Omega_1,\Omega_2,\ldots$ be a sequence of regions in $\Omega$, each of finite measure, which together fill all of $\Omega$; for example, let $\Omega_N$ be the ball of radius $N$ about the origin. Let $f=f(q_1\cdots q_k)$ be any element of $F_\Omega$. For every $N=1,2,\ldots$ define $f_N=f_N(q_1\cdots q_k)$ by
\[
f_N(q_1\cdots q_k)=
\begin{cases}
f(q_1\cdots q_k),&q_1\cdots q_k\in\Omega_N\text{ and }|f(q_1\cdots q_k)|\leq N,\\
0,&\text{otherwise}.
\end{cases}
\]
As $N\to\infty$, $f_N(q_1\cdots q_k)\to f(q_1\cdots q_k)$; indeed, equality holds from a certain $N$ onward. Thus $|f-f_N|^2\to0$. Furthermore, $f-f_N$ equals either $0$ or $f$, so $|f-f_N|^2\leq|f|^2$. The integrals $\int_\Omega|f-f_N|^2$ thus have the fixed, finite majorant $\int_\Omega|f|^2$. Since the integrands tend to zero, so do the integrals, by the convergence theorem just cited: $\int_\Omega|f-f_N|^2\to0$, $\|f-f_N\|\to0$.

Call $G$ the class of all functions $g=g(q_1\cdots q_k)$ for which the set of points where $g\ne0$ has finite measure and which satisfy an inequality $|g|\leq C$ throughout the space, with arbitrary but fixed $C$. All the $f_N$ above belong to $G$, so $G$ is everywhere dense in $F_\Omega$.

Let $g$ belong to $G$, and let $\varepsilon>0$. Let $M$ be the measure of the set where $g\ne0$, and $C$ an upper bound for $|g|$. Choose a chain of rational numbers $-C<\rho_1<\rho_2<\cdots<\rho_t<C$ such that $\rho_1<-C+\varepsilon$, $\rho_2<\rho_1+\varepsilon$, $\ldots$, $\rho_t<\rho_{t-1}+\varepsilon$, $C<\rho_t+\varepsilon$; this is easily done. Replace each value of $\operatorname{Re}g(q_1\cdots q_k)$ by the nearest $\rho_s$ ($s=1,2,\ldots,t$), leaving only $0$ unchanged. This gives a new function $h_1(q_1\cdots q_k)$ differing everywhere from $\operatorname{Re}g$ by less than $\varepsilon$. Similarly form $h_2(q_1\cdots q_k)$ from $\operatorname{Im}g$. Then, for $h=h_1+ih_2$,
\begin{align*}
\int_\Omega|g-h|^2
&=\int_\Omega|\operatorname{Re}g-h_1|^2+\int_\Omega|\operatorname{Im}g-h_2|^2\\
&\leq M\varepsilon^2+M\varepsilon^2=2M\varepsilon^2,\\
\|g-h\|&\leq\sqrt{2M}\,\varepsilon.
\end{align*}
Given $\delta>0$, choose $\varepsilon<\delta/\sqrt{2M}$; then $\|g-h\|<\delta$.

% Source PDF page 42 / original page 35.
\hypertarget{origpage-35}{}%
Call $H$ the class of all functions $h=h(q_1\cdots q_k)$ taking only finitely many distinct values, all of the form $\rho+i\sigma$ with rational $\rho,\sigma$, and taking each nonzero value only on a set of finite measure. The functions $h$ just constructed belong to $H$, so $H$ is everywhere dense in $G$, and hence in $F_\Omega$.

For a set $\Pi$ of finite Lebesgue measure define a function $f_\Pi=f_\Pi(q_1\cdots q_k)$ by
\[
f_\Pi(q_1\cdots q_k)=\begin{cases}1,&\text{in }\Pi,\\0,&\text{otherwise}.\end{cases}
\]
The class $H$ evidently consists of all
\[
\sum_{s=1}^t(\rho_s+i\sigma_s)f_{\Pi_s}
\qquad(t=1,2,\ldots;\ \rho_s,\sigma_s\text{ rational}).
\]
We now seek a sequence of $\Pi$-sets, $\Pi^{(1)},\Pi^{(2)},\ldots$, with the following property: for every $\Pi$-set and every $\varepsilon>0$ there is a $\Pi^{(n)}$ such that the measure of the set of all points belonging to $\Pi$ but not to $\Pi^{(n)}$, or to $\Pi^{(n)}$ but not to $\Pi$, is less than $\varepsilon$. This set is called the symmetric difference of $\Pi$ and $\Pi^{(n)}$. If we have such a sequence, the functions
\[
\sum_{s=1}^t(\rho_s+i\sigma_s)f_{\Pi^{(n_s)}}
\qquad(t=1,2,\ldots;\ \rho_s,\sigma_s\text{ rational};\ n_s=1,2,\ldots)
\]
are everywhere dense in $H$. For if we choose $\Pi^{(n_s)}$ for each of the earlier $\Pi_s$ as just described, we obtain
\begin{align*}
&\int_\Omega\left|\sum_{s=1}^t(\rho_s+\sigma_s)f_{\Pi_s}-\sum_{s=1}^t(\rho_s+i\sigma_s)f_{\Pi^{(n_s)}}\right|^2\\
&\quad\leq\sum_{s=1}^t\int_\Omega\left|(\rho_s+i\sigma_s)f_{\Pi_s}-(\rho_s+i\sigma_s)f_{\Pi^{(n_s)}}\right|^2\\
&\quad=\sum_{s=1}^t(\rho_s^2+\sigma_s^2)\int_\Omega|f_{\Pi_s}-f_{\Pi^{(n_s)}}|^2\\
&\quad=\sum_{s=1}^t(\rho_s^2+\sigma_s^2)\,\text{measure of the symmetric difference of }\Pi_s,\Pi^{(n_s)}\\
&\quad<\sum_{s=1}^t(\rho_s^2+\sigma_s^2)\,\varepsilon.
\end{align*}
Given $\delta>0$, the choice $\varepsilon=\delta^2/\sum_{s=1}^t(\rho_s^2+\sigma_s^2)$ already yields
\[
\left\|\sum_{s=1}^t(\rho_s+i\sigma_s)f_{\Pi_s}-\sum_{s=1}^t(\rho_s+i\sigma_s)f_{\Pi^{(n_s)}}\right\|<\delta.
\]
The functions $\sum_{s=1}^t(\rho_s+\sigma_s)f_{\Pi^{(n_s)}}$ form a sequence if ordered appropriately.\tnote{The source omits $i$ from $\rho_s+i\sigma_s$ both here and in the first term of the preceding integral estimate. Both omissions are retained; the surrounding construction uses complex coefficients $\rho_s+i\sigma_s$.} For let $\tau$ be a common denominator of all $\rho_1,\sigma_1,\ldots,\rho_t,\sigma_t$, and $\rho'_1,\sigma'_1,\ldots,\rho'_t,\sigma'_t$ their numerators. The functions in question are then
\[
\frac1\tau\sum_{s=1}^t(\rho'_s+i\sigma'_s)f_{\Pi^{(n_s)}},
\]

% Source PDF page 43 / original page 36.
\hypertarget{origpage-36}{}%
where $t,\tau=1,2,\ldots$; $\rho'_s,\sigma'_s=0,\pm1,\pm2,\ldots$; and $n_s=1,2,\ldots$, for $s=1,\ldots,t$. Arranging these functions in a sequence is the same task as arranging their numerical parameters $t,\tau,\rho'_1,\sigma'_1,\ldots,\rho'_t,\sigma'_t,n_1,\ldots,n_t$. Order these tuples by increasing values of
\[
t+\tau+|\rho'_1|+|\sigma'_1|+\cdots+|\rho'_t|+|\sigma'_t|+n_1+\cdots+n_t.
\]
Only finitely many of the tuples correspond to each given value of this sum. If we also put each of these finite collections in any order, we obtain a simple sequence.

To give the required sequence of sets $\Pi^{(1)},\Pi^{(2)},\ldots$, we use the fact that for every set $\Pi$ of finite Lebesgue measure $M$ and every $\delta>0$ there is an open point set $\Pi'$ containing $\Pi$ whose measure exceeds that of $\Pi$ by less than $\delta$; see the references in \hyperlink{note-52}{note 52}, and also \hyperlink{note-45}{note 45}, where the concept of an open point set is defined. For an open $\Pi'$ and an $\varepsilon>0$, however, there evidently exists a set $\Pi''$ composed of finitely many cubes, contained in $\Pi'$, whose measure falls short of that of $\Pi''$ by less than $\delta$.\tnote{The source compares the measure of $\Pi''$ with itself; the comparison required here is with $\Pi'$. It also introduces $\varepsilon$ where the estimate uses $\delta$. These printed slips are retained.} The edge lengths and center coordinates of the cubes may all be chosen rational. It is now easy to see that the symmetric difference defined above, of $\Pi$ and $\Pi''$, has measure less than $\delta+\delta=2\delta$, and hence less than $\varepsilon$ when $\delta=\varepsilon/2$. We have therefore reached our goal if we can arrange the sets of cubes just described in a sequence.

These sets are characterized by the number $n=1,2,\ldots$ of their cubes, together with their edge lengths $\varkappa^{(\nu)}$ and center coordinates $\xi_1^{(\nu)},\ldots,\xi_k^{(\nu)}$ ($\nu=1,\ldots,n$). The $\varkappa^{(\nu)},\xi_1^{(\nu)},\ldots,\xi_k^{(\nu)}$ are rational; let their common denominator be $\eta=1,2,\ldots$, and their numerators $\varkappa'^{(\nu)}=1,2,\ldots$, $\xi_1'^{(\nu)},\ldots,\xi_k'^{(\nu)}=0,\pm1,\pm2,\ldots$. Thus our sets of cubes are characterized by the tuples
\[
n,\eta,\varkappa'^{(1)},\xi_1'^{(1)},\ldots,\xi_k'^{(1)},\ldots,\varkappa'^{(n)},\xi_1'^{(n)},\ldots,\xi_k'^{(n)}.
\]
Ordering them by increasing sums
\[
n+\eta+\varkappa'^{(1)}+|\xi_1'^{(1)}|+\cdots+|\xi_k'^{(1)}|+\cdots+\varkappa'^{(n)}+|\xi_1'^{(n)}|+\cdots+|\xi_k'^{(n)}|,
\]
we obtain a simple sequence, exactly as in the earlier analogous example of linear combinations of functions.

Before proceeding, let us answer the following question. Given a space $\mathfrak{R}$ satisfying \hyperref[ax:II.A]{A}--\hyperref[ax:II.E]{E}, with \hyperref[ax:II.Cinfty]{$C^{(\infty)}$}, in which subsets $\mathfrak{M}$ of $\mathfrak{R}$ do \hyperref[ax:II.A]{A}--\hyperref[ax:II.E]{E} hold again, with the definitions of $af$, $f\pm g$, and $(f,g)$ unchanged?

% Source PDF page 44 / original page 37.
\hypertarget{origpage-37}{}%
For \hyperref[ax:II.A]{A} to hold, $\mathfrak{M}$ must be a linear manifold. \hyperref[ax:II.B]{B} holds automatically. We postpone \hyperref[ax:II.Cn]{C} for a moment: some \hyperref[ax:II.Cn]{$C^{(n)}$} or \hyperref[ax:II.Cinfty]{$C^{(\infty)}$} holds in any event. \hyperref[ax:II.D]{D} says that if a sequence in $\mathfrak{M}$ satisfies Cauchy's convergence condition, it has a limit in $\mathfrak{M}$. Since such a sequence certainly has a limit in $\mathfrak{R}$, all that is required is that this limit also belong to $\mathfrak{M}$; that is, $\mathfrak{M}$ must be closed. \hyperref[ax:II.E]{E} always holds, as we saw in the proof of \hyperref[thm:II.9]{Theorem 9}. Thus, to summarize, $\mathfrak{M}$ must be a closed linear manifold. Let $\varphi_1,\varphi_2,\ldots$ be the orthonormal system spanning $\mathfrak{M}$ (\hyperref[thm:II.9]{Theorem 9}). If it is infinite, \hyperref[ax:II.Cinfty]{$C^{(\infty)}$} evidently holds, and $\mathfrak{M}$ is isomorphic to $\mathfrak{R}_\infty$, hence to $\mathfrak{R}$ itself. If it terminates at $\varphi_n$, so that the system is $\varphi_1,\ldots,\varphi_n$, then \hyperref[ax:II.Cn]{$C^{(n)}$} holds, for example by \hyperref[thm:II.3n]{Theorem $3^{(n)}$}; that is, $\mathfrak{M}$ is isomorphic to $\mathfrak{R}_n$.

But \hyperref[ax:II.D]{D} and \hyperref[ax:II.E]{E} hold in $\mathfrak{M}$ in any event, and therefore in every $\mathfrak{R}_n$; thus they also follow from \hyperref[ax:II.A]{A}--\hyperref[ax:II.Cn]{$C^{(n)}$}.

As we see, we have used logical devices to avoid directly verifying \hyperref[ax:II.A]{A}--\hyperref[ax:II.E]{E}, with \hyperref[ax:II.Cn]{$C^{(n)}$} or \hyperref[ax:II.Cinfty]{$C^{(\infty)}$}, in $\mathfrak{R}_n$ or $\mathfrak{R}_\infty$. Such a verification presents no essential difficulty, however, and is left to the reader.

It remains to show the independence of \hyperref[ax:II.D]{D} and \hyperref[ax:II.E]{E} from \hyperref[ax:II.A]{A}--\hyperref[ax:II.Cinfty]{$C^{(\infty)}$}. As we have just seen, every linear manifold in $\mathfrak{R}_\infty$ satisfies \hyperref[ax:II.A]{A}, \hyperref[ax:II.B]{B}, \hyperref[ax:II.E]{E}, and either \hyperref[ax:II.Cn]{$C^{(n)}$} or \hyperref[ax:II.Cinfty]{$C^{(\infty)}$}; if it is not closed, however, it fails to satisfy \hyperref[ax:II.D]{D}. In that case \hyperref[ax:II.Cinfty]{$C^{(\infty)}$} must hold, since \hyperref[ax:II.Cn]{$C^{(n)}$} implies \hyperref[ax:II.D]{D}. It is easy to give such a manifold. Let $\varphi_1,\varphi_2,\ldots$ be an orthonormal system. The elements $\sum_{\nu=1}^N x_\nu\varphi_\nu$, with $N=1,2,\ldots$ and arbitrary $x_1,\ldots,x_N$, form a linear manifold, but not a closed one. For $\sum_{\nu=1}^\infty\frac1\nu\varphi_\nu$ exists, since $\sum_{\nu=1}^\infty(1/\nu)^2$ is finite, and is a limit point of this manifold but not an element of it:
\[
\sum_{\nu=1}^N\frac1\nu\varphi_\nu\longrightarrow\sum_{\nu=1}^\infty\frac1\nu\varphi_\nu\qquad(N\to\infty).
\]
Thus \hyperref[ax:II.D]{D} is independent of \hyperref[ax:II.A]{A}--\hyperref[ax:II.Cinfty]{$C^{(\infty)}$} and \hyperref[ax:II.E]{E}.

Next consider all complex functions $x(\alpha)$ of a continuously varying parameter $\alpha$, $-\infty<\alpha<+\infty$. Require that the points where $x(\alpha)$ is nonzero can be listed in a sequence and that the sum $\sum_\alpha|x(\alpha)|^2$ over those points is finite.\origNote{55} All these functions $x(\alpha)$ form a space $\mathfrak{R}_{\mathrm{kont.}}$. For any two points $x(\alpha),y(\alpha)$ of this space, $x(\alpha)$ and $y(\alpha)$ are nonzero only on two sequences of $\alpha$'s, respectively, and we can combine these into a single sequence. Outside some sequence $\alpha_1,\alpha_2,\ldots$, therefore, $x(\alpha)=y(\alpha)=0$. We therefore need consider only the values $x_n=x(\alpha_n)$, $y_n=y(\alpha_n)$ for $n=1,2,\ldots$. As long as we are dealing with only two points of $\mathfrak{R}_{\mathrm{kont.}}$, everything behaves just as in $\mathfrak{R}_\infty$. Consequently \hyperref[ax:II.A]{A} and \hyperref[ax:II.B]{B} hold in $\mathfrak{R}_{\mathrm{kont.}}$ exactly as in $\mathfrak{R}_\infty$.\origNote{56}

% Source PDF page 45 / original page 38.
\hypertarget{origpage-38}{}%
The same is true for $k$ points of $\mathfrak{R}_{\mathrm{kont.}}$, $k=1,2,\ldots$, so \hyperref[ax:II.Cinfty]{$C^{(\infty)}$} also holds. It remains true even for a sequence of points of $\mathfrak{R}_{\mathrm{kont.}}$. Let these be $x_1(\alpha),x_2(\alpha),\ldots$. For each $n=1,2,\ldots$, the $\alpha$'s for which $x_n(\alpha)\ne0$ form a sequence $\alpha_1^{(n)},\alpha_2^{(n)},\ldots$. All these sequences together form a double sequence $\alpha_m^{(n)}$ ($n,m=1,2,\ldots$), which can also be written as the simple sequence $\alpha_1^{(1)},\alpha_2^{(1)},\alpha_1^{(2)},\alpha_3^{(1)},\alpha_2^{(2)},\alpha_1^{(3)},\ldots$. Hence \hyperref[ax:II.D]{D} also holds in $\mathfrak{R}_{\mathrm{kont.}}$ just as in $\mathfrak{R}_\infty$. \hyperref[ax:II.E]{E} is different: there all points of $\mathfrak{R}$ are involved, since all must be limit points of a suitable sequence. We therefore cannot infer it for $\mathfrak{R}_{\mathrm{kont.}}$ from $\mathfrak{R}_\infty$. Indeed, it is not satisfied, for one of its consequences fails: there is an orthonormal system which cannot be written as a sequence, contrary to \hyperref[thm:II.3infty]{Theorem $3^{(\infty)}$}.

Put
\[
x_\beta(\alpha)=\begin{cases}1,&\alpha=\beta,\\0,&\alpha\ne\beta.\end{cases}
\]
For every $\beta$, $x_\beta(\alpha)$ is an element of $\mathfrak{R}_{\mathrm{kont.}}$, and the $x_\beta(\alpha)$ form an orthonormal system. They could be written as a sequence only if all the $\beta$'s, $-\infty<\beta<+\infty$, could be so written, which, as is well known, is impossible.\origNote{57} Thus \hyperref[ax:II.E]{E} too is independent of \hyperref[ax:II.A]{A}--\hyperref[ax:II.Cinfty]{$C^{(\infty)}$} and \hyperref[ax:II.D]{D}.

Notice, incidentally, the fundamental difference between the function space of the $f(x)$ with finite $\int_{-\infty}^{\infty}|f(x)|^2\,dx$ and that of the $x(\alpha)$ with finite $\sum_\alpha|x(\alpha)|^2$. We could equally well describe the former as the space of all $x(\alpha)$ with finite $\int_{-\infty}^{\infty}|x(\alpha)|^2\,d\alpha$. The entire difference consists in replacing $\int_{-\infty}^{\infty}\cdots\,d\alpha$ by $\sum_\alpha\cdots$. Nevertheless the former space is an $F_\Omega$, thus satisfies \hyperref[ax:II.A]{A}--\hyperref[ax:II.E]{E} and is isomorphic to $\mathfrak{R}_\infty$, whereas the latter, $\mathfrak{R}_{\mathrm{kont.}}$, violates \hyperref[ax:II.E]{E} and is essentially different from $\mathfrak{R}_\infty$. Yet the two spaces are identical; only the definitions of norm in them differ!

\section{Closed Linear Manifolds}\label{sec:II.4}
Section \hyperref[sec:II.2]{II.2} is important for us not only because of the proof of isomorphism, but also because several theorems on orthonormal systems were proved there. We now wish to pursue the geometrical analysis of Hilbert space and examine more closely the closed linear manifolds, which play a role in $\mathfrak{R}_\infty$ analogous to that played in $\mathfrak{R}_n$ by straight lines, planes, and so forth, that is, the $\mathfrak{R}_m$, $m\leq n$.

First recall the notation of \hyperref[def:II.2]{Definitions 2} and \hyperref[def:II.5]{5}. If $\mathfrak{A}$ is any set in $\mathfrak{R}$, then $\{\mathfrak{A}\}$ and $[\mathfrak{A}]$ are, respectively, the linear manifold and the closed linear manifold spanned by $\mathfrak{A}$, that is, the smallest object of the respective kind containing $\mathfrak{A}$.

% Source PDF page 46 / original page 39.
\hypertarget{origpage-39}{}%
We extend this notation as follows. If $\mathfrak{A},\mathfrak{B},\ldots$ are arbitrary subsets and $f,g,\ldots$ elements of $\mathfrak{R}$, then $\{\mathfrak{A},\mathfrak{B},\ldots,f,g,\ldots\}$ and $[\mathfrak{A},\mathfrak{B},\ldots,f,g,\ldots]$ denote, respectively, the linear manifold and the closed linear manifold spanned by the set obtained by combining the $\mathfrak{A},\mathfrak{B},\ldots$ and the $f,g,\ldots$.

In particular, if $\mathfrak{M},\mathfrak{N},\ldots$, finitely or infinitely many, are closed linear manifolds, we denote the closed linear manifold $[\mathfrak{M},\mathfrak{N},\ldots]$ by $\mathfrak{M}+\mathfrak{N}+\cdots$. Clearly $\{\mathfrak{M},\mathfrak{N},\ldots\}$ consists of all sums $f+g+\cdots$, with $f$ ranging over $\mathfrak{M}$, $g$ over $\mathfrak{N}$, and so on; $[\mathfrak{M},\mathfrak{N},\ldots]=\mathfrak{M}+\mathfrak{N}+\cdots$ is obtained by adjoining the limit points. If there are only finitely many sets $\mathfrak{M},\mathfrak{N},\ldots$, and every element of each is orthogonal to every element of the others, then, as we shall soon see, these two constructions give the same result; this need not hold in general.

If $\mathfrak{M}$ is a subset of $\mathfrak{N}$, consider also the totality of elements of $\mathfrak{N}$ orthogonal to all elements of $\mathfrak{M}$. This too is evidently a closed linear manifold; denote it by $\mathfrak{N}-\mathfrak{M}$. \hyperref[thm:II.14]{Theorem 14} will make clear the reasons for calling this subtraction. Especially important is $\mathfrak{R}-\mathfrak{M}$, the set of all $f$ orthogonal to the whole of $\mathfrak{M}$; it is called the closed linear manifold complementary to $\mathfrak{M}$.

Finally, we mention three particularly simple closed linear manifolds: first, $\mathfrak{R}$ itself; second, the set consisting of $0$ alone, $\{0\}=[0]$; and third, the set of all $af$, where $f$ is a given element of $\mathfrak{R}$ and $a$ is variable, which is plainly a closed linear manifold and hence equals both $\{f\}$ and $[f]$.

We now introduce the concept of projection, entirely analogous to that of Euclidean geometry.

\leavevmode\phantomsection\textbf{Theorem 10.}\label{thm:II.10} Let $\mathfrak{M}$ be a closed linear manifold. Every $f$ can then be decomposed in one and only one way into two summands $f=g+h$, with $g$ in $\mathfrak{M}$ and $h$ in $\mathfrak{R}-\mathfrak{M}$.

\emph{Remark.} We call $g$ the projection of $f$ onto $\mathfrak{M}$, and $h$, which is orthogonal to the whole of $\mathfrak{M}$, the perpendicular from $f$ to $\mathfrak{M}$. We introduce the notation $P_{\mathfrak{M}}f$ for $g$.

\emph{Proof.} Let $\varphi_1,\varphi_2,\ldots$ be the orthonormal system which exists by \hyperref[thm:II.9]{Theorem 9} and spans the closed linear manifold $\mathfrak{M}$. Put $g=\sum_n(f,\varphi_n)\varphi_n$. By \hyperref[thm:II.6]{Theorem 6}, the series converges when it has infinitely many terms; its sum $g$ plainly belongs to $\mathfrak{M}$. Furthermore, by \hyperref[thm:II.6]{Theorem 6}, $h=f-g$ is orthogonal to all of $\varphi_1,\varphi_2,\ldots$. But the vectors orthogonal to $h$ form a closed linear manifold, so not only $\varphi_1,\varphi_2,\ldots$, but all of $\mathfrak{M}$ is orthogonal to $h$; that is, $h$ belongs to $\mathfrak{R}-\mathfrak{M}$.

% Source PDF page 47 / original page 40.
\hypertarget{origpage-40}{}%
If there were another such decomposition $f=g'+h'$, with $g'$ in $\mathfrak{M}$ and $h'$ in $\mathfrak{R}-\mathfrak{M}$, then $g+h=g'+h'$, and $g-g'=h'-h=j$. Thus $j$ would belong both to $\mathfrak{M}$ and to $\mathfrak{R}-\mathfrak{M}$, and would consequently be orthogonal to itself: $(j,j)=0$, $j=0$, and hence $g=g'$, $h=h'$.

Thus $P_{\mathfrak{M}}f$ is an operation assigning to every $f$ in $\mathfrak{R}$ its projection $P_{\mathfrak{M}}f$ onto $\mathfrak{M}$. In the next section we shall define an operator $R$ as a function defined on a subset of $\mathfrak{R}$ and taking values in $\mathfrak{R}$; that is, a correspondence assigning certain $Rf$ in $\mathfrak{R}$ to certain $f$ in $\mathfrak{R}$. It need not assign a value to every $f$; for other $f$ in $\mathfrak{R}$ it may be undefined, or ``meaningless.'' Thus $P_{\mathfrak{M}}$ is an operator defined everywhere in $\mathfrak{R}$; call it the projection operator of $\mathfrak{M}$.

\leavevmode\phantomsection\textbf{Theorem 11.}\label{thm:II.11} The operator $P_{\mathfrak{M}}$ has the following properties:
\begin{align*}
P_{\mathfrak{M}}(a_1f_1+\cdots+a_nf_n)&=a_1P_{\mathfrak{M}}f_1+\cdots+a_nP_{\mathfrak{M}}f_n,\\
(P_{\mathfrak{M}}f,g)&=(f,P_{\mathfrak{M}}g),\\
P_{\mathfrak{M}}(P_{\mathfrak{M}}f)&=P_{\mathfrak{M}}f.
\end{align*}
$\mathfrak{M}$ is the set of all values of $P_{\mathfrak{M}}$, that is, the set of all $P_{\mathfrak{M}}f$; it can also be characterized as the set of all solutions of $P_{\mathfrak{M}}f=f$, whereas $\mathfrak{R}-\mathfrak{M}$ is the set of all solutions of $P_{\mathfrak{M}}f=0$.

\emph{Remark.} In the next section we shall see that the first property characterizes the so-called linear operators, and the second the so-called Hermitian operators. The third says that applying $P_{\mathfrak{M}}$ twice has the same effect as applying it once. The usual symbolic notation for this is $P_{\mathfrak{M}}P_{\mathfrak{M}}=P_{\mathfrak{M}}$, or $P_{\mathfrak{M}}^2=P_{\mathfrak{M}}$.

\emph{Proof.} From
\[
f_1=g_1+h_1,\quad\ldots,\quad f_n=g_n+h_n,
\qquad(g_1,\ldots,g_n\in\mathfrak{M};\ h_1,\ldots,h_n\in\mathfrak{R}-\mathfrak{M}),
\]
we obtain
\[
a_1f_1+\cdots+a_nf_n=(a_1g_1+\cdots+a_ng_n)+(a_1h_1+\cdots+a_nh_n),
\]
where $a_1g_1+\cdots+a_ng_n$ belongs to $\mathfrak{M}$ and $a_1h_1+\cdots+a_nh_n$ to $\mathfrak{R}-\mathfrak{M}$. Hence
\begin{align*}
P_{\mathfrak{M}}(a_1f_1+\cdots+a_nf_n)&=a_1g_1+\cdots+a_ng_n\\
&=a_1P_{\mathfrak{M}}f_1+\cdots+a_nP_{\mathfrak{M}}f_n.
\end{align*}
This is the first formula. Second, let
\[
f=g'+h',\quad g=g''+h'',\qquad
(g',g''\in\mathfrak{M};\ h',h''\in\mathfrak{R}-\mathfrak{M}).
\]
Then $g',g''$ are orthogonal to $h',h''$, so
\[
(g',g)=(g',g''+h'')=(g',g'')=(g'+h',g'')=(f,g''),
\]
that is, $(P_{\mathfrak{M}}f,g)=(f,P_{\mathfrak{M}}g)$. This is the second formula. Finally, $P_{\mathfrak{M}}f$ belongs to $\mathfrak{M}$, so $P_{\mathfrak{M}}f=P_{\mathfrak{M}}f+0$ is the decomposition of $P_{\mathfrak{M}}f$ guaranteed by \hyperref[thm:II.10]{Theorem 10}; that is, $P_{\mathfrak{M}}(P_{\mathfrak{M}}f)=P_{\mathfrak{M}}f$. This is the third formula.

% Source PDF page 48 / original page 41.
\hypertarget{origpage-41}{}%
The equations $P_{\mathfrak{M}}f=f$ and $P_{\mathfrak{M}}f=0$ mean, respectively, that in the decomposition $f=g+h$, with $g$ in $\mathfrak{M}$ and $h$ in $\mathfrak{R}-\mathfrak{M}$ (\hyperref[thm:II.10]{Theorem 10}), we must have $f=g$, $h=0$, or $g=0$, $f=h$; that is, $f$ belongs to $\mathfrak{M}$ or to $\mathfrak{R}-\mathfrak{M}$, respectively. These are the fifth and sixth assertions. By definition all $P_{\mathfrak{M}}f$ belong to $\mathfrak{M}$, and each $f'$ in $\mathfrak{M}$ equals some $P_{\mathfrak{M}}f$, for example, by what we have just said, $P_{\mathfrak{M}}f'$. This is the fourth assertion.

Observe also that the second and third formulas of this theorem give
\[
(P_{\mathfrak{M}}f,P_{\mathfrak{M}}g)=(f,P_{\mathfrak{M}}P_{\mathfrak{M}}g)=(f,P_{\mathfrak{M}}g)=(P_{\mathfrak{M}}f,g).
\]
We now wish to define the projection operators $P_{\mathfrak{M}}$ independently of the $\mathfrak{M}$.

\leavevmode\phantomsection\textbf{Theorem 12.}\label{thm:II.12} An everywhere defined operator $E$, in the sense explained before \hyperref[thm:II.11]{Theorem 11}, is a projection operator, that is, $E=P_{\mathfrak{M}}$ for a suitable closed linear manifold $\mathfrak{M}$, if and only if it has the properties
\[
(Ef,g)=(f,Eg),\qquad E^2=E
\]
(see the remark to \hyperref[thm:II.11]{Theorem 11}). Moreover, $\mathfrak{M}$ is then uniquely determined by $E$, by \hyperref[thm:II.11]{Theorem 11}.

\emph{Proof.} The necessity of these conditions and the determination of $\mathfrak{M}$ by $E$ are clear from \hyperref[thm:II.11]{Theorem 11}. We need only show that if $E$ has the stated properties, there is a closed linear manifold $\mathfrak{M}$ with $E=P_{\mathfrak{M}}$.

Let $\mathfrak{M}$ be the closed linear manifold spanned by all $Ef$. The vector $g-Eg$ is orthogonal to all $Ef$, since
\[
(Ef,g-Eg)=(Ef,g)-(Ef,Eg)=(Ef,g)-(E^2f,g)=0.
\]
The elements of $\mathfrak{R}$ orthogonal to $g-Eg$ form a closed linear manifold containing every $Ef$, and therefore all of $\mathfrak{M}$. Hence $g-Eg$ belongs to $\mathfrak{R}-\mathfrak{M}$. Thus the decomposition of $g$ associated with $\mathfrak{M}$ is $g=Eg+(g-Eg)$, so $P_{\mathfrak{M}}g=Eg$, where $g$ was arbitrary. This proves everything.

If $\mathfrak{M}=\mathfrak{R}$ or $[0]$, then $\mathfrak{R}-\mathfrak{M}=[0]$ or $\mathfrak{R}$, respectively. Thus $f=f+0$ or $f=0+f$ is the decomposition according to \hyperref[thm:II.11]{Theorem 11}, and $P_{\mathfrak{M}}f=f$ or $0$. We call the everywhere defined operator specified by $Rf=f$ the operator $1$, and that specified by $Rf=0$ the operator $0$. Thus $P_{\mathfrak{R}}=1$, $P_{[0]}=0$. Furthermore, the decomposition $f=g+h$ associated with $\mathfrak{M}$, where $g$ is in $\mathfrak{M}$ and $h$ in $\mathfrak{R}-\mathfrak{M}$, can clearly also be used for $\mathfrak{R}-\mathfrak{M}$ in the form $f=h+g$. For since $g$ belongs to $\mathfrak{M}$, it is orthogonal to every element of $\mathfrak{R}-\mathfrak{M}$, and hence belongs to $\mathfrak{R}-(\mathfrak{R}-\mathfrak{M})$. Consequently $P_{\mathfrak{M}}f=g$ and $P_{\mathfrak{R}-\mathfrak{M}}f=h=f-g$; that is, $P_{\mathfrak{R}-\mathfrak{M}}f=f-P_{\mathfrak{M}}f$. We express this fact, $P_{\mathfrak{R}-\mathfrak{M}}f=1f-P_{\mathfrak{M}}f$, symbolically as $P_{\mathfrak{R}-\mathfrak{M}}=1-P_{\mathfrak{M}}$. For addition, subtraction, and multiplication of operators, see the discussion of \hyperref[thm:II.14]{Theorem 14}.

% Source PDF page 49 / original page 42.
\hypertarget{origpage-42}{}%
Notice that we readily saw above that $\mathfrak{M}$ is a subset of $\mathfrak{R}-(\mathfrak{R}-\mathfrak{M})$; a direct proof that the two sets are equal would probably be more involved. It follows immediately, however, from
\[
P_{\mathfrak{R}-(\mathfrak{R}-\mathfrak{M})}=1-P_{\mathfrak{R}-\mathfrak{M}}=1-(1-P_{\mathfrak{M}})=P_{\mathfrak{M}}.
\]
The preceding discussion also shows that if $E$ is a projection operator, so is $1-E$; and since $1-(1-E)=E$, the converse holds as well.

\leavevmode\phantomsection\textbf{Theorem 13.}\label{thm:II.13} We always have
\[
\|Ef\|^2=(Ef,f),\qquad\|Ef\|\leq\|f\|.
\]
The conditions $\|Ef\|=0$ and $\|Ef\|=\|f\|$ characterize the $f$ in $\mathfrak{R}-\mathfrak{M}$ and $\mathfrak{M}$, respectively.

\emph{Remark.} In particular,
\[
\|Ef-Eg\|=\|E(f-g)\|\leq\|f-g\|;
\]
that is, the operator $E$ is continuous. Compare the discussion after \hyperref[thm:II.2]{Theorem 2} in \hyperref[sec:II.1]{II.1}.

\emph{Proof.} By the observation after \hyperref[thm:II.11]{Theorem 11},
\[
\|Ef\|^2=(Ef,Ef)=(Ef,f).
\]
Since $1-E$ is also a projection operator,
\begin{align*}
\|Ef\|^2+\|f-Ef\|^2&=\|Ef\|^2+\|(1-E)f\|^2\\
&=(Ef,f)+((1-E)f,f)=(f,f)=\|f\|^2.
\end{align*}
Both summands are nonnegative, so both are at most $\|f\|^2$; in particular $\|Ef\|^2\leq\|f\|^2$, $\|Ef\|\leq\|f\|$. We know from \hyperref[thm:II.11]{Theorem 11} that $\|Ef\|=0$, or $Ef=0$, expresses membership in $\mathfrak{R}-\mathfrak{M}$. By the relation above, $\|Ef\|=\|f\|$ means $\|f-Ef\|=0$, $Ef=f$, and hence, by \hyperref[thm:II.11]{Theorem 11}, that $f$ belongs to $\mathfrak{M}$.

For two operators $R,S$, we understand by $R\pm S$, $aR$, with $a$ a complex number, and $RS$ the operators defined by
\[
(R\pm S)f=Rf\pm Sf,\qquad(aR)f=a\cdot Rf,\qquad(RS)f=R(Sf),
\]
and naturally put
\[
R^0=1,\quad R^1=R,\quad R^2=RR,\quad R^3=RRR,\quad\ldots.
\]
The rules governing these operations are easy to establish. For $R\pm S$ and $aR$, all the elementary laws valid for numbers are readily verified, but this is not so for $RS$. The distributive laws $(R\pm S)T=RT\pm ST$ and $R(S\pm T)=RS\pm RT$ do hold, as can be checked; the latter, however, requires $R$ to be linear, as in the remark to \hyperref[thm:II.11]{Theorem 11} and the discussion in the next section. The associative law $(RS)T=R(ST)=RST$ also holds, but the commutative law $RS=SR$ does not hold in general: $(RS)f=R(Sf)$ and $(SR)f=S(Rf)$ need not be equal.

% Source PDF page 50 / original page 43.
\hypertarget{origpage-43}{}%
If it does hold for two particular operators $R,S$, they are said to commute. Thus, for example, $0$ and $1$ commute with all everywhere defined $R$:
\[
R0=0R=0,\qquad R1=1R=R;
\]
so do $R^m$ and $R^n$, since $R^mR^n=R^{m+n}$ and thus does not depend on the order of $m,n$.

\leavevmode\phantomsection\textbf{Theorem 14.}\label{thm:II.14} Let $E,F$ be projection operators, associated with the closed linear manifolds $\mathfrak{M},\mathfrak{N}$. Their product $EF$ is also a projection operator if and only if $E,F$ commute, that is, $EF=FE$. In that case it is the projection onto the closed linear manifold $\mathfrak{P}$ consisting of the common elements of $\mathfrak{M},\mathfrak{N}$. The sum $E+F$ is a projection operator if and only if $EF=0$, or equivalently $FE=0$. This means that all of $\mathfrak{M}$ is orthogonal to all of $\mathfrak{N}$; $E+F$ is then the projection onto $\mathfrak{M}+\mathfrak{N}=[\mathfrak{M},\mathfrak{N}]$, which in this case equals $\{\mathfrak{M},\mathfrak{N}\}$. The difference $E-F$ is a projection operator if and only if $EF=F$, or equivalently $FE=F$. This means that $\mathfrak{N}$ is a subset of $\mathfrak{M}$; $E-F$ is then the projection onto $\mathfrak{M}-\mathfrak{N}$.

\emph{Proof.} For $EF$, we must check the two conditions of \hyperref[thm:II.12]{Theorem 12}:
\[
(EFf,g)=(f,EFg),\qquad(EF)^2=EF.
\]
Since $(EFf,g)=(Ff,Eg)=(f,FEg)$, the first condition means
\[
(f,EFg)=(f,FEg),\qquad(f,(EF-FE)g)=0.
\]
Since this holds for all $f$, $(EF-FE)g=0$; and since this holds for all $g$, $EF-FE=0$, $EF=FE$. Commutativity is thus necessary and sufficient already for the first condition, but it also implies the second:
\[
(EF)^2=EFEF=EEFF=E^2F^2=EF.
\]
Since $E+F$ always satisfies $((E+F)f,g)=(f,(E+F)g)$, because $E,F$ do, only the second condition, $(E+F)^2=E+F$, needs checking. In view of
\[
(E+F)^2=E^2+F^2+EF+FE=(E+F)+(EF+FE),
\]
it simply means $EF+FE=0$. Now if $EF=0$, then $EF$ is a projection operator, and what was proved above gives $EF=FE$, hence $EF+FE=0$. Conversely, $EF+FE=0$ implies
\begin{align*}
E(EF+FE)&=E^2F+EFE=EF+EFE=0,\\
E(EF+FE)E&=E^2FE+EFE^2=EFE+EFE=2\cdot EFE=0.
\end{align*}
Thus $EFE=0$, and therefore $EF=0$. Consequently $EF=0$ is necessary and sufficient, as is $FE=0$, since $E,F$ play identical roles.

The operator $E-F$ is a projection operator if and only if $1-(E-F)=(1-E)+F$ is one. Since $1-E$ and $F$ are projection operators, what we have just proved shows that this is characterized by $(1-E)F=0$, $F-EF=0$, $EF=F$, or equally by $F(1-E)=0$, $F-FE=0$, $FE=F$.

% Source PDF page 51 / original page 44.
\hypertarget{origpage-44}{}%
We still have to prove the assertions about $\mathfrak{M},\mathfrak{N}$, where $E=P_{\mathfrak{M}}$, $F=P_{\mathfrak{N}}$. First let $EF=FE$. Every $EFf=FEf$ then belongs to both $\mathfrak{M}$ and $\mathfrak{N}$, hence to $\mathfrak{P}$. For every $g$ in $\mathfrak{P}$ we have $Eg=Fg=g$, so $EFg=Eg=g$; that is, $g$ has the form $EFf$. Thus $\mathfrak{P}$ is the range of $EF$, and \hyperref[thm:II.11]{Theorem 11} gives $EF=P_{\mathfrak{P}}$. Second, let $EF=0$, hence also $FE=0$. Every $(E+F)f=Ef+Ff$ belongs to $\{\mathfrak{M},\mathfrak{N}\}$, and every $g$ in $\{\mathfrak{M},\mathfrak{N}\}$ equals $h+j$, with $h$ in $\mathfrak{M}$ and $j$ in $\mathfrak{N}$. Hence $Eh=h$, $Fh=FEh=0$, $Fj=j$, $Ej=EFj=0$, so
\[
(E+F)(h+j)=Eh+Fh+Ej+Fj=h+j,\qquad(E+F)g=g.
\]
Thus $g$ has the form $(E+F)f$. Consequently $\{\mathfrak{M},\mathfrak{N}\}$ is the range of $E+F$; but $E+F$ is a projection operator, so $\{\mathfrak{M},\mathfrak{N}\}$ is its associated closed linear manifold (\hyperref[thm:II.11]{Theorem 11}). Since $\{\mathfrak{M},\mathfrak{N}\}$ is closed, it equals $[\mathfrak{M},\mathfrak{N}]=\mathfrak{M}+\mathfrak{N}$. Third, let $EF=F$, hence also $FE=F$. Then $E=P_{\mathfrak{M}}$ and $1-F=P_{\mathfrak{R}-\mathfrak{N}}$, so $E-F=E-EF=E(1-F)$ equals $P_{\mathfrak{P}}$, where $\mathfrak{P}$ is the common part of $\mathfrak{M}$ and $\mathfrak{R}-\mathfrak{N}$, that is, $\mathfrak{M}-\mathfrak{N}$.

Finally, $EF=0$ means that $(EFf,g)=0$, and hence $(Ff,Eg)=0$, always holds; that is, the whole of $\mathfrak{M}$ is orthogonal to the whole of $\mathfrak{N}$. And $EF=F$ means $F(1-E)=0$: all of $\mathfrak{N}$ is orthogonal to $\mathfrak{R}-\mathfrak{M}$, or equivalently, $\mathfrak{N}$ is a subset of $\mathfrak{R}-(\mathfrak{R}-\mathfrak{M})=\mathfrak{M}$.

If $\mathfrak{N}$ is a subset of $\mathfrak{M}$, we shall also say that $F=P_{\mathfrak{N}}$ is a part of $E=P_{\mathfrak{M}}$, and write $E\geq F$ or $F\leq E$. This means $EF=F$, or equivalently $FE=F$, and entails commutativity. Either by considering $\mathfrak{M},\mathfrak{N}$ or by direct calculation, one sees that $0\leq E\leq1$ always holds; that $E\leq F$, $F\leq E$ imply $E=F$; and that $E\leq F$, $F\leq G$ imply $E\leq G$. Thus our relation $\leq$ has the properties of an ordering by magnitude. Notice also that $E\leq F$, $1-E\geq1-F$, and orthogonality of $E$ to $1-F$ are all equivalent. Furthermore, orthogonality of $E,F$ implies that of $E',F'$ if $E'\leq E$, $F'\leq F$. If $\mathfrak{M},\mathfrak{N}$ are orthogonal, we also call $E,F$ orthogonal; this means $EF=0$, or equivalently $FE=0$. Conversely, if $E,F$ commute, we shall also say that their $\mathfrak{M},\mathfrak{N}$ commute.

\leavevmode\phantomsection\textbf{Theorem 15.}\label{thm:II.15} The relation $E\leq F$ is equivalent to $\|Ef\|\leq\|Ff\|$ for every $f$.

\emph{Proof.} The relation $E\leq F$ implies $E=EF$, so $\|Ef\|=\|EFf\|\leq\|Ff\|$ by \hyperref[thm:II.13]{Theorem 13}. Conversely, this inequality has the consequence that if $Ff=0$, then $\|Ef\|\leq\|Ff\|=0$, $Ef=0$. Since $F(1-F)f=(F-F^2)f=0$, we therefore have identically $E(1-F)f=0$, that is, $E(1-F)=E-EF=0$, $E=EF$, hence $E\leq F$.

\leavevmode\phantomsection\textbf{Theorem 16.}\label{thm:II.16} Let $E_1,\ldots,E_k$ be projection operators. Their sum $E_1+\cdots+E_k$ is also a projection operator if and only if all $E_m,E_l$, $m,l=1,\ldots,k$, $m\ne l$, are orthogonal to one another.

% Source PDF page 52 / original page 45.
\hypertarget{origpage-45}{}%
Another necessary and sufficient condition is
\[
\|E_1f\|^2+\cdots+\|E_kf\|^2\leq\|f\|^2
\]
for every $f$. Moreover, $E_1+\cdots+E_k$ is then the projection operator of $\mathfrak{M}_1+\cdots+\mathfrak{M}_k=[\mathfrak{M}_1,\ldots,\mathfrak{M}_k]$, which in this case equals $\{\mathfrak{M}_1,\ldots,\mathfrak{M}_k\}$, where $E_1=P_{\mathfrak{M}_1},\ldots,E_k=P_{\mathfrak{M}_k}$.

\emph{Proof.} The last assertion follows by repeated application of \hyperref[thm:II.14]{Theorem 14}, as does the sufficiency of the first criterion. If the second criterion holds, so does the first. For $m\ne l$ and $E_mf=f$, we have
\[
\begin{aligned}
\|f\|^2+\|E_lf\|^2&=\|E_mf\|^2+\|E_lf\|^2\\
&\leq\|E_1f\|^2+\cdots+\|E_kf\|^2\leq\|f\|^2,
\end{aligned}
\qquad\|E_lf\|^2=0,\quad E_lf=0.
\]
But $E_m(E_mf)=E_mf$ holds identically, so $E_l(E_mf)=0$, that is, $E_lE_m=0$. Finally, the second condition is necessary: if $E_1+\cdots+E_k$ is a projection operator, \hyperref[thm:II.13]{Theorem 13} gives
\begin{align*}
\|E_1f\|^2+\cdots+\|E_kf\|^2&=(E_1f,f)+\cdots+(E_kf,f)\\
&=((E_1+\cdots+E_k)f,f)\\
&=\|(E_1+\cdots+E_k)f\|^2\leq\|f\|^2.
\end{align*}
We therefore have the following logical scheme:
\begin{gather*}
E_1+\cdots+E_k\text{ is a projection operator}\ \Longrightarrow\ \text{second criterion}\\
\Longrightarrow\ \text{first criterion}\ \Longrightarrow\ E_1+\cdots+E_k\text{ is a projection operator}.
\end{gather*}
Thus all three are equivalent.

Lastly, we prove a theorem on the convergence of projection operators.

\leavevmode\phantomsection\textbf{Theorem 17.}\label{thm:II.17} Let $E_1,E_2,\ldots$ be an increasing or decreasing sequence of projection operators: $E_1\leq E_2\leq\cdots$ or $E_1\geq E_2\geq\cdots$. They then converge to a projection operator $E$, in the sense that $E_nf\to Ef$ for every $f$; moreover, all $E_n\leq E$ or all $E_n\geq E$, respectively.

\emph{Proof.} It suffices to examine the second case, since the first is reduced to it by replacing $E_1,E_2,\ldots,E$ by $1-E_1,1-E_2,\ldots,1-E$. Let, then, $E_1\geq E_2\geq\cdots$.

By \hyperref[thm:II.15]{Theorem 15}, $\|E_1f\|^2\geq\|E_2f\|^2\geq\cdots\geq0$, so $\lim_{m\to\infty}\|E_mf\|^2$ exists. Hence for every $\varepsilon>0$ there is an $N=N(\varepsilon)$ such that $|\|E_mf\|^2-\|E_lf\|^2|<\varepsilon$ for $m,l\geq N$. For $m\leq l$ we have $E_m\geq E_l$, and $E_m-E_l$ is a projection operator. Thus
\begin{align*}
\|E_mf\|^2-\|E_lf\|^2&=(E_mf,f)-(E_lf,f)=((E_m-E_l)f,f)\\
&=\|(E_m-E_l)f\|^2=\|E_mf-E_lf\|^2,
\end{align*}
which gives $\|E_mf-E_lf\|<\sqrt\varepsilon$. The sequence $E_1f,E_2f,\ldots$ therefore satisfies Cauchy's convergence condition, and hence has a limit $f^*$, by \hyperref[ax:II.D]{D} of \hyperref[sec:II.1]{II.1}. We thus define an everywhere defined operator by $Ef=f^*$.

% Source PDF page 53 / original page 46.
\hypertarget{origpage-46}{}%
Passing to the limit in $(E_nf,g)=(f,E_ng)$ gives $(Ef,g)=(f,Eg)$, and passing to the limit in $(E_nf,E_ng)=(E_nf,g)$ gives $(Ef,Eg)=(Ef,g)$; hence $(E^2f,g)=(Ef,g)$, $E^2=E$. Thus $E$ is a projection operator. For $l\geq m$ we have $\|E_mf\|\geq\|E_nf\|$; letting $l\to\infty$ gives $\|E_mf\|\geq\|Ef\|$, so $E_m\geq E$ by \hyperref[thm:II.15]{Theorem 15}.\tnote{The source writes $E_n$ in the inequality preceding the limit $l\to\infty$; its index should be $l$. The printed index is retained.}

If $E_1,E_2,\ldots$ are pairwise orthogonal projection operators, then $E_1,E_1+E_2,E_1+E_2+E_3,\ldots$ are all projection operators (\hyperref[thm:II.16]{Theorem 16}), and plainly form an increasing sequence. By \hyperref[thm:II.17]{Theorem 17} they therefore converge to a projection operator which is greater than or equal to all of them and which we may denote by $E_1+E_2+\cdots$. Let $E_1=P_{\mathfrak{M}_1}$, $E_2=P_{\mathfrak{M}_2}$, $\ldots$, $E_1+E_2+\cdots=P_{\mathfrak{M}}$. Since all $E_m\leq E$, $\mathfrak{M}_m$ is a subset of $\mathfrak{M}$, so $\mathfrak{M}$ also contains $[\mathfrak{M}_1,\mathfrak{M}_2,\ldots]=\mathfrak{M}_1+\mathfrak{M}_2+\cdots=\mathfrak{M}'$. Conversely, all $\mathfrak{M}_m$ are subsets of $\mathfrak{M}'$, so $E_m\leq P_{\mathfrak{M}'}=E'$. By continuity, therefore, $E\leq E'$ (compare the reasoning in the proof above), and $\mathfrak{M}$ is a subset of $\mathfrak{M}'$. Hence $\mathfrak{M}=\mathfrak{M}'$, $E=E'$; that is, $\mathfrak{M}=\mathfrak{M}_1+\mathfrak{M}_2+\cdots$, or, in another notation,
\[
P_{\mathfrak{M}_1+\mathfrak{M}_2+\cdots}=P_{\mathfrak{M}_1}+P_{\mathfrak{M}_2}+\cdots.
\]
This concludes our discussion of projection operators.

\section{Operators in Hilbert Space}\label{sec:II.5}
We now understand the geometry of the infinite-dimensional Hilbert space $\mathfrak{R}_\infty$ well enough to turn to its linear operators, that is, the linear mappings of $\mathfrak{R}_\infty$ into itself. For this purpose we must introduce several concepts, some of which have already been anticipated in the preceding section.

We have already dealt with operators in the preceding section. In accordance with what was said there before \hyperref[thm:II.11]{Theorem 11}, we define them as follows.

\leavevmode\phantomsection\textbf{Definition 6.}\label{def:II.6} An operator $R$ is a function defined on a subset of $\mathfrak{R}$ and taking values in $\mathfrak{R}$; that is, a correspondence assigning certain elements $Rf$ of $\mathfrak{R}$ to certain elements $f$ of $\mathfrak{R}$.

Here we admit the spaces $\mathfrak{R}_n$ as well as $\mathfrak{R}_\infty$. Notice that if $\mathfrak{R}_\infty$ is an $F_\Omega$, the operator $R$ is defined for elements of $F_\Omega$, that is, ordinary configuration-space functions, and its values are functions of the same kind. The operators are then so-called ``functions of functions,'' or ``functionals.'' Compare the examples in \hyperref[sec:I.2]{I.2} and \hyperref[sec:I.4]{I.4}. The class of $f$ for which $Rf$ is meaningful, the domain of $R$, need not include all of $\mathfrak{R}$; if it does, $R$ is called everywhere defined. Moreover, the class of $Rf$, the range of $R$, which is the image of its domain under $R$, need by no means be contained in the domain. Thus the existence of $Rf$ need not imply that of $R(Rf)=R^2f$.\origNote{58}

% Source PDF page 54 / original page 47.
\hypertarget{origpage-47}{}%
We have already stated in the preceding section what is meant by $R\pm S$, $aR$, $RS$, $R^m$, where $R,S$ are operators, $a$ a complex number, and $m=0,1,2,\ldots$:
\begin{gather*}
(R\pm S)f=Rf\pm Sf,\qquad(aR)f=a\cdot Rf,\qquad(RS)f=R(Sf),\\
R^0=1,\quad R^1=R,\quad R^2=RR,\quad R^3=RRR,\quad\ldots.
\end{gather*}
In specifying the domains, we must of course remember that the left sides, that is, the operators $R\pm S$, $aR$, $RS$, are meaningful only when the right sides are meaningful. For example, $R\pm S$ is defined only on the common part of the domains of $R$ and $S$, and so forth. If $Rf$ takes different values for distinct arguments, the operator has an inverse $R^{-1}$: $R^{-1}f$ is defined whenever $Rg=f$ has a solution $g$, and is equal to that $g$. The rules of calculation valid for $R\pm S$, $aR$, and $RS$ were already discussed in the preceding section. Here we add only the following observation about domains. The operators declared equal there also have identical domains, whereas operator equations such as $0\cdot R=0$ do not hold as regards domains: $0f$ is always meaningful, but by definition $(0\cdot R)f$ is meaningful only when $Rf$ is, although whenever both sides are meaningful, both equal zero. In contrast, $1\cdot R=R\cdot1=R$, and also $R^m\cdot R^l=R^{m+l}$, hold for the domains as well.

If $R,S$ possess inverses, then $RS$ also has one, and, as is readily seen, $(RS)^{-1}=S^{-1}R^{-1}$. Furthermore, for $a\ne0$, $(aR)^{-1}=\frac1a R^{-1}$. If $R^{-1}$ exists, we can also form the remaining negative powers of $R$:
\[
R^{-2}=R^{-1}R^{-1},\qquad R^{-3}=R^{-1}R^{-1}R^{-1},\quad\ldots.
\]
After these general remarks, we turn to a closer examination of the special classes of operators that will be particularly important for us.

\leavevmode\phantomsection\textbf{Definition 7.}\label{def:II.7} An operator $A$ is called linear if its domain is a linear manifold, that is, if together with $f_1,\ldots,f_k$ it contains $a_1f_1+\cdots+a_kf_k$, and if
\[
A(a_1f_1+\cdots+a_kf_k)=a_1Af_1+\cdots+a_kAf_k.
\]
Henceforth we shall consider only linear operators, and specifically those whose domains are everywhere dense.

For some purposes, this last requirement can replace the requirement that an operator be defined everywhere, which we must abandon in quantum mechanics. The point is important enough to merit closer examination. Consider, for example, the configuration space of Schr\"odinger's wave mechanics, assumed for simplicity to be one-dimensional: $-\infty<q<+\infty$.

% Source PDF page 55 / original page 48.
\hypertarget{origpage-48}{}%
The wave functions are the $\varphi(q)$ with finite $\int_{-\infty}^{+\infty}|\varphi(q)|^2\,dq$; these form a Hilbert space (see \hyperref[sec:II.3]{II.3}). Consider also the operators $q\cdots$ and $\frac{h}{2\pi i}\frac{d}{dq}\cdots$. They are plainly linear operators, but their domains by no means comprise the whole Hilbert space. For $q\cdots$, this is because
\[
\int_{-\infty}^{+\infty}|q\varphi(q)|^2\,dq=\int_{-\infty}^{+\infty}q^2|\varphi(q)|^2\,dq
\]
may very well be infinite even when $\int_{-\infty}^{+\infty}|\varphi(q)|^2\,dq$ is finite, so $q\varphi(q)$ no longer belongs to Hilbert space. For $\frac{h}{2\pi i}\frac{d}{dq}$, this is because there are nondifferentiable functions, as well as functions for which $\int_{-\infty}^{+\infty}|\varphi(q)|^2\,dq$ is finite but
\[
\int_{-\infty}^{+\infty}\left|\frac{h}{2\pi i}\frac{d}{dq}\varphi(q)\right|^2\,dq
=\frac{h^2}{4\pi^2}\int_{-\infty}^{+\infty}\left|\frac{d}{dq}\varphi(q)\right|^2\,dq
\]
is not, for example $|q|^{1/2}e^{-q^2}$ or $e^{-q^2}\sin(e^{q^2})$. The domains are nevertheless everywhere dense. Both operators certainly apply to every $\varphi(q)$ which is nonzero only in a finite interval $-C\leq q\leq C$ and is continuously differentiable everywhere; and this set of functions is everywhere dense.\origNote{59}

We further define:

\leavevmode\phantomsection\textbf{Definition 8.}\label{def:II.8} Two operators $A,A^*$ are called adjoint to one another if they have the same domain and the following relations hold throughout that domain:
\[
(Af,g)=(f,A^*g),\qquad(A^*f,g)=(f,Ag).
\]
Interchanging $f,g$ and taking the complex conjugates of both sides shows that each of these two relations follows from the other. It is also clear that the relation between $A,A^*$ is symmetric, that is, $A^*,A$ are adjoint as well. Thus $A^{**}=A$.

We also observe that $A$ can have only one adjoint $A^*$: if $A$ is adjoint both to $A_1^*$ and to $A_2^*$, then $A_1^*=A_2^*$. Indeed, for every $g$ for which $Ag$ is meaningful,
\[
(A_1^*f,g)=(f,Ag)=(A_2^*f,g).
\]
Since these $g$ form an everywhere dense set, $A_1^*f=A_2^*f$; since this holds generally, $A_1^*=A_2^*$. Thus $A$ determines $A^*$ uniquely, and likewise $A^*$ determines $A$.

It is readily seen that $0$, $1$, and indeed all projection operators $E$ are adjoint to themselves (compare \hyperref[thm:II.12]{Theorem 12}): $0^*$, $1^*$, and $E^*$ exist and equal $0$, $1$, and $E$, respectively. Furthermore, $(aA)^*=\overline aA^*$, and, insofar as $A\pm B$ can be formed at all, that is, their domain is everywhere dense, $(A\pm B)^*=A^*\pm B^*$. With readily determined restrictions concerning the domains, we finally have $(AB)^*=B^*A^*$, since $(ABf,g)=(Bf,A^*g)=(f,B^*A^*g)$, as well as $(A^{-1})^*=A^{*-1}$, since
\[
(A^{-1}f,g)=(A^{-1}f,A^*A^{*-1}g)=(AA^{-1}f,A^{*-1}g)=(f,A^{*-1}g).
\]

% Source PDF page 56 / original page 49.
\hypertarget{origpage-49}{}%
Return now to the example of Schr\"odinger's wave mechanics, this time with a $k$-dimensional configuration space. Hilbert space consists of the $\varphi(q_1,\ldots,q_k)$ with finite $\int_{-\infty}^{\infty}\cdots\int_{-\infty}^{\infty}|\varphi(q_1\cdots q_k)|^2\,dq_1\cdots dq_k$. For the operators $q_l\cdots$ and $\frac{h}{2\pi i}\frac{\partial}{\partial q_l}\cdots$ we have
\[
(q_l)^*=q_l,\qquad
\left(\frac{h}{2\pi i}\frac{\partial}{\partial q_l}\right)^*
=\frac{h}{2\pi i}\frac{\partial}{\partial q_l}.
\]
The first is clear from
\begin{align*}
&\int_{-\infty}^{\infty}\cdots\int_{-\infty}^{\infty}q_l\varphi(q_1\cdots q_k)\overline{\psi(q_1\cdots q_k)}\,dq_1\cdots dq_k\\
&\qquad=\int_{-\infty}^{\infty}\cdots\int_{-\infty}^{\infty}\varphi(q_1\cdots q_k)\overline{q_l\psi(q_1\cdots q_k)}\,dq_1\cdots dq_k.
\end{align*}
The second says that
\begin{align*}
&\int_{-\infty}^{\infty}\cdots\int_{-\infty}^{\infty}\frac{h}{2\pi i}\frac{\partial}{\partial q_l}\varphi(q_1\cdots q_k)\overline{\psi(q_1\cdots q_k)}\,dq_1\cdots dq_k\\
&\qquad=\int_{-\infty}^{\infty}\cdots\int_{-\infty}^{\infty}\varphi(q_1\cdots q_k)\overline{\frac{h}{2\pi i}\frac{\partial}{\partial q_l}\psi(q_1\cdots q_k)}\,dq_1\cdots dq_k,
\end{align*}
that is,
\begin{align*}
\int_{-\infty}^{\infty}\cdots\int_{-\infty}^{\infty}
\biggl\{&\frac{\partial}{\partial q_l}\varphi(q_1\cdots q_k)\overline{\psi(q_1\cdots q_k)}\\
&+\varphi(q_1\cdots q_k)\frac{\partial}{\partial q_l}\overline{\psi(q_1\cdots q_k)}\biggr\}\,dq_1\cdots dq_k=0,
\end{align*}
or
\begin{align*}
\lim_{A\to+\infty,\,B\to+\infty}
\int_{-\infty}^{\infty}\cdots\int_{-\infty}^{\infty}
\bigl[\varphi(q_1\cdots q_k)\overline{\psi(q_1\cdots q_k)}\bigr]_{q_l=-B}^{q_l=+A}\\
{}\times dq_1\cdots dq_{l-1}\,dq_{l+1}\cdots dq_k=0.
\end{align*}
The limit must exist, since the convergence of all the integrals
\[
\int_{-\infty}^{\infty}\cdots\int_{-\infty}^{\infty}\cdots\,dq_1\cdots dq_k
\]
is assured: $\varphi,\psi,\frac{\partial}{\partial x_l}\varphi,\frac{\partial}{\partial x_l}\psi$ belong to Hilbert space. It remains only to show that the limit vanishes. If it were nonzero, the certainly existing limit of
\[
\int_{-\infty}^{\infty}\cdots\int_{-\infty}^{\infty}\varphi(q_1\cdots q_k)\overline{\psi(q_1\cdots q_k)}\,dq_1\cdots dq_{l-1}\,dq_{l+1}\cdots dq_k
\]
as $q_l\to+\infty$ or as $q_l\to-\infty$ would be nonzero. This is incompatible with the absolute convergence of
\[
\int_{-\infty}^{\infty}\cdots\int_{-\infty}^{\infty}\varphi(q_1\cdots q_k)\overline{\psi(q_1\cdots q_k)}\,dq_1\cdots dq_{l-1}\,dq_l\,dq_{l+1}\cdots dq_k,
\]
since $\varphi,\psi$ belong to Hilbert space.

% Source PDF page 57 / original page 50.
\hypertarget{origpage-50}{}%
If $A$ is the integral operator
\[
A\varphi(q_1\cdots q_k)=\int_{-\infty}^{\infty}\cdots\int_{-\infty}^{\infty}K(q_1\cdots q_k;q'_1\cdots q'_k)\varphi(q'_1\cdots q'_k)\,dq'_1\cdots dq'_k,
\]
one readily sees that $A^*$ too is an integral operator, but with kernel
\[
\overline{K(q'_1\cdots q'_k;q_1\cdots q_k)}
\quad\text{instead of}\quad K(q_1\cdots q_k;q'_1\cdots q'_k).
\]

Now consider matrix theory, where Hilbert space consists of all sequences $x_1,x_2,\ldots$ with finite $\sum_{n=1}^\infty|x_n|^2$. A linear operator $A$ transforms $\{x_1,x_2,\ldots\}$ into $\{y_1,y_2,\ldots\}$:
\[
A\{x_1,x_2,\ldots\}=\{y_1,y_2,\ldots\}.
\]
Because $A$ is linear, $y_1,y_2,\ldots$ must depend linearly on $x_1,x_2,\ldots$:\origNote{60}
\[
y_\mu=\sum_{\nu=1}^\infty a_{\mu\nu}x_\nu.
\]
Thus $A$ is characterized by specifying the matrix $a_{\mu\nu}$. It is immediately apparent that the matrix associated with $A^*$ is $\overline{a_{\nu\mu}}$, the conjugate transpose matrix.\origNote{60}

The analogy with matrix theory just explained suggests the definition of a Hermitian operator which we shall now give. At the same time we introduce two other concepts important for our later purposes.

\leavevmode\phantomsection\textbf{Definition 9.}\label{def:II.9} The operator $A$ is called Hermitian if $A^*=A$. It is also called definite if, in addition, $(Af,f)\geq0$ always holds.\origNote{61}\tnote{The historical term \emph{definite} (German \emph{definit}) permits zero values here. In modern terminology this condition is usually called \emph{nonnegative} or \emph{positive semidefinite}; strict positive definiteness is not intended.} The operator $U$ is called unitary if $UU^*=U^*U=1$.\origNote{62}

For unitary operators we therefore have $U^*=U^{-1}$. By definition,
\[
(Uf,Ug)=(U^*Uf,g)=(f,g),
\]
so in particular, for $f=g$, $\|Uf\|=\|f\|$. Conversely, this last property implies unitarity if $U$ is defined everywhere and takes every element of the space as a value (see \hyperlink{note-62}{note 62}). This is proved as follows. First,
\[
\|Uf\|=\|f\|,\qquad\text{that is, }(Uf,Uf)=(f,f),\quad(U^*Uf,f)=(f,f).
\]
Replacing $f$ by $(f+g)/2$ and $(f-g)/2$, respectively, and subtracting gives, as is readily calculated, $\operatorname{Re}(Uf,Ug)=\operatorname{Re}(f,g)$.

% Source PDF page 58 / original page 51.
\hypertarget{origpage-51}{}%
Replacing $f$ here by $if$ gives $\operatorname{Im}$ instead of $\operatorname{Re}$. Thus, in general,
\[
(Uf,Ug)=(f,g),\qquad\text{that is, }(U^*Uf,g)=(f,g).
\]
For fixed $f$ this holds for all $g$, so $U^*Uf=f$; since it holds for all $f$, $U^*U=1$. It remains to show $UU^*=1$. For every $f$ there is a $g$ with $Ug=f$, and then
\[
UU^*f=UU^*\cdot Ug=U\cdot U^*Ug=Ug=f,\qquad\text{hence }UU^*=1.
\]
Since linearity gives
\[
\|Uf-Ug\|=\|U(f-g)\|=\|f-g\|,
\]
every unitary $U$ is continuous. Hermitian operators need by no means be continuous; for example, precisely the operators important for quantum mechanics, $q\cdots$ and $\frac{h}{2\pi i}\frac{d}{dq}\cdots$, are discontinuous.\origNote{63}

Our formal rules for $A^*$ give immediately the following results. If $U,V$ are unitary, then so are $U^{-1}$ and $UV$, and therefore all powers of $U$. If $A,B$ are Hermitian, then so are $A\pm B$; but $aA$ is Hermitian only for real $a$, unless $A=0$, and $AB$ only if $A,B$ commute, that is, $AB=BA$. We also know that all projection operators, in particular $0$ and $1$, are Hermitian, as are the operators $q_l\cdots$, $\frac{h}{2\pi i}\frac{\partial}{\partial q_l}\cdots$ of Schr\"odinger's theory. If $A$ is Hermitian, so are all its powers, including $A^{-1}$ if it exists, and all polynomials in $A$ with real coefficients. It is also noteworthy that if $A$ is Hermitian and $X$ arbitrary, $XAX^*$ is Hermitian:
\[
(XAX^*)^*=X^{**}A^*X^*=XAX^*.
\]
Examples are all $XX^*$, with $A=1$, and $X^*X$, with $X^*$ in place of $X$. For unitary $U$, $UAU^{-1}$ is Hermitian, since $U^{-1}=U^*$.

Continuity is fundamental for operators, just as it is for the scalar-valued functions familiar from analysis. We therefore give several criteria for the continuity of a linear operator.

\leavevmode\phantomsection\textbf{Theorem 18.}\label{thm:II.18} A linear operator $R$ is continuous everywhere if it is continuous at $f=0$. A necessary and sufficient condition for continuity is the existence of a constant $C$ such that $\|Rf\|\leq C\|f\|$ holds generally. This condition is itself equivalent to the general validity of
\[
|(Rf,g)|\leq C\|f\|\|g\|.
\]
For Hermitian $R$, this need only be required for $f=g$: $|(Rf,f)|\leq C\|f\|^2$, or, since $(Rf,f)$ is real (\hyperlink{note-61}{note 61}),
\[
-C\|f\|^2\leq(Rf,f)\leq C\|f\|^2.
\]

% Source PDF page 59 / original page 52.
\hypertarget{origpage-52}{}%
\emph{Remark.} The concept of continuity for operators originates with Hilbert,\origNote{64} who called it ``boundedness'' and defined it by our penultimate criterion. If only one of the two inequalities in the last criterion holds generally, $R$ is called semibounded, from below or from above, respectively. Thus, for example, every definite $R$ is semibounded from below, with $C=0$.

\emph{Proof.} Continuity at $f=0$ means that for every $\varepsilon>0$ there is a $\delta>0$ such that $\|f\|<\delta$ implies $\|Rf\|<\varepsilon$. Then
\[
\|f-f_0\|<\delta\quad\Longrightarrow\quad\|R(f-f_0)\|=\|Rf-Rf_0\|<\varepsilon;
\]
that is, $R$ is also continuous at $f=f_0$, hence everywhere.

If $\|Rf\|\leq C\|f\|$, with, of course, $C>0$, we have continuity: we may put $\delta=\varepsilon/C$. Conversely, if continuity holds, we may determine the $\delta$ corresponding to $\varepsilon=1$ and put $C=2/\delta$. For $\|Rf\|\leq C\|f\|$ holds in any event for $f=0$. For $f\ne0$, hence $\|f\|>0$, put $g=\frac{\delta/2}{\|f\|}f$. Then $\|g\|=\delta/2$, so
\[
\|Rg\|=\frac{\delta/2}{\|f\|}\|Rf\|<1,
\qquad\|Rf\|<\frac{\|f\|}{\delta/2}=C\|f\|.
\]
From $\|Rf\|\leq C\|f\|$ follows $|(Rf,g)|\leq\|Rf\|\|g\|\leq C\|f\|\|g\|$. Conversely, setting $g=Rf$ in $|(Rf,g)|\leq C\|f\|\|g\|$ gives $\|Rf\|^2\leq C\|f\|\|Rf\|$, hence $\|Rf\|\leq C\|f\|$. For Hermitian $R$, it remains to show that
\[
|(Rf,f)|\leq C\|f\|^2\quad\Longrightarrow\quad |(Rf,g)|\leq C\|f\|\|g\|.
\]
Substitution of $(f+g)/2$ and $(f-g)/2$ for $f$ gives\origNote{65}
\begin{align*}
|\operatorname{Re}(Rf,g)|
&=\left|\left(R\frac{f+g}{2},\frac{f+g}{2}\right)-\left(R\frac{f-g}{2},\frac{f-g}{2}\right)\right|\\
&\leq C\left(\left\|\frac{f+g}{2}\right\|^2+\left\|\frac{f-g}{2}\right\|^2\right)
=C\frac{\|f\|^2+\|g\|^2}{2}.
\end{align*}
Now, as in the proof of \hyperref[thm:II.1]{Theorem 1}, replace $f,g$ by $af,a^{-1}g$, with $a>0$. Minimizing the right side gives $|\operatorname{Re}(Bf,g)|\leq C\|f\|\|g\|$.\tnote{The source prints $B$ here; the operator in the surrounding argument is $R$. The printed letter is retained.} Next replace $f$ by $e^{i\alpha}f$, with real $\alpha$; the maximum of the left side gives
\[
|(Rf,g)|\leq C\|f\|\|g\|.
\]
So far, however, this has been proved only when $Rg$ is defined. Such $g$ form an everywhere dense set, and $Rg$ does not occur in the final result, which therefore extends to all $g$ by continuity.

We prove one more theorem on definite operators.

% Source PDF page 60 / original page 53.
\hypertarget{origpage-53}{}%
\leavevmode\phantomsection\textbf{Theorem 19.}\label{thm:II.19} If $R$ is Hermitian and definite, then
\[
|(Rf,g)|\leq\sqrt{(Rf,f)(Rg,g)}
\]
always holds. In that case $(Rf,f)=0$ implies $Rf=0$.

\emph{Proof.} The inequality follows from the general validity of $(Rf,f)\geq0$, that is, definiteness, in precisely the same way that Schwarz's inequality $|(f,g)|\leq\sqrt{(f,f)(g,g)}$, or $|(f,g)|\leq\|f\|\|g\|$, was proved from $(f,f)\geq0$ in \hyperref[thm:II.1]{Theorem 1}. If now $(Rf,f)=0$, this inequality also gives $(Rf,g)=0$ whenever $Rg$ is meaningful. Thus it holds on an everywhere dense set of $g$'s, and by continuity for all $g$. Hence $Rf=0$.

In conclusion, we call attention once more to the important concept of commutativity of two operators $R,S$, that is, the relation $RS=SR$.

From $RS=SR$ follows
\[
S\cdots SSR=S\cdots SRS=S\cdots RSS=\cdots=RS\cdots SS;
\]
that is, $R$ and $S^n$ commute, $n=1,2,\ldots$. Since $R1=1R=R$ and $S^0=1$, this holds also for $n=0$. If $S^{-1}$ exists, then $S^{-1}\cdot SR\cdot S^{-1}=S^{-1}\cdot RS\cdot S^{-1}$. In view of
\begin{align*}
S^{-1}\cdot SR\cdot S^{-1}&=S^{-1}S\cdot RS^{-1}=RS^{-1},\\
S^{-1}\cdot RS\cdot S^{-1}&=S^{-1}R\cdot SS^{-1}=S^{-1}R,
\end{align*}
we therefore also have $RS^{-1}=S^{-1}R$. Hence $n=-1$ is admissible too, and consequently all $n=-2,-3,\ldots$. Thus $R$ commutes with all powers of $S$. A second application shows that every power of $R$ commutes with every power of $S$. If $R$ commutes with $S,T$, it plainly also commutes with all $aS$, and with $S\pm T$ and $ST$. Together with the preceding result, this implies that if $R,S$ commute, then all polynomials in $R$ commute with all polynomials in $S$. In particular, putting $R=S$ shows that all polynomials in $R$ commute with one another.

\section{The Eigenvalue Problem}\label{sec:II.6}
We can now address in abstract Hilbert space the problem that was central to quantum mechanics in the special realizations $F_Z$ and $F_\Omega$: solving equations $\hyperref[eq:I.3.E1]{\mathbf{E}_1}$ and $\hyperref[eq:I.3.E2]{\mathbf{E}_2}$ of \hyperref[sec:I.3]{I.3}. This is called the eigenvalue problem; we must give it a new, unified formulation.

In \hyperref[sec:I.3]{I.3}, both $\hyperref[eq:I.3.E1]{\mathbf{E}_1}$ and $\hyperref[eq:I.3.E2]{\mathbf{E}_2}$ required us to find all solutions $\varphi\ne0$ of
\begin{equation}
\mathsf H\varphi=\lambda\varphi,\tag{E}\label{eq:II.6.E}
\end{equation}
where $\mathsf H$ is the Hermitian operator corresponding to the Hamiltonian function (see that section), $\varphi$ is an element of Hilbert space, and $\lambda$ a real number: $\mathsf H$ is given, while $\varphi,\lambda$ are sought.

% Source PDF page 61 / original page 54.
\hypertarget{origpage-54}{}%
There were, however, requirements on how many solutions had to be found. We needed enough solutions to ensure that:

1. In matrix theory, the solutions
\[
\varphi_1=\{s_{11},s_{12},\ldots\},\quad\varphi_2=\{s_{21},s_{22},\ldots\},\quad\ldots
\]
---we are in $F_Z$---could be assembled into a matrix $\mathsf S=\{s_{\mu\nu}\}$ possessing an inverse $\mathsf S^{-1}$; see \hyperref[sec:I.3]{I.3}.

2. In wave theory, every wave function $\varphi(q_1\cdots q_f)$, which need not itself be a solution, could be expanded in a series of the solutions
\[
\varphi_1=\varphi_1(q_1\cdots q_f),\quad\varphi_2=\varphi_2(q_1\cdots q_f),\quad\ldots,
\]
where $\varphi_1,\varphi_2,\ldots$ may belong to different $\lambda$'s:
\[
\varphi(q_1\cdots q_f)=\sum_{n=1}^\infty c_n\varphi_n(q_1\cdots q_f).
\]
Although this latter point was not mentioned in \hyperref[sec:I.3]{I.3}, the requirement is indispensable for the further development of wave theory, especially for Schr\"odinger's ``perturbation theory.''\origNote{66}

Now 1 amounts to the same thing as 2. For the matrix $\mathsf S$ transforms $\{1,0,0,\ldots\}$, $\{0,1,0,\ldots\}$, $\ldots$ into
\[
\{s_{11},s_{12},s_{13},\ldots\},\quad\{s_{21},s_{22},s_{23},\ldots\},\quad\ldots,
\]
respectively, and therefore transforms the whole Hilbert space $\mathfrak{R}_\infty$ into the closed linear manifold spanned by $\varphi_1,\varphi_2,\ldots$. For $\mathsf S^{-1}$ to exist, this latter manifold must therefore also equal $\mathfrak{R}_\infty$. But 2 says precisely the same thing: it too requires that every $\varphi$ be approximated arbitrarily closely by linear combinations of $\varphi_1,\varphi_2,\ldots$.\origNote{67} Let us make the implications of this condition precise and reexamine the properties of equation \hyperref[eq:II.6.E]{E} with the formal apparatus now available.

First, since $\varphi\ne0$ is required and $a\varphi$ is a solution whenever $\varphi$ is, it suffices to consider solutions with $\|\varphi\|=1$. Second, we need not require $\lambda$ to be real; this follows from $\mathsf H\varphi=\lambda\varphi$:
\[
(\mathsf H\varphi,\varphi)=(\lambda\varphi,\varphi)=\lambda(\varphi,\varphi)=\lambda
\]
(see \hyperref[sec:II.5]{II.5}, \hyperlink{note-61}{note 61}). Third, solutions $\varphi_1,\varphi_2$ belonging to different $\lambda_1,\lambda_2$ are orthogonal:
\[
(\mathsf H\varphi_1,\varphi_2)=\lambda_1(\varphi_1,\varphi_2),\qquad
(\mathsf H\varphi_1,\varphi_2)=(\varphi_1,\mathsf H\varphi_2)=\lambda_2(\varphi_1,\varphi_2).
\]
Thus $\lambda_1(\varphi_1,\varphi_2)=\lambda_2(\varphi_1,\varphi_2)$ and $\lambda_1\ne\lambda_2$ give $(\varphi_1,\varphi_2)=0$.

Let $\lambda_1,\lambda_2,\ldots$ now be all the distinct $\lambda$'s for which \hyperref[eq:II.6.E]{E} is solvable. For each $\lambda$ with solvable $\mathsf H\varphi=\lambda\varphi$, choose a solution $\varphi_\lambda$ of norm one. By what has just been said, these $\varphi_\lambda$ form an orthonormal system. By \hyperref[sec:II.2]{II.2}, \hyperref[thm:II.3infty]{Theorem $3^{(\infty)}$}, this system is either finite or an infinite sequence, so the $\lambda$'s too can be written as a sequence, possibly terminating.

% Source PDF page 62 / original page 55.
\hypertarget{origpage-55}{}%
For each $\lambda=\lambda_\rho$, all solutions of $\mathsf H\varphi=\lambda\varphi$ form a linear manifold, and in fact a closed one.\origNote{68} By \hyperref[thm:II.9]{Theorem 9}, therefore, there is an orthonormal system $\varphi_{\rho,1},\ldots,\varphi_{\rho,\nu_\rho}$ of such solutions spanning precisely this closed linear manifold. The number $\nu_\rho$ is plainly the maximum number of linearly independent solutions for $\lambda=\lambda_\rho$, and is therefore also called the multiplicity of the eigenvalue $\lambda_\rho$. Here $\nu=1,2,\ldots,\infty$; $\nu=\infty$ can occur, for example for $\mathsf H=1$, $\nu=1$.\tnote{The source prints $\nu=1$ at the end of this example. For the identity operator, the eigenvalue is $\lambda=1$ and its multiplicity is infinite. The printed symbol is retained.} By the preceding discussion, the $\varphi_{\rho,1},\ldots,\varphi_{\rho,\nu_\rho}$ for different $\rho$ are also orthogonal to one another, so all the
\[
\varphi_{\rho,\nu}\qquad(\rho=1,2,\ldots;\ \nu=1,\ldots,\nu_\rho)
\]
together form an orthonormal system. Its construction shows that it spans the same closed linear manifold as all solutions $\varphi$ of \hyperref[eq:II.6.E]{E}.

Enumerate the $\varphi_{\rho,\nu}$ in any order as $\psi_1,\psi_2,\ldots$, and let the corresponding $\lambda_\rho$ be $\lambda^{(1)},\lambda^{(2)},\ldots$. The condition formulated above, that all solutions of \hyperref[eq:II.6.E]{E} should span $\mathfrak{R}_\infty$ as a closed linear manifold, therefore means that $\psi_1,\psi_2,\ldots$, which are only a subsystem of the solutions, must already do so. By \hyperref[thm:II.7]{Theorem 7}, $\alpha)$, this means that the orthonormal system is complete.

Solving the eigenvalue problem in the sense of quantum mechanics would thus require finding enough solutions
\[
\varphi=\psi_1,\psi_2,\ldots\qquad\text{and}\qquad\lambda=\lambda_1,\lambda_2,\ldots
\]
of \hyperref[eq:II.6.E]{E} to form a complete orthonormal system from them. In general, however, this is by no means possible. In wave theory, for example, one finds that some of the solutions of \hyperref[eq:II.6.E]{E}, that is, $\hyperref[eq:I.3.E2]{\mathbf{E}_2}$ of \hyperref[sec:I.3]{I.3}, do not have a finite integral of their squared absolute value,\origNote{69} and hence do not belong to Hilbert space. Yet all these solutions are needed to expand every wave function in terms of solutions, as explained above. Thus there is no complete orthonormal system of solutions in Hilbert space, the only space admitted in \hyperref[eq:II.6.E]{E}.

On the other hand, Hilbert's theory of the eigenvalue problem shows that this phenomenon is by no means exceptional in the behavior of operators, even continuous ones.\origNote{70} We must therefore deal with the situation that arises when it occurs. Its physical meaning will be considered later; see \hyperref[sec:III.3]{III.3}. When it occurs, that is, when the orthonormal system selected from the solutions of \hyperref[eq:II.6.E]{E} is not complete, one says that a ``continuous spectrum of $\mathsf H$'' exists. The $\lambda_1,\lambda_2,\ldots$ form the ``point spectrum of $\mathsf H$.''

% Source PDF page 63 / original page 56.
\hypertarget{origpage-56}{}%
Since the formulation based on \hyperref[eq:II.6.E]{E} has proved inadequate, our next task is to find another formulation of the eigenvalue problem for Hermitian operators and then apply it to quantum mechanics. First we shall state and explain this formulation, following Hilbert; see the reference in \hyperlink{note-70}{note 70}.

\section{Continuation}\label{sec:II.7}
The equation
\[
\mathsf H\varphi=\lambda\varphi,
\]
together with the requirement that a complete orthonormal system can be formed from its solutions, was obtained by analogy with the finite-dimensional spaces $\mathfrak{R}_n$.

In $\mathfrak{R}_n$, $\mathsf H$ is a matrix $\{h_{\mu\nu}\}$, $\mu,\nu=1,\ldots,n$, $h_{\mu\nu}=\overline{h_{\nu\mu}}$; it is a familiar algebraic fact that the solutions $\varphi=\{x_1,\ldots,x_n\}$ of $\mathsf H\varphi=\lambda\varphi$, that is,
\[
\sum_{\nu=1}^n h_{\mu\nu}x_\nu=\lambda x_\mu\qquad(\mu=1,\ldots,n),
\]
contain a complete orthonormal system.\origNote{71}

As we have seen, letting $n\to\infty$ does not carry this property of $\mathfrak{R}_n$ over to $\mathfrak{R}_\infty$, so the eigenvalue problem must be formulated differently in $\mathfrak{R}_\infty$. We shall now see that it has a new formulation which is equivalent to the old one in $\mathfrak{R}_n$ and does extend to $\mathfrak{R}_\infty$. Both formulations express the same thing in every $\mathfrak{R}_n$, $n=1,2,\ldots$, namely the possibility of transforming Hermitian matrices to principal axes; but one extends to $\mathfrak{R}_\infty$, while the other does not.

Let $\{x_{11},\ldots,x_{1n}\},\ldots,\{x_{n1},\ldots,x_{nn}\}$ be the complete orthonormal system of solutions of the eigenvalue equation, and $\lambda_1,\ldots,\lambda_n$ the corresponding $\lambda$'s. Thus the vectors $\{x_{11},\ldots,x_{1n}\},\ldots,\{x_{n1},\ldots,x_{nn}\}$ form a Cartesian coordinate system in $\mathfrak{R}_n$. The transformation formulas from the coordinates $\mathfrak{x}_1,\ldots,\mathfrak{x}_n$ in this coordinate system to coordinates $\xi_1,\ldots,\xi_n$ are then
\[
\{\xi_1,\ldots,\xi_n\}=\mathfrak{x}_1\{x_{11},\ldots,x_{1n}\}+\cdots+\mathfrak{x}_n\{x_{n1},\ldots,x_{nn}\},
\]
that is,
\[
\xi_1=\sum_{\mu=1}^n x_{\mu1}\mathfrak{x}_\mu,\quad\ldots,\quad
\xi_n=\sum_{\mu=1}^n x_{\mu n}\mathfrak{x}_\mu,
\]
and inversely,
\[
\mathfrak{x}_1=\sum_{\mu=1}^n\overline{x_{1\mu}}\xi_\mu,\quad\ldots,\quad
\mathfrak{x}_n=\sum_{\mu=1}^n\overline{x_{n\mu}}\xi_\mu.
\]

% Source PDF page 64 / original page 57.
\hypertarget{origpage-57}{}%
Using the variables $\mathfrak{x}_1,\ldots,\mathfrak{x}_n$ and a new set $\eta_1,\ldots,\eta_n$, together with the associated $\mathfrak{y}_1,\ldots,\mathfrak{y}_n$ given by the formula above, we may write the conditions $\sum_{\nu=1}^n h_{\mu\nu}x_{\rho\nu}=\lambda_\rho x_{\rho\mu}$ as
\[
\sum_{\rho,\mu=1}^n\left(\sum_{\nu=1}^n h_{\mu\nu}x_{\rho\nu}\right)\mathfrak{x}_\rho\overline{\eta_\mu}
=\sum_{\rho,\mu=1}^n\lambda_\rho x_{\rho\mu}\mathfrak{x}_\rho\overline{\eta_\mu},
\]
that is,
\begin{equation}
\sum_{\mu,\nu=1}^n h_{\mu\nu}\xi_\mu\overline{\eta_\nu}
=\sum_{\rho=1}^n\lambda_\rho
\left(\sum_{\mu=1}^n\overline{x_{\rho\mu}}\xi_\mu\right)
\overline{\left(\sum_{\mu=1}^m\overline{x_{\rho\mu}}\eta_\mu\right)}.
\tag{D}\label{eq:II.7.D}
\end{equation}
The Cartesian character of our coordinate system, on the other hand, is expressed by
\begin{equation}
\sum_{\mu=1}^n\xi_\mu\overline{\eta_\mu}
=\sum_{\mu=1}^n
\left(\sum_{\rho=1}^n\overline{x_{\rho\mu}}\xi_\mu\right)
\overline{\left(\sum_{\rho=1}^n\overline{x_{\rho\mu}}\eta_\mu\right)}.
\tag{O}\label{eq:II.7.O}
\end{equation}
Finding a matrix $\{x_{\mu\nu}\}$ with properties \hyperref[eq:II.7.D]{D} and \hyperref[eq:II.7.O]{O} is thus equivalent to solving the eigenvalue problem in $\mathfrak{R}_n$; and it was in this form that the transfer to $\mathfrak{R}_\infty$ failed. This failure is not surprising, however, for the following reason. Conditions \hyperref[eq:II.7.D]{D} and \hyperref[eq:II.7.O]{O} do not completely determine the unknowns $\lambda_\rho,x_{\mu\nu}$. Although the theory of these principal-axis transformations shows that the $\lambda_\rho$ are uniquely determined up to order (see the reference in \hyperlink{note-71}{note 71}), the situation is much worse for the $x_{\mu\nu}$. Each row $x_{\rho1},\ldots,x_{\rho n}$ may plainly be multiplied by a factor $\theta_\rho$ of absolute value one. If some $\lambda_\rho$ coincide, even an arbitrary unitary transformation among the corresponding rows $x_{\rho1},\ldots,x_{\rho n}$ is admissible! To attempt the difficult passage to the limit $n\to\infty$ with quantities that are not uniquely determined is hopeless. How can the process converge when the freedom in choosing $\lambda_\rho,x_{\mu\nu}$ allows them to undergo arbitrarily large fluctuations along the way?

This gives us a clue to the proper approach. We must first try to replace conditions \hyperref[eq:II.7.D]{D}, \hyperref[eq:II.7.O]{O} and the unknowns $\lambda_\rho,x_{\mu\nu}$ by new conditions and quantities with the required uniqueness. We shall find that passage to the limit then presents fewer difficulties.

If $l$ is any value assumed by one or more $\lambda_\rho$, then
\[
\sum_{\lambda_\rho=l}
\left(\sum_{\mu=1}^n\overline{x_{\rho\mu}}\xi_\mu\right)
\overline{\left(\sum_{\mu=1}^n\overline{x_{\rho\mu}}\eta_\mu\right)}
\]
is invariant under the changes of $\lambda_\rho,x_{\mu\nu}$ listed above, which are compatible with \hyperref[eq:II.7.D]{D} and \hyperref[eq:II.7.O]{O}. If $l$ differs from all the $\lambda_\rho$, the expression is zero and is therefore certainly invariant. Consequently the Hermitian form
\[
E(l;\xi,\eta)=\sum_{\lambda_\rho\leq l}
\left(\sum_{\rho=1}^n\overline{x_{\rho\mu}}\xi_\mu\right)
\overline{\left(\sum_{\rho=1}^n\overline{x_{\rho\mu}}\eta_\mu\right)},
\]
where $\xi,\eta$ stand for $\xi_1,\ldots,\xi_n$ and $\eta_1,\ldots,\eta_n$, respectively, is likewise invariant, for arbitrary $l$.\tnote{The source prints $m$ as the last upper summation limit in \hyperref[eq:II.7.D]{D} and mixes the summation indices $\mu,\rho$ in \hyperref[eq:II.7.O]{O} and in the displayed definition of $E(l;\xi,\eta)$. These indices are retained. The preceding coordinate transformation supplies the consistent interpretation.}

% Source PDF page 65 / original page 58.
\hypertarget{origpage-58}{}%
Knowing the $E(l;\xi,\eta)$, that is, their coefficients, makes it easy to recover the $\lambda_\rho,x_{\mu\nu}$. If we therefore formulate the eigenvalue problem, \hyperref[eq:II.7.D]{D} and \hyperref[eq:II.7.O]{O}, so that only the $E(l;\xi,\eta)$ occur in place of the $\lambda_\rho,x_{\mu\nu}$, we have achieved the desired unique formulation.

Let $E(l)$ be the matrix of the Hermitian form $E(l;\xi,\eta)$.\origNote{72} What do \hyperref[eq:II.7.D]{D} and \hyperref[eq:II.7.O]{O} say about the family of matrices $E(l)$?

\hyperref[eq:II.7.O]{O} means that for sufficiently large $l$, namely larger than every $\lambda_\rho$, $E(l)=1$, the unit matrix. The nature of $E(l)$ shows that for sufficiently small $l$, namely smaller than every $\lambda_\rho$, $E(l)=0$; and as $l$ increases from $-\infty$ to $+\infty$, $E(l)$ remains constant except at finitely many points, where it changes by jumps. These points are the distinct values among $\lambda_1,\ldots,\lambda_n$; denote them by $l_1<l_2<\cdots<l_m$, $m\leq n$. At each such point the function is discontinuous only from the left, since $\sum_{\lambda_\rho\leq l}$ is continuous from the right as a function of $l$; with $\sum_{\lambda_\rho<l}$ the reverse would hold. Finally, as we shall show, for $l'\leq l''$,
\[
E(l')E(l'')=E(l'')E(l')=E(l'),
\]
where the products are matrix products.

It is more convenient to prove this for $E(l';\xi,\eta)$ and $E(l'';\xi,\eta)$ in the coordinate system of $\mathfrak{x}_1,\ldots,\mathfrak{x}_n$ and $\mathfrak{y}_1,\ldots,\mathfrak{y}_n$. Introducing these variables changes $E(l';\xi,\eta)$ and $E(l'';\xi,\eta)$ into
\[
\sum_{\lambda_\rho\leq l'}\mathfrak{x}_\rho\overline{\mathfrak{y}_\rho}
\qquad\text{and}\qquad
\sum_{\lambda_\rho\leq l''}\mathfrak{x}_\rho\overline{\mathfrak{y}_\rho},
\]
respectively. Their matrices thus have zeros everywhere off the diagonal; in the $\rho$th diagonal position they have $1$ if $\lambda_\rho\leq l'$ or $\lambda_\rho\leq l''$, respectively, and $0$ otherwise. For these matrices the assertion above is clear.

We now reformulate \hyperref[eq:II.7.D]{D} as well. It evidently says
\[
\sum_{\mu,\nu=1}^n h_{\mu\nu}\xi_\mu\overline{\eta_\nu}
=\sum_{\tau=1}^m l_\tau\{E(l_\tau;\xi,\eta)-E(l_{\tau-1};\xi,\eta)\},
\]
where $l_0$ is any number less than $l_1$. But $E(l;\xi,\eta)$ is constant on each of the intervals
\[
-\infty<l<l_1;\quad l_1\leq l<l_2;\quad\ldots;\quad l_{m-1}\leq l<l_m;\quad l_m\leq l<+\infty.
\]
Thus for every system of numbers $\Lambda_0<\Lambda_1<\Lambda_2<\cdots<\Lambda_k$, provided $l_1,\ldots,l_m$ occur among $\Lambda_1,\ldots,\Lambda_k$,
\[
\sum_{\mu,\nu=1}^n h_{\mu\nu}\xi_\mu\overline{\eta_\nu}
=\sum_{\tau=1}^k\Lambda_\tau\{E(\Lambda_\tau;\xi,\eta)-E(\Lambda_{\tau-1};\xi,\eta)\}.
\]

% Source PDF page 66 / original page 59.
\hypertarget{origpage-59}{}%
Using Stieltjes's concept of the integral,\origNote{73} this can also be written
\[
\sum_{\mu,\nu=1}^n h_{\mu\nu}\xi_\mu\overline{\eta_\nu}
=\int_{-\infty}^{+\infty}\lambda\,dE(\lambda;\xi,\eta).
\]
The integral $\int_{-\infty}^{+\infty}$ can plainly be replaced by any $\int_a^b$ with $a<l_1$, $b>l_m$. Alternatively, since this equation holds for every coefficient, we can write it as a matrix equation:
\[
\mathsf H=\int_{-\infty}^{+\infty}\lambda\,dE(\lambda),\qquad \mathsf H=\{h_{\mu\nu}\}.
\]
Our discussion has thus led to the following task. Given a Hermitian matrix $\mathsf H=\{h_{\mu\nu}\}$, find a family of Hermitian matrices $E(\lambda)$, $-\infty<\lambda<+\infty$, with these properties:

\leavevmode\phantomsection\textbf{$S_1$.}\label{cond:II.S1} For sufficiently small $\lambda$, $E(\lambda)=0$; for sufficiently large $\lambda$, $E(\lambda)=1$. As a function of $\lambda$, $E(\lambda)$ is constant everywhere except at finitely many points. At those points it changes by jumps, with discontinuity only from the left.

\leavevmode\phantomsection\textbf{$S_2$.}\label{cond:II.S2} We always have $E(\lambda')E(\lambda'')=E(\operatorname{Min}(\lambda',\lambda''))$.\origNote{74}

\leavevmode\phantomsection\textbf{$S_3$.}\label{cond:II.S3} Using the Stieltjes integral,
\[
\mathsf H=\int_{-\infty}^{+\infty}\lambda\,dE(\lambda).
\]
Although it would be easy to recover the solutions \hyperref[eq:II.7.D]{D}, \hyperref[eq:II.7.O]{O} from $\hyperref[cond:II.S1]{S_1}$--$\hyperref[cond:II.S3]{S_3}$, we shall not do so here, since this is the only form of the eigenvalue problem that we shall need in quantum mechanics. Instead, we proceed at once to generalize $\hyperref[cond:II.S1]{S_1}$--$\hyperref[cond:II.S3]{S_3}$ from finitely many to infinitely many variables, that is, from $\mathfrak{R}_n$ to $\mathfrak{R}_\infty$.

In $\mathfrak{R}_\infty$, $\mathsf H$ and the $E(\lambda)$ must evidently be understood as Hermitian operators. Given $\mathsf H$, we shall therefore try to determine a family $E(\lambda)$ standing in a certain relation to it, modeled on $\hyperref[cond:II.S1]{S_1}$--$\hyperref[cond:II.S3]{S_3}$. We must find the $\mathfrak{R}_\infty$ analogues of $\hyperref[cond:II.S1]{S_1}$--$\hyperref[cond:II.S3]{S_3}$.

We retain $\hyperref[cond:II.S2]{S_2}$ unchanged, since the number of dimensions of $\mathfrak{R}_n$ plays no role here. We shall, however, reformulate it slightly using our results on projection operators (\hyperref[sec:II.4]{II.4}). First, for $\lambda'=\lambda''=\lambda$, it states $E(\lambda)^2=E(\lambda)$, so the $E(\lambda)$ must be projection operators. Then $\hyperref[cond:II.S2]{S_2}$ means that $\lambda'\leq\lambda''$ implies $E(\lambda')\leq E(\lambda'')$; we may restrict ourselves to $\lambda'\leq\lambda''$, since the corresponding result follows for $\lambda'\geq\lambda''$. See \hyperref[thm:II.14]{Theorem 14} and the text following it in \hyperref[sec:II.4]{II.4}.

% Source PDF page 67 / original page 60.
\hypertarget{origpage-60}{}%
Some caution is required with $\hyperref[cond:II.S3]{S_3}$: $A=\int_{-\infty}^{\infty}\lambda\,dE(\lambda)$ is meaningless, since the Stieltjes integral is defined for numbers, not for operators. It is easy, however, to replace $\mathsf H,E(\lambda)$ by numbers that convey the intended operator relation. We require
\[
(\mathsf Hf,g)=\int_{-\infty}^{+\infty}\lambda\,d(E(\lambda)f,g)
\]
for every $f,g$ in $\mathfrak{R}_\infty$, insofar as $\mathsf Hf$ is meaningful; $\mathsf H=\int_{-\infty}^{+\infty}\lambda\,dE(\lambda)$ is to be understood only symbolically, as an abbreviation for this.

Finally, $\hyperref[cond:II.S1]{S_1}$ is substantially affected when the number of dimensions becomes infinite. In $\mathfrak{R}_n$, the points beyond which $E(\lambda)$ equals $0$ or $1$, or at which it makes its discontinuous jumps, are the eigenvalues of $\mathsf H$, while its intervals of constancy are the intervals free of eigenvalues. As $n\to\infty$, various things can happen. The smallest and largest eigenvalues may tend to $-\infty$ and $+\infty$, respectively, while the others, becoming ever more numerous, may become more and more densely situated, so that the intervals of constancy may shrink to points. In Hilbert's theory, this last possibility in particular may signal the appearance of the so-called continuous spectrum.\origNote{75} We must therefore alter $\hyperref[cond:II.S1]{S_1}$ quite substantially in passing from the $\mathfrak{R}_n$ to $\mathfrak{R}_\infty$, allowing for the possible disappearance of the discrete, stepwise character of the variation.

From this perspective it is quite plausible to drop the requirement that $E(\lambda)$ eventually attain the values $0,1$, and instead require merely convergence to $0$ and $1$ as $\lambda\to-\infty$ and $\lambda\to+\infty$, respectively. Likewise, instead of requiring constancy on intervals and jumps at points, we admit continuous increase as well. The less restrictive requirement that, at any point of discontinuity, the function be discontinuous only from the left may, on the other hand, be provisionally retained. We therefore formulate $\hyperref[cond:II.S1]{S_1}$ as follows: as $\lambda\to-\infty$, $E(\lambda)\to0$; as $\lambda\to+\infty$, $E(\lambda)\to1$; and as $\lambda\to\lambda_0$, $\lambda\geq\lambda_0$, $E(\lambda)\to E(\lambda_0)$.\origNote{76}

One more observation is needed concerning $\hyperref[cond:II.S3]{S_3}$. In the finite-dimensional $\mathfrak{R}_n$, we had $A=\sum_{\tau=1}^m l_\tau F_\tau$, where $F_\tau$ denotes the matrix $E(l_\tau)-E(l_{\tau-1})$. By $\hyperref[cond:II.S1]{S_1}$,
\begin{align*}
\text{for }\sigma\geq\tau:\quad
F_\tau E(l_\sigma)&=E(l_\tau)E(l_\sigma)-E(l_{\tau-1})E(l_\sigma)\\
&=E(l_\tau)-E(l_{\tau-1})=F_\tau,\\
\text{for }\sigma\leq\tau-1:\quad
F_\tau E(l_\sigma)&=E(l_\tau)E(l_\sigma)-E(l_{\tau-1})E(l_\sigma)\\
&=E(l_\sigma)-E(l_\sigma)=0.
\end{align*}

% Source PDF page 68 / original page 61.
\hypertarget{origpage-61}{}%
Hence, since $F_\sigma=E(l_\sigma)-E(l_{\sigma-1})$,
\[
F_\tau F_\sigma=\begin{cases}F_\tau,&\tau=\sigma,\\0,&\tau\ne\sigma.\end{cases}
\]
Consequently,
\[
\mathsf H^2=\left(\sum_{\tau=1}^m l_\tau F_\tau\right)^2
=\sum_{\tau,\sigma=1}^m l_\tau l_\sigma F_\tau F_\sigma
=\sum_{\tau=1}^m l_\tau^2F_\tau.
\]
Similarly, $\mathsf H^p=\sum_{\tau=1}^m l_\tau^pF_\tau$. The same transformation as for $\mathsf H$ itself therefore gives
\[
\mathsf H^2=\int_{-\infty}^{+\infty}\lambda^2\,dE(\lambda).
\]
In $\mathfrak{R}_\infty$ we therefore expect the analogous symbolic equation, which in scalar form reads
\[
(\mathsf H^2f,g)=\int_{-\infty}^{+\infty}\lambda^2\,d(E(\lambda)f,g).
\]
This will be confirmed below. For $f=g$, since
\[
(\mathsf H^2f,f)=(\mathsf Hf,\mathsf Hf)=\|\mathsf Hf\|^2,\qquad(E(\lambda)f,f)=\|E(\lambda)f\|^2,
\]
it follows that
\[
\|\mathsf Hf\|^2=\int_{-\infty}^{+\infty}\lambda^2\,d(\|E(\lambda)f\|^2).
\]
This formula suggests that the $E(\lambda)$ determine not only the value of $\mathsf Hf$ when it is meaningful, but also when it is meaningful. For the integral $\int_{-\infty}^{+\infty}\lambda^2\,d(\|E(\lambda)f\|^2)$ has an integrand which is never negative, $\lambda^2\geq0$, and a nondecreasing expression after the $d$ sign, $\|E(\lambda)f\|^2$; compare $\hyperref[cond:II.S2]{S_2}$ and \hyperref[thm:II.15]{Theorem 15} in \hyperref[sec:II.4]{II.4}. Thus by its nature the integral either converges to zero or to a positive finite value, or diverges properly, that is, to $+\infty$.\origNote{77} This holds independently of its relation to $\mathsf H$, whether $\mathsf Hf$ is meaningful or not. We may therefore expect $\mathsf Hf$ to be meaningful, that is, to exist in $\mathfrak{R}_\infty$, if and only if the proposed value of $\|\mathsf Hf\|^2$, namely the expression $\int_{-\infty}^{+\infty}\lambda^2\,d\|E(\lambda)f\|^2$, which is meaningful for every $f$, is finite.

Our reformulation of $\hyperref[cond:II.S1]{S_1}$--$\hyperref[cond:II.S3]{S_3}$ therefore reads as follows. Given a Hermitian operator $\mathsf H$, we seek a family of projection operators $E(\lambda)$, $-\infty<\lambda<+\infty$, with the following properties:

\leavevmode\phantomsection\textbf{$\overline{S}_1$.}\label{cond:II.Sbar1} As $\lambda\to-\infty$ or $\lambda\to+\infty$, $E(\lambda)f\to0$ or $f$, respectively; as $\lambda\to\lambda_0$ with $\lambda\geq\lambda_0$, $E(\lambda)f\to E(\lambda_0)f$, for every $f$.

\leavevmode\phantomsection\textbf{$\overline{S}_2$.}\label{cond:II.Sbar2} If $\lambda'\leq\lambda''$, then $E(\lambda')\leq E(\lambda'')$.

% Source PDF page 69 / original page 62.
\hypertarget{origpage-62}{}%
\leavevmode\phantomsection\textbf{$\overline{S}_3$.}\label{cond:II.Sbar3} The integral
\[
\int_{-\infty}^{+\infty}\lambda^2\,d(\|E(\lambda)f\|^2),
\]
which by its nature either converges to zero or a positive finite value, or diverges properly, determines the domain of $\mathsf H$: $\mathsf Hf$ is defined if and only if this integral is finite. Then, for every $g$,
\[
(\mathsf Hf,g)=\int_{-\infty}^{+\infty}\lambda\,d(E(\lambda)f,g).
\]
The latter integral converges absolutely whenever the former is finite.\origNote{78}

$\mathsf H$ itself does not enter properties $\hyperref[cond:II.Sbar1]{\overline{S}_1}$, $\hyperref[cond:II.Sbar2]{\overline{S}_2}$ at all. A family of projection operators $E(\lambda)$ with these two properties is called a resolution of the identity. A resolution of the identity standing in relation $\hyperref[cond:II.Sbar3]{\overline{S}_3}$ to $\mathsf H$ is said to belong to $\mathsf H$.

Thus the eigenvalue problem of $\mathfrak{R}_\infty$ is this: does every given Hermitian operator $\mathsf H$ have resolutions of the identity belonging to it, and how many? The desired answer would be that there is always exactly one. We must also examine how our definition of the eigenvalue problem relates to the methods generally used in quantum mechanics, especially in wave theory, to determine the eigenvalues of Hermitian operators.

\section{Preliminary Considerations on the Eigenvalue Problem}\label{sec:II.8}
The first question prompted by our definition of the eigenvalue problem is this. Conditions $\hyperref[cond:II.Sbar1]{\overline{S}_1}$--$\hyperref[cond:II.Sbar3]{\overline{S}_3}$ look completely different from the problem with which we began the preceding section; it is no longer apparent what they have to do with one another. In $\mathfrak{R}_n$, to be sure, we derived $\hyperref[cond:II.S1]{S_1}$--$\hyperref[cond:II.S3]{S_3}$ from those conditions. In $\mathfrak{R}_\infty$, however, the situation is essentially different: the two formulations are no longer equivalent here, although, as mentioned, they were equivalent in $\mathfrak{R}_n$. The entire question is therefore reopened. We must determine to what extent the new formulation includes the old one, that is, when and how our $E(\lambda)$ determine the earlier $\lambda_1,\lambda_2,\ldots$ and $\varphi_1,\varphi_2,\ldots$.

Let the resolution of the identity $E(\lambda)$ belong to the Hermitian operator $A$. When is the equation
\[
A\varphi=\lambda_0\varphi
\]
solvable? The equation $A\varphi=\lambda_0\varphi$ is equivalent to $(A\varphi,g)-\lambda_0(\varphi,g)=0$ for every $g$; that is,

% Source PDF page 70 / original page 63.
\hypertarget{origpage-63}{}%
\begin{align*}
0&=\int_{-\infty}^{\infty}\lambda\,d(E(\lambda)f,g)-\lambda_0(\varphi,g)\\
&=\int_{-\infty}^{\infty}\lambda\,d(E(\lambda)f,g)-\lambda_0\int_{-\infty}^{\infty}d(E(\lambda)f,g)\\
&=\int_{-\infty}^{\infty}(\lambda-\lambda_0)\,d(E(\lambda)f,g).
\end{align*}
First put $g=E(\lambda_0)f$. Then
\begin{align*}
0&=\int_{-\infty}^{\infty}(\lambda-\lambda_0)\,d(E(\lambda)f,E(\lambda_0)f)\\
&=\int_{-\infty}^{\infty}(\lambda-\lambda_0)\,d(E(\lambda_0)E(\lambda)f,f)\\
&=\int_{-\infty}^{\infty}(\lambda-\lambda_0)\,d(E(\operatorname{Min}(\lambda,\lambda_0))f,f)\\
&=\int_{-\infty}^{\infty}(\lambda-\lambda_0)\,d(\|E(\operatorname{Min}(\lambda,\lambda_0))f\|^2).
\end{align*}
We can split $\int_{-\infty}^{\infty}$ into $\int_{-\infty}^{\lambda_0}+\int_{\lambda_0}^{\infty}$. In the first integral, $\operatorname{Min}(\lambda,\lambda_0)$ can be replaced by $\lambda$, and in the second by $\lambda_0$. The expression after the $d$ sign in the latter integral is therefore constant, so that integral vanishes. The first integral therefore gives\tnote{The source prints the limits $\lambda_0,\infty$ in the following integral, although the preceding split requires $-\infty,\lambda_0$. The printed limits are retained.}
\[
\int_{\lambda_0}^{\infty}(\lambda-\lambda_0)\,d(\|E(\lambda)f\|^2)=0.
\]
Second, put $g=f$. Then
\[
0=\int_{-\infty}^{\infty}(\lambda-\lambda_0)\,d(E(\lambda)f,f)
=\int_{-\infty}^{\infty}(\lambda-\lambda_0)\,d(\|E(\lambda)f\|^2).
\]
Subtracting the first equation from this, and also reversing the sign of the integrand, we obtain
\[
\int_{-\infty}^{\lambda_0}(\lambda_0-\lambda)\,d(\|E(\lambda)f\|^2)=0.
\]
Now examine the integrals $\int_{-\infty}^{\lambda_0}(\lambda_0-\lambda)\,d(\|E(\lambda)f\|^2)$ and $\int_{\lambda_0}^{\infty}(\lambda-\lambda_0)\,d(\|E(\lambda)f\|^2)$ more closely. In both, the integrand is nonnegative and a monotonically nondecreasing function of $\lambda$ stands after the $d$ sign. Thus for every $\varepsilon>0$,
\begin{align*}
\int_{-\infty}^{\lambda_0}(\lambda_0-\lambda)\,d(\|E(\lambda)f\|^2)
&\geq\int_{-\infty}^{\lambda_0-\varepsilon}(\lambda_0-\lambda)\,d(\|E(\lambda)f\|^2)\\
&\geq\int_{-\infty}^{\lambda_0-\varepsilon}\varepsilon\,d(\|E(\lambda)f\|^2)\\
&=\varepsilon\|E(\lambda_0-\varepsilon)f\|^2,
\end{align*}
and
\begin{align*}
\int_{\lambda_0}^{\infty}(\lambda-\lambda_0)\,d(\|E(\lambda)f\|^2)
&\geq\int_{\lambda_0+\varepsilon}^{\infty}(\lambda-\lambda_0)\,d(\|E(\lambda)f\|^2)\\
&\geq\int_{\lambda_0+\varepsilon}^{\infty}\varepsilon\,d(\|E(\lambda)f\|^2)\\
&=\varepsilon(\|f\|^2-\|E(\lambda_0+\varepsilon)f\|^2)\\
&=\varepsilon\|f-E(\lambda_0+\varepsilon)f\|^2.
\end{align*}

% Source PDF page 71 / original page 64.
\hypertarget{origpage-64}{}%
The right sides are therefore nonpositive; since they are identically nonnegative, they must vanish. Hence
\[
E(\lambda_0-\varepsilon)f=0,\qquad E(\lambda_0+\varepsilon)f=f.
\]
By the right continuity of $E(\lambda)$, we may let $\varepsilon\to0$ in the second equation, obtaining $E(\lambda_0)f=f$. For $\lambda\geq\lambda_0$, the second equation, with $\varepsilon=\lambda-\lambda_0\geq0$, then gives $E(\lambda)f=f$; for $\lambda<\lambda_0$, the first, with $\varepsilon=\lambda_0-\lambda>0$, gives $E(\lambda)f=0$. Thus
\[
E(\lambda)f=\begin{cases}f,&\lambda\geq\lambda_0,\\0,&\lambda<\lambda_0.\end{cases}
\]
This necessary condition is also sufficient, for it implies
\[
(Af,g)=\int_{-\infty}^{\infty}\lambda\,d(E(\lambda)f,g)=\lambda_0(f,g),
\]
as follows by recalling the definition of the Stieltjes integral in \hyperlink{note-73}{note 73}. Hence $(Af-\lambda_0f,g)=0$ for every $g$, that is, $Af=\lambda_0f$.

How should this condition be interpreted? First, it implies that $E(\lambda)$ is discontinuous at $\lambda=\lambda_0$. By \hyperref[thm:II.17]{Theorem 17} in \hyperref[sec:II.4]{II.4}, $E(\lambda)$ must tend to a projection operator both as $\lambda\to\lambda_0$, $\lambda<\lambda_0$, and as $\lambda\to\lambda_0$, $\lambda>\lambda_0$; denote these limits by $E^{(1)}(\lambda_0)$ and $E^{(2)}(\lambda_0)$, respectively.\origNote{79} By $\hyperref[cond:II.S1]{S_1}$, $E^{(2)}(\lambda_0)=E(\lambda_0)$, but in the case of discontinuity, $E^{(1)}(\lambda_0)\ne E(\lambda_0)$. Furthermore, since $E(\lambda)\leq E(\lambda_0)$ for $\lambda<\lambda_0$ by $\hyperref[cond:II.S2]{S_2}$, we always have $E^{(1)}(\lambda_0)\leq E(\lambda_0)$. Thus $E(\lambda_0)-E^{(1)}(\lambda_0)$ is a projection operator, and discontinuity is characterized by its being nonzero.

The condition $E(\lambda)f=0$ for every $\lambda<\lambda_0$ implies $E^{(1)}(\lambda_0)f=0$, and is also a consequence of it, since $E(\lambda)\leq E^{(1)}(\lambda_0)$. The condition $E(\lambda)f=f$ for every $\lambda\geq\lambda_0$ follows from $E(\lambda_0)f=f$: we have $E(\lambda_0)\leq E(\lambda)$, $E(\lambda)E(\lambda_0)=E(\lambda_0)$, and hence $E(\lambda)f=E(\lambda)E(\lambda_0)f=E(\lambda_0)f=f$. Thus $Af=\lambda_0f$ is characterized by $E^{(1)}(\lambda_0)f=0$, $E(\lambda_0)f=f$, or, by \hyperref[thm:II.14]{Theorem 14} in \hyperref[sec:II.4]{II.4}, by $(E(\lambda_0)-E^{(1)}(\lambda_0))f=f$. If we put $E(\lambda_0)-E^{(1)}(\lambda_0)=P_{\mathfrak{M}_{\lambda_0}}$, this says that $f$ belongs to $\mathfrak{M}_{\lambda_0}$.

We have therefore shown that $Af=\lambda f$ has a nonzero solution $f$ only at points of discontinuity of $E(\lambda)$; the solutions form the closed linear manifold $\mathfrak{M}_{\lambda_0}$ just defined.

The complete orthonormal system of solutions, possibly belonging to different $\lambda$'s, sought in \hyperref[sec:II.6]{II.6} therefore exists if and only if the $\mathfrak{M}_{\lambda_0}$, $-\infty<\lambda_0<\infty$, together span the closed linear manifold $\mathfrak{R}_\infty$. We discussed in \hyperref[sec:II.6]{II.6} how the system is then to be constructed. The mutual orthogonality of the $\mathfrak{M}_{\lambda_0}$ can be seen afresh: $\lambda_0<\mu_0$ implies
\[
P_{\mathfrak{M}_{\lambda_0}}P_{\mathfrak{M}_{\mu_0}}
=(E(\lambda_0)-E^{(1)}(\lambda_0))(E(\mu_0)-E^{(1)}(\mu_0))=0,
\]
because
\[
E(\lambda_0)-E^{(1)}(\lambda_0)\leq E(\lambda_0)\leq E^{(1)}(\mu_0),\qquad
E(\mu_0)-E^{(1)}(\mu_0)\leq1-E^{(1)}(\mu_0).
\]

% Source PDF page 72 / original page 65.
\hypertarget{origpage-65}{}%
Without determining precisely when such a complete system exists, we can at least show that it cannot exist if $E(\lambda)$ increases continuously on some interval $\mu_1,\mu_2$: that is, if $\mu_1<\mu_2$, $E(\lambda)$ is continuous on $\mu_1\leq\lambda\leq\mu_2$, and $E(\mu_1)\ne E(\mu_2)$. For when $\lambda\leq\mu_1$, $E(\lambda)-E^{(1)}(\lambda)\leq E(\lambda)\leq E(\mu_1)$; when $\mu_1<\lambda\leq\mu_2$, continuity gives $E(\lambda)-E^{(1)}(\lambda)=0$; and when $\mu_2<\lambda$, $E(\lambda)-E^{(1)}(\lambda)\leq1-E^{(1)}(\lambda)\leq1-E(\mu_2)$. Thus $E(\lambda)-E^{(1)}(\lambda)$ is always orthogonal to $E(\mu_2)-E(\mu_1)$. Put $E(\mu_2)-E(\mu_1)=P_{\mathfrak{N}}$; then all $\mathfrak{M}_\lambda$ are orthogonal to $\mathfrak{N}$. If a complete orthonormal system could be selected from them, $\mathfrak{N}$ would contain only $0$, that is, $E(\mu_2)-E(\mu_1)=0$, contrary to the assumption.

We call the points of discontinuity of $E(\lambda)$ the point spectrum of $A$. They are the $\lambda$ for which $Af=\lambda f$, $f\ne0$, is solvable. Choosing an $f$ with $\|f\|=1$ from each nonzero $\mathfrak{M}_\lambda$ produces an orthonormal system, by the orthogonality of the $\mathfrak{M}_\lambda$. By \hyperref[thm:II.3infty]{Theorem 3} in \hyperref[sec:II.2]{II.2}, this system is finite or a sequence, so the $\lambda$'s of the point spectrum form at most a sequence as well.

The spectrum of $A$ consists of all points in every neighborhood of which $E(\lambda)$ is nonconstant. We have seen that if there are intervals containing points of the spectrum of $A$ but none of its point spectrum, that is, intervals of continuity of $E(\lambda)$ on which it is not constant, the eigenvalue problem is certainly not solvable in the sense formulated at the beginning of \hyperref[sec:II.6]{II.6}. We shall not investigate the precise conditions of this insolubility further. The problem may also be insoluble when the point spectrum meets every interval containing points of the spectrum. Separating the point spectrum from the rest is then considerably more laborious and falls outside the scope of our subject. The reader will find these investigations in Hilbert's cited papers.

We shall, however, show how to construct $E(\lambda)$ when there is a complete orthonormal system $\varphi_1,\varphi_2,\ldots$ of solutions of $A\varphi=\lambda\varphi$, with $\lambda=\lambda_1,\lambda_2,\ldots$ for $\varphi=\varphi_1,\varphi_2,\ldots$, respectively---the case of a pure point spectrum, as we shall call it. We have\origNote{80}
\[
E(\lambda)=\sum_{\lambda_\rho\leq\lambda}P_{[\varphi_\rho]}.
\]
The sum may have no summands, in which case $E(\lambda)=0$; a positive finite number, when its meaning is clear; or infinitely many, when it converges by the concluding discussion of \hyperref[sec:II.4]{II.4}.

Indeed, $\hyperref[cond:II.Sbar2]{\overline{S}_2}$ is evident, for if $\lambda'\leq\lambda''$, then
\[
E(\lambda'')-E(\lambda')=\sum_{\lambda'<\lambda_\rho\leq\lambda''}P_{[\varphi_\rho]}
\]
is a projection operator, so $E(\lambda')\leq E(\lambda'')$ by \hyperref[thm:II.14]{Theorem 14}.

% Source PDF page 73 / original page 66.
\hypertarget{origpage-66}{}%
To prove $\hyperref[cond:II.Sbar1]{\overline{S}_1}$, note that for every $f$,\origNote{81}
\[
\sum_\rho\|P_{[\varphi_\rho]}f\|^2=\sum_\rho|(f,\varphi_\rho)|^2=\|f\|^2\,
\]
by \hyperref[thm:II.7]{Theorem 7}; thus $\sum_\rho\|P_{[\varphi_\rho]}f\|^2$ converges. For every $\varepsilon>0$, we can therefore specify finitely many $\rho$'s such that the sum over those alone already exceeds $\|f\|^2-\varepsilon$, and consequently every sum $\sum_\rho'$ from which they are absent is less than $\varepsilon$. Then also
\[
\left\|\sum_p'P_{[\varphi_\rho]}f\right\|^2=\sum_p'\|P_{[\varphi_\rho]}f\|^2<\varepsilon.
\]
In particular,
\[
\left\|\sum_{\lambda_\rho\leq\lambda}P_{[\varphi_\rho]}f\right\|^2<\varepsilon,\quad
\left\|\sum_{\lambda_\rho>\lambda}P_{[\varphi_\rho]}f\right\|^2<\varepsilon,\quad
\left\|\sum_{\lambda_0<\lambda_\rho\leq\lambda}P_{[\varphi_\rho]}f\right\|^2<\varepsilon,
\]
when $\lambda$ is chosen sufficiently small, sufficiently large, or sufficiently close to $\lambda_0$ and $\geq\lambda$, respectively. Thus indeed\origNote{82}
\begin{alignat*}{2}
E(\lambda)f&=\sum_{\lambda_\rho\leq\lambda}P_{[\varphi_\rho]}f\longrightarrow0&\quad &\text{as }\lambda\to-\infty,\\
f-E(\lambda)f&=\sum_{\lambda_\rho>\lambda}P_{[\varphi_\rho]}f\longrightarrow0&\quad &\text{as }\lambda\to+\infty,\\
E(\lambda)f-E(\lambda_0)f&=\sum_{\lambda_0<\lambda_\rho\leq\lambda}P_{[\varphi_\rho]}f\longrightarrow0&\quad &\text{as }\lambda\to\lambda_0,\ \lambda\geq\lambda_0.
\end{alignat*}
Thus $\hyperref[cond:II.Sbar1]{\overline{S}_1}$ holds.

To establish $\hyperref[cond:II.Sbar3]{\overline{S}_3}$, put $f=x_1\varphi_1+x_2\varphi_2+\cdots$, so $Af=\lambda_1x_1\varphi_1+\lambda_2x_2\varphi_2+\cdots$. For $Af$ to be meaningful, $\sum_{\rho=1}^\infty\lambda_\rho^2|x_\rho|^2$ must therefore be finite. But\origNote{83}
\[
\begin{aligned}
\int_{-\infty}^{\infty}\lambda^2\,d\|E(\lambda)f\|^2
&=\int_{-\infty}^{\infty}\lambda^2\,d\left(\sum_{\lambda_\rho\leq\lambda}|x_\rho|^2\right)
=\sum_{\rho=1}^\infty\lambda_\rho^2|x_\rho|^2,
\end{aligned}
\]
and\origNote{83}
\[
\begin{aligned}
\int_{-\infty}^{\infty}\lambda^2\,d(E(\lambda)f,g)
&=\int_{-\infty}^{\infty}\lambda\,d\left(\sum_{\lambda_\rho\leq\lambda}x_\rho\overline{v_\rho}\right)
=\sum_{\rho=1}^\infty\lambda_\rho x_\rho\overline{y_\rho}=(Af,g).
\end{aligned}
\]
Thus $\hyperref[cond:II.Sbar3]{\overline{S}_3}$ is satisfied as well.\tnote{The source prints $\lambda^2$ on the left of the second integral identity, where the right side requires $\lambda$, and changes $v_\rho$ to $y_\rho$ for the same coefficients. It also prints $\lambda\geq\lambda$ above where $\lambda\geq\lambda_0$ is intended. The printed forms are retained.}

Let us also consider two cases of a pure continuous spectrum, that is, cases in which there is no point spectrum at all. Let $\mathfrak{R}_\infty$ be the space of all functions $f(q_1\cdots q_l)$ with finite
\[
\int_{-\infty}^{\infty}\cdots\int_{-\infty}^{\infty}|f(q_1\cdots q_l)|^2\,dq_1\cdots dq_l,
\]
and let $A$ be the operator $q_j\cdots$, whose Hermitian character is evident.

The equation $Af=\lambda f$ means $(q_j-\lambda)f(q_1,\ldots,q_l)=0$, that is, $f(q_1\cdots q_l)=0$ everywhere, with the possible exception of the $(l-1)$-dimensional plane $q_j=0$. But, by what was said about condition \hyperref[ax:II.B]{B} in \hyperref[sec:II.3]{II.3}, this plane is inessential because its Lebesgue measure, that is, volume, is zero; hence $f\equiv0$.\origNote{84}\tnote{The source prints $q_j=0$ for the exceptional plane; the preceding equation gives $q_j=\lambda$, as the following paragraph also states. The printed zero is retained.}

% Source PDF page 74 / original page 67.
\hypertarget{origpage-67}{}%
Thus there is never a nonzero solution of $Af=\lambda f$. We can nevertheless see, by a nonrigorous argument, where to look for the solution. The equation $(q_j-\lambda)f(q_1,\ldots,q_l)=0$ says that $f$ may be nonzero only for $q_j=\lambda$. A linear combination of solutions for several values, say $\lambda=\lambda',\lambda'',\ldots,\lambda^{(s)}$, would therefore be an $f$ nonzero only for
\[
q_j=\lambda',\lambda'',\ldots,\lambda^{(s)}.
\]
As a linear combination of all solutions with $\lambda\leq\lambda_0$, we may therefore regard an $f$ which is nonzero only for $\lambda\leq\lambda_0$.\tnote{The source writes $\lambda\leq\lambda_0$ here. Since the passage concerns the support of $f$, this should read $q_j\leq\lambda_0$. The original wording is retained.} In the case of a pure point spectrum, however, we had
\[
E(\lambda_0)=\sum_{\lambda_\rho\leq\lambda_0}P_{[\varphi_\rho]}=P_{\mathfrak{N}_{\lambda_0}},
\qquad\mathfrak{N}_{\lambda_0}=[\varphi_{\rho(\lambda_\rho\leq\lambda_0)}];
\]
that is, $\mathfrak{N}_\lambda$ consisted of the linear combinations of all $\varphi_\rho$ with $\lambda_\rho\leq\lambda_0$, hence of all solutions of $Af=\lambda f$ with $\lambda\leq\lambda_0$. We may therefore expect, nonrigorously and heuristically of course, that now $E(\lambda_0)=P_{\mathfrak{N}_{\lambda_0}}$, where $\mathfrak{N}_{\lambda_0}$ consists of the $f$ which are nonzero only for $q_j\leq\lambda_0$. Then $\mathfrak{R}_\infty-\mathfrak{N}_{\lambda_0}$ plainly consists of the $f$ which always vanish for $q_j\leq\lambda_0$. Consequently,
\[
E(\lambda_0)f(q_1\cdots q_l)=
\begin{cases}f(q_1\cdots q_l),&q_j\leq\lambda_0,\\0,&q_j>\lambda_0.\end{cases}
\]
In this wholly nonrigorous way we have found a family of projection operators $E(\lambda)$ which may be expected to satisfy $\hyperref[cond:II.Sbar1]{\overline{S}_1}$--$\hyperref[cond:II.Sbar3]{\overline{S}_3}$ for our $A$. Indeed, $\hyperref[cond:II.Sbar1]{\overline{S}_1}$ and $\hyperref[cond:II.Sbar2]{\overline{S}_2}$ hold trivially; in $\hyperref[cond:II.Sbar1]{\overline{S}_1}$, the case $\lambda\to\lambda_0$ holds even without the restriction $\lambda\geq\lambda_0$, so $E(\lambda)$ is continuous everywhere in $\lambda$. To see that $\hyperref[cond:II.Sbar3]{\overline{S}_3}$ also holds, it suffices to establish the equations
\begin{align*}
&\int_{-\infty}^{\infty}\lambda^2\,d\|E(\lambda)f\|^2\\
&=\int_{-\infty}^{\infty}\lambda^2\,d\left(\int_{-\infty}^{\infty}\cdots\int_{-\infty}^{\lambda}\cdots\int_{-\infty}^{\infty}|f(q_1\cdots q_j\cdots q_l)|^2\,dq_1\cdots dq_j\cdots dq_l\right)\\
&=\int_{-\infty}^{\infty}\lambda^2\left(\int_{-\infty}^{\infty}\cdots\int_{-\infty}^{\infty}|f(q_1\cdots q_{j-1}\lambda q_{j+1}\cdots q_l)|^2\,dq_1\cdots dq_{j-1}\,dq_{j+1}\cdots dq_l\right)d\lambda\\
&=\int_{-\infty}^{\infty}\cdots\int_{-\infty}^{\infty}q_j^2|f(q_1\cdots q_{j-1}q_jq_{j+1}\cdots q_l)|^2\,dq_1\cdots dq_{j-1}\,dq_j\,dq_{j+1}\cdots dq_l
=\|Af\|^2,
\end{align*}
and
\begin{align*}
&\int_{-\infty}^{\infty}\lambda\,d(E(\lambda)f,g)\\
&=\int_{-\infty}^{\infty}\lambda\,d\left(\begin{aligned}
&\int_{-\infty}^{\infty}\cdots\int_{-\infty}^{\lambda}\cdots\int_{-\infty}^{\infty}f(q_1\cdots q_j\cdots q_l)\\
&\quad\times\overline{g(q_1\cdots q_j\cdots q_l)}\,dq_1\cdots dq_j\cdots dq_l
\end{aligned}\right)\\
&=\int_{-\infty}^{\infty}\lambda\left(\begin{aligned}
&\int_{-\infty}^{\infty}\cdots\int_{-\infty}^{\infty}f(q_1\cdots q_{j-1}\lambda q_{j+1}\cdots q_l)\\
&\quad\times\overline{g(q_1\cdots q_{j-1}\lambda q_{j+1}\cdots q_l)}\\
&\quad\times dq_1\cdots dq_{j-1}\,dq_{j+1}\cdots dq_l
\end{aligned}\right)d\lambda\\
&=\int_{-\infty}^{\infty}\cdots\int_{-\infty}^{\infty}q_jf(q_1\cdots q_{j-1}q_jq_{j+1}\cdots q_l)\overline{g(q_1\cdots q_{j-1}q_jq_{j+1}\cdots q_l)}\\
&\hspace{4em}\times dq_1\cdots dq_{j-1}\,dq_j\,dq_{j+1}\cdots dq_l=(Af,g).
\end{align*}

% Source PDF page 75 / original page 68.
\hypertarget{origpage-68}{}%
Once again, the approach based on the point spectrum, or equivalently on the old definition of the eigenvalue problem, must fail, since $E(\lambda)$ increases continuously everywhere.

This example also indicates how to find the $E(\lambda)$ in the continuous spectrum in the general case. Determine (without mathematical rigor!) the solutions of $Af=\lambda f$; since $\lambda$ lies in the continuous spectrum, these $f$ do not belong to $\mathfrak{R}_\infty$ at all. Form their linear combinations for all $\lambda\leq\lambda_0$. Some of these again belong to $\mathfrak{R}_\infty$ and may form a closed linear manifold $\mathfrak{N}_{\lambda_0}$. Then put $E(\lambda_0)=P_{\mathfrak{N}_{\lambda_0}}$. If this has been done properly, one can afterward verify $\hyperref[cond:II.Sbar1]{\overline{S}_1}$--$\hyperref[cond:II.Sbar3]{\overline{S}_3}$ for $A$ and these $E(\lambda)$, thus giving a rigorous justification of the heuristic argument after the fact.\origNote{85}

The second example we wish to consider is the other important operator of wave mechanics, $\frac{h}{2\pi i}\frac{\partial}{\partial q_j}$. To avoid inessential complications, put $l=j=1$; the other cases proceed similarly. We must therefore examine the operator
\[
A'f(q)=\frac{h}{2\pi i}\frac{\partial}{\partial q}f(q).
\]
When $q$ ranges over $-\infty<q<+\infty$, this is a Hermitian operator, as we saw in \hyperref[sec:II.5]{II.5}. For a finite range $a\leq q\leq b$, however, this is not so:
\begin{align*}
(A'f,g)-(f,A'g)
&=\int_a^b\frac{h}{2\pi i}f'(q)\overline{g(q)}\,dq-\int_a^b f(q)\overline{\frac{h}{2\pi i}g'(q)}\,dq\\
&=\frac{h}{2\pi i}\int_a^b\{f'(q)\overline{g(q)}+f(q)\overline{g(q)}'\}\,dq\\
&=\frac{h}{2\pi i}[f(q)\overline{g(q)}]_a^b
=\frac{h}{2\pi i}[f(b)\overline{g(b)}-f(a)\overline{g(a)}].
\end{align*}
For this to vanish, the domain of $\frac{h}{2\pi i}\frac{\partial}{\partial q}$ must be restricted so that two arbitrarily chosen $f,g$ from it satisfy $f(a)\overline{g(a)}=f(b)\overline{g(b)}$, that is, $f(a):f(b)=\overline{g(b)}:\overline{g(a)}$. Varying $f$ with $g$ fixed, we see that $f(a):f(b)$ must have the same value $\theta$ throughout the domain, where $\theta$ might also be $0$ or $\infty$. Substituting for $f,g$ then gives $\theta=1/\overline\theta$, that is, $|\theta|=1$. Thus, to make $\frac{h}{2\pi i}\frac{\partial}{\partial q}$ Hermitian, we must impose an additional boundary condition of the form
\[
f(a):f(b)=\theta,
\]
with $\theta$ any fixed number of absolute value one.

First take the interval $-\infty<q<+\infty$.

% Source PDF page 76 / original page 69.
\hypertarget{origpage-69}{}%
The solutions of $A\varphi=\lambda\varphi$, that is, $\frac{h}{2\pi i}\varphi'(q)=\lambda\varphi(q)$, are the functions
\[
\varphi(q)=c e^{\frac{2\pi i}{h}\lambda q}.
\]
These cannot, however, be used directly for our purpose, since $\int_{-\infty}^{\infty}|\varphi(q)|^2\,dq=\int_{-\infty}^{\infty}|c|^2\,dq=+\infty$, unless $c=0$, $\varphi\equiv0$. Notice that in the first example we found the solution $\delta(q-\lambda)$, a fictitious, nonexistent function (see \hyperlink{note-84}{note 84}); here we find $e^{\frac{2\pi i}{h}\lambda q}$, a perfectly ordinary function, which nevertheless does not belong to $\mathfrak{R}_\infty$ because the integral of its squared absolute value is infinite. From our standpoint the two situations amount to the same thing: whatever does not belong to $\mathfrak{R}_\infty$ does not exist for us.\origNote{86}

As in the first case, we now form linear combinations of the solutions belonging to $\lambda\leq\lambda_0$, that is, functions
\[
f(q)=\int_{-\infty}^{\lambda_0}c(\lambda)e^{\frac{2\pi i}{h}\lambda q}\,d\lambda.
\]
We hope that among these there will be functions belonging to $\mathfrak{R}_\infty$, that these will form a closed linear manifold $\mathfrak{N}'_{\lambda_0}$, and finally that the projection operators $E'(\lambda_0)=P_{\mathfrak{N}'_{\lambda_0}}$ will form the resolution of the identity belonging to $A'$. An example supporting the first conjecture is obtained by putting
\[
c(\lambda)=\begin{cases}1,&\lambda\geq\lambda_1,\\0,&\lambda<\lambda_1,\end{cases}
\qquad\lambda_1<\lambda_0.
\]
Then
\[
f(q)=\int_{\lambda_1}^{\lambda_0}e^{\frac{2\pi i}{h}\lambda q}\,d\lambda
=\frac{e^{\frac{2\pi i}{h}\lambda_0q}-e^{\frac{2\pi i}{h}\lambda_1q}}{\frac{2\pi i}{h}q}.
\]
This $f(q)$ is regular at every finite point and decreases like $1/q$ as $q\to\pm\infty$, so $\int_{-\infty}^{\infty}|f(q)|^2\,dq$ is finite. The other conjectures are also correct, as follows from the theory of Fourier integrals. That theory states the following.\origNote{87}

Given any function $f(x)$ for which $\int_{-\infty}^{\infty}|f(x)|^2\,dx$ is finite, we can define the function
\[
Lf(x)=F(y)=\frac1{\sqrt{2\pi}}\int_{-\infty}^{\infty}e^{ixy}f(x)\,dx.
\]
The integral $\int_{-\infty}^{\infty}|F(y)|^2\,dy$ is again finite and in fact equals $\int_{-\infty}^{\infty}|f(x)|^2\,dx$. Moreover, $LLf(x)=f(-x)$. This is the so-called Laplace transformation, which also plays a major role elsewhere in the theory of differential equations.

% Source PDF page 77 / original page 70.
\hypertarget{origpage-70}{}%
Replacing $x,y$ by $\sqrt{2\pi/h}\,q$, $\sqrt{2\pi/h}\,p$, we obtain the transformation
\[
Mf(q)=F(p)=\frac1{\sqrt h}\int_{-\infty}^{\infty}e^{\frac{2\pi i}{h}pq}f(q)\,dq,
\]
which has the same properties. It therefore maps $\mathfrak{R}_\infty$ onto itself: $Mf(q)=g(p)$ is solvable for every $g(p)$ in $\mathfrak{R}_\infty$, with $f(q)=Mg(-p)$. It leaves $\|f\|$ invariant and is linear; by \hyperref[sec:II.5]{II.5}, this $M$ is unitary. Also $M^2f(q)=f(-q)$, so $M^{-1}f(q)=M^*f(q)=Mf(-q)$, and $M$ commutes with $M^2$, that is, with $f(q)\mapsto f(-q)$.

Our proposed definition of $\mathfrak{N}'_{\lambda_0}$ was therefore that $f(q)$ belongs to it if $F(p)=M^{-1}f(q)$ vanishes for every $p>\lambda_0$. Here
\[
F(p)=\sqrt h\,c(p),
\]
with the $c(\lambda)$ used above. These $F(p)$ form, as we know, the closed linear manifold $\mathfrak{N}_{\lambda_0}$; hence its image under $M$, $\mathfrak{N}'_{\lambda_0}$, is also a closed linear manifold. The operator $E'(\lambda_0)$ is obtained from $E(\lambda_0)$ just as $\mathfrak{N}'_{\lambda_0}$ is obtained from $\mathfrak{N}_{\lambda_0}$, by transforming all of $\mathfrak{R}_\infty$ with $M$. Therefore $E'(\lambda_0)=ME(\lambda_0)M^{-1}$. Consequently $E'(\lambda)$, like $E(\lambda)$, has properties $\hyperref[cond:II.Sbar1]{\overline{S}_1}$, $\hyperref[cond:II.Sbar2]{\overline{S}_2}$. It remains to prove $\hyperref[cond:II.Sbar3]{\overline{S}_3}$, that is, that the resolution of the identity $E(\lambda)$ belongs to $A'$.

We confine ourselves to showing that if $f(q)$ can be differentiated without particular convergence difficulties and $\int_{-\infty}^{\infty}|\frac{h}{2\pi i}f'(q)|^2\,dq$ is finite, then $\int_{-\infty}^{\infty}\lambda^2\,d\|E'(\lambda)f\|^2$ is finite and $(A'f,g)=\int_{-\infty}^{\infty}\lambda\,d(E'(\lambda)f,g)$.\origNote{88} Indeed, with $M^{-1}f(q)=F(p)$,
\begin{align*}
A'f(q)&=\frac{h}{2\pi i}f'(q)
=\frac{h}{2\pi i}\frac{\partial}{\partial q}(MF(p))\\
&=\frac{h}{2\pi i}\frac{\partial}{\partial q}\left(\frac1{\sqrt h}\int F(p)e^{\frac{2\pi i}{h}pq}\,dp\right)\\
&=\frac{\sqrt h}{2\pi i}\int F(p)\frac{\partial}{\partial q}\left(e^{\frac{2\pi i}{h}pq}\right)dp\\
&=\frac1{\sqrt h}\int F(p)\,p\,e^{\frac{2\pi i}{h}pq}\,dp=M(pF(p)).
\end{align*}
Thus for the specified $f$, $A'=MAM^{-1}$, where $A$ is the operator $q\cdots$, or $p\cdots$ since we use $p$ as the variable here. Since the assertions above hold for $A,E(\lambda)$, they remain valid after transforming $\mathfrak{R}_\infty$ by $M$; hence they hold for $A'=MAM^{-1}$ and $E'(\lambda)=ME(\lambda)M^{-1}$.

% Source PDF page 78 / original page 71.
\hypertarget{origpage-71}{}%
The situation is essentially different for $\frac{h}{2\pi i}\frac{\partial}{\partial q}$ on the interval $a\leq q\leq b$, with finite $a<b$. As we know, a boundary condition
\[
f(a):f(b)=\theta,\qquad|\theta|=1,
\]
is also necessary to preserve its Hermitian character. Again $\frac{h}{2\pi i}\frac{\partial}{\partial q}f(q)=\lambda f(q)$ is solved by
\[
f(q)=c e^{\frac{2\pi i}{h}\lambda q},
\]
but now $\int_a^b|f(q)|^2\,dq=\int_a^b|c|^2\,dq=(b-a)|c|^2$ is finite, so $f(q)$ always belongs to $\mathfrak{R}_\infty$. The boundary condition must still be satisfied:
\[
f(a):f(b)=e^{\frac{2\pi i}{h}\lambda(a-b)}=\theta.
\]
Putting $\theta=e^{-i\alpha}$, $0\leq\alpha<2\pi$, this becomes
\[
\frac{2\pi i}{h}\lambda(a-b)=-i\alpha-2k\pi i\qquad(k=0,\pm1,\pm2,\ldots),
\qquad\lambda=\frac h{b-a}\left(\frac\alpha{2\pi}+k\right).
\]
Thus there is a point spectrum. The normalized solutions are characterized by $(b-a)|c|^2=1$, for example $c=1/\sqrt{b-a}$:
\[
\varphi_k(q)=\frac1{\sqrt{b-a}}e^{\frac{2\pi i}{h}\lambda q}
=\frac1{\sqrt{b-a}}e^{\frac{2\pi i}{b-a}(\frac\alpha{2\pi}+k)q},
\qquad k=0,\pm1,\pm2,\ldots.
\]
This is an orthonormal system, and it is also complete. For if $f(q)$ is orthogonal to every $\varphi_k(q)$, then $e^{\frac{\alpha i}{b-a}q}f(q)$ is orthogonal to every $e^{\frac{2\pi i}{b-a}kq}$, and therefore $e^{\frac{\alpha i}{2\pi}x}f(\frac{b-a}{2\pi}x)$ is orthogonal to every $e^{ikx}$, that is, to $1,\cos x,\sin x,\cos2x,\sin2x,\ldots$. It is defined on the interval
\[
a\leq\frac{b-a}{2\pi}x\leq b,
\]
whose length is $2\pi$, so familiar theorems show that it must vanish.\origNote{89} Hence $f(q)\equiv0$ as well.

Thus there is a pure point spectrum, a case we dealt with in general at the beginning of this section. Notice how the boundary condition, that is, $\theta$ or $\alpha$, affects the eigenvalues and eigenfunctions.

Finally, we may also consider a half-infinite interval, say $0\leq q<+\infty$.

% Source PDF page 79 / original page 72.
\hypertarget{origpage-72}{}%
Again we must first check the Hermitian character. We have
\begin{align*}
(A'f,g)-(f,A'g)&=\frac{h}{2\pi i}\int_0^{+\infty}\bigl(f'(q)\overline{g(q)}+f(q)\overline{g'(q)}\bigr)\,dq\\
&=\frac{h}{2\pi i}[f(q)\overline{g(q)}]_0^{+\infty}.
\end{align*}
That $f(q)\overline{g(q)}\to0$ as $q\to+\infty$ is shown just as it was in \hyperref[sec:II.5]{II.5} for the interval $-\infty,+\infty$. Thus $f(0)\overline{g(0)}=0$ is required. Putting $f=g$ shows that the boundary condition must be $f(0)=0$.

Serious difficulties arise in this case. The solutions of $A'\varphi=\lambda\varphi$ are the same as on the interval $-\infty<q<+\infty$, namely $ce^{\frac{2\pi i}{h}\lambda q}$. They do not belong to $\mathfrak{R}_\infty$ and violate the boundary condition. The failure to satisfy the boundary condition is already cause for concern. Stranger still, the procedure outlined earlier would lead to the same $E(\lambda)$ as for $-\infty<q<+\infty$, since the improper solutions, that is, those not belonging to $\mathfrak{R}_\infty$, are the same. How can this be reconciled with the fact that the operator is different? Moreover, the resulting operators are not the ones we need. For if in the Hilbert space $F_{\Omega'}$ of the $f(q)$, $0\leq q<+\infty$, with finite $\int_0^{+\infty}|f(q)|^2\,dq$, we again define $M,M^{-1}$ by
\begin{align*}
Mf(q)=F(p)&=\frac1{\sqrt h}\int_0^{\infty}e^{\frac{2\pi i}{h}pq}f(q)\,dq,\\
M^{-1}F(p)=f(q)&=\frac1{\sqrt h}\int_{-\infty}^{\infty}e^{-\frac{2\pi i}{h}pq}F(p)\,dp
\quad(=MF(-p)),
\end{align*}
then $M$ maps the Hilbert space $F_{\Omega'}$ of all $f(q)$, $0\leq q<\infty$, with finite $\int_0^{\infty}|f(q)|^2\,dq$, into a different Hilbert space: $F_{\Omega''}$ of all $F(p)$, $-\infty<p<\infty$, with finite $\int_{-\infty}^{\infty}|F(p)|^2\,dp$.

Although $\|Mf(q)\|=\|f(q)\|$ naturally still holds---by the theorems mentioned earlier, if we extend $f(q)$ by zero for $-\infty<q<0$---the equality $\|M^{-1}F(p)\|=\|F(p)\|$ by no means holds in general. For by those theorems, if we define $f(q)$ also for $q<0$ by $f(q)=\frac1{\sqrt h}\int_{-\infty}^{\infty}e^{-\frac{2\pi i}{h}pq}F(p)\,dp$, then
\[
\|F\|^2=\int_{-\infty}^{\infty}|F(p)|^2\,dp=\int_{-\infty}^{\infty}|f(q)|^2\,dq,
\qquad
\|M^{-1}F\|^2=\|f\|^2=\int_0^{\infty}|f(q)|^2\,dq.
\]

% Source PDF page 80 / original page 73.
\hypertarget{origpage-73}{}%
Thus $\|M^{-1}F\|<\|F\|$ unless the function $f(q)$ defined above happens to vanish for every $q<0$. Consequently $E'(\lambda)=ME(\lambda)M^{-1}$ is not a resolution of the identity at all;\origNote{90} the method has failed.

We shall soon see, in \hyperlink{note-105}{note 105}, that this is inherent in the situation, because no resolution of the identity belongs to this operator.

Before concluding these preliminary considerations, we give several formal rules for calculating with operators represented, through their resolutions of the identity, in the symbolic form
\[
A=\int_{-\infty}^{\infty}\lambda\,dE(\lambda).
\]
First let $F$ be a projection operator commuting with every $E(\lambda)$. Then for all $\lambda'<\lambda''$,
\begin{align*}
\|E(\lambda'')Ff-E(\lambda')Ff\|^2
&=\|(E(\lambda'')-E(\lambda'))Ff\|^2\\
&=\|F(E(\lambda'')-E(\lambda'))f\|^2
\leq\|(E(\lambda'')-E(\lambda'))f\|^2.
\end{align*}
Since $E(\lambda'')$, $E(\lambda')$, $E(\lambda'')-E(\lambda')$, and also $E(\lambda'')F$, $E(\lambda')F$, $E(\lambda'')F-E(\lambda')F=(E(\lambda'')-E(\lambda'))F$ are projection operators, this gives
\[
\|E(\lambda'')Ff\|^2-\|E(\lambda')Ff\|^2
\leq\|E(\lambda'')f\|^2-\|E(\lambda')f\|^2.
\]
Therefore $\int_{-\infty}^{\infty}\lambda^2\,d\|E(\lambda)f\|^2\geq\int_{-\infty}^{\infty}\lambda^2\,d\|E(\lambda)Ff\|^2$, and by $\hyperref[cond:II.S3]{S_3}$, $AFf$ is meaningful whenever $Af$ is. Furthermore,\origNote{91}
\[
AF=\int_{-\infty}^{\infty}\lambda\,d(E(\lambda)F)
=\int_{-\infty}^{\infty}\lambda\,d(FE(\lambda))=FA.
\]
Thus $A,F$ also commute. In particular, by $\hyperref[cond:II.Sbar2]{\overline{S}_2}$, we may take any $E(\lambda)$ for $F$. Then
\begin{align*}
AE(\lambda)&=\int_{-\infty}^{\infty}\lambda'\,d(E(\lambda')E(\lambda))
=\int_{-\infty}^{\infty}\lambda'\,d(E(\operatorname{Min}(\lambda,\lambda')))\\
&=\int_{-\infty}^{\lambda}+\int_\lambda^{\infty}
=\int_{-\infty}^{\lambda}\lambda'\,dE(\lambda')+\int_\lambda^{\infty}\lambda'\,dE(\lambda).
\end{align*}
The last integral vanishes because the expression after $d$ is constant. Hence
\[
AE(\lambda)=E(\lambda)A=\int_{-\infty}^{\lambda}\lambda'\,dE(\lambda').
\]
It follows further that\origNote{92}
\[
A^2=\int_{-\infty}^{\infty}\lambda\,d(E(\lambda)A)
=\int_{-\infty}^{\infty}\lambda\,d\left(\int_{-\infty}^{\lambda}\lambda'\,dE(\lambda')\right)
=\int_{-\infty}^{\infty}\lambda^2\,dE(\lambda).
\]

% Source PDF page 81 / original page 74.
\hypertarget{origpage-74}{}%
In general,
\[
A^n=\int_{-\infty}^{\infty}\lambda^n\,dE(\lambda),
\]
for induction from $n-1$ to $n$ gives\origNote{92}
\begin{align*}
A^n=A^{n-1}A
&=\int_{-\infty}^{\infty}\lambda^{n-1}\,d(E(\lambda)A)\\
&=\int_{-\infty}^{\infty}\lambda^{n-1}\,d\left(\int_{-\infty}^{\lambda}\lambda'\,dE(\lambda')\right)\\
&=\int_{-\infty}^{\infty}\lambda^{n-1}\lambda\,dE(\lambda)
=\int_{-\infty}^{\infty}\lambda^n\,dE(\lambda).
\end{align*}
Thus, for any polynomial $p(x)=a_0+a_1x+\cdots+a_nx^n$,
\[
p(A)=\int_{-\infty}^{\infty}p(\lambda)\,dE(\lambda).
\]
Here we naturally understand $p(A)$ to mean $a_01+a_1A+\cdots+a_nA^n$; the identity $\int_{-\infty}^{\infty}dE(\lambda)=1$ follows from $\hyperref[cond:II.Sbar1]{\overline{S}_1}$.

We further observe that if $r(\lambda),s(\lambda)$ are arbitrary functions and we define two operators $B,C$ symbolically by\origNote{93}
\[
B=\int_{-\infty}^{\infty}r(\lambda)\,dE(\lambda),\qquad
C=\int_{-\infty}^{\infty}s(\lambda)\,dE(\lambda),
\]
then
\[
BC=\int_{-\infty}^{\infty}r(\lambda)s(\lambda)\,dE(\lambda).
\]
The proof proceeds exactly as in the special case $B=C=A$:
\begin{align*}
BE(\lambda)&=\int_{-\infty}^{\infty}r(\lambda')\,d(E(\lambda')E(\lambda))
=\int_{-\infty}^{\infty}r(\lambda')\,d(E(\operatorname{Min}(\lambda,\lambda')))\\
&=\int_{-\infty}^{\lambda}+\int_\lambda^{\infty}
=\int_{-\infty}^{\lambda}r(\lambda')\,dE(\lambda')+\int_\lambda^{\infty}r(\lambda')\,dE(\lambda)\\
&=\int_{-\infty}^{\lambda}r(\lambda')\,dE(\lambda'),\\
CB&=\int_{-\infty}^{\infty}s(\lambda)\,d(BE(\lambda))
=\int_{-\infty}^{\infty}s(\lambda)\,d\left(\int_{-\infty}^{\lambda}r(\lambda')\,dE(\lambda')\right)\\
&=\int_{-\infty}^{\infty}s(\lambda)\cdot r(\lambda)\,dE(\lambda)
=\int_{-\infty}^{\infty}s(\lambda)r(\lambda)\,dE(\lambda).
\end{align*}
The following relations are trivial to verify:
\begin{gather*}
B^*=\int_{-\infty}^{\infty}\overline{r(\lambda)}\,dE(\lambda),\qquad
aB=\int_{-\infty}^{\infty}ar(\lambda)\,dE(\lambda),\\
B\pm C=\int_{-\infty}^{\infty}(r(\lambda)\pm s(\lambda))\,dE(\lambda).
\end{gather*}
There is therefore no formal obstacle to writing $B=r(A)$ for such functions $r(\lambda)$ as well.\origNote{94}

% Source PDF page 82 / original page 75.
\hypertarget{origpage-75}{}%
Especially noteworthy are the discontinuous functions
\[
e_\lambda(\lambda')=\begin{cases}1,&\lambda'\leq\lambda,\\0,&\lambda'>\lambda.\end{cases}
\]
For these, by $\hyperref[cond:II.Sbar1]{\overline{S}_1}$,
\[
e_\lambda(A)=\int_{-\infty}^{\infty}e_\lambda(\lambda')\,dE(\lambda)
=\int_{-\infty}^{\lambda}dE(\lambda')=E(\lambda).
\]
At the beginning of this section we discussed the operator $A=q_j\cdots$; its $E(\lambda)$ was multiplication by $1$ or $0$ according as $q_j\leq\lambda$ or $q_j>\lambda$, that is, multiplication by $e_\lambda(q)$. Thus $e_\lambda(q_j\cdots)=e_\lambda(q_j)\cdots$, which serves to illustrate these concepts.

\section{Digression on the Uniqueness and Solvability of the Eigenvalue Problem}\label{sec:II.9}
The preceding section gave only a qualitative survey, guided by particular examples, of the question of which resolutions of the identity $E(\lambda)$ belong to a given Hermitian operator $A$. We still lack a systematic investigation. A mathematically complete treatment is beyond the scope of this book. We must confine ourselves to proving some points and reporting the rest, all the more because a detailed knowledge of these matters is not absolutely necessary for understanding quantum mechanics.\origNote{95}

\hyperref[thm:II.18]{Theorem 18} showed that continuity of linear operations is expressed by
\begin{equation}
\|Af\|\leq C\|f\|,\tag{St}\label{eq:II.9.St}
\end{equation}
with arbitrary but fixed $C$. By \hyperref[thm:II.18]{Theorem 18} there are several equivalent forms of condition \hyperref[eq:II.9.St]{St}, namely
\begin{equation}
|(Af,g)|\leq C\|f\|\|g\|,\tag{$\mathrm{St}_1$}\label{eq:II.9.St1}
\end{equation}
\begin{equation}
|(Af,f)|\leq C\|f\|^2.\tag{$\mathrm{St}_2$}\label{eq:II.9.St2}
\end{equation}
The latter applies only to Hermitian $A$.

Condition $\hyperref[eq:II.9.St1]{\mathrm{St}_1}$, equivalent to continuity, is Hilbert's concept of boundedness. Hilbert formulated and solved the eigenvalue problem for bounded, that is, continuous, Hermitian operators; see \hyperlink{note-70}{note 70}. Before discussing this, however, we must introduce one further concept.

A Hermitian operator $A$ is called closed if it has the following property: whenever $f_1,f_2,\ldots$ is a sequence of points, all $Af_n$ are meaningful, and $f_n\to f$, $Af_n\to f^*$, then $Af$ is meaningful and equals $f^*$. Notice that continuity could be defined in a closely related way: if all $Af_n$ and $Af$ are meaningful and $f_n\to f$, then also $Af_n\to Af$. The difference is that closedness assumes the existence of a limit $f^*$ of the $Af_n$, and only under that assumption asserts its equality to $Af$, whereas continuity also asserts the existence of $f^*$.

% Source PDF page 83 / original page 76.
\hypertarget{origpage-76}{}%
Here are some examples. Let $\mathfrak{R}_\infty$ again be the space of all $f(q)$ with finite $\int_{-\infty}^{\infty}|f(q)|^2\,dq$, $-\infty<q<\infty$; let $A$ be the operator $q\cdots$, defined for all $f(q)$ with finite $\int_{-\infty}^{\infty}|f(q)|^2\,dq$ and $\int_{-\infty}^{\infty}q^2|f(q)|^2\,dq$; and let $A'$ be the operator $\frac{h}{2\pi i}\frac{\partial}{\partial q}$, defined for all everywhere differentiable functions with finite $\int_{-\infty}^{\infty}|f(q)|^2\,dq$ and $\int_{-\infty}^{\infty}|\frac{h}{2\pi i}f'(q)|^2\,dq$. Both, as we know, are Hermitian. The operator $A$ is closed. For suppose $f_n\to f$, $Af_n\to f^*$, that is,
\[
\int_{-\infty}^{\infty}|f_n(q)-f(q)|^2\,dq\to0,\qquad
\int_{-\infty}^{\infty}|qf_n(q)-f^*(q)|^2\,dq\to0.
\]
By the reasoning used to prove \hyperref[ax:II.D]{D} in \hyperref[sec:II.3]{II.3}, there is a subsequence $f_{n_1},f_{n_2},\ldots$ of $f_1,f_2,\ldots$ converging to a limit everywhere except on a set of $q$'s of measure zero: $f_{n_\nu}(q)\to g(q)$. Hence
\[
\int_{-\infty}^{\infty}|g(q)-f(q)|^2\,dq=0,\qquad
\int_{-\infty}^{\infty}|qg(q)-f^*(q)|^2\,dq=0.
\]
Thus, except on a set of measure zero, $g(q)=f(q)$ and $qg(q)=f^*(q)$, hence also $qf(q)=f^*(q)$; that is, $f^*(q)$ and $qf(q)$ are not essentially different. But $f^*(q)$ belongs to $\mathfrak{R}_\infty$ by assumption, and therefore so does $qf(q)$. Consequently $Af(q)$ is meaningful and equals $qf(q)=f^*(q)$.

In contrast, $A'$ is not closed. Put $f_n(q)=e^{-\sqrt{q^2+1/n}}$, $f(q)=e^{-|q|}$. Clearly all $Af_n$ are meaningful but $Af$ is not, since $f$ is nondifferentiable at $q=0$. Nevertheless, as is readily calculated, $f_n\to f$, $Af_n\to f^*$ if $f^*(q)=-\operatorname{sgn}(q)e^{-|q|}$, where $\operatorname{sgn}(q)=-1,0,+1$ for $q<0$, $q=0$, $q>0$, respectively.\tnote{In this example the source writes $A$ where the differential operator $A'$ is meant and omits its factor $h/(2\pi i)$ from $f^*$. The printed notation and formula are retained.}

We now show that, unlike continuity, closedness is a property that can always be achieved for Hermitian operators with little effort, by the process of extension: leave the operator unchanged at all points of $\mathfrak{R}_\infty$ where it is defined, and define it at some additional points where it was not yet defined.

Let $A$ be any Hermitian operator. Define an operator $\widetilde A$ as follows: $\widetilde Af$ is meaningful if there is a sequence $f_1,f_2,\ldots$ with all $Af_n$ meaningful, such that $f$ is the limit of $f_n$ and the $Af_n$ also have a limit $f^*$; then put $\widetilde Af=f^*$. This definition is admissible only if it is unambiguous, that is, if $f_n\to f$, $g_n\to f$, $Af_n\to f^*$, $Ag_n\to g^*$ imply $f^*=g^*$. Indeed, whenever $Ag$ is meaningful,
\begin{align*}
(f^*,g)&=\lim(Af_n,g)=\lim(f_n,Ag)=(f,Ag),\\
(g^*,g)&=\lim(Ag_n,g)=\lim(g_n,Ag)=(f,Ag).
\end{align*}

% Source PDF page 84 / original page 77.
\hypertarget{origpage-77}{}%
Thus $(f^*,g)=(g^*,g)$. The set of these $g$ is everywhere dense, so $f^*=g^*$, and $\widetilde A$ is well defined. This $\widetilde A$ is an extension of $A$: whenever $Af$ is meaningful, $\widetilde Af$ is meaningful too and equals $Af$. It suffices to put all $f_n=f$ and $f^*=Af$. By passage to the limit, the linearity and Hermitian character of $A$ imply the same properties for $\widetilde A$. Finally, $\widetilde A$ is closed. Suppose all $\widetilde Af_n$ are meaningful and $f_n\to f$, $\widetilde Af_n\to f^*$. There are sequences $f_{n,1},f_{n,2},\ldots$ with $Af_{n,m}$ meaningful, $f_{n,m}\to f_n$, $Af_{n,m}\to f_n^*$, and $\widetilde Af_n=f_n^*$. For each $n$ there is an $N_n$ such that for $m\geq N_n$,
\[
\|f_{n,m}-f_n\|\leq\frac1n,\qquad\|Af_{n,m}-f_n^*\|\leq\frac1n.
\]
Hence $f_{n,N_n}-f_n\to0$, $Af_{n,N_n}-\widetilde Af_n\to0$, so $f_{n,N_n}-f\to0$, $Af_{n,N_n}-f^*\to0$. By definition, $\widetilde Af=f^*$.

Notice that a discontinuous operator can never be made continuous by extension.

If an operator $B$ extends an operator $A$, that is, whenever $Af$ is meaningful, $Bf$ is meaningful and equals $Af$, we write $B\succ A$ or $A\prec B$. We have just proved $A\prec\widetilde A$, with $\widetilde A$ Hermitian and closed. It is immediately clear that for every closed $B$ with $A\prec B$, we must also have $\widetilde A\prec B$. Thus $\widetilde A$ is the smallest closed extension of $A$, and $\widetilde{\widetilde A}=\widetilde A$.

The close relation between $A$ and $\widetilde A$ suggests replacing $A$ throughout by $\widetilde A$, since this extends the domain of $A$ to its natural extent; viewed in relation to the domain of $\widetilde A$, the original domain is unnecessarily restricted. We make this replacement, and may henceforth assume all Hermitian operators to be closed.

Consider also a continuous Hermitian operator $A$. For it, closedness is equivalent to closedness of its domain. The condition $\|Af\|\leq C\|f\|$, characteristic of continuity, plainly also holds for $\widetilde A$, so $\widetilde A$ too is continuous. Its domain is then closed, but must also be everywhere dense, and consequently equals $\mathfrak{R}_\infty$. Thus $\widetilde A$ is everywhere defined, as is every closed continuous operator. The converse holds as well: if a closed operator is everywhere defined, it is continuous. This is Toeplitz's theorem,\origNote{96} whose proof we cannot discuss here.

Hilbert's result is that every continuous Hermitian operator has a unique resolution of the identity; see the reference in \hyperlink{note-70}{note 70}. Since a continuous operator is defined everywhere, $\int\lambda^2\,d\|E(\lambda)f\|^2$ must always be finite. Since it also equals $\|Af\|^2$, and hence is at most $C^2\|f\|^2$ by \hyperref[eq:II.9.St]{St}, we have

% Source PDF page 85 / original page 78.
\hypertarget{origpage-78}{}%
\begin{align*}
0&\geq\|Af\|^2-C^2\|f\|^2\\
&=\int_{-\infty}^{\infty}\lambda^2\,d\|E(\lambda)f\|^2-C^2\int_{-\infty}^{\infty}d\|E(\lambda)f\|^2\\
&=\int_{-\infty}^{\infty}(\lambda^2-C^2)\,d\|E(\lambda)f\|^2.
\end{align*}
Now put $f=E(-C-\varepsilon)g$. Then $E(\lambda)f=E(\operatorname{Min}(\lambda,-C-\varepsilon))g$, which is constant for $\lambda\geq-C-\varepsilon$, so we need consider only $\int_{-\infty}^{-C-\varepsilon}$. There $E(\lambda)f=E(\lambda)g$ and $\lambda^2-C^2\geq(C+\varepsilon)^2-C^2>2C\varepsilon$. Hence
\begin{gather*}
0\geq2C\varepsilon\int_{-\infty}^{-C-\varepsilon}d\|E(\lambda)g\|^2
=2C\varepsilon\|E(-C-\varepsilon)g\|^2,\\
\|E(-C-\varepsilon)g\|^2\leq0,\qquad E(-C-\varepsilon)g=0.
\end{gather*}
Similarly, taking $f=g-E(C+\varepsilon)g$, one proves $g-E(C+\varepsilon)g=0$.

Thus for every $\varepsilon>0$, $E(-C-\varepsilon)=0$, $E(C+\varepsilon)=1$; that is, $E(\lambda)=0$ for $\lambda<-C$ and $E(\lambda)=1$ for $\lambda>C$. By $\hyperref[cond:II.Sbar2]{\overline{S}_2}$, the latter also holds for $\lambda=C$. Thus $E(\lambda)$ varies only in $-C\leq\lambda\leq C$. Conversely, this implies continuity of $A$:
\[
\begin{aligned}
\|Af\|^2&=\int_{-\infty}^{\infty}\lambda^2\,d\|E(\lambda)f\|^2
=\int_{-C}^C\lambda^2\,d\|E(\lambda)f\|^2\\
&\leq C^2\int_{-C}^C d\|Ef\|^2
=C^2\int_{-\infty}^{\infty}d\|E(\lambda)f\|^2=C^2\|f\|^2,
\end{aligned}
\qquad\|Af\|\leq C\|f\|.
\]
We see, then, that resolutions of the identity varying only within finite $\lambda$-intervals account for precisely all continuous $A$. But what of the remaining, discontinuous Hermitian operators? All resolutions of the identity which vary at arbitrarily large $\lambda$ are still available. Do they account for all these Hermitian operators?

We must first understand the significance of the fact that these operators cannot be defined everywhere.

It is conceivable that a Hermitian operator is undefined at points of Hilbert space where it could reasonably be defined. For example, our operator $A'=\frac{h}{2\pi i}\frac{\partial}{\partial q}$ was undefined for $f(q)=e^{-|q|}$; and we could also have restricted $\frac{h}{2\pi i}\frac{\partial}{\partial q}$ to analytic functions on $-\infty<q<+\infty$, with real $q$,\origNote{97} and so forth. Our requirement that its domain be everywhere dense does protect it from excessively arbitrary restrictions, and we can also confine ourselves to closed operators. Yet even this is not effective enough.

% Source PDF page 86 / original page 79.
\hypertarget{origpage-79}{}%
For example, take the operator $A'=\frac{h}{2\pi i}\frac{\partial}{\partial q}$ on the interval $0\leq q\leq1$. Thus $f(q)$ is assumed to be differentiable everywhere, with $\int_0^1|f(q)|^2\,dq$ and $\int_0^1|f'(q)|^2\,dq$ finite. To make $A'$ Hermitian we must impose a further boundary condition, $f(0):f(1)=e^{-i\alpha}$, $0\leq\alpha<2\pi$. Denote the domain of these $f(q)$ by $\mathfrak A_\alpha$, and $A'$ itself in this case by $A'_\alpha$. We also consider the boundary condition $f(1)=f(0)=0$; denote the corresponding domain of $f(q)$ by $\mathfrak A^0$, and $A'$ itself by $A'{}^0$. All the $A'_\alpha$ are extensions of $A'{}^0$, which is therefore Hermitian and has an everywhere dense domain.\origNote{98} Consequently their closures $\widetilde A'_\alpha$ are also extensions of $\widetilde A'{}^0$. All the $\widetilde A'_\alpha$ are distinct from one another and from $\widetilde A'{}^0$. For the manifestly unitary operation $f(q)\mapsto e^{i\beta q}f(q)$ carries $A'$ into $A'+\frac{h\beta}{2\pi}1$, and $\mathfrak A_\alpha$ into $\mathfrak A_{\alpha+\beta}$, $\mathfrak A^0$ into $\mathfrak A^0$; thus it carries $A'_\alpha$ into $A'_{\alpha-\beta}+\frac{h\beta}{2\pi}1$, $A'{}^0$ into $A'{}^0+\frac{h\beta}{2\pi}1$, and therefore $\widetilde A'_\alpha$ into $\widetilde A'_{\alpha-\beta}+\frac{h\beta}{2\pi}1$, $\widetilde A'{}^0$ into $\widetilde A'{}^0+\frac{h\beta}{2\pi}1$. Hence $\widetilde A'_\alpha=\widetilde A'{}^0$ would imply $\widetilde A'_{\alpha-\beta}=\widetilde A'{}^0$, so that all the $\widetilde A'_\gamma$ would be equal. It therefore suffices to show that $\widetilde A'_\alpha\ne\widetilde A'_\gamma$ for $\alpha\ne\gamma$. This is the case if $A'_\alpha,A'_\gamma$ have no common Hermitian extension, that is, if $A'$ is not Hermitian on the union of $\mathfrak A_\alpha$ and $\mathfrak A_\gamma$. Since $e^{i\alpha q}$ belongs to $\mathfrak A_\alpha$ and $e^{i\gamma q}$ to $\mathfrak A_\gamma$, and
\begin{align*}
(A'e^{i\alpha q},e^{i\gamma q})-(e^{i\alpha q},A'e^{i\gamma q})
&=i\alpha\int_0^1e^{i(\alpha-\gamma)q}\,dq-i\gamma\int_0^1e^{i(\alpha-\gamma)q}\,dq\\
&=\int_0^1e^{i(\alpha-\gamma)q}i(\alpha-\gamma)\,dq
=e^{i(\alpha-\gamma)}-1\ne0,
\end{align*}
this is indeed the case. Thus the closed Hermitian operator $\widetilde A'{}^0$ is defined on too narrow a domain: it has proper closed Hermitian extensions, namely the $\widetilde A'_\alpha$, where ``proper'' means different from itself. Moreover, there are infinitely many possible choices of extension: every $\widetilde A'_\alpha$ is admissible, and each gives a different solution of the eigenvalue problem. The spectrum is always a pure point spectrum, but depends on $\alpha$, being the set $h(\frac{\alpha}{2\pi}+k)$, $k=0,\pm1,\pm2,\ldots$. For the operator $\widetilde A'{}^0$ itself, on the other hand, we cannot expect any reasonable solution of the eigenvalue problem at all. Indeed, later in this section we shall show that a Hermitian operator belonging to a resolution of the identity---that is, one for which the eigenvalue problem is solvable---has no proper extensions. An operator with no proper extensions, which is therefore already defined at every point where it could reasonably be defined, without destroying its Hermitian character, will be called \emph{maximal}. We have thus seen that a resolution of the identity can belong only to maximal operators.

% Source PDF page 87 / original page 80.
\hypertarget{origpage-80}{}%
It can be shown that every Hermitian operator can be extended to a maximal one; for a nonmaximal closed operator this can always be done in infinitely many ways. Thus the only uniquely determined stage in the extension is the passage from $A$ to its closure $\widetilde A$; cf.\ the end of \hyperlink{note-95}{note 95}. The most we can therefore hope for is that precisely one resolution of the identity belongs to each maximal Hermitian operator. Closed continuous Hermitian operators are defined everywhere, and are consequently maximal.

We now have to examine whether this expectation is justified. This involves two questions: the existence and the uniqueness of the resolution of the identity. The result will be that there is either no such resolution or precisely one; and the former case actually occurs. Thus the eigenvalue problem is uniquely determined but is sometimes insoluble. In a certain sense, however, insolubility is exceptional. We shall briefly indicate the method of proof.

If $f(\lambda)$ is a rational function and $A$ is a finite-dimensional matrix that can be transformed into diagonal form by a unitary transformation, then the eigenvectors of $A$ remain eigenvectors of $f(A)$, while the corresponding eigenvalues $\lambda_1,\ldots,\lambda_n$ become $f(\lambda_1),\ldots,f(\lambda_n)$.\origNote{99} If $f(\lambda)$ maps the real axis onto the unit circle, it transforms matrices with real eigenvalues into matrices whose eigenvalues have absolute value one: that is, it transforms Hermitian matrices into unitary matrices.\origNote{100} Such a function is $f(\lambda)=\frac{\lambda+i}{\lambda-i}$; the corresponding transformation, called the \emph{Cayley transformation}, is
\[
U=\frac{A-i1}{A+i1},\qquad A=-i\frac{U+1}{U-1}.
\]
We shall apply this transformation to our Hermitian operator $A$.\tnote{The scalar function printed above is the reciprocal of the function used in the displayed formula for $U$. The latter convention is used consistently in the subsequent argument; both printed forms have been retained.} Its meaning is that $Uf$ is defined if and only if $f$ can be written as $(A+i1)\varphi=A\varphi+i\varphi$, and is then $(A-i1)\varphi=A\varphi-i\varphi$. We hope that this determines a unique value of $Uf$ for every $f$, and that $U$ is unitary. The finite-dimensional proof cannot be used here, since it assumes the possibility of diagonalization, that is, a pure point spectrum. If the assertion is true, however, we can solve the eigenvalue problem by the following method.

For a unitary operator $U$, the eigenvalue problem has a unique solution: there is precisely one family of projections $E(\sigma)$, $0\leq\sigma\leq1$, satisfying the following conditions.

% Source PDF page 88 / original page 81.
\hypertarget{origpage-81}{}%
\begin{description}
\item[$\overline{\overline S}_1$.]\leavevmode\phantomsection\label{cond:II.Sdouble1} $E(0)=0$, $E(1)=1$; and as $\sigma\to\sigma_0$ with $\sigma\geq\sigma_0$, $E(\sigma)f\to E(\sigma_0)f$.
\item[$\overline{\overline S}_2$.]\leavevmode\phantomsection\label{cond:II.Sdouble2} If $\sigma'\leq\sigma''$, then $E(\sigma')\leq E(\sigma'')$.
\item[$\overline{\overline S}_3$.]\leavevmode\phantomsection\label{cond:II.Sdouble3} $(Uf,g)=\int_0^1e^{2\pi i\sigma}\,d(E(\sigma)f,g)$.
\end{description}
Here $Uf$ is always meaningful, and the integral always converges absolutely.\origNote{101} This follows from Hilbert's theory, since a unitary operator is continuous; cf.\ \hyperlink{note-70}{notes 70} and \hyperlink{note-101}{101}. The analogy with $\hyperref[cond:II.Sbar1]{\overline S_1}$--$\hyperref[cond:II.Sbar3]{\overline S_3}$ for Hermitian operators is evident. There are only two differences. Instead of the real integrand $\lambda$ ranging over the real axis, the complex integrand $e^{2\pi i\sigma}$ runs around the unit circle. And the description of the domain in $\hyperref[cond:II.Sbar3]{\overline S_3}$ is unnecessary here, since unitary operators are defined everywhere. Already in $\mathfrak{R}_n$, the relation between Hermitian and unitary matrices was largely analogous to that between the real axis and the unit circle; cf.\ \hyperlink{note-100}{note 100}.

By $\hyperref[cond:II.Sdouble1]{\overline{\overline S}_1}$, $E(\sigma)f\to0$ as $\sigma\to0$, $\sigma\geq0$. It need not follow, however, that $E(\sigma)f\to f$ as $\sigma\to1$, $\sigma\leq1$: $E(\sigma)$ may be discontinuous at $1$. By \hyperref[thm:II.17]{Theorem 17} in \hyperref[sec:II.4]{II.4}, the limit from the left is a projection $E'$; cf.\ \hyperlink{note-79}{note 79}. If $E'\ne1$, there are nonzero solutions of $E'f=0$. Since $E(\sigma)\leq E'$ for $\sigma<1$, this is equivalent to $E(\sigma)f=0$ for every $\sigma<1$. As at the beginning of \hyperref[sec:II.8]{II.8}, we obtain $(Uf,g)=(f,g)$ for every $g$, and hence $Uf=f$. Conversely, if $Uf=f$, then
\[
\int_0^1e^{2\pi i\sigma}\,d(E(\sigma)f,f)=(Uf,f)=(f,f),
\]
and consequently
\[
\mathfrak R\int_0^1e^{2\pi i\sigma}\,d(E(\sigma)f,f)=(f,f),
\qquad
\int_0^1(1-\cos2\pi\sigma)\,d(E(\sigma)f,f)=0,
\]
or
\[
\int_0^1(1-\cos2\pi\sigma)\,d\|E(\sigma)f\|^2=0.
\]
Thus $E(\sigma)f=0$ for every $0<\sigma<1$, and the same is automatic for $\sigma=0$. Discontinuity at $1$ is therefore equivalent to the existence of a nonzero solution of $Uf=f$. But in the Cayley transformation, $\varphi=Af+if$, $U\varphi=Af-if$; hence $U\varphi=\varphi$ implies $f=0$ and $\varphi=0$. In this case, therefore, $E(\sigma)f\to f$ also as $\sigma\to1$ from below.

% Source PDF page 89 / original page 82.
\hypertarget{origpage-82}{}%
Now
\[
\lambda=-i\frac{e^{2\pi i\sigma}+1}{e^{2\pi i\sigma}-1}=-\cot\pi\sigma,
\qquad \sigma=-\frac1\pi\operatorname{arccot}\lambda
\]
maps the interval $0<\sigma<1$ onto the whole real axis $-\infty<\lambda<+\infty$ in a one-to-one, monotonically increasing manner. Consequently, if we put
\begin{equation}
F(\lambda)=E\left(-\frac1\pi\operatorname{arccot}\lambda\right),\qquad
E=F(-\cot\pi\sigma),
\tag{C}\label{eq:II.9.C}
\end{equation}
then $F(\lambda)$ satisfies $\hyperref[cond:II.Sbar1]{\overline S_1},\hyperref[cond:II.Sbar2]{\overline S_2}$. We shall show that $F(\lambda)$ satisfies $\hyperref[cond:II.Sbar3]{\overline S_3}$ for $A$ if and only if $E(\sigma)$ satisfies $\hyperref[cond:II.Sdouble3]{\overline{\overline S}_3}$ for $U$. This reduces the existence and uniqueness questions for the possibly discontinuous Hermitian operator $A$ to the corresponding questions for the unitary operator $U$, where, as already noted, both have affirmative answers.

Suppose first that $A$ is Hermitian and that its Cayley transform $U$ is unitary. Then $E(\sigma)$ exists and satisfies $\hyperref[cond:II.Sdouble1]{\overline{\overline S}_1}$--$\hyperref[cond:II.Sdouble3]{\overline{\overline S}_3}$, while $F(\lambda)$, defined by \hyperref[eq:II.9.C]{(C)}, satisfies $\hyperref[cond:II.Sbar1]{\overline S_1},\hyperref[cond:II.Sbar2]{\overline S_2}$. If $Af$ is defined, put $\varphi=Af+if$, so that $U\varphi=Af-if$, $f=(\varphi-U\varphi)/(2i)$, and $Af=(\varphi+U\varphi)/2$. In a partly symbolic notation,\origNote{102} we obtain
\begin{align*}
f&=\frac1{2i}(\varphi-U\varphi)
=\frac1{2i}\left(\varphi-\int_0^1e^{2\pi i\sigma}\,dE(\sigma)\varphi\right)\\
&=\int_0^1\frac{1-e^{2\pi i\sigma}}{2i}\,dE(\sigma)\varphi,\\
E(\sigma)f&=\int_0^1\frac{1-e^{2\pi i\sigma'}}{2i}\,d(E(\sigma)E(\sigma')\varphi)\\
&=\int_0^1\frac{1-e^{2\pi i\sigma'}}{2i}\,d(E(\operatorname{Min}(\sigma,\sigma'))\varphi)\\
&=\int_0^\sigma\frac{1-e^{2\pi i\sigma'}}{2i}\,dE(\sigma')\varphi.
\end{align*}
Hence
\begin{align*}
\|E(\sigma)f\|^2&=(E(\sigma)f,f)
=\int_0^\sigma\frac{1-e^{2\pi i\sigma'}}{2i}\,d(E(\sigma')\varphi,f)\\
&=\int_0^\sigma\frac{1-e^{2\pi i\sigma'}}{2i}\,d\overline{(E(\sigma')f,\varphi)}\\
&=\int_0^\sigma\frac{1-e^{2\pi i\sigma'}}{2i}\,
d\left(\int_0^{\sigma'}\frac{1-e^{-2\pi i\sigma''}}{-2i}\,
d\overline{(E(\sigma'')\varphi,\varphi)}\right)\\
&=\int_0^\sigma\frac{1-e^{2\pi i\sigma'}}{2i}
\frac{1-e^{-2\pi i\sigma'}}{-2i}\,d\overline{(E(\sigma')\varphi,\varphi)}\\
&=\int_0^\sigma\frac{(1-e^{2\pi i\sigma'})(1-e^{-2\pi i\sigma'})}{4}\,
d\|E(\sigma')\varphi\|^2\\
&=\int_0^\sigma\sin^2\pi\sigma'\,d\|E(\sigma')\varphi\|^2.
\end{align*}

% Source PDF page 90 / original page 83.
\hypertarget{origpage-83}{}%
For the integral in $\hyperref[cond:II.Sbar3]{\overline S_3}$ this gives
\begin{align*}
\int_{-\infty}^{+\infty}\lambda^2\,d\|F(\lambda)f\|^2
&=\int_0^1\cot^2\pi\sigma\,d\|E(\sigma)f\|^2\\
&=\int_0^1\cot^2\pi\sigma\,
d\left(\int_0^\sigma\sin^2\pi\sigma'\,d\|E(\sigma')\varphi\|^2\right)\\
&=\int_0^1\cot^2\pi\sigma\sin^2\pi\sigma\,d\|E(\sigma)\varphi\|^2\\
&=\int_0^1\cos^2\pi\sigma\,d\|E(\sigma)\varphi\|^2.
\end{align*}
This integral is finite, since it is at most $\int_0^1d\|E(\sigma)\varphi\|^2=\|\varphi\|^2$. Furthermore,
\begin{align*}
Af&=\frac12(\varphi+U\varphi)
=\frac12\left(\varphi+\int_0^1e^{2\pi i\sigma}\,dE(\sigma)\varphi\right)\\
&=\int_0^1\frac{1+e^{2\pi i\sigma}}2\,dE(\sigma)\varphi\\
&=\int_0^1-i\frac{e^{2\pi i\sigma}+1}{e^{2\pi i\sigma}-1}
\frac{1-e^{2\pi i\sigma}}{2i}\,dE(\sigma)\varphi\\
&=\int_0^1-\cot\pi\sigma\,
d\left(\int_0^\sigma\frac{1-e^{2\pi i\sigma'}}{2i}\,dE(\sigma')\varphi\right)\\
&=\int_0^1-\cot\pi\sigma\,dE(\sigma)f
=\int_{-\infty}^{+\infty}\lambda\,dF(\lambda)f.
\end{align*}
Thus the final assertion in $\hyperref[cond:II.Sbar3]{\overline S_3}$ also holds. In other words, $A$ is an extension of the operator belonging to $F(\lambda)$. Since the latter is maximal, as we shall prove, it must coincide with $A$.\origNote{103}\tnote{The preceding calculation shows, for every $f$ in the domain of $A$, that $Af$ equals the spectral integral defined by $F$. It therefore directly shows that $A$ is a restriction of that operator. The source reverses the direction of extension here; its statement is retained.}

Conversely, suppose that $F(\lambda)$ belongs to $A$, that is, that $\hyperref[cond:II.Sbar1]{\overline S_1}$--$\hyperref[cond:II.Sbar3]{\overline S_3}$ hold. Define $E(\sigma)$ by \hyperref[eq:II.9.C]{(C)}; then $\hyperref[cond:II.Sdouble1]{\overline{\overline S}_1},\hyperref[cond:II.Sdouble2]{\overline{\overline S}_2}$ hold. For arbitrary $\varphi$, put
\begin{align*}
f&=\int_{-\infty}^{+\infty}\frac1{\lambda+i}\,dF(\lambda)\varphi
=\int_0^1\frac1{-\cot\pi\sigma+i}\,dE(\sigma)\varphi\\
&=\int_0^1\frac{1-e^{2\pi i\sigma}}{2i}\,dE(\sigma)\varphi.
\end{align*}
This integral converges, since its integrand is bounded. As before,
\begin{align*}
F(\lambda)f=E(\sigma)f
&=\int_0^1\frac{1-e^{2\pi i\sigma'}}{2i}\,dE(\sigma)E(\sigma')\varphi\\
&=\int_0^1\frac{1-e^{2\pi i\sigma'}}{2i}\,dE(\operatorname{Min}(\sigma,\sigma'))\varphi\\
&=\int_0^\sigma\frac{1-e^{2\pi i\sigma'}}{2i}\,dE(\sigma')\varphi.
\end{align*}

% Source PDF page 91 / original page 84.
\hypertarget{origpage-84}{}%
Thus
\begin{align*}
Af&=\int_{-\infty}^{+\infty}\lambda\,dF(\lambda)f
=\int_0^1-\cot(\pi\sigma)\,dE(\sigma)f\\
&=\int_0^1-i\frac{e^{2\pi i\sigma}+1}{e^{2\pi i\sigma}-1}\,
d\left(\int_0^\sigma\frac{1-e^{2\pi i\sigma'}}{2i}\,dE(\sigma')\varphi\right)\\
&=\int_0^1-i\frac{e^{2\pi i\sigma}+1}{e^{2\pi i\sigma}-1}
\frac{1-e^{2\pi i\sigma}}{2i}\,dE(\sigma)\varphi
=\int_0^1\frac{1+e^{2\pi i\sigma}}2\,dE(\sigma)\varphi,
\end{align*}
and hence
\[
Af+if=\int_0^1dE(\sigma)\varphi=\varphi,\qquad
Af-if=\int_0^1e^{2\pi i\sigma}\,dE(\sigma)\varphi.
\]
Thus $U\varphi$ is meaningful and equals $\int_0^1e^{2\pi i\sigma}\,dE(\sigma)\varphi$. Since $\varphi$ was arbitrary, $U$ is meaningful everywhere. By taking the inner product with an arbitrary $\psi$ and then the complex conjugate, we also see that $U^*\psi$ equals $\int_0^1e^{-2\pi i\sigma}\,dE(\sigma)\psi$. The final calculation in \hyperref[sec:II.8]{II.8} then shows that $UU^*=U^*U=1$: $U$ is unitary and belongs to $E(\sigma)$.

The solvability of the eigenvalue problem for $A$ is therefore equivalent to the unitarity of its Cayley transform $U$, and uniqueness is then assured. It remains only to decide whether we can always form $U$, and when it is unitary. To settle these questions, we again start from a closed Hermitian operator $A$.

We defined $U$ as follows: $U\varphi$ is meaningful if and only if $\varphi=Af+if$, and is then equal to $Af-if$. First, however, we must show that this prescription is unambiguous: each $\varphi$ can arise from at most one $f$. In other words, $Af+if=Ag+ig$ must imply $f=g$; or, by the linearity of $A$, $Af+if=0$ must imply $f=0$. We have
\begin{align*}
\|Af\pm if\|^2&=(Af\pm if,Af\pm if)\\
&=(Af,Af)\pm(Af,if)\pm(if,Af)+(if,if)\\
&=\|Af\|^2\mp i(Af,f)\pm i(Af,f)+\|f\|^2
=\|Af\|^2+\|f\|^2.
\end{align*}
Thus $Af+if=0$ implies $\|f\|^2\leq\|Af+if\|^2=0$, and hence $f=0$, which justifies our definition. Secondly, $\|Af+if\|=\|Af-if\|$, that is, $\|U\varphi\|=\|\varphi\|$. Thus $U$ is continuous wherever it is defined. Furthermore, let $\mathfrak E$ be the domain of $U$, the set of all $Af+if$, and $\mathfrak F$ its range, the set of all $U\varphi$, or equivalently of all $Af-if$. Since $A$ and $U$ are linear, $\mathfrak E,\mathfrak F$ are linear manifolds; they are also closed. For if $\varphi$ is a limit point of $\mathfrak E$ or $\mathfrak F$, respectively, there is a sequence $\varphi_1,\varphi_2,\ldots$ in the corresponding manifold with $\varphi_n\to\varphi$.

% Source PDF page 92 / original page 85.
\hypertarget{origpage-85}{}%
Thus $\varphi_n=Af_n\pm if_n$. Since the $\varphi_n$ converge, they satisfy Cauchy's convergence condition; cf.\ \hyperref[ax:II.D]{\textbf{D}} in \hyperref[sec:II.1]{II.1}. By
\[
\|f_m-f_n\|\leq\|A(f_m-f_n)\pm i(f_m-f_n)\|=\|\varphi_m-\varphi_n\|,
\]
the $f_n$ certainly satisfy this condition as well; and by
\[
\|Af_m-Af_n\|=\|A(f_m-f_n)\|
\leq\|A(f_m-f_n)\pm i(f_m-f_n)\|=\|\varphi_m-\varphi_n\|,
\]
so do the $Af_m$. Thus both $f_1,f_2,\ldots$ and $Af_1,Af_2,\ldots$ converge, by \hyperref[ax:II.D]{\textbf{D}} in \hyperref[sec:II.1]{II.1}: $f_n\to f$, $Af_n\to f^*$. Since $A$ is closed, $Af$ is meaningful and equals $f^*$. Hence
\[
\varphi_n=Af_n\pm if_n\to f^*\pm if=Af\pm if,\qquad
\varphi_n\to\varphi,
\]
so that $\varphi=Af\pm if$; thus $\varphi$ too belongs to $\mathfrak E$ or $\mathfrak F$, respectively.

Consequently $U$ is defined on the closed linear manifold $\mathfrak E$ and maps it onto the closed linear manifold $\mathfrak F$. It is linear and, since $\|Uf-Ug\|=\|U(f-g)\|=\|f-g\|$, leaves all distances unchanged: we shall call it \emph{length-preserving}. Thus $f\ne g$ implies $Uf\ne Ug$, so the mapping is one-to-one. We also have $(f,g)=(Uf,Ug)$, as is proved in exactly the same way as for unitary operators in \hyperref[sec:II.5]{II.5}. Hence $U$ also leaves all inner products unchanged. But $U$ is plainly unitary if and only if $\mathfrak E=\mathfrak F=\mathfrak R_\infty$.

Now let $A,B$ be two closed Hermitian operators, $U,V$ their Cayley transforms, and let the sets above be $\mathfrak E,\mathfrak F$ and $\mathfrak G,\mathfrak H$, respectively. We immediately see that if $B$ is a proper extension of $A$, then $V$ is a proper extension of $U$: $\mathfrak E$ is a proper subset of $\mathfrak G$, and $\mathfrak F$ a proper subset of $\mathfrak H$. Hence $\mathfrak E\ne\mathfrak R_\infty$, $\mathfrak F\ne\mathfrak R_\infty$. Then $U$ is not unitary, and the eigenvalue problem for $A$ is insoluble. This proves the theorem we have cited several times: if the eigenvalue problem for $A$ is soluble, then $A$ has no proper extensions; that is, $A$ is maximal.

Let us return to the single closed Hermitian operator $A$, with its $\mathfrak E,\mathfrak F,U$. If $Af$ is meaningful, then for $Af+if=\varphi$, $U\varphi$ is meaningful and equals $Af-if$. Hence $f=\frac1{2i}(\varphi-U\varphi)$, $Af=\frac12(\varphi+U\varphi)$, or, on putting $\psi=\varphi/(2i)$, $f=\psi-U\psi$, $Af=i(\psi+U\psi)$. Conversely, if $f=\psi-U\psi$, then $Af$ is certainly meaningful: since $U\psi$ is meaningful, $\psi=Af'+if'$ with $Af'$ meaningful, and $U\psi=Af'-if'$, so that $f=\psi-U\psi=2if'$. The domain of $A$ is therefore the set of all $\psi-U\psi$, and for $f=\psi-U\psi$ we have $Af=i(\psi+U\psi)$. Thus $A$, too, is uniquely determined by $U$ and its $\mathfrak E,\mathfrak F$. At the same time, we see that the $\psi-U\psi$ must be everywhere dense, since they form the domain of $A$.

Conversely, let us start from two closed linear manifolds $\mathfrak E,\mathfrak F$ and a linear, length-preserving mapping $U$ from $\mathfrak E$ onto $\mathfrak F$. Does there then exist a Hermitian $A$ whose Cayley transform is this $U$?

% Source PDF page 93 / original page 86.
\hypertarget{origpage-86}{}%
At any rate, the $\psi-U\psi$ must be everywhere dense, so let us assume this. By what we have just said, the proposed $A$ is then uniquely specified. The remaining questions are whether this definition is possible, whether $A$ really is Hermitian, and whether $U$ really is its Cayley transform. The first point is settled once we show that any $f$ of the form $f=\varphi-U\varphi$ determines $\varphi$ uniquely; that is, if $\varphi-U\varphi=\psi-U\psi$ implies $\varphi=\psi$, or if $\varphi-U\varphi=0$ implies $\varphi=0$. Indeed, suppose $\varphi-U\varphi=0$. For $g=\psi-U\psi$ we then have
\begin{align*}
(\varphi,g)&=(\varphi,\psi)-(\varphi,U\psi)\\
&=(U\varphi,U\psi)-(\varphi,U\psi)
=(U\varphi-\varphi,U\psi)=0;
\end{align*}
since the set of these $g$ is everywhere dense, $\varphi=0$.

For the second point we must prove $(Af,g)=(f,Ag)$: interchanging $f,g$ must turn $(Af,g)$ into its complex conjugate. Let $f=\varphi-U\varphi$, $g=\psi-U\psi$; then $Af=i(\varphi+U\varphi)$ and
\begin{align*}
(Af,g)&=(i(\varphi+U\varphi),\psi-U\psi)\\
&=i(\varphi,\psi)+i(U\varphi,\psi)-i(\varphi,U\psi)-i(U\varphi,U\psi)\\
&=i\big[(U\varphi,\psi)-\overline{(U\psi,\varphi)}\big]
=i(U\varphi,\psi)+\overline{i(U\psi,\varphi)}.
\end{align*}
This plainly has the required behavior when $f,g$, or equivalently $\varphi,\psi$, are interchanged. For the third point, we argue as follows. Let $V$ be the Cayley transform of $A$. Its domain is the set of all
\[
Af+if=i(\varphi+U\varphi)+i(\varphi-U\varphi)=2i\varphi,
\]
that is, the domain of $U$; and there
\begin{align*}
V(2i\varphi)&=V(Af+if)=Af-if\\
&=i(\varphi+U\varphi)-i(\varphi-U\varphi)=2iU\varphi.
\end{align*}
Thus $V\varphi=U\varphi$, and hence $V=U$.

Assigning to each closed Hermitian operator $A$ its Cayley transform $U$ therefore gives a one-to-one correspondence with the linear, length-preserving mappings $U$ for which the set of all $\varphi-U\varphi$ is everywhere dense.\origNote{104} We can now describe all Hermitian extensions $B$ of $A$, since all length-preserving extensions $V$ of $U$ are easy to find. The $\varphi-V\varphi$ are automatically everywhere dense, because their subset consisting of the $\varphi-U\varphi$ already is. For $A$ to be maximal, $U$ must be maximal, and conversely. If $U$ is not maximal, then $\mathfrak E\ne\mathfrak R_\infty$, $\mathfrak F\ne\mathfrak R_\infty$; conversely, these inequalities imply that $U$ is not maximal. For then $\mathfrak R_\infty-\mathfrak E\ne0$, $\mathfrak R_\infty-\mathfrak F\ne0$. We can therefore choose a nonzero $\varphi_0$ from $\mathfrak R_\infty-\mathfrak E$ and a nonzero $\psi_0$ from $\mathfrak R_\infty-\mathfrak F$; by replacing them with $\varphi_0/\|\varphi_0\|$, $\psi_0/\|\psi_0\|$, we may even arrange $\|\varphi_0\|=\|\psi_0\|=1$. Now define an operator $V$ on $[\mathfrak E,\varphi_0]$ by putting $Vf=U\varphi+a\psi_0$ for $f=\varphi+a\varphi_0$, with $\varphi$ in $\mathfrak E$ and $a$ a number. Clearly $V$ is linear; since $\varphi$ is orthogonal to $\varphi_0$ and $U\varphi$ to $\psi_0$, we have

% Source PDF page 94 / original page 87.
\hypertarget{origpage-87}{}%
\[
\|f\|^2=\|\varphi\|^2+|a|^2,\qquad
\|Vf\|^2=\|U\varphi\|^2+|a|^2,
\]
and hence $\|Vf\|=\|f\|$: $V$ is length-preserving. Finally, $V$ is a proper extension of $U$. Thus the characteristic condition for $A$ to be maximal is $\mathfrak E=\mathfrak R_\infty$ or $\mathfrak F=\mathfrak R_\infty$.

If, on the other hand, $A$ is not maximal, both closed linear manifolds $\mathfrak R_\infty-\mathfrak E$, $\mathfrak R_\infty-\mathfrak F$ are nonzero. Let orthonormal systems spanning them be $\varphi_1,\ldots,\varphi_p$ and $\psi_1,\ldots,\psi_q$, respectively; cf.\ \hyperref[thm:II.9]{Theorem 9} in \hyperref[sec:II.2]{II.2}. Here $p,q$ may have any of the values $1,2,\ldots,\infty$; for $p=\infty$ or $q=\infty$, the respective $\varphi$ or $\psi$ sequence does not terminate. Put $r=\operatorname{Min}(p,q)$ and define $V$ on $[\mathfrak E,\varphi_1,\ldots,\varphi_r]$ as follows: for $f=\varphi+\sum_{\nu=1}^r a_\nu\varphi_\nu$, where $\varphi$ belongs to $\mathfrak E$ and $a_1,\ldots,a_r$ are numbers, put $Vf=U\varphi+\sum_{\nu=1}^r a_\nu\psi_\nu$. Again it is easy to see that $V$ is linear and length-preserving, and a proper extension of $U$. Its domain is $[\mathfrak E,\varphi_1,\ldots,\varphi_r]$, which for $r=p$ equals $[\mathfrak E,\mathfrak R_\infty-\mathfrak E]=\mathfrak R_\infty$; its range is $[\mathfrak F,\psi_1,\ldots,\psi_r]$, which for $r=q$ equals $[\mathfrak F,\mathfrak R_\infty-\mathfrak F]=\mathfrak R_\infty$. Thus at least one of them certainly equals $\mathfrak R_\infty$. Let $V$ be the Cayley transform of the Hermitian operator $B$. By what we have said, $B$ is an extension of $A$ and is maximal. Notice that the $\varphi$ and $\psi$ can be chosen in infinitely many ways: for example, we can replace $\psi_1$ by any $\theta\psi_1$ with $|\theta|=1$. The same is therefore true of $V$ and $B$.

This completes our discussion of the eigenvalue problem, with the following result. Whenever it is soluble, it has only one solution; for nonmaximal operators it is certainly insoluble. Nonmaximal operators can always be extended to maximal ones in infinitely many different ways. Throughout, we are speaking of closed Hermitian operators. The condition for maximality is not quite the same as that for solvability of the eigenvalue problem: the former is $\mathfrak E=\mathfrak R_\infty$ \emph{or} $\mathfrak F=\mathfrak R_\infty$, the latter $\mathfrak E=\mathfrak R_\infty$ \emph{and} $\mathfrak F=\mathfrak R_\infty$.

We shall not investigate more closely here those operators for which the former condition holds but the latter does not. For these operators the eigenvalue problem is insoluble, and maximality prevents us from remedying this by passing to a proper extension. They are characterized by $\mathfrak E=\mathfrak R_\infty$, $\mathfrak F\ne\mathfrak R_\infty$, or $\mathfrak E\ne\mathfrak R_\infty$, $\mathfrak F=\mathfrak R_\infty$. Such operators do indeed exist, and all arise from two simple normal forms, so that, compared with maximal operators having a soluble eigenvalue problem, they can be regarded as exceptional cases. The reader will find further details in the author's paper cited in \hyperlink{note-95}{note 95}. At any rate, such operators must for the present be excluded from quantum-mechanical considerations. As we shall see below, the resolution of the identity belonging to a Hermitian operator enters so essentially into all the concepts of quantum mechanics
% Source PDF page 95 / original page 88.
\hypertarget{origpage-88}{}%
that we cannot dispense with its existence, and hence with the solvability of the eigenvalue problem.\origNote{105} Accordingly, we shall as a rule admit only Hermitian operators whose eigenvalue problem is soluble. Since this property strengthens maximality, these operators are called \emph{hypermaximal}.\origNote{106} They are precisely the operators that correspond one-to-one to resolutions of the identity.

Finally, let us mention two classes of closed Hermitian operators that are certainly hypermaximal. First, the continuous operators: these are meaningful everywhere and hence maximal; and since their eigenvalue problem is soluble by Hilbert's theory, cf.\ the reference in \hyperlink{note-70}{note 70}, they are even hypermaximal. Second, the operators that are real in some realization of $\mathfrak R_\infty$, provided they are maximal. For the only difference in the definitions of $\mathfrak E,\mathfrak F$ was the sign of $i$, which cannot matter when everything else is real. Thus $\mathfrak E=\mathfrak R_\infty$ implies $\mathfrak F=\mathfrak R_\infty$, and conversely: maximality implies hypermaximality. Without assuming maximality, we can at any rate say that $\mathfrak R_\infty-\mathfrak E$ and $\mathfrak R_\infty-\mathfrak F$ have the same number of dimensions. Thus, in the terminology used above in investigating extensions, $p=q$, and hence $r=p=q$ and
\begin{align*}
[\mathfrak E,\varphi_1,\ldots,\varphi_r]&=[\mathfrak E,\mathfrak R_\infty-\mathfrak E]=\mathfrak R_\infty,\\
[\mathfrak F,\psi_1,\ldots,\psi_r]&=[\mathfrak F,\mathfrak R_\infty-\mathfrak F]=\mathfrak R_\infty.
\end{align*}
The extension obtained there is therefore hypermaximal. Real operators consequently always have hypermaximal extensions. The paper cited in \hyperlink{note-95}{note 95} shows that the same holds for all definite operators.

\section{Commuting Operators}\label{sec:II.10}
According to the definition given in \hyperref[sec:II.4]{II.4}, two operators $R,S$ commute if $RS=SR$. If the operators are not both defined everywhere, equality also requires the two sides to have the same domain. We shall initially consider only Hermitian operators that are defined everywhere, and hence continuous, so that this difficulty concerning domains does not arise. Together with $R,S$ we consider their resolutions of the identity $E(\lambda),F(\lambda)$.

The commutativity of $R,S$ means $(RSf,g)=(SRf,g)$ for all $f,g$, or $(Sf,Rg)=(Rf,Sg)$. Furthermore, commutativity of $R,S$ implies that of $R^n,S$, $n=0,1,2,\ldots$, and therefore that of $p(R),S$ for every polynomial $p(x)=a_0+a_1x+\cdots+a_nx^n$. Now, symbolically,
\[
R=\int_{-C}^C\lambda\,dE(\lambda),\qquad
s(R)=\int_{-C}^Cs(\lambda)\,dE(\lambda).
\]

% Source PDF page 96 / original page 89.
\hypertarget{origpage-89}{}%
Here $C$ is the constant introduced in \hyperref[sec:II.9]{II.9} for the continuous operator $R$, there called $A$; and $s(x)$ is any function; cf.\ \hyperref[sec:II.8]{II.8}, especially \hyperlink{note-94}{note 94}. For polynomials $s(x)$, $(s(R)f,Sg)=(Sf,s(R)g)$, and therefore
\begin{equation}
\int_{-C}^Cs(\lambda)\,d(E(\lambda)f,Sg)
=\int_{-C}^Cs(\lambda)\,d(Sf,E(\lambda)g).
\tag{*}\label{eq:II.10.star}
\end{equation}
Since every continuous function $s(x)$ can be approximated arbitrarily closely by polynomials, uniformly on $-C\leq x\leq C$, \hyperref[eq:II.10.star]{(*)} also holds for continuous $s(x)$. Now put
\[
s(x)=\begin{cases}\lambda_0-x,&x\leq\lambda_0,\\0,&x\geq\lambda_0.\end{cases}
\]
Then \hyperref[eq:II.10.star]{(*)} gives
\[
\int_{-C}^{\lambda_0}(\lambda_0-\lambda)\,d(E(\lambda)f,Sg)
=\int_{-C}^{\lambda_0}(\lambda_0-\lambda)\,d(Sf,E(\lambda)g).
\]
Replacing $\lambda_0$ by $\lambda_0+\varepsilon$, $\varepsilon>0$, subtracting, and dividing by $\varepsilon$, we obtain
\begin{align*}
&\int_{-C}^{\lambda_0}d(E(\lambda)f,Sg)
+\int_{\lambda_0}^{\lambda_0+\varepsilon}\frac{\lambda-\lambda_0}{\varepsilon}\,d(E(\lambda)f,Sg)\\
&\qquad=\int_{-C}^{\lambda_0}d(Sf,E(\lambda)g)
+\int_{\lambda_0}^{\lambda_0+\varepsilon}\frac{\lambda-\lambda_0}{\varepsilon}\,d(Sf,E(\lambda)g).
\end{align*}
Letting $\varepsilon\to0$, and noting $\hyperref[cond:II.Sbar1]{\overline S_1}$, we get
\[
\int_{-C}^{\lambda_0}d(E(\lambda)f,Sg)=\int_{-C}^{\lambda_0}d(Sf,E(\lambda)g),
\qquad (E(\lambda_0)f,Sg)=(Sf,E(\lambda_0)g).
\]
Thus all $E(\lambda_0)$ with $-C\leq\lambda_0\leq C$ commute with $S$. The others certainly do as well, since for $\lambda_0<-C$ or $\lambda_0>C$, respectively, $E(\lambda_0)=0$ or $1$. Hence, if $R$ commutes with $S$, so do all the $E(\lambda)$. Conversely, if every $E(\lambda)$ commutes with $S$, then \hyperref[eq:II.10.star]{(*)} holds for every function $s(x)$, and thus every $s(R)$ commutes with $S$. We may conclude, first, that $R$ commutes with $S$ if and only if every $E(\lambda)$ does; and second, that in this case all functions of $R$, the $s(R)$, commute with $S$ as well.

But an $E(\lambda)$ commutes with $S$ if and only if it commutes with every $F(\mu)$: apply our theorem to $S,E(\lambda)$ in place of $R,S$. Thus another characteristic condition for $R,S$ to commute is that every $E(\lambda)$ commute with every $F(\mu)$. Furthermore, as shown above, commutativity of $R,S$ implies that of $r(R),S$, and this, replacing $R,S$ by $S,r(R)$, implies that of $r(R),s(S)$.

% Source PDF page 97 / original page 90.
\hypertarget{origpage-90}{}%
If the Hermitian operators $R,S$ are not required to be continuous, the situation becomes more complicated, since the domains of $RS$ and $SR$ may be hard to determine. For example, $R\cdot0$ is always meaningful: $0f=0$ and $R\cdot0f=R(0f)=R0=0$. But $0\cdot R$ is meaningful only where $R$ is; cf.\ the discussion in \hyperref[sec:II.5]{II.5}. Thus, if $R$ is not meaningful everywhere, the different domains imply $R\cdot0\ne0\cdot R$, so that strictly speaking $R$ and $0$ do not commute. This state of affairs is quite inconvenient for our later purposes: $0$ ought to commute with all Hermitian operators, not merely continuous ones.\origNote{107} We shall therefore define commutativity differently for discontinuous $R,S$, restricting ourselves to the hypermaximal $R,S$, which by \hyperref[sec:II.9]{II.9} are in any case the only ones of interest. Our definition is this: $R,S$ shall commute in the new sense if all the $E(\lambda)$ commute in the old sense with all the $F(\mu)$, these again being their respective resolutions of the identity. For continuous $R,S$, the new definition is identical with the old one, as we know; but if $R$ or $S$, or both, are discontinuous, it need not be. An example is $R,0$: they did not commute in the old sense, but do in the new, since for $0$ every $F(\mu)$ is either $0$ or $1$,\origNote{108} and therefore commutes with the $E(\lambda)$.

We proved above that if $R,S$ are two commuting continuous Hermitian operators, then every function $r(R)$ of $R$ commutes with every function $s(S)$ of $S$. Since the premise is always satisfied for $R=S$, two functions $r(R),s(R)$ of the same operator always commute. This also follows from the multiplication formula at the end of \hyperref[sec:II.8]{II.8}: $r(R)s(R)=t(R)$, where $r(x)s(x)\equiv t(x)$. Moreover, if $r(x),s(x)$ are real, $r(R),s(R)$ are Hermitian: by \hyperref[sec:II.8]{II.8}, if $r(x)$ is real, then $(r(R))^*=\overline r(R)=r(R)$.

The converse also holds. If $A,B$ are two commuting Hermitian operators, there exists a Hermitian operator $R$ of which both are functions: $A=r(R)$, $B=s(R)$. Indeed, something stronger holds: given any finite or infinite set of commuting Hermitian operators $A,B,C,\ldots$, there is a Hermitian operator $R$ of which all the $A,B,C,\ldots$ are functions. We cannot prove this theorem here and must refer the reader to the literature.\origNote{109} For our purposes, the theorem is needed only for finitely many $A,B,C,\ldots$ with pure point spectra. We shall prove it below for this case; for the general case we can offer only a few preliminary remarks.

Thus let $A,B,C,\ldots$ be finitely many Hermitian operators with pure point spectra. For any number $\lambda$, denote by $\mathfrak L_\lambda$ the closed linear manifold spanned by all solutions of $Af=\lambda f$, and by $E_\lambda$ its projection operator.

% Source PDF page 98 / original page 91.
\hypertarget{origpage-91}{}%
Then $\lambda$ is a point eigenvalue of $A$ if and only if there are nonzero solutions $f$, that is, if $\mathfrak L_\lambda\ne(0)$, or $E_\lambda\ne0$. Correspondingly, form $\mathfrak M_\lambda,F_\lambda$ for $B$, $\mathfrak N_\lambda,G_\lambda$ for $C$, and so on. The equation $Af=\lambda f$ implies $ABf=BAf=B(\lambda f)=\lambda Bf$, so $Bf$, like $f$, belongs to $\mathfrak M_\lambda$. Since $E_\lambda f$ always belongs to $\mathfrak M_\lambda$, so does $BE_\lambda f$, and hence $E_\lambda BE_\lambda f=BE_\lambda f$. This holds identically, so $E_\lambda BE_\lambda=BE_\lambda$. Applying $*$ gives $E_\lambda BE_\lambda=E_\lambda B$, and therefore $E_\lambda B=BE_\lambda$. Just as the commutativity of $A,B$ has yielded that of $B,E_\lambda$, so the commutativity of $B,E_\lambda$ yields that of $E_\lambda,F_\mu$. Since $A,B$ were in no way distinguished among $A,B,C,\ldots$, we may say that all the $E_\lambda,F_\mu,G_\nu,\ldots$ commute with one another. Thus $K(\lambda\mu\nu\cdots)=E_\lambda F_\mu G_\nu\cdots$ is also a projection operator; denote its closed linear manifold by $\mathfrak K(\lambda\mu\nu\cdots)$. By \hyperref[thm:II.14]{Theorem 14} in \hyperref[sec:II.4]{II.4}, $\mathfrak K(\lambda\mu\nu\cdots)$ is the intersection of $\mathfrak L_\lambda,\mathfrak M_\mu,\mathfrak N_\nu,\ldots$, that is, the set of all simultaneous solutions of
\[
Af=\lambda f,\qquad Bf=\mu f,\qquad Cf=\nu f,\qquad\ldots.
\]
Let $\lambda,\mu,\nu,\ldots$ and $\lambda',\mu',\nu',\ldots$ be two distinct systems of numbers: $\lambda\ne\lambda'$, or $\mu\ne\mu'$, or $\nu\ne\nu'$, and so forth. Let $f$ belong to $\mathfrak K(\lambda\mu\nu\cdots)$ and $f'$ to $\mathfrak K(\lambda'\mu'\nu'\cdots)$. Then $f,f'$ are orthogonal: if $\lambda\ne\lambda'$, because $Af=\lambda f$, $Af'=\lambda'f'$; if $\mu\ne\mu'$, because $Bf=\mu f$, $Bf'=\mu'f'$; and so on. Thus all of $\mathfrak K(\lambda\mu\nu\cdots)$ is orthogonal to all of $\mathfrak K(\lambda'\mu'\nu'\cdots)$.

Since $A$ has a pure point spectrum, the $\mathfrak L_\lambda$ span all of $\mathfrak R_\infty$ as a closed linear manifold. Hence a nonzero $f$ cannot be orthogonal to all the $\mathfrak L_\lambda$: for at least one $\mathfrak L_\lambda$ its projection onto $\mathfrak L_\lambda$ must be nonzero, or $E_\lambda f\ne0$. Similarly, there must be a $\mu$ with $F_\mu f\ne0$, a $\nu$ with $G_\nu f\ne0$, and so on. Consequently, for every nonzero $f$, we can find a $\lambda$ with $E_\lambda f\ne0$, then a $\mu$ with $F_\mu(E_\lambda f)\ne0$, then a $\nu$ with $G_\nu(F_\mu(E_\lambda f))\ne0$, and so forth. Finally we obtain $\cdots G_\nu F_\mu E_\lambda f\ne0$, $E_\lambda F_\mu G_\nu\cdots f\ne0$, $K(\lambda\mu\nu\cdots)f\ne0$: $f$ is not orthogonal to $\mathfrak K(\lambda\mu\nu\cdots)$. In other words, an $f$ orthogonal to every $\mathfrak K(\lambda\mu\nu\cdots)$ is zero. Thus the $\mathfrak K(\lambda\mu\nu\cdots)$ together span all of $\mathfrak R_\infty$ as a closed linear manifold.

Now let $\varphi^{(1)}_{(\lambda\mu\nu\cdots)},\varphi^{(2)}_{(\lambda\mu\nu\cdots)},\ldots$ be an orthonormal system spanning the closed linear manifold $\mathfrak K(\lambda\mu\nu\cdots)$. This sequence is finite or infinite according to the dimension of $\mathfrak K(\lambda\mu\nu\cdots)$; when $\mathfrak K(\lambda\mu\nu\cdots)=0$, the sequence is empty. Every $\varphi^{(n)}_{(\lambda\mu\nu\cdots)}$ belongs to some $\mathfrak K(\lambda\mu\nu\cdots)$ and is therefore an eigenfunction of every $A,B,C,\ldots$. Any two distinct $\varphi^{(n)}_{(\lambda\mu\nu\cdots)}$ are orthogonal: if they have the same system $\lambda,\mu,\nu,\ldots$, by their construction; if they have different systems, because they belong to different $\mathfrak K(\lambda\mu\nu\cdots)$. All the $\varphi^{(n)}_{(\lambda\mu\nu\cdots)}$ span the same closed linear manifold as all the $\mathfrak K(\lambda\mu\nu\cdots)$, namely $\mathfrak R_\infty$. Thus they form a complete orthonormal system.

% Source PDF page 99 / original page 92.
\hypertarget{origpage-92}{}%
We have therefore constructed a complete orthonormal system consisting entirely of common eigenfunctions of $A,B,C,\ldots$. From now on we shall denote its elements by $\psi_1,\psi_2,\ldots$ and write the corresponding eigenvalue equations as
\[
A\psi_m=\lambda_m\psi_m,\qquad
B\psi_m=\mu_m\psi_m,\qquad
C\psi_m=\nu_m\psi_m,\qquad\ldots.
\]
Now choose any system of mutually distinct numbers $\varkappa_1,\varkappa_2,\varkappa_3,\ldots$, and form a Hermitian operator $R$ with pure point spectrum $\varkappa_1,\varkappa_2,\ldots$ and corresponding eigenfunctions $\psi_1,\psi_2,\ldots$. Thus\origNote{110}
\[
R\left(\sum_{m=1}^\infty x_m\psi_m\right)
=\sum_{m=1}^\infty x_m\varkappa_m\psi_m.
\]
Let $F(\varkappa)$ be a function defined for $-\infty<\varkappa<+\infty$, with $F(\varkappa_m)=\lambda_m$, $m=1,2,\ldots$; at all other points $\varkappa$ it may be arbitrary. Similarly, let $G(\varkappa_m)=\mu_m$, $H(\varkappa_m)=\nu_m$, and so on. We wish to show that
\[
A=F(R),\qquad B=G(R),\qquad C=H(R),\qquad\ldots.
\]
We must show that if $R$ has pure point spectrum $\varkappa_1,\varkappa_2,\ldots$ with eigenfunctions $\psi_1,\psi_2,\ldots$, then $F(R)$ has pure point spectrum $F(\varkappa_1),F(\varkappa_2),\ldots$ with the same eigenfunctions $\psi_1,\psi_2,\ldots$. Since these already form a complete orthonormal system, it suffices to prove $F(R)\psi_m=F(\varkappa_m)\cdot\psi_m$.

As in \hyperref[sec:II.8]{II.8}, let $E(\lambda)=\sum_{\varkappa_m\leq\lambda}P_{[\psi_m]}$ be the resolution of the identity belonging to $R$. Then, as we know, symbolically
\[
R=\int\lambda\,dE(\lambda),
\]
and by definition
\[
F(R)=\int F(\lambda)\,dE(\lambda).
\]
Furthermore,
\[
E(\lambda)\psi_m=\begin{cases}\psi_m,&\varkappa_m\leq\lambda,\\0,&\varkappa_m>\lambda.\end{cases}
\]
It follows that
\[
(F(R)\psi_m,g)=\int F(\lambda)\,d(E(\lambda)\psi_m,g)
=F(\varkappa_m)\cdot(\psi_m,g)
\]
for every $g$, and hence indeed $F(R)\psi_m=F(\varkappa_m)\cdot\psi_m$.

This disposes of the case of pure point spectra, as promised. For continuous spectra we must content ourselves with the reference in \hyperlink{note-109}{note 109}; let us merely single out one particularly characteristic case.

Let $\mathfrak R_\infty$ be the space of all $f(q_1q_2)$ with finite $\iint|f(q_1q_2)|^2\,dq_1\,dq_2$, the range of $q_1,q_2$ being the square $0\leq q_1,q_2\leq1$. Form the operators $A=q_1\cdots$, $B=q_2\cdots$. They are Hermitian, continuous for this range of $q_1,q_2$, though not for $-\infty<q_1,q_2<+\infty$, and they commute. Hence both must be functions of an $R$.

% Source PDF page 100 / original page 93.
\hypertarget{origpage-93}{}%
This $R$ therefore commutes with $A,B$, from which it follows, though we shall not explain the details here, that $R$ has the form $s(q_1,q_2)\cdots$, with $s(q_1,q_2)$ a bounded function. Thus $R^n$, $n=0,1,2,\ldots$, equals $(s(q_1,q_2))^n\cdots$, and $F(R)$ equals $F(s_1(q_1,q_2))\cdots$ if $F(\varkappa)$ is a polynomial. This formula can be extended to all $F(\varkappa)$, though again we cannot enter into the details. Consequently $F(R)=A$, $G(R)=B$ imply\origNote{111}
\[
F(s(q_1,q_2))=q_1,\qquad G(s(q_1,q_2))=q_2.
\]
The mutually inverse mappings $s(q_1,q_2)=\varkappa$ and $F(\varkappa)=q_1$, $G(\varkappa)=q_2$ would therefore have to map the square $0\leq q_1,q_2\leq1$ one-to-one onto a one-dimensional set of numbers $\varkappa$, contrary to geometric intuition.

Our cited proof nevertheless tells us that this must be possible; and in fact a mapping of the desired kind is supplied by the so-called Peano curve.\origNote{112} A closer examination of the proof cited in \hyperlink{note-109}{note 109} does indeed show that, in the present case, it leads to the Peano curve or to related constructions.

\section{The Trace}\label{sec:II.11}
We shall here define some important invariants of operators.

For a matrix $\{a_{\mu\nu}\}$ in $\mathfrak R_n$, the trace $\sum_{\mu=1}^n a_{\mu\mu}$ is such an invariant. It is unitarily invariant: it does not change when $\{a_{\mu\nu}\}$ is transformed to another Cartesian coordinate system.\origNote{113} If, however, we replace the matrix $\{a_{\mu\nu}\}$ by the corresponding operator
\[
A\{x_1,\ldots,x_n\}=\{y_1,\ldots,y_n\},\qquad
y_\mu=\sum_{\nu=1}^n a_{\mu\nu}x_\nu,
\]
then the $a_{\mu\nu}$ can be expressed in terms of $A$ as follows. The vectors $\varphi_1=\{1,0,\ldots,0\}$, $\varphi_2=\{0,1,\ldots,0\}$, \ldots, $\varphi_n=\{0,0,\ldots,1\}$ form a complete orthonormal system, and evidently $a_{\mu\nu}=(A\varphi_\nu,\varphi_\mu)$; cf.\ \hyperref[sec:II.5]{II.5}, especially \hyperlink{note-60}{note 60}. The trace is thus $\sum_{\mu=1}^n(A\varphi_\mu,\varphi_\mu)$, and its unitary invariance means that its value is the same for every complete orthonormal system.

We can immediately form the analogous concept in $\mathfrak R_\infty$. Let $A$ be a linear operator. Choose any complete orthonormal system $\varphi_1,\varphi_2,\ldots$ for which all the $A\varphi_\mu$ are meaningful. This is certainly possible if the domain of $A$ is everywhere dense: it suffices to orthogonalize a sequence $f_1,f_2,\ldots$ dense in that domain, by \hyperref[thm:II.8]{Theorem 8} in \hyperref[sec:II.2]{II.2}. Put $\operatorname{Spur}(A)=\sum_{\mu=1}^\infty(A\varphi_\mu,\varphi_\mu)$. We have to show that this really depends only on $A$, not on the $\varphi_\mu$.

% Source PDF page 101 / original page 94.
\hypertarget{origpage-94}{}%
For this purpose we first introduce two complete orthonormal systems $\varphi_1,\varphi_2,\ldots$ and $\psi_1,\psi_2,\ldots$, and put
\[
\operatorname{Spur}(A;\varphi,\psi)
=\sum_{\mu,\nu=1}^\infty(A\varphi_\mu,\psi_\nu)(\psi_\nu,\varphi_\mu).
\]
By \hyperref[thm:II.7]{Theorem 7}\,$\gamma$ in \hyperref[sec:II.2]{II.2}, this equals $\sum_{\mu=1}^\infty(A\varphi_\mu,\varphi_\mu)$, so its dependence on the $\psi_\nu$ is only apparent. Furthermore,
\begin{align*}
\sum_{\mu,\nu=1}^\infty(A\varphi_\mu,\psi_\nu)(\psi_\nu,\varphi_\mu)
&=\sum_{\mu,\nu=1}^\infty(\varphi_\mu,A^*\psi_\nu)(\psi_\nu,\varphi_\mu)\\
&=\overline{\sum_{\mu,\nu=1}^\infty(A^*\psi_\nu,\varphi_\mu)(\varphi_\mu,\psi_\nu)}.
\end{align*}
Thus $\operatorname{Spur}(A;\varphi,\psi)=\overline{\operatorname{Spur}(A^*;\psi,\varphi)}$. Since, as shown above, the right-hand side depends only apparently on the $\varphi_\mu$, the same holds for the left-hand side. Its dependence on both the $\varphi_\mu$ and the $\psi_\nu$ is therefore only apparent; in reality it depends solely on $A$. We may accordingly denote $\operatorname{Spur}(A;\varphi,\psi)$ by $\operatorname{Spur}(A)$. Since it equals $\sum_{\mu=1}^\infty(A\varphi_\mu,\varphi_\mu)$, the required invariance is proved. The last equation also gives $\operatorname{Spur}(A)=\overline{\operatorname{Spur}(A^*)}$.

The relations
\[
\operatorname{Spur}(aA)=a\operatorname{Spur}(A),\qquad
\operatorname{Spur}(A\pm B)=\operatorname{Spur}(A)\pm\operatorname{Spur}(B)
\]
are evident. Furthermore, even for noncommuting $A,B$,
\[
\operatorname{Spur}(AB)=\operatorname{Spur}(BA).
\]
This is shown as follows:
\begin{align*}
\operatorname{Spur}(AB)
&=\sum_{\mu=1}^\infty(AB\varphi_\mu,\varphi_\mu)
=\sum_{\mu=1}^\infty(B\varphi_\mu,A^*\varphi_\mu)\\
&=\sum_{\mu,\nu=1}^\infty(B\varphi_\mu,\psi_\nu)(\psi_\nu,A^*\varphi_\mu)\\
&=\sum_{\mu,\nu=1}^\infty(B\varphi_\mu,\psi_\nu)(A\psi_\nu,\varphi_\mu).
\end{align*}
Here $\varphi_1,\varphi_2,\ldots$ and $\psi_1,\psi_2,\ldots$ may be any two complete orthonormal systems. The symmetry of this expression in $A,B$, with simultaneous interchange of $\varphi,\psi$, is evident. Thus for Hermitian operators $A,B$,
\begin{align*}
\operatorname{Spur}(AB)
&=\overline{\operatorname{Spur}((AB)^*)}
=\overline{\operatorname{Spur}(B^*A^*)}\\
&=\overline{\operatorname{Spur}(BA)}
=\overline{\operatorname{Spur}(AB)}.
\end{align*}
Hence $\operatorname{Spur}(AB)$ is real; of course this holds all the more for $\operatorname{Spur}(A)$.

If $\mathfrak M$ is a closed linear manifold and $E$ its projection operator, $\operatorname{Spur}(E)$ is determined as follows. Let $\psi_1,\ldots,\psi_k$ be an orthonormal system spanning $\mathfrak M$, and $\chi_1,\ldots,\chi_l$ one spanning $\mathfrak R_\infty-\mathfrak M$. Of course $k$ or $l$, or both, are infinite.

% Source PDF page 102 / original page 95.
\hypertarget{origpage-95}{}%
Then $\psi_1,\ldots,\psi_k,\chi_1,\ldots,\chi_l$ together span $\mathfrak R_\infty$ and hence form a complete orthonormal system, by \hyperref[thm:II.7]{Theorem 7}\,$\alpha$ in \hyperref[sec:II.2]{II.2}. Thus
\begin{align*}
\operatorname{Spur}(E)
&=\sum_{\mu=1}^k(E\psi_\mu,\psi_\mu)+\sum_{\mu=1}^l(E\chi_\mu,\chi_\mu)\\
&=\sum_{\mu=1}^k(\psi_\mu,\psi_\mu)+\sum_{\mu=1}^l(0,\chi_\mu)
=\sum_{\mu=1}^k1=k.
\end{align*}
That is, $\operatorname{Spur}(E)$ is the number of dimensions of $\mathfrak M$.

If $A$ is definite, all $(A\varphi_\mu,\varphi_\mu)\geq0$, and hence $\operatorname{Spur}(A)\geq0$. If $\operatorname{Spur}(A)=0$, all the $(A\varphi_\mu,\varphi_\mu)$ must vanish, so $A\varphi_\mu=0$ by \hyperref[thm:II.19]{Theorem 19} in \hyperref[sec:II.5]{II.5}. If $\|\varphi\|=1$, we can find a complete orthonormal system $\varphi_1,\varphi_2,\ldots$ with $\varphi_1=\varphi$: take an everywhere dense sequence $f_1,f_2,\ldots$ and ``orthogonalize'' $\varphi,f_1,f_2,\ldots$, as in the proof of \hyperref[thm:II.7]{Theorem 7} in \hyperref[sec:II.2]{II.2}.\tnote{The source refers to Theorem 7 here; the orthogonalization procedure is proved in \hyperref[thm:II.8]{Theorem 8}. The source reference is retained.} This produces a complete orthonormal system beginning with $\varphi$. Thus $A\varphi=0$. For arbitrary $f$, $Af=0$ is clear if $f=0$; if $f\ne0$, it follows from the preceding argument with $\varphi=f/\|f\|$. Hence $A=0$. In other words, if $A$ is definite and nonzero, then $\operatorname{Spur}(A)>0$.

For all their brevity and simplicity, our considerations concerning the trace are not mathematically unassailable. We have considered the series
\[
\sum_{\mu,\nu=1}^\infty(A\varphi_\mu,\psi_\nu)(\psi_\nu,\varphi_\mu)
\quad\text{and}\quad
\sum_{\mu=1}^\infty(A\varphi_\mu,\varphi_\mu)
\]
without regard to their convergence, and transformed them into one another by rearranging their sums---in short, we have done everything that one ought properly not to do. Admittedly, such carelessness occurs elsewhere in theoretical physics, and the present instance will cause no harm in our quantum-mechanical applications; nevertheless, we must acknowledge that it is carelessness. It is all the more essential to emphasize that in the fundamental statistical assertions of quantum mechanics the trace is used only for operators $AB$ with both $A,B$ definite, and that this concept can also be established with complete precision. In the remainder of this section we shall therefore collect those facts concerning the trace that can be proved with absolute mathematical rigor.

We first consider the trace of $A^*A$, where $A$ is arbitrary. By \hyperref[sec:II.4]{II.4}, $A^*A$ is Hermitian, and it is definite since $(A^*Af,f)=(Af,Af)\geq0$. We have
\[
\operatorname{Spur}(A^*A)
=\sum_{\mu=1}^\infty(A^*A\varphi_\mu,\varphi_\mu)
=\sum_{\mu=1}^\infty(A\varphi_\mu,A\varphi_\mu)
=\sum_{\mu=1}^\infty\|A\varphi_\mu\|^2.
\]
Since every term of this series is nonnegative, it either converges or diverges to $+\infty$, and in either case is meaningful. Independently of the preceding discussion, we shall now show that its sum is independent of the choice of $\varphi_1,\varphi_2,\ldots$.


% Consolidated source: chapter-II-s11.tex
% Continuation only: II.11 begins on original page 93 and its section heading
% and original pages 93--95 are supplied by the Chapter II main translator.
% The sentence crossing original pages 95--96 is completed in the preceding file.
% Source PDF page 103 / original page 96
\hypertarget{origpage-96}{}%
Only series with terms $\geq0$ will occur, so everything will be meaningful and every rearrangement of summation permissible.

Let $\varphi_1,\varphi_2,\ldots$ and $\psi_1,\psi_2,\ldots$ be two complete orthonormal systems. We define
\[
\Sigma(A;\varphi_\mu,\psi_\nu)
=\sum_{\mu,\nu=1}^{\infty}|(A\varphi_\mu,\psi_\nu)|^2.
\]
By \hyperref[thm:II.7]{Theorem~7}, $\gamma$, in \hyperref[sec:II.2]{II.2}, this equals $\sum_{\mu=1}^{\infty}\|A\varphi_\mu\|^2$; that is, $\Sigma(A;\varphi_\mu,\psi_\nu)$ depends only apparently on the $\psi_\nu$. Further, provided that the $A\varphi_\mu$ and the $A^*\psi_\nu$ are meaningful,
\begin{align*}
\Sigma(A;\varphi_\mu,\psi_\nu)
&=\sum_{\mu,\nu=1}^{\infty}|(A\varphi_\mu,\psi_\nu)|^2
=\sum_{\mu,\nu=1}^{\infty}|(\varphi_\mu,A^*\psi_\nu)|^2\\
&=\sum_{\mu,\nu=1}^{\infty}|(A^*\psi_\nu,\varphi_\mu)|^2
=\Sigma(A^*;\psi_\nu,\varphi_\mu).
\end{align*}
Its dependence on the $\varphi_\mu$ is likewise only apparent: by the same argument, the expression on the right is independent of them. Consequently $\Sigma(A;\varphi_\mu,\psi_\nu)$ depends on $A$ alone, and we call it $\Sigma(A)$. What we have proved gives
\[
\Sigma(A)=\sum_{\mu=1}^{\infty}\|A\varphi_\mu\|^2
=\sum_{\mu,\nu=1}^{\infty}|(A\varphi_\mu,\psi_\nu)|^2
\]
and $\Sigma(A)=\Sigma(A^*)$. This provides a correct new definition of $\operatorname{Spur}(A^*A)$ as $\Sigma(A)$.

We shall also prove independently some properties of $\Sigma(A)$ that follow as well from the general properties of $\operatorname{Spur}(A)$ already derived.

The definition gives $\Sigma(A)\geq0$ in general; and if $\Sigma(A)=0$, all the $A\varphi_\mu$ must be zero, whence $A=0$ follows just as before. Thus, for $A\ne0$, $\Sigma(A)>0$.

Evidently $\Sigma(aA)=|a|^2\Sigma(A)$. If $A^*B=0$, then
\begin{align*}
&\|(A+B)\varphi_\mu\|^2-\|A\varphi_\mu\|^2-\|B\varphi_\mu\|^2\\
&\qquad=(A\varphi_\mu,B\varphi_\mu)+(B\varphi_\mu,A\varphi_\mu)\\
&\qquad=2\operatorname{Re}(A\varphi_\mu,B\varphi_\mu)
=2\operatorname{Re}(\varphi_\mu,A^*B\varphi_\mu)=0,
\end{align*}
so, on summing over $\mu=1,2,\ldots$,
\[
\Sigma(A+B)=\Sigma(A)+\Sigma(B).
\]
This relation is unchanged when $A,B$ are interchanged; it therefore holds for $B^*A=0$ as well. We may also replace $A,B$ in it by $A^*,B^*$: thus $AB^*=0$ or $BA^*=0$ is likewise sufficient. If $A$ (or $B$) is Hermitian, we may therefore write $AB=0$ or $BA=0$.

If $E$ is the projection operator of the closed linear manifold $\mathfrak M$, then for the $\psi_1,\ldots,\psi_k,\chi_1,\ldots,\chi_l$ considered in determining $\operatorname{Spur}(E)$,
% Source PDF page 104 / original page 97
\hypertarget{origpage-97}{}%
\begin{align*}
\Sigma(E)&=\sum_{\mu=1}^{k}\|E\psi_\mu\|^2
+\sum_{\mu=1}^{l}\|E\chi_\mu\|^2\\
&=\sum_{\mu=1}^{k}\|\psi_\mu\|^2
+\sum_{\mu=1}^{l}\|0\|^2
=\sum_{\mu=1}^{k}1=k.
\end{align*}
Thus $\Sigma(E)$ too is the dimension of $\mathfrak M$; since $E^*E=EE=E$, this was only to be expected.

For two definite (Hermitian) operators $A,B$, $\operatorname{Spur}(AB)$ can now be reduced to $\Sigma$. For there exist two operators $A',B'$ of the same kind with $A'^2=A$, $B'^2=B$;\origNote{114} we call them $\sqrt A,\sqrt B$. Purely formally,
\begin{align*}
\operatorname{Spur}(AB)
&=\operatorname{Spur}(\sqrt A\sqrt A\sqrt B\cdot\sqrt B)\\
&=\operatorname{Spur}(\sqrt B\cdot\sqrt A\sqrt A\sqrt B)\\
&=\operatorname{Spur}\bigl((\sqrt A\sqrt B)^*(\sqrt A\sqrt B)\bigr)
=\Sigma(\sqrt A\sqrt B).
\end{align*}
On the basis of its own definition, and without reference to its connection with the trace, this $\Sigma(\sqrt A\sqrt B)$ has all the properties we expect of $\operatorname{Spur}(AB)$, namely:
\begin{align*}
\Sigma(\sqrt A\sqrt B)&=\Sigma(\sqrt B\sqrt A),\\
\Sigma(\sqrt A\sqrt{B+C})&=\Sigma(\sqrt A\sqrt B)+\Sigma(\sqrt A\sqrt C),\\
\Sigma(\sqrt{A+B}\sqrt C)&=\Sigma(\sqrt A\sqrt C)+\Sigma(\sqrt B\sqrt C).
\end{align*}
The first follows from the symmetry of $\Sigma(XY)$ in $X,Y$:
\[
\Sigma(XY)=\sum_{\mu,\nu=1}^{\infty}|(XY\varphi_\mu,\psi_\nu)|^2
=\sum_{\mu,\nu=1}^{\infty}|(Y\varphi_\mu,X\psi_\nu)|^2;
\]
by the first, the second follows from the third. Thus only the third remains to be proved: that is, that $\Sigma(\sqrt A\sqrt B)$ is additive in $A$. This is seen by writing $\Sigma(\sqrt A\sqrt B)$ as follows:
\begin{align*}
\Sigma(\sqrt A\sqrt B)
&=\sum_{\mu=1}^{\infty}\|\sqrt A\sqrt B\varphi_\mu\|^2\\
&=\sum_{\mu=1}^{\infty}(\sqrt A\sqrt B\varphi_\mu,\sqrt A\sqrt B\varphi_\mu)\\
&=\sum_{\mu=1}^{\infty}(\sqrt A\cdot\sqrt A\sqrt B\varphi_\mu,\sqrt B\varphi_\mu)\\
&=\sum_{\mu=1}^{\infty}(A\sqrt B\varphi_\mu,\sqrt B\varphi_\mu).
\end{align*}
This establishes the concept of trace rigorously to the extent sought above.

The last formula also permits the following conclusion: if $A,B$ are definite, then $\operatorname{Spur}(AB)=0$ implies $AB=0$. For the former says $\Sigma(\sqrt A\sqrt B)=0$, so $\sqrt A\sqrt B=0$ (cf.\ what was said on \hyperlink{origpage-96}{p.~96}, or the further discussion of $\Sigma$ below), and hence $AB=\sqrt A\cdot\sqrt A\sqrt B\cdot\sqrt B=0$.

For a definite Hermitian operator $A$, calculation with the trace is correct even in its original form. For let $\varphi_1,\varphi_2,\ldots$ be a complete orthonormal system. Then
% Source PDF page 105 / original page 98
\hypertarget{origpage-98}{}%
$\sum_{\mu=1}^{\infty}(A\varphi_\mu,\varphi_\mu)$---the sum that was to define the trace!---has only terms $\geq0$, and therefore either converges or diverges to $+\infty$. Two cases are now possible. Either this sum is infinite for every choice of $\varphi_1,\varphi_2,\ldots$; in that case the trace is indeed defined independently of $\varphi_1,\varphi_2,\ldots$, and equals $+\infty$. Or it is finite for at least one choice of $\varphi_1,\varphi_2,\ldots$, say $\bar\varphi_1,\bar\varphi_2,\ldots$. Then, because
\begin{align*}
\left(\sum_{\mu=1}^{\infty}(A\bar\varphi_\mu,\bar\varphi_\mu)\right)^2
&=\sum_{\mu,\nu=1}^{\infty}(A\bar\varphi_\mu,\bar\varphi_\mu)
(A\bar\varphi_\nu,\bar\varphi_\nu)\\
&\geq\sum_{\mu,\nu=1}^{\infty}|(A\bar\varphi_\mu,\bar\varphi_\nu)|^2
=\Sigma(A),
\end{align*}
$\Sigma(A)$ is finite as well, say equal to $C^2$. If $\psi_1,\psi_2,\ldots$ is any complete orthonormal system, then
\[
\Sigma(A)=\sum_{\mu=1}^{\infty}\|A\psi_\mu\|^2=C^2,
\qquad\|A\psi_1\|^2\leq C^2,
\qquad\|A\psi_1\|\leq C.
\]
Since every $\psi$ with $\|\psi\|=1$ can be chosen as $\psi_1$ of such a system, $\|\psi\|=1$ implies $\|A\psi\|\leq C$. Thus $\|Af\|\leq C\|f\|$ holds generally: for $f=0$ this is evident, and for $f\ne0$ it is enough to put $\psi=\frac1{\|f\|}\cdot f$. Hence $A$ satisfies condition \hyperref[eq:II.9.St]{St.} of \hyperref[sec:II.9]{II.9}; it is a continuous operator. But considerably more holds.

Because $\Sigma(A)$ is finite, $A$ belongs to the class of so-called completely continuous operators, for which Hilbert showed that the eigenvalue problem is soluble in its original form: that is, there exists a complete orthonormal system $\psi_1,\psi_2,\ldots$ with $A\psi_\mu=\lambda_\mu\psi_\mu$ (and $\lambda_\mu\to0$ as $\mu\to\infty$).\origNote{115} Definiteness gives $\lambda_\mu=(A\psi_\mu,\psi_\mu)\geq0$; furthermore,
\[
\sum_{\mu=1}^{\infty}\lambda_\mu^2
=\sum_{\mu=1}^{\infty}\|A\psi_\mu\|^2=\Sigma(A)=C^2.
\]
If $\varphi_1,\varphi_2,\ldots$ is any complete orthonormal system, then
\begin{align*}
\sum_{\mu=1}^{\infty}(A\varphi_\mu,\varphi_\mu)
&=\sum_{\mu=1}^{\infty}\left(\sum_{\nu=1}^{\infty}
(A\varphi_\mu,\psi_\nu)(\psi_\nu,\varphi_\mu)\right)\\
&=\sum_{\mu=1}^{\infty}\left(\sum_{\nu=1}^{\infty}
(\varphi_\mu,A\psi_\nu)(\psi_\nu,\varphi_\mu)\right)\\
&=\sum_{\mu=1}^{\infty}\left(\sum_{\nu=1}^{\infty}
\lambda_\nu(\varphi_\mu,\psi_\nu)(\psi_\nu,\varphi_\mu)\right)\\
&=\sum_{\mu=1}^{\infty}\left(\sum_{\nu=1}^{\infty}
\lambda_\mu|(\varphi_\mu,\psi_\nu)|^2\right).
\end{align*}
Since everything is $\geq0$, we may rearrange the summation.\tnote{The source prints $\lambda_\mu$ in the final inner sum of the preceding display, although the preceding line and the rearranged sum below require $\lambda_\nu$. The printed subscript is retained.} Thus:
\begin{align*}
\sum_{\mu=1}^{\infty}(A\varphi_\mu,\varphi_\mu)
&=\sum_{\mu,\nu=1}^{\infty}\lambda_\nu|(\varphi_\mu,\psi_\nu)|^2\\
&=\sum_{\nu=1}^{\infty}\lambda_\nu
\left(\sum_{\mu=1}^{\infty}|(\varphi_\mu,\psi_\nu)|^2\right)\\
&=\sum_{\nu=1}^{\infty}\lambda_\nu\|\psi_\nu\|^2
=\sum_{\nu=1}^{\infty}\lambda_\nu.
\end{align*}

% Source PDF page 106 / original page 99
\hypertarget{origpage-99}{}%
In this case, therefore, $\sum_{\mu=1}^{\infty}(A\varphi_\mu,\varphi_\mu)$ is again independent of $\varphi_1,\varphi_2,\ldots$, and equals the sum of the eigenvalues. Since it is finite for $\bar\varphi_1,\bar\varphi_2,\ldots$, it is therefore always finite. Thus $\operatorname{Spur}(A)$ is again uniquely determined, but this time finite.

Calculation with the trace is therefore justified in both cases.

We give a few more estimates concerning $\operatorname{Spur}(A)$ and $\Sigma(A)$. For every $A$ with finite $\Sigma(A)$, we had $\|Af\|\leq\sqrt{\Sigma(A)}\cdot\|f\|$; for every definite (Hermitian) $A$ with finite $\operatorname{Spur}(A)$,
\[
\|Af\|\leq\operatorname{Spur}(A)\cdot\|f\|.
\]
Now let $A$ be definite, $\operatorname{Spur}(A)=1$, and suppose that, for a suitable $\varphi$ with $\|\varphi\|=1$, $\|A\varphi\|^2\geq1-\varepsilon$ or $(A\varphi,\varphi)\geq1-\varepsilon$. Since $(A\varphi,\varphi)\leq\|A\varphi\|\cdot\|\varphi\|=\|A\varphi\|$, the second implies the first [with $(1-\varepsilon)^2\geq1-2\varepsilon$ in place of $1-\varepsilon$, hence $2\varepsilon$ in place of $\varepsilon$], so it is enough to consider the first.

Let $\psi$ be orthogonal to $\varphi$, with $\|\psi\|=1$. We can then find a complete orthonormal system $\chi_1,\chi_2,\ldots$ with $\chi_1=\varphi$, $\chi_2=\psi$, so
\begin{align*}
\sum_{\mu=1}^{\infty}\|A\chi_\mu\|^2
&=\Sigma(A)\leq(\operatorname{Spur}(A))^2=1,\\
\sum_{\mu=1}^{\infty}\|A\chi_\mu\|^2
&\geq\|A\varphi\|^2+\|A\psi\|^2
\geq1-2\varepsilon+\|A\psi\|^2,\\
\|A\psi\|^2&\leq2\varepsilon,\qquad\|A\psi\|\leq\sqrt{2\varepsilon}.
\end{align*}
For any $f$ orthogonal to $\varphi$, it follows that $\|Af\|\leq\sqrt{2\varepsilon}\cdot\|f\|$ (clear for $f=0$; otherwise put $\psi=\frac1{\|f\|}\cdot f$). Taking account also of $(Af,g)=(f,Ag)$, we may say that $|(Af,g)|\leq\sqrt{2\varepsilon}\cdot\|f\|\cdot\|g\|$ if either $f$ or $g$ is orthogonal to $\varphi$.

Now let $f,g$ be arbitrary. Then
\[
f=\alpha\varphi+f',\qquad g=\beta\varphi+g',
\]
where $f',g'$ are orthogonal to $\varphi$, and $\alpha=(f,\varphi)$, $\beta=(g,\varphi)$. Consequently
\[
(Af,g)=\alpha\bar\beta(A\varphi,\varphi)
+\alpha(A\varphi,g')+\bar\beta(Af',\varphi)+(Af',g'),
\]
so, on putting $(A\varphi,\varphi)=c$,
\[
|(Af,g)-c\alpha\bar\beta|
\leq|\alpha|\cdot|(A\varphi,g')|
+|\beta|\cdot|(Af',\varphi)|+|(Af',g')|,
\]
and, by the estimates above,
\begin{align*}
|(Af,g)-c\alpha\bar\beta|
&\leq\sqrt{2\varepsilon}\cdot
\bigl(|\alpha|\cdot\|g'\|+|\beta|\cdot\|f'\|+\|f'\|\cdot\|g'\|\bigr)\\
&\leq\sqrt{2\varepsilon}\cdot(|\alpha|+\|f'\|)(|\beta|+\|g'\|)\\
&\leq2\sqrt{2\varepsilon}\cdot
\sqrt{|\alpha|^2+\|f'\|^2}\sqrt{|\beta|^2+\|g'\|^2}\\
&=2\sqrt{2\varepsilon}\cdot\|f\|\cdot\|g\|.
\end{align*}
On the other hand,
\begin{align*}
(Af,g)-c\alpha\bar\beta
&=(Af,g)-c(f,\varphi)(\varphi,g)\\
&=((A-cP_{[\varphi]})f,g).
\end{align*}
Thus $|((A-cP_{[\varphi]})f,g)|\leq2\sqrt{2\varepsilon}\cdot\|f\|\cdot\|g\|$ holds generally, and hence, as we know from \hyperref[sec:II.9]{II.9}, so does
% Source PDF page 107 / original page 100
\hypertarget{origpage-100}{}%
\[
\|(A-cP_{[\varphi]})f\|\leq2\sqrt{2\varepsilon}\cdot\|f\|.
\]
For $f=\varphi$ this says (where $c=(A\varphi,\varphi)$ is real and $\geq0$):
\[
\|A\varphi-c\varphi\|\leq2\sqrt{2\varepsilon},
\]
\begin{align*}
c=\|c\varphi\|&\leq\|A\varphi-c\varphi\|+\|A\varphi\|
\leq2\sqrt{2\varepsilon}+1,\\
c=\|c\varphi\|&\geq-\|A\varphi-c\varphi\|+\|A\varphi\|
\geq-2\sqrt{2\varepsilon}+(1-\varepsilon),
\end{align*}
\[
1-(\varepsilon+2\sqrt{2\varepsilon})\leq c\leq1+2\sqrt{2\varepsilon}.
\]
Consequently
\begin{align*}
\|(A-P_{[\varphi]})f\|
&\leq\|(A-cP_{[\varphi]})f\|+\|(c-1)P_{[\varphi]}f\|\\
&\leq2\sqrt{2\varepsilon}\cdot\|f\|
+(\varepsilon+2\sqrt{2\varepsilon})\|P_{[\varphi]}f\|\\
&\leq(\varepsilon+4\sqrt{2\varepsilon})\cdot\|f\|.
\end{align*}
Thus $A$ converges uniformly to $P_{[\varphi]}$ as $\varepsilon\to0$.

Finally, we consider $\operatorname{Spur}(A)$ and $\Sigma(A)$ in the realizations $F_Z$ and $F_\Omega$ of $\mathfrak R_\infty$ (cf.\ \hyperref[sec:I.4]{I.4} and \hyperref[sec:II.3]{II.3}), since these are the realizations we shall use in physical applications.

In $F_Z$, the set of all $\{x_1,x_2,\ldots\}$ with finite $\sum_{\mu=1}^{\infty}|x_\mu|^2$, $A$ can be described by a matrix $\{a_{\mu\nu}\}$:
\[
A\{x_1,x_2,\ldots\}=\{y_1,y_2,\ldots\},
\qquad y_\mu=\sum_{\nu=1}^{\infty}a_{\mu\nu}x_\nu.
\]
The sequences $\{1,0,0,\ldots\}$, $\{0,1,0,\ldots\}$, $\ldots$ form a complete orthonormal system $\varphi_1,\varphi_2,\ldots$, and
\[
A\varphi_\mu=\{a_{1\mu},a_{2\mu},\ldots\}
=\sum_{\varrho=1}^{\infty}a_{\varrho\mu}\varphi_\varrho,
\]
so $(A\varphi_\mu,\varphi_\mu)=a_{\mu\mu}$ and $\|A\varphi_\mu\|^2=\sum_{\varrho=1}^{\infty}|a_{\varrho\mu}|^2$. This immediately gives
\[
\operatorname{Spur}(A)=\sum_{\mu=1}^{\infty}a_{\mu\mu},
\qquad\Sigma(A)=\sum_{\mu,\nu=1}^{\infty}|a_{\mu\nu}|^2.
\]
In $F_\Omega$, the set of all $f(P)$ defined in $\Omega$ with finite $\int_\Omega|f(P)|^2\,dv$, we consider only the integral operators
\[
Af(P)=\int_\Omega a(P,P')f(P')\,dv'
\]
($a(P,P')$ being a function of two variables defined in $\Omega$, the ``integral kernel'' of $A$; cf.\ \hyperref[sec:I.4]{I.4}). Let $\varphi_1(P),\varphi_2(P),\ldots$ be any complete orthonormal system. Then
\begin{align*}
\operatorname{Spur}(A)
&=\sum_{\mu=1}^{\infty}(A\varphi_\mu(P),\varphi_\mu(P))\\
&=\sum_{\mu=1}^{\infty}\int_\Omega
\left[\int_\Omega a(P,P')\varphi_\mu(P')\,dv'\right]
\overline{\varphi_\mu(P)}\,dv,
\end{align*}
% Source PDF page 108 / original page 101
\hypertarget{origpage-101}{}%
and because in general [\hyperref[thm:II.7]{Theorem~7}, $\beta$, in \hyperref[sec:II.2]{II.2}, applied to $\overline{g(P)}$]
\[
\sum_{\mu=1}^{\infty}\left(\int_\Omega\overline{g(P')}\,
\overline{\varphi_\mu(P')}\,dv'\right)\varphi_\mu(P)
=\overline{g(P)},
\]
\[
\sum_{\mu=1}^{\infty}\left(\int_\Omega g(P')\varphi_\mu(P')\,dv'\right)
\overline{\varphi_\mu(P)}=g(P),
\]
we have
\[
\operatorname{Spur}(A)=\int_\Omega a(P,P)\,dv.
\]
Furthermore,
\[
\Sigma(A)=\sum_{\mu=1}^{\infty}\int_\Omega
\left|\int_\Omega a(P,P')\varphi_\mu(P')\,dv'\right|^2\,dv,
\]
and hence, since (\hyperref[thm:II.7]{Theorem~7}, $\gamma$, in \hyperref[sec:II.2]{II.2})
\begin{align*}
\sum_{\mu=1}^{\infty}\left|\int_\Omega g(P')\varphi_\mu(P')\,dv'\right|^2
&=\sum_{\mu=1}^{\infty}\left|\int_\Omega
\overline{g(P')}\,\overline{\varphi_\mu(P')}\,dv'\right|^2\\
&=\int_\Omega|\overline{g(P')}|^2\,dv'
=\int_\Omega|g(P')|^2\,dv',
\end{align*}
we obtain
\[
\Sigma(A)=\int_\Omega\int_\Omega|a(P,P')|^2\,dv\,dv'.
\]

We see that $\operatorname{Spur}(A)$ and $\Sigma(A)$ accomplish what was attempted in \hyperref[sec:I.4]{I.4} by mathematically rather questionable devices: in passing from $F_Z$ to $F_\Omega$, $\sum_{\mu=1}^{\infty}\cdots$ is replaced by $\int_\Omega\cdots\,dv$.

This concludes our mathematical discussion of Hermitian operators. Readers with a mathematical interest will find further information on the subject in the relevant literature.\origNote{116}


% Consolidated source: chapter-III.tex
% Source PDF page 108 / original page 101; original-page target belongs to the end of Chapter II.
\chapter{Quantum-Mechanical Statistics}\label{chap:III}
\section{The Statistical Assertions of Quantum Mechanics}\label{sec:III.1}
Let us return to the analysis of quantum-mechanical theories that we interrupted for the mathematical discussion of \hyperref[chap:II]{Part~II}. We had considered only how quantum mechanics makes it possible to determine all possible values of one particular physical quantity, energy: these are the eigenvalues of the energy operator $\mathsf H$, that is, the numbers in its spectrum. We did not discuss what it can say about the values of other quantities, or about causal or statistical relations between the values of several quantities. We shall now examine the theory's assertions on these questions. Here we shall use the wave-mechanical mode of description, since the equivalence of the two theories has been established.

% Source PDF page 109 / original page 102
\hypertarget{origpage-102}{}%
In this description it is clear that everything we wish to say about the state of a system must be obtained from its wave function $\varphi(q_1,\ldots,q_k)$. (Let the system have $k$ degrees of freedom, and let $q_1,\ldots,q_k$ be the coordinates of its configuration space.) We do not restrict ourselves to the stationary states of the system---the quantum orbits, for which $\varphi$ is an eigenfunction of $\mathsf H$: $\mathsf H\varphi=\lambda\varphi$, cf.\ \hyperref[sec:I.3]{I.3}---but admit all its states, that is, wave functions $\varphi$ varying according to Schrödinger's time-dependent differential equation $\mathsf H\varphi=-\frac h{2\pi i}\frac\partial{\partial t}\varphi$; cf.\ \hyperref[sec:I.2]{I.2}. What assertions, then, can be made about a system in the state $\varphi$?

First we recall that $\varphi$ was normalized in \hyperref[sec:I.3]{I.3} by
\[
\int_{-\infty}^{\infty}\cdots\int_{-\infty}^{\infty}
|\varphi(q_1\cdots q_k)|^2\,dq_1\cdots dq_k=1.
\]
In our present terminology, we regard $\varphi$ as a point of the Hilbert space $\mathfrak R_\infty$ consisting of all $f(q_1\cdots q_k)$ for which
\[
\int_{-\infty}^{\infty}\cdots\int_{-\infty}^{\infty}
|f(q_1\cdots q_k)|^2\,dq_1\cdots dq_k
\]
is finite (a space of type $F_\Omega$). The normalization is then $\|\varphi\|=1$: $\varphi$ must lie on the surface of the unit sphere in Hilbert space.\origNote{117} We already know that a constant factor in $\varphi$, independent of $q_1,\ldots,q_k$, has no physical significance. (That is, replacing $\varphi$ by $a\varphi$, where $a$ is a complex number; because of the normalization $\|\varphi\|=1$, we must have $|a|=1$.) We should also point out that, besides the coordinates $q_1,\ldots,q_k$ of our system's configuration space, $\varphi$ depends on the time $t$. Hilbert space is formed only with respect to $q_1,\ldots,q_k$, however, since the normalization also referred to these alone. Dependence on $t$ is not to be included here; $t$ is to be regarded as a parameter. Consequently, as a point of $\mathfrak R_\infty$, $\varphi$ depends on $t$ but is independent of $q_1,\ldots,q_k$: as a point of $\mathfrak R_\infty$ it represents the entire functional dependence on $q_1,\ldots,q_k$. For this reason we shall sometimes indicate the parameter $t$ when treating $\varphi$ as a point of $\mathfrak R_\infty$, writing $\varphi_t$.

Consider now the state $\varphi=\varphi(q_1\cdots q_k)$. The statistical assertion that can then be made is this: the system is at the configuration-space point $q_1,\ldots,q_k$ with probability density $|\varphi(q_1\cdots q_k)|^2$. Thus the probability that it lies in a region $V$ of configuration space is
\[
\int\cdots\int_V|\varphi(q_1\cdots q_k)|^2\,dv.
\]
(This is one of the first and simplest examples in which the statistical character of quantum mechanics was recognized;\origNote{118}
% Source PDF page 110 / original page 103
\hypertarget{origpage-103}{}%
the connection between this assertion and Schrödinger's charge-distribution assumption, cf.\ \hyperref[sec:I.2]{I.2}, is also clear.) Furthermore, if the energy of the system has the operator $\mathsf H$, and this operator has eigenvalues $\lambda_1,\lambda_2,\ldots$ and eigenfunctions $\varphi_1,\varphi_2,\ldots$, the probability of the energy value $\lambda_n$ in the state $\varphi$ is
\[
\left|\int\cdots\int\varphi(q_1\cdots q_k)
\overline{\varphi_n(q_1\cdots q_k)}\,dq_1\cdots dq_k\right|^2
\]
(see the papers cited in \hyperlink{note-118}{Note~118}). We shall now recast these two assertions in a common form.

Let $V$ be the $k$-dimensional rectangular parallelepiped
\[
q'_1<q_1\le q''_1,\quad\ldots,\quad q'_k<q_k\le q''_k,
\]
and call the intervals $\{q'_1,q''_1\},\ldots,\{q'_k,q''_k\}$ respectively $I_1,\ldots,I_k$. The operators of $q_1,\ldots,q_k$ are respectively $q_1\cdots,\ldots,q_k\cdots$. Their resolutions of the identity are defined as follows (cf.\ \hyperref[sec:II.8]{II.8}): call the one belonging to $q_j$ ($j=1,\ldots,k$) $E_j(\lambda)$; then
\[
E_j(\lambda)f(q_1\cdots q_k)=
\begin{cases}f(q_1\cdots q_k),&q_j\le\lambda,\\0,&q_j>\lambda.\end{cases}
\]
We introduce the following general notation: if $F(\lambda)$ is a resolution of the identity and $I$ is an interval $\{\lambda',\lambda''\}$, write $F(I)=F(\lambda'')-F(\lambda')$. This is a projection operator, since $\lambda'\le\lambda''$ and $F(\lambda')\le F(\lambda'')$. The probability that the system lies in the above $V$, that is, that $q_1$ lies in $I_1$, $\ldots$, $q_k$ in $I_k$, is
\begin{align*}
&\int_{q'_1}^{q''_1}\cdots\int_{q'_k}^{q''_k}|\varphi(q_1\cdots q_k)|^2\,dq_1\cdots dq_k\\
&\qquad=\int\cdots\int|E_1(I_1)\cdots E_k(I_k)\varphi(q_1\cdots q_k)|^2\,dq_1\cdots dq_k,
\end{align*}
because $E_1(I_1)\cdots E_k(I_k)\varphi(q_1\cdots q_k)$ equals $\varphi(q_1\cdots q_k)$ for $q_1$ in $I_1$, $\ldots$, $q_k$ in $I_k$, and zero otherwise. Thus it is
\[
\|E_1(I_1)\cdots E_k(I_k)\varphi\|^2.
\]
In the second case we consider the probability that the energy lies in the interval $I=\{\lambda',\lambda''\}$. The resolution of the identity $E(\lambda)$ belonging to $\mathsf H$ is defined (cf.\ \hyperref[sec:II.8]{II.8}) by $E(\lambda)=\sum_{\lambda_n\le\lambda}P_{[\varphi_n]}$, so
\[
E(I)=E(\lambda'')-E(\lambda')=\sum_{\lambda'<\lambda_n\le\lambda''}P_{[\varphi_n]}.
\]
Since only $\lambda_1,\lambda_2,\ldots$ are possible energy values, the probability in question is the sum of the probabilities of all $\lambda_n$ with $\lambda'<\lambda_n\le\lambda''$, hence
\begin{align*}
&\sum_{\lambda'<\lambda_n\le\lambda''}
\left|\int\cdots\int\varphi(q_1\cdots q_k)\overline{\varphi_n(q_1\cdots q_k)}\,dq_1\cdots dq_k\right|^2\\
&\qquad=\sum_{\lambda'<\lambda_n\le\lambda''}|(\varphi,\varphi_n)|^2
=\sum_{\lambda'<\lambda_n\le\lambda''}(P_{[\varphi_n]}\varphi,\varphi)\\
&\qquad=\left(\left\{\sum_{\lambda'<\lambda_n\le\lambda''}P_{[\varphi_n]}\right\}\varphi,\varphi\right)
=(E(I)\varphi,\varphi)=\|E(I)\varphi\|^2.
\end{align*}
% Source PDF page 111 / original page 104
\hypertarget{origpage-104}{}%

In both cases, then, we have obtained a result that can be formulated as follows:

\leavevmode\phantomsection\label{post:III.1.W}\textbf{(W.)} The probability that, in the state $\varphi$, the quantities with respective operators $R_1,\ldots,R_l$\origNote{119} take values in the respective intervals $I_1,\ldots,I_l$ is
\[
\|E_1(I_1)\cdots E_l(I_l)\varphi\|^2,
\]
where $E_1(\lambda),\ldots,E_l(\lambda)$ are the resolutions of the identity belonging respectively to $R_1,\ldots,R_l$.

The first case corresponded to $l=k$, $R_1=q_1\cdots,\ldots,R_k=q_k\cdots$; the second to $l=1$, $R_1=\mathsf H$. We shall now assume this assertion \hyperref[post:III.1.W]{W} to be generally valid. It does indeed contain all the statistical assertions of quantum mechanics made so far.

Its validity must nevertheless be restricted in one respect. Since the order of $R_1,\ldots,R_l$ is entirely immaterial in the question, it must also be immaterial in the result. Thus the $E_1(I_1),\ldots,E_l(I_l)$, or equivalently all the $E_1(\lambda_1),\ldots,E_l(\lambda_l)$, must commute. According to \hyperref[sec:II.10]{II.10}, this means that $R_1,\ldots,R_l$ commute with one another. This condition is fulfilled for $q_1\cdots,\ldots,q_k\cdots$; for $l=1$, $R_1=\mathsf H$, it becomes irrelevant altogether.

We therefore postulate \hyperref[post:III.1.W]{W} generally for commuting $R_1,\ldots,R_l$. Then $E_1(I_1),\ldots,E_l(I_l)$ commute, so $E_1(I_1)\cdots E_l(I_l)$ is a projection operator (\hyperref[thm:II.14]{Theorem~14 in II.4}), and the probability is
\[
W=\|E_1(I_1)\cdots E_l(I_l)\varphi\|^2
=(E_1(I_1)\cdots E_l(I_l)\varphi,\varphi)
\]
(\hyperref[thm:II.12]{Theorem~12 in II.4}).

Before proceeding, we must verify certain properties of \hyperref[post:III.1.W]{W} that every meaningful statistical theory must possess.

1. The order of the assertions is immaterial.

2. Vacuous assertions may be added without changing $W$: these are assertions for which the interval $I_j$ is $\{-\infty,+\infty\}$, since they merely introduce a factor
\[
E_j(I)=E(+\infty)-E(-\infty)=1-0=1.
\]

3. The addition theorem of probabilities holds. That is, if an interval $I_j$ is divided into two intervals $I'_j,I''_j$, the old probability is the sum of the two new ones. For let $I_j,I'_j,I''_j$ be respectively $\{\lambda',\lambda''\}$, $\{\lambda',\lambda\}$, $\{\lambda,\lambda''\}$; then
\[
E(\lambda'')-E(\lambda')=(E(\lambda)-E(\lambda'))+(E(\lambda'')-E(\lambda)),
\]
or $E(I_j)=E(I'_j)+E(I''_j)$. In view of the second form of $W$ given above, which is linear in $E_1(I_1)\cdots E_j(I_j)\cdots E_l(I_l)$, this gives additivity of probabilities.

% Source PDF page 112 / original page 105
\hypertarget{origpage-105}{}%
4. For absurd assertions (one $I_j$ empty), $W=0$, since the corresponding $E_j(I_j)$ is zero. For self-evident assertions (all $I_j$ equal to $\{-\infty,+\infty\}$), $W=1$, since all $E_j(I_j)=1$ and $W=\|\varphi\|^2=1$. Always $0\le W\le1$, by \hyperref[thm:II.13]{Theorem~13 in II.4}.

Finally, \hyperref[post:III.1.W]{W} contains the assertion that a quantity $R_j$ can take only its eigenvalues, that is, the numbers in its spectrum, as values. For if the interval $I_j=\{\lambda',\lambda''\}$ lies outside the spectrum, $E_j(\lambda)$ is constant in it, so
\[
E_j(I_j)=E_j(\lambda'')-E_j(\lambda')=0,
\]
and $W=0$ follows.

Let us now put $l=1$ and denote $R_1$ by $R$. Call the physical quantity to which $R$ is assigned $\mathfrak R$ (cf.\ \hyperlink{note-119}{Note~119}). Let $F(\lambda)$ be any function; we wish to calculate the expectation value of $F(\mathfrak R)$.

For this purpose divide the interval $\{-\infty,+\infty\}$ into a sequence of subintervals $\{\lambda_n,\lambda_{n+1}\}$, $n=0,\pm1,\pm2,\ldots$. The probability that $\mathfrak R$ lies in $\{\lambda_n,\lambda_{n+1}\}$ is
\[
(\{E(\lambda_{n+1})-E(\lambda_n)\}\varphi,\varphi)
=(E(\lambda_{n+1})\varphi,\varphi)-(E(\lambda_n)\varphi,\varphi),
\]
and the expectation value of $F(\mathfrak R)$ is therefore
\[
\sum_{n=-\infty}^{+\infty}F(\lambda'_n)
\{(E(\lambda_{n+1})\varphi,\varphi)-(E(\lambda_n)\varphi,\varphi)\},
\]
where $\lambda'_n$ is a suitable intermediate value in $\{\lambda_n,\lambda_{n+1}\}$. As the division $\ldots,\lambda_{-2},\lambda_{-1},\lambda_0,\lambda_1,\lambda_2,\ldots$ is made finer and finer, these sums converge to the Stieltjes integral
\[
\int_{-\infty}^{\infty}F(\lambda)\,d(E(\lambda)\varphi,\varphi).
\]
The expectation value in question is therefore also equal to this integral. By the general definition of functions of operators in \hyperref[sec:II.8]{II.8}, however, the integral equals $(F(R)\varphi,\varphi)$. Hence:

\leavevmode\phantomsection\label{post:III.1.E1}\textbf{($E_1$.)} Let $\mathfrak R$ be any physical quantity, $R$ its operator (cf.\ \hyperlink{note-119}{Note~119}), and $F(\lambda)$ an arbitrary function. The expectation value of $f(\mathfrak R)$ in the state $\varphi$ is then given by
\[
\operatorname{Erw}(F(\mathfrak R);\varphi)=(F(R)\varphi,\varphi).
\]
In particular, putting $F(\lambda)=\lambda$ gives:

\leavevmode\phantomsection\label{post:III.1.E2}\textbf{($E_2$.)} Let $\mathfrak R,R$ be as above. The expectation value of $\mathfrak R$ in the state $\varphi$ is then given by
\[
\operatorname{Erw}(\mathfrak R;\varphi)=(R\varphi,\varphi).
\]
We shall next investigate the relations between \hyperref[post:III.1.W]{W}, \hyperref[post:III.1.E1]{$E_1$}, and \hyperref[post:III.1.E2]{$E_2$}.
% Source PDF page 113 / original page 106
\hypertarget{origpage-106}{}%

We deduced \hyperref[post:III.1.E1]{$E_1$} from \hyperref[post:III.1.W]{W}, and \hyperref[post:III.1.E2]{$E_2$} from \hyperref[post:III.1.E1]{$E_1$}.

If we denote the operator of $f(\mathfrak R)$ by $S$, comparison of \hyperref[post:III.1.E1]{$E_1$},\hyperref[post:III.1.E2]{$E_2$} gives
\[
(S\varphi,\varphi)=(F(R)\varphi,\varphi)
\]
for all states $\varphi$, that is, all $\varphi$ with $\|\varphi\|=1$. Thus, in general,
\[
(Sf,f)=(F(R)f,f)
\]
(clear for $f=0$; otherwise put $\varphi=f/\|f\|$), and hence also
\[
(Sf,g)=(F(R)f,g).
\]
For replacing $f$ by $(f+g)/2$ and $(f-g)/2$ and subtracting gives equality of the real parts; using $if,g$ instead of $f,g$ gives equality of the imaginary parts. Thus $S=F(R)$. We formulate this important result separately:

\leavevmode\phantomsection\label{post:III.1.F}\textbf{(F.)} If the quantity $\mathfrak R$ has the operator $R$, the quantity $F(\mathfrak R)$ must have the operator $F(R)$.

Assuming \hyperref[post:III.1.F]{F}, however, \hyperref[post:III.1.E2]{$E_2$} plainly follows from \hyperref[post:III.1.E1]{$E_1$}.\tnote{The source repeats the implication from $E_1$ to $E_2$ here. The equivalence asserted next requires the converse implication, from $E_2$ together with $F$ to $E_1$. The printed direction is retained.}

Thus, under the assumption \hyperref[post:III.1.F]{F}, \hyperref[post:III.1.E1]{$E_1$} and \hyperref[post:III.1.E2]{$E_2$} are equivalent assertions. We shall now show that they are also equivalent to \hyperref[post:III.1.W]{W}. Since they follow from \hyperref[post:III.1.W]{W}, we need only prove that \hyperref[post:III.1.W]{W} follows from \hyperref[post:III.1.E1]{$E_1$} or \hyperref[post:III.1.E2]{$E_2$}.

Let $R_1,\ldots,R_l$ be commuting operators belonging respectively to the quantities $\mathfrak R_1,\ldots,\mathfrak R_n$. According to \hyperref[sec:II.10]{II.10}, they are functions of a Hermitian operator $R$:
\[
R_1=F_1(R),\quad\ldots,\quad R_l=F_l(R).
\]
We assume that $R$ too belongs to a quantity $\mathfrak R$. (We therefore assume that to every quantity $\mathfrak R$ there belongs a [hypermaximal] Hermitian operator $R$, and conversely. See \hyperlink{note-119}{Note~119} and \hyperref[sec:IV.2]{IV.2}.) Then by \hyperref[post:III.1.F]{F},
\[
\mathfrak R_1=F_1(\mathfrak R),\quad\ldots,\quad\mathfrak R_l=F_l(\mathfrak R).
\]
Let $I_1,\ldots,I_l$ be the intervals occurring in \hyperref[post:III.1.W]{W}, and put
\[
G_j(\lambda)=\begin{cases}1,&\lambda\text{ in }I_j,\\0,&\text{otherwise},\end{cases}
\qquad(j=1,\ldots,l).
\]
We set
\[
H(\lambda)=G_1(F_1(\lambda))\cdots G_l(F_l(\lambda))
\]
and form the quantity
\[
\mathfrak S=H(\mathfrak R).
\]
If $\mathfrak R_j$ lies in $I_j$, that is, $F_j(\mathfrak R)$ lies in $I_j$, then $G_j(F_j(\mathfrak R))=1$; otherwise it is zero. Thus $\mathfrak S=H(\mathfrak R)$ equals $1$ when all the $\mathfrak R_j$ lie in their respective $I_j$ ($j=1,\ldots,l$), and zero otherwise. The expectation value of $\mathfrak S$
% Source PDF page 114 / original page 107
\hypertarget{origpage-107}{}%
therefore equals the probability $W$ that $\mathfrak R_1$ lies in $I_1$, $\ldots$, $\mathfrak R_l$ in $I_l$. Hence
\begin{align*}
W&=\operatorname{Erw}(\mathfrak S,\varphi)=(H(R)\varphi,\varphi)\\
&=(G_1(F_1(R))\cdots G_l(F_l(R))\varphi,\varphi)
=(G_1(R_1)\cdots G_l(R_l)\varphi,\varphi).
\end{align*}
Again denote by $E_j(\lambda)$ the resolution of the identity belonging to $R_j$, and let $I_j$ be the interval $\{\lambda'_j,\lambda''_j\}$. By the discussion at the end of \hyperref[sec:II.8]{II.8}, and using its notation,
\begin{align*}
G_j(\lambda)&=e_{\lambda''_j}(\lambda)-e_{\lambda'_j}(\lambda),\\
G_j(R_j)&=e_{\lambda''_j}(R_j)-e_{\lambda'_j}(R_j)
=E_j(\lambda''_j)-E_j(\lambda'_j)=E_j(I_j),
\end{align*}
and therefore
\[
W=(E_1(I_1)\cdots E_l(I_l)\varphi,\varphi).
\]
But this is precisely \hyperref[post:III.1.W]{W}.

Because of their simple form, \hyperref[post:III.1.E2]{$E_2$} and \hyperref[post:III.1.F]{F} are particularly suitable as the foundations on which the whole theory is built. We have seen that the most general possible probability assertion, \hyperref[post:III.1.W]{W}, follows from them. Two features of \hyperref[post:III.1.W]{W} are striking, however:

1. \hyperref[post:III.1.W]{W} is statistical, not causal: it does not specify which values $\mathfrak R_1,\ldots,\mathfrak R_l$ possess in the state $\varphi$, but only the probabilities with which they take all possible values.

2. The question posed in \hyperref[post:III.1.W]{W} cannot be answered for arbitrary quantities $\mathfrak R_1,\ldots,\mathfrak R_l$, but only for those whose operators $R_1,\ldots,R_l$ commute with one another.

Our next task is to discuss the significance of these two facts.

\section{The Statistical Interpretation}\label{sec:III.2}
Classical mechanics is a causal discipline: exact knowledge of the state of a system determines a unique numerical value for every physical quantity (energy, angular momentum, and so forth). For a system with $k$ degrees of freedom, this requires $2k$ numbers: the $k$ configuration-space coordinates $q_1,\ldots,q_k$ and their $k$ time derivatives $\frac{\partial q_1}{\partial t},\ldots,\frac{\partial q_k}{\partial t}$, or, in place of the latter, the $k$ momenta $p_1,\ldots,p_k$. Nevertheless, a statistical treatment of classical mechanics also exists. It is, so to speak, a luxury, an addition: if not all $2k$ determining data $(q_1,\ldots,q_k,p_1,\ldots,p_k)$ are known, but only some of them, and perhaps even these not quite exactly, then by averaging in some manner over the data that remain unknown one can at least make statistical assertions about all physical quantities. The same applies to the earlier or
% Source PDF page 115 / original page 108
\hypertarget{origpage-108}{}%
later states of the system. If $q_1,\ldots,q_k,p_1,\ldots,p_k$ are known at time $t=t_0$, the equations of motion of classical mechanics enable us to calculate them causally at any other time $t$. If only some of them are known, however, we must average over the others and can make only statistical assertions about other times.\origNote{120}

The statistical assertions we encountered in quantum mechanics have a different character. Here, for $k$ degrees of freedom, the state is described by the wave function $\varphi(q_1\cdots q_k)$, that is, by a point $\varphi$ in a suitably realized $\mathfrak R_\infty$ ($\|\varphi\|=1$; a numerical factor of absolute value $1$ is immaterial). Thus, although we believe that specifying $\varphi$ gives complete knowledge of the state, it nevertheless permits only statistical assertions about the values of physical quantities.

This statistical character is restricted to predicting the values of physical quantities. The earlier and later states $\varphi_t$ can be calculated causally from $\varphi_{t_0}=\varphi$. Schrödinger's time-dependent differential equation (cf.\ \hyperref[sec:I.2]{I.2}) makes this possible, for
\[
\varphi_{t_0}=\varphi,\qquad
\frac h{2\pi i}\frac\partial{\partial t}\varphi_t=-\mathsf H\varphi_t
\]
determine the entire course of $\varphi_t$. This differential equation can even be solved explicitly:
\[
\varphi_t=e^{-\frac{2\pi i}h(t-t_0)\mathsf H}\varphi
\]
($e^{-\frac{2\pi i}h(t-t_0)\mathsf H}$ is unitary).\origNote{121} In this formula $\mathsf H$ was assumed independent of time, but $\varphi_t$ is uniquely determined even for a time-dependent $\mathsf H$, since the differential equation is of first order; there is then no equally simple explicit formula for the solution.

If one wishes to explain the acausal connection between $\varphi$ and the values of physical quantities on the model of classical mechanics, the following view plainly suggests itself. In reality $\varphi$ does not determine the state exactly at all: further numerical data are needed to know it completely. Besides $\varphi$, there are further data that determine the system, further coordinates. If all these were known, the values of all physical quantities could be specified exactly and definitely; with $\varphi$ alone, however, only statistical assertions are possible, just as in classical mechanics when only some of the $q_1,\ldots,q_k,p_1,\ldots,p_k$ are given. This view is of course only hypothetical. It is an attempt whose value depends on whether the additional coordinates supplementing $\varphi$ can actually be found and used to construct a causal theory that agrees with experience and, when only
% Source PDF page 116 / original page 109
\hypertarget{origpage-109}{}%
$\varphi$ is specified and the remaining coordinates are averaged over, again yields the statistical assertions of quantum mechanics.

These hypothetical additional coordinates are generally called ``hidden coordinates'' or ``hidden parameters'', since they would have to play a rather hidden role alongside $\varphi$, which alone has been revealed by research so far. In classical physics, explanation by hidden parameters has already reduced much apparently statistical behaviour to the causal foundation of mechanics; the kinetic theory of gases is a characteristic example (cf.\ \hyperlink{note-120}{Note~120}).

Whether such an explanation by hidden parameters is possible for quantum mechanics is a much-discussed question. The view that it will one day be answered affirmatively still has eminent advocates. If justified, it would stamp the present form of the theory as provisional, since the $\varphi$-description of states would then be essentially incomplete.

We shall show later (\hyperref[sec:IV.2]{IV.2}) that hidden parameters certainly cannot be introduced without substantially changing the present theory. For the moment we merely point out that $\varphi$ differs quite essentially from partial sets of the $q_1,\ldots,q_k,p_1,\ldots,p_k$ in classical mechanics in that its time dependence is causal, not statistical: as we saw above, $\varphi_{t_0}$ determines every $\varphi_t$ uniquely.

Until a closer analysis of the assertions of quantum mechanics enables us to examine objectively the possibility of introducing hidden parameters, as will be done at the place cited above, we shall dispense with this possible explanation. We therefore adopt the opposite standpoint: we accept the fact that the laws of nature governing elementary processes, that is, the laws of quantum mechanics, are statistical in character. (At most, causality in the macroscopic world can be simulated by the levelling effect of the ``law of large numbers'', which becomes appreciable when many elementary processes act together simultaneously. See the remark at the end of \hyperlink{note-120}{Note~120} and \hyperlink{note-175}{Note~175}.) Accordingly, we accept \hyperref[post:III.1.W]{W}, or equivalently \hyperref[post:III.1.E2]{$E_2$}, as the most far-reaching assertion concerning elementary processes.

This view of quantum mechanics, which accepts its statistical assertions as the actual form of the laws of nature and abandons the principle of causality, is the so-called statistical interpretation. It is due to M. Born.\origNote{122} It is the only interpretation of quantum mechanics---that is, of the sum of our experience concerning elementary processes---that can be consistently carried through today. We shall adopt it in what follows, until we can undertake a detailed discussion of the situation.

% Source PDF page 117 / original page 110
\hypertarget{origpage-110}{}%
\section{Simultaneous Measurability and Measurability in General}\label{sec:III.3}
The second circumstance we noted as striking at the end of \hyperref[sec:III.1]{III.1} was connected with the fact that \hyperref[post:III.1.W]{W} gave information not only about the probabilities with which one quantity $\mathfrak R$ takes given numerical values, but also about probability relations between several quantities $\mathfrak R_1,\ldots,\mathfrak R_l$. \hyperref[post:III.1.W]{W} gave the probabilities that these simultaneously take certain given values, or more precisely lie in certain intervals $I_1,\ldots,I_l$, all in a given state $\varphi$. But these quantities $\mathfrak R_1,\ldots,\mathfrak R_l$ were subject to a peculiar restriction: their operators $R_1,\ldots,R_l$ had to commute. For noncommuting $R_1,\ldots,R_l$, \hyperref[post:III.1.W]{W} gave no information at all about the mutual probability dependence of $\mathfrak R_1,\ldots,\mathfrak R_l$; it could be used only to determine the probability distribution of each quantity separately, without regard to the others.

The most natural assumption would be that this is an imperfection of \hyperref[post:III.1.W]{W} and that a more general formula, covering these cases too, must exist. For even if quantum mechanics gives only statistical information about nature, the least we may expect of it is that it should describe not only the statistics of individual quantities but also the correlations between several quantities.

Contrary to this view, which at first seems reasonable, we shall soon see that no such generalization of \hyperref[post:III.1.W]{W} is possible. Besides formal reasons arising from the structure of the theory's mathematical apparatus, there are also weighty physical reasons for such a restriction. The necessity of this restriction and its physical interpretation will even give us essential insights into the nature of elementary processes.

To clarify this, we must examine more closely what the process of measuring a quantity $\mathfrak R$, about which \hyperref[post:III.1.W]{W} makes a probability assertion, means in the quantum-mechanical description.

First we draw attention to an important experiment carried out by Compton and Simons even before quantum mechanics was established.\origNote{123} In this experiment light was scattered by electrons, and the scattering process was checked quantitatively by subsequently intercepting the scattered light and electrons and measuring their energy and momentum. Thus collisions took place between light quanta and electrons, and the
% Source PDF page 118 / original page 111
\hypertarget{origpage-111}{}%
observer could test whether the laws of elastic collision were satisfied by measuring their paths after the collision. (Only elastic collisions need be considered, since neither an electron nor a light quantum was believed capable of absorbing energy in any form other than kinetic energy. In the light of all experience, both appear to have a uniquely fixed structure! The collision calculation must of course be relativistic.\origNote{123}) Such a numerical check was indeed possible: the paths before the collision were known and those after it were observed, so the collision problem was overdetermined. To determine it mechanically, two of these four paths and the ``central line'' of the collision---the direction of momentum transfer---suffice; thus at most three paths are needed, and the fourth provides the check. The experiment fully confirmed the mechanical laws of collision.

If we accept the validity of the collision laws and regard the paths before the collision as known, this result can also be formulated as follows. Measuring either the light quantum's path after the collision or the electron's path suffices to calculate the location and central line of the collision. The Compton--Simons experiment shows that these two observations give the same result.

More generally: the same physical quantity, namely any coordinate of the collision point or of the direction of the central line, is measured in two different ways, by intercepting the light quantum and by intercepting the electron; the result is always the same.

These two measurements are not quite simultaneous. The light quantum and the electron do not arrive at the same time, and by suitably arranging the measuring apparatus one can intercept whichever one wishes first, though the difference is only $10^{-9}$--$10^{-10}$ seconds. Call the earlier measurement $M_1$, the later $M_2$, and the measured quantity $\mathfrak R$. The essential point is this: although the entire arrangement is such that before measurement only statistical predictions can be made about $\mathfrak R$, that is, about $M_1,M_2$ (see the reference in \hyperlink{note-123}{Note~123}), the statistical correlation between $M_1,M_2$ is perfectly sharp, or causal: the $\mathfrak R$-value of $M_1$ is certainly equal to that of $M_2$. Before $M_1,M_2$, therefore, both results are largely indeterminate; after $M_1$ has occurred but before $M_2$, the result of $M_2$ is already fixed causally and uniquely.

As a matter of principle, three degrees of causality or acausality are conceivable. First, the value of $\mathfrak R$ might be entirely statistical, so that the result of a measurement could be predicted only statistically; and if a second measurement were made immediately after the first, its result might again show dispersion regardless of the value found in the first---possibly just as much dispersion as the first.\origNote{124} Second, it is conceivable that
% Source PDF page 119 / original page 112
\hypertarget{origpage-112}{}%
the value of $\mathfrak R$ shows dispersion in the first measurement, but every measurement immediately following is compelled to give a result agreeing with that of the first. Third, $\mathfrak R$ might be causally determined from the outset.

The Compton--Simons experiment shows that only the second case is possible for a statistical theory. If, therefore, the system is initially in a state in which the value of $\mathfrak R$ cannot be predicted with certainty, a measurement $M$ of $\mathfrak R$ ($M_1$ in the example above) transforms this state into another: one in which the value of $\mathfrak R$ is uniquely fixed. The new state into which $M$ puts the system depends not only on the arrangement of $M$ but also on its result, which could not be predicted causally in the old state; for the value of $\mathfrak R$ in the new state must be precisely equal to this result of $M$.

Let $\mathfrak R$ now be a quantity whose operator $R$ has a pure point spectrum $\lambda_1,\lambda_2,\ldots$, with respective eigenfunctions $\varphi_1,\varphi_2,\ldots$ forming a complete orthonormal system. Suppose every eigenvalue is simple, that is, of multiplicity $1$ (cf.\ \hyperref[sec:II.6]{II.6}), so $\lambda_\mu\ne\lambda_\nu$ for $\mu\ne\nu$. Suppose we have measured $\mathfrak R$ and found a value $\lambda^*$. What is the state of the system after the measurement?

By what was said above, this state $\varphi$ must be such that a new measurement of $\mathfrak R$ yields $\lambda^*$ with certainty. (Naturally this measurement must be made immediately, since after $\tau$ seconds $\varphi$ has changed to $e^{-\frac{2\pi i}h\tau\mathsf H}\varphi$. See \hyperref[sec:III.2]{III.2}; $\mathsf H$ is the energy operator.)

We shall answer more generally the question of when measurement of $\mathfrak R$ in the state $\varphi$ gives $\lambda^*$ with certainty, without restrictive assumptions on $R$.

Let $E(\lambda)$ be the resolution of the identity belonging to $R$, and $I$ an interval $\{\lambda',\lambda''\}$. Our assumption can also be expressed by saying that $\mathfrak R$ lies in $I$ with probability $0$ if $I$ does not contain $\lambda^*$, or with probability $1$ if $\lambda^*$ belongs to $I$, that is, if $\lambda'<\lambda^*\le\lambda''$.

According to \hyperref[post:III.1.W]{W} this means $\|E(I)\varphi\|^2=1$, or, since $\|\varphi\|=1$, $\|E(I)\varphi\|=\|\varphi\|$. Since $E(I)$ and $1-E(I)$ are projection operators, \hyperref[thm:II.13]{Theorem~13 of II.4} gives
\begin{gather*}
\|\varphi-E(I)\varphi\|^2=\|\varphi\|^2-\|E(I)\varphi\|^2=0,\\
\varphi-E(I)\varphi=0,\\
E(\lambda'')\varphi-E(\lambda')\varphi=E(I)\varphi=\varphi.
\end{gather*}
Letting $\lambda'\to-\infty$ gives $E(\lambda'')\varphi=\varphi$; letting $\lambda''\to+\infty$ gives $E(\lambda')\varphi=0$.
% Source PDF page 120 / original page 113
\hypertarget{origpage-113}{}%
(Cf.\ $S_1$ in \hyperref[sec:II.7]{II.7}.) Thus
\[
E(\lambda)\varphi=\begin{cases}\varphi,&\lambda\ge\lambda^*,\\0,&\lambda<\lambda^*.\end{cases}
\]
According to \hyperref[sec:II.8]{II.8}, this is precisely what characterizes $R\varphi=\lambda^*\varphi$.

Another way of proving $R\varphi=\lambda^*\varphi$ rests on \hyperref[post:III.1.E1]{$E_1$}, or equivalently \hyperref[post:III.1.E2]{$E_2$}. That $\mathfrak R$ definitely has the value $\lambda^*$ means that $(\mathfrak R-\lambda^*)^2$ has expectation zero. Its operator is $F(R)$ with $F(\lambda)=(\lambda-\lambda^*)^2$, hence $(R-\lambda^*\cdot1)^2$. We must therefore have
\begin{align*}
((R-\lambda^*\cdot1)^2\varphi,\varphi)
&=((R-\lambda^*\cdot1)\varphi,(R-\lambda^*\cdot1)\varphi)\\
&=\|(R-\lambda^*\cdot1)\varphi\|^2
=\|R\varphi-\lambda^*\varphi\|^2=0,
\end{align*}
that is, $R\varphi=\lambda^*\varphi$.

In the special case considered above, then, we see that $R\varphi=\lambda^*\varphi$ is necessary. As discussed in \hyperref[sec:II.6]{II.6}, this implies that $\lambda^*$ equals some $\lambda_\mu$ (since $\|\varphi\|=1$, $\varphi\ne0$), and that $\varphi=a\varphi_\mu$. Since $\|\varphi\|=\|\varphi_\mu\|=1$, we must also have $|a|=1$, so $a$ may be omitted without changing the state. Thus $\lambda^*=\lambda_\mu$, $\varphi=\varphi_\mu$ for some $\mu=1,2,\ldots$. (The assertion about $\lambda^*$ could also have been deduced directly from \hyperref[post:III.1.W]{W}, but the assertion about $\varphi$ could not.)

Under the above assumptions on $R$, a measurement of $\mathfrak R$ consequently transforms any state $\psi$ into one of the states $\varphi_1,\varphi_2,\ldots$, associated respectively with the results $\lambda_1,\lambda_2,\ldots$. The probabilities of these transformations therefore equal the probabilities of measuring $\lambda_1,\lambda_2,\ldots$ and can be calculated from \hyperref[post:III.1.W]{W}.

According to \hyperref[post:III.1.W]{W}, the probability that the value of $\mathfrak R$ lies in $I$ is $\|E(I)\psi\|^2$. Since by \hyperref[sec:II.8]{II.8}, $E(I)=\sum_{\lambda_n\text{ in }I}P_{[\varphi_n]}$, we obtain
\begin{align*}
W&=\|E(I)\psi\|^2=(E(I)\psi,\psi)
=\sum_{\lambda_n\text{ in }I}(P_{[\varphi_n]}\psi,\psi)\\
&=\sum_{\lambda_n\text{ in }I}|(\psi,\varphi_n)|^2.
\end{align*}
One therefore expects the probability of $\lambda_n$ to be $|(\psi,\varphi_n)|^2$. If we can choose $I$ so that the only $\lambda_m$ in it is $\lambda_n$, this follows directly from the formula above. Otherwise, when the other $\lambda_m$ accumulate at $\lambda_n$, we may argue as follows. Let $F(\lambda)$ be $1$ for $\lambda=\lambda_n$ and zero otherwise. The required probability $W_n$ is the expectation value of $F(\mathfrak R)$, hence, by \hyperref[post:III.1.E2]{$E_2$} or \hyperref[post:III.1.E1]{$E_1$}, it is $(F(R)\psi,\psi)$. By definition (\hyperref[sec:II.8]{II.8}),
\[
(F(R)\psi,\psi)=\int_{-\infty}^{\infty}F(\lambda)\,d(\|E(\lambda)\psi\|^2).
\]
Recalling the definition of the Stieltjes integral, we readily see that this is zero if $E(\lambda)\psi$ is continuous
% Source PDF page 121 / original page 114
\hypertarget{origpage-114}{}%
in $\lambda$ at $\lambda=\lambda_n$, and in general equals the jump of the monotonically nondecreasing function $\|E(\lambda)\psi\|^2$ at $\lambda=\lambda_n$. This jump equals $\|P_{\mathfrak M}\psi\|^2$, where $\mathfrak M$ is the closed linear manifold spanned by all solutions of $R\psi=\lambda_n\psi$ (cf.\ \hyperref[sec:II.8]{II.8}). In the present case $\mathfrak M=[\varphi_n]$, so
\[
W_n=\|P_{[\varphi_n]}\psi\|^2=|(\psi,\varphi_n)|^2.
\]
This answers the question of what happens when a quantity $\mathfrak R$ is measured, under the above assumptions on its operator $R$. For the moment, of course, the ``how'' remains unexplained. This discontinuous, jump-like transition from $\psi$ into one of the states $\varphi_1,\varphi_2,\ldots$ is certainly not of the kind described by Schrödinger's time-dependent differential equation. These states are independent of $\psi$, which enters only into the respective probabilities $W_n=|(\psi,\varphi_n)|^2$, $n=1,2,\ldots$, of the jumps. The differential equation always produces a continuous change of $\psi$, with a uniquely determined final result depending on $\psi$ (cf.\ \hyperref[sec:III.2]{III.2}). Only later shall we attempt to bridge this gap (see \hyperref[chap:VI]{VI}).\origNote{125}

Continue to assume that $R$ has a pure point spectrum, but no longer require all eigenvalues to be simple. We can again form $\varphi_1,\varphi_2,\ldots$ and $\lambda_1,\lambda_2,\ldots$, but equal values may now occur among the $\lambda_n$. After a measurement of $\mathfrak R$, the state $\varphi$ certainly satisfies $R\varphi=\lambda^*\varphi$, where $\lambda^*$ is the measurement result. Thus $\lambda^*$ equals some $\lambda_n$, but concerning $\varphi$ we can say only this: if the $\lambda_n$ equal to $\lambda^*$ are $\lambda_{n_1},\lambda_{n_2},\ldots$, in finite or infinite number, then
\[
\varphi=\sum_\nu a_\nu\varphi_{n_\nu}.
\]
(If there are infinitely many $n_\nu$, $\sum_\nu|a_\nu|^2$ must be finite.) Two such $\varphi$ represent the same state only if they differ by a numerical factor, that is, only if the ratio $a_1:a_2:\cdots$ is the same. As soon as more than one $n_\nu$ exists, then, that is, when the eigenvalue $\lambda^*$ is multiple, even knowledge of the measurement result does not uniquely determine the state $\varphi$ after the measurement.

The probability of $\lambda^*$ is calculated from \hyperref[post:III.1.W]{W}, or \hyperref[post:III.1.E1]{$E_1$} or \hyperref[post:III.1.E2]{$E_2$}, exactly as before:
\[
W(\lambda^*)=\sum_{\lambda_n=\lambda^*}|(\psi,\varphi_n)|^2
=\sum_\nu|(\psi,\varphi_{n_\nu})|^2.
\]
If $R$ does not have a pure point spectrum, the situation is as follows. All solutions $f$ of $Rf=\lambda f$ span a closed linear manifold $\mathfrak M_\lambda$, and all the $\mathfrak M_\lambda$ together span a further one $\overline{\mathfrak M}$. The absence of a pure point spectrum is characterized by $\overline{\mathfrak M}\ne\mathfrak R_\infty$, that is, $\overline{\mathfrak N}=\mathfrak R_\infty-\overline{\mathfrak M}\ne(0)$. (For this and what follows, see \hyperref[sec:II.8]{II.8}.) There are at most countably many values of $\lambda$ for which $\mathfrak M_\lambda\ne(0)$; these
% Source PDF page 122 / original page 115
\hypertarget{origpage-115}{}%
form the point spectrum of $R$. If we measure $\mathfrak R$ in the state $\psi$, the probability that the result will be $\lambda^*$ is
\[
W(\lambda^*)=\|P_{\mathfrak M_{\lambda^*}}\psi\|^2
=(P_{\mathfrak M_{\lambda^*}}\psi,\psi).
\]
The most convenient derivation is the argument used above, based on \hyperref[post:III.1.E2]{$E_2$} (or \hyperref[post:III.1.E1]{$E_1$}) and the function
\[
F(\lambda)=\begin{cases}1,&\lambda=\lambda^*,\\0,&\lambda\ne\lambda^*.\end{cases}
\]
Correspondingly, the probability that the value of $\mathfrak R$ will be some $\lambda^*$ in the point spectrum $\mathfrak P$ of $R$ is
\[
W=\sum_{\lambda^*\text{ in }\mathfrak P}(P_{\mathfrak M_{\lambda^*}}\psi,\psi)
=(P_{\overline{\mathfrak M}}\psi,\psi)=\|P_{\overline{\mathfrak M}}\psi\|^2.
\]
This can also be seen directly by means of the function
\[
F(\lambda)=\begin{cases}1,&\lambda^*\text{ in }\mathfrak P,\\0,&\text{otherwise}.\end{cases}
\]
But if we measure $\mathfrak R$ exactly, the subsequent state $\varphi$ must satisfy $R\varphi=\lambda^*\varphi$, so the result $\lambda^*$ must belong to $\mathfrak P$. The probability that an exact measurement succeeds is therefore at most $\|P_{\overline{\mathfrak M}}\psi\|^2$. This number is not always $1$; for the $\psi$ in $\overline{\mathfrak N}$ it is even zero. Thus an exact measurement is not always possible.

We have seen that a quantity $\mathfrak R$ can always, that is, in every state $\psi$, be measured exactly if and only if it has a pure point spectrum. If it does not, it can be measured only with limited accuracy: the number line $-\infty<\lambda<+\infty$ can be divided into intervals $\ldots,I^{(-2)},I^{(-1)},I^{(0)},I^{(1)},I^{(2)},\ldots$, and one can decide in which interval $\mathfrak R$ lies. Let the division points be $\ldots,\lambda^{(-2)},\lambda^{(-1)},\lambda^{(0)},\lambda^{(1)},\lambda^{(2)},\ldots$, with $I^{(n)}=\{\lambda^{(n)},\lambda^{(n+1)}\}$. The maximum interval length $\varepsilon=\operatorname{Max}(\lambda^{(n+1)}-\lambda^{(n)})$, or mesh of the division, is the measurement accuracy. We can analyze this latter process mathematically. Let $F(\lambda)$ be the following function, where $\lambda'_n$ is an intermediate value in $I^{(n)}$, chosen arbitrarily for each $n=0,\pm1,\pm2,\ldots$ but then kept fixed:
\[
F(\lambda)=\lambda'_n\quad\text{if }\lambda\text{ lies in }I^{(n)}.
\]
The approximate measurement of $\mathfrak R$ is then equivalent to the exact measurement of $F(\mathfrak R)$. Now
\begin{align*}
F(R)&=\int_{-\infty}^{\infty}F(\lambda)\,dE(\lambda)
=\sum_{n=-\infty}^{+\infty}\int_{\lambda^{(n)}}^{\lambda^{(n+1)}}F(\lambda)\,dE(\lambda)\\
&=\sum_{n=-\infty}^{+\infty}\lambda'_n\int_{\lambda^{(n)}}^{\lambda^{(n+1)}}dE(\lambda)
=\sum_{n=-\infty}^{+\infty}\lambda'_nE(I^{(n)}).
\end{align*}
% Source PDF page 123 / original page 116
\hypertarget{origpage-116}{}%
Plainly $F(R)f=\lambda'_nf$ for every $f$ in the closed linear manifold belonging to $E(I^{(n)})$. Thus, for $F(R)$, the manifold $\mathfrak M_{\lambda'_n}$ contains this manifold. Consequently $P_{\mathfrak M_{\lambda'_n}}\ge E(I^{(n)})$, so
\begin{align*}
P_{\overline{\mathfrak M}}
&\ge\sum_{n=-\infty}^{+\infty}P_{\mathfrak M_{\lambda'_n}}
\ge\sum_{n=-\infty}^{+\infty}E(I^{(n)})\\
&=\sum_{n=-\infty}^{+\infty}(E(\lambda^{(n+1)})-E(\lambda^{(n)}))=1-0=1.
\end{align*}
It follows that
\[
\sum_{n=-\infty}^{+\infty}P_{\mathfrak M_{\lambda'_n}}=P_{\overline{\mathfrak M}}=1,
\qquad P_{\mathfrak M_{\lambda'_n}}=E(I_n),
\]
that is, $F(R)$ has a pure point spectrum consisting of the $\lambda'_n$.

Thus $F(\mathfrak R)$ really is exactly measurable, and the probability that its value is $\lambda'_n$, or equivalently that the value of $\mathfrak R$ lies in $I^{(n)}$, is
\[
\|P_{\mathfrak M_{\lambda'_n}}\psi\|^2=\|E(I^{(n)})\psi\|^2,
\]
in agreement with the assertion of \hyperref[post:III.1.W]{W} for $\mathfrak R$.

This result can also be interpreted physically and shows good agreement between the theory and ordinary physical intuition.

In the classical-mechanical view, without any quantum conditions, every quantity $\mathfrak R$ is assigned a perfectly definite value in every state. At the same time it is recognized that every conceivable measuring apparatus can determine this value only within some margin of error that never vanishes, owing to the imperfections of human observation: a pointer can be read, or the darkening of a photographic plate localized, only with limited accuracy. By sufficiently refining the measuring arrangement this margin of error can be reduced arbitrarily close to zero, but it never becomes exactly zero. We should expect the same to hold in quantum theory for quantities $\mathfrak R$ that are not quantized according to the intuitive picture one used to form of these things, especially before quantum mechanics was discovered. Examples are the Cartesian coordinates of an electron, which can take any value between $-\infty$ and $+\infty$ and whose operators have continuous spectra. For quantities that are ``quantized'' according to that intuitive picture, the situation is reversed. Since they can take only discrete values, it suffices to observe them accurately enough to leave no doubt which of the ``quantized'' values is involved; one then knows with certainty the exact value taken. For example, if a hydrogen atom is known to contain less energy than is needed for the second-lowest energy level, its energy content is known with absolute exactness: it is the lowest energy level.

% Source PDF page 124 / original page 117
\hypertarget{origpage-117}{}%
As the situation in ``matrix theory'' already shows (cf.\ \hyperref[sec:I.2]{I.2} and \hyperref[sec:II.6]{II.6}), this division into ``quantized'' and ``unquantized'' quantities corresponds to the division into quantities $\mathfrak R$ whose operator $R$ has a pure point spectrum and those whose operator does not. And precisely for the former, and only for these, we found absolutely exact observation possible; the latter could be observed with arbitrarily great accuracy, but never with absolute accuracy.\origNote{126}

(Incidentally, it is just here that the introduction of ``improper'' eigenfunctions, or eigenfunctions not belonging to Hilbert space, mentioned in the \hyperref[chap:introduction]{Introduction} and in \hyperref[sec:I.3]{I.3}---see also \hyperref[sec:II.8]{II.8}, especially \hyperlink{note-84}{Note~84} and \hyperlink{note-86}{Note~86}---represents reality less well than our procedure. It creates the appearance that states exist in which quantities with a continuous spectrum take certain values exactly, although this is precisely what never occurs. Besides their mathematical untenability, we believe this is a further reason for rejecting such idealizations, despite their having been proposed on several occasions.)

We have thus brought the question of the processes occurring when one quantity is measured to a provisional conclusion, and can turn to the simultaneous measurement of several quantities.

First let $\mathfrak R,\mathfrak S$ be two quantities with respective operators $R,S$, and assume that they are simultaneously measurable. What follows from this?

Initially we shall require absolutely exact measurements, so both $R,S$ must have pure point spectra, respectively $\lambda_1,\lambda_2,\ldots$ and $\mu_1,\mu_2,\ldots$. Let the corresponding complete orthonormal systems of eigenfunctions be $\varphi_1,\varphi_2,\ldots$ and $\psi_1,\psi_2,\ldots$.

To dispose of the simplest case first, suppose one of the two operators, say $R$, has only simple eigenvalues: $\lambda_m\ne\lambda_n$ for $m\ne n$.

If we measure $\mathfrak R,\mathfrak S$ simultaneously, the subsequent state is one in which both $\mathfrak R$ and $\mathfrak S$ definitely have the values just measured, say $\lambda_{\bar m},\mu_{\bar n}$. The subsequent state $\psi$ must therefore satisfy $R\psi=\lambda_{\bar m}\psi$, $S\psi=\mu_{\bar n}\psi$. The first equation gives $\psi=\varphi_{\bar m}$, apart from a numerical factor that may be omitted. The second gives $\psi=\sum_\nu a_\nu\psi_{n_\nu}$, where $\mu_{n_1},\mu_{n_2},\ldots$ are all the $\mu_n$ equal to $\mu_{\bar n}$. If the initial state was $\varphi$, then $\lambda_{\bar m},\varphi_{\bar m}$ have probability $|(\varphi,\varphi_{\bar m})|^2$. For $\varphi=\varphi_m$, therefore, $\bar m=m$ becomes certain. Thus for every $m$ we can say that $\varphi_m$ is a sum $\sum_\nu a_\nu\psi_{n_\nu}$ with all the $\mu_{n_\nu}$ equal, or $S\varphi_m=\bar\mu\varphi_m$, where $\bar\mu=\mu_{n_1}=\mu_{n_2}=\cdots$. Hence $RSf=SRf$ for $f=\varphi_m$, both sides being $\lambda_m\bar\mu\cdot\varphi_m$. The same therefore holds for their linear combinations and, if $R,S$ are continuous, for their limit points, that is, for every $f$. Thus $R,S$ commute.
% Source PDF page 125 / original page 118
\hypertarget{origpage-118}{}%

If they are not continuous, we argue as follows. The resolutions of the identity $E(\lambda),F(\mu)$ belonging to $R,S$ are defined by
\[
E(\lambda)=\sum_{\lambda_m\le\lambda}P_{[\varphi_m]},\qquad
F(\mu)=\sum_{\mu_n\le\mu}P_{[\psi_n]}.
\]
Thus $F(\mu)\varphi_m=\varphi_m$ or $0$, according as $\mu\ge\bar\mu$ or $\mu<\bar\mu$, with $\bar\mu$ as above. Also $E(\lambda)\varphi_m=\varphi_m$ or $0$, according as $\lambda\ge\lambda_m$ or $\lambda<\lambda_m$. In every case, therefore, $E(\lambda)F(\mu)\varphi_m=F(\mu)E(\lambda)\varphi_m$ for all $\varphi_m$. As before, this gives commutativity of $E(\lambda),F(\mu)$, and hence, by \hyperref[sec:II.10]{II.10}, of $R,S$.

According to \hyperref[sec:II.10]{II.10}, however, there exists a complete orthonormal system of common eigenfunctions of $R,S$, so we may assume $\varphi_m=\psi_m$. Since $m\ne n$ implies $\lambda_m\ne\lambda_n$, we can construct a function
\[
F(\lambda)=\begin{cases}\mu_n,&\lambda=\lambda_n,\ n=1,2,\ldots,\\
\text{arbitrary},&\text{otherwise},\end{cases}
\]
and obtain $S=F(R)$, that is, $\mathfrak S=F(\mathfrak R)$. Thus $\mathfrak R,\mathfrak S$ are not merely simultaneously measurable: every measurement of $\mathfrak R$ is itself a measurement of $\mathfrak S$, since $\mathfrak S$ is a function of $\mathfrak R$, and hence causally determined by it.\origNote{127}

We now turn to the other, more general case, in which nothing is assumed about the multiplicities of the eigenvalues of $R,S$. Here we use an essentially different method.

First consider the quantity $\mathfrak R+\mathfrak S$. A simultaneous measurement of $\mathfrak R,\mathfrak S$ is also a measurement of $\mathfrak R+\mathfrak S$, since adding the results gives the value of $\mathfrak R+\mathfrak S$. Consequently, in every state $\psi$, the expectation value of $\mathfrak R+\mathfrak S$ is the sum of the expectation values of $\mathfrak R$ and $\mathfrak S$. Notice that this holds regardless of whether $\mathfrak R,\mathfrak S$ are statistically independent, whether correlations exist between them and of what kind, or even whether one causally determines the other; for, as is well known, the law
\[
\text{expectation of the sum}=\text{sum of the expectations}
\]
holds generally. If $T$ is the operator of $\mathfrak R+\mathfrak S$, this expectation is on the one hand $(T\psi,\psi)$ and on the other hand
\[
(R\psi,\psi)+(S\psi,\psi)=((R+S)\psi,\psi).
\]
Thus for every $\psi$,
\[
(T\psi,\psi)=((R+S)\psi,\psi),
\]
so $T=R+S$. Therefore $\mathfrak R+\mathfrak S$ has the operator $R+S$.\origNote{128} Similarly, $a\mathfrak R+b\mathfrak S$, for real $a,b$, has the operator $aR+bS$. This also follows from the first formula by substituting $\mathfrak R,\mathfrak S$ and $R,S$ in the functions $F(\lambda)=a\lambda$, $G(\mu)=b\mu$.

% Source PDF page 126 / original page 119
\hypertarget{origpage-119}{}%
A simultaneous measurement of $\mathfrak R,\mathfrak S$ also measures
\begin{gather*}
\frac{\mathfrak R+\mathfrak S}2,\quad
\left(\frac{\mathfrak R+\mathfrak S}2\right)^2,\quad
\frac{\mathfrak R-\mathfrak S}2,\quad
\left(\frac{\mathfrak R-\mathfrak S}2\right)^2,\\
\left(\frac{\mathfrak R+\mathfrak S}2\right)^2-
\left(\frac{\mathfrak R-\mathfrak S}2\right)^2=\mathfrak R\cdot\mathfrak S.
\end{gather*}
Their operators are therefore the following, using also the fact that if $T$ is the operator of $\mathfrak T$, then $F(\mathfrak T)$ has the operator $F(T)$, and thus $\mathfrak T^2$ has the operator $T^2$:
\begin{gather*}
\frac{R+S}2,\quad
\left(\frac{R+S}2\right)^2=\frac{R^2+S^2+RS+SR}4,\\
\frac{R-S}2,\quad
\left(\frac{R-S}2\right)^2=\frac{R^2+S^2-RS-SR}4,\\
\left(\frac{R+S}2\right)^2-\left(\frac{R-S}2\right)^2=\frac{RS+SR}2.
\end{gather*}
Thus $\mathfrak R\cdot\mathfrak S$ has the operator $(RS+SR)/2$. This also applies to all $F(\mathfrak R),G(\mathfrak S)$, which are measured at the same time, so $F(\mathfrak R)\cdot G(\mathfrak S)$ has the operator
\[
\frac{F(R)G(S)+G(S)F(R)}2.
\]
Let $E(\lambda),F(\mu)$ be the resolutions of the identity belonging to $R,S$, and let
\[
F(\lambda)=\begin{cases}1,&\lambda\le\bar\lambda,\\0,&\lambda>\bar\lambda,\end{cases}
\qquad
G(\mu)=\begin{cases}1,&\mu\le\bar\mu,\\0,&\mu>\bar\mu.\end{cases}
\]
As we know, $F(R)=E(\bar\lambda)$, $G(S)=F(\bar\mu)$; hence $F(\mathfrak R)\cdot G(\mathfrak S)$ has the operator $(EF+FE)/2$, where we abbreviate $E(\bar\lambda),F(\bar\mu)$ to $E,F$. Since $F(\mathfrak R)$ is always $0$ or $1$, $F(\mathfrak R)^2=F(\mathfrak R)$, so
\[
F(\mathfrak R)\cdot(F(\mathfrak R)\cdot G(\mathfrak S))=F(\mathfrak R)\cdot G(\mathfrak S).
\]
Applying our multiplication formula to $F(\mathfrak R)$ and $F(\mathfrak R)\cdot G(\mathfrak S)$, all of which are simultaneously measurable, gives this product the operator
\[
\frac{E\frac{EF+FE}2+\frac{EF+FE}2E}2
=\frac{E^2F+2EFE+FE^2}4
=\frac{EF+FE+2EFE}4.
\]
This must equal $(EF+FE)/2$, whence
\[
EF+FE=2\cdot EFE.
\]
Left multiplication by $E$ gives
\[
E^2F+EFE=2E^2FE,\quad EF+EFE=2EFE,\quad EF=EFE;
\]
right multiplication gives
\[
EFE+FE^2=2EFE^2,\quad EFE+FE=2EFE,\quad FE=EFE.
\]
Therefore $EF=FE$: all $E(\bar\lambda),F(\bar\mu)$ commute, so $R,S$ again commute.

% Source PDF page 127 / original page 120
\hypertarget{origpage-120}{}%
According to \hyperref[sec:II.10]{II.10}, the condition that $R,S$ commute is equivalent to requiring the existence of a Hermitian operator $T$ of which $R,S$ are functions: $R=F(T)$, $S=G(T)$. If $T$ belongs to the quantity $\mathfrak T$, then $\mathfrak R=F(\mathfrak T)$, $\mathfrak S=G(\mathfrak T)$. This condition is also sufficient for simultaneous measurability, since a measurement of $\mathfrak T$---an absolutely exact one, for $\mathfrak T$ has a pure point spectrum, cf.\ \hyperref[sec:II.10]{II.10}---simultaneously measures its functions $\mathfrak R,\mathfrak S$. Commutativity of $R,S$ is therefore the necessary and sufficient condition.

If several quantities $\mathfrak R,\mathfrak S,\ldots$, in finite number, are given with operators $R,S,\ldots$, and absolutely exact measurement is again required, the criterion for simultaneous measurability is as follows. If all the quantities $\mathfrak R,\mathfrak S,\ldots$ are simultaneously measurable, so must every pair formed from them be; thus all the operators $R,S,\ldots$ commute with one another. Conversely, if $R,S,\ldots$ commute, \hyperref[sec:II.10]{II.10} gives an operator $T$ of which all are functions: $R=F(T)$, $S=G(T)$, $\ldots$. Thus, for the corresponding quantity $\mathfrak T$, $\mathfrak R=F(\mathfrak T)$, $\mathfrak S=G(\mathfrak T)$, $\ldots$. An exact measurement of $\mathfrak T$, which again has a pure point spectrum (cf.\ \hyperref[sec:II.10]{II.10}), is consequently a simultaneous measurement of $\mathfrak R,\mathfrak S,\ldots$. Thus commutativity of $R,S,\ldots$ is necessary and sufficient for simultaneous measurability of $\mathfrak R,\mathfrak S,\ldots$.

Now consider measurements that are not absolutely exact but have a prescribed, arbitrarily high, accuracy. The $R,S,\ldots$ need no longer have pure point spectra.

Measurements of $\mathfrak R,\mathfrak S,\ldots$ with limited accuracy achieve the same result as absolutely exact measurements of $F(\mathfrak R),G(\mathfrak S),\ldots$. Here $F(\lambda),G(\lambda),\ldots$ are certain functions constructed as described at the beginning of this section in the discussion of measurements with limited accuracy, where naturally only $F(\lambda)$ was specified. Therefore $\mathfrak R,\mathfrak S,\ldots$ are certainly simultaneously measurable if all $F(\mathfrak R),G(\mathfrak S),\ldots$ are simultaneously measurable with absolute accuracy. The latter is equivalent to commutativity of $F(R),G(S),\ldots$, which follows from commutativity of $R,S,\ldots$. Thus commutativity of $R,S,\ldots$ is in any event sufficient.

Conversely, suppose $\mathfrak R,\mathfrak S,\ldots$ are simultaneously measurable. A sufficiently accurate measurement of $\mathfrak R$ allows us to decide whether its value is $>\bar\lambda$ or $\le\bar\lambda$ (cf.\ our definition of ``limited accuracy'', discussed especially in \hyperlink{note-126}{Note~126}). If $F(\lambda)$ is defined to be $1$ for $\lambda\le\bar\lambda$ and $0$ for $\lambda>\bar\lambda$, then $F(\mathfrak R)$ can be measured with absolute accuracy. Correspondingly, with $G(\mu)=1$ for $\mu\le\bar\mu$ and $0$ for $\mu>\bar\mu$, $G(\mathfrak S)$ can be measured with absolute accuracy, and the two quantities can moreover be measured simultaneously. Thus $F(R),G(S)$ commute. Now let
% Source PDF page 128 / original page 121
\hypertarget{origpage-121}{}%
$E(\lambda),F(\mu)$ be the resolutions of the identity belonging to $R,S$. Then $F(R)=E(\bar\lambda)$, $G(S)=F(\bar\mu)$, so $E(\bar\lambda),F(\bar\mu)$ commute for all $\bar\lambda,\bar\mu$. Hence $R,S$ commute; and since the same must hold for every pair in $R,S,\ldots$, all these operators commute with one another. The condition is therefore necessary as well.

We thus see that simultaneous measurability of any finite number of quantities $\mathfrak R,\mathfrak S,\ldots$ is characterized by commutativity of their operators $R,S,\ldots$. This holds both for absolutely exact measurements and for measurements of arbitrarily high accuracy; in the former case one must additionally require the condition characterizing the possibility of absolutely exact measurement, namely that the operators have pure point spectra.

This supplies the mathematical proof that \hyperref[post:III.1.W]{W} makes the most far-reaching assertion possible in this theory, that is, in a theory containing \hyperref[post:III.1.W]{W}. It assumes only commutativity of the operators $R_1,\ldots,R_l$; without this, nothing could in any event be said about simultaneous measurement results for $\mathfrak R_1,\ldots,\mathfrak R_l$, since simultaneous measurements of these quantities would not be possible at all.

\section{Uncertainty Relations}\label{sec:III.4}
The preceding sections have given us essential information about the measurement process for one quantity or for several simultaneously measurable quantities. We must now examine the behaviour of quantities that cannot be measured simultaneously when we are interested in their statistics for the same system in the same state $\varphi$.

Let two such quantities $\mathfrak R,\mathfrak S$ and their noncommuting operators $R,S$ be given. Despite this assumption, there may be states $\varphi$ in which both quantities have sharp values, that is, zero dispersions: these are common eigenfunctions of both operators. These common eigenfunctions cannot, however, form a complete orthogonal system, since $R,S$ would then commute. (Consider the construction in \hyperref[sec:II.8]{II.8} of the corresponding resolutions of the identity $E(\lambda),F(\lambda)$. If $\varphi_1,\varphi_2,\ldots$ is the complete orthogonal system in question, both the $E(\lambda)$ and the $F(\lambda)$ are sums of $P_{[\varphi_\varrho]}$ and therefore commute because the $P_{[\varphi_\varrho]}$ commute.) The meaning of this is easily seen: the closed linear manifold $\mathfrak M$ spanned by these $\varphi$ must be smaller than $\mathfrak R_\infty$. Were it equal to $\mathfrak R_\infty$, the required complete orthonormal system could be constructed just as was done for a single operator at the beginning of \hyperref[sec:II.6]{II.6}.

For states in $\mathfrak M$, $\mathfrak R,\mathfrak S$ are simultaneously measurable. The simplest proof is to give a model of this simultaneous
% Source PDF page 129 / original page 122
\hypertarget{origpage-122}{}%
measurement. Since the common eigenfunctions $\varphi$ of $R,S$ span $\mathfrak M$, there is an orthonormal system of such $\varphi$, say $\varphi_1,\varphi_2,\ldots$, that likewise spans $\mathfrak M$, that is, is complete in $\mathfrak M$. This too is obtained by the construction in \hyperref[sec:II.6]{II.6} mentioned above. Extend $\varphi_1,\varphi_2,\ldots$ to a complete system $\varphi_1,\varphi_2,\ldots,\psi_1,\psi_2,\ldots$ by adjoining an orthonormal system $\psi_1,\psi_2,\ldots$ spanning $\mathfrak R_\infty-\mathfrak M$. Let $\lambda_1,\lambda_2,\ldots,\mu_1,\mu_2,\ldots$ be mutually distinct numbers, define $T$ by
\[
T\left(\sum_nx_n\varphi_n+\sum_ny_n\psi_n\right)
=\sum_n\lambda_nx_n\varphi_n+\sum_n\mu_ny_n\psi_n,
\]
and let $\mathfrak T$ be its corresponding quantity.

As we know from \hyperref[sec:III.3]{III.3}, a measurement of $\mathfrak T$ produces one of the states $\varphi_1,\varphi_2,\ldots,\psi_1,\psi_2,\ldots$. If it produces a $\varphi_n$, which is recognized by the result being a $\lambda_n$, the values of $\mathfrak R$ and $\mathfrak S$ are also known. By assumption $\mathfrak R,\mathfrak S$ have sharp values in $\varphi_n$, and one can predict with certainty that an immediately following measurement of $\mathfrak R$ or $\mathfrak S$ will find the respective value. If a $\psi_n$ is produced, however, nothing of the kind is known: $\psi_n$ does not lie in $\mathfrak M$, so $\mathfrak R,\mathfrak S$ are not sharp in $\psi_n$. As we know, the probability of finding $\psi_n$ is $(P_{[\psi_n]}\varphi,\varphi)$, and the probability of finding any $\psi_n$ ($n=1,2,\ldots$) is
\begin{align*}
\sum_n(P_{[\psi_n]}\varphi,\varphi)
&=(P_{\mathfrak R_\infty-\mathfrak M}\varphi,\varphi)
=\|P_{\mathfrak R_\infty-\mathfrak M}\varphi\|^2\\
&=\|\varphi-P_{\mathfrak M}\varphi\|^2.
\end{align*}
It is therefore zero, that is, $\mathfrak R,\mathfrak S$ are certainly measured simultaneously, when $\varphi=P_{\mathfrak M}\varphi$, or when $\varphi$ belongs to $\mathfrak M$.\origNote{129}

Since we are now interested in quantities that are not simultaneously measurable, we shall assume the extreme case $\mathfrak M=(0)$. Thus we assume that $\mathfrak R,\mathfrak S$ are not simultaneously measurable in any state at all, or that no common eigenfunction of $R,S$ exists.

If $R,S$ have resolutions of the identity $E(\lambda),F(\lambda)$, then, as we know from \hyperref[sec:III.1]{III.1}, their expectation values in the state $\varphi$ are
\[
\varrho=(R\varphi,\varphi),\qquad\sigma=(S\varphi,\varphi),
\]
and their dispersions, that is, the expectation values of $(\mathfrak R-\varrho)^2$, $(\mathfrak S-\sigma)^2$ (cf.\ the discussion of absolutely exact measurement in \hyperref[sec:III.3]{III.3}), are
\begin{align*}
\varepsilon^2&=((R-\varrho\cdot1)^2\varphi,\varphi)
=\|(R-\varrho\cdot1)\varphi\|^2=\|R\varphi-\varrho\varphi\|^2,\\
\eta^2&=((S-\sigma\cdot1)^2\varphi,\varphi)
=\|(S-\sigma\cdot1)\varphi\|^2=\|S\varphi-\sigma\varphi\|^2.
\end{align*}
A familiar transformation gives\origNote{130}
\[
\varepsilon^2=\|R\varphi\|^2-(R\varphi,\varphi)^2,\qquad
\eta^2=\|S\varphi\|^2-(S\varphi,\varphi)^2.
\]
Since $\|\varphi\|=1$, Schwarz's inequality, \hyperref[thm:II.1]{Theorem~1 of II.1}, already shows that the left-hand sides are $\ge0$. The question now arises: since we can never have $\varepsilon=0$ and $\eta=0$ together, whereas $\varepsilon$ by itself can be made arbitrarily small,
% Source PDF page 130 / original page 123
\hypertarget{origpage-123}{}%
and likewise $\eta$ (for $\mathfrak R,\mathfrak S$ are separately measurable with arbitrary, perhaps even absolute, accuracy), there must be relations between $\varepsilon,\eta$ preventing them from becoming small simultaneously. What are these relations?

Their existence was discovered by Heisenberg.\origNote{131} They are of great importance for understanding the uncertainties that quantum mechanics introduces into the description of nature, and are therefore called uncertainty relations. We shall first derive the most important such relation mathematically, and then return to its fundamental significance and its relation to experience.

In matrix theory, operators $P,Q$ with the commutation property
\[
PQ-QP=\frac h{2\pi i}1
\]
played an important role. They were assigned, for example, to a coordinate and its conjugate momentum (cf.\ \hyperref[sec:I.2]{I.2}), or more generally to any two quantities canonically conjugate in classical mechanics (see, for example, the papers cited in \hyperlink{note-2}{Note~2}). More generally, consider two Hermitian operators $P,Q$ with
\[
PQ-QP=a\cdot1.
\]
Since $(PQ-QP)^*=QP-PQ$, we have $(a\cdot1)^*=\bar a\cdot1=-a\cdot1$, so $\bar a=-a$: $a$ is purely imaginary. Of course this operator equality is not valid as an assertion about the domains; $PQ-QP$ need not be everywhere defined. For every $\varphi$ we then have
\begin{align*}
2\operatorname{Im}(P\varphi,Q\varphi)
&=-i[(P\varphi,Q\varphi)-(Q\varphi,P\varphi)]\\
&=-i[(QP\varphi,\varphi)-(PQ\varphi,\varphi)]\\
&=(i\{PQ-QP\}\varphi,\varphi)=ia\cdot\|\varphi\|^2.
\end{align*}
If $a\ne0$, \hyperref[thm:II.1]{Theorem~1 of II.1} gives
\[
\|\varphi\|^2=-\frac{2i}a\operatorname{Im}(P\varphi,Q\varphi)
\le\frac2{|a|}|(P\varphi,Q\varphi)|
\le\frac2{|a|}\|P\varphi\|\cdot\|Q\varphi\|,
\]
and hence, for $\|\varphi\|=1$,
\[
\|P\varphi\|\cdot\|Q\varphi\|\ge\frac{|a|}2.
\]
Since $P-\varrho\cdot1$, $Q-\sigma\cdot1$ also have the commutation property above, we likewise have
\[
\|P\varphi-\varrho\varphi\|\cdot\|Q\varphi-\sigma\varphi\|\ge\frac{|a|}2.
\]
Introducing the means and dispersions,
\begin{alignat*}{2}
\varrho&=(P\varphi,\varphi),&\quad \varepsilon^2&=\|P\varphi-\varrho\varphi\|^2,\\
\sigma&=(Q\varphi,\varphi),&\quad \eta^2&=\|Q\varphi-\sigma\varphi\|^2,
\end{alignat*}
we obtain
\begin{equation}
\varepsilon\eta\ge\frac{|a|}2.\tag{U}\label{eq:III.4.U}
\end{equation}
% Source PDF page 131 / original page 124
\hypertarget{origpage-124}{}%
For equality to hold, it is necessary and sufficient that equality hold in each of the inequalities used in the derivation. With $P'=P-\varrho\cdot1$, $Q'=Q-\sigma\cdot1$, we therefore have
\[
-\frac{i|a|}a\operatorname{Im}(P'\varphi,Q'\varphi)
=|(P'\varphi,Q'\varphi)|=\|P'\varphi\|\cdot\|Q'\varphi\|.
\]
By \hyperref[thm:II.1]{Theorem~1 of II.1}, the latter equality means that $P'\varphi,Q'\varphi$ differ only by a constant factor. Since $\|P'\varphi\|\cdot\|Q'\varphi\|\ge|a|/2>0$ implies $P'\varphi\ne0$, $Q'\varphi\ne0$, we must have $P'\varphi=cQ'\varphi$, $c'\ne0$.\tnote{The source writes $c'\ne0$, although the proportionality factor is denoted by $c$.} The former equality says that $(P'\varphi,Q'\varphi)=c\|Q'\varphi\|^2$ is purely imaginary, and its coefficient of $i$ has the same sign as $-i|a|/a$, which is real; that is, the opposite sign to the coefficient of $i$ in $a$. Thus $c=i\gamma$, where $\gamma$ is real and $\gamma\gtrless0$ for $a/i\lessgtr0$. Consequently
\begin{equation}
P'\varphi=i\gamma Q'\varphi,\qquad
\gamma\text{ real, and }\gamma\lessgtr0\text{ for }ia\lessgtr0.
\tag{Gl.}\label{eq:III.4.Gl}
\end{equation}
The definition of $\varrho,\sigma$ further requires $(P'\varphi,\varphi)=0$, $(Q'\varphi,\varphi)=0$. But \hyperref[eq:III.4.Gl]{(Gl.)} gives $(P'\varphi,\varphi)=i\gamma(Q'\varphi,\varphi)$, with a real quantity on the left and a purely imaginary one on the right, so these equations hold automatically. We also determine $\varepsilon,\eta$. We have
\[
\varepsilon:\eta=\|P'\varphi\|:\|Q'\varphi\|=|c|=|\gamma|,
\qquad\varepsilon\eta=\frac{|a|}2.
\]
Since $\varepsilon,\eta$ are both positive,
\[
\varepsilon=\sqrt{\frac{|a|\cdot|\gamma|}2},\qquad
\eta=\sqrt{\frac{|a|}{2|\gamma|}}.
\]
For the quantum-mechanical case $a=h/(2\pi i)$, \hyperref[eq:III.4.U]{(U)} becomes
\begin{equation}
\varepsilon\eta\ge\frac h{4\pi}.\tag{$U'$}\label{eq:III.4.Uprime}
\end{equation}
We can also discuss \hyperref[eq:III.4.Gl]{(Gl.)} when $P,Q$ are, for example, the operators of Schrödinger's theory: $P=\frac h{2\pi i}\frac\partial{\partial q}\cdots$, $Q=q\cdots$. (Cf.\ \hyperref[sec:I.2]{I.2}; we assume a mechanical system with one degree of freedom, whose sole coordinate is $q$.) Then \hyperref[eq:III.4.Gl]{(Gl.)} reads
\[
\left(\frac h{2\pi i}\frac\partial{\partial q}-\varrho\right)\varphi
=i\gamma(q-\sigma)\varphi;
\]
since $ia=h/(2\pi)>0$, $\gamma>0$. Hence
\[
\frac\partial{\partial q}\varphi
=\left\{-\frac{2\pi}h\gamma q+\frac{2\pi}h\gamma\sigma+\frac{2\pi}h\varrho i\right\}\varphi,
\]
that is,
\begin{align*}
\varphi&=e^{\int^q\{-\frac{2\pi}h\gamma q+\frac{2\pi}h\gamma\sigma+\frac{2\pi}h\varrho i\}\,dq}\\
&=Ce^{-\frac{\pi\gamma}hq^2+\frac{2\pi\gamma}h\sigma q+\frac{2\pi\varrho}hiq}
=C'e^{-\frac{\pi\gamma}h(q-\sigma)^2+\frac{2\pi\varrho}hiq}.
\end{align*}
% Source PDF page 132 / original page 125
\hypertarget{origpage-125}{}%
Since $\gamma>0$, $\|\varphi\|^2=\int_{-\infty}^{+\infty}|\varphi(q)|^2\,dq$ is indeed finite, and $C'$ is determined by $\|\varphi\|=1$:
\begin{align*}
\|\varphi\|^2&=\int_{-\infty}^{+\infty}|\varphi(q)|^2\,dq
=|C'|^2\int_{-\infty}^{+\infty}e^{-\frac{2\pi\gamma}h(q-\sigma)^2}\,dq\\
&=|C'|^2\sqrt{\frac h{2\pi\gamma}}\int_{-\infty}^{+\infty}e^{-x^2}\,dx\\
&=|C'|^2\sqrt{\frac h{2\pi\gamma}}\sqrt\pi
=|C'|^2\sqrt{\frac h{2\gamma}}=1,\\
|C'|&=\left(\frac{2\gamma}h\right)^{1/4}.
\end{align*}
Omitting the physically meaningless factor of absolute value $1$, we therefore have $C'=(2\gamma/h)^{1/4}$, that is,
\[
\varphi=\varphi(q)=\left(\frac{2\gamma}h\right)^{1/4}
e^{-\frac{\pi\gamma}h(q-\sigma)^2+\frac{2\pi\varrho}hiq}.
\]
We already know $\varepsilon,\eta$:
\[
\varepsilon=\sqrt{\frac{h\gamma}{4\pi}},\qquad
\eta=\sqrt{\frac h{4\pi\gamma}}.
\]
Apart from the condition $\varepsilon\eta=h/(4\pi)$, they are arbitrary as $\gamma$ varies from $0$ to $+\infty$. Thus every quadruple $\varrho,\sigma,\varepsilon,\eta$ satisfying $\varepsilon\eta=h/(4\pi)$ is realized by exactly one $\varphi$. These $\varphi$ were first studied by Heisenberg and used to illustrate quantum-mechanical situations. They are especially well suited to that purpose, because they represent the closest possible approximation within quantum mechanics to classical-mechanical conditions, where neither $p$ nor $q$ has dispersion, and because one can assign prescribed values to $\varepsilon$ or $\eta$ at will. See the reference in \hyperlink{note-131}{Note~131}.

The foregoing considerations have covered only one aspect of the uncertainty relations, namely the formal one. A full understanding requires us to consider them from another standpoint as well: that of direct physical experience. Their relation to such experience is simpler and more readily surveyed than that of many of the facts on which quantum mechanics was built, and so the wholly formal derivation above does not do them full justice. An intuitive discussion is all the more necessary because at first glance one might even think they contradict intuition: it is not immediately apparent why position and velocity, that is,
% Source PDF page 133 / original page 126
\hypertarget{origpage-126}{}%
coordinate and momentum, of a material body could not be measured simultaneously with arbitrarily high accuracy if sufficiently fine measuring instruments were available. We must therefore make clear, by an exact analysis of the most refined measurement processes, perhaps practicable only as thought experiments, that this is not so. Rather, the familiar laws of wave optics, electrodynamics, and elementary atomic processes place insuperable obstacles in the way of accurate measurement precisely where the uncertainty relations require them. This can already be recognized by examining these processes in a purely classical way, without quantum theory. This is of fundamental importance: it shows that in the classical domain, whose validity is established independently of the correctness of quantum mechanics---where quantum phenomena do not yet compel a substantial correction of the older view, and which is in fact the only domain directly accessible to our intuition\origNote{132}---we encounter no conflict with the paradoxical-sounding uncertainty relations of quantum mechanics.

What must be shown is this. If $p,q$ are canonically conjugate quantities and a system is in a state in which the value of $p$ can be specified with accuracy $\varepsilon$, that is, after a measurement of $p$ with margin of error $\varepsilon$, then $q$ cannot be known with any greater accuracy than $\eta=\frac h{4\pi}:\varepsilon$. Equivalently, a measurement of $p$ with accuracy $\varepsilon$ must introduce an uncertainty $\eta=\frac h{4\pi}:\varepsilon$ into the value of $q$.

Naturally, these very qualitative considerations cannot recover every detail exactly: instead of $\varepsilon\eta=h/(4\pi)$, we shall be able to show only $\varepsilon\eta\sim h$ for the most accurate possible measurement, meaning that $\varepsilon\eta$ is of the order of magnitude of $h$. As a characteristic example, consider the conjugate pair position (coordinate)--momentum of a particle $\boldsymbol{T}$.\origNote{133}

We first examine the determination of position. This is done by looking at $\boldsymbol{T}$, that is, illuminating it and absorbing the scattered light in the eye. Thus a light quantum $\boldsymbol{L}$ is fired from the light source $\boldsymbol{l}$ towards $\boldsymbol{T}$, deflected by its collision with $\boldsymbol{T}$ from its straight path $\beta\beta_1$ into $\beta\beta_2$, and annihilated by absorption at the end of its path on the screen $Sch$, which replaces the eye or a photographic plate (Fig.~\ref{fig:1}).\tnote{In the reproduced figure, $Sch$ abbreviates the German \emph{Schirm}, ``screen''.} The measurement is made by establishing that $\boldsymbol{L}$ has struck the screen $Sch$ not at $1$, the end of its natural path $\beta\beta_1$, but at $2$, the end of $\beta\beta_2$. To determine the location of the collision, and thus of $\boldsymbol{T}$, from this, the directions of $\beta$ and $\beta_2$ must also be known, that is, the direction of $\boldsymbol{L}$ before and
% Source PDF page 134 / original page 127
\hypertarget{origpage-127}{}%
after the collision. This is achieved by inserting suitable systems of diaphragms $ss$ and $s's'$. Strictly speaking, this does not measure the coordinate of $\boldsymbol{T}$ but decides whether it has a certain value or not: the value corresponding to the intersection of directions $\beta$ and $\beta_2$, which can be placed arbitrarily by suitably arranging the diaphragms. Only the combination of several such decisions, that is, the installation of several diaphragms $s's'$, amounts to a measurement. What, then, is the accuracy of this position measurement?

It has a fundamental limit in the laws of optical imaging. With light of wavelength $\lambda$ it is impossible to form a sharp image of objects smaller than $\lambda$, or even to reduce the blurring enough to obtain a distorted image. Admittedly, we did not require an optical image, since the mere fact that $\boldsymbol{L}$ is deflected suffices to establish the position of $\boldsymbol{T}$. But the diaphragms $ss$ and $s's'$ cannot be narrower than $\lambda$, for otherwise $\boldsymbol{L}$ cannot pass through them undisturbed: a set of interference fringes appears, so the direction $\beta$ or $\beta_2$ of the light ray can no longer be inferred from the direction of the successive diaphragms $ss$ or $s's'$. Thus with this projectile $\boldsymbol{L}$ it is in no case possible to aim and hit with greater accuracy than $\lambda$.

\begin{figure}[tbp]
\centering
\includegraphics[width=.35\linewidth,keepaspectratio]{figures/figure-1.png}
\caption{}\label{fig:1}
\end{figure}

The wavelength $\lambda$ is therefore also the measure of the coordinate-measurement error: $\lambda\sim\varepsilon$. Further characteristics of $\boldsymbol{L}$ are its frequency $\nu$, energy $\overline E$, and momentum $\bar p$; the familiar relations are
\[
\nu=\frac c\lambda,\qquad
\overline E=h\nu=\frac{hc}\lambda,\qquad
\bar p=\frac{\overline E}c=\frac{h\nu}c=\frac h\lambda,
\]
where $c$ is the velocity of light.\origNote{134} Thus $\bar p\sim h/\varepsilon$. During the collision between $\boldsymbol{L}$ and $\boldsymbol{T}$, whose details are not exactly known, a momentum transfer plainly of the order of $\bar p$, or $h/\varepsilon$, takes place. This produces in $\boldsymbol{T}$ a momentum uncertainty $\eta\sim h/\varepsilon$.

This would prove $\varepsilon\eta\sim h$, had one circumstance not been overlooked. In fact, more is known about the collision process than this suggests:
% Source PDF page 135 / original page 128
\hypertarget{origpage-128}{}%
we know the directions of motion of $\boldsymbol{L}$ before and after it, $\beta$ and $\beta_2$, and thus its momenta; these determine how much momentum is transferred to $\boldsymbol{T}$. Consequently, $\bar p$ is not a measure of $\eta$; the measure is instead the possible uncertainty in the ray directions $\beta,\beta_2$. To determine more precisely the relation between ``aiming'' at the small object $\boldsymbol{T}$ and the associated uncertainty of direction, it is useful to employ a more perfect aiming device than the slit $ss$, namely a lens. We must therefore take account of the familiar theory of the microscope, which says the following: to illuminate a patch of linear extent $\varepsilon$, that is, to hit $\boldsymbol{T}$ with $\boldsymbol{L}$ to accuracy $\varepsilon$, a wavelength $\lambda$ and lens aperture $\varphi$ are required satisfying $\lambda/(2\sin(\varphi/2))\sim\varepsilon$ (Fig.~\ref{fig:2}).\origNote{135}

The uncertainty in the $tt$-component of the momentum of $\boldsymbol{L}$ arises because its direction lies between $-\varphi/2$ and $+\varphi/2$ but is otherwise unknown. It therefore amounts to
\[
2\sin\frac\varphi2\cdot\bar p=\frac\lambda\varepsilon\cdot\frac h\lambda=\frac h\varepsilon.
\]
This is the correct measure of $\eta$, so again $\eta\sim h/\varepsilon$, or $\varepsilon\eta\sim h$.

This example shows the mechanism of the uncertainty principles quite clearly. To aim accurately we need large eyes, that is, a large aperture $\varphi$, and light of very short wavelength. This means very indeterminate and large momenta of the light quantum, which cause collisions with the observed object $\boldsymbol{T}$ whose details cannot be followed (Compton effects), thereby ``smearing out'' its momentum.

\begin{figure}[tbp]
\centering
\includegraphics[width=.35\linewidth,keepaspectratio]{figures/figure-2.png}
\caption{}\label{fig:2}
\end{figure}

Let us also consider the reverse measurement process: measurement of velocity, or momentum. The natural way to measure the velocity of $\boldsymbol{T}$ is to measure its position at two different times, say $0$ and $t$, and divide the change in coordinate by $t$. The velocity must, however, be constant throughout the interval $0,t$. If it changes, the change is a measure of the deviation between the mean velocity just calculated and the actual velocity, say at time $t$, and hence of the uncertainty in the measurement. The same applies to momentum. If the coordinate measurements are made to accuracy $\varepsilon$, this does not affect the accuracy with which the mean momentum is measured, since $t$ can be chosen arbitrarily large. It does, however, cause momentum changes $\sim h/\varepsilon$, and hence an
% Source PDF page 136 / original page 129
\hypertarget{origpage-129}{}%
uncertainty $\eta\sim h/\varepsilon$ in the final momentum. Thus $\varepsilon\eta\sim h$ results in any event. Only momentum measurements unconnected with any position measurement could therefore yield something new. Such measurements are quite possible and are frequently used in astronomy, for example; they rest on the Doppler effect. Let us turn to this.

The familiar Doppler effect is as follows. A body $\boldsymbol{T}$ moving with velocity $v$ emits light of frequency $\nu_0$, as measured at the moving body. A stationary observer perceives a different frequency $\nu$, given by
\[
\frac{\nu-\nu_0}{\nu_0}=\frac vc\cos\theta.
\]
Here $\theta$ is the angle between the directions of motion and emission. This formula is nonrelativistic, valid only for small $v/c$, but that could easily be corrected. The velocity can therefore be determined if $\nu$ is observed and $\nu_0$ is known, perhaps because it is a particular spectral line of a known element. More precisely, one measures the velocity component in the direction of observation, or light emission,\tnote{The source prints $\nu$ in the denominator of the following velocity formula; the preceding Doppler relation and the subsequent momentum formula instead give $\nu_0$. The printed denominator is retained.}
\[
v\cos\theta=\frac{c(\nu-\nu_0)}\nu,
\]
or the corresponding momentum component $p'=p\cos\theta=mc(\nu-\nu_0)/\nu_0$, where $m$ is the mass of $\boldsymbol{T}$. The dispersion $\eta$ of $p'$ is plainly related to the dispersion $\Delta\nu$ of $\nu$, so $\eta\sim mc\Delta\nu/\nu_0\sim mc\Delta\nu/\nu$. The momentum of $\boldsymbol{T}$ is admittedly changed by its emission of a light quantum of frequency $\nu$ and thus momentum $\bar p=h\nu/c$, but under normal circumstances the uncertainty of this quantity, $h\Delta\nu/c$, can be neglected compared with $mc\Delta\nu/\nu$.\origNote{136}

The frequency $\nu$ is measured by some interference phenomenon, which naturally gives an absolutely sharp $\nu$ only for a purely monochromatic light-wave train. Such a train has the form
\[
a\sin\left(2\pi\left(\frac q\lambda-\nu t\right)+\alpha\right),
\]
where $q$ is the coordinate, $t$ time, $a$ amplitude, and $\alpha$ phase; the expression represents some component of the electric or magnetic field strength. It therefore extends infinitely in both space and time. To avoid this we must replace the expression, which because $\lambda=c/\nu$ can also be written $a\sin(2\pi\nu(q/c-t)+\alpha)$, by another function $F(q/c-t)$ that is nonzero only in a finite interval of its argument. If the light wave has this form, it must, as is well known, be Fourier-analyzed:
\[
F(x)=\int_0^{+\infty}a_\nu\sin(2\pi\nu x+\alpha_\nu)\,d\nu,
\]
% Source PDF page 137 / original page 130
\hypertarget{origpage-130}{}%
and the interference pattern displays all frequencies $\nu$ for which $a_\nu\ne0$, with relative intensity $a_\nu^2\,d\nu$ in the frequency interval $\nu,\nu+d\nu$. The dispersion of $\nu$, namely $\Delta\nu$, is to be calculated from this distribution.

If our wave train has length $\tau$ in $x$, and therefore length $\tau$ in $t$ or $c\tau$ in $q$, this frequency dispersion is seen to be $\sim1/\tau$.\origNote{137} In this measurement method an uncertainty of position arises because each light emission gives $\boldsymbol{T}$ a recoil $h\nu/c$ in the direction of observation, that is, a velocity change $h\nu/(mc)$. Since the emission process takes a time $\tau$, however, we cannot locate the instant of this velocity change more accurately than $\tau$. Hence there is a positional uncertainty $\varepsilon\sim(h\nu/(mc))\tau$. Thus
\[
\varepsilon\sim\frac{h\nu}{mc}\tau,\qquad
\eta\sim\frac{mc\Delta\nu}\nu=\frac{mc}\nu\frac1\tau,\qquad
\varepsilon\eta\sim h.
\]
Again, therefore, $\varepsilon\eta\sim h$.

If $\boldsymbol{T}$, instead of emitting its own light as assumed here, scatters light from elsewhere, that is, is illuminated, the calculation is entirely analogous.

\section{Projection Operators as Propositions}\label{sec:III.5}
As in \hyperref[sec:III.1]{III.1}, consider a physical system $\boldsymbol{S}$ with $k$ degrees of freedom, whose configuration space is therefore described by the $k$ coordinates $q_1,\ldots,q_k$ (cf.\ also \hyperref[sec:I.2]{I.2}). In the classical-mechanical view, all physical quantities $\mathfrak R$ that can be formed in the system $\boldsymbol{S}$ are functions of $q_1,\ldots,q_k$ and the conjugate momenta $p_1,\ldots,p_k$:
\[
\mathfrak R=\mathfrak R(q_1,\ldots,q_k,p_1,\ldots,p_k).
\]
For example, energy is the Hamiltonian function $H(q_1,\ldots,q_k,p_1,\ldots,p_k)$. In quantum mechanics, on the other hand, as already indicated in \hyperref[sec:III.1]{III.1}, the quantities $\mathfrak R$ are uniquely associated with the hypermaximal Hermitian operators $R$. In particular, $q_1,\ldots,q_k$ correspond to the operators $Q_1=q_1\cdots,\ldots,Q_k=q_k\cdots$, and $p_1,\ldots,p_k$ to
\[
P_1=\frac h{2\pi i}\frac\partial{\partial q_1}\cdots,\quad\ldots,\quad
P_k=\frac h{2\pi i}\frac\partial{\partial q_k}\cdots.
\]
The example of the Hamiltonian function in \hyperref[sec:I.2]{I.2} already made clear that, because $Q_l,P_l$ do not commute, it is in general impossible to define $R=\mathfrak R(Q_1,\ldots,Q_k,P_1,\ldots,P_k)$. Although we cannot state anything definitive about the relation between the $\mathfrak R(q_1,\ldots,q_k,p_1,\ldots,p_k)$ and the $R$, we did establish the following special facts in \hyperref[sec:III.1]{III.1} and \hyperref[sec:III.3]{III.3}:

\leavevmode\phantomsection\label{post:III.5.L}\textbf{L.} If the simultaneously observable quantities $\mathfrak R,\mathfrak S$ correspond to the operators $R,S$, then the quantity $a\mathfrak R+b\mathfrak S$, for real $a,b$, corresponds to the operator $aR+bS$.

% Source PDF page 138 / original page 131
\hypertarget{origpage-131}{}%
\leavevmode\phantomsection\label{post:III.5.F}\textbf{F.} If the quantity $\mathfrak R$ corresponds to the operator $R$, then the quantity $F(\mathfrak R)$, where $F(\lambda)$ is any real function, corresponds to the operator $F(R)$.

\hyperref[post:III.5.L]{L} and \hyperref[post:III.5.F]{F} admit a certain further generalization. The generalization of \hyperref[post:III.5.F]{F} follows necessarily and reads:

\leavevmode\phantomsection\label{post:III.5.Fstar}\textbf{$F^*$.} If the simultaneously observable quantities $\mathfrak R,\mathfrak S,\ldots$ correspond to the operators $R,S,\ldots$, which therefore commute and whose number is finite, then the quantity $F(\mathfrak R,\mathfrak S,\ldots)$ corresponds to the operator $F(R,S,\ldots)$.

Here we shall assume $F(\lambda,\mu,\ldots)$ to be a real polynomial in $\lambda,\mu,\ldots$, so that $F(R,S,\ldots)$ has a clear meaning, since $R,S,\ldots$ commute. The assertion \hyperref[post:III.5.Fstar]{$F^*$} could also be justified for arbitrary $F(\lambda,\mu,\ldots)$; for the definition of general $F(R,S,\ldots)$ see the reference in \hyperlink{note-94}{Note~94}. Every polynomial is obtained by repeating the three operations $a\lambda$, $\lambda+\mu$, $\lambda\mu$, so it suffices to consider these. Since $\lambda\mu=\frac14((\lambda+\mu)^2-(\lambda-\mu)^2)$, that is,
\[
\frac14\cdot(\lambda+\mu)^2+\left(-\frac14\right)\cdot(\lambda+(-1)\cdot\mu)^2,
\]
we can replace these three operations by $a\lambda$, $\lambda+\mu$, $\lambda^2$. The first two fall under \hyperref[post:III.5.L]{L} and the last under \hyperref[post:III.5.F]{F}. This proves \hyperref[post:III.5.Fstar]{$F^*$}.

\hyperref[post:III.5.L]{L}, on the other hand, is extended in quantum mechanics even to the case where $\mathfrak R,\mathfrak S$ are not simultaneously measurable. We shall discuss this question only later, in \hyperref[sec:IV.1]{IV.1}; for now we merely observe that it is by no means clear what $a\mathfrak R+b\mathfrak S$ is to mean for quantities $\mathfrak R,\mathfrak S$ that are not simultaneously measurable.

Besides the physical quantities $\mathfrak R$, there is something else that physics studies: the properties of the states of the system $\boldsymbol{S}$. It is a property, for example, that a certain quantity $\mathfrak R$ takes the value $\lambda$; or that the value of $\mathfrak R$ is positive; or that the values of two simultaneously observable quantities $\mathfrak R,\mathfrak S$ are respectively $\lambda,\mu$; or that the sum of the squares of these values is $>1$; and so on. We denoted quantities by $\mathfrak R,\mathfrak S,\ldots$; we shall denote properties by $\mathfrak E,\mathfrak F,\ldots$. Quantities corresponded to the hypermaximal Hermitian operators $R,S,\ldots$, as just discussed; what corresponds to properties?

We can associate with every property $\mathfrak E$ a quantity defined as follows: every measurement that decides whether $\mathfrak E$ is present or absent is to count as a measurement of this quantity; its value is $1$ if $\mathfrak E$ is present and $0$ otherwise. We shall denote this quantity associated with $\mathfrak E$ by $\mathfrak E$ as well.

It takes only the values $0$ and $1$. Conversely, every quantity $\mathfrak R$ capable of only these two values is associated with a property $\mathfrak E$, plainly the property ``the value of $\mathfrak R$ is $\ne0$''. Thus the
% Source PDF page 139 / original page 132
\hypertarget{origpage-132}{}%
quantities $\mathfrak E$ associated with properties are characterized by the condition above.

That $\mathfrak E$ takes only the values $0,1$ can also be formulated by saying that it makes the polynomial $F(\lambda)=\lambda-\lambda^2$ vanish identically. If $\mathfrak E$ has the operator $E$, then $F(\mathfrak E)$ has the operator $F(E)=E-E^2$; thus the condition is $E-E^2=0$, or $E=E^2$. In other words, the operator $E$ of $\mathfrak E$ is a projection operator.

Properties $\mathfrak E$ are therefore associated with projection operators $E$, through the corresponding quantities $\mathfrak E$ just defined. If, besides $E=P_{\mathfrak M}$, we also consider the corresponding closed linear manifold $\mathfrak M$, we may equally say that properties are associated with closed linear manifolds $\mathfrak M$.

We shall now examine more closely the calculus of the corresponding $\mathfrak E,E,\mathfrak M$.

By definition, deciding whether the property $\mathfrak E$ is present in the state $\varphi$ is equivalent to measuring the quantity $\mathfrak E$ and determining whether its value is $1$ or $0$. The probability of the first case, that $\mathfrak E$ is present, therefore equals the expectation value of $\mathfrak E$, namely
\[
(E\varphi,\varphi)=\|E\varphi\|^2=\|P_{\mathfrak M}\varphi\|^2,
\]
and the probability of the second case, that $\mathfrak E$ is absent, equals the expectation value of $1-\mathfrak E$, namely
\[
((1-E)\varphi,\varphi)=\|(1-E)\varphi\|^2=\|\varphi-P_{\mathfrak M}\varphi\|^2.
\]
Their sum is of course $(\varphi,\varphi)$, that is, $1$. Hence $\mathfrak E$ is certainly present or certainly absent when the second or first probability, respectively, is zero: when $P_{\mathfrak M}\varphi=\varphi$ or $P_{\mathfrak M}\varphi=0$. This means that $\varphi$ belongs to $\mathfrak M$ or is orthogonal to $\mathfrak M$, respectively; equivalently, that $\varphi$ belongs to $\mathfrak M$ or to $\mathfrak R_\infty-\mathfrak M$.

Thus $\mathfrak M$ can also be defined as the set of all $\varphi$ that certainly possess the property $\mathfrak E$. Strictly speaking, this gives only those elements of $\mathfrak M$ for which $\|\varphi\|=1$; $\mathfrak M$ itself is obtained by multiplying these elements by positive constants and adjoining zero.

If we call the property opposite to $\mathfrak E$, its denial, ``non-$\mathfrak E$'', it follows at once that when $E,\mathfrak M$ belong to $\mathfrak E$, then $1-E,\mathfrak R_\infty-\mathfrak M$ belong to non-$\mathfrak E$.

For properties as for quantities, the question of simultaneous measurability, or more properly simultaneous decidability, arises. Clearly $\mathfrak E,\mathfrak F$ are simultaneously decidable if and only if the corresponding quantities $\mathfrak E,\mathfrak F$ are simultaneously measurable. Whether the accuracy is arbitrarily high or absolute is immaterial, since they can take only the values $0,1$. Thus the condition is that $E,F$ commute. The same applies to several properties $\mathfrak E,\mathfrak F,\mathfrak G,\ldots$.

% Source PDF page 140 / original page 133
\hypertarget{origpage-133}{}%
From simultaneously decidable properties $\mathfrak E,\mathfrak F$ we can form the further properties ``$\mathfrak E$ and $\mathfrak F$'' and ``$\mathfrak E$ or $\mathfrak F$''. The quantity belonging to ``$\mathfrak E$ and $\mathfrak F$'' is $1$ when the quantities belonging to $\mathfrak E$ and $\mathfrak F$ are both $1$, and zero when either is zero; it is therefore their product. By \hyperref[post:III.5.Fstar]{$F^*$}, its operator is the product of the operators of $\mathfrak E$ and $\mathfrak F$, namely $EF$. By \hyperref[thm:II.14]{Theorem~14 of II.4}, the corresponding closed linear manifold is the common part $\mathfrak P$ of $\mathfrak M,\mathfrak N$.

The property ``$\mathfrak E$ or $\mathfrak F$'', on the other hand, can be written
\[
\text{``non-((non-}\mathfrak E\text{) and (non-}\mathfrak F\text{))''},
\]
so its operator is
\[
1-(1-E)(1-F)=E+F-EF,
\]
which by its construction is of course a projection operator. Since $F-EF$ is a projection operator, the linear manifold belonging to $E+F-EF$ is $\mathfrak M+(\mathfrak N-\mathfrak P)$ (\hyperref[thm:II.14]{Theorem~14, II.4}). It is part of $\{\mathfrak M,\mathfrak N\}$ and plainly contains $\mathfrak M$, hence by symmetry also $\mathfrak N$, and thus $\{\mathfrak M,\mathfrak N\}$. It therefore equals $\{\mathfrak M,\mathfrak N\}$, which, being consequently closed, equals $[\mathfrak M,\mathfrak N]$.

If $\mathfrak E$ is a property that is always present, that is, a vacuous one, its quantity is identically $1$, so $E=1$, $\mathfrak M=\mathfrak R_\infty$. If $\mathfrak E$ is never present, that is, is impossible, its quantity is identically zero, so $E=0$, $\mathfrak M=(0)$. If two properties $\mathfrak E,\mathfrak F$ are incompatible, they must in any event be simultaneously decidable and ``$\mathfrak E$ and $\mathfrak F$'' must be impossible: $E,F$ must commute and $EF=0$. But $EF=0$ itself implies commutativity (\hyperref[thm:II.14]{Theorem~14, II.4}), so it alone characterizes incompatibility. If $E,F$ are assumed to commute, $EF=0$ says only that the common part of $\mathfrak M,\mathfrak N$ consists solely of zero; this condition alone, however, does not imply commutativity of $E,F$. In general, $EF=0$ means for $\mathfrak M,\mathfrak N$ that the whole of $\mathfrak M$ is orthogonal to the whole of $\mathfrak N$ (\hyperref[thm:II.14]{Theorem~14, II.4}).

If $\mathfrak R$ is a quantity with operator $R$ and resolution of the identity $E(\lambda)$, the operator of the property ``$\mathfrak R$ lies in the interval $I=\{\lambda',\mu'\}$'', $\lambda'\le\mu'$, is plainly $E(\mu')-E(\lambda')$. This follows, for example, because the probability of the property is $((E(\mu')-E(\lambda'))\varphi,\varphi)$ (cf.\ \hyperref[post:III.1.W]{W} in \hyperref[sec:III.1]{III.1}); or because the quantity belonging to the property is $\mathfrak E=F(\mathfrak R)$, with
\[
F(\lambda)=\begin{cases}1,&\lambda'<\lambda\le\mu',\\0,&\text{otherwise},\end{cases}
\]
and $F(R)=E(\mu')-E(\lambda')$ (cf.\ \hyperref[sec:II.8]{II.8} or \hyperref[sec:III.1]{III.1}). We also called this operator $E(I)$ in \hyperref[sec:III.1]{III.1}.

To collect the results, we have established the following relations between properties $\mathfrak E$, their projection operators $E$, and the corresponding closed linear manifolds $\mathfrak M$:

% Source PDF page 141 / original page 134
\hypertarget{origpage-134}{}%
\leavevmode\phantomsection\label{post:III.5.alpha}$\alpha)$ In the state $\varphi$, the property $\mathfrak E$ is present or absent with the respective probabilities
\[
(E\varphi,\varphi)=\|E\varphi\|^2=\|P_{\mathfrak M}\varphi\|^2
\]
and
\[
((1-E)\varphi,\varphi)=\|(1-E)\varphi\|^2=\|\varphi-P_{\mathfrak M}\varphi\|^2.
\]

\leavevmode\phantomsection\label{post:III.5.beta}$\beta)$ The property $\mathfrak E$ is certainly present or absent for the $\varphi$ in $\mathfrak M$ or $\mathfrak R_\infty-\mathfrak M$, respectively, and only for these.

\leavevmode\phantomsection\label{post:III.5.gamma}$\gamma)$ Simultaneous decidability of several properties $\mathfrak E,\mathfrak F,\ldots$ is characterized by commutativity of their operators $E,F,\ldots$.

\leavevmode\phantomsection\label{post:III.5.delta}$\delta)$ If $E,\mathfrak M$ belong to $\mathfrak E$, then $1-E,\mathfrak R_\infty-\mathfrak M$ belong to non-$\mathfrak E$.

\leavevmode\phantomsection\label{post:III.5.epsilon}$\varepsilon)$ If $E,\mathfrak M$ belong to $\mathfrak E$ and $F,\mathfrak N$ to $\mathfrak F$, and $\mathfrak E,\mathfrak F$ are simultaneously decidable, then ``$\mathfrak E$ and $\mathfrak F$'' corresponds to $EF$ and the common part of $\mathfrak M,\mathfrak N$, while ``$\mathfrak E$ or $\mathfrak F$'' corresponds to $E+F-EF$ and $\{\mathfrak M,\mathfrak N\}$, which equals $[\mathfrak M,\mathfrak N]$.

\leavevmode\phantomsection\label{post:III.5.eta}$\eta)$ The property $\mathfrak E$ always holds if $E=1$, or equivalently $\mathfrak M=\mathfrak R_\infty$; it never holds if $E=0$, or equivalently $\mathfrak M=(0)$.

\leavevmode\phantomsection\label{post:III.5.theta}$\vartheta)$ The properties $\mathfrak E,\mathfrak F$ are incompatible if $EF=0$, or equivalently if the whole of $\mathfrak M$ is orthogonal to the whole of $\mathfrak N$.

\leavevmode\phantomsection\label{post:III.5.zeta}$\zeta)$ Let $\mathfrak R$ be a quantity, $R$ its operator, and $I$ an interval. Let $E(\lambda)$ be the resolution of the identity belonging to $R$, $I=\{\lambda',\mu'\}$ with $\lambda'\le\mu'$, and $E(I)=E(\mu')-E(\lambda')$ (cf.\ \hyperref[sec:III.1]{III.1}). Then the property ``$\mathfrak R$ lies in $I$'' has the operator $E(I)$.

From \hyperref[post:III.5.alpha]{$\alpha)$}--\hyperref[post:III.5.zeta]{$\zeta)$} we can derive the earlier probability assertions \hyperref[post:III.1.W]{W}, \hyperref[post:III.1.E1]{$E_1$}, \hyperref[post:III.1.E2]{$E_2$}, and the assertions of \hyperref[sec:III.3]{III.3} about simultaneous measurability. That the latter are equivalent to \hyperref[post:III.5.gamma]{$\gamma)$} is clear; \hyperref[post:III.1.W]{W} follows from \hyperref[post:III.5.alpha]{$\alpha)$}, \hyperref[post:III.5.epsilon]{$\varepsilon)$}, \hyperref[post:III.5.zeta]{$\zeta)$}, and implies \hyperref[post:III.1.E1]{$E_1$}, \hyperref[post:III.1.E2]{$E_2$}.

As we see, the relation between the properties of a physical system and projection operators enables a kind of logical calculus with the latter. Compared with ordinary logic, however, this calculus is enriched by the concept of ``simultaneous decidability'', characteristic of quantum mechanics.

The calculus of propositions based on projection operators has a further advantage over the calculus of quantities based on the totality of all hypermaximal Hermitian operators: the concept of ``simultaneous decidability'' refines that of ``simultaneous measurability''. For example, let $\mathfrak R,\mathfrak S$ have operators $R,S$ with respective resolutions of the identity $E(\lambda),F(\mu)$, and let $I=\{\lambda',\lambda''\}$, $J=\{\mu',\mu''\}$. For the questions ``Does $\mathfrak R$ lie in $I$?'' and ``Does $\mathfrak S$ lie in $J$?'' to be simultaneously decidable, \hyperref[post:III.5.gamma]{$\gamma)$} and \hyperref[post:III.5.zeta]{$\zeta)$} require only that the operators $E(I)=E(\lambda'')-E(\lambda')$ and $F(J)=F(\mu'')-F(\mu')$ commute. Simultaneous measurability of $\mathfrak R,\mathfrak S$, by contrast, requires commutativity of every $E(\lambda)$ with every $F(\mu)$.


% Consolidated source: chapter-III-6.tex
% Source PDF page 142 / original page 135.
\hypertarget{origpage-135}{}%
\section{Theory of Light}\label{sec:III.6}
We have recovered all the statistical assertions of quantum mechanics cited in \hyperref[sec:I.2]{I.2}, and have even substantially generalized and systematically developed them, with a single exception. We still lack \textsc{Heisenberg}'s expression for the transition probability from one stationary state of a quantized system to another, although this very expression played an important role in the origin of quantum mechanics (cf.\ the discussion in \hyperref[sec:I.2]{I.2}). Following \textsc{Dirac},\origNote{138} we shall show how these transition probabilities can be derived from the ordinary statistical assertions of quantum mechanics, that is, from the theory just presented. This is all the more important because it will give us a deeper insight into how transitions between stationary states occur and into the connections underlying the \textsc{Einstein--Bohr} energy--frequency relations. Indeed, the theory of light due to \textsc{Dirac} which we are about to present is probably one of the most beautiful achievements in quantum mechanics.

Let $\boldsymbol{S}$ be a system (say, a quantized atom) whose energy corresponds to the Hermitian operator $\mathsf{H}_0$. For brevity, denote by $\xi$ the coordinates describing the configuration space of $\boldsymbol{S}$ (if, for example, $\boldsymbol{S}$ consists of $l$ particles, these are $3l$ Cartesian coordinates: $x_1=q_1,y_1=q_2,z_1=q_3,\ldots,x_l=q_{3l-2},y_l=q_{3l-1},z_l=q_{3l}$). For simplicity, let $\mathsf{H}_0$ have a pure point spectrum: eigenvalues $W_1,W_2,\ldots$ and eigenfunctions $\varphi_1(\xi),\varphi_2(\xi),\ldots$ (some of the $W_n$ may coincide). An arbitrary state of $\boldsymbol{S}$, that is, a wave function $\varphi(\xi)$, evolves according to the time-dependent \textsc{Schrödinger} differential equation (cf.\ \hyperref[sec:III.2]{III.2}):
\[
 \frac{h}{2\pi i}\frac{\partial}{\partial t}\varphi_t(\xi)=-\mathsf{H}_0\varphi_t(\xi).
\]
Thus, if at $t=t_0$ we have $\varphi_t(\xi)=\varphi(\xi)=\sum_{k=1}^{\infty}a_k\varphi_k(\xi)$, then in general
\[
 \varphi_t(\xi)=\sum_{k=1}^{\infty}a_ke^{-\frac{2\pi i}{h}W_k(t-t_0)}\varphi_k(\xi).
\]
A state $\varphi(\xi)=\varphi_k(\xi)$ therefore passes into $e^{-\frac{2\pi i}{h}W_k(t-t_0)}\varphi_k(\xi)$, that is, since the factor $e^{-\frac{2\pi i}{h}W_k(t-t_0)}$ is inessential, into itself: the $\varphi_k(\xi)$ are precisely stationary. Thus at first we find no transitions at all from one $\varphi_k(\xi)$ to another. How, then, can such transitions be spoken of? The explanation is simple: we have disregarded the agency which causes these transitions---light. Even in the original \textsc{Bohr} theory, the stationary quantum orbits decay only through the emission of light (cf.\ the reference in \hyperlink{note-5}{Note~5}); if this is disregarded, as in the formulation just given, it is very reasonable that absolute and eternal stability should result.
% Source PDF page 143 / original page 136.
\hypertarget{origpage-136}{}%
We must therefore enlarge the system under investigation by including the light which $\boldsymbol{S}$ may emit, and hence all the light which can possibly interact with $\boldsymbol{S}$. Let the system formed by the light (that is, the electromagnetic field of classical theory, excluding the static fields due to electron and nuclear charges) be denoted by $\boldsymbol{L}$. We have to investigate $\boldsymbol{S}+\boldsymbol{L}$.

The following must therefore be found:

\leavevmode\phantomsection\label{problem:III.6.1}\textbf{1.} A formulation of the quantum-mechanical description of $\boldsymbol{L}$, that is, the configuration space of $\boldsymbol{L}$ must be specified.

\leavevmode\phantomsection\label{problem:III.6.2}\textbf{2.} The energy operator of $\boldsymbol{S}+\boldsymbol{L}$. This divides into three parts.

\leavevmode\phantomsection\label{problem:III.6.2alpha}$\alpha)$ The energy of $\boldsymbol{S}$ present independently of $\boldsymbol{L}$, that is, the unperturbed $\boldsymbol{S}$-energy: for the system $\boldsymbol{S}$ this was the operator $\mathsf{H}_0$.

\leavevmode\phantomsection\label{problem:III.6.2beta}$\beta)$ The energy of $\boldsymbol{L}$ present independently of $\boldsymbol{S}$, that is, the unperturbed $\boldsymbol{L}$-energy. Denote its operator by $\mathsf{H}_l$.

\leavevmode\phantomsection\label{problem:III.6.2gamma}$\gamma)$ The remaining energy, that is, the interaction energy of $\boldsymbol{S}$ and $\boldsymbol{L}$. Denote its operator by $\mathsf{H}_w$.

As we see, these are questions which, according to the basic principles of quantum mechanics, must first be answered classically so that the result can then be rewritten in terms of operators (cf.\ \hyperref[sec:I.2]{I.2}). We therefore initially adopt a purely classical standpoint regarding the nature of light: in the sense of the electromagnetic theory of light, we regard it as an oscillatory state of the electromagnetic field.\origNote{139}

To avoid unnecessary complications (loss of light in infinite space, and so on), we imagine $\boldsymbol{S}$ and $\boldsymbol{L}$ enclosed in a very large cavity $\boldsymbol H$ of volume $\mathcal V$, with, say, reflecting walls. As is well known, the state of the electromagnetic field in $\boldsymbol H$ is described by the electric and magnetic field strengths $\mathfrak E=\{\mathfrak E_x,\mathfrak E_y,\mathfrak E_z\}$ and $\mathfrak H=\{\mathfrak H_x,\mathfrak H_y,\mathfrak H_z\}$. All the quantities $\mathfrak E_x,\ldots,\mathfrak H_z$ are functions of the Cartesian coordinates $x,y,z$ of an arbitrary point of $\boldsymbol H$ and of the time $t$. We emphasize that we shall now repeatedly consider real spatial vectors $\mathfrak a=\{\mathfrak a_x,\mathfrak a_y,\mathfrak a_z\}$, $\mathfrak b=\{\mathfrak b_x,\mathfrak b_y,\mathfrak b_z\}$, and so on (for example, $\mathfrak E,\mathfrak H$), and operations on them such as the inner product
\[
 [\mathfrak a,\mathfrak b]=\mathfrak a_x\mathfrak b_x+\mathfrak a_y\mathfrak b_y+\mathfrak a_z\mathfrak b_z.
\]
There is little danger of confusion with the inner product $(\varphi,\psi)$ in $\mathfrak R_\infty$. We shall denote the differential operator $\frac{\partial^2}{\partial x^2}+\frac{\partial^2}{\partial y^2}+\frac{\partial^2}{\partial z^2}$ by $\Delta$, and the familiar vector differential operations by $\operatorname{div},\operatorname{grad},\operatorname{rot}$. In the empty cavity $\boldsymbol H$, the vectors $\mathfrak E,\mathfrak H$ satisfy \textsc{Maxwell}'s equations
% Source PDF page 144 / original page 137.
\hypertarget{origpage-137}{}%
\begin{align*}
 \operatorname{div}\mathfrak H&=0,&
 \operatorname{rot}\mathfrak E+\frac1c\frac{\partial}{\partial t}\mathfrak H&=0,\\
 \operatorname{div}\mathfrak E&=0,&
 \operatorname{rot}\mathfrak H-\frac1c\frac{\partial}{\partial t}\mathfrak E&=0.
\end{align*}

The first equation of the first row is satisfied by setting $\mathfrak H=\operatorname{rot}\mathfrak A$ ($\mathfrak A=\{\mathfrak A_x,\mathfrak A_y,\mathfrak A_z\}$ is the so-called vector potential, whose components also depend on $x,y,z,t$), and the second then by $\mathfrak E=-\frac1c\frac{\partial}{\partial t}\mathfrak A$. The equations of the second row now require
\begin{equation}
 \operatorname{div}\mathfrak A=0,\qquad
 \Delta\mathfrak A-\frac1{c^2}\frac{\partial^2}{\partial t^2}\mathfrak A=0.
 \tag{A}\label{eq:III.6.A}
\end{equation}
[The vector potential is usually introduced somewhat differently to increase the symmetry between space and time. Most treatments of \textsc{Maxwell}'s theory show that this formulation in terms of $\mathfrak A$ gives the most general solution of \textsc{Maxwell}'s equations---noting in particular that the first equation of the second row really gives only $\frac{\partial}{\partial t}\operatorname{div}\mathfrak A=0$, that is, $\operatorname{div}\mathfrak A=f(x,y,z)$---and so we shall not discuss this further here. Cf.\ the reference in \hyperlink{note-139}{Note~139}.]

Our starting point in what follows is \hyperref[eq:III.6.A]{(A)}. The fact that the walls of $\boldsymbol H$ are reflecting is expressed by the requirement that, at the boundary of $\boldsymbol H$, $\mathfrak A$ be perpendicular to its walls. The familiar method of finding all such $\mathfrak A$ is this: since $t$ does not enter the whole problem explicitly, the most general $\mathfrak A$ is a linear combination of solutions which are products of a vector depending on $x,y,z$ and a scalar depending on $t$:
\[
 \mathfrak A=\mathfrak A(xyzt)=\overline{\mathfrak A}(xyz)\,\widetilde q(t).
\]
Thus \hyperref[eq:III.6.A]{(A)} gives
\begin{align}
 \operatorname{div}\overline{\mathfrak A}&=0,\quad
 \Delta\overline{\mathfrak A}=\eta\overline{\mathfrak A},\quad
 \overline{\mathfrak A}\text{ perpendicular at the boundary of }\boldsymbol H,
 \tag{$A_1$}\label{eq:III.6.A1}\\
 \frac{\partial^2}{\partial t^2}\widetilde q(t)&=c^2\eta\widetilde q(t),
 \tag{$A_2$}\label{eq:III.6.A2}
\end{align}
where, by \hyperref[eq:III.6.A1]{($A_1$)}, $\eta$ depends only on $x,y,z$, and by \hyperref[eq:III.6.A2]{($A_2$)} only on $t$, so it is constant.

Thus \hyperref[eq:III.6.A1]{($A_1$)} is an eigenvalue problem, $\eta$ being the eigenvalue parameter and $\overline{\mathfrak A}$ the general eigenfunction. The theory of this problem has been worked out completely; we give only the results here.\origNote{140} The problem \hyperref[eq:III.6.A1]{($A_1$)} has a pure point spectrum; all eigenvalues $\eta_1,\eta_2,\ldots$ (let the corresponding $\overline{\mathfrak A}$ be $\overline{\mathfrak A}_1,\overline{\mathfrak A}_2,\ldots$) are negative and tend to $-\infty$ as $n\longrightarrow+\infty$. We can normalize a complete system $\mathfrak A_1,\mathfrak A_2,\ldots$ by
\[
 \iiint_{\boldsymbol H}[\overline{\mathfrak A}_m,\overline{\mathfrak A}_n]\,dx\,dy\,dz
 =\begin{cases}4\pi c^2,&m=n,\\0,&m\ne n.\end{cases}
\]
% Source PDF page 145 / original page 138.
\hypertarget{origpage-138}{}%
(We choose $4\pi c^2$ instead of the customary $1$ because it will later prove somewhat more convenient.) Writing $\eta_n$ ($<0$) as $-4\pi^2\varrho_n^2/c^2$ ($\varrho_n>0$, and $\varrho_n\longrightarrow+\infty$ as $n\longrightarrow+\infty$), \hyperref[eq:III.6.A2]{($A_2$)} gives
\[
 \widetilde q_n(t)=\gamma\cos2\pi\varrho_n(t-\tau)
 \qquad(\gamma,\tau\text{ arbitrary}).
\]
Thus the most general solution $\mathfrak A$ is\tnote{The following expansion contains two factors $\gamma_n$ in each term in the source; both are retained.}
\begin{align*}
 \mathfrak A=\mathfrak A(xyzt)
 &=\sum_{n=1}^{\infty}\overline{\mathfrak A}_n(xyz)\widetilde q_n(t)\\
 &=\sum_{n=1}^{\infty}\gamma_n\overline{\mathfrak A}_n(xyz)\gamma_n
 \cos2\pi\varrho_n(t-\tau_n)
\end{align*}
($\gamma_1,\gamma_2,\ldots,\tau_1,\tau_2,\ldots$ arbitrary real constants).

The energy of the arbitrary field $\mathfrak A=\sum_{n=1}^{\infty}\overline{\mathfrak A}_n(xyz)\widetilde q_n(t)$ [$\mathfrak A$ is not assumed to solve \hyperref[eq:III.6.A]{(A)}, so the $\widetilde q_n(t)$ are arbitrary] is\tnote{The second line of the energy calculation omits the factor $1/c^2$ before the time-derivative inner product in the source; the next line includes it. The formulas retain this printed discrepancy.}
\begin{align*}
 E&=\frac1{8\pi}\iiint_{\boldsymbol H}
 \bigl([\mathfrak E,\mathfrak E]+[\mathfrak H,\mathfrak H]\bigr)\,dx\,dy\,dz\\
 &=\frac1{8\pi}\iiint_{\boldsymbol H}
 \left(\left[\frac{\partial\mathfrak A}{\partial t},\frac{\partial\mathfrak A}{\partial t}\right]
 +[\operatorname{rot}\mathfrak A,\operatorname{rot}\mathfrak A]\right)\,dx\,dy\,dz\\
 &=\frac1{8\pi}\sum_{m,n=1}^{\infty}\iiint_{\boldsymbol H}
 \left(\frac1{c^2}\frac{\partial\widetilde q_m(t)}{\partial t}
 \frac{\partial\widetilde q_n(t)}{\partial t}
 [\overline{\mathfrak A}_m,\overline{\mathfrak A}_n]\right.\\
 &\qquad\left.{}+\widetilde q_m(t)\widetilde q_n(t)
 [\operatorname{rot}\overline{\mathfrak A}_m,\operatorname{rot}\overline{\mathfrak A}_n]\right)\,dx\,dy\,dz.
\end{align*}
Integration by parts gives:\origNote{141}
\begin{align*}
 &\iiint_{\boldsymbol H}[\operatorname{rot}\overline{\mathfrak A}_m,
 \operatorname{rot}\overline{\mathfrak A}_n]\,dx\,dy\,dz
 =\iiint_{\boldsymbol H}[\operatorname{rot}\operatorname{rot}\overline{\mathfrak A}_m,
 \overline{\mathfrak A}_n]\,dx\,dy\,dz\\
 &\quad=\iiint_{\boldsymbol H}[-\Delta\overline{\mathfrak A}_m
 +\operatorname{grad}\operatorname{div}\overline{\mathfrak A}_m,
 \overline{\mathfrak A}_n]\,dx\,dy\,dz\\
 &\quad=\frac{4\pi^2\varrho_m^2}{c^2}
 \iiint_{\boldsymbol H}[\overline{\mathfrak A}_m,\overline{\mathfrak A}_n]\,dx\,dy\,dz.
\end{align*}
Consequently,
\begin{align*}
 E&=\frac1{8\pi}\sum_{m,n=1}^{\infty}
 \left(\frac1{c^2}\frac{\partial\widetilde q_m(t)}{\partial t}
 \frac{\partial\widetilde q_n(t)}{\partial t}
 +\frac{4\pi^2\varrho_m^2}{c^2}\widetilde q_m(t)\widetilde q_n(t)\right)\\
 &\qquad{}\times\iiint_{\boldsymbol H}
 [\overline{\mathfrak A}_m,\overline{\mathfrak A}_n]\,dx\,dy\,dz\\
 &=\frac12\sum_{m=1}^{\infty}
 \left[\left(\frac{\partial\widetilde q_m(t)}{\partial t}\right)^2
 +4\pi^2\varrho_m^2(\widetilde q_m(t))^2\right].
\end{align*}
But we can regard $\widetilde q_1,\widetilde q_2,\ldots$ as the coordinates describing the instantaneous state of the field, that is, as the coordinates of the configuration space of $\boldsymbol{L}$.
% Source PDF page 146 / original page 139.
\hypertarget{origpage-139}{}%
From the formula
\[
 E=\frac12\sum_{n=1}^{\infty}
 \left[\left(\frac{\partial\widetilde q_n}{\partial t}\right)^2
 +4\pi^2\varrho_n^2\widetilde q_n^2\right]
\]
the conjugate momentum $\widetilde p_n$ is determined, in classical mechanics, as
\[
 \widetilde p_n=\frac{\partial E}{\partial(\partial\widetilde q_n/\partial t)}
 =\frac{\partial\widetilde q_n}{\partial t},\qquad
 E=\frac12\sum_{n=1}^{\infty}(\widetilde p_n^2+4\pi^2\varrho_n^2\widetilde q_n^2)
\]
(cf.\ \hyperref[sec:I.2]{I.2}). From this, classical mechanics gives the equations of motion
\begin{align*}
 \frac{\partial\widetilde p_n}{\partial t}
 &=-\frac{\partial E}{\partial\widetilde q_n}=-4\pi^2\varrho_n^2\widetilde q_n,\\
 \frac{\partial\widetilde q_n}{\partial t}
 &=\frac{\partial E}{\partial\widetilde p_n}=\widetilde p_n,
\end{align*}
that is,
\[
 \frac{\partial^2\widetilde q_n}{\partial t^2}+4\pi^2\varrho_n^2\widetilde q_n=0,
\]
exactly \hyperref[eq:III.6.A2]{($A_2$)}, which followed from \textsc{Maxwell}'s equations. Thus \textsc{Jeans}'s \textbf{theorem} holds:

The radiation field $\boldsymbol{L}$ admits a purely classical mechanical description. Its coordinates $\widetilde q_1,\widetilde q_2,\ldots$ are related to the instantaneous vector potential $\mathfrak A$ describing the field by
\[
 \mathfrak A=\mathfrak A(xyz)=\sum_{n=1}^{\infty}\widetilde q_n\overline{\mathfrak A}_n(xyz)
\]
and its energy (Hamiltonian function) is
\[
 E=\tfrac12\sum_{n=1}^{\infty}(\widetilde p_n^2+4\pi^2\varrho_n^2\widetilde q_n^2).
\]

A point of mass $1$, constrained to a straight line and having coordinate $\widetilde q$, in the potential field $C\widetilde q^2$, $C=2\pi^2\varrho^2$, has energy $\frac12[(\frac{\partial\widetilde q}{\partial t})^2+4\pi^2\varrho^2\widetilde q^2]$, or, since again $\widetilde p=\frac{\partial\widetilde q}{\partial t}$, $\frac12(\widetilde p^2+4\pi^2\varrho^2\widetilde q^2)$. Its equation of motion is therefore $\frac{\partial^2\widetilde q}{\partial t^2}+4\pi^2\varrho^2\widetilde q=0$, with solution $\widetilde q=\gamma\cos2\pi\varrho(t-\tau)$ ($\gamma,\tau$ arbitrary). In view of the form of its motion, this system is called a ``linear oscillator of frequency $\varrho$.'' Thus $\boldsymbol{L}$ is to be regarded as the combination of a sequence of linear oscillators whose frequencies are the natural frequencies of the cavity $\boldsymbol H$: $\varrho_1,\varrho_2,\ldots$.

This ``mechanical'' description of the electromagnetic field is important because it can immediately be reinterpreted quantum-mechanically by the usual method. The configuration space of $\boldsymbol{L}$ is described by $\widetilde q_1,\widetilde q_2,\ldots$; in the expression for $E$, $\widetilde p_n,\widetilde q_n$ are to be replaced by $\frac{h}{2\pi i}\frac{\partial}{\partial\widetilde q_n}\cdots$ and $\widetilde q_n\cdots$, respectively. Call these operators $\widetilde P_n$ and $\widetilde Q_n$.
% Source PDF page 147 / original page 140.
\hypertarget{origpage-140}{}%
This answers questions~\hyperref[problem:III.6.1]{\textbf{1}} and \hyperref[problem:III.6.2beta]{\textbf{2}$\beta)$}; in particular,
\[
 \mathsf{H}_l=\tfrac12\sum_{n=1}^{\infty}(\widetilde P_n^2+4\pi^2\varrho_n^2\widetilde Q_n^2)
\]
is the operator required in \hyperref[problem:III.6.2beta]{\textbf{2}$\beta)$}. Question~\hyperref[problem:III.6.2alpha]{\textbf{2}$\alpha)$} was solved from the outset, since we assumed $\mathsf{H}_0$ known. Only \hyperref[problem:III.6.2gamma]{\textbf{2}$\gamma)$} remains, and it now presents no difficulty either.

In classical electrodynamics, the interaction of $\boldsymbol{S}$ with $\boldsymbol{L}$ is calculated as follows. Let $\boldsymbol{S}$ consist of $l$ particles (say, protons and electrons), with respective charges and masses $e_1,m_1,\ldots,e_l,m_l$ and Cartesian coordinates $x_1=q_1,y_1=q_2,z_1=q_3,\ldots,x_l=q_{3l-2},y_l=q_{3l-1},z_l=q_{3l}$ (these constitute the earlier symbol $\xi$); let the corresponding momenta be $p_1^x,p_1^y,p_1^z,\ldots,p_l^x,p_l^y,p_l^z$. The interaction energy is, to sufficient approximation,\origNote{142}
\begin{multline*}
 \sum_{\nu=1}^{l}\frac{e_\nu}{cm_\nu}
 \bigl(p_\nu^x\mathfrak A_x(x_\nu,y_\nu,z_\nu)
 +p_\nu^y\mathfrak A_y(x_\nu,y_\nu,z_\nu)\\
 {}+p_\nu^z\mathfrak A_z(x_\nu,y_\nu,z_\nu)\bigr).
\end{multline*}
The corresponding quantum-mechanical operator is obtained by replacing $p_\nu^x,p_\nu^y,p_\nu^z,x_\nu,y_\nu,z_\nu$ ($\nu=1,\ldots,l$) by the operators
\[
 \frac{h}{2\pi i}\frac{\partial}{\partial x_\nu}\cdots,\quad
 \frac{h}{2\pi i}\frac{\partial}{\partial y_\nu}\cdots,\quad
 \frac{h}{2\pi i}\frac{\partial}{\partial z_\nu}\cdots,\quad
 x_\nu\cdots,\quad y_\nu\cdots,\quad z_\nu\cdots,
\]
which we denote by $P_\nu^x,P_\nu^y,P_\nu^z,Q_\nu^x,Q_\nu^y,Q_\nu^z$. Taking account also of\tnote{The following expansion is printed with $\varrho_n$ as coefficient; the earlier field expansion uses $\widetilde q_n$.}
\[
 \mathfrak A(xyz)=\sum_{n=1}^{\infty}\varrho_n\overline{\mathfrak A}_n(xyz),
\]
we obtain the desired $\mathsf{H}_w$:
\begin{align*}
 \mathsf{H}_w=\sum_{n=1}^{\infty}\sum_{\nu=1}^{l}\frac{e_\nu}{cm_\nu}\widetilde Q_n
 \bigl\{&P_\nu^x\overline{\mathfrak A}_{n,x}(Q_\nu^x,Q_\nu^y,Q_\nu^z)\\
 &{}+P_\nu^y\overline{\mathfrak A}_{n,y}(Q_\nu^x,Q_\nu^y,Q_\nu^z)
 +P_\nu^z\overline{\mathfrak A}_{n,z}(Q_\nu^x,Q_\nu^y,Q_\nu^z)\bigr\}.
\end{align*}
Here, however, it must be noted that in replacing the products $p_\nu^x\overline{\mathfrak A}_{n,x}(x_\nu,y_\nu,z_\nu),\ldots$ by operators, we arbitrarily fixed the order $P_\nu^x\overline{\mathfrak A}_{n,x}(Q_\nu^x,Q_\nu^y,Q_\nu^z),\ldots$, although we might equally well have taken $\overline{\mathfrak A}_{n,x}(Q_\nu^x,Q_\nu^y,Q_\nu^z)P_\nu^x,\ldots$; in fact, to ensure that the resulting operator is Hermitian, a symmetric form such as
\[
 \tfrac12\bigl(P_\nu^x\overline{\mathfrak A}_{n,x}(Q_\nu^x,Q_\nu^y,Q_\nu^z)
 +\overline{\mathfrak A}_{n,x}(Q_\nu^x,Q_\nu^y,Q_\nu^z)P_\nu^x\bigr),\ldots
\]
would be required. Fortunately, all this makes no difference, since\origNote{143}
\begin{align*}
 &[P_\nu^x\overline{\mathfrak A}_{n,x}(Q_\nu^x,Q_\nu^y,Q_\nu^z)+\cdots]
 -[\overline{\mathfrak A}_{n,x}(Q_\nu^x,Q_\nu^y,Q_\nu^z)P_\nu^x+\cdots]\\
 &\quad=[P_\nu^x\overline{\mathfrak A}_{n,x}(Q_\nu^x,Q_\nu^y,Q_\nu^z)
 -\overline{\mathfrak A}_{n,x}(Q_\nu^x,Q_\nu^y,Q_\nu^z)P_\nu^x]+\cdots\\
 &\quad=\frac{h}{2\pi i}\frac{\partial}{\partial x}
 \overline{\mathfrak A}_{n,x}(Q_\nu^x,Q_\nu^y,Q_\nu^z)+\cdots\\
 &\quad=\frac{h}{2\pi i}\operatorname{div}\overline{\mathfrak A}_n(Q_\nu^x,Q_\nu^y,Q_\nu^z)=0.
\end{align*}
% Source PDF page 148 / original page 141.
\hypertarget{origpage-141}{}%
The total energy of our system $\boldsymbol{S}+\boldsymbol{L}$, or its operator
\[
 \mathsf{H}=\mathsf{H}_0+\mathsf{H}_l+\mathsf{H}_w,
\]
has thus been specified exactly. Before transforming $\mathsf{H}$ further, however, we make the following remark. The configuration space of $\boldsymbol{S}+\boldsymbol{L}$ is described by the coordinates $\xi$ (that is, $q_1,\ldots,q_{3l}$ or $x_1,y_1,z_1,\ldots,x_l,y_l,z_l$) and $\widetilde q_1,\widetilde q_2,\ldots$; the wave function therefore also depends on these. Now it is formally inconvenient and questionable to admit systems with infinitely many degrees of freedom, or wave functions with infinitely many arguments; our prescriptions have always referred to finitely many coordinates. We shall therefore initially consider only the first $N$ of $\widetilde q_1,\widetilde q_2,\ldots$, namely $\widetilde q_1,\ldots,\widetilde q_N$ (that is, restrict $\mathfrak A$ to linear combinations of $\overline{\mathfrak A}_1,\ldots,\overline{\mathfrak A}_N$), and only in the completed result perform the physically necessary limiting passage $N\longrightarrow+\infty$. Thus
\begin{align*}
 \mathsf{H}={}&\mathsf{H}_0+\tfrac12\sum_{n=1}^{N}
 (\widetilde P_n^2+4\pi^2\varrho_n^2\widetilde Q_n^2)\\
 &{}+\sum_{n=1}^{N}\sum_{\nu=1}^{l}\frac{e_\nu}{cm_\nu}\widetilde Q_n
 \bigl\{P_\nu^x\overline{\mathfrak A}_{n,x}(Q_\nu^x,Q_\nu^y,Q_\nu^z)\\
 &\qquad{}+P_\nu^y\overline{\mathfrak A}_{n,y}(Q_\nu^x,Q_\nu^y,Q_\nu^z)
 +P_\nu^z\overline{\mathfrak A}_{n,z}(Q_\nu^x,Q_\nu^y,Q_\nu^z)\bigr\}.
\end{align*}
It proves convenient to introduce, in place of $\widetilde P_n,\widetilde Q_n$, the non-Hermitian operator $\widetilde R_n$ and its adjoint $\widetilde R_n^*$:
\begin{align*}
 \widetilde R_n&=\frac1{\sqrt{2h\varrho_n}}(2\pi\varrho_n\widetilde Q_n+i\widetilde P_n),\\
 \widetilde R_n^*&=\frac1{\sqrt{2h\varrho_n}}(2\pi\varrho_n\widetilde Q_n-i\widetilde P_n).
\end{align*}
Then $\widetilde Q_n=\frac1{2\pi}\sqrt{\frac{h}{2\varrho_n}}(\widetilde R_n+\widetilde R_n^*)$, and, since $\widetilde P_n\widetilde Q_n-\widetilde Q_n\widetilde P_n=\frac{h}{2\pi i}1$,
\begin{align*}
 \widetilde R_n\widetilde R_n^*
 &=\frac1{2h\varrho_n}(\widetilde P_n^2+4\pi^2\varrho_n^2\widetilde Q_n^2)+\tfrac12\,1,\\
 \widetilde R_n^*\widetilde R_n
 &=\frac1{2h\varrho_n}(\widetilde P_n^2+4\pi^2\varrho_n^2\widetilde Q_n^2)-\tfrac12\,1,
\end{align*}
and in particular $\widetilde R_n\widetilde R_n^*-\widetilde R_n^*\widetilde R_n=1$. The energy formula is\tnote{In the square root in the interaction term, the source prints $\varrho_\nu$; the preceding substitution gives $\varrho_n$. The printed subscript is retained.}
\begin{align*}
 \mathsf{H}={}&\mathsf{H}_0+\sum_{n=1}^{N}h\varrho_n\widetilde R_n^*\widetilde R_n\\
 &{}+\sum_{n=1}^{N}\sum_{\nu=1}^{l}\frac{e_\nu}{2\pi cm_\nu}
 \sqrt{\frac{h}{2\varrho_\nu}}(\widetilde R_n+\widetilde R_n^*)\\
 &\quad{}\times\bigl[P_\nu^x\overline{\mathfrak A}_{n,x}(Q_\nu^x,Q_\nu^y,Q_\nu^z)
 +P_\nu^y\overline{\mathfrak A}_{n,y}(Q_\nu^x,Q_\nu^y,Q_\nu^z)\\
 &\qquad{}+P_\nu^z\overline{\mathfrak A}_{n,z}(Q_\nu^x,Q_\nu^y,Q_\nu^z)\bigr]+C,
\end{align*}
where $C=\text{constant}=\frac12\sum_{n=1}^{N}h\varrho_n$. Since an additive constant in the energy expression is meaningless, we may omit $C$; this is all the more appropriate because $C$ becomes infinite as $N\longrightarrow+\infty$, thus obstructing the proper completion of the theory.

\hypertarget{origpage-142}{}%
The Hermitian operator $\widetilde R_n^*\widetilde R_n$ is hypermaximal; indeed, it has a pure point spectrum consisting of the numbers $0,1,2,\ldots$. Let the corresponding eigenfunctions be called $\psi_0^n(\widetilde q_n),\psi_1^n(\widetilde q_n),\psi_2^n(\widetilde q_n),\ldots$.

[If, instead of the variable $\widetilde q_n$, we write $\frac1{2\pi}\sqrt{h/\varrho_n}\,q$, then
\begin{gather*}
 \frac1{\sqrt{2h\varrho_n}}2\pi\varrho_n\widetilde q_n
 =2\pi\sqrt{\frac{\varrho_n}{2h}}\,\widetilde q_n,\\
 \frac1{\sqrt{2h\varrho_n}}\frac h{2\pi i}\frac\partial{\partial\widetilde q_n}
 =\frac1{2\pi}\sqrt{\frac h{2\varrho_n}}\frac1i\frac\partial{\partial\widetilde q_n}
\end{gather*}
become $q/\sqrt2$ and $\frac1{\sqrt2}\frac1i\frac\partial{\partial q}$, respectively, so that
\begin{align*}
 \widetilde R_n&=\frac1{\sqrt2}\left(q+\frac\partial{\partial q}\right),&
 \widetilde R_n^*&=\frac1{\sqrt2}\left(q-\frac\partial{\partial q}\right),\\
 \widetilde R_n\widetilde R_n^*&=-\frac12\frac{\partial^2}{\partial q^2}+\frac12q^2+\frac12,\\
 \widetilde R_n^*\widetilde R_n&=-\frac12\frac{\partial^2}{\partial q^2}+\frac12q^2-\frac12.
\end{align*}
The reader will find the eigenvalue theory of these operators in many accounts, e.g., Courant--Hilbert, p.~261, formulas (42), (43) and the discussion following them, and p.~76, formulas (60), (61); or Weyl, \emph{Gruppentheorie und Quantenmechanik}, pp.~74 ff.]

Since $\psi_1(\xi),\psi_2(\xi),\ldots$ form a complete orthogonal system in $\xi$-space, and $\psi_0^n(\widetilde q_n),\psi_1^n(\widetilde q_n),\ldots$ one in $\widetilde q_n$-space, the functions
\[
 \Phi_{kM_1\cdots M_N}(\xi\widetilde q_1\cdots\widetilde q_N)
 =\psi_k(\xi)\psi_{M_1}^1(\widetilde q_1)\cdots\psi_{M_N}^N(\widetilde q_N),
\]
$k=1,2,\ldots$, $M_1,\ldots,M_N=0,1,2,\ldots$, form a complete orthogonal system in $\xi,\widetilde q_1,\ldots,\widetilde q_N$-space, i.e., in configuration space. We can therefore expand every wave function $\Phi=\Phi(\xi,\widetilde q_1,\ldots,\widetilde q_N)$ as follows:
\begin{align*}
 \varphi(\xi,\widetilde q_1,\ldots,\widetilde q_N)
 &=\sum_{k=1}^{\infty}\sum_{M_1=0}^{\infty}\cdots\sum_{M_N=0}^{\infty}
 a_{kM_1\cdots M_N}\Phi_{kM_1\cdots M_N}(\xi\widetilde q_1\cdots\widetilde q_N)\\
 &=\sum_{k=1}^{\infty}\sum_{M_1=0}^{\infty}\cdots\sum_{M_N=0}^{\infty}
 a_{kM_1\cdots M_N}\psi_k(\xi)\psi_{M_1}^1(\widetilde q_1)\cdots\psi_{M_N}^N(\widetilde q_N).
\end{align*}
It is immaterial that we enumerate the complete orthogonal system and the expansion coefficients with $N+1$ indices $k,M_1,\ldots,M_N$ instead of one. On the basis of the considerations of \hyperref[sec:II.2]{II.2}, we can state that the Hilbert space of the wave functions $\Phi$ can also be realized as a space of $(N+1)$-indexed sequences $a_{kM_1\cdots M_N}$ (with finite $\sum_{k=1}^{\infty}\sum_{M_1=0}^{\infty}\cdots\sum_{M_N=0}^{\infty}|a_{kM_1\cdots M_N}|^2$). What form does $\mathsf{H}$ take in this realization of Hilbert space? To determine this, we first calculate the $\mathsf{H}\Phi_{kM_1\cdots M_N}$. Since $\mathsf{H}_0$ acts on $\xi$ alone and $\psi_k(\xi)$ is an eigenfunction of $\mathsf{H}_0$ with eigenvalue $W_k$, and since $\widetilde R_n^*\widetilde R_n$ acts on $\widetilde q_n$ alone and $\psi_{M_n}^n(\widetilde q_n)$ \hypertarget{origpage-143}{}is an eigenfunction of $\widetilde R_n^*\widetilde R_n$ with eigenvalue $M_n$, we obtain
\begin{align*}
 \mathsf{H}\Phi_{kM_1\cdots M_N}
 ={}&\left(W_k+\sum_{n=1}^{N}h\varrho_n M_n\right)\Phi_{kM_1\cdots M_N}\\
 &{}+\sum_{n=1}^{N}\sum_{\nu=1}^{l}\frac{e_\nu}{2\pi cm_\nu}\sqrt{\frac h{2\varrho_n}}\\
 &\quad{}\times\bigl[P_\nu^x\overline{\mathfrak A}_{n,x}(Q_\nu^x,Q_\nu^y,Q_\nu^z)
 +P_\nu^y\overline{\mathfrak A}_{n,y}(Q_\nu^x,Q_\nu^y,Q_\nu^z)\\
 &\qquad{}+P_\nu^z\overline{\mathfrak A}_{n,z}(Q_\nu^x,Q_\nu^y,Q_\nu^z)\bigr]\psi_k(\xi)\\
 &\quad{}\times\psi_{M_1}^1(\widetilde q_1)\cdots
 (\widetilde R_n+\widetilde R_n^*)\psi_{M_n}^n(\widetilde q_n)\cdots\psi_{M_N}^N(\widetilde q_N).
\end{align*}
For all operators $A$ that, like the expression in square brackets, affect only the variable $\xi$, we can use the expansion
\[
 A\psi_k(\xi)=\sum_{j=1}^{\infty}(A\psi_k,\psi_j)\psi_j(\xi)
 =\sum_j(A)_{kj}\psi_j(\xi),
\]
where $(A)_{kj}=(A\psi_k,\psi_j)$ is thus defined. Furthermore, as shown in the passage cited above,
\[
 \widetilde R_n\psi_M^n(\widetilde q_n)=\sqrt M\,\psi_{M-1}^n(\widetilde q_n),\qquad
 \widetilde R_n^*\psi_M^n(\widetilde q_n)=\sqrt{M+1}\,\psi_{M+1}^n(\widetilde q_n)
\]
(for $M=0$ the right-hand side of the first formula is to be set equal to $0$, regardless of the meaningless $\psi_{-1}^n$ occurring in it). Hence
\begin{align*}
 \mathsf{H}\Phi_{kM_1\cdots M_N}
 ={}&\left(W_k+\sum_{n=1}^{N}h\varrho_n M_n\right)\Phi_{kM_1\cdots M_N}\\
 &{}+\sum_{j=1}^{\infty}\sum_{n=1}^{N}\sqrt{\frac h{2\varrho_n}}
 \left(\sum_{\nu=1}^{l}\frac{e_\nu}{2\pi cm_\nu}
 \bigl(P_\nu^x\overline{\mathfrak A}_{n,x}(Q_\nu^x,Q_\nu^y,Q_\nu^z)+\cdots\bigr)_{kj}\right)\\
 &\quad{}\times\left(\sqrt{M_n+1}\,\Phi_{kM_1\cdots M_n+1\cdots M_N}
 +\sqrt{M_n}\,\Phi_{kM_1\cdots M_n-1\cdots M_N}\right).
\end{align*}
We can now determine the form of $\mathsf{H}$ as an operation on $a_{kM_1\cdots M_N}$:\tnote{In the preceding source formula, the sum is over $j$, but the two functions $\Phi$ retain the initial subscript $k$. The displayed indices are reproduced as printed.}
\[
 \mathsf{H}\sum_{kM_1\cdots M_N}a_{kM_1\cdots M_N}\Phi_{kM_1\cdots M_N}
 =\sum_{kM_1\cdots M_N}a'_{kM_1\cdots M_N}\Phi_{kM_1\cdots M_N},
\]
with
\begin{align*}
 &\mathsf{H}a_{kM_1\cdots M_N}=a'_{kM_1\cdots M_N}\\
 &\quad={}\left(W_k+\sum_{n=1}^{N}h\varrho_n M_n\right)a_{kM_1\cdots M_N}\\
 &{}+\sum_{j=1}^{\infty}\sum_{n=1}^{N}\sqrt{\frac h{2\varrho_n}}
 \left(\sum_{\nu=1}^{l}\frac{e_\nu}{2\pi cm_\nu}
 \bigl(P_\nu^x\overline{\mathfrak A}_{n,x}(Q_\nu^x,Q_\nu^y,Q_\nu^z)+\cdots\bigr)_{kj}\right)\\
 &\quad{}\times\left(\sqrt{M_n}\,a_{jM_1\cdots M_n-1\cdots M_N}
 +\sqrt{M_n+1}\,a_{jM_1\cdots M_n+1\cdots M_N}\right).
\end{align*}
Our discussion of $\mathsf{H}$ has now reached the point where it is appropriate to take the limit $N\longrightarrow\infty$. Since the indexing of the $a_{kM_1\cdots M_N}$ changes in this process, an entirely new operator $\mathsf{H}$ results. We must introduce components $a_{kM_1M_2\cdots}$ with infinitely many indices $M_1,M_2,\ldots$; but merely to preserve the finiteness of the sum $\sum_{n=1}^{\infty}h\varrho_n M_n$ occurring in $\mathsf{H}$, we must restrict ourselves to
\hypertarget{origpage-144}{}sequences $M_1,M_2,\ldots$ containing only finitely many nonzero numbers. Henceforth, therefore, Hilbert space is formed by all sequences $a_{kM_1M_2\cdots}$ with finite $\sum_{k=1}^{\infty}\sum_{M_1=0}^{\infty}\sum_{M_2=0}^{\infty}\cdots|a_{kM_1M_2\cdots}|^2$, where the indices $k,M_1,M_2,\ldots$ range over the following domain: $k=1,2,\ldots$, $M_1,M_2,\ldots=0,1,2,\ldots$, with only finitely (but arbitrarily) many $M_n\ne0$.\origNote{144} And the final form of $\mathsf{H}$ is
\begin{align*}
 &\mathsf{H}a_{kM_1M_2\cdots}=a'_{kM_1M_2\cdots}\\
 &\quad={}\left(W_k+\sum_{n=1}^{\infty}h\varrho_n M_n\right)a_{kM_1M_2\cdots}\\
 &{}+\sum_{j=1}^{\infty}\sum_{n=1}^{\infty}w_{kj}^n
 \left(\sqrt{M_n+1}\,a_{jM_1M_2\cdots M_n+1\cdots}
 +\sqrt{M_n}\,a_{jM_1M_2\cdots M_n-1\cdots}\right),
\end{align*}
where $w_{kj}^n$ is defined by
\[
 w_{kj}^n=\sqrt{\frac h{2\varrho_n}}\sum_{\nu=1}^{l}\frac{e_\nu}{2\pi cm_\nu}
 \bigl(P_\nu^x\overline{\mathfrak A}_{n,x}(Q_\nu^x,Q_\nu^y,Q_\nu^z)+\cdots\bigr)_{kj}.
\]
Before drawing from this result the physical conclusions that interest us, let us recall that it was obtained on the basis of the electrodynamic theory of light. We shall first want to ascertain whether the routine application of the quantum-mechanical transformation that we performed has sufficed to reproduce the deviations of light from the wave model---in particular, its actually discrete, corpuscular nature. (It is conceivable that, to comprehend this, one would have to start directly from a corpuscular model of light, instead of ``quantizing'' the electromagnetic field as we have done here.)

Our expression for $\mathsf{H}$ already represents something like light corpuscles. Its second summand is a kind of perturbation which, as we shall soon see, causes the quantum jumps of the system $\boldsymbol{S}$ from one ``stationary'' state to another. These jumps are the phenomenon that really interests us, but the corresponding contribution is considerably smaller than those of the material system $\boldsymbol{S}$ and the radiation already present, which, as we shall see, are represented by the first summand of $\mathsf{H}$. If we omit the second summand, there remains
\[
 \mathsf{H}_1a_{kM_1M_2\cdots}=\left(W_k+\sum_{n=1}^{\infty}h\varrho_n M_n\right)a_{kM_1M_2\cdots}.
\]
But this energy expression can be interpreted as follows: the energy $W_k$ of the system $\boldsymbol{S}$, increased by the energies $h\varrho_n M_n$ ($n=1,2,\ldots$). Here it is natural to regard the numbers $M_n=0,1,2,\ldots$ as numbers of particles with the respective energies $h\varrho_n$; and $h\varrho_n$ is precisely the energy that, according to Einstein, is to be ascribed to a light quantum of frequency $\varrho_n$ (cf.\ \hyperlink{note-134}{Note~134}). Thus the structure of $\mathsf{H}_1$ suggests that the electromagnetic \hypertarget{origpage-145}{}field existing in $\boldsymbol H$ (less the electrostatic part), i.e., $\boldsymbol{L}$, actually consists of light quanta with frequencies $\varrho_1,\varrho_2,\ldots$ and energies $h\varrho_1,h\varrho_2,\ldots$, whose respective numbers are the indices $M_1,M_2,\ldots$ ($=0,1,2,\ldots$). That no frequencies other than precisely $\varrho_1,\varrho_2,\ldots$ are taken into account can be made plausible by the fact that these are the natural frequencies of $\boldsymbol H$: the vector potentials $\overline{\mathfrak A}_n(xyz)\gamma\cos2\pi\varrho_n(t-\tau)$ represent the only stationary electromagnetic oscillations possible in $\boldsymbol H$.

These speculations and interpretations, however, are only of heuristic value. We have a satisfactory and definitive answer to our question only if we succeed in arriving at the energy expression $\mathsf{H}$ by starting from the light-quantum model for the radiation $\boldsymbol{L}$. Our reason for carrying out the classical treatment first was that no expression was available for the interaction energy of a light quantum with matter (at this point the reinterpretation of classical electrodynamics was not feasible). Now, however, if the result we shall derive with a general interaction energy agrees in form with $\mathsf{H}$, we shall be able to determine that energy by comparing coefficients.

What is the configuration space of $\boldsymbol{L}$ (\hyperref[problem:III.6.1]{question~1}) on the basis of the light-quantum hypothesis? Let a single light quantum (in the cavity $\boldsymbol H$) be characterized by certain coordinates, the totality of which we shall replace by the symbol $u$.\origNote{145} Let its stationary states (in $\boldsymbol H$) have the wave functions $\psi_1(u),\psi_2(u),\ldots$ (forming a complete orthonormal system) and the energies $E_1,E_2,\ldots$; these correspond to the electromagnetic normal oscillations $\overline{\mathfrak A}_1,\overline{\mathfrak A}_2,\ldots$ with frequencies $\varrho_1,\varrho_2,\ldots$. (In Einstein's conception, $E_n=h\varrho_n$ must hold, which we shall also prove.) The following must also be observed here. Even in the electromagnetic argument we chose the zero of the light energy so that its minimum value was $0$; this corresponded to the indices $M_1=M_2=\cdots=0$. We thereby also recognized nonexistence as a possible state of light, which is justified by the facts. Indeed, light quanta can be emitted and absorbed, i.e., created and annihilated. Such a conception, however, is wholly foreign to quantum mechanics: every particle contributes its coordinates to the configuration space of the system and thus enters so intimately into the formal description of the total system that it is absolutely indestructible. Even after its annihilation, we must ascribe to it a kind of latent existence, since its coordinates still belong to configuration space. Consequently, one of the states $\psi_n(u)$, with energy $E_n=0$, must correspond to the nonexistence of the light quantum. We prefer to denote it by $\psi_0(u)$ ($E_0=0$), so that $\psi_1(u),\psi_2(u),\ldots$
\hypertarget{origpage-146}{}correspond to an existing light quantum, but only $\psi_0(u),\psi_1(u),\psi_2(u),\ldots$ form a complete orthogonal system.

We now pass to $\boldsymbol{L}$, the system of all light quanta. Since we also count light quanta that do not exist, $\boldsymbol{L}$ consists of so many light quanta that no more than this number can ever arise: that is, of infinitely many. But since it is inexpedient to operate from the outset with infinitely many constituents in $\boldsymbol{L}$, we first proceed as though only $S$ light quanta existed ($S=1,2,\ldots$), and only at the end take the limit $S\longrightarrow+\infty$.\origNote{146} We give these $S$ light quanta the numbers $1,\ldots,S$; let their coordinates be called $u_1,\ldots,u_S$. The configuration space of $\boldsymbol{L}$ is thus described by $u_1,\ldots,u_S$, and that of $\boldsymbol{S}+\boldsymbol{L}$ by $\xi,u_1,\ldots,u_S$. The most general wave function for $\boldsymbol{S}+\boldsymbol{L}$ is therefore $f(\xi,u_1,\ldots,u_S)$; the functions $\varphi_k(\xi)\psi_{n_1}(u_1)\cdots\psi_{n_S}(u_S)$, $k=1,2,\ldots$, $u_1,\ldots,u_S=0,1,2,\ldots$, form a complete orthogonal system.\tnote{The source prints $u_1,\ldots,u_S$ in this index range; the indices of the displayed wave functions are $n_1,\ldots,n_S$.}

Light quanta have the fundamental property of being completely identical with one another, i.e., there is no means in the world of distinguishing two light quanta with the same coordinate $u$. Or again: a state in which light quanta Nos.~$m$ and $n$ have the respective $u$-values $u_m=u'$, $u_n=u''$ does not differ from the state in which $u_m=u''$, $u_n=u'$ (this is the classical, not the quantum-mechanical, manner of description, for we specified the value of $u$ and not the wave function $\varphi(u)$!). In quantum-mechanical terms this means that the states belonging to the wave functions $f(\xi,u_1,\ldots,u_m,\ldots,u_n,\ldots,u_S)$ and $f(\xi,u_1,\ldots,u_n,\ldots,u_m,\ldots,u_S)$ are indistinguishable. That is, every physical quantity $\mathfrak R$ has the same expectation value in them (and therefore, since this also applies to the $F(\mathfrak R)$, the same statistics; cf.\ the discussion of \hyperref[post:III.1.E1]{$E_1$} and \hyperref[post:III.1.E2]{$E_2$} in \hyperref[sec:III.1]{III.1}). If we denote the functional operation that interchanges $u_m,u_n$ by $O_{mn}$ ($O_{mn}$ is, as one sees immediately, both Hermitian and unitary, $O_{mn}^2=1$), this means: $\mathfrak R$ has the same expectation value for $f$ as for $O_{mn}f$, i.e.,
\[
 (Rf,f)=(RO_{mn}f,O_{mn}f)=(O_{mn}RO_{mn}f,f),
\]
hence $R=O_{mn}RO_{mn}$, or equivalently $O_{mn}R=RO_{mn}$.
This means that, in the present case, only those operators $R$ are admissible that commute with all $O_{mn}$ ($m,n=1,\ldots,S$, $m\ne n$), i.e., into which, in view of the definition of the $O_{mn}$, all the coordinates $u_1,\ldots,u_S$ enter symmetrically.

A wave function $f$ that is symmetric in all variables $u_1,\ldots,u_S$, i.e., for which $O_{mn}f=f$ ($m,n=1,\ldots,S$, $m\ne n$), is transformed by such an operator $R$ into another of the same kind: $O_{mn}Rf=RO_{mn}f=Rf$. These $f$ form a closed linear manifold, hence a Hilbert subspace $\overline{\mathfrak R}_{\infty}^{(S)}$ in the Hilbert \hypertarget{origpage-147}{}space $\mathfrak R_{\infty}^{(S)}$ of all $f$; and the $R$ map elements of $\overline{\mathfrak R}_{\infty}^{(S)}$ into elements of the same space, i.e., they can be regarded as operators in the Hilbert space $\overline{\mathfrak R}_{\infty}^{(S)}$. Consequently, $\overline{\mathfrak R}_{\infty}^{(S)}$ is just as usable for the purposes of quantum mechanics as the originally envisaged $\mathfrak R_{\infty}^{(S)}$, and the question arises whether, in view of the symmetry of $\boldsymbol{L}$ under interchanges of the light quanta, we should restrict ourselves to symmetric wave functions, i.e., replace $\mathfrak R_{\infty}^{(S)}$ by $\overline{\mathfrak R}_{\infty}^{(S)}$. We shall do so; success, i.e., the complete agreement to be obtained with the electromagnetically derived expression for $\mathsf{H}$, will justify us afterward.\origNote{147}

The $\varphi_k(\xi)\psi_{n_1}(u_1)\cdots\psi_{n_S}(u_S)$ formed a complete orthonormal system in $\mathfrak R_{\infty}^{(S)}$; we now wish to use it to specify one in $\overline{\mathfrak R}_{\infty}^{(S)}$. Let $M_0,M_1,\ldots$ be any numbers $=0,1,2,\ldots$ with $M_0+M_1+\cdots=S$ (thus only finitely many are nonzero!). Denote by $[M_1,M_2,\ldots]$ the totality of all index systems $u_1,\ldots,u_S$ in which $0$ occurs $M_0$ times, $1$ occurs $M_1$ times, etc.; there are exactly $M_0!M_1!\cdots$ different ones. We set
\[
 \Phi_{M_0M_1\cdots}(u_1\cdots u_S)
 =\sum_{n_1\cdots n_S\text{ in }[M_0M_1\cdots]}
 \psi_{n_1}(u_1)\cdots\psi_{n_S}(u_S).
\]
Since $\Phi_{M_0M_1\cdots}$ is the sum of $M_0!M_1!\cdots$ pairwise orthogonal summands of norm $1$, its squared norm is the sum of $M_0!M_1!\cdots$ ones, and thus its norm is $\sqrt{M_0!M_1!\cdots}$. Two different $\Phi_{M_0M_1\cdots}$ have pairwise orthogonal summands and are therefore orthogonal. Thus the functions
\[
 \psi_{M_0M_1\cdots}(u_1\cdots u_S)
 =\frac1{\sqrt{M_0!M_1!\cdots}}\Phi_{M_0M_1\cdots}(u_1\cdots u_S)
\]
form an orthonormal system.\tnote{The counting and normalization factors on this source page are reproduced as printed. For distinct index arrangements with the stated multiplicities, their number is $S!/(M_0!M_1!\cdots)$, rather than the product of factorials printed here. The source also alternates between $[M_1,M_2,\ldots]$ and $[M_0,M_1,\ldots]$, and between $u_j$ and $n_j$ when naming the indices.}
Let $f(\xi,u_1,\ldots,u_S)$ be symmetric in $u_1,\ldots,u_S$. Its inner product with every summand of
\[
 \varphi_k(\xi)\Phi_{M_0M_1\cdots}(u_1\cdots u_S)
\]
is the same. Consequently, if $f$ is orthogonal to this sum, it is orthogonal to each summand. Orthogonality to the sum is equivalent to orthogonality to its normalized counterpart,
\[
 \varphi_k(\xi)\psi_{M_0M_1\cdots}(u_1\cdots u_S).
\]
If $f$ is orthogonal to every such normalized product, it is therefore orthogonal to every function
\[
 \varphi_k(\xi)\psi_{n_1}(u_1)\cdots\psi_{n_S}(u_S),
\]
and hence $f=0$. The normalized products themselves belong to $\overline{\mathfrak R}_{\infty}^{(S)}$, so they form a complete orthonormal system in that space.

Let us now consider the energies occurring in the system $\boldsymbol{S}+\boldsymbol{L}$. First \hyperref[problem:III.6.2alpha]{[2a]}, there is the energy of $\boldsymbol{S}$, whose operator for $\boldsymbol{S}$ is defined by $\mathsf{H}_0\varphi_k(\xi)=W_k\varphi_k(\xi)$, and for $\boldsymbol{S}+\boldsymbol{L}$ therefore by
\[
 \mathsf{H}_0\varphi_k(\xi)\psi_{M_0M_1\cdots}(u_1\cdots u_S)
 =W_k\varphi_k(\xi)\psi_{M_0M_1\cdots}(u_1\cdots u_S).
\]
Second \hyperref[problem:III.6.2beta]{[2b]}, every light quantum $l'$ has an energy operator satisfying $\mathsf{H}_{l'}\psi_n(u)=E_n\psi_n(u)$,
\hypertarget{origpage-148}{}so the $m$th one in $\boldsymbol{S}+\boldsymbol{L}$ ($m=1,\ldots,S$) has
\begin{align*}
 &\mathsf{H}_{l_m}\varphi_k(\xi)\psi_{n_1}(u_1)\cdots\psi_{n_m}(u_m)\cdots\psi_{n_S}(u_S)\\
 &\quad=E_{n_m}\varphi_k(\xi)\psi_{n_1}(u_1)\cdots\psi_{n_m}(u_m)\cdots\psi_{n_S}(u_S),
\end{align*}
and we must form $\mathsf{H}_l=\mathsf{H}_{l_1}+\cdots+\mathsf{H}_{l_S}$. Finally \hyperref[problem:III.6.2gamma]{[2c]}, let the interaction energy of a light quantum $l'$ with $\boldsymbol{S}$ be described by an operator $V$, for the moment not known more precisely, which we characterize by its matrix:
\[
 V_{l'}\varphi_k(\xi)\psi_n(u)=\sum_{j=1}^{\infty}\sum_{p=0}^{\infty}V_{kn/jp}\varphi_j(\xi)\psi_p(u).
\]
In $\boldsymbol{S}+\boldsymbol{L}$, this is, for the $m$th light quantum,
\begin{align*}
 &V_{l_m}\varphi_k(\xi)\psi_{n_1}(u_1)\cdots\psi_{n_m}(u_m)\cdots\psi_{n_S}(u_S)\\
 &\quad=\sum_{j=1}^{\infty}\sum_{p=0}^{\infty}V_{kn_m/jp}\varphi_j(\xi)
 \psi_{n_1}(u_1)\cdots\psi_p(u_m)\cdots\psi_{n_S}(u_S)\\
 &\quad=\sum_{j=1}^{\infty}\sum_{p_1,\ldots,p_S=0}^{\infty}
 \delta(n_1-p_1)\cdots V_{kn_m/jp_m}\cdots\delta(n_S-p_S)\\
 &\qquad{}\times\varphi_j(\xi)\psi_{p_1}(u_1)\cdots\psi_{p_m}(u_m)\cdots\psi_{p_S}(u_S)
\end{align*}
($\delta(n)$ is $1$ for $n=0$, $0$ for $n\ne0$), and we must form $\mathsf{H}_w=V_{l_1}+\cdots+V_{l_S}$.

Together, then, we obtain
\begin{align*}
 &\mathsf{H}\varphi_k(\xi)\psi_{n_1}(u_1)\cdots\psi_{n_S}(u_S)\\
 &\quad=(W_k+E_{n_1}+\cdots+E_{n_S})\varphi_k(\xi)\psi_{n_1}(u_1)\cdots\psi_{n_S}(u_S)\\
 &\qquad{}+\sum_{j=1}^{\infty}\sum_{p_1,\ldots,p_S=0}^{\infty}\sum_{m=1}^{S}
 \delta(n_1-p_1)\cdots V_{kn_m/jp_m}\cdots\delta(n_S-p_S)\\
 &\qquad{}\times\varphi_j(\xi)\psi_{p_1}(u_1)\cdots\psi_{p_S}(u_S).
\end{align*}
A simple rearrangement gives
\begin{align*}
 &\mathsf{H}\varphi_k(\xi)\Phi_{M_0M_1\cdots}(u_1\cdots u_S)\\
 &\quad=\left(W_k+\sum_{n=0}^{\infty}M_nE_n\right)\varphi_k(\xi)\Phi_{M_0M_1\cdots}(u_1\cdots u_S)\\
 &\qquad{}+\sum_{j=1}^{\infty}\sum_{n,p=0}^{\infty}M_nV_{kn/jp}\varphi_j(\xi)
 \Phi_{M_0M_1\cdots M_n-1\cdots M_p+1\cdots}(u_1\cdots u_S)
\end{align*}
(for $n=p$, $\cdots M_n-1\cdots M_p+1\cdots$ is to be replaced by $\cdots M_n\cdots$), and thus, for the orthonormal functions,
\begin{align*}
 &\mathsf{H}\varphi_k(\xi)\psi_{M_0M_1\cdots}(u_1\cdots u_S)\\
 &\quad=\left(W_k+\sum_{n=0}^{\infty}M_nE_n\right)\varphi_k(\xi)\psi_{M_0M_1\cdots}(u_1\cdots u_S)\\
 &\qquad{}+\sum_{j=1}^{\infty}\sum_{n,p=0}^{\infty}\sqrt{M_n(M_p+1-\delta(n-p))}\,V_{kn/jp}\varphi_j(\xi)\\
 &\qquad{}\times\psi_{M_0M_1\cdots M_n-1\cdots M_p+1\cdots}(u_1\cdots u_S).
\end{align*}
We can expand the general $f(\xi,u_1\cdots u_S)$ of $\overline{\mathfrak R}_{\infty}^{(S)}$ in these orthogonal \hypertarget{origpage-149}{}functions:\tnote{The source uses capital $\Psi$ on this page for the normalized functions denoted by lower-case $\psi$ on the preceding two pages. The printed capitalization is retained.}
\[
 f(\xi u_1\cdots u_S)=\sum_{k=1}^{\infty}
 \sum_{\substack{M_0,M_1,\ldots\ge0\\M_0+M_1+\cdots=S}}
 a_{kM_0M_1\cdots}\varphi_k(\xi)\Psi_{M_0M_1\cdots}(u_1\cdots u_S).
\]
Hence $\overline{\mathfrak R}_{\infty}^{(S)}$ can also be regarded as the Hilbert space of sequences $a_{kM_0M_1\cdots}$, $k=1,2,\ldots$, $M_0,M_1,\ldots=0,1,2,\ldots$, $M_0+M_1+\cdots=S$, with finite $\sum_{kM_0M_1\cdots}|a_{kM_0M_1\cdots}|^2$. Here
\[
 \mathsf{H}a_{kM_0M_1\cdots}=a'_{kM_0M_1\cdots}
\]
is defined by
\begin{align*}
 &\mathsf{H}\sum_{k=1}^{\infty}\sum_{\substack{M_0,M_1,\ldots\ge0\\M_0+M_1+\cdots=S}}
 a_{kM_0M_1\cdots}\varphi_k(\xi)\Psi_{M_0M_1\cdots}(u_1\cdots u_S)\\
 &\quad=\sum_{k=1}^{\infty}\sum_{\substack{M_1,M_2,\ldots\ge0\\M_1+M_2+\cdots=S}}
 a'_{kM_0M_1\cdots}\varphi_k(\xi)\Psi_{M_0M_1\cdots}(u_1\cdots u_S),
\end{align*}
so that
\begin{align*}
 &\mathsf{H}a_{kM_0M_1\cdots}=a'_{kM_0M_1\cdots}\\
 &\quad={}\left(W_k+\sum_{n=0}^{\infty}M_nE_n\right)a_{kM_0M_1\cdots}\\
 &{}+\sum_{j=1}^{\infty}\sum_{n,p=0}^{\infty}\sqrt{M_n(M_p+1-\delta(n-p))}\,
 \overline{V_{kn/jp}}\,a_{jM_0M_1\cdots M_n-1\cdots M_p+1\cdots}.
\end{align*}
($k,j$ and $n,p$ have exchanged their roles as compared with the formula for $\varphi_k(\xi)\Psi_{M_0M_1\cdots}(u_1\cdots u_S)$; in place of $V_{jp/kn}$ we have written $\overline{V_{kn/jp}}$ in view of the Hermitian character of $V$.)

We must now prepare for the limiting passage $S\longrightarrow+\infty$. Since $M_0$ is determined by $M_1,M_2,\ldots$, namely $M_0=S-M_1-M_2-\cdots$, we may write $a_{kM_1M_2\cdots}$ instead of $a_{kM_0M_1\cdots}$; the indices are then restricted by the conditions $k=1,2,\ldots$, $M_1,M_2,\ldots=0,1,2,\ldots$, $M_1+M_2+\cdots\le S$. If we observe that $E_0=0$ and introduce the notations
\[
 SV_{k0/j0}=V_{k/j},\qquad \sqrt S\,V_{k0/jn}=V_{k/jn},\qquad
 \sqrt S\,V_{kn/j0}=\overline{V_{j/kn}}
\]
($V_{kn/jp}$ is Hermitian!), we obtain
\begin{align*}
 &\mathsf{H}a_{kM_1M_2\cdots}=a'_{kM_1M_2\cdots}\\
 &\quad={}\left(W_k+\sum_{n=1}^{\infty}M_nE_n\right)a_{kM_1M_2\cdots}\\
 &{}+\sum_{j=1}^{\infty}V_{k/j}a_{jM_1M_2\cdots}\\
 &{}+\sum_{j=1}^{\infty}\sum_{n=1}^{\infty}\sqrt{M_n}
 \sqrt{\frac{S-M_1-M_2-\cdots+1}{S}}\,V_{j/kn}a_{jM_1M_2\cdots M_n-1\cdots}\\
 &{}+\sum_{j=1}^{\infty}\sum_{n=1}^{\infty}\sqrt{M_n+1}
 \sqrt{\frac{S-M_1-M_2-\cdots}{S}}\,\overline{V_{k/jn}}\,a_{jM_1M_2\cdots M_n+1\cdots}\\
 &{}+\sum_{j=1}^{\infty}\sum_{n,p=1}^{\infty}\sqrt{M_n(M_p+1)}\,
 \overline{V_{kn/jp}}\,a_{jM_1M_2\cdots M_n-1\cdots M_p+1\cdots}.
\end{align*}
We can now carry out $S\longrightarrow+\infty$: the $a_{kM_1M_2\cdots}$ are then again to be extended over all sequences $kM_1M_2\cdots$ with $k=1,2,\ldots$, $M_1,M_2,\ldots=0,1,2,\ldots$,
\hypertarget{origpage-150}{}with only finitely (but arbitrarily) many $M_n\ne0$ (cf.\ \hyperlink{note-144}{Note~144}), and $\mathsf{H}$ becomes
\begin{align*}
 &\mathsf{H}a_{kM_1M_2\cdots}=a'_{kM_1M_2\cdots}\\
 &\quad={}\left(W_k+\sum_{n=1}^{\infty}M_nE_n\right)a_{kM_1M_2\cdots}
 +\sum_{j=1}^{\infty}V_{k/j}a_{jM_1M_2\cdots}\\
 &{}+\sum_{j=1}^{\infty}\sum_{n=1}^{\infty}
 \bigl(V_{j/kn}\sqrt{M_n+1}\,a_{jM_1M_2\cdots M_n+1\cdots}\\
 &\qquad{}+\overline{V_{k/jn}}\sqrt{M_n}\,a_{jM_1M_2\cdots M_n-1\cdots}\bigr)\\
 &{}+\sum_{j=1}^{\infty}\sum_{n,p=1}^{\infty}\overline{V_{kn/jp}}\sqrt{M_n(M_p+1)}\,
 a_{jM_1M_2\cdots M_n-1\cdots M_p+1\cdots}.
\end{align*}
The analogy with the equation derived from the electromagnetic theory of light is immediately evident:\tnote{In passing from the preceding finite-$S$ formula to this limit, the source interchanges the $M_n+1$ and $M_n-1$ terms attached to the two coefficients. Both formulas are retained as printed.} to obtain complete agreement we need only set
\[
 E_n=h\varrho_n,\qquad V_{k/j}=0,\qquad V_{k/jn}=w_{jk}^n=\overline{w_{kj}^n},\qquad V_{kn/jp}=0.
\]
We see that the light-quantum conception proves identical with the classical electromagnetic one if
\begin{enumerate}
 \item the latter is rewritten according to the general quantum-mechanical scheme;
 \item the energy of the light quanta is set equal, in accordance with Einstein's rule, to $h$ times their frequency;
 \item the interaction energy of the light quanta with matter is assigned correctly (cf.\ the above expressions for $V$).
\end{enumerate}
This brilliantly resolves one of the gravest paradoxes of the earlier forms of quantum theory: the dual nature of light (electromagnetic waves or discrete corpuscles, light quanta).\origNote{148} Admittedly, it is difficult to find a direct, intuitive interpretation for the interaction energy $V$ between light and matter just calculated, all the more so because its only nonzero matrix elements $V_{kn/jp}$ (those with $n=0,p\ne0$ or $n\ne0,p=0$) depend on the number $S$ of all possible light quanta (they are proportional to $1/\sqrt S$), although ultimately we must take $S\longrightarrow+\infty$. Nevertheless, this can be accepted in view of the fact that every description by means of a model is only an approximation, whereas the exact content of the theory is conveyed precisely by the expression for the operator $\mathsf{H}$.

Let us now return to our actual task, the determination of the \emph{transition probabilities}. In accordance with the time-dependent Schr\"odinger differential equation, the changes of $a_{kM_1M_2\cdots}=a_{kM_1M_2\cdots}(t)$ are determined by
\begin{align*}
 &\frac h{2\pi i}\frac\partial{\partial t}a_{kM_1M_2\cdots}=-\mathsf{H}a_{kM_1M_2\cdots}\\
 &\quad=-\left(W_k+\sum_{n=1}^{\infty}h\varrho_n M_n\right)a_{kM_1M_2\cdots}\\
 &{}-\sum_{j=1}^{\infty}\sum_{n=1}^{\infty}w_{kj}^n
 \left(\sqrt{M_n+1}\,a_{jM_1M_2\cdots M_n+1\cdots}
 +\sqrt{M_n}\,a_{jM_1M_2\cdots M_n-1\cdots}\right).
\end{align*}
\hypertarget{origpage-151}{}%
Since the principal change in the $a_{kM_1M_2\cdots}$ is caused by the first summand of this expression, it is expedient to separate it off by setting
\[
 a_{kM_1M_2\cdots}(t)=e^{-\frac{2\pi i}{h}(W_k+\sum_{n=1}^{\infty}h\varrho_n M_n)t}\,b_{kM_1M_2\cdots}(t).
\]
We then obtain\tnote{The source prints a minus sign between the two terms in this differential equation, although the preceding equation has a plus sign.}
\begin{align*}
 \frac\partial{\partial t}b_{kM_1M_2\cdots}
 ={}&-\frac{2\pi i}{h}\sum_{j=1}^{\infty}\sum_{n=1}^{\infty}w_{kj}^n
 \bigl(e^{-\frac{2\pi i}{h}(W_j-W_k+h\varrho_n)t}\sqrt{M_n+1}\,b_{jM_1M_2\cdots M_n+1\cdots}\\
 &\qquad{}-e^{-\frac{2\pi i}{h}(W_j-W_k-h\varrho_n)t}\sqrt{M_n}\,b_{jM_1M_2\cdots M_n-1\cdots}\bigr).
\end{align*}
The physical meaning of the $a_{kM_1M_2\cdots}$ and $b_{kM_1M_2\cdots}$ is fixed by their origin. For finite $\overline M_0+\overline M_1+\overline M_2+\cdots=S$, $\varphi_{\bar k}(\xi)\Psi_{\overline M_0\overline M_1\cdots}(u_1\cdots u_S)$ was the state in which $\boldsymbol{S}$ is in the $\bar k$th quantum orbit, with respectively $\overline M_0,\overline M_1,\overline M_2,\ldots$ light quanta in the respective states $\psi_0,\psi_1,\psi_2,\ldots$; that is, $\overline M_0$ in the state of ``nonexistence,'' and $\overline M_1,\overline M_2,\ldots$ in the states belonging to the respective normal oscillations $\overline{\mathfrak A}_1,\overline{\mathfrak A}_2,\ldots$. The $a_{kM_1M_2\cdots}$ belonging to this wave function are thus\tnote{The source prints $M_0$ in the first occupation-number factor of the following formula. The same configuration is represented by $M_1$ in the probability formula below and in the later initial condition. The printed index is retained.}
\[
 a_{kM_1M_2\cdots}=\delta(k-\bar k)\delta(M_0-\overline M_1)\delta(M_2-\overline M_2)\cdots.
\]
(Only finitely many factors differ from $1$, since, with finitely many exceptions, $M_n=\overline M_n=0$ always holds.) This must of course be retained after $S\longrightarrow\infty$. For an arbitrary state $a_{kM_1M_2\cdots}$ of $\boldsymbol{S}+\boldsymbol{L}$, the configuration just described therefore has the probability (if it is measured; cf.\ what was said about the simple pure point spectrum in \hyperref[sec:III.3]{III.3})
\begin{align*}
 &\left|\sum_{k,M_1,M_2,\ldots}a_{kM_1M_2\cdots}
 \delta(k-\bar k)\delta(M_1-\overline M_1)\delta(M_2-\overline M_2)\cdots\right|^2\\
 &\qquad=|a_{\bar k\overline M_1\overline M_2\cdots}|^2
 =|b_{\bar k\overline M_1\overline M_2\cdots}|^2.
\end{align*}
In particular, the total probability that $\boldsymbol{S}$ is in the $\bar k$th quantum orbit is
\[
 \vartheta_{\bar k}=\sum_{\overline M_1,\overline M_2,\ldots}
 |b_{\bar k\overline M_1\overline M_2\cdots}|^2.
\]
Let the atom at the initial instant ($t=0$) be in the $\bar k$th state, with respectively $\overline M_1,\overline M_2,\ldots$ light quanta in the states $\overline{\mathfrak A}_1,\overline{\mathfrak A}_2,\ldots$ present; i.e.,
\[
 b_{kM_1M_2\cdots}=a_{kM_1M_2\cdots}
 =\delta(k-\bar k)\delta(M_1-\overline M_1)\delta(M_2-\overline M_2)\cdots.
\]
According to the above differential equation, in a first approximation (i.e., for times $t$ so short that the right-hand sides may still be regarded as
\hypertarget{origpage-152}{}constants), the only $\frac\partial{\partial t}b_{kM_1M_2\cdots}$ that differ from $0$ will be those for which either $M_1,M_2,\ldots,M_n+1,\ldots$ or $M_1,M_2,\ldots,M_n-1,\ldots$ coincides with $\overline M_1,\overline M_2,\ldots$: i.e., all $k,\overline M_1,\overline M_2,\ldots,\overline M_n\pm1,\ldots$. For these, integration gives\tnote{The denominator of the second source formula contains $-h\varrho_n$, while its exponential contains $+h\varrho_n$; the subsequent absorption probability uses the latter sign.}
\begin{align*}
 b_{k\overline M_1\overline M_2\cdots\overline M_n+1\cdots}
 &=w_{\bar k k}^n\frac{1-e^{-\frac{2\pi i}{h}(W_{\bar k}-W_k-h\varrho_n)t}}
 {W_{\bar k}-W_k-h\varrho_n}\sqrt{\overline M_n+1},\\
 b_{k\overline M_1\overline M_2\cdots\overline M_n-1\cdots}
 &=w_{\bar k k}^n\frac{1-e^{-\frac{2\pi i}{h}(W_{\bar k}-W_k+h\varrho_n)t}}
 {W_{\bar k}-W_k-h\varrho_n}\sqrt{\overline M_n}.
\end{align*}
All other $b_{kM_1M_2\cdots}$ are $0$ in this approximation. (Except for $b_{\bar k\overline M_1\overline M_2\cdots}$, which should, in this approximation, i.e., apart from terms in $t^2$, equal its initial value $1$. However, the conclusion $\frac\partial{\partial t}b_{\bar k\overline M_1\overline M_2\cdots}=0$ becomes questionable because the right-hand side of our differential equation contains, in this case, infinitely many $b_{k\overline M_1\overline M_2\cdots\overline M_n\pm1\cdots}$ that do not vanish in our approximation. Thus the smallness of each of these summands for small $t$ does not yet allow us to infer the smallness of their sum. Indeed, calculation of the next approximation would show that $b_{\bar k\overline M_1\overline M_2\cdots}$ differs from $1$ by terms in $t$, not in $t^2$.\origNote{149} However, since
\[
 \sum_{k,M_1,M_2,\ldots}|b_{kM_1M_2\cdots}|^2
 =\sum_{k,M_1,M_2,\ldots}|a_{kM_1M_2\cdots}|^2=1,
\]
and hence
\[
 |b_{\bar k\overline M_1\overline M_2\cdots}|^2
 =1-\sum_{kM_1M_2\cdots\ne\bar k\overline M_1\overline M_2\cdots}|b_{kM_1M_2\cdots}|^2,
\]
a direct determination of this $b_{\bar k\overline M_1\overline M_2\cdots}$ is unnecessary.)

The qualitative nature of the process is clearly recognizable here: a $b_{k\overline M_1\overline M_2\cdots\overline M_n+1\cdots}$, corresponding to the emission of an $\overline{\mathfrak A}_n$ light quantum (of frequency $\varrho_n$), grows larger as the denominator $W_{\bar k}-W_k-h\varrho_n$ becomes smaller, i.e., as the natural frequency $\varrho_n$ of light approaches the ``Bohr frequency'' $(W_{\bar k}-W_k)/h$.\origNote{150} Likewise, the $b_{k\overline M_1\overline M_2\cdots\overline M_n-1\cdots}$ corresponding to absorption grows larger as $\varrho_n$ approaches $(W_k-W_{\bar k})/h$. We see that Bohr's frequency relation does not hold exactly (not all frequencies are available among the $\varrho_n$), but it does hold with overwhelming probability when the time $t$ is short and the $\varrho_n$ lie very densely (as is the case for a large cavity $\boldsymbol H$). Moreover, the $w_{\bar k k}^n$ also increase the rates at which these processes occur; they will soon be identifiable with the transition probabilities.

\hypertarget{origpage-153}{}%
Our formulas for $b_{k\overline M_1\overline M_2\cdots\overline M_n\pm1\cdots}$ further give\origNote{151}
\begin{align*}
 |b_{k\overline M_1\overline M_2\cdots\overline M_n+1\cdots}|^2
 &=\frac2{h^2}(\overline M_n+1)|w_{\bar k k}^n|^2
 \frac{1-\cos 2\pi\left(\varrho_n-\frac{W_{\bar k}-W_k}{h}\right)t}
 {\left(\varrho_n-\frac{W_{\bar k}-W_k}{h}\right)^2},
\end{align*}
and\origNote{151}
\begin{align*}
 |b_{k\overline M_1\overline M_2\cdots\overline M_n-1\cdots}|^2
 &=\frac2{h^2}\overline M_n|w_{\bar k k}^n|^2
 \frac{1-\cos 2\pi\left(\varrho_n-\frac{W_k-W_{\bar k}}{h}\right)t}
 {\left(\varrho_n-\frac{W_k-W_{\bar k}}{h}\right)^2},\\
 |b_{kM_1M_2\cdots}|^2&=0
 \quad\text{for }kM_1M_2\cdots\ne\bar k\overline M_1\overline M_2\cdots,\;
 k\overline M_1\overline M_2\cdots\overline M_n\pm1\cdots.
\end{align*}
This gives, for $\vartheta_k$, $k\ne\bar k$,
\begin{align*}
 \vartheta_k={}&\sum_{n=1}^{\infty}\frac2{h^2}(\overline M_n+1)|w_{\bar k k}^n|^2
 \frac{1-\cos 2\pi\left(\varrho_n-\frac{W_{\bar k}-W_k}{h}\right)t}
 {\left(\varrho_n-\frac{W_{\bar k}-W_k}{h}\right)^2}\\
 &{}+\sum_{n=1}^{\infty}\frac2{h^2}\overline M_n|w_{\bar k k}^n|^2
 \frac{1-\cos 2\pi\left(\varrho_n-\frac{W_k-W_{\bar k}}{h}\right)t}
 {\left(\varrho_n-\frac{W_k-W_{\bar k}}{h}\right)^2}.
\end{align*}
(The first sum corresponds to emissions, the second to absorptions.) To express these $\vartheta_k$ in closed form, we must now make simplifying assumptions: on the one hand, we assume $\boldsymbol H$ very large (i.e., its volume $\mathcal V\longrightarrow\infty$); on the other, we treat the normal oscillations $\overline{\mathfrak A}_n$ of $\boldsymbol H$ statistically. For this purpose, in each of the two sums we collect all terms belonging to $\varrho_n$ between $\varrho$ and $\varrho+d\varrho$ (we substitute the value of $w_{\bar k k}^n$, and assume $d\varrho\ll\varrho$):
\begin{multline*}
 \frac1{4\pi^2c^2h\varrho}
 \left[\sum_{\varrho\le\varrho_n<\varrho+d\varrho}
 \left|\sum_{\nu=1}^{l}\frac{e_\nu}{m_\nu}
 \bigl(P_\nu^x\overline{\mathfrak A}_{n,x}(Q_\nu^x,Q_\nu^y,Q_\nu^z)+\cdots\bigr)_{k\bar k}\right|^2
 (\overline M_n+1)\right]\\
 {}\times\frac{1-\cos2\pi\left(\varrho-\frac{W_{\bar k}-W_k}{h}\right)t}
 {\left(\varrho-\frac{W_{\bar k}-W_k}{h}\right)^2},
\end{multline*}
and the same expression again, but with $\overline M_n$ in place of $\overline M_n+1$, and $(W_k-W_{\bar k})/h$ in place of $(W_{\bar k}-W_k)/h$. We must now evaluate the expression in square brackets.

Usually, however, the description gives much less information than the individual values of $M_1,M_2,\ldots$: it specifies only the intensities. That is, the light energy per unit volume in the spectral interval from $\varrho$ to
\hypertarget{origpage-154}{}$\varrho+d\varrho$ is $I(\varrho)d\varrho$. This means
\begin{gather*}
 \sum_{\varrho\le\varrho_n<\varrho+d\varrho}h\varrho_n\overline M_n
 \approx h\varrho\sum_{\varrho\le\varrho_n<\varrho+d\varrho}\overline M_n
 =\mathcal V I(\varrho)d\varrho,\\
 \sum_{\varrho\le\varrho_n<\varrho+d\varrho}\overline M_n
 =\frac{\mathcal V I(\varrho)}{h\varrho}\,d\varrho.
\end{gather*}
According to a generally valid asymptotic formula of Weyl (cf.\ the work cited in \hyperlink{note-140}{Note~140}), the number of $\varrho_n$ in the interval $\varrho\le\varrho_n<\varrho+d\varrho$ is $\frac{8\pi\mathcal V\varrho^2}{c^2}d\varrho$; hence\tnote{The source prints $c^2$ in Weyl\textquotesingle s mode-counting formula, but $c^3$ in the expression immediately following it. Both powers are retained.}
\[
 \sum_{\varrho\le\varrho_n<\varrho+d\varrho}(\overline M_n+1)
 \approx\frac{\mathcal V\left(I(\varrho)+\frac{8\pi h\varrho^3}{c^3}\right)}{h\varrho}\,d\varrho.
\]
Thus, if
\[
 \left|\sum_{\nu=1}^{l}\frac{e_\nu}{m_\nu}
 \bigl(P_\nu^x\overline{\mathfrak A}_{n,x}(Q_\nu^x,Q_\nu^y,Q_\nu^z)+\cdots\bigr)_{k\bar k}\right|^2
\]
undergoes sufficiently rapid fluctuations in the interval $\varrho\le\varrho_n<\varrho+d\varrho$ about a mean value, which we shall call $w_{k\bar k}(\varrho)$, the bracketed expressions obtained above have the values
\[
 w_{k\bar k}(\varrho)\frac{\mathcal V\left(I(\varrho)+\frac{8\pi h\varrho^3}{c^3}\right)}{h\varrho}\,d\varrho
 \quad\text{and}\quad
 w_{k\bar k}(\varrho)\frac{\mathcal V I(\varrho)}{h\varrho}\,d\varrho.
\]
If we further write $\nu_{\bar k k}$ for $(W_{\bar k}-W_k)/h$ and $\nu_{k\bar k}$ for $(W_k-W_{\bar k})/h$, we obtain for our sums
\begin{align*}
 \vartheta_k=\frac{\mathcal V}{4\pi^2c^2h^2}\int_0^{\infty}
 \biggl\{&\left(I(\varrho)+\frac{8\pi h}{c^3}\varrho^3\right)
 \frac{1-\cos2\pi(\varrho-\nu_{\bar k k})t}{(\varrho-\nu_{\bar k k})^2}\\
 &{}+I(\varrho)\frac{1-\cos2\pi(\varrho-\nu_{k\bar k})t}{(\varrho-\nu_{k\bar k})^2}
 \biggr\}\frac{w_{k\bar k}(\varrho)}{\varrho^2}\,d\varrho.
\end{align*}
For small $t$, this integral evidently has order of magnitude $t^2$ (because $1-\cos2\pi ct$ does), except in those parts of the domain of integration where the denominators $(\varrho-\nu_{\bar k k})^2$ or $(\varrho-\nu_{k\bar k})^2$ are small. Here contributions large compared with $t^2$ may arise; if this is the case, these contributions give the asymptotic expression for $\vartheta_k$. It will indeed turn out to be so, for we shall obtain contributions of order $t$; these are therefore what we must determine.

Since $\nu_{\bar k k}=-\nu_{k\bar k}=(W_{\bar k}-W_k)/h$, only the denominator of the first term becomes small for $W_{\bar k}>W_k$, and only that of the second term for $W_{\bar k}<W_k$. We therefore retain only the first or second term for $W_{\bar k}>W_k$ or $W_{\bar k}<W_k$, respectively. Furthermore, write $\bar\nu_{k\bar k}=|W_{\bar k}-W_k|/h$ for the relevant positive frequency, $\nu_{\bar k k}$ or $\nu_{k\bar k}$. Values of $\varrho$ farther from this frequency give only $t^2$ contributions to the integral. Hence we can replace the remaining factor in the \hypertarget{origpage-155}{}integrand by its value for $\varrho=\bar\nu_{k\bar k}$, i.e., $I w_{k\bar k}(\bar\nu_{k\bar k})/\bar\nu_{k\bar k}^{2}$, where $I=I(\bar\nu_{k\bar k})+\frac{8\pi h}{c^3}\bar\nu_{k\bar k}^{3}$ or $I=I(\bar\nu_{k\bar k})$, respectively. Thus
\[
 \vartheta_k=\frac{\mathcal V I w_{k\bar k}(\bar\nu_{k\bar k})}{4\pi^2c^2h^2\bar\nu_{k\bar k}^{2}}
 \int_0^{\infty}\frac{1-\cos2\pi(\varrho-\bar\nu_{k\bar k})t}{(\varrho-\bar\nu_{k\bar k})^2}\,d\varrho.
\]
Since the added interval again contributes only a term of order $t^2$, we may replace $\int_0^{\infty}$ by $\int_{-\infty}^{\infty}$ and then introduce the new variable $x=2\pi(\varrho-\bar\nu_{k\bar k})t$. Since\origNote{152}
\[
 \int_{-\infty}^{\infty}\frac{1-\cos2\pi(\varrho-\bar\nu_{k\bar k})t}{(\varrho-\bar\nu_{k\bar k})^2}\,d\varrho
 =2\pi t\int_{-\infty}^{\infty}\frac{1-\cos x}{x^2}\,dx=2\pi^2t,
\]
we finally have
\[
 \vartheta_k=\frac{\mathcal V I w_{k\bar k}(\bar\nu_{k\bar k})}{2h^2\bar\nu_{k\bar k}^{2}}\,t,
\]
which also proves that $\vartheta_k$ is of order $t$.\tnote{The factor $c^2$ present in the preceding denominator is absent in this source formula.}

To determine $w_{k\bar k}(\bar\nu_{k\bar k})$, we must find an expression free of $\overline{\mathfrak A}_n$ for
\[
 \left|\sum_{\nu=1}^{l}\frac{e_\nu}{m_\nu}
 \bigl(P_\nu^x\overline{\mathfrak A}_{n,x}(Q_\nu^x,Q_\nu^y,Q_\nu^z)+\cdots\bigr)_{k\bar k}\right|^2.
\]
We obtain it by replacing $\overline{\mathfrak A}_n$, in view of its rapid fluctuations, by a randomly oriented vector of constant length. Because it is constant, i.e., independent of $Q_\nu^x,Q_\nu^y,Q_\nu^z$, it is therefore a numerical vector times the matrix $1$. The constant length $\gamma_n$ is to be determined from the normalization
\[
 \iiint_{\boldsymbol H}[\overline{\mathfrak A}_n,\overline{\mathfrak A}_n]\,dx\,dy\,dz=4\pi c^2.
\]
Thus
\[
 \mathcal V\gamma_n^2=4\pi c^2,\qquad \gamma_n^2=\frac{4\pi c^2}{\mathcal V}.
\]
The squared $x$-component $\overline{\mathfrak A}_{n,x}^{2}$ contributes, on average, one third of $[\overline{\mathfrak A}_n,\overline{\mathfrak A}_n]=\overline{\mathfrak A}_{n,x}^{2}+\overline{\mathfrak A}_{n,y}^{2}+\overline{\mathfrak A}_{n,z}^{2}=\gamma_n^2$, so its mean is $\frac13\gamma_n^2=\frac{4\pi c^2}{3\mathcal V}$; likewise for $\overline{\mathfrak A}_{n,y}^{2}$ and $\overline{\mathfrak A}_{n,z}^{2}$. Consequently,
\begin{align*}
 w_{k\bar k}(\varrho)
 &=\mathop{\operatorname{mean}}_{\varrho\le\varrho_n<\varrho+d\varrho}
 \left|\sum_{\nu=1}^{l}\frac{e_\nu}{m_\nu}
 \bigl(P_\nu^x\overline{\mathfrak A}_{n,x}(Q_\nu^x,Q_\nu^y,Q_\nu^z)+\cdots\bigr)_{k\bar k}\right|^2\\
 &\approx\frac{4\pi c^2}{3\mathcal V}
 \left(\left|\left(\sum_{\nu=1}^{l}\frac{e_\nu}{m_\nu}P_\nu^x\right)_{k\bar k}\right|^2+\cdots\right).
\end{align*}

\hypertarget{origpage-156}{}%
Since $\mathsf{H}_0$, the energy of the system $\boldsymbol{S}$ considered by itself, equals kinetic energy plus potential energy, and thus has the form
\[
 \mathsf{H}_0=\sum_{\nu=1}^{l}\frac1{2m_\nu}\bigl((P_\nu^x)^2+(P_\nu^y)^2+(P_\nu^z)^2\bigr)
 +V(Q_1^x,Q_1^y,Q_1^z,\ldots,Q_l^x,Q_l^y,Q_l^z),
\]
we have\origNote{153}
\[
 \mathsf{H}_0Q_\nu^x-Q_\nu^x\mathsf{H}_0=\frac h{2\pi i}\frac1{m_\nu}P_\nu^x.
\]
And since $\mathsf{H}_0$ is a diagonal matrix with diagonal elements $W_1,W_2,\ldots$ [$(\mathsf{H}_0)_{kj}=W_k\delta_{kj}$], it follows for the matrix elements that
\begin{align*}
 (P_\nu^x)_{k\bar k}
 &=\frac{2\pi i m_\nu}{h}(\mathsf{H}_0Q_\nu^x-Q_\nu^x\mathsf{H}_0)_{k\bar k}
 =\frac{2\pi i m_\nu}{h}(W_k-W_{\bar k})(Q_\nu^x)_{k\bar k}\\
 &=\pm i\,2\pi m_\nu\bar\nu_{k\bar k}(Q_\nu^x)_{k\bar k}.
\end{align*}
Thus\tnote{The source retains an additional $h^2$ in the denominator of this expression for $w_{k\bar k}(\varrho)$, despite the cancellation of $h$ in the preceding momentum identity.}
\[
 w_{k\bar k}(\varrho)=\frac{16\pi^3c^2}{3\mathcal Vh^2}\bar\nu_{k\bar k}^{2}
 \left(\left|\left(\sum_{\nu=1}^{l}e_\nu Q_\nu^x\right)_{k\bar k}\right|^2+\cdots\right).
\]
Substitution in the formula for $\vartheta_k$ gives
\[
 \vartheta_k=\frac{8\pi^3}{3h^2}
 \left(\left|\left(\sum_{\nu=1}^{l}e_\nu Q_\nu^x\right)_{k\bar k}\right|^2+\cdots\right)It.
\]
If we also set $w_{k\bar k}=\left|\left(\sum_{\nu=1}^{l}e_\nu Q_\nu^x\right)_{k\bar k}\right|^2+\cdots$, this result is evidently to be interpreted as follows: the atom $\boldsymbol{S}$, in the $k$th state, makes the following transitions (quantum jumps).
\begin{enumerate}
 \item It passes to higher states $\bar k$ ($W_{\bar k}>W_k$), doing so
 \[
 \frac{8\pi^3}{3h^2}w_{k\bar k}I\left(\frac{W_{\bar k}-W_k}{h}\right)
 \]
 times per second; i.e., in proportion to the intensity of the radiation field at the corresponding Bohr frequency, $(W_{\bar k}-W_k)/h$.
 \item It passes to lower states $\bar k$ ($W_{\bar k}<W_k$), doing so
 \[
 \frac{8\pi^3}{3h^2}w_{k\bar k}I\left(\frac{W_k-W_{\bar k}}{h}\right)
 \]
 times per second; i.e., in proportion to the intensity of the radiation field at the corresponding Bohr frequency, $(W_k-W_{\bar k})/h$.
 \item It also passes to lower states $\bar k$ ($W_{\bar k}<W_k$), doing so
 \[
 \frac{64\pi^4}{3hc^3}w_{k\bar k}\left(\frac{W_k-W_{\bar k}}{h}\right)^3
 \]
 times per second; i.e., wholly independently of the radiation field present.
\end{enumerate}
Item~1 corresponds to absorption processes from the radiation field; item~2 to emission processes induced by the radiation field; but item~3 to spontaneous emission processes, which the atom will always perform until it has found permanent rest in its lowest stationary state (minimum $W_k$!).

The three transition mechanisms 1--3 had already been discovered thermodynamically by Einstein before the discovery of quantum \hypertarget{origpage-157}{}mechanics;\origNote{154} only the value of the ``transition probability'' $w_{k\bar k}$ was missing. As we have already mentioned, the above value
\[
 w_{k\bar k}=\left|\left(\sum_{\nu=1}^{l}e_\nu Q_\nu^x\right)_{k\bar k}\right|^2
 +\left|\left(\sum_{\nu=1}^{l}e_\nu Q_\nu^y\right)_{k\bar k}\right|^2
 +\left|\left(\sum_{\nu=1}^{l}e_\nu Q_\nu^z\right)_{k\bar k}\right|^2
\]\nopagebreak[4]
is contained in Heisenberg's first attempt at an interpretation. We have now justified it (following Dirac) on the basis of the general theory.


% Consolidated source: chapters-IV-VI.tex
% Source PDF page 164 / original page 157. The page anchor is in the preceding chapter.
\chapter{Deductive Construction of the Theory}\label{chap:IV}
\section{Foundations of the Statistical Theory}\label{sec:IV.1}

In Part~\hyperref[chap:III]{III} we succeeded in reducing all the assertions of quantum mechanics to the statistical formula (there denoted by \hyperref[post:III.1.E2]{$E_2$})
\begin{equation}\tag{$\overline E$}\label{eq:IV.1.E}
\operatorname{Erw}(\mathfrak R,\varphi)=(R\varphi,\varphi).
\end{equation}
Here $\operatorname{Erw}(\mathfrak R,\varphi)$ is the expectation value of the quantity $\mathfrak R$ in the state $\varphi$, and $R$ is the operator corresponding to $\mathfrak R$. We shall now show how this formula itself can be derived from a few general qualitative assumptions. At the same time, the entire construction of quantum mechanics on this basis, carried out in Part~\hyperref[chap:III]{III}, will be examined once more. Before doing so, however, we must make one further remark.

In the state $\varphi$, the quantity $\mathfrak R$ has the expectation value $\varrho=(R\varphi,\varphi)$ and, as its dispersion $\varepsilon^2$, the expectation value of the quantity $(\mathfrak R-\varrho)^2$, namely
\[
((R-\varrho\mathbf1)^2\varphi,\varphi)=\|R\varphi\|^2-(R\varphi,\varphi)^2.
\]
See note~\hyperlink{note-130}{130}; everything here is calculated from~\hyperref[eq:IV.1.E]{$\overline E$}. This dispersion is in general positive; it vanishes only when $R\varphi=\varrho\varphi$, as in~\hyperref[sec:III.3]{III.3}. Thus, as we have observed several times, even a uniform state $\varphi$ gives only a statistical description. This statistical character becomes even more pronounced if we do not even know which state is actually present: for example, several states $\varphi_1,\varphi_2,\ldots$ may be possible, with respective probabilities $w_1,w_2,\ldots$, where $w_1\geq0,w_2\geq0,\ldots$ and $w_1+w_2+\cdots=1$. Then, according to the universally valid rules of probability theory, the expectation value of $\mathfrak R$ is
\[
\varrho'=\sum_nw_n(R\varphi_n,\varphi_n).
\]
Now, in general, $(R\varphi,\varphi)=\operatorname{Spur}(P_{[\varphi]}R)$. For if we choose a complete orthonormal system $\psi_1,\psi_2,\ldots$ with $\psi_1=\varphi$, so that $\psi_2,\psi_3,\ldots$ are orthogonal to $\varphi$, then
\[
P_{[\varphi]}\psi_n=\begin{cases}\varphi,&n=1,\\0,&\text{otherwise},\end{cases}
\]
and hence
\begin{align*}
\operatorname{Spur}(P_{[\varphi]}R)
&=\sum_{m,n}(P_{[\varphi]}\psi_n,\psi_m)(R\psi_m,\psi_n)\\
&=\sum_n(\varphi,\psi_m)(R\psi_m,\varphi)=(R\varphi,\varphi).
\end{align*}

% Source PDF page 165 / original page 158
\hypertarget{origpage-158}{}%
Thus our $\varrho'$ is $\operatorname{Spur}(\{\sum_nw_nP_{[\varphi_n]}\}R)$.\tnote{The source uses $n$ as the summation index in the second line, although the summand is indexed by $m$. The printed index is retained.} The operator
\[
U=\sum_nw_nP_{[\varphi_n]}
\]
is definite, because all the $P_{[\varphi_n]}$ are definite and $w_n\geq0$. Since $\operatorname{Spur}P_{[\varphi_n]}=1$, its trace is $\sum_nw_n=1$. Moreover, it completely characterizes the statistical properties of the mixture of states just described:
\[
\varrho'=\operatorname{Spur}(UR).
\]
Noting that, in addition to states, we shall also have to consider these mixtures, we turn to the general investigation.

Let us forget the whole of quantum mechanics and retain the following. We are given a system $\boldsymbol{S}$\origNote{155}, characterized for the experimenter by specifying all the quantities that can actually be measured on it and their functional relations to one another. A quantity is properly understood as an instruction for measuring it, together with an instruction for reading or calculating its value from the positions of the pointers of the measuring instruments. If $\mathfrak R$ is a quantity and $f(x)$ any function, the quantity $f(\mathfrak R)$ is defined as follows: to measure $f(\mathfrak R)$, measure $\mathfrak R$; if the value $a$ is found for $\mathfrak R$, then $f(\mathfrak R)$ has the value $f(a)$. As we see, all the quantities $f(\mathfrak R)$, with $\mathfrak R$ fixed and $f(x)$ arbitrary, are in this way measured at once, together with one another and with $\mathfrak R$: this is a first example of simultaneously measurable quantities.

In general, we call two or more quantities $\mathfrak R,\mathfrak S$ simultaneously measurable if there is an arrangement that measures both at the same time on the same system, their respective values merely being calculated in different ways from the readings. In classical mechanics, as is well known, all quantities are simultaneously measurable; in quantum mechanics this is not so, as we saw in~\hyperref[sec:III.3]{III.3}. For such quantities, and for a function $f(x,y)$ of two variables, we can also define the quantity $f(\mathfrak R,\mathfrak S)$: it is measured by measuring $\mathfrak R,\mathfrak S$ simultaneously, and if the values $a,b$ are found, the value of $f(\mathfrak R,\mathfrak S)$ is $f(a,b)$. We must realize, however, that it is entirely meaningless to try to form $f(\mathfrak R,\mathfrak S)$ when $\mathfrak R,\mathfrak S$ are not simultaneously measurable: there is then no way of specifying the corresponding measuring arrangement.

The investigation of physical quantities need not be confined to a single object $\boldsymbol{S}$, especially when it is doubtful whether several quantities can be measured simultaneously. In such cases we are led to consider large statistical ensembles consisting of many systems $\boldsymbol{S}_1,\ldots,\boldsymbol{S}_N$, that is, $N$ copies of $\boldsymbol{S}$, with $N$ large\origNote{156}. On such an ensemble $[\boldsymbol{S}_1,\ldots,\boldsymbol{S}_N]$ one naturally measures not the ``value'' of a quantity $\mathfrak R$ but its distribution of values: for each interval $a'<a\leq a''$,
% Source PDF page 166 / original page 159
\hypertarget{origpage-159}{}%
with given $a',a''$, $a'\leq a''$, one determines the number of systems among $\boldsymbol{S}_1,\ldots,\boldsymbol{S}_N$ for which the value of $\mathfrak R$ lies in that interval. Dividing this number by $N$ gives the probability function $w(a',a'')=w(a'')-w(a')$\origNote{157}. The great advantage of considering such ensembles is that:

\leavevmode\phantomsection\textbf{1.}\label{advantage:IV.1.1} Even if measuring a quantity $\mathfrak R$ considerably alters the measured system $\boldsymbol{S}$---as indeed happens in quantum mechanics, and as we saw in~\hyperref[sec:III.4]{III.4} must in principle happen in the physics of elementary processes, since the intervention of measurement is of the same order of magnitude as the system or its observed parts---a statistical determination of the probability distribution of $\mathfrak R$ in the ensemble $[\boldsymbol{S}_1,\ldots,\boldsymbol{S}_N]$ nevertheless changes the ensemble arbitrarily little, provided $N$ is sufficiently large.

\leavevmode\phantomsection\textbf{2.}\label{advantage:IV.1.2} Even if two or more quantities $\mathfrak R,\mathfrak S$ are not simultaneously measurable on the individual system $\boldsymbol{S}$, their probability distributions can nevertheless be determined simultaneously, to arbitrary accuracy, in one and the same ensemble $[\boldsymbol{S}_1,\ldots,\boldsymbol{S}_N]$, provided $N$ is sufficiently large.

For an ensemble consisting of $N$ elements, it is sufficient to perform the statistical observations concerning the distribution of values of $\mathfrak R$ not on all $N$ elements $\boldsymbol{S}_1,\ldots,\boldsymbol{S}_N$, but on any subsystem of $M\leq N$ elements, say $[\boldsymbol{S}_1,\ldots,\boldsymbol{S}_M]$, provided both $M,N$ are large; $M$ may nevertheless be small compared with $N$\origNote{158}. The changes caused by measurement then affect only the fraction $M/N$ of the ensemble, and hence an arbitrarily small fraction if $M/N$ is chosen sufficiently small. For sufficiently large $N$, this is possible even with large $M$, as asserted in \hyperref[advantage:IV.1.1]{\textbf{1.}} To measure two or more quantities $\mathfrak R,\mathfrak S$ simultaneously, we need two subsystems, say $[\boldsymbol{S}_1,\ldots,\boldsymbol{S}_M]$ and $[\boldsymbol{S}_{M+1},\ldots,\boldsymbol{S}_{2M}]$, with $2M\leq N$, the first being used to determine the statistics of $\mathfrak R$ and the second those of $\mathfrak S$. The two measurements then do not interfere with each other, although they are made on the same ensemble $[\boldsymbol{S}_1,\ldots,\boldsymbol{S}_N]$. Indeed, they change only an arbitrarily small fraction of it if $2M/N$ is sufficiently small; for sufficiently large $N$, this is possible even with large $M$, as asserted in \hyperref[advantage:IV.1.2]{\textbf{2.}}

We see that the use of statistical ensembles, that is, of methods of probability theory, is called for by the possibility that measurement influences the individual system and by the possible failure of simultaneous measurability for several quantities. A general theory must take these circumstances into account: their occurrence in elementary processes was always to be expected\origNote{159}, and today a precise analysis has made it unquestionable; see~\hyperref[sec:III.4]{III.4}. Statistical ensembles remove these difficulties
% Source PDF page 167 / original page 160
\hypertarget{origpage-160}{}%
and thereby make an objective description possible once more: a description independent both of chance and of whether, in a given state, one measures one or the other of two quantities that cannot be measured simultaneously.

In such ensembles it is not surprising if a physical quantity $\mathfrak R$ has no sharp value: its distribution of values need not consist of a single value $a_0$\origNote{160}; several values or intervals of values may instead be possible, with a positive dispersion\textsuperscript{\hyperlink{note-160}{160}}. Nevertheless, two different reasons for this behavior are conceivable:

\leavevmode\phantomsection\textbf{I.}\label{case:IV.1.I} The individual systems $\boldsymbol{S}_1,\ldots,\boldsymbol{S}_N$ of our ensemble may be in different states, so that the ensemble $[\boldsymbol{S}_1,\ldots,\boldsymbol{S}_N]$ is defined by their relative frequencies. Our failure to obtain sharp values for the physical quantities is then due to our ignorance: we do not know the state of the system being measured, and therefore cannot say what the result will be.

\leavevmode\phantomsection\textbf{II.}\label{case:IV.1.II} All the individual systems $\boldsymbol{S}_1,\ldots,\boldsymbol{S}_N$ are in the same state, but the laws of nature are not causal. The dispersions are then caused not by our ignorance, but by nature itself, which has disregarded the ``principle of sufficient reason.''

Case~\hyperref[case:IV.1.I]{\textbf{I}} is universally familiar; case~\hyperref[case:IV.1.II]{\textbf{II}}, on the other hand, is important and new. Admittedly, one will at first be skeptical about the possibility of its occurrence, but we shall find an objective criterion that allows us to decide whether or not it occurs. At first sight, one might raise serious objections to its conceivability and meaningfulness. We do not believe these objections to be sound, however, and \hyperref[case:IV.1.II]{\textbf{II}} is the only way out of certain difficulties, for example in quantum mechanics. We therefore turn to a discussion of the conceptual difficulties of~\hyperref[case:IV.1.II]{\textbf{II}}.

The following objection may be raised against~\hyperref[case:IV.1.II]{\textbf{II}}: nature cannot violate the ``principle of sufficient reason,'' that is, causality, because this is really more a definition of sameness. In other words, the assertion that two identical objects $\boldsymbol{S}_1,\boldsymbol{S}_2$---two copies of the system $\boldsymbol{S}$ in the same state---will behave alike under every conceivable intervention is true because it says nothing. For if $\boldsymbol{S}_1,\boldsymbol{S}_2$ behaved differently under the same intervention, for instance if measuring a quantity $\mathfrak R$ gave different values for them, they would not be called identical. Thus, in an ensemble $[\boldsymbol{S}_1,\ldots,\boldsymbol{S}_N]$ exhibiting dispersion with respect to a quantity $\mathfrak R$, the individual systems $\boldsymbol{S}_1,\ldots,\boldsymbol{S}_N$ cannot, by definition, all be in the same state. Applied to quantum mechanics, this argument would run as follows: measuring the same quantity $\mathfrak R$ on several systems, all in the state with wave function $\varphi$, gives different
% Source PDF page 168 / original page 161
\hypertarget{origpage-161}{}%
values whenever $\varphi$ is not an eigenfunction of the operator $R$ belonging to $\mathfrak R$\origNote{161}. Hence these systems are simply not ``identical'': the description by the wave function is incomplete. There would therefore have to be further determining elements, the ``hidden parameters'' mentioned in~\hyperref[sec:III.2]{III.2}. We shall soon see that this is not so readily possible.

Thus, as long as any quantity $\mathfrak R$ still exhibits dispersion in a large statistical ensemble, it must be possible to divide the ensemble into several differently constituted parts, according to the different states of its elements. This seems all the more plausible because a simple method of making such a division appears actually to exist: one can divide the ensemble according to the different values that $\mathfrak R$ has assumed in it. A genuinely uniform ensemble should be obtained only after this subdivision or decomposition has been carried out with respect to all existing quantities $\mathfrak R,\mathfrak S,\mathfrak T,\ldots$. At the end, none of these quantities would any longer exhibit dispersion in any of the resulting subensembles.

First of all, the assertions in these last sentences are false, because they overlook the fact that measurement changes the measured system. Suppose we measure $\mathfrak R$, which for simplicity can assume only two values $a_1,a_2$, on all the objects $\boldsymbol{S}_1,\ldots,\boldsymbol{S}_N$. Let the value $a_1$ be found on $\boldsymbol{S}'_1,\ldots,\boldsymbol{S}'_{N_1}$ and $a_2$ on $\boldsymbol{S}''_1,\ldots,\boldsymbol{S}''_{N-N_1}$. Afterwards $\mathfrak R$ exhibits dispersion neither in $[\boldsymbol{S}'_1,\ldots,\boldsymbol{S}'_{N_1}]$ nor in $[\boldsymbol{S}''_1,\ldots,\boldsymbol{S}''_{N-N_1}]$: it is always $a_1$ in the former and always $a_2$ in the latter. This, however, is not merely a division of $[\boldsymbol{S}_1,\ldots,\boldsymbol{S}_N]$ into these two parts, since measuring $\mathfrak R$ has changed the individual systems.

We did find in \hyperref[advantage:IV.1.1]{\textbf{1}} a method of determining the distribution of values of $\mathfrak R$ while changing $[\boldsymbol{S}_1,\ldots,\boldsymbol{S}_N]$ only slightly: we measured only on $\boldsymbol{S}_1,\ldots,\boldsymbol{S}_M$, with $M$ large and $M/N$ small. But this procedure does not produce a decomposition, since for most of the systems $\boldsymbol{S}_1,\ldots,\boldsymbol{S}_N$, namely $\boldsymbol{S}_{M+1},\ldots,\boldsymbol{S}_N$, it does not establish the value of $\mathfrak R$ in each individual system. We now see why the proposed method of producing completely uniform ensembles fails. Measure next a second quantity $\mathfrak S$, likewise capable of only two values $b_1,b_2$, in $[\boldsymbol{S}'_1,\ldots,\boldsymbol{S}'_{N_1}]$ and $[\boldsymbol{S}''_1,\ldots,\boldsymbol{S}''_{N-N_1}]$. Let $b_1$ be found in $\boldsymbol{S}'''_1,\ldots,\boldsymbol{S}'''_{N_{11}}$ and $\boldsymbol{S}^{\mathrm V}_1,\ldots,\boldsymbol{S}^{\mathrm V}_{N_{12}}$, and $b_2$ in $\boldsymbol{S}^{\mathrm{IV}}_1,\ldots,\boldsymbol{S}^{\mathrm{IV}}_{N_1-N_{11}}$ and $\boldsymbol{S}^{\mathrm{VI}}_1,\ldots,\boldsymbol{S}^{\mathrm{VI}}_{N-N_1-N_{12}}$. Then $\mathfrak S$ no longer exhibits dispersion in
\begin{gather*}
[\boldsymbol{S}'''_1,\ldots,\boldsymbol{S}'''_{N_{11}}],\quad
[\boldsymbol{S}^{\mathrm{IV}}_1,\ldots,\boldsymbol{S}^{\mathrm{IV}}_{N_1-N_{11}}],\\
[\boldsymbol{S}^{\mathrm V}_1,\ldots,\boldsymbol{S}^{\mathrm V}_{N_{12}}],\quad
[\boldsymbol{S}^{\mathrm{VI}}_1,\ldots,\boldsymbol{S}^{\mathrm{VI}}_{N-N_1-N_{12}}]:
\end{gather*}
its values are always $b_1,b_2,b_1,b_2$, respectively. But although the first two ensembles are parts of $[\boldsymbol{S}'_1,\ldots,\boldsymbol{S}'_{N_1}]$ and the last two parts of $[\boldsymbol{S}''_1,\ldots,\boldsymbol{S}''_{N-N_1}]$, in which $\mathfrak R$ had no dispersion, $\mathfrak R$ may exhibit dispersion in every one of them, because measuring $\mathfrak S$ has changed the individual systems of which they consist! Thus no progress is made, because each step destroys the result of the preceding ones\origNote{162}; and no accumulation of successive
% Source PDF page 169 / original page 162
\hypertarget{origpage-162}{}%
measurements can bring causal order into this confusion. In the atomic domain we are at the boundary of the physical world, where every measurement is an intervention of the same order of magnitude as the measured object and therefore affects it substantially; in essence, this is a consequence of the uncertainty relations.

Thus there may be no method that always permits ensembles with dispersion to be further decomposed without changing their elements, still less one that leads all the way to uniform ensembles with no dispersion at all, which we are accustomed to think of as consisting of causally determined individual objects identical with one another. Nevertheless, one might try to maintain the fiction that every ensemble with dispersion can be divided into two or more parts different both from one another and from the original ensemble, without changing their elements: that is, in such a way that mixing the two component ensembles reproduces the original ensemble.

We see, then, that the attempt to establish causality as a definition of sameness has turned into a factual question, one that can and must be answered, and may perhaps be answered in the negative. The question is: can every ensemble $[\boldsymbol{S}_1,\ldots,\boldsymbol{S}_N]$ in which some quantities $\mathfrak R$ exhibit dispersion really be represented by mixing two or more ensembles different from one another and from it? More than two, say $n=3,4,\ldots$, can be reduced to two by considering the first and the mixture of the remaining $n-1$.

If, for example, $[\boldsymbol{S}_1,\ldots,\boldsymbol{S}_N]$ were a mixture of $[\boldsymbol{S}'_1,\ldots,\boldsymbol{S}'_P]$ and $[\boldsymbol{S}''_1,\ldots,\boldsymbol{S}''_Q]$, then for every quantity $\mathfrak R$ its probability function $w_{\mathfrak R}(a)$, as in note~\hyperlink{note-157}{157}, could be expressed in terms of the probability functions $w'_{\mathfrak R}(a),w''_{\mathfrak R}(a)$ of the last two ensembles as follows:
\begin{equation}\tag{$M_1$}\label{eq:IV.1.M1}
w_{\mathfrak R}(a)=\alpha w'_{\mathfrak R}(a)+\beta w''_{\mathfrak R}(a),\qquad
\alpha>0,\quad\beta>0,\quad\alpha+\beta=1.
\end{equation}
Here $\alpha=P/N$, $\beta=Q/N$, and $N=P+Q$; these coefficients do not depend on $\mathfrak R$. Ultimately, therefore, this is a mathematical question: if an ensemble with probability functions $w_{\mathfrak R}(a)$ contains quantities $\mathfrak R$ with dispersion---the corresponding property of $w_{\mathfrak R}(a)$ is stated in note~\hyperlink{note-160}{160}---are there two other ensembles, with probability functions $w'_{\mathfrak R}(a)$ and $w''_{\mathfrak R}(a)$, such that~\hyperref[eq:IV.1.M1]{$M_1$} holds for all $\mathfrak R$? This can also be formulated somewhat differently if we characterize an ensemble not by the probability functions $w_{\mathfrak R}(a)$ of the quantities $\mathfrak R$, but by their expectation values
\[
\operatorname{Erw}(\mathfrak R)=\int_{-\infty}^{+\infty}a\,dw_{\mathfrak R}(a).
\]
Our question then takes the following form. An ensemble is
% Source PDF page 170 / original page 163
\hypertarget{origpage-163}{}%
dispersion-free if, for every $\mathfrak R$, the expression
\[
\operatorname{Erw}\bigl([\mathfrak R-\operatorname{Erw}(\mathfrak R)]^2\bigr)
=\operatorname{Erw}(\mathfrak R^2)-[\operatorname{Erw}(\mathfrak R)]^2
\]
is zero; see note~\hyperlink{note-160}{160}. That is,
\begin{equation}\tag{$\mathrm{Str}_1$}\label{eq:IV.1.Str1}
\operatorname{Erw}(\mathfrak R^2)=[\operatorname{Erw}(\mathfrak R)]^2.
\end{equation}
If this is not the case, is it always possible to find two other ensembles, with $\operatorname{Erw}'(\mathfrak R)$ and $\operatorname{Erw}''(\mathfrak R)$, where
\[
\operatorname{Erw}(\mathfrak R)\not\equiv\operatorname{Erw}'(\mathfrak R)
\not\equiv\operatorname{Erw}''(\mathfrak R),
\]
such that
\begin{equation}\tag{$M_2$}\label{eq:IV.1.M2}
\begin{gathered}
\operatorname{Erw}(\mathfrak R)=\alpha\operatorname{Erw}'(\mathfrak R)+\beta\operatorname{Erw}''(\mathfrak R),\\
\alpha>0,\quad\beta>0,\quad\alpha+\beta=1
\end{gathered}
\end{equation}
always holds, with $\alpha,\beta$ independent of $\mathfrak R$? [Observe that for a given quantity $\mathfrak R$, the number $\operatorname{Erw}(\mathfrak R)$ is no substitute for the function $w_{\mathfrak R}(a)$. Knowledge of all the $\operatorname{Erw}(\mathfrak R)$, however, is equivalent to knowledge of all the $w_{\mathfrak R}(a)$. For if
\[
f_a(x)=\begin{cases}1,&x\leq a,\\0,&x>a,\end{cases}
\]
then $w_{\mathfrak R}(a)=\operatorname{Erw}(f_a(\mathfrak R))$.]

For the mathematical treatment of this question, it is more convenient to consider the corresponding $\operatorname{Erw}(\mathfrak R)$ instead of the ensembles $[\boldsymbol{S}_1,\ldots,\boldsymbol{S}_N]$. Every ensemble has such a function, defined for all physical quantities $\mathfrak R$ in $\boldsymbol{S}$ and taking real numbers as its values; conversely, this function completely characterizes all the statistical properties of the ensemble. See the preceding remarks on the connection between $\operatorname{Erw}(\mathfrak R)$ and $w_{\mathfrak R}(a)$. We must, of course, still discover what properties a function of $\mathfrak R$ must possess in order to be the $\operatorname{Erw}(\mathfrak R)$ of a suitable ensemble. Once we have done so, however, we can define:

\noindent\phantomsection$\alpha)$\label{def:IV.1.alpha} An expectation function $\operatorname{Erw}(\mathfrak R)$ is called \emph{dispersion-free} if it satisfies condition~\hyperref[eq:IV.1.Str1]{$\mathrm{Str}_1$}.

\noindent\phantomsection$\beta)$\label{def:IV.1.beta} An expectation function $\operatorname{Erw}(\mathfrak R)$ is called \emph{uniform} or \emph{pure} if~\hyperref[eq:IV.1.M2]{$M_2$} implies
\[
\operatorname{Erw}(\mathfrak R)\equiv\operatorname{Erw}'(\mathfrak R)
\equiv\operatorname{Erw}''(\mathfrak R).
\]
It is clear from their meaning that every dispersion-free $\operatorname{Erw}(\mathfrak R)$ function is pure, and we shall shortly prove this. Our question, however, is: is every pure $\operatorname{Erw}(\mathfrak R)$ function dispersion-free?

Every $\operatorname{Erw}(\mathfrak R)$ function must evidently have the following properties:

\noindent\phantomsection\textbf{A.}\label{cond:IV.1.A} If the quantity $\mathfrak R$ is identically $1$---that is, if the ``measurement instruction'' for $\mathfrak R$ says that nothing need be measured, since $\mathfrak R$ always has the value $1$---then $\operatorname{Erw}(\mathfrak R)=1$.

\noindent\phantomsection\textbf{B.}\label{cond:IV.1.B} For every $\mathfrak R$ and every real number $a$, $\operatorname{Erw}(a\mathfrak R)=a\operatorname{Erw}(\mathfrak R)$.\origNote{163}

% Source PDF page 171 / original page 164
\hypertarget{origpage-164}{}%
\noindent\phantomsection\textbf{C.}\label{cond:IV.1.C} If the quantity $\mathfrak R$ is by its nature never negative---if, for example, it is the square of another quantity $\mathfrak S$\textsuperscript{\hyperlink{note-163}{163}}---then $\operatorname{Erw}(\mathfrak R)\geq0$ as well.

\noindent\phantomsection\textbf{D.}\label{cond:IV.1.D} If the quantities $\mathfrak R,\mathfrak S,\ldots$ are simultaneously measurable, then
\[
\operatorname{Erw}(\mathfrak R+\mathfrak S+\cdots)=\operatorname{Erw}(\mathfrak R)+\operatorname{Erw}(\mathfrak S)+\cdots.
\quad\text{\textsuperscript{\hyperlink{note-163}{163}}}
\]
For $\mathfrak R,\mathfrak S,\ldots$ that are not simultaneously measurable, $\mathfrak R+\mathfrak S+\cdots$ is undefined; see above.

All this follows directly from the definitions of the quantities under consideration---that is, their measurement instructions---and from the definition of expectation as the arithmetic mean of all the measurement results in a sufficiently large statistical ensemble. Property \hyperref[cond:IV.1.D]{D} rests on the theorem of probability theory that the expectation of a sum is always the sum of the expectations of its individual terms, regardless of any statistical dependence between them. This differs, for example, from the corresponding rule for a product. That we formulated it only for simultaneously measurable $\mathfrak R,\mathfrak S,\ldots$ is natural: otherwise $\mathfrak R+\mathfrak S+\cdots$ is meaningless.

Quantum mechanics, however, contains a further arithmetic operation beyond those discussed so far: the addition of any two quantities, which need not be simultaneously observable. This rests on the fact that the sum $R+S$ of two Hermitian operators $R,S$ is again Hermitian even when $R,S$ do not commute, whereas their product $RS$, for example, is Hermitian only if they commute; see~\hyperref[sec:II.5]{II.5}. In every state $\varphi$, the expectations add:
\[
(R\varphi,\varphi)+(S\varphi,\varphi)=((R+S)\varphi,\varphi);
\]
see \hyperref[post:III.1.E2]{$E_2$} in~\hyperref[sec:III.1]{III.1}. The same holds for several terms. We now incorporate this fact into our general approach, which has not yet been specialized to quantum mechanics at all:

\noindent\phantomsection\textbf{E.}\label{cond:IV.1.E} If $\mathfrak R,\mathfrak S,\ldots$ are any quantities, there is a further quantity $\mathfrak R+\mathfrak S+\cdots$, independent of the choice of the $\operatorname{Erw}(\mathfrak R)$ function, such that
\[
\operatorname{Erw}(\mathfrak R+\mathfrak S+\cdots)=\operatorname{Erw}(\mathfrak R)+\operatorname{Erw}(\mathfrak S)+\cdots.
\]
If $\mathfrak R,\mathfrak S,\ldots$ are simultaneously measurable, \hyperref[cond:IV.1.D]{D} requires this $\mathfrak R+\mathfrak S+\cdots$ to be the ordinary sum. In general, however, it is characterized only implicitly by \hyperref[cond:IV.1.E]{E}, and we can hardly combine the measurement instructions for $\mathfrak R,\mathfrak S,\ldots$ into one for $\mathfrak R+\mathfrak S+\cdots$.\origNote{164}

We add the following to what has been said. We wish to admit not only $\operatorname{Erw}(\mathfrak R)$ functions representing expectations, but also those corresponding to relative expectations; that is, we drop the normalizing condition \hyperref[cond:IV.1.A]{A}. If $\operatorname{Erw}(1)$, which is $\geq0$ by \hyperref[cond:IV.1.C]{C}, is finite and nonzero, this makes no essential difference, since everything remains as before for $\operatorname{Erw}(\mathfrak R)/\operatorname{Erw}(1)$. But $\operatorname{Erw}(1)=\infty$ corresponds to an essentially different possibility, for whose sake we made this
% Source PDF page 172 / original page 165
\hypertarget{origpage-165}{}%
extension, and which is best illustrated by a simple example. There are cases in which it is better to work with relative probabilities rather than normalized probabilities, in particular when the total relative probability is infinite; $\operatorname{Erw}(1)$ corresponds to the total probability. One such case is the following. Let the system under consideration be a particle moving in one dimension, with a statistical distribution such that it is equally likely to be anywhere on an infinite straight line. Every finite interval on this line then has probability zero. The equal probability of all positions, however, is expressed not by this fact, but by the fact that the probabilities of two finite intervals are in the ratio of their lengths. Since $0/0$ is meaningless, we can express this only by introducing their lengths as their relative probabilities; the total relative probability then becomes $\infty$.

Taking all this into account, we obtain the following final form of our conditions. \hyperref[cond:IV.Aprime]{A$'$} corresponds to \hyperref[cond:IV.1.C]{C}; \hyperref[cond:IV.Bprime]{B$'$} corresponds to \hyperref[cond:IV.1.B]{B}, \hyperref[cond:IV.1.D]{D}, and \hyperref[cond:IV.1.E]{E}.

\noindent\phantomsection\textbf{A$'$.}\label{cond:IV.Aprime} If the quantity $\mathfrak R$ is by its nature never negative---if, for example, it is the square of another quantity $\mathfrak S$---then $\operatorname{Erw}(\mathfrak R)\geq0$.

\noindent\phantomsection\textbf{B$'$.}\label{cond:IV.Bprime} If $\mathfrak R,\mathfrak S,\ldots$ are any quantities and $a,b,\ldots$ are real numbers, then
\[
\operatorname{Erw}(a\mathfrak R+b\mathfrak S+\cdots)=a\operatorname{Erw}(\mathfrak R)+b\operatorname{Erw}(\mathfrak S)+\cdots.
\]
We emphasize further:

\noindent\phantomsection1.\label{remark:IV.1.1} Since we are considering relative expectations, the functions $\operatorname{Erw}(\mathfrak R)$ and $c\operatorname{Erw}(\mathfrak R)$, with $c>0$ a constant, are to be regarded as not essentially different.

\noindent\phantomsection2.\label{remark:IV.1.2} $\operatorname{Erw}(\mathfrak R)\equiv0$ for all $\mathfrak R$ yields no assertions; this function is therefore to be excluded.

\noindent\phantomsection3.\label{remark:IV.1.3} Absolute, that is, properly normalized, expectations occur when $\operatorname{Erw}(1)=1$. By \hyperref[cond:IV.Aprime]{A$'$}, $\operatorname{Erw}(1)$ is in any case $\geq0$; as long as it is finite and nonzero, \hyperref[remark:IV.1.1]{1} with $c=1/\operatorname{Erw}(1)$ restores the proper normalization. When $\operatorname{Erw}(1)=0$, as we shall show, \hyperref[remark:IV.1.2]{2} applies, and this case is therefore excluded. When $\operatorname{Erw}(1)=\infty$, the statistics are essentially unnormalized, that is, relative.

It remains to return to our definitions \hyperref[def:IV.1.alpha]{$\alpha)$}, \hyperref[def:IV.1.beta]{$\beta)$}. In view of \hyperref[remark:IV.1.1]{1},~\hyperref[eq:IV.1.M2]{$M_2$} can be replaced by the simpler condition
\begin{equation}\tag{$M_3$}\label{eq:IV.1.M3}
\operatorname{Erw}(\mathfrak R)=\operatorname{Erw}'(\mathfrak R)+\operatorname{Erw}''(\mathfrak R).
\end{equation}
For~\hyperref[eq:IV.1.Str1]{$\mathrm{Str}_1$}, we must observe that the calculation there assumed $\operatorname{Erw}(1)=1$. When $\operatorname{Erw}(1)=\infty$, the property of being dispersion-free cannot be defined at all, since it means $\operatorname{Erw}((\mathfrak R-\varrho)^2)=0$, where $\varrho$ is the absolute
% Source PDF page 173 / original page 166
\hypertarget{origpage-166}{}%
expectation of $\mathfrak R$, namely $\operatorname{Erw}(\mathfrak R)/\operatorname{Erw}(1)$, which here may be $\infty/\infty$ and hence meaningless.\origNote{165} Thus \hyperref[def:IV.1.alpha]{$\alpha)$}, \hyperref[def:IV.1.beta]{$\beta)$} now read as follows:

\noindent\phantomsection$\alpha')$\label{def:IV.1.alphaprime} An expectation function $\operatorname{Erw}(\mathfrak R)$ is called \emph{dispersion-free} if $\operatorname{Erw}(1)$ is nonzero and finite and, after normalization to $\operatorname{Erw}(1)=1$ as allowed by \hyperref[remark:IV.1.1]{1}, it satisfies~\hyperref[eq:IV.1.Str1]{$\mathrm{Str}_1$}.

\noindent\phantomsection$\beta')$\label{def:IV.1.betaprime} An expectation function $\operatorname{Erw}(\mathfrak R)$ is called \emph{uniform} or \emph{pure} if~\hyperref[eq:IV.1.M3]{$M_3$} implies
\[
\operatorname{Erw}'(\mathfrak R)=c'\operatorname{Erw}(\mathfrak R),\qquad
\operatorname{Erw}''(\mathfrak R)=c''\operatorname{Erw}(\mathfrak R),
\]
where $c',c''$ are constants, necessarily with $c'+c''=1$ and, by \hyperref[cond:IV.Aprime]{A$'$} and \hyperref[remark:IV.1.1]{1}, \hyperref[remark:IV.1.2]{2}, with $c'>0$, $c''>0$.

On the basis of \hyperref[cond:IV.Aprime]{A$'$}, \hyperref[cond:IV.Bprime]{B$'$} and \hyperref[def:IV.1.alphaprime]{$\alpha')$}, \hyperref[def:IV.1.betaprime]{$\beta')$}, we can now decide the question of causality conclusively as soon as we know the physical quantities $\mathfrak R,\mathfrak S,\ldots$ in $\boldsymbol{S}$ and the functional relations between them. The next section will do this for quantum mechanics.

We close this section with two further remarks.

The first concerns the case $\operatorname{Erw}(1)=0$. \hyperref[cond:IV.Bprime]{B$'$} then gives $\operatorname{Erw}(c)=0$. If a quantity $\mathfrak R$ is always $\geq c'$ and $\leq c''$, \hyperref[cond:IV.Aprime]{A$'$} gives $\operatorname{Erw}(c''-\mathfrak R)\geq0$ and $\operatorname{Erw}(\mathfrak R-c')\geq0$, whence \hyperref[cond:IV.Bprime]{B$'$} gives $\operatorname{Erw}(c')\leq\operatorname{Erw}(\mathfrak R)\leq\operatorname{Erw}(c'')$, so $\operatorname{Erw}(\mathfrak R)=0$. Now let $\mathfrak R$ be arbitrary, and let $f_1(x),f_2(x),\ldots$ be a sequence of bounded functions with
\[
f_1(x)+f_2(x)+\cdots=x
\]
% Source formula retained verbatim; see tmp/translation/review-IV-V.txt.
for example,
\[
f_1(x)=\frac{\sin x}{x},\qquad
f_n(x)=\frac{\sin nx}{nx}-\frac{\sin(n-1)x}{(n-1)x},\quad n=2,3,\ldots.
\]
Then $\operatorname{Erw}(f_n(\mathfrak R))=0$ for $n=1,2,\ldots$, and \hyperref[cond:IV.Bprime]{B$'$} therefore gives $\operatorname{Erw}(\mathfrak R)=0$ as well. Thus, as asserted earlier, $\operatorname{Erw}(1)=0$ is to be excluded by \hyperref[remark:IV.1.2]{2}.\tnote{The printed example does not have the stated sum: its $N$th partial sum is $\sin(Nx)/(Nx)$, which tends to $0$ for $x\ne0$, rather than to $x$. The formulas are retained as printed.}

Second, it is striking that~\hyperref[eq:IV.1.Str1]{$\mathrm{Str}_1$} characterizes the absence of dispersion by $\operatorname{Erw}(\mathfrak R^2)=[\operatorname{Erw}(\mathfrak R)]^2$, although in this case $\operatorname{Erw}(\mathfrak R)$ is simply the value of $\mathfrak R$, and $\operatorname{Erw}(f(\mathfrak R))$ the value of $f(\mathfrak R)$, so that for every function $f(x)$ we must have
\begin{equation}\tag{$\mathrm{Str}_2$}\label{eq:IV.1.Str2}
\operatorname{Erw}(f(\mathfrak R))=f(\operatorname{Erw}(\mathfrak R)).
\end{equation}
Condition~\hyperref[eq:IV.1.Str1]{$\mathrm{Str}_1$} is the special case $f(x)=x^2$ of~\hyperref[eq:IV.1.Str2]{$\mathrm{Str}_2$}; how is it that this suffices? The answer is that if~\hyperref[eq:IV.1.Str2]{$\mathrm{Str}_2$} holds for $f(x)=x^2$, it automatically holds for all $f(x)$. In place of $x^2$, one could even have used any other continuous convex function of $x$, that is, any function for which
\[
f\left(\frac{x+y}{2}\right)<\frac{f(x)+f(y)}2\qquad(x\ne y).
\]
We do not enter into the proof.

% Source PDF page 174 / original page 167
\hypertarget{origpage-167}{}%
\section{Proof of the Statistical Formulas}\label{sec:IV.2}
As we know, the physical quantities of a quantum-mechanical system are uniquely associated with hypermaximal Hermitian operators; see, for example, the discussion in~\hyperref[sec:III.5]{III.5}. It is convenient to assume that this correspondence is one-to-one, that is, that every hypermaximal Hermitian operator really corresponds to a physical quantity. We occasionally made use of this in~\hyperref[sec:III.3]{III.3} as well. The following laws hold; see \hyperref[post:III.5.F]{F}, \hyperref[post:III.5.L]{L} in~\hyperref[sec:III.5]{III.5} and the remarks at the end of~\hyperref[sec:IV.1]{IV.1}:

\noindent\phantomsection\textbf{I.}\label{cond:IV.2.I} If the quantity $\mathfrak R$ has the operator $R$, then the quantity $f(\mathfrak R)$ has the operator $f(R)$.

\noindent\phantomsection\textbf{II.}\label{cond:IV.2.II} If the quantities $\mathfrak R,\mathfrak S,\ldots$ have the operators $R,S,\ldots$, then the quantity $\mathfrak R+\mathfrak S+\cdots$ has the operator $R+S+\cdots$. Simultaneous measurability of $\mathfrak R,\mathfrak S,\ldots$ is not assumed; see the remarks at the place just cited.

\hyperref[cond:IV.Aprime]{A$'$}, \hyperref[cond:IV.Bprime]{B$'$}, \hyperref[def:IV.1.alphaprime]{$\alpha')$}, \hyperref[def:IV.1.betaprime]{$\beta')$} and \hyperref[cond:IV.2.I]{I}, \hyperref[cond:IV.2.II]{II} form the mathematical starting point of our analysis.

Let $\varphi_1,\varphi_2,\ldots$ be a complete orthonormal system. In place of each operator $R$, we consider the matrix $a_{\mu\nu}=(R\varphi_\mu,\varphi_\nu)$. We also form the Hermitian operators with the respective matrices
\[
e_{\mu\nu}^{(n)}=\begin{cases}1,&\mu=\nu=n,\\0,&\text{otherwise},\end{cases}
\]
\[
f_{\mu\nu}^{(mn)}=\begin{cases}1,&\mu=m,\ \nu=n,\\1,&\mu=n,\ \nu=m,\\0,&\text{otherwise},\end{cases}
\qquad
g_{\mu\nu}^{(mn)}=\begin{cases}i,&\mu=m,\ \nu=n,\\-i,&\mu=n,\ \nu=m,\\0,&\text{otherwise}.\end{cases}
\]
These operators are, incidentally,\tnote{The source prints $P_{[\varphi_m]}$ here, although the first matrix is indexed by $n$ and corresponds to $P_{[\varphi_n]}$. The printed subscript is retained.}
\[
P_{[\varphi_m]},\qquad
P_{[(\varphi_m+\varphi_n)/\sqrt2]}-P_{[(\varphi_m-\varphi_n)/\sqrt2]},\qquad
P_{[(\varphi_m+i\varphi_n)/\sqrt2]}-P_{[(\varphi_m-i\varphi_n)/\sqrt2]}.
\]
Let the corresponding quantities be $\mathfrak U^{(n)},\mathfrak V^{(mn)},\mathfrak W^{(mn)}$. Since $a_{nm}=\overline{a}_{mn}$, we evidently have
\[
a_{\mu\nu}=\sum_n a_{nn}e_{\mu\nu}^{(n)}
+\sum_{m<n}\operatorname{Re}a_{mn}f_{\mu\nu}^{(mn)}
+\sum_{m<n}\operatorname{Im}a_{mn}g_{\mu\nu}^{(mn)},
\]
and hence
\[
\mathfrak R=\sum_n a_{nn}\mathfrak U^{(n)}
+\sum_{m<n}\operatorname{Re}a_{mn}\mathfrak V^{(mn)}
+\sum_{m<n}\operatorname{Im}a_{mn}\mathfrak W^{(mn)}.
\]
By \hyperref[cond:IV.2.II]{II} and \hyperref[cond:IV.Bprime]{B$'$},
\begin{align*}
\operatorname{Erw}(\mathfrak R)={}&\sum_n a_{nn}\operatorname{Erw}(\mathfrak U^{(n)})
+\sum_{m<n}\operatorname{Re}a_{mn}\operatorname{Erw}(\mathfrak V^{(mn)})\\
&+\sum_{m<n}\operatorname{Im}a_{mn}\operatorname{Erw}(\mathfrak W^{(mn)}).
\end{align*}
% Source PDF page 175 / original page 168
\hypertarget{origpage-168}{}%
Thus, if we put
\begin{alignat*}{2}
u_{nn}&=\operatorname{Erw}(\mathfrak U^{(n)}),\\
u_{mn}&=\tfrac12\operatorname{Erw}(\mathfrak V^{(mn)})+\tfrac i2\operatorname{Erw}(\mathfrak W^{(mn)}),&\quad &m<n,\\
u_{nm}&=\tfrac12\operatorname{Erw}(\mathfrak V^{(mn)})-\tfrac i2\operatorname{Erw}(\mathfrak W^{(mn)}),&\quad &m<n,
\end{alignat*}
we obtain
\[
\operatorname{Erw}(\mathfrak R)=\sum_{m,n}u_{nm}a_{mn}.
\]
Since $u_{nm}=\overline{u}_{mn}$, we can define a Hermitian operator $U$ by $(U\varphi_m,\varphi_n)=u_{mn}$,\origNote{166} and the right-hand side is $\operatorname{Spur}(UR)$; see~\hyperref[sec:II.11]{II.11}. Thus the final formula is
\begin{equation}\tag{Sp.}\label{eq:IV.2.Sp}
\operatorname{Erw}(\mathfrak R)=\operatorname{Spur}(UR),
\end{equation}
where $U$ is a Hermitian operator independent of $R$,\origNote{167} and is therefore determined by the ensemble itself.

In view of \hyperref[cond:IV.2.II]{II},~\hyperref[eq:IV.2.Sp]{Sp.} satisfies \hyperref[cond:IV.Bprime]{B$'$} for every choice of $U$. We therefore need only determine what restriction \hyperref[cond:IV.Aprime]{A$'$} places on $U$. If $\|\varphi\|=1$ but $\varphi$ is otherwise arbitrary, the quantity $\mathfrak R$ belonging to $P_{[\varphi]}$ satisfies $\mathfrak R^2=\mathfrak R$ by $P_{[\varphi]}^2=P_{[\varphi]}$ and \hyperref[cond:IV.2.I]{I}, so that \hyperref[cond:IV.Aprime]{A$'$} gives $\operatorname{Erw}(\mathfrak R)\geq0$. Thus
\[
\operatorname{Spur}(UP_{[\varphi]})=(U\varphi,\varphi)\geq0.
\]
If $f$ is arbitrary and $f\ne0$, we may choose $\varphi=f/\|f\|$; then $(U\varphi,\varphi)=(Uf,f)/\|f\|^2$, so again $(Uf,f)\geq0$. For $f=0$ this holds in any case. Thus $U$ is definite. The definiteness of $U$, which therefore follows from \hyperref[cond:IV.Aprime]{A$'$}, is also sufficient for the validity of \hyperref[cond:IV.Aprime]{A$'$}.

\hyperref[cond:IV.Aprime]{A$'$} says that every $\operatorname{Erw}(\mathfrak S^2)$ is to be $\geq0$, and nothing more. For if $\mathfrak R$ can take only nonnegative values, then $f(\mathfrak R)=\mathfrak R$ for $f(x)=|x|$; and since $g(x)=\sqrt{|x|}$ satisfies $(g(x))^2=f(x)$ identically, we have $(g(\mathfrak R))^2=f(\mathfrak R)$, hence $\mathfrak R=\mathfrak S^2$ with $\mathfrak S=g(\mathfrak R)$.\origNote{168} If $S$ is the operator of $\mathfrak S$, we therefore require $\operatorname{Spur}(US^2)\geq0$. Now $S^2$ is definite, since
\[
(S^2f,f)=(Sf,Sf)\geq0.
\]
Writing $A,B$ in place of $U,S^2$, we therefore have to prove the theorem: if $A,B$ are Hermitian and definite, then $\operatorname{Spur}(AB)\geq0$. We proved this in~\hyperref[sec:II.11]{II.11}, using a general theorem about definite operators; see note~\hyperlink{note-114}{114}.\origNote{169}

We have thus completely specified the functions $\operatorname{Erw}(\mathfrak R)$: they correspond to the definite Hermitian operators $U$, and the connection is given by~\hyperref[eq:IV.2.Sp]{Sp.} We shall call $U$ the \emph{statistical operator} of the ensemble in question.

Remarks \hyperref[remark:IV.1.1]{1}, \hyperref[remark:IV.1.2]{2}, \hyperref[remark:IV.1.3]{3} of~\hyperref[sec:IV.1]{IV.1} are now easy to discuss. Correspondingly, we obtain:

\noindent\phantomsection1.\label{remark:IV.2.1} From the standpoint of relative probabilities and expectations, $U$ and $cU$ are not essentially different, where $c>0$ is a constant.

% Source PDF page 176 / original page 169
\hypertarget{origpage-169}{}%
\noindent\phantomsection2.\label{remark:IV.2.2} $U=0$ yields no assertions whatever and is therefore to be excluded.

\noindent\phantomsection3.\label{remark:IV.2.3} We obtain absolute, that is, properly normalized, probabilities and expectations when $\operatorname{Spur}U=1$. As long as $\operatorname{Spur}U$ is finite, we can normalize $U$, according to \hyperref[remark:IV.2.1]{1}, by multiplying it by $c=1/\operatorname{Spur}U$. The definiteness of $U$ gives $\operatorname{Spur}U\geq0$; in fact, $\operatorname{Spur}U>0$, since $\operatorname{Spur}U=0$ implies $U=0$, the case excluded by \hyperref[remark:IV.2.2]{2}. This was shown generally at the end of~\hyperref[sec:IV.1]{IV.1}, and for our case also follows from~\hyperref[sec:II.11]{II.11}. Only when $\operatorname{Spur}U$ is infinite do we have essentially relative probabilities and expectations.

Finally, we must investigate \hyperref[def:IV.1.alpha]{$\alpha)$}, \hyperref[def:IV.1.beta]{$\beta)$} of~\hyperref[sec:IV.1]{IV.1}, that is, find the dispersion-free and the uniform operators among the $U$.

First consider the dispersion-free ones. We must assume $U$ to be properly normalized; see~\hyperref[sec:IV.1]{IV.1}. We would then require $\operatorname{Erw}(\mathfrak R^2)=[\operatorname{Erw}(\mathfrak R)]^2$ to hold for every $\mathfrak R$, that is, $\operatorname{Spur}(UR^2)=[\operatorname{Spur}(UR)]^2$. For $R=P_{[\varphi]}$ we have $R^2=R=P_{[\varphi]}$ and $\operatorname{Spur}(UP_{[\varphi]})=(U\varphi,\varphi)$; hence we should have $(U\varphi,\varphi)=(U\varphi,\varphi)^2$, or $(U\varphi,\varphi)=0$ or $1$. If $\|\varphi'\|=1$, $\|\varphi''\|=1$, we can vary $\varphi$ continuously, beginning at $\varphi'$ and ending at $\varphi''$, with $\|\varphi\|=1$ throughout.\origNote{170} Then $(U\varphi,\varphi)$ also varies continuously, and since it can be only $0$ or $1$, it is constant. Thus $(U\varphi',\varphi')=(U\varphi'',\varphi'')$. Consequently $(U\varphi,\varphi)$ is either always $0$ or always $1$, from which we immediately infer $U=0$ or $U=1$. The case $U=0$ is excluded by \hyperref[remark:IV.2.2]{2}; $U=1$ cannot be normalized, since $\operatorname{Spur}1$ is the dimension of the space, namely $\infty$, and, as is also easily seen directly, it is not dispersion-free. Thus there are no dispersion-free ensembles.

We now turn to the uniform ones. By \hyperref[def:IV.1.beta]{$\beta)$} and \hyperref[eq:IV.2.Sp]{Sp.}, $U$ is uniform if
\[
U=V+W,
\]
with $V,W$, like $U$, definite and Hermitian, implies $V=c'U$, $W=c''U$.\origNote{171} We assert that this property belongs to the $U=P_{[\varphi]}$, $\|\varphi\|=1$, and to these alone.

First let $U$ have the stated property. Since $U\ne0$, there is an $f_0$ with $Uf_0\ne0$; hence $f_0\ne0$, and therefore $(Uf_0,f_0)>0$; see~\hyperref[sec:II.5]{II.5}, \hyperref[thm:II.19]{Theorem~19}. Form the two Hermitian operators
\[
Vf=\frac{(f,Uf_0)}{(Uf_0,f_0)}Uf_0,\qquad
Wf=Uf-\frac{(f,Uf_0)}{(Uf_0,f_0)}Uf_0.
\]
Then
\[
(Vf,f)=\frac{|(f,Uf_0)|^2}{(Uf_0,f_0)}\geq0,\qquad
(Wf,f)=\frac{(Uf,f)(Uf_0,f_0)-|(f,Uf_0)|^2}{(Uf_0,f_0)}\geq0;
\]
see~\hyperref[sec:II.5]{II.5}, \hyperref[thm:II.19]{Theorem~19}. Thus $V,W$ are definite, and evidently $U=V+W$. It follows that $V=c'U$, and since $Vf_0=Uf_0\ne0$, we have $c'=1$, hence $U=V$. If we now put
\[
\varphi=\frac{Uf_0}{\|Uf_0\|},\qquad\|\varphi\|=1,\qquad
c=\frac{\|Uf_0\|^2}{(Uf_0,f_0)}>0,
\]
% Source PDF page 177 / original page 170
\hypertarget{origpage-170}{}%
then $Uf=Vf=c(f,\varphi)\varphi=cP_{[\varphi]}f$, so $U=cP_{[\varphi]}$; by \hyperref[remark:IV.2.1]{1}, $U$ is therefore essentially $P_{[\varphi]}$.

Conversely, let $U=P_{[\varphi]}$, $\|\varphi\|=1$. If $U=V+W$ with $V,W$ definite, we argue as follows. From $Uf=0$ we obtain
\[
0\leq(Vf,f)\leq(Vf,f)+(Wf,f)=(Uf,f)=0,
\]
so $(Vf,f)=0$ and hence $Vf=0$; see the place just cited. But $(f,\varphi)=0$ implies $Uf=P_{[\varphi]}f=0$, and hence also $Vf=0$. Thus for every $g$, $(f,Vg)=(Vf,g)=0$. That is, everything orthogonal to $\varphi$ is also orthogonal to $Vg$, so $Vg=c_g\varphi$, where $c_g$ is a number depending on $g$. We need only the case $g=\varphi$, namely $V\varphi=c'\varphi$. Every $f$ has the form $(f,\varphi)\varphi+f'$, with $f'$ orthogonal to $\varphi$. Therefore
\[
Vf=(f,\varphi)V\varphi+Vf'=(f,\varphi)c'\varphi=c'P_{[\varphi]}f=c'Uf.
\]
Thus $V=c'U$ and $W=U-V=(1-c')U$, which proves everything.

The uniform ensembles therefore correspond to $U=P_{[\varphi]}$, $\|\varphi\|=1$; in this case~\hyperref[eq:IV.2.Sp]{Sp.} becomes formula \hyperref[post:III.1.E2]{$E_2$} of~\hyperref[sec:III.1]{III.1}:
\begin{equation}\tag{$E_2$}\label{eq:IV.2.E2}
\operatorname{Erw}(\mathfrak R)=\operatorname{Spur}(P_{[\varphi]}R)=(R\varphi,\varphi).
\end{equation}
Observe that $\operatorname{Erw}(1)=\operatorname{Spur}(P_{[\varphi]})=1$, since $P_{[\varphi]}$ belongs to the one-dimensional manifold $[\varphi]$, or by~\hyperref[eq:IV.2.E2]{$E_2$}. Thus the present form of $U$ is properly normalized. Finally, let us determine when $P_{[\varphi]}$ and $P_{[\psi]}$ have the same statistics, that is, when $P_{[\varphi]}=cP_{[\psi]}$ for a constant $c>0$; see \hyperref[remark:IV.2.1]{1}. Since $\operatorname{Spur}(P_{[\varphi]})=\operatorname{Spur}(P_{[\psi]})=1$, we then have $c=1$, hence $P_{[\varphi]}=P_{[\psi]}$, $[\varphi]=[\psi]$, and therefore $\varphi=a\psi$. From $\|\varphi\|=\|\psi\|=1$ it follows that the constant $a$ satisfies $|a|=1$; this is evidently sufficient as well.

We can therefore summarize the result as follows. There are no dispersion-free ensembles. Uniform ensembles do exist, and they correspond precisely to the $U=P_{[\varphi]}$, $\|\varphi\|=1$. For these $U$, \hyperref[eq:IV.2.Sp]{Sp.} becomes \hyperref[eq:IV.2.E2]{$E_2$}; the normalization is correct, and $U$ remains unchanged when $\varphi$ is replaced by $a\varphi$, with $a$ constant and $|a|=1$, whereas every other change in $\varphi$ changes it essentially; see \hyperref[remark:IV.2.1]{1}. Thus the uniform ensembles correspond to the states of quantum mechanics as previously characterized: by a $\varphi$ in Hilbert space with $\|\varphi\|=1$, where a constant factor of absolute value $1$ is immaterial; see, for example,~\hyperref[sec:III.2]{III.2}. Statistical assertions are then made by \hyperref[eq:IV.2.E2]{$E_2$}.\origNote{172}

We have derived all this from the purely qualitative conditions \hyperref[cond:IV.Aprime]{A$'$}, \hyperref[cond:IV.Bprime]{B$'$}, \hyperref[def:IV.1.alpha]{$\alpha)$}, \hyperref[def:IV.1.beta]{$\beta)$}, \hyperref[cond:IV.2.I]{I}, \hyperref[cond:IV.2.II]{II}.

Within the framework of our conditions, the question has therefore been decided against causality: all ensembles exhibit dispersion, even the uniform ones.

It remains to discuss the question of ``hidden parameters'' raised in~\hyperref[sec:III.2]{III.2}, that is, whether the dispersions of
% Source PDF page 178 / original page 171
\hypertarget{origpage-171}{}%
the uniform ensembles characterized by the wave functions $\varphi$, that is, by \hyperref[eq:IV.2.E2]{$E_2$}, might arise because these are not the true states but merely mixtures of several states. Characterizing the actual state would then require, besides the wave function $\varphi$, further specifications---the ``hidden parameters''---which together determine everything causally, that is, lead to dispersion-free ensembles. The statistics of the uniform ensemble, $U=P_{[\varphi]}$, $\|\varphi\|=1$, would then result from averaging over all the actual states of which it is composed, that is, over the range of values of the ``hidden parameters'' realized in those states. This is impossible for two reasons. First, the uniform ensemble in question could then be represented as a mixture of two different ensembles,\origNote{173} contrary to its definition. Second, the dispersion-free ensembles that would have to correspond to the ``actual'' states---ensembles consisting entirely of systems in the same ``actual'' state---do not exist at all. Observe that we did not have to enter into the details of the mechanism of the ``hidden parameters'': the established results of quantum mechanics can in no circumstances be recovered with their aid. Indeed, if determining factors other than the wave function, the ``hidden parameters,'' are to exist, it is even impossible for the same physical quantities with the same relations to be present, that is, for \hyperref[cond:IV.2.I]{I}, \hyperref[cond:IV.2.II]{II} to hold.

It would not suffice for further, hitherto undiscovered physical quantities to exist in addition to the known ones represented in quantum mechanics by operators. For even for the latter, known quantities, the relations assumed by quantum mechanics, namely \hyperref[cond:IV.2.I]{I}, \hyperref[cond:IV.2.II]{II}, would have to fail. Thus this is not, as is often supposed, a question of interpreting quantum mechanics. Rather, quantum mechanics would have to be objectively false for any behavior of the elementary processes other than statistical behavior to be possible.

The following circumstance is also worth mentioning. At first sight, the uncertainty relations bear a certain resemblance to the fundamental postulates of relativity theory. The latter asserts that it is impossible in principle to establish the simultaneity of two events occurring at points a distance $r$ apart more accurately than to within a time interval of length $r/c$, where $c$ is the velocity of light. The uncertainty relations assert that it is impossible in principle to specify the position of a material point in phase space more accurately than to within a region of volume $(h/4\pi)^3$.%
% Source PDF page 179 / original page 172
\hypertarget{origpage-172}{}%
\origNote{174} Yet there is a fundamental difference. Relativity theory does deny the possibility of an objectively precise measurement of simultaneity; nevertheless, by introducing a Galilean frame of reference, one can lay down a coordinate system in the world that permits a definition of simultaneity sufficiently consistent with our ordinary concepts in this respect. This definition of simultaneity is denied objective significance only because the coordinate system can be chosen in infinitely many different ways, producing infinitely many different, equally good definitions of simultaneity. Thus behind the impossibility of measurement lies an infinite multiplicity of possible theoretical definitions. In quantum mechanics it is otherwise. There it is altogether impossible to describe a system with wave function $\varphi$ by points of phase space, even if new, hypothetical and unobserved coordinates, the ``hidden parameters,'' are introduced, since this would lead to dispersion-free ensembles. Thus not only is measurement impossible; so is any reasonable theoretical definition---one that, though empirically unprovable, as in relativity theory, would at least also be empirically irrefutable. The impossibility of measurement in principle thus arises because in one case there are infinitely many ways of defining the concepts in question that do not contradict immediate experience, or the general fundamental assumptions of the theory, while in the other there is none at all.

The position of causality in present-day physics can therefore be characterized as follows. In the macroscopic domain there is no empirical evidence for causality, nor can there be any. For the apparent causal order of the macroscopic world---the world of objects perceptible to the naked eye---certainly has no cause other than the ``law of large numbers,'' quite independently of whether the laws of nature governing the elementary processes, the actual laws, are causal or not.\origNote{175} That macroscopically identical objects behave macroscopically alike has little to do with causality. They are not identical at all, since the coordinates that precisely determine the states of their atoms almost never agree; the macroscopic viewpoint averages over these coordinates, which here are ``hidden parameters.'' But their number is very large, about $10^{25}$ in $1\,\mathrm g$ of matter, so that by familiar theorems of probability theory the averaging produces a far-reaching reduction of all dispersions. This holds, of course, only in general: in special cases, such as Brownian motion, unstable states, and others, this apparent macroscopic causality fails. Only in the atomic domain, in the elementary processes themselves, could the question of
% Source PDF page 180 / original page 173
\hypertarget{origpage-173}{}%
causality really be tested. Here, however, at the present stage of our knowledge, everything speaks against it. The only formal theory now available that orders and summarizes our experience in a reasonably satisfactory way, namely quantum mechanics, is in compelling logical contradiction with it. It would admittedly be an exaggeration to claim that causality has thereby been disposed of. Quantum mechanics in its present form is certainly incomplete, and it may even be false, although this is rather unlikely in view of its astonishing effectiveness in understanding general problems and calculating results for particular ones. Despite its brilliant agreement with experience, and despite the insight it has given us into a qualitatively new aspect of the world, we can never say of a theory that experience has proved it, but only that it is the best known summary of that experience. Yet with all these reservations duly observed, we may still say that at present there is no occasion and no excuse for speaking of causality in nature. No empirical evidence supports its existence, since macroscopic observations are in principle unsuited to doing so; and the only known theory compatible with our experience of elementary processes, quantum mechanics, contradicts it.

It is certainly a deeply rooted way of looking at things, shared by all people, but by no means a necessity of thought. Among other things, the very fact that the statistical theory could be constructed has shown this. Anyone who approaches the subject without preconceived opinions has no reason to cling to it. Under these circumstances, is there any justification for sacrificing a reasonable physical theory to it?

\section{Inferences from Experiments}\label{sec:IV.3}
The preceding section taught us that the most general statistical ensemble compatible with our qualitative fundamental assumptions is characterized by a definite operator $U$, according to the law~\hyperref[eq:IV.2.Sp]{Sp.} The special ensembles we called uniform were characterized by $U=P_{[\varphi]}$, $\|\varphi\|=1$. Since these are the actual states of the system $\boldsymbol{S}$, which cannot be decomposed further, we also call them states, calling $U=P_{[\varphi]}$ the state $\varphi$.

If $U$ has a pure point spectrum, say with eigenvalues $w_1,w_2,\ldots$ and eigenfunctions $\varphi_1,\varphi_2,\ldots$ forming a complete orthonormal system, then, by~\hyperref[sec:II.8]{II.8},
\[
U=\sum_n w_nP_{[\varphi_n]}.
\]
% Source PDF page 181 / original page 174
\hypertarget{origpage-174}{}%
Since $U$ is definite, all $w_n\geq0$: indeed, $U\varphi_n=w_n\varphi_n$, so $(U\varphi_n,\varphi_n)=w_n\geq0$. Moreover,
\[
\sum_nw_n=\sum_n(U\varphi_n,\varphi_n)=\operatorname{Spur}U;
\]
see also the beginning of~\hyperref[sec:IV.1]{IV.1}. Thus $\sum_nw_n=1$ when $U$ is properly normalized. By what was said at the beginning of~\hyperref[sec:IV.1]{IV.1}, $U$ can therefore be regarded as a mixture of the states $\varphi_1,\varphi_2,\ldots$ with the respective relative weights $w_1,w_2,\ldots$. If $U$ is properly normalized, these weights also give the absolute probabilities.

But by~\hyperref[sec:II.11]{II.11}, especially note~\hyperlink{note-115}{115}, a properly normalized $U$, with $\operatorname{Spur}U=1$, is completely continuous and therefore has a pure point spectrum; the same naturally holds whenever $\operatorname{Spur}U$ is finite. Infinite $\operatorname{Spur}U$ may be regarded as a limiting case, into which we shall not enter here. In the case of real interest, therefore, the ensemble under consideration can be represented as a mixture of states, which we have, moreover, chosen pairwise orthogonal. We shall therefore call the general ensembles \emph{mixtures}, in contrast to the uniform ensembles, the states.

If all the eigenvalues of $U$ are simple, that is, $w_1,w_2,\ldots$ are different from one another, then, as we know, $\varphi_1,\varphi_2,\ldots$ are uniquely determined up to constant factors of absolute value $1$, so that the corresponding states, and the $P_{[\varphi_1]},P_{[\varphi_2]},\ldots$, are completely determined. The weights $w_1,w_2,\ldots$ are likewise uniquely determined, apart, of course, from their order. In this case, therefore, one can specify uniquely the pairwise orthogonal states from which the mixture $U$ is constructed. The situation is essentially different if $U$ also has multiple eigenvalues, ``degeneracies.'' In~\hyperref[sec:II.8]{II.8} we discussed in detail how $\varphi_1,\varphi_2,\ldots$ may be chosen: there are infinitely many essentially different ways, although $w_1,w_2,\ldots$ remain uniquely determined even then. List the distinct values among $w_1,w_2,\ldots$ as $w',w'',\ldots$; for each, specify the closed linear manifold of its eigenfunctions, the solutions of $Uf=wf$, as $\mathfrak M_{w'},\mathfrak M_{w''},\ldots$. Then proceed as follows. From each of $\mathfrak M_{w'},\mathfrak M_{w''},\ldots$, choose at will an orthonormal system spanning it, respectively $\chi'_1,\chi'_2,\ldots$, $\chi''_1,\chi''_2,\ldots$, and so on. Then $\chi'_1,\chi'_2,\ldots,\chi''_1,\chi''_2,\ldots,\ldots$ are the $\varphi_1,\varphi_2,\ldots$, and the corresponding eigenvalues $w',w',\ldots,w'',w'',\ldots,\ldots$ are the $w_1,w_2,\ldots$. As soon as one $\mathfrak M_w$ has more than one dimension, that is, a multiple eigenvalue exists, its $\chi_1,\chi_2,\ldots$ are no longer determined up to constant factors of absolute value $1$. For example, $\chi_1$ can be any normalized element of $\mathfrak M_w$. The corresponding states are therefore no longer uniquely determined either.

This phenomenon can also be formulated as follows. If the states $\chi_1,\chi_2,\ldots$ are pairwise orthogonal---that is, if they form an orthonormal system, which may be finite or infinite---and we mix them in such a way that they all receive the same
% Source PDF page 182 / original page 175
\hypertarget{origpage-175}{}%
weight, say the relative weights $1:1:\cdots$, then the mixture thus produced depends only on the closed linear manifold $\mathfrak M$ spanned by $\chi_1,\chi_2,\ldots$. Indeed,
\[
U=P_{[\chi_1]}+P_{[\chi_2]}+\cdots=P_{\mathfrak M}.
\]
If there are finitely many $\chi_1,\chi_2,\ldots$, say $s$ of them, $\chi_1,\ldots,\chi_s$, this $U$ can also be regarded as a mixture of all the normalized elements of $\mathfrak M$, that is, all the states in $\mathfrak M$. These are
\[
\chi=x_1\chi_1+\cdots+x_s\chi_s,\qquad |x_1|^2+\cdots+|x_s|^2=1.
\]
Indeed, put $x_1=u_1+iv_1,\ldots,x_s=u_s+iv_s$, and let $K$ denote the $(2s-1)$-dimensional unit sphere
\[
u_1^2+v_1^2+\cdots+u_s^2+v_s^2=1,
\]
corresponding to $|x_1|^2+\cdots+|x_s|^2=1$, with surface element $do$. For
\[
U'=\int\!\!\int_K\cdots\int\!\!\int P_{[\chi]}\,do,
\]
we then have
\begin{align*}
(U'f,g)
&=\int\!\!\int_K\cdots\int\!\!\int(P_{[\chi]}f,g)\,do
=\int\!\!\int_K\cdots\int\!\!\int(f,\chi)\overline{(g,\chi)}\,do\\
&=\int\!\!\int_K\cdots\int\!\!\int
\left(f,\sum_{\mu=1}^s(u_\mu+iv_\mu)\chi_\mu\right)
\overline{\left(g,\sum_{\mu=1}^s(u_\mu+iv_\mu)\chi_\mu\right)}\,do\\
&=\int\!\!\int_K\cdots\int\!\!\int\sum_{\mu,\nu=1}^s
(f,\chi_\mu)\overline{(g,\chi_\nu)}(u_\mu-iv_\mu)(u_\nu+iv_\nu)\,do\\
&=\sum_{\mu,\nu=1}^s(f,\chi_\mu)\overline{(g,\chi_\nu)}
\int\!\!\int_K\cdots\int\!\!\int
[(u_\mu u_\nu+v_\mu v_\nu)+i(u_\mu v_\nu-u_\nu v_\mu)]\,do.
\end{align*}
By symmetry, all the integrals of $u_\mu v_\nu,u_\nu v_\mu$, and those of $u_\mu u_\nu,v_\mu v_\nu$ for $\mu\ne\nu$, are zero;\origNote{176} for $\mu=\nu$, the latter are all $C/(2s)$, with $C>0$.\textsuperscript{\hyperlink{note-176}{176}} Thus
\begin{align*}
(U'f,g)&=\frac Cs\sum_{\mu=1}^s(f,\chi_\mu)\overline{(g,\chi_\mu)}
=\frac Cs\sum_{\mu=1}^s(P_{[\chi_\mu]}f,g)\\
&=\left(\left\{\frac Cs\sum_{\mu=1}^sP_{[\chi_\mu]}\right\}f,g\right).
\end{align*}
Consequently
\[
U'=\frac Cs\sum_{\mu=1}^sP_{[\chi_\mu]}=\frac CsU,
\]
so that $U',U$ are not essentially different.

These matters are of great importance for the character of quantum-mechanical statistics, and we therefore repeat them:

\noindent1. If a mixture consists of mutually orthogonal states with exactly equal weights, it is no longer possible to determine what those states were. Equivalently, by mixing in exactly equal
% Source PDF page 183 / original page 176
\hypertarget{origpage-176}{}%
proportions, one can produce the same mixture from different sets of mutually orthogonal components.

\noindent2. If there are finitely many components, the mixture thus obtained is identical with a mixture of all the states that are linear combinations of these components.

The simplest example is this: mixing orthogonal $\varphi,\psi$ in the ratio $1:1$ produces the same result as mixing, for example, $(\varphi+\psi)/\sqrt2$ and $(\varphi-\psi)/\sqrt2$ in the ratio $1:1$, or mixing all the $x\varphi+y\psi$ with $|x|^2+|y|^2=1$. If two nonorthogonal $\varphi,\psi$ are mixed, in a ratio that need not even be $1:1$, it is even less possible to identify them from the resulting mixture, since the latter can certainly also be regarded as a mixture of orthogonal states.

We defer further examination of the nature of mixtures until the thermodynamic discussions in~\hyperref[sec:V.2]{V.2} and the following sections.

Formula~\hyperref[eq:IV.2.Sp]{Sp.} in~\hyperref[sec:IV.2]{IV.2} specified how to calculate the expectation of the quantity $\mathfrak R$ with operator $R$ in a mixture with statistical operator $U$: it is $\operatorname{Spur}(UR)$. The probability that the value $a$ of $R$ lies in the interval $a'<a\leq a''$, with given $a',a''$ and $a'\leq a''$, is found as in~\hyperref[sec:III.1]{III.1} or~\hyperref[sec:III.5]{III.5}. Form the quantity $F(\mathfrak R)$ with
\[
F(x)=\begin{cases}1,&a'<x\leq a'',\\0,&\text{otherwise};\end{cases}
\]
its expectation is the probability in question. By \hyperref[cond:IV.2.I]{I} in~\hyperref[sec:IV.2]{IV.2}, $F(\mathfrak R)$ has operator $F(R)$; if $E(\lambda)$ is the resolution of the identity belonging to $R$, then, as we have calculated several times, $F(R)=E(a'')-E(a')$. The required probability is
\[
w(a',a'')=\operatorname{Spur}\bigl(U(E(a'')-E(a'))\bigr).
\]
Thus the probability function describing the statistics of $\mathfrak R$ is $w(a)=\operatorname{Spur}(UE(a))$; see~\hyperref[sec:IV.1]{IV.1}, note~\hyperlink{note-175}{175}.\tnote{The source cites note~\hyperlink{note-175}{175}; the definition of the probability function is in note~\hyperlink{note-157}{157}.} For states, that is, $U=P_{[\varphi]}$, we again obtain $w(a)=\operatorname{Spur}(P_{[\varphi]}E(a))=(E(a)\varphi,\varphi)$. Naturally, these probabilities are only relative if $U$ is not properly normalized.

The question of when the quantity $\mathfrak R$ with operator $R$ assumes the value $\lambda^*$ with certainty in the mixture with statistical operator $U$ can be answered directly with the aid of $w(a)$. We must have $w(a)=0$ for $a<\lambda^*$ and $w(a)=1$ for $a\geq\lambda^*$. If $U$ is not properly normalized, the latter condition becomes $w(a)=\operatorname{Erw}(1)=\operatorname{Spur}U$. Thus\origNote{177}
\[
\operatorname{Spur}(UE(a))=0\quad(a<\lambda^*),\qquad
\operatorname{Spur}(U(1-E(a)))=0\quad(a\geq\lambda^*).
\]
For definite operators $A,B$, however, $\operatorname{Spur}(AB)=0$ implies $AB=0$; see~\hyperref[sec:II.11]{II.11}. Hence $UE(a)=0$ for $a<\lambda^*$ and $UE(a)=U$ for $a\geq\lambda^*$, or equivalently $E(a)U=0$ and $E(a)U=U$, respectively, since the factors must commute when their product is Hermitian. That is, for $f=Ug$,
\[
E(a)f=\begin{cases}f,&a\geq\lambda^*,\\0,&a<\lambda^*;\end{cases}
\]
according to the discussion in~\hyperref[sec:II.8]{II.8},
% Source PDF page 184 / original page 177
\hypertarget{origpage-177}{}%
this means $Rf=\lambda^*f$, or $RUg=\lambda^*Ug$ identically in $g$. The final condition is therefore $RU=\lambda^*U$. Alternatively, if $\mathfrak M$ denotes the closed linear manifold formed by all solutions $h$ of $Rh=\lambda^*h$, the condition is that $Uf$ always lies in $\mathfrak M$.

The same result could also have been obtained from the vanishing of the dispersion, that is, of the expectation, possibly relative, of $(\mathfrak R-\lambda^*)^2$.

In~\hyperref[sec:III.3]{III.3} we answered the following questions, with $\mathfrak R,\mathfrak S,\ldots$ physical quantities and $R,S,\ldots$ their respective operators:

\noindent\phantomsection1.\label{criterion:IV.3.1} When is $\mathfrak R$ measurable with absolute precision? Answer: when $R$ has a pure point spectrum.

\noindent\phantomsection2.\label{criterion:IV.3.2} When are $\mathfrak R,\mathfrak S$ simultaneously measurable with absolute precision? Answer: when $R,S$ have pure point spectra and commute.

\noindent\phantomsection3.\label{criterion:IV.3.3} When are several quantities $\mathfrak R,\mathfrak S,\ldots$ simultaneously measurable with absolute precision? Answer: when $R,S,\ldots$ have pure point spectra and all commute.

\noindent\phantomsection4.\label{criterion:IV.3.4} When are several quantities $\mathfrak R,\mathfrak S,\ldots$ simultaneously measurable with arbitrary precision? Answer: when $R,S,\ldots$ all commute.

We used the following principle, abstracted from the result of the Compton--Simon experiment:

\noindent\phantomsection\textbf{(M.)}\label{post:IV.3.M} If the physical quantity $\mathfrak R$ is measured twice in immediate succession on a system $\boldsymbol{S}$, the same value is obtained both times. This is the case even though $\mathfrak R$ may have dispersion in the original state of $\boldsymbol{S}$, and the measurement of $\mathfrak R$ may also alter the state of $\boldsymbol{S}$.

We discussed the physical significance of \hyperref[post:IV.3.M]{M} in detail in~\hyperref[sec:III.3]{III.3}. Further assumptions used in answering \hyperref[criterion:IV.3.1]{1}--\hyperref[criterion:IV.3.4]{4} were: the statistical formula \hyperref[post:III.1.E2]{$E_2$} of~\hyperref[sec:III.3]{III.3} for states; assumption \hyperref[post:III.1.F]{F} there, according to which $F(\mathfrak R)$ has operator $F(R)$ when $\mathfrak R$ has operator $R$; and the assumption that $R+S$ has the operator $\mathfrak R+\mathfrak S$ when the simultaneously observable quantities $\mathfrak R,\mathfrak S$ have the respective operators $R,S$.\tnote{The source reverses quantity and operator in this sentence: $\mathfrak R+\mathfrak S$ is the quantity, and $R+S$ its operator. The printed order is retained.}
% The last operator/quantity order follows the source's apparent transposition, original p.177.

Since these three assumptions are again at our disposal---the first follows from~\hyperref[eq:IV.2.Sp]{Sp.} in~\hyperref[sec:IV.2]{IV.2}, while the other two correspond to \hyperref[cond:IV.2.I]{I}, \hyperref[cond:IV.2.II]{II} there---and \hyperref[post:IV.3.M]{M} is also to be accepted as valid, because we have recognized it as indispensable to the conceptual construction of quantum mechanics, the proofs given in~\hyperref[sec:III.3]{III.3} for \hyperref[criterion:IV.3.1]{1}--\hyperref[criterion:IV.3.4]{4} apply here as well. The stated answers are therefore valid again.

In~\hyperref[sec:III.5]{III.5} we examined the physical quantities that take only two values, $0,1$. These corresponded one-to-one to properties $\mathfrak E$. For if $\mathfrak E$ was given, the quantity could be defined as follows: it is measured by deciding whether $\mathfrak E$
% Source PDF page 185 / original page 178
\hypertarget{origpage-178}{}%
holds or not, and its value is then $1$ or $0$, respectively. Conversely, if the quantity was given, $\mathfrak E$ was the property that this quantity has the value $1$, that is, not $0$. It followed from \hyperref[post:III.5.F]{F} there, that is, \hyperref[cond:IV.2.I]{I} in~\hyperref[sec:IV.2]{IV.2}, that the corresponding operators $E$ are precisely the projection operators. The probability that $\mathfrak E$ holds was equal to the expectation of the quantity defined above. In~\hyperref[sec:III.5]{III.5} it was calculated only for states $\varphi$, that is, $U=P_{[\varphi]}$, $\|\varphi\|=1$; but we can determine it generally from \hyperref[eq:IV.2.Sp]{Sp.}: it is $\operatorname{Spur}(UE)$, relative, and absolute only if $U$ is properly normalized, $\operatorname{Spur}U=1$.

Since we have established \hyperref[criterion:IV.3.1]{1}--\hyperref[criterion:IV.3.4]{4}, the statements \hyperref[post:III.5.alpha]{$\alpha)$}--\hyperref[post:III.5.zeta]{$\zeta)$} derived from them in~\hyperref[sec:III.5]{III.5} are also established. We must observe, however, that \hyperref[post:III.5.alpha]{$\alpha)$} there gave information only for states, whereas we have just extended it to all mixtures:

\noindent\phantomsection$\alpha')$\label{prop:IV.3.alphaprime} In the mixture with statistical operator $U$, the property $\mathfrak E$ holds or fails to hold with probabilities
\[
\operatorname{Spur}(UE)\qquad\text{and}\qquad\operatorname{Spur}(U(1-E)),
\]
respectively. These are relative probabilities, absolute only if $U$ is properly normalized, $\operatorname{Spur}U=1$.

Suppose several quantities $\mathfrak R_1,\ldots,\mathfrak R_l$ are investigated, with respective operators $R_1,\ldots,R_l$ and resolutions of the identity $E_1(\lambda),\ldots,E_l(\lambda)$. Let $l$ intervals
\[
I_1:\lambda'_1<\lambda\leq\lambda''_1,\quad\ldots,\quad
I_l:\lambda'_l<\lambda\leq\lambda''_l
\]
be given, and put $E_1(I_1)=E_1(\lambda''_1)-E_1(\lambda'_1),\ldots,E_l(I_l)=E_l(\lambda''_l)-E_l(\lambda'_l)$. The properties ``$\mathfrak R_1$ lies in $I_1$,'' \ldots, ``$\mathfrak R_l$ lies in $I_l$'' then have the respective projection operators $E_1(I_1),\ldots,E_l(I_l)$; see \hyperref[post:III.5.zeta]{$\zeta)$}. Their simultaneous decidability is characterized by the commutativity of $E_1(I_1),\ldots,E_l(I_l)$; see \hyperref[post:III.5.gamma]{$\gamma)$}. The projection operator for their simultaneous occurrence is $E=E_1(I_1)\cdots E_l(I_l)$; see \hyperref[post:III.5.epsilon]{$\varepsilon)$}. Its probability is therefore $\operatorname{Spur}(UE)$; see \hyperref[prop:IV.3.alphaprime]{$\alpha')$}.

Let us now take the reverse path. Suppose we do not know the state of a system $\boldsymbol{S}$, but have made certain measurements on $\boldsymbol{S}$ and know their results. This is, in fact, always the situation in reality, since only measurement results can tell us anything about the state of $\boldsymbol{S}$. Strictly speaking, states are only a theoretical construction; only measurement results are actually available, and the task of physics is to establish relations between past and future measurement results. Admittedly, this is always done by introducing the auxiliary concept of ``state.'' But physical theory must then tell us, on the one hand, how to infer the present state from past measurement results and, on the other, how to infer future measurement results from the present state. So far we have concerned ourselves only
% Source PDF page 186 / original page 179
\hypertarget{origpage-179}{}%
with the second question, and must now turn to the first.

If the past measurements do not suffice to determine the present state uniquely, they may nevertheless tell us the probabilities with which the individual states occur. This holds in causal theories, such as classical mechanics, just as it does in quantum mechanics. The actual task is therefore this: for given measurement results, find a mixture whose statistics are the same as those to be expected for a system $\boldsymbol{S}$ of which we know only that the specified measurements were performed on it and yielded the specified results. We must, however, state more precisely what it means to ``know only this'' about $\boldsymbol{S}$, and nothing more, and how this is to lead to statistics.

The connection with statistics must in any case be the following. If these measurements yielded the specified results for many systems $\boldsymbol{S}'_1,\ldots,\boldsymbol{S}'_M$, all specimens of $\boldsymbol{S}$, then the ensemble $[\boldsymbol{S}'_1,\ldots,\boldsymbol{S}'_M]$ agrees in all its statistical properties with the mixture associated above with the measurement results. By \hyperref[post:IV.3.M]{M}, we may imagine the identical results for all $\boldsymbol{S}'_1,\ldots,\boldsymbol{S}'_M$ to have arisen as follows. There was originally a large ensemble $[\boldsymbol{S}_1,\ldots,\boldsymbol{S}_N]$ on which the measurements were made; those elements for which the desired results appeared were then collected into a new ensemble, namely $[\boldsymbol{S}'_1,\ldots,\boldsymbol{S}'_M]$. Everything now depends, of course, on how $[\boldsymbol{S}_1,\ldots,\boldsymbol{S}_N]$ is chosen. This initial ensemble specifies, so to speak, the a priori probabilities of the individual states of the system $\boldsymbol{S}$. The situation is naturally familiar from general probability theory: to infer states from measurement results, that is, causes from effects, or to calculate a posteriori probabilities, one must know the a priori probabilities. In general these can be chosen in many ways, so that our problem, too, has no unique answer. We shall see, however, that under the special circumstances of quantum mechanics one particular choice of the initial ensemble $[\boldsymbol{S}_1,\ldots,\boldsymbol{S}_N]$, and hence of the a priori probabilities, is singled out.

The situation is different if enough measurement results are available to determine the state of $\boldsymbol{S}$ completely: then every question must have a unique answer. We shall shortly see how this fact comes into play.

Finally, we mention the following. Instead of saying that several measurement results concerning $\boldsymbol{S}$ are known, we can also say that $\boldsymbol{S}$ has been tested for a certain property $\mathfrak E$ and that its presence
% Source PDF page 187 / original page 180
\hypertarget{origpage-180}{}%
has been established. We know from \hyperref[post:III.5.alpha]{$\alpha)$}--\hyperref[post:III.5.zeta]{$\zeta)$} how these matters are connected. For example, if the simultaneously ascertainable measurement results say that the values of $\mathfrak R_1,\ldots,\mathfrak R_l$ lie in the respective intervals $I_1,\ldots,I_l$, then, in the notation used earlier, the projection operator of $\mathfrak E$ is $E=E_1(I_1)\cdots E_l(I_l)$.

Our knowledge of $\boldsymbol{S}$ therefore always amounts to the presence of a certain property $\mathfrak E$, formally characterized by specifying the projection operator $E$. We seek $U$, the statistical operator of the equivalent ensemble $[\boldsymbol{S}'_1,\ldots,\boldsymbol{S}'_M]$, and $U_0$, the statistical operator of the general initial ensemble $[\boldsymbol{S}_1,\ldots,\boldsymbol{S}_N]$. What are the mathematical relations between $E,U,U_0$?

By \hyperref[post:IV.3.M]{M}, $\mathfrak E$ is certainly present in $[\boldsymbol{S}'_1,\ldots,\boldsymbol{S}'_M]$ as well; that is, the quantity belonging to $\mathfrak E$ has the value $1$. As we saw at the beginning of this section, this means $EU=U$, or that $Uf$ always lies in $\mathfrak M$, where $\mathfrak M$ is the set of all $f$ with $Ef=f$, the closed linear manifold belonging to $E$.

Instead of $EU=U$, we may also write $UE=U$, $U(1-E)=0$. Thus $Ug=0$ for all $g=(1-E)f$, that is, for all $g$ in the closed linear manifold belonging to $1-E$, or all $g$ in $\mathfrak R-\mathfrak M$. Consequently $Uf=0$ for all $f$ in $\mathfrak R-\mathfrak M$, while for $f$ in $\mathfrak M$, $Uf$ also lies in $\mathfrak M$. Nothing more can be said about $U$ in this way.

This determines $U$ essentially, that is, up to a constant factor, if and only if $\mathfrak M$ has zero or one dimension. Indeed, for $\mathfrak M=[0]$ it gives $U=0$, which is impossible by~\hyperref[sec:III.2]{III.2}, remark~1.\tnote{The source gives this reference as III.2, remark~1. The exclusion of $U=0$ is stated in \hyperref[remark:IV.2.2]{IV.2, remark~2}.} For $\mathfrak M=[\varphi]$, with $\varphi\ne0$ so that we may assume $\|\varphi\|=1$, we have $U\varphi=c\varphi$. Thus $Uf=cf$ for all $f$ in $\mathfrak M$, since these are of the form $a\varphi$; hence in general $Uf=UEf=cEf$, so $U=cE=cP_{[\varphi]}$. Since $c>0$, because $U$ is definite and nonzero, $U$ is essentially $E=P_{[\varphi]}$. If $\mathfrak M$ has dimension $\geq2$, we can choose normalized orthogonal $\varphi,\psi$ in $\mathfrak M$; then $P_{[\varphi]},P_{[\psi]}$ are two essentially different $U$ satisfying our condition. Thus $E=0$ is impossible; for $E=P_{[\varphi]}$, $\|\varphi\|=1$, we have $U=E=P_{[\varphi]}$; otherwise $U$ is not uniquely determined.

The fact that no $U$ at all can be found for $E=0$ would be troublesome if a system $\boldsymbol{S}$ could possess such a property $\mathfrak E$. But this is excluded by \hyperref[post:III.5.eta]{$\eta)$}: $\mathfrak E$ is never present, and its probability is always $0$. One-dimensional $\mathfrak M$, that is, $E=P_{[\varphi]}$ with $\|\varphi\|=1$, determines $U$ uniquely as the state $\varphi$. This, then, is the measurement that completely determines the state of $\boldsymbol{S}$, namely as $\varphi$, if its outcome is positive.\origNote{178} All other measurements are incomplete: they do not suffice to determine the state.

In the general case we proceed as follows. If we also call the quantity belonging to $\mathfrak E$ by the name $\mathfrak E$, then $U$ arises by
% Source PDF page 188 / original page 181
\hypertarget{origpage-181}{}%
measuring $\mathfrak E$ on the ensemble belonging to $U_0$, namely $[\boldsymbol{S}_1,\ldots,\boldsymbol{S}_N]$, and collecting all the elements for which the value $1$ is obtained into the ensemble of $U$, namely $[\boldsymbol{S}'_1,\ldots,\boldsymbol{S}'_M]$. The measurement of $\mathfrak E$ can be performed in many different ways, for example by measuring another quantity $\mathfrak R$ of which $\mathfrak E$ is a known function, $\mathfrak E=F(\mathfrak R)$. For example, let $\varphi_1,\varphi_2,\ldots$ be an orthonormal system spanning $\mathfrak M$, and $\psi_1,\psi_2,\ldots$ a corresponding system for $\mathfrak R-\mathfrak M$. Then $\varphi_1,\varphi_2,\ldots,\psi_1,\psi_2,\ldots$ spans $\mathfrak M+(\mathfrak R-\mathfrak M)=\mathfrak R$ and is therefore complete. Let $\lambda_1,\lambda_2,\ldots,\mu_1,\mu_2,\ldots$ be distinct real numbers, and define the operator $R$ by
\[
R\left(\sum_nx_n\varphi_n+\sum_nx_n\psi_n\right)
=\sum_n\lambda_nx_n\varphi_n+\sum_n\mu_ny_n\psi_n.
\]
% Source has x_n, not y_n, in the second sum on the left.
Evidently $R$ has the pure point spectrum $\lambda_1,\lambda_2,\ldots,\mu_1,\mu_2,\ldots$, with respective eigenfunctions $\varphi_1,\varphi_2,\ldots,\psi_1,\psi_2,\ldots$, and all eigenvalues are simple.\tnote{The source prints $x_n$ in the second sum on the left, although the corresponding sum on the right contains $y_n$. The formula is retained as printed.} If $F(x)$ is any function with
\[
F(\lambda_n)=1,\qquad F(\mu_n)=0,
\]
then $F(R)$ has eigenvalue $1$ for $\varphi_1,\varphi_2,\ldots$, and therefore for every $f$ in $\mathfrak M$; it has eigenvalue $0$ for $\psi_1,\psi_2,\ldots$, and therefore for every $f$ in $\mathfrak R-\mathfrak M$. Thus $E=F(R)$. If $R$ belongs to $\mathfrak R$, it follows that $\mathfrak E=F(\mathfrak R)$. The measurement of $\mathfrak E$ can therefore be carried out as a measurement of $\mathfrak R$.

In this case we can calculate the relation between $U_0$ and $U$. After the measurement of $\mathfrak R$, each system is in one of the states $\varphi_1,\varphi_2,\ldots,\psi_1,\psi_2,\ldots$, according as the value found was respectively $\lambda_1,\lambda_2,\ldots,\mu_1,\mu_2,\ldots$. The probabilities are respectively
\begin{gather*}
\operatorname{Spur}(U_0P_{[\varphi_1]})=(U_0\varphi_1,\varphi_1),\quad
\operatorname{Spur}(U_0P_{[\varphi_2]})=(U_0\varphi_2,\varphi_2),\quad\ldots,\\
\operatorname{Spur}(U_0P_{[\psi_1]})=(U_0\psi_1,\psi_1),\quad
\operatorname{Spur}(U_0P_{[\psi_2]})=(U_0\psi_2,\psi_2),\quad\ldots;
\end{gather*}
see the considerations of~\hyperref[sec:III.3]{III.3}, whose validity we have established. These fractions of the $U_0$ ensemble therefore pass into the ensembles $P_{[\varphi_1]},P_{[\varphi_2]},\ldots,P_{[\psi_1]},P_{[\psi_2]},\ldots$. Since $\mathfrak E=1$ corresponds to $\mathfrak R=\lambda_1,\lambda_2,\ldots$, the $U$ ensemble is obtained by combining the first group. Thus
\[
U=\sum_n(U_0\varphi_n,\varphi_n)P_{[\varphi_n]}.
\]
Now every $P_{[\varphi_n]}$ commutes with $R$,\origNote{179} so $U$ must do so too. If, therefore, $U$ does not commute with every $R$ formed as described above, certain measurement procedures, those based on the corresponding $\mathfrak R$, are excluded as ways of producing it from $U_0$. We then know more about $U$ than merely that it arose through a measurement of $\mathfrak E$. But since $U$ is to represent precisely this state of our knowledge, we try to retain this condition: if there are $U$ for which no measurement procedure for $\mathfrak E$ need be excluded, we give them preference. Let us therefore investigate whether such $U$ exist and what they are.

% Source PDF page 189 / original page 182
\hypertarget{origpage-182}{}%
As we saw, $U$ would have to commute with every $R$ produced in the stated manner. This implies
\[
RU\varphi_n=UR\varphi_n=U(\lambda_n\varphi_n)=\lambda_nU\varphi_n,
\]
so that $U\varphi_n$ is an eigenfunction of $R$ with eigenvalue $\lambda_n$, hence $U\varphi_n=a_n\varphi_n$. In particular, $U\varphi_1=a_1\varphi_1$. Given any $\varphi$ in $\mathfrak M$ with $\|\varphi\|=1$, we can choose $\varphi_1,\varphi_2,\ldots,\psi_1,\psi_2,\ldots$ so that $\varphi_1=\varphi$. Every such $\varphi$ is therefore an eigenfunction of $U$. Moreover, all these $\varphi$ belong to the same eigenvalue. For if $\varphi,\psi$ belonged to different eigenvalues, they would have to be orthogonal. Then $(\varphi+\psi)/\sqrt2$ would also be an eigenfunction, and since
\[
\left(\frac{\varphi+\psi}{\sqrt2},\varphi\right)=\frac{(\varphi,\varphi)}{\sqrt2}=\frac1{\sqrt2},\qquad
\left(\frac{\varphi+\psi}{\sqrt2},\psi\right)=\frac{(\psi,\psi)}{\sqrt2}=\frac1{\sqrt2},
\]
it would be orthogonal to neither and would therefore belong to the same eigenvalue as both, which is impossible since they were to belong to different ones. Thus $U\varphi=a\varphi$ with constant $a$. We can evidently drop $\|\varphi\|=1$, so that $Uf=af$ for all $f$ in $\mathfrak M$. Hence $UEg=aEg$ always, or $UE=aE$; but $U=UE$, so $U=aE$. Both $U,E$ are definite and nonzero, so $a>0$. Without changing $U$ essentially, we may therefore put $U=E$.

This $U$ actually solves the problem for every $R$, that is, every system $\varphi_1,\varphi_2,\ldots,\psi_1,\psi_2,\ldots$, if $U_0$ is suitably chosen. For $U_0=1$,
\[
\sum_n(U_0\varphi_n,\varphi_n)P_{[\varphi_n]}
=\sum_n(\varphi_n,\varphi_n)P_{[\varphi_n]}
=\sum_nP_{[\varphi_n]}=P_{\mathfrak M}=E=U.
\]
Thus $U=E$ is justified in the sense of the program outlined above. We can also determine $U_0$ if we assume it to be universal, independent of $E$ and $R$. The choice $U_0=1$ meets this requirement, and no other choice does. For
\begin{align*}
(U\varphi_m,\varphi_m)&=(E\varphi_m,\varphi_m)=(\varphi_m,\varphi_m)=1,\\
(U\varphi_m,\varphi_m)&=\sum_n(U_0\varphi_n,\varphi_n)(P_{[\varphi_n]}\varphi_m,\varphi_m)\\
&=\sum_n(U_0\varphi_n,\varphi_n)|(\varphi_n,\varphi_m)|^2
=(U_0\varphi_m,\varphi_m).
\end{align*}
Thus $(U_0\varphi_m,\varphi_m)=1$. Since every $\varphi$ in $\mathfrak M$ with $\|\varphi\|=1$ can be made $\varphi_1$, we have $(U_0\varphi,\varphi)=1$, whence $(U_0f,f)=(f,f)$ for all $f$ in $\mathfrak M$. Since $\mathfrak M$ is arbitrary, this holds for all $f$ whatsoever, and consequently $U_0=1$.

If two properties $\mathfrak E,\mathfrak F$ are given, not necessarily simultaneously ascertainable, then the probability that a system $\boldsymbol{S}$ in which the presence of $\mathfrak E$ has just been established also possesses $\mathfrak F$ on an immediately subsequent test is, by the above,
\[
\operatorname{Spur}(EF)=\sum(EF),
\]
where $E,F$ are the operators of $\mathfrak E,\mathfrak F$. The first formula holds because $U=E$; the second follows from $E^2=E$, $F^2=F$ and~\hyperref[sec:II.11]{II.11}. Incidentally, these probabilities
% Source PDF page 190 / original page 183
\hypertarget{origpage-183}{}%
are relative, with $\mathfrak E$ regarded as fixed and $\mathfrak F$ as variable. If $\operatorname{Spur}(E)=\sum(E)$, the dimension of $\mathfrak M$, is finite, they can be normalized by division by this number.

Instead of $\mathfrak E,\mathfrak F$, we can also consider physical quantities. Let $\mathfrak R_1,\ldots,\mathfrak R_j$ be simultaneously measurable quantities, and likewise $\mathfrak S_1,\ldots,\mathfrak S_l$. Let their respective operators be $R_1,\ldots,R_j,S_1,\ldots,S_l$, with respective resolutions of the identity $E_1(\lambda),\ldots,E_j(\lambda),F_1(\lambda),\ldots,F_l(\lambda)$. Let
\begin{gather*}
I_1:\lambda'_1<\lambda\leq\lambda''_1,\quad\ldots,\quad I_j:\lambda'_j<\lambda\leq\lambda''_j,\\
J_1:\mu'_1<\lambda\leq\mu''_1,\quad\ldots,\quad J_l:\mu'_l<\lambda\leq\mu''_l
\end{gather*}
be intervals, and put
\begin{gather*}
E_1(I_1)=E_1(\lambda''_1)-E_1(\lambda'_1),\quad\ldots,\quad E_j(I_j)=E_j(\lambda''_j)-E_j(\lambda'_j),\\
F_1(J_1)=F_1(\mu''_1)-F_1(\mu'_1),\quad\ldots,\quad F_l(J_l)=F_l(\mu''_l)-F_l(\mu'_l).
\end{gather*}
The question is this: $\mathfrak R_1,\ldots,\mathfrak R_j$ were measured on $\boldsymbol{S}$, and their values lay respectively in $I_1,\ldots,I_j$; how probable is it that an immediately subsequent measurement will find the values of $\mathfrak S_1,\ldots,\mathfrak S_l$ respectively in $J_1,\ldots,J_l$? Evidently we must put $E=E_1(I_1)\cdots E_j(I_j)$, $F=F_1(J_1)\cdots F_l(J_l)$; see \hyperref[post:III.5.epsilon]{$\varepsilon)$}, \hyperref[post:III.5.zeta]{$\zeta)$}. The required probability is then
\begin{align*}
&\operatorname{Spur}\bigl(E_1(I_1)\cdots E_j(I_j)F_1(J_1)\cdots F_l(J_l)\bigr)\\
&\qquad=\sum\bigl(E_1(I_1)\cdots E_j(I_j)F_1(J_1)\cdots F_l(J_l)\bigr).
\end{align*}

In conclusion, let us once more point out the significance of the general initial ensemble $U_0=1$. We obtained $U$ from it by separating it into two parts during the measurement of $\mathfrak R$. Had we not separated it---had we measured $\mathfrak R$ on every one of its elements and then reunited them all into an ensemble---the result would again have been $U_0=1$. This is easily checked directly, or proved by choosing $E=1$: the $\mu_1,\mu_2,\ldots$ and $\psi_1,\psi_2,\ldots$ then disappear, while $\lambda_1,\lambda_2,\ldots$ and $\varphi_1,\varphi_2,\ldots$ form a complete system. Thus although the measurement of $\mathfrak R$ may alter individual elements, all these changes must compensate one another exactly, since the ensemble as a whole does not change. This property, moreover, is characteristic of $U_0=1$. For if
\[
U_0=\sum_{n=1}^{\infty}(U_0\varphi_n,\varphi_n)P_{[\varphi_n]}
\]
for every complete orthonormal system $\varphi_1,\varphi_2,\ldots$, then $U_0$ commutes with $P_{[\varphi_1]}$, since each $P_{[\varphi_n]}$ does. Thus $U_0$ commutes with every $P_{[\varphi]}$, $\|\varphi\|=1$. Consequently
\[
U_0\varphi=U_0P_{[\varphi]}\varphi=P_{[\varphi]}U_0\varphi=(U_0\varphi,\varphi)\varphi,
\]
so that $\varphi$ is an eigenfunction of $U_0$. From this, $U_0=1$ follows exactly as $U=E$ was inferred earlier from the corresponding relation with $\mathfrak M,E$ in place of $\mathfrak R,1$.

In $U_0=1$, therefore, all possible states are in the most complete equilibrium possible; no intervention by measurement can alter it. For every
% Source PDF page 191 / original page 184
\hypertarget{origpage-184}{}%
complete orthonormal system $\varphi_1,\varphi_2,\ldots$,
\[
U_0=1=\sum_{n=1}^{\infty}P_{[\varphi_n]},
\]
that is, the $1:1:\cdots$ mixture of all the states $\varphi_1,\varphi_2,\ldots$. We see from this that $U_0=1$ corresponds to the thermodynamic assumption customary in the older quantum theory, the ``a priori equal probability of all simple quantum orbits.'' It will also play an important role in our thermodynamic considerations, to which the next sections are devoted.

\chapter{General Considerations}\label{chap:V}
\section{Measurement and Reversibility}\label{sec:V.1}
We can predict what happens to a mixture with statistical operator $U$ when a quantity $\mathfrak R$ with operator $R$ is measured in it, insofar as this question admits a unique answer at all. This means measuring the quantity on every element of the ensemble and then collecting all the individual systems that have undergone the measurement into an ensemble again.

First, let $R$ have a pure, simple point spectrum; let $\varphi_1,\varphi_2,\ldots$ be the complete orthonormal system of eigenfunctions, and $\lambda_1,\lambda_2,\ldots$ the corresponding eigenvalues, all distinct by assumption. After the measurement, the situation is this: in the fraction $(U\varphi_m,\varphi_m)$ of the original ensemble, $\mathfrak R$ has the value $\lambda_n$, $n=1,2,\ldots$.\tnote{The source uses $m$ in this fraction but $n$ for the corresponding value and state. The indices are retained as printed.} This fraction therefore forms an ensemble in which $\mathfrak R$ certainly has the value $\lambda_n$, by \hyperref[post:IV.3.M]{M} in~\hyperref[sec:IV.3]{IV.3}. It is consequently in the state $\varphi_n$, with the properly normalized statistical operator $P_{[\varphi_n]}$. Collecting the systems therefore produces a mixture with statistical operator
\[
U'=\sum_{n=1}^{\infty}(U\varphi_n,\varphi_n)P_{[\varphi_n]}.
\]

Second, let $R$ again have a pure point spectrum, with $\varphi_1,\varphi_2,\ldots$ and $\lambda_1,\lambda_2,\ldots$ having the same meanings as before, but now allow repeated eigenvalues: some $\lambda_n$ are equal. The measurement process for $\mathfrak R$ is then not uniquely defined, as was the case, for example, for $\mathfrak E$ in~\hyperref[sec:IV.3]{IV.3}. Indeed, let $\mu_1,\mu_2,\ldots$ be distinct real numbers, $S$ the operator corresponding to $\varphi_1,\varphi_2,\ldots$ and $\mu_1,\mu_2,\ldots$, and $\mathfrak S$ the associated quantity. If $F(x)$ is a function with
\[
F(\lambda_n)=\mu_n\qquad(n=1,2,\ldots),
\]
then $F(S)=R$, and therefore $F(\mathfrak S)=\mathfrak R$.\tnote{The source prints $F(\lambda_n)=\mu_n$. For the subsequent identity $F(S)=R$, with $S$ having eigenvalues $\mu_n$ and $R$ eigenvalues $\lambda_n$, the relation would instead be $F(\mu_n)=\lambda_n$. The printed formula is retained.} The measurement of $\mathfrak S$ is thus also to be regarded as a measurement of $\mathfrak R$. It transforms $U$ into the $U'$ given above,
% Source PDF page 192 / original page 185
\hypertarget{origpage-185}{}%
which is independent of the entirely arbitrary $\mu_1,\mu_2,\ldots$, but not of $\varphi_1,\varphi_2,\ldots$, even though the latter are not uniquely determined because of the multiplicity of the eigenvalues of $R$. In~\hyperref[sec:IV.2]{IV.2}, following~\hyperref[sec:II.8]{II.8}, we stated what can be said about $\varphi_1,\varphi_2,\ldots$. Let $\lambda',\lambda'',\ldots$ be the distinct values among $\lambda_1,\lambda_2,\ldots$, and $\mathfrak M_{\lambda'},\mathfrak M_{\lambda''},\ldots$ the sets of $f$ with $Rf=\lambda'f$, $Rf=\lambda''f$, and so on. If $\chi'_1,\chi'_2,\ldots$, $\chi''_1,\chi''_2,\ldots$, and so on are arbitrary orthonormal systems spanning the respective $\mathfrak M_{\lambda'},\mathfrak M_{\lambda''},\ldots$, then $\chi'_1,\chi'_2,\ldots,\chi''_1,\chi''_2,\ldots,\ldots$ is the most general $\varphi_1,\varphi_2,\ldots$ system. Thus, according to the choice of $\mathfrak S$, that is, the actual measurement arrangement, any
\[
U'=\sum_n(U\chi'_n,\chi'_n)P_{[\chi'_n]}
+\sum_n(U\chi''_n,\chi''_n)P_{[\chi''_n]}+\cdots
\]
may occur. This is unique only in special cases.

Let us also characterize this special case. Each separate sum in the preceding expression must then be uniquely determined. For every eigenvalue $\lambda$, if $\mathfrak M_\lambda$ is the set of $f$ with $Rf=\lambda f$, the sum $\sum_n(U\chi_n,\chi_n)P_{[\chi_n]}$ must be the same for every orthonormal system spanning $\mathfrak M_\lambda$. Call this sum $V$. Repeating verbatim the considerations of~\hyperref[sec:IV.3]{IV.3}, replacing the $U_0,U,\mathfrak M$ there by $U,V,\mathfrak M_\lambda$, shows that $V=c_\lambda P_{\mathfrak M_\lambda}$, with $c_\lambda$ constant and positive, and that this is equivalent to $(Uf,f)=c_\lambda(f,f)$ for all $f$ in $\mathfrak M_\lambda$. Since these $f$ are precisely the $P_{\mathfrak M_\lambda}g$ for all $g$, we require
\[
(UP_{\mathfrak M_\lambda}g,P_{\mathfrak M_\lambda}g)
=c_\lambda(P_{\mathfrak M_\lambda}g,P_{\mathfrak M_\lambda}g),
\]
or $(P_{\mathfrak M_\lambda}UP_{\mathfrak M_\lambda}g,g)=c_\lambda(P_{\mathfrak M_\lambda}g,g)$, that is,
\[
P_{\mathfrak M_\lambda}UP_{\mathfrak M_\lambda}=c_\lambda P_{\mathfrak M_\lambda}
\]
for every eigenvalue $\lambda$ of $R$. If this condition, which evidently restricts $U$ considerably, is not satisfied, different arrangements for measuring $\mathfrak R$ can actually transform $U$ into different $U'$. Nevertheless, in~\hyperref[sec:V.4]{V.4} we shall be able to say something further, on a thermodynamic basis, about the result of a measurement of $\mathfrak R$ in general.

Third, let $R$ not have a pure point spectrum. Then by~\hyperref[sec:III.3]{III.3}, or criterion~\hyperref[criterion:IV.3.1]{1} in~\hyperref[sec:IV.3]{IV.3}, it cannot be measured with absolute precision at all. As we discussed there, measurements of $\mathfrak R$ with limited precision are equivalent to measurements of quantities with pure point spectra.

The time-dependent Schr\"odinger differential equation describes a different kind of intervention in material systems. Whereas experiments or measurements act discontinuously, noncausally, and instantaneously, this equation describes how the system changes continuously and causally in the course of time when its total energy is known. For states $\varphi$, if $\mathsf{H}$ is the energy operator, it reads
\begin{equation}\tag{$Z_1$}\label{eq:V.1.Z1}
\frac{\partial}{\partial t}\varphi_t=-\frac{2\pi i}{h}\mathsf{H}\varphi_t.
\end{equation}

% Source PDF page 193 / original page 186
\hypertarget{origpage-186}{}%
For the statistical operator of the state $\varphi_t$, namely $U_t=P_{[\varphi_t]}$, this means
\begin{align*}
\left(\frac{\partial U_t}{\partial t}\right)f
&=\frac{\partial}{\partial t}(U_tf)
=\frac{\partial}{\partial t}\bigl((f,\varphi_t)\varphi_t\bigr)\\
&=\left(f,\frac{\partial\varphi_t}{\partial t}\right)\varphi_t
+(f,\varphi_t)\frac{\partial\varphi_t}{\partial t}\\
&=-\left(f,\frac{2\pi i}{h}\mathsf{H}\varphi_t\right)\varphi_t
-(f,\varphi_t)\frac{2\pi i}{h}\mathsf{H}\varphi_t\\
&=\frac{2\pi i}{h}\bigl((\mathsf{H}f,\varphi_t)\varphi_t-(f,\varphi_t)\mathsf{H}\varphi_t\bigr)
=\frac{2\pi i}{h}(U_t\mathsf{H}-\mathsf{H}U_t)f.
\end{align*}
\begin{equation}\tag{$Z_2$}\label{eq:V.1.Z2}
\frac{\partial U_t}{\partial t}=\frac{2\pi i}{h}(U_t\mathsf{H}-\mathsf{H}U_t).
\end{equation}
If $U_t$ is not a state but a mixture of several states, say $P_{[\varphi_t^{(1)}]},P_{[\varphi_t^{(2)}]},\ldots$ with respective weights $w_1,w_2,\ldots$, its evolution is determined by that of these $P_{[\varphi_t^{(1)}]},P_{[\varphi_t^{(2)}]},\ldots$. Adding the corresponding equations \hyperref[eq:V.1.Z2]{$Z_2$} shows that \hyperref[eq:V.1.Z2]{$Z_2$} holds for these $U_t$ as well. Since all $U$ are such mixtures or limiting cases of them---for example, every $U$ with finite $\operatorname{Spur}U$ is such a mixture---we may claim general validity for \hyperref[eq:V.1.Z2]{$Z_2$}.

In \hyperref[eq:V.1.Z2]{$Z_2$}, as in Schr\"odinger's differential equation \hyperref[eq:V.1.Z1]{$Z_1$}, $\mathsf{H}$ may also depend on $t$. If it does not, however, we can give the explicit solutions: for \hyperref[eq:V.1.Z1]{$Z_1$}, as we already know,
\begin{equation}\tag{$Z'_1$}\label{eq:V.1.Z1prime}
\varphi_t=e^{-\frac{2\pi i}{h}t\mathsf{H}}\varphi_0,
\end{equation}
and for \hyperref[eq:V.1.Z2]{$Z_2$},
\begin{equation}\tag{$Z'_2$}\label{eq:V.1.Z2prime}
U_t=e^{-\frac{2\pi i}{h}t\mathsf{H}}U_0e^{\frac{2\pi i}{h}t\mathsf{H}}.
\end{equation}
It is easy to verify that these are solutions and that they can be deduced from one another. That there is only one solution for given initial values $\varphi_0$ or $U_0$ is clear: the differential equations \hyperref[eq:V.1.Z1]{$Z_1$}, \hyperref[eq:V.1.Z2]{$Z_2$} are of first order in $t$.

We thus have two fundamentally different kinds of intervention that can be made on a system $\boldsymbol{S}$ or an ensemble $[\boldsymbol{S}_1,\ldots,\boldsymbol{S}_N]$. First, changes brought about by deliberate acts of measurement, represented by the formula
\begin{equation}\tag{1}\label{eq:V.1.1}
U\longrightarrow U'=\sum_{n=1}^{\infty}(U\varphi_n,\varphi_n)P_{[\varphi_n]},
\end{equation}
where $\varphi_1,\varphi_2,\ldots$ is a complete orthonormal system; see above. Second, automatic changes through the passage of time, represented by the formula
\begin{equation}\tag{2}\label{eq:V.1.2}
U\longrightarrow U_t=e^{-\frac{2\pi i}{h}t\mathsf{H}}Ue^{\frac{2\pi i}{h}t\mathsf{H}},
\end{equation}
% Source PDF page 194 / original page 187
\hypertarget{origpage-187}{}%
where $\mathsf{H}$ is the energy operator, $t$ the time, and $\mathsf{H}$ is independent of $t$. If $\mathsf{H}$ depends on $t$, divide the time interval under consideration into small intervals in which $\mathsf{H}$ does not change, or changes only very little, and apply~\hyperref[eq:V.1.2]{2} to each separately; composing these successive evolutions gives the final result.

We must now examine carefully these two kinds of intervention, their nature, and their relation to one another.

First, it is striking that~\hyperref[eq:V.1.2]{2} permits $\mathsf{H}$ to depend on time, in the manner just described. One would therefore expect~\hyperref[eq:V.1.2]{2} alone to suffice for describing the intervention produced by a measurement. A physical intervention can be nothing other than the temporary introduction of certain energy couplings into the system under consideration, that is, the introduction of a suitable time dependence of $\mathsf{H}$, prescribed by the observer. Why, then, do we need the special process~\hyperref[eq:V.1.1]{1} for measurement? The reason is this. During measurement we cannot consider the system $\boldsymbol{S}$ by itself. To follow its interaction with the measuring apparatus $\boldsymbol{M}$ by calculation, we must investigate the system $\boldsymbol{S}+\boldsymbol{M}$. The theory of measurement concerns $\boldsymbol{S}+\boldsymbol{M}$, since it must specify how the state of $\boldsymbol{S}$ is related to certain properties of the state of $\boldsymbol{M}$, namely the positions of certain pointers that the observer reads. It is, moreover, rather arbitrary whether we also include the observer in $\boldsymbol{M}$, and replace the relation between the state of $\boldsymbol{S}$ and the pointer positions in $\boldsymbol{M}$ by the relation to the chemical changes in the observer's eye, or even in the brain, that is, to what has been ``seen'' or ``apperceived.'' We shall examine this more closely in~\hyperref[sec:VI.1]{VI.1}. In any case,~\hyperref[eq:V.1.2]{2} is applicable only to $\boldsymbol{S}+\boldsymbol{M}$. We must, however, show that for $\boldsymbol{S}$ it yields the same result as applying~\hyperref[eq:V.1.1]{1} directly to $\boldsymbol{S}$. Only when this has succeeded shall we have obtained a unified view of the physical world on a quantum-mechanical basis. We defer the discussion of this question until~\hyperref[sec:VI.3]{VI.3}.

Second, concerning~\hyperref[eq:V.1.1]{1}, we have repeatedly pointed out that a measurement of the kind described there must be instantaneous: performed in so short a time that the change in $U$ due to~\hyperref[eq:V.1.2]{2} is not yet noticeable. Attempting to correct for this by calculating the altered $U_t$ according to~\hyperref[eq:V.1.2]{2} would gain nothing. For to use any $U_t$, we must know $t$, the instant of measurement, precisely; thus the duration of the measurement would still have to be short. This raises a fundamental difficulty, since it is well known that there is a quantity
% Source PDF page 195 / original page 188
\hypertarget{origpage-188}{}%
that in classical mechanics is canonically conjugate to time: energy.\origNote{180} We must therefore expect uncertainty relations for the canonically conjugate pair time--energy similar to those for the pair Cartesian coordinate--momentum.\origNote{181} Special relativity also shows that there must be a far-reaching analogy: the three spatial coordinates together with time form a ``four-vector,'' just as the three spatial momenta together with energy do. Such an uncertainty relation would mean that a very accurate energy measurement cannot be performed in a very short time. Specifically, for the measurement error $\varepsilon$ and the duration $\tau$, one should expect a relation of the form
\[
\varepsilon\tau\sim h.
\]
A physical discussion similar to that carried out in~\hyperref[sec:III.4]{III.4} for $p,q$ does indeed lead to this result.\textsuperscript{\hyperlink{note-181}{181}} Without entering into details, we mention only the case of a light quantum. By Bohr's frequency relation, its energy uncertainty $\varepsilon$ equals $h$ times its frequency uncertainty, $h\Delta\nu$. But, as discussed in note~\hyperlink{note-137}{137}, $\Delta\nu$ is at best the reciprocal of its duration, $1/\tau$, so $\varepsilon\gtrsim h/\tau$. To establish the monochromaticity of the light quantum throughout the interval $\tau$, the measurement must extend over the whole interval. The case of a light quantum is characteristic, since atomic energy levels are usually determined from the frequencies of the corresponding spectral lines. Since energy behaves in this way, other quantities $\mathfrak R$, too, may exhibit a connection between the precision and duration of measurement. How, then, can our assumption of instantaneous measurements be justified?

Above all, we must concede that this objection touches an essential weakness, indeed the principal weakness, of quantum mechanics: its nonrelativistic character, which singles out time $t$ from the three spatial coordinates $x,y,z$ and presupposes an objective concept of simultaneity. Whereas all other quantities, particularly $x,y,z$, which the Lorentz transformation links most closely with $t$, are represented by operators, time is represented by an ordinary numerical parameter $t$, just as in classical mechanics. Again, a system consisting of two particles has a wave function depending on their $2\times3=6$ spatial coordinates but only on one time $t$, although the Lorentz transformation would make two times desirable. That we could ignore the minimum duration of measurements imposed by natural law may therefore be connected with the nonrelativistic character of quantum mechanics. This is an explanation, but not a gratifying one.

% Source PDF page 196 / original page 189
\hypertarget{origpage-189}{}%
Closer examination shows, however, that the situation is not so bad after all. What we really need is not that $t$ be small, but only that, in calculating the probabilities $(U\varphi_n,\varphi_n)$, and hence in forming $U'=\sum_{n=1}^{\infty}(U\varphi_n,\varphi_n)P_{[\varphi_n]}$, it make little difference whether we start from $U$ itself or from some $U_t=e^{-\frac{2\pi i}{h}t\mathsf{H}}Ue^{\frac{2\pi i}{h}t\mathsf{H}}$. Since
\[
(U_t\varphi_n,\varphi_n)
=\left(e^{-\frac{2\pi i}{h}t\mathsf{H}}Ue^{\frac{2\pi i}{h}t\mathsf{H}}\varphi_n,\varphi_n\right)
=\left(Ue^{\frac{2\pi i}{h}t\mathsf{H}}\varphi_n,e^{\frac{2\pi i}{h}t\mathsf{H}}\varphi_n\right),
\]
we can achieve this, for example, by modifying $\mathsf{H}$ through suitable perturbing energies so that $e^{\frac{2\pi i}{h}t\mathsf{H}}\varphi_n$ differs from $\varphi_n$ only by a constant factor of absolute value $1$. In other words, the state $\varphi_n$ is to remain essentially constant under~\hyperref[eq:V.1.2]{2}, to be a stationary state. Equivalently, $\mathsf{H}\varphi_n$ is to equal a real constant times $\varphi_n$: $\varphi_n$ is to be an eigenfunction of $\mathsf{H}$. At first sight, such a modification of the energy operator $\mathsf{H}$, making the eigenfunctions of $R$ stationary and hence eigenfunctions of $\mathsf{H}$, so that $R,\mathsf{H}$ commute, may seem implausible. In fact it is not, and one can even see that typical measurement arrangements aim precisely at this modification of $\mathsf{H}$.

Every measurement ends with a light quantum or a mass point flying off with a certain energy in a certain direction, and indicating the result by these characteristics, that is, by its momenta; or with a mass point, such as a pointer on a scale, coming to rest, its Cartesian coordinates indicating the result. For a light quantum, in the terminology of~\hyperref[sec:III.6]{III.6}, specifying which $M_n$ is $1$, with the others $0$, and hence specifying all the values $M_1,M_2,\ldots$, is equivalent to specifying the quantity to be measured. For the mass point flying away, this is the specification of its three momentum components $P^x,P^y,P^z$; for the mass point at rest, of its three Cartesian coordinates $x,y,z$, or, as we called their operators, $Q^x,Q^y,Q^z$. But the measurement is genuinely complete only when the light quantum or mass point has really flown ``away'': when the light is no longer exposed to absorption, or the mass point to deflection by potential energies; or when the resting mass point really rests, which requires a large mass.\origNote{182} The uncertainty relation alone requires the latter, since the velocity must be close to zero and therefore have small dispersion, although its product with the mass, the momentum, has large dispersion because the coordinate dispersion is small. Pointers are indeed macroscopic objects, that is, enormous. Now the energy operator $\mathsf{H}$, as far as
% Source PDF page 197 / original page 190
\hypertarget{origpage-190}{}%
the light quantum is concerned, is, by~\hyperref[sec:III.6]{III.6},
\begin{align*}
&\sum_{n=1}^{\infty}h\varrho_nM_n\\
&\quad+\sum_{j=1}^{\infty}\sum_{n=1}^{\infty}w_{kj}^{n}
\left(\sqrt{M_n+1}\begin{pmatrix}k\to j\\M_n\to M_n+1\end{pmatrix}
+\sqrt{M_n}\begin{pmatrix}k\to j\\M_n\to M_n-1\end{pmatrix}\right);
\end{align*}
for the two mass-point examples it is
\[
\frac{(P^x)^2+(P^y)^2+(P^z)^2}{2m}+V(Q^x,Q^y,Q^z),
\]
where $m$ is the mass and $V$ the potential energy. Our criteria require, respectively, that the $w_{kj}^n$ vanish, that $V$ be constant, or that $m$ be very large. These are precisely what make the $M_n$, or $P^x,P^y,P^z$, or $Q^x,Q^y,Q^z$, commute with the $\mathsf{H}$ given above.

To close this discussion, we mention that making the states of actual interest stationary---here $\varphi_1,\varphi_2,\ldots$---also plays a role elsewhere in theoretical physics. The assumption, often unavoidable in physicochemical ``thought experiments,'' that chemical reactions can be interrupted, ``poisoned,'' is of this kind.\origNote{183}

The two interventions~\hyperref[eq:V.1.1]{1} and~\hyperref[eq:V.1.2]{2} differ fundamentally. That both are formally deterministic, and hence causal, is immaterial. Since we are considering the statistical properties of mixtures, it is not surprising that every change, even a statistical one, produces a causal change in probabilities and expectations. This is precisely why statistical ensembles and probabilities are introduced. What matters is that~\hyperref[eq:V.1.2]{2} does not increase the statistical uncertainty present in $U$, whereas~\hyperref[eq:V.1.1]{1} does. The former carries states into states,
\[
P_{[\varphi]}\quad\text{into}\quad P_{[e^{-\frac{2\pi i}{h}t\mathsf{H}}\varphi]},
\]
whereas~\hyperref[eq:V.1.1]{1} can very well transform states into mixtures. In this sense, the development of a state according to~\hyperref[eq:V.1.1]{1} is statistical, but according to~\hyperref[eq:V.1.2]{2} causal.

Furthermore, for fixed $\mathsf{H},t$,~\hyperref[eq:V.1.2]{2} is simply a unitary transformation of all $U$:
\[
U_t=AUA^{-1},\qquad A=e^{-\frac{2\pi i}{h}t\mathsf{H}}\quad\text{unitary}.
\]
Thus $Uf=g$ means $U_t(Af)=Ag$, so that $U_t$ arises from $U$ through the unitary transformation $A$ of Hilbert space: an isomorphism leaving all our fundamental geometrical concepts invariant; see the principles explained in~\hyperref[sec:I.4]{I.4}. It is therefore reversible: one need only replace $A$ by $A^{-1}$. This is possible, since the considerable freedom in choosing $\mathsf{H},t$ permits $A,A^{-1}$ to be regarded as entirely arbitrary unitary operators. Just as classical
% Source PDF page 198 / original page 191
\hypertarget{origpage-191}{}%
mechanics does, then, process~\hyperref[eq:V.1.2]{2} fails to reproduce one of the most essential and striking properties of the real world: its irreversibility, the fundamental difference between the temporal directions ``future'' and ``past.''

Process~\hyperref[eq:V.1.1]{1} behaves fundamentally differently. The transition
\[
U\longrightarrow U'=\sum_{n=1}^{\infty}(U\varphi_n,\varphi_n)P_{[\varphi_n]}
\]
cannot simply be reversed. We shall soon see that it is altogether irreversible, in the sense that in general it is not even possible, by repeatedly applying any processes~\hyperref[eq:V.1.1]{1} and~\hyperref[eq:V.1.2]{2}, ever to return from a given $U'$ to its $U$.

We have thus reached the point at which it is appropriate to introduce the thermodynamic viewpoint, since it alone permits a proper understanding of the difference between~\hyperref[eq:V.1.1]{1} and~\hyperref[eq:V.1.2]{2}, into which questions of reversibility evidently enter.

\section{Thermodynamic Considerations}\label{sec:V.2}
We shall investigate the thermodynamics of quantum-mechanical ensembles from two different standpoints. First, we assume the validity of the two laws of thermodynamics, that is, the impossibility of perpetual motion of the first and second kinds, the energy and entropy principles,\origNote{184} and calculate from them the entropy of every ensemble. We use the ordinary methods of phenomenological thermodynamics; quantum mechanics enters only insofar as our thermodynamic considerations concern objects whose behavior is governed by its laws, our ensembles and their statistical operators $U$. The validity of the two laws, however, is assumed, not proved. Subsequently we shall prove their validity in quantum mechanics. Since the energy principle holds in any case, only the entropy principle is at issue: we shall show that interventions~\hyperref[eq:V.1.1]{1} and~\hyperref[eq:V.1.2]{2} never decrease the entropy calculated by the first method. This order may seem somewhat unnatural, but it is justified because only the phenomenological thermodynamic discussion gives us the overview of the problem required for the second kind of consideration.

We therefore begin with the phenomenological treatment, which will also enable us to resolve a familiar paradox of classical thermodynamics. We point out at once that the far-fetched nature of our ``thought experiments''---their infeasibility in practice---does not detract
% Source PDF page 199 / original page 192
\hypertarget{origpage-192}{}%
from their force as proofs. In phenomenological thermodynamics, every conceivable process has evidential force if it does not violate the laws of thermodynamics.

The aim is to determine the entropy of an ensemble $[\boldsymbol{S}_1,\ldots,\boldsymbol{S}_N]$ with statistical operator $U$, assumed properly normalized, $\operatorname{Spur}U=1$. In the language of classical statistical mechanics, we are dealing with a Gibbs ensemble. Statistics and thermodynamics refer not to the interacting constituents of a single, very complicated mechanical system with many only imperfectly known degrees of freedom,\origNote{185} but to an ensemble of very many mechanical systems, each of which may, though need not, itself have an unmanageably large number of degrees of freedom, and which are to be regarded as completely separate and noninteracting.\origNote{186} Because the systems $\boldsymbol{S}_1,\ldots,\boldsymbol{S}_N$ are completely separate, and because we wish to apply to them the purely enumerative methods of probability theory, there can be no doubt that ordinary probability theory is to be used. The different counting methods of Bose--Einstein and Fermi--Dirac, appropriate to certain ensembles of indistinguishable and interacting individuals, namely light quanta, or electrons and protons respectively---see~\hyperref[sec:III.6]{III.6}, especially note~\hyperlink{note-147}{147}---do not come into consideration here.

The method introduced by Einstein for the thermodynamic treatment of such ensembles $[\boldsymbol{S}_1,\ldots,\boldsymbol{S}_N]$ is as follows.\origNote{187} Enclose each system $\boldsymbol{S}_1,\ldots,\boldsymbol{S}_N$ in a respective box $\boldsymbol{K}_1,\ldots,\boldsymbol{K}_N$ whose walls block the transmission of any physical influence, which is permissible because the systems do not interact. Each box is also to have a very large mass, so that any changes in state, and hence in energy and mass, of $\boldsymbol{S}_1,\ldots,\boldsymbol{S}_N$ affect its mass only slightly. This will also make their velocities in the thought experiments so small that we may calculate nonrelativistically. Place these boxes $\boldsymbol{K}_1,\ldots,\boldsymbol{K}_N$ in a very large box $\overline{\boldsymbol{K}}$: its volume $\mathcal V$ is to be much larger than the sum of the volumes of $\boldsymbol{K}_1,\ldots,\boldsymbol{K}_N$. For simplicity, let there be no force fields in $\overline{\boldsymbol{K}}$; in particular, it is to be far from all gravitational fields and so large that the masses of $\boldsymbol{K}_1,\ldots,\boldsymbol{K}_N$ have no effect. We may therefore regard $\boldsymbol{K}_1,\ldots,\boldsymbol{K}_N$, containing the respective $\boldsymbol{S}_1,\ldots,\boldsymbol{S}_N$, as molecules of a gas enclosed in the large box $\overline{\boldsymbol{K}}$. If we put $\overline{\boldsymbol{K}}$ in contact with a very large heat reservoir at temperature $\mathbf T$, its walls will also assume this temperature, and their actual molecules will enter the corresponding Brownian motion. They will therefore
% Source PDF page 200 / original page 193
\hypertarget{origpage-193}{}%
transfer momenta to the nearest of the $\boldsymbol{K}_1,\ldots,\boldsymbol{K}_N$, setting them in motion, and these will in turn transfer momenta to the others. Soon all $\boldsymbol{K}_1,\ldots,\boldsymbol{K}_N$ will be in motion, exchanging momenta with the wall molecules at the walls of $\overline{\boldsymbol{K}}$, and with one another through collisions inside $\overline{\boldsymbol{K}}$. A stationary equilibrium of these motions is reached only when $\boldsymbol{K}_1,\ldots,\boldsymbol{K}_N$ have acquired the velocity distribution in equilibrium with the Brownian motion of the wall molecules at temperature $\mathbf T$: the Maxwell velocity distribution of a gas at temperature $\mathbf T$, of which $\boldsymbol{K}_1,\ldots,\boldsymbol{K}_N$ are to be regarded as the ``molecules.''\origNote{188} We may therefore say that the $[\boldsymbol{S}_1,\ldots,\boldsymbol{S}_N]$ gas has assumed temperature $\mathbf T$. For brevity, we shall henceforth call the ensemble $[\boldsymbol{S}_1,\ldots,\boldsymbol{S}_N]$ with statistical operator $U$ the $U$ ensemble, and the corresponding gas the $U$ gas.

We consider such a gas because, by definition, the entropy difference between the $U$ ensemble and the $V$ ensemble---$U,V$ definite operators of trace $1$, with respective ensembles $[\boldsymbol{S}_1,\ldots,\boldsymbol{S}_N]$ and $[\boldsymbol{S}'_1,\ldots,\boldsymbol{S}'_N]$---must be determined by transforming the former reversibly into the latter.\origNote{189} This is best accomplished by the detour through the $U$ and $V$ gases. We assert that the entropy difference between the $U$ and $V$ ensembles is exactly the same as that between the $U$ and $V$ gases, provided both are considered at the same temperature $\mathbf T$, which may otherwise be arbitrary. If $\mathbf T$ is very close to zero, this is evidently true with arbitrary precision. For the difference between the $U$ ensemble and the $U$ gas at temperature zero vanishes, since the latter's $\boldsymbol{K}_1,\ldots,\boldsymbol{K}_N$ have no motion of their own, and the presence of the stationary boxes $\boldsymbol{K}_1,\ldots,\boldsymbol{K}_N,\overline{\boldsymbol{K}}$ is thermodynamically immaterial; likewise for $V$. We therefore attain our aim if we show that a given change in $\mathbf T$ changes the entropy of the $U$ gas by the same amount as that of the $V$ gas. The entropy change of a gas heated from $\mathbf T_1$ to $\mathbf T_2$ depends only on its caloric equation of state, more precisely on its heat capacities.\origNote{190} Of course we must not assume the gas to be ideal when, as here, $\mathbf T_1$ must be chosen close to zero.\origNote{191} But the two gases, $U$ and $V$, certainly have the same equation of state and the same heat capacities: from the standpoint of gas kinetics, the boxes $\boldsymbol{K}_1,\ldots,\boldsymbol{K}_N$ dominate and completely conceal the systems $\boldsymbol{S}_1,\ldots,\boldsymbol{S}_N$ or $\boldsymbol{S}'_1,\ldots,\boldsymbol{S}'_N$ enclosed in them. The difference between $U$ and $V$ is therefore unnoticeable during this heating process, and the two entropy differences agree, as asserted. Henceforth
% Source PDF page 201 / original page 194
\hypertarget{origpage-194}{}%
we shall compare only the $U$ and $V$ gases, choosing $\mathbf T$ so high that they may be regarded as ideal gases.\origNote{192} We can then account fully for their gas-kinetic behavior and turn to the actual task: transforming the $U$ gas reversibly into the $V$ gas. Unlike the processes used so far, this will require us to act on the $\boldsymbol{S}_1,\ldots,\boldsymbol{S}_N$ inside $\boldsymbol{K}_1,\ldots,\boldsymbol{K}_N$: we shall have to ``open'' the boxes.

First we show that all states $U=P_{[\varphi]}$ have the same entropy: the $P_{[\varphi]}$ ensemble can be reversibly transformed into the $P_{[\psi]}$ ensemble without consuming or producing heat energy. Mechanical energy must, of course, be consumed or produced if the expectation of the energy in $P_{[\varphi]}$ differs from that in $P_{[\psi]}$; see note~\hyperlink{note-185}{185}. For this we shall not even need the gases just considered: the transformation can also be made at temperature zero, with the ensembles themselves. Once this has been proved, we can and shall normalize the entropies of the $U$ ensembles so that all states have entropy zero.

The transformation of $P_{[\varphi]}$ into $P_{[\psi]}$ described above need not even be reversible. If it is not, the entropy difference must be $\geq$ the expression given in note~\hyperlink{note-185}{185}; see that note. It is therefore $\geq0$. Interchanging $P_{[\varphi]},P_{[\psi]}$ likewise shows that it is $\leq0$, and hence it is nevertheless zero.

The simplest procedure would be to use the time-dependent Schr\"odinger differential equation, our process~\hyperref[eq:V.1.2]{2}. We would have to find an energy operator $\mathsf{H}$ and a numerical value of $t$ such that the unitary operator $e^{-\frac{2\pi i}{h}t\mathsf{H}}$ takes $\varphi$ into $\psi$. In $t$ seconds, $P_{[\varphi]}$ would then pass of itself into $P_{[\psi]}$; the process is even reversible, and no heat has been involved; see~\hyperref[sec:V.1]{V.1}. We prefer, however, to avoid assumptions about the possible forms of energy operators $\mathsf{H}$ and to use only process~\hyperref[eq:V.1.1]{1}, interventions of the nature of measurements. The simplest possibility would be to measure on the $P_{[\varphi]}$ ensemble a quantity $\mathfrak R$ whose operator $R$ has a pure point spectrum of simple eigenvalues $\lambda_1,\lambda_2,\ldots$, with $\psi$ among its eigenfunctions $\psi_1,\psi_2,\ldots$, say $\psi_1=\psi$. Since this measurement transforms $\varphi$ into a mixture of the states $\psi_1,\psi_2,\ldots$, $\psi$ is among the states produced. This is not expedient, however, since $\psi_1=\psi$ occurs only with probability $|(\varphi,\psi)|^2$, while the fraction $1-|(\varphi,\psi)|^2$ passes into other states; for orthogonal $\varphi,\psi$, this is the whole ensemble. Nevertheless, a related
% Source PDF page 202 / original page 195
\hypertarget{origpage-195}{}%
procedure leads to the goal. By successively performing a large number of different measurements, we shall transform $P_{[\varphi]}$ into an ensemble differing arbitrarily little from $P_{[\psi]}$. As discussed above, it is immaterial that all these operations are, or at least may be, irreversible.

We assume $\varphi,\psi$ orthogonal; otherwise we could choose a $\chi$ orthogonal to both, with $\|\chi\|=1$, and pass from $\varphi$ to $\chi$ and then from $\chi$ to $\psi$. Let $k=1,2,\ldots$ be a number to be chosen later, and put
\[
\psi^{(\nu)}=\cos\frac{\pi\nu}{2k}\,\varphi+\sin\frac{\pi\nu}{2k}\,\psi
\qquad(\nu=0,1,\ldots,k).
\]
Evidently $\psi^{(0)}=\varphi$, $\psi^{(k)}=\psi$, and always $\|\psi^{(\nu)}\|=1$. Extend each $\psi^{(\nu)}$, $\nu=1,\ldots,k$, to a complete orthonormal system $\psi_1^{(\nu)},\psi_2^{(\nu)},\ldots$, with $\psi_1^{(\nu)}=\psi^{(\nu)}$. Let $R^{(\nu)}$ be an operator with a pure point spectrum of distinct eigenvalues, say $\lambda_1^{(\nu)},\lambda_2^{(\nu)},\ldots$, whose eigenfunctions are $\psi_1^{(\nu)},\psi_2^{(\nu)},\ldots$, and let $\mathfrak R^{(\nu)}$ be the corresponding quantity. Observe also that
\begin{align*}
(\psi^{(\nu-1)},\psi^{(\nu)})
&=\cos\frac{\pi(\nu-1)}{2k}\cos\frac{\pi\nu}{2k}
+\sin\frac{\pi(\nu-1)}{2k}\sin\frac{\pi\nu}{2k}\\
&=\cos\left(\frac{\pi\nu}{2k}-\frac{\pi(\nu-1)}{2k}\right)=\cos\frac\pi{2k}.
\end{align*}

Now measure $\mathfrak R^{(1)}$ on the ensemble $U^{(0)}=P_{[\psi^{(0)}]}=P_{[\varphi]}$, producing, say, $U^{(1)}$; then measure $\mathfrak R^{(2)}$ on $U^{(1)}$, producing $U^{(2)}$; and so on, finally measuring $\mathfrak R^{(k)}$ on $U^{(k-1)}$, producing $U^{(k)}$. That $U^{(k)}$ is arbitrarily close to $P_{[\psi^{(k)}]}=P_{[\psi]}$ for sufficiently large $k$ can be understood intuitively as follows. Measuring $\mathfrak R^{(\nu)}$ on $\psi^{(\nu-1)}$ carries the fraction
\[
|(\psi^{(\nu-1)},\psi^{(\nu)})|^2=\left(\cos\frac\pi{2k}\right)^2
\]
into $\psi^{(\nu)}$. Thus the successive measurement of $\mathfrak R^{(1)},\mathfrak R^{(2)},\ldots,\mathfrak R^{(k)}$ carries at least the fraction $(\cos(\pi/2k))^{2k}$ from $\psi^{(0)}=\varphi$, through $\psi^{(1)},\psi^{(2)},\ldots,\psi^{(k-1)}$, into $\psi=\psi^{(k)}$. Since $(\cos(\pi/2k))^{2k}\to1$ as $k\to\infty$, the resulting ensemble consists almost entirely of systems in state $\psi$ if $k$ is large enough. The exact proof is as follows. Process~\hyperref[eq:V.1.1]{1} does not change the trace; since $\operatorname{Spur}U^{(0)}=\operatorname{Spur}P_{[\varphi]}=1$, we have $\operatorname{Spur}U^{(1)}=\operatorname{Spur}U^{(2)}=\cdots=\operatorname{Spur}U^{(j)}=1$. On the other hand,
\begin{align*}
(U^{(\nu)}f,f)
&=\sum_n(U^{(\nu-1)}\psi_n^{(\nu)},\psi_n^{(\nu)})(P_{[\psi_n^{(\nu)}]}f,f)\\
&=\sum_n(U^{(\nu-1)}\psi_n^{(\nu)},\psi_n^{(\nu)})|(\psi_n^{(\nu)},f)|^2.
\end{align*}
For $\nu=1,\ldots,k-1$, take $f=\psi_1^{(\nu+1)}=\psi^{(\nu+1)}$, and for $\nu=k$, take $f=\psi_1^{(k)}=\psi^{(k)}=\psi$. Then
\begin{align*}
(U^{(\nu)}\psi^{(\nu+1)},\psi^{(\nu+1)})
&\geq(U^{(\nu-1)}\psi^{(\nu)},\psi^{(\nu)})|(\psi^{(\nu)},\psi^{(\nu+1)})|^2\\
&=\left(\cos\frac\pi{2k}\right)^2(U^{(\nu-1)}\psi^{(\nu)},\psi^{(\nu)}),\\
(U^{(k)}\psi^{(k)},\psi^{(k)})&=(U^{(k-1)}\psi^{(k)},\psi^{(k)}).
\end{align*}
% Source PDF page 203 / original page 196
\hypertarget{origpage-196}{}%
Together with
\[
(U^{(0)}\psi^{(1)},\psi^{(1)})=(P_{[\psi^{(0)}]}\psi^{(1)},\psi^{(1)})
=|(\psi^{(0)},\psi^{(1)})|^2=\left(\cos\frac\pi{2k}\right)^2,
\]
this gives\tnote{Here the source prints $U^{(j)}$; the preceding argument and the following limit concern $U^{(k)}$.}
\[
(U^{(j)}\psi,\psi)\geq\left(\cos\frac\pi{2k}\right)^{2k}.
\]
Since $\operatorname{Spur}U^{(k)}=1$ and $(\cos(\pi/2k))^{2k}\to1$ as $k\to\infty$, we can apply the result obtained in~\hyperref[sec:II.11]{II.11}: $U^{(k)}$ converges to $P_{[\psi]}$. Our goal is thereby attained.

We must next ask to what extent one of the principal devices used in phenomenological thermodynamic ``thought experiments''---the so-called semipermeable wall---can also be used for quantum-mechanical systems. Phenomenological thermodynamics admits the proposition that if I and II are two different states of the same system $\boldsymbol{S}$, one may assume the existence of a wall completely permeable to I and completely impermeable to II.\origNote{193} This is, so to speak, the thermodynamic definition of the difference, and hence also the identity, of two states. To what extent is such an assumption permissible in quantum mechanics?

We first show that if $\varphi_1,\varphi_2,\ldots,\psi_1,\psi_2,\ldots$ is an orthonormal system, there is a semipermeable wall that lets systems $\boldsymbol{S}$ in any of the states $\varphi_1,\varphi_2,\ldots$ pass unhindered, and reflects them unchanged in any of the states $\psi_1,\psi_2,\ldots$. Systems in other states, however, may even be altered on collision with the wall.

We may assume $\varphi_1,\varphi_2,\ldots,\psi_1,\psi_2,\ldots$ complete, since otherwise it could be completed by further $\chi_1,\chi_2,\ldots$, which could then be included among the $\varphi_1,\varphi_2,\ldots$. Choose an operator $R$ with a pure point spectrum of simple eigenvalues $\lambda_1,\lambda_2,\ldots,\mu_1,\mu_2,\ldots$, with the respective eigenfunctions $\varphi_1,\varphi_2,\ldots,\psi_1,\psi_2,\ldots$, and with $\lambda_n<0$, $\mu_n>0$. Let $\mathfrak R$ belong to $R$. Build many windows into the wall, each arranged as follows. Every molecule $\boldsymbol{K}_1,\ldots,\boldsymbol{K}_N$ of our gas---we again consider $U$ gases at temperatures $\mathbf T>0$---is held there and opened; $\mathfrak R$ is measured on the system $\boldsymbol{S}_1$, \ldots, or $\boldsymbol{S}_N$ inside it; the box is then closed again and, according as the measured value of $\mathfrak R$ was $<0$ or $>0$, the box and its contents are allowed through the window with unchanged momentum, or reflected. This arrangement clearly serves the intended purpose. It remains only to discuss what changes remain in it after such collisions, and
% Source PDF page 204 / original page 197
\hypertarget{origpage-197}{}%
to what extent it is or is not related to the so-called ``Maxwell demon'' of thermodynamics.\origNote{194}

On the first point: since measuring $\mathfrak R$ may alter the state of $\boldsymbol{S}$, and possibly its energy expectation, the measuring apparatus must supply or absorb this difference in mechanical energy in accordance with the first law. For example, it may contain a spring that is released or tensioned, or something similar. Since this is a purely automatic measuring mechanism, and only mechanical energy, not heat, is converted, certainly no entropy changes occur here; this alone is important for us now. If $\boldsymbol{S}$ is in one of the states $\varphi_1,\varphi_2,\ldots,\psi_1,\psi_2,\ldots$, the measurement of $\mathfrak R$ does not alter $\boldsymbol{S}$ at all, and no compensating changes remain in the measuring apparatus.

The second point is more troublesome. Our arrangement is rather similar to a ``Maxwell demon,'' a semipermeable wall that transmits molecules arriving from the right and reflects those arriving from the left. If such a wall is inserted in the middle of a vessel containing gas, all the gas soon lies on the left-hand side: the volume is halved without expenditure of entropy. This means an uncompensated decrease in the gas's entropy; by the second law, such a wall therefore cannot exist. Our semipermeable wall nevertheless differs essentially from this thermodynamically inadmissible one. It distinguishes the ``molecules'' $\boldsymbol{K}_1,\ldots,\boldsymbol{K}_N$ only by their internal properties---the states of the enclosed $\boldsymbol{S}_1$, \ldots, or $\boldsymbol{S}_N$---rather than by external properties, such as whether a molecule arrives from the right or the left. This is precisely the crucial point. A detailed analysis is possible on the basis of L.~Szilard's investigations, which clarify the nature of semipermeable walls, the ``Maxwell demon,'' and, generally, the role of ``interventions of intelligent beings in thermodynamic systems.'' We cannot enter into these matters here, especially since the reader will find a detailed account in the work cited in note~\hyperlink{note-194}{194}.

The foregoing shows, in particular, that two states $\varphi,\psi$ of the system $\boldsymbol{S}$ can certainly be separated by semipermeable walls if they are orthogonal. We now prove the converse: if $\varphi,\psi$ are not orthogonal, assuming such a semipermeable wall contradicts the second law. Thus the necessary and sufficient condition for separability by semipermeable walls is $(\varphi,\psi)=0$, and not, as in classical theory, $\varphi\ne\psi$; we write $\varphi,\psi$ in place of I, II used above. This resolves an old paradox of the classical form of thermodynamics, the troublesome discontinuity in using
% Source PDF page 205 / original page 198
\hypertarget{origpage-198}{}%
semipermeable walls: states differing arbitrarily little are still completely separable, while absolutely identical states are not separable at all. We now have a continuous transition: complete separation is possible only for $(\varphi,\psi)=0$, and separation becomes progressively less effective as $(\varphi,\psi)$ increases. Only when $(\varphi,\psi)$ reaches its maximum, namely $|(\varphi,\psi)|=1$, are the states $\varphi,\psi$ identical and separation completely impossible. For $\|\varphi\|=\|\psi\|=1$, so $|(\varphi,\psi)|=1$ implies $\varphi=c\psi$, with $c$ constant and $|c|=1$.

To carry out the announced argument, we must anticipate the final result of this section, the entropy of the $U$ ensemble. Naturally, we shall not use the present result in deriving that formula.

Assume, then, that a semipermeable wall separating $\varphi$ and $\psi$ exists; we are to deduce $(\varphi,\psi)=0$. Consider a $\tfrac12(P_{[\varphi]}+P_{[\psi]})$ gas, consisting of $N/2$ systems in state $\varphi$ and $N/2$ in state $\psi$; the operator has trace $1$. Choose $\mathcal V$, that is, $\overline{\boldsymbol{K}}$, and $\mathbf T$ so that the gas is ideal. Let $\overline{\boldsymbol{K}}$ have the longitudinal section $12341$ indicated in Fig.~\ref{fig:3}. Insert the semipermeable wall at one end, $aa$, and push it to the middle, $bb$. Keep the gas's temperature constant by contact at the other end, $23$, with a large heat reservoir $\boldsymbol{W}$ at temperature $\mathbf T$.
\begin{figure}[tbp]
\centering
\includegraphics[width=0.55\linewidth]{figures/figure-3.png}
\caption{}\label{fig:3}
\end{figure}
Nothing happens to the $\varphi$ molecules in this process, while the $\psi$ molecules are pushed into the right half of $\overline{\boldsymbol{K}}$, between $bb$ and $23$. Thus the $\tfrac12(P_{[\varphi]}+P_{[\psi]})$ gas is a $1:1$ mixture of a $P_{[\varphi]}$ gas and a $P_{[\psi]}$ gas. Nothing happens to the former; the latter is compressed isothermally to half its volume. The equation of state of the ideal gas shows that the mechanical work performed is $(N/2)\varkappa\mathbf T\ln2$, where $N/2$ is the number of molecules of the $P_{[\psi]}$ gas and $\varkappa$ is Boltzmann's constant.\origNote{195} Since the gas's energy does not change in this isothermal process,\origNote{196} this amount of energy is taken up by the heat reservoir $\boldsymbol{W}$. The entropy increase of the reservoir is therefore $\mathbf Q/\mathbf T=N\varkappa\cdot\tfrac12\ln2$; see note~\hyperlink{note-186}{186}.

After this process, half the original $P_{[\varphi]}$ gas, $N/4$ molecules, is to the left of $bb$; to its right are the
% Source PDF page 206 / original page 199
\hypertarget{origpage-199}{}%
other half of the original $P_{[\varphi]}$ gas, $N/4$ molecules, and the whole $P_{[\psi]}$ gas, $N/2$ molecules: altogether $3N/4$ molecules of a $\tfrac13P_{[\varphi]}+\tfrac23P_{[\psi]}$ gas. Compress the former gas and expand the latter to volumes $\mathcal V/4$ and $3\mathcal V/4$, respectively. The mechanical work is again taken from or supplied to the heat reservoir $\boldsymbol{W}$, respectively; it amounts to $(N/4)\varkappa\mathbf T\ln2$ and $(3N/4)\varkappa\mathbf T\ln(3/2)$; see note~\hyperlink{note-191}{191}. The entropy increase of the reservoir is therefore $N\varkappa\cdot\tfrac14\ln2$ and $-N\varkappa\cdot\tfrac34\ln(3/2)$, respectively. Altogether,
\[
N\varkappa\left(\tfrac12\ln2+\tfrac14\ln2-\tfrac34\ln\tfrac32\right)
=N\varkappa\cdot\tfrac34\ln\tfrac43.
\]

At the end we have a $P_{[\varphi]}$ gas of $N/4$ molecules in volume $\mathcal V/4$, and a $\tfrac13P_{[\varphi]}+\tfrac23P_{[\psi]}$ gas of $3N/4$ molecules in volume $3\mathcal V/4$. Originally there was a $\tfrac12P_{[\varphi]}+\tfrac12P_{[\psi]}$ gas of $N$ molecules in volume $\mathcal V$, or, if we wish, two such gases of $N/4$ and $3N/4$ molecules in volumes $\mathcal V/4$ and $3\mathcal V/4$, respectively. The change produced by the whole process is thus this: $N/4$ molecules in volume $\mathcal V/4$ passed from a $\tfrac12P_{[\varphi]}+\tfrac12P_{[\psi]}$ gas into a $P_{[\varphi]}$ gas; $3N/4$ molecules in volume $3\mathcal V/4$ passed from a $\tfrac12P_{[\varphi]}+\tfrac12P_{[\psi]}$ gas into a $\tfrac13P_{[\varphi]}+\tfrac23P_{[\psi]}$ gas; and the entropy of $\boldsymbol{W}$ increased by $N\varkappa\cdot\tfrac34\ln(4/3)$. Since the process was reversible, the total entropy increase must be zero: the two gas entropy changes must exactly compensate that of $\boldsymbol{W}$. We therefore still need to know the entropy changes of the gases.

As we shall see, a $U$ gas of $M$ molecules has entropy $-M\varkappa\operatorname{Spur}(U\ln U)$ if the entropy of a $P_{[\chi]}$ gas of the same volume and temperature is counted as zero; see above. If $U$ has a pure point spectrum with eigenvalues $w_1,w_2,\ldots$, this is $-M\varkappa\sum_{n=1}^{\infty}w_n\ln w_n$, taking $x\ln x=0$ at $x=0$. As is easily calculated, the operators
\[
P_{[\varphi]},\qquad\tfrac12P_{[\varphi]}+\tfrac12P_{[\psi]},\qquad
\tfrac13P_{[\varphi]}+\tfrac23P_{[\psi]}
\]
have the respective eigenvalues
\[
1,0;\qquad \frac{1+\alpha}2,\frac{1-\alpha}2,0;\qquad
\frac{3+\sqrt{1+8\alpha^2}}6,\frac{3-\sqrt{1+8\alpha^2}}6,0,
\]
where $\alpha=|(\varphi,\psi)|$, so $0\leq\alpha\leq1$. The eigenvalue zero always has infinite multiplicity, while the others are simple.\origNote{197} The entropy of the gases has therefore increased by
\begin{align*}
&-\frac N4\varkappa\cdot0
-\frac{3N}4\varkappa\Biggl[
\frac{3+\sqrt{1+8\alpha^2}}6\ln\frac{3+\sqrt{1+8\alpha^2}}6\\
&\qquad\qquad\qquad+\frac{3-\sqrt{1+8\alpha^2}}6\ln\frac{3-\sqrt{1+8\alpha^2}}6\Biggr]\\
&\qquad+N\varkappa\left(\frac{1+\alpha}2\ln\frac{1+\alpha}2
+\frac{1-\alpha}2\ln\frac{1-\alpha}2\right).
\end{align*}
% Source PDF page 207 / original page 200
\hypertarget{origpage-200}{}%
When the entropy increase $N\varkappa\cdot\tfrac34\ln(4/3)$ of $\boldsymbol{W}$ is added, the total must be zero. Dividing by $N\varkappa/4$, we obtain
\begin{align*}
&-\frac{3+\sqrt{1+8\alpha^2}}2\ln\frac{3+\sqrt{1+8\alpha^2}}6
-\frac{3-\sqrt{1+8\alpha^2}}2\ln\frac{3-\sqrt{1+8\alpha^2}}6\\
&\qquad+2(1+\alpha)\ln\frac{1+\alpha}2
+2(1-\alpha)\ln\frac{1-\alpha}2+3\ln\frac43=0,
\end{align*}
with $0\leq\alpha\leq1$.

It is easy to see that the left-hand side increases continually as $\alpha$ increases from $0$ to $1$,\origNote{198} from $0$ to $3\ln(4/3)$. Hence $\alpha=0$: if $\alpha\ne0$, the inverse of the described process would decrease entropy, contrary to the second law. Thus $(\varphi,\psi)=0$ is proved.

After these preparations we turn to determining the entropy of a $U$ gas of $N$ molecules in volume $\mathcal V$ at temperature $\mathbf T$, or, more precisely, its excess entropy over a $P_{[\varphi]}$ gas under the same conditions. By our earlier remarks and the normalization given above, this is the entropy of a $U$ ensemble of $N$ individual systems. As stated above, let $\operatorname{Spur}U=1$.
\begin{figure}[tbp]
\centering
\includegraphics[width=0.55\linewidth]{figures/figure-4.png}
\caption{}\label{fig:4}
\end{figure}
Then, as we know, $U$ has a pure point spectrum $w_1,w_2,\ldots$, with $w_1\geq0,w_2\geq0,\ldots$, $w_1+w_2+\cdots=1$. Let the corresponding eigenfunctions be $\varphi_1,\varphi_2,\ldots$; then $U=\sum_{n=1}^{\infty}w_nP_{[\varphi_n]}$; see~\hyperref[sec:IV.3]{IV.3}. Thus our $U$ gas is a mixture of $P_{[\varphi_1]},P_{[\varphi_2]},\ldots$ gases, consisting respectively of $w_1N,w_2N,\ldots$ molecules, all in the same volume $\mathcal V$. Again let $\mathbf T,\mathcal V$ be such that all these gases are ideal, and let $\overline{\boldsymbol{K}}$ have a rectangular cross section. Perform the following reversible interventions to separate the $\varphi_1,\varphi_2,\ldots$ molecules from one another; see Fig.~\ref{fig:4}. Attach to $\overline{\boldsymbol{K}}$, $23452$, a second rectangular box $\overline{\boldsymbol{K}}'$, $12561$, of the same size. Replace the common wall $25$ by two walls immediately adjacent to one another: one, $25$, fixed and semipermeable, transmitting $\varphi_1$ but reflecting $\varphi_2,\varphi_3,\ldots$; the other, $bb$, movable but an ordinary, completely impermeable wall. At $dd$, immediately adjacent to $34$, insert another semipermeable wall, transmitting $\varphi_2,\varphi_3,\ldots$ and reflecting $\varphi_1$. Then, keeping their distance apart fixed, move $bb$ and $dd$ to $aa$ and $cc$, respectively, immediately next to $16$ and $25$. The
% Source PDF page 208 / original page 201
\hypertarget{origpage-201}{}%
$\varphi_2,\varphi_3,\ldots$ molecules are entirely unaffected, while the $\varphi_1$ molecules are compelled to remain between the moving walls $bb,dd$. Since the distance between these walls is constant, we need perform no work against the gas pressure, and consequently no heat is produced. Finally, replace the walls $25,cc$ by a fixed, completely impermeable wall $25$, and remove $aa$, restoring the boxes $\overline{\boldsymbol{K}},\overline{\boldsymbol{K}}'$. All the $\varphi_1$ molecules are now in $\overline{\boldsymbol{K}}'$. We have thus ``drawn them off'' reversibly from $\overline{\boldsymbol{K}}$ into the equally large box $\overline{\boldsymbol{K}}'$, without any work, heat production, or change of temperature.\origNote{199}

Similarly, draw off the $\varphi_2,\varphi_3,\ldots$ molecules into equal boxes $\overline{\boldsymbol{K}}'',\overline{\boldsymbol{K}}''',\ldots$. We end with $P_{[\varphi_1]},P_{[\varphi_2]},\ldots$ gases, consisting of $w_1N,w_2N,\ldots$ molecules respectively, each in volume $\mathcal V$. Now compress them isothermally to volumes $w_1\mathcal V,w_2\mathcal V,\ldots$. As compensation, we must supply to them, from a large heat reservoir at temperature $\mathbf T$ so that the process is reversible, the quantities of heat $w_1N\varkappa\mathbf T\ln w_1$, $w_2N\varkappa\mathbf T\ln w_2$, and so on. These quantities are all negative: the mechanical compression work on each gas is their negative; see note~\hyperlink{note-191}{191}. The entropy change in these processes is therefore $\sum_{n=1}^{\infty}w_nN\varkappa\ln w_n$. Finally, transform all the $P_{[\varphi_1]},P_{[\varphi_2]},\ldots$ gases into a $P_{[\varphi]}$ gas, with $\varphi$ any chosen state. We now have only $P_{[\varphi]}$ gases, of $w_1N,w_2N,\ldots$ molecules in respective volumes $w_1\mathcal V,w_2\mathcal V,\ldots$. Since they are all identical and have the same density $N:\mathcal V$, we may mix them, again reversibly. This produces a $P_{[\varphi]}$ gas of $N$ molecules in volume $\mathcal V$, since $\sum_{n=1}^{\infty}w_n=1$.

We have therefore performed the desired reversible transition. The entropy changed by $N\varkappa\sum_{n=1}^{\infty}w_n\ln w_n$. Since its final value is zero under our normalization, its initial value was
\[
-N\varkappa\sum_{n=1}^{\infty}w_n\ln w_n.
\]
Since $U$ has eigenfunctions $\varphi_1,\varphi_2,\ldots$ with eigenvalues $w_1,w_2,\ldots$, $U\ln U$ has the same eigenfunctions with eigenvalues $w_1\ln w_1,w_2\ln w_2,\ldots$. Thus $\operatorname{Spur}(U\ln U)=\sum_{n=1}^{\infty}w_n\ln w_n$. Observe that $0\leq w_n\leq1$, so $w_n\ln w_n\leq0$, with equality only for $w_n=0,1$. We must take $w_n\ln w_n=0$ at $w_n=0$, because the vanishing $w_n$ need not be considered in our argument; continuity considerations give the same result.

% Source PDF page 209 / original page 202
\hypertarget{origpage-202}{}%
We have thus determined the entropy of a $U$ ensemble consisting of $N$ individual systems as
\[
-N\varkappa\operatorname{Spur}(U\ln U).
\]
What was just said about the $w_n\ln w_n$ shows that it is always $\geq0$. For it to be zero, every $w_n$ must be $0$ or $1$. Since $\operatorname{Spur}U=1$, exactly one $w_n$ is then $1$ and all the others are zero, so $U=P_{[\varphi_n]}$. Thus the $U=P_{[\varphi]}$, the states, have entropy zero; every other mixture has positive entropy.

\section{Questions of Reversibility and Equilibrium}\label{sec:V.3}
We can now prove the irreversibility of the measurement process asserted in~\hyperref[sec:V.1]{V.1}. For example, if $U$ is a state, $U=P_{[\varphi]}$, measuring a quantity $\mathfrak R$ whose operator $R$ has eigenfunctions $\varphi_1,\varphi_2,\ldots$ transforms it into the ensemble
\[
U'=\sum_{n=1}^{\infty}(P_{[\varphi]}\varphi_n,\varphi_n)P_{[\varphi_n]}
=\sum_{n=1}^{\infty}|(\varphi,\varphi_n)|^2P_{[\varphi_n]}.
\]
If $U'$ is not a state, entropy has increased: the entropy of $U$ was zero, while that of $U'$ is positive. The process is therefore irreversible. For $U'$ to be a state, and hence some $P_{[\varphi_{\bar n}]}$ since the $\varphi_n$ are its eigenfunctions, all $|(\varphi,\varphi_n)|^2$ must vanish except one, which must then equal $1$. Thus $\varphi$ is orthogonal to all $\varphi_n$ for $n\ne\bar n$. But then $\varphi=c\varphi_{\bar n}$ with $|c|=1$, so $P_{[\varphi]}=P_{[\varphi_{\bar n}]}$, $U=U'$. Thus every measurement on a state is irreversible except when the state is an eigenfunction of the measured quantity, that is, when this quantity has a sharp value in the given state; then the measurement leaves the state wholly unchanged. As we see, noncausal behavior is also unambiguously linked with certain accompanying thermodynamic phenomena.

We now investigate, in full generality, when process~\hyperref[eq:V.1.1]{1}, $U\to U'=\sum_{n=1}^{\infty}(U\varphi_n,\varphi_n)P_{[\varphi_n]}$, increases entropy.

The entropy of $U$ is $-N\varkappa\operatorname{Spur}(U\ln U)$. If $w_1,w_2,\ldots$ are its eigenvalues and $\psi_1,\psi_2,\ldots$ its eigenfunctions, this equals
\[
-N\varkappa\sum_{n=1}^{\infty}w_n\ln w_n
=-N\varkappa\sum_{n=1}^{\infty}(U\psi_n,\psi_n)\ln(U\psi_n,\psi_n).
\]
The eigenvalues of $U'$ are $(U\varphi_1,\varphi_1),(U\varphi_2,\varphi_2),\ldots$, so its entropy is
\[
-N\varkappa\sum_{n=1}^{\infty}(U\varphi_n,\varphi_n)\ln(U\varphi_n,\varphi_n).
\]
Thus the entropy of $U$ is $\lesseqgtr$ that of $U'$ according as
\begin{equation}\tag{*}\label{eq:V.3.star}
\sum_{n=1}^{\infty}(U\psi_n,\psi_n)\ln(U\psi_n,\psi_n)
\gtreqless\sum_{n=1}^{\infty}(U\varphi_n,\varphi_n)\ln(U\varphi_n,\varphi_n).
\end{equation}
We first show that $\geq$ always holds in~\hyperref[eq:V.3.star]{*}, so that $U\to U'$ never decreases entropy. This is thermodynamically
% Source PDF page 210 / original page 203
\hypertarget{origpage-203}{}%
clear, but for our later purposes it is important to have a purely mathematical proof as well. We proceed by keeping $U$, and with it $\psi_1,\psi_2,\ldots$, fixed, while allowing $\varphi_1,\varphi_2,\ldots$ to range over all complete orthonormal systems.

By continuity, we may first restrict ourselves to systems $\varphi_1,\varphi_2,\ldots$ in which only finitely many $\varphi_n$ differ from the corresponding $\psi_n$. Suppose, for example, $\varphi_n=\psi_n$ for $n>M$. The $\varphi_n$ with $n\leq M$ are then linear combinations of the $\psi_n$ with $n\leq M$, and conversely. Thus
\[
\psi_m=\sum_{n=1}^Mx_{mn}\varphi_n\qquad(m=1,\ldots,M),
\]
and the $M$-dimensional matrix $\{x_{mn}\}$ is evidently unitary.\tnote{Here and in the following proof, the source interchanges the roles of $\varphi$ and $\psi$ relative to the preceding definitions: the $\varphi_n$ are now treated as eigenfunctions of $U$. The subsequent notation is retained as printed. The following identity also prints $N$ as the upper limit and $w=1,\ldots,M$ rather than $m=1,\ldots,M$.} We have $(U\varphi_m,\varphi_m)=w_m$ and, as is easily calculated, $(U\psi_m,\psi_m)=\sum_{n=1}^Nw_n|x_{mn}|^2$, $w=1,\ldots,M$. We therefore have to prove
\[
\sum_{m=1}^Mw_m\ln w_m\geq
\sum_{m=1}^M\left(\sum_{n=1}^Mw_n|x_{mn}|^2\right)
\ln\left(\sum_{n=1}^Mw_n|x_{mn}|^2\right).
\]
The right-hand side is a continuous function of the $M^2$ bounded variables $x_{mn}$, so it has a maximum and attains it, subject to $\{x_{mn}\}$ being unitary. Since the left-hand side is its value at $x_{mn}=1$ for $m=n$ and $x_{mn}=0$ for $m\ne n$, we have to show that the maximum occurs at this system of $x_{mn}$.

Let $x^0_{mn}$, $m,n=1,\ldots,M$, be a system of values at which the maximum is attained. Multiplying the matrix $\{x^0_{mn}\}$ by the unitary matrix
\[
\begin{pmatrix}
\alpha&\beta&0&\cdots&0\\
-\overline\beta&\overline\alpha&0&\cdots&0\\
0&0&1&\cdots&0\\
\vdots&\vdots&\vdots&\ddots&\vdots\\
0&0&0&\cdots&1
\end{pmatrix},\qquad |\alpha|^2+|\beta|^2=1,
\]
produces a unitary matrix $\{x'_{mn}\}$, another admissible system. Put $\alpha=\sqrt{1-\varepsilon^2}$, $\beta=\theta\varepsilon$, with $\varepsilon$ real and $|\theta|=1$. We take $\varepsilon$ small, retaining terms in $1,\varepsilon,\varepsilon^2$ and neglecting $\varepsilon^3,\varepsilon^4,\ldots$. Then $\alpha\approx1-\tfrac12\varepsilon^2$, and in the new matrix
\begin{align*}
 x'_{1n}&\approx(1-\tfrac12\varepsilon^2)x^0_{1n}+\theta\varepsilon x^0_{2n},\\
 x'_{2n}&\approx-\overline\theta\varepsilon x^0_{1n}+(1-\tfrac12\varepsilon^2)x^0_{2n},\\
 x'_{mn}&=x^0_{mn}\qquad(m\geq3).
\end{align*}
% Source PDF page 211 / original page 204
\hypertarget{origpage-204}{}%
Consequently,
\begin{align*}
\sum_{n=1}^Mw_n|x'_{1n}|^2
&\approx\sum_{n=1}^Mw_n|x^0_{1n}|^2
+\sum_{n=1}^M2w_n\operatorname{Re}(\overline\theta x^0_{1n}\overline{x^0_{2n}})\varepsilon\\
&\quad+\sum_{n=1}^Mw_n(-|x^0_{1n}|^2+|x^0_{2n}|^2)\varepsilon^2,\\
\sum_{n=1}^Mw_n|x'_{2n}|^2
&\approx\sum_{n=1}^Mw_n|x^0_{2n}|^2
-\sum_{n=1}^M2w_n\operatorname{Re}(\overline\theta x^0_{1n}\overline{x^0_{2n}})\varepsilon\\
&\quad-\sum_{n=1}^Mw_n(-|x^0_{1n}|^2+|x^0_{2n}|^2)\varepsilon^2,\\
\sum_{n=1}^Mw_n|x'_{mn}|^2&=\sum_{n=1}^Mw_n|x^0_{mn}|^2\qquad(m\geq3).
\end{align*}
Substitute these expressions into $f(x)=x\ln x$, observing that $f'(x)=\ln x+1$, $f''(x)=1/x$, and add them all. This gives
\begin{align*}
&\sum_{m=1}^M\left(\sum_{n=1}^Mw_n|x'_{mn}|^2\right)\ln\left(\sum_{n=1}^Mw_n|x'_{mn}|^2\right)\\
&\quad\approx\sum_{m=1}^M\left(\sum_{n=1}^Mw_n|x^0_{mn}|^2\right)\ln\left(\sum_{n=1}^Mw_n|x^0_{mn}|^2\right)\\
&\qquad+\left[\ln\left(\sum_{n=1}^Mw_n|x^0_{1n}|^2\right)-\ln\left(\sum_{n=1}^Mw_n|x^0_{2n}|^2\right)\right]\\
&\qquad\qquad\cdot\sum_{n=1}^M2w_n\operatorname{Re}(\overline\theta x^0_{1n}\overline{x^0_{2n}})\varepsilon\\
&\qquad+\Biggl\{-\left[\ln\left(\sum_{n=1}^Mw_n|x^0_{1n}|^2\right)-\ln\left(\sum_{n=1}^Mw_n|x^0_{2n}|^2\right)\right]\\
&\qquad\qquad\cdot\left[\sum_{n=1}^Mw_n|x^0_{1n}|^2-\sum_{n=1}^Mw_n|x^0_{2n}|^2\right]\\
&\qquad\qquad+\frac12\left[\frac1{\sum_{n=1}^Mw_n|x^0_{1n}|^2}+\frac1{\sum_{n=1}^Mw_n|x^0_{2n}|^2}\right]\\
&\qquad\qquad\cdot\left[\sum_{n=1}^M2\operatorname{Re}(\overline\theta x^0_{1n}\overline{x^0_{2n}})\right]^2\Biggr\}\varepsilon^2.
\end{align*}
For the first term on the right to be the maximum, the coefficient of $\varepsilon$ must be zero and that of $\varepsilon^2$ must be $\leq0$.\tnote{In the last squared sum in this expansion, the source omits the weight $w_n$ that occurs in the preceding first-order term. The displayed formula is retained as printed.} The former has two factors,
\[
\ln\left(\sum_{n=1}^Mw_n|x^0_{1n}|^2\right)-\ln\left(\sum_{n=1}^Mw_n|x^0_{2n}|^2\right),
\qquad\sum_{n=1}^M2w_n\operatorname{Re}(\overline\theta x^0_{1n}\overline{x^0_{2n}}).
\]
If the first is zero, the first term in the coefficient of $\varepsilon^2$ vanishes; this term is always $\leq0$. The second term, which is evidently always $\geq0$, must then vanish for the whole coefficient to be $\leq0$. This means $\sum_{n=1}^M2w_n\operatorname{Re}(\overline\theta x^0_{1n}\overline{x^0_{2n}})=0$. Thus the second factor in the coefficient of $\varepsilon$ is zero in every case, which can also be written
\[
2\operatorname{Re}\left(\overline\theta\sum_{n=1}^Mw_nx^0_{1n}\overline{x^0_{2n}}\right)=0.
\]
A suitable choice of $\theta$ turns this into the absolute value of the sum, so the sum must vanish:
\[
\sum_{n=1}^Mw_nx^0_{1n}\overline{x^0_{2n}}=0.
\]

% Source PDF page 212 / original page 205
\hypertarget{origpage-205}{}%
Since we can replace $1,2$ by any two distinct $k,j=1,\ldots,M$, we obtain
\[
\sum_{n=1}^Mw_nx^0_{kn}\overline{x^0_{jn}}=0\qquad(k\ne j).
\]
Thus the unitary coordinate transformation with matrix $\{x^0_{mn}\}$ takes the diagonal matrix with diagonal elements $w_1,\ldots,w_n$ into diagonal form again. Since the diagonal elements are the multipliers, or eigenvalues, of the matrix, they do not change under a coordinate transformation; at most they can be permuted. Before the transformation they were $w_m$, $m=1,\ldots,M$; afterward they are $\sum_{n=1}^Mw_n|x^0_{mn}|^2$, $m=1,\ldots,N$. The sums
\[
\sum_{n=1}^Mw_n\ln w_n,\qquad
\sum_{m=1}^M\left(\sum_{n=1}^Mw_n|x^0_{mn}|^2\right)\ln\left(\sum_{n=1}^Mw_n|x^0_{mn}|^2\right)
\]
therefore have the same value. A maximum thus also occurs at $x_{mn}=1$ for $m=n$, $x_{mn}=0$ for $m\ne n$, as asserted.

Let us now determine when equality holds in~\hyperref[eq:V.3.star]{*}. If it holds, then
\[
\sum_{n=1}^{\infty}(U\chi_n,\chi_n)\ln(U\chi_n,\chi_n)
\]
attains its maximum not only for $\chi_n=\varphi_n$, $n=1,2,\ldots$, but also for $\chi_n=\psi_n$, $n=1,2,\ldots$, as $\chi_1,\chi_2,\ldots$ ranges over all complete orthonormal systems. In particular this holds when only the first $M$ of the $\psi_n$ are unitarily transformed among themselves, so that $\chi_n=\psi_n$ for $n>M$. Put $u_{mn}=(U\psi_m,\psi_n)$, $m,n=1,\ldots,M$. Let $v_1,\ldots,v_N$ be the eigenvalues of the finite matrix $\{u_{mn}\}$, also Hermitian and definite, and let $\{\alpha_{mn}\}$, $m,n=1,\ldots,M$, transform it into diagonal form. Let this transformation take $\psi_1,\ldots,\psi_M$ into $\omega_1,\ldots,\omega_M$, with
\[
\psi_m=\sum_{n=1}^M\alpha_{mn}\omega_n\qquad(m=1,\ldots,M).
\]
Then $U\omega_n=v_n\omega_n$, hence
\[
(U\omega_m,\omega_n)=\begin{cases}v_n,&m=n,\\0,&m\ne n.\end{cases}
\]
For $\xi_m=\sum_{n=1}^Mx_{mn}\omega_n$, $m=1,\ldots,M$, where $\{x_{mn}\}$ is also unitary, we therefore have $(U\xi_k,\xi_j)=\sum_{n=1}^Mv_nx_{kn}\overline{x_{jn}}$. By the assumption about $\psi_1,\ldots,\psi_N$,
\[
\sum_{n=1}^M\left(\sum_{n=1}^Mv_n|x_{mn}|^2\right)\ln\left(\sum_{n=1}^Mv_n|x_{mn}|^2\right)
\]
therefore attains its maximum at $x_{mn}=\alpha_{mn}$. What was proved above then gives $\sum_{n=1}^Mv_n\alpha_{kn}\overline{\alpha_{jn}}=0$ for $k\ne j$, that is, $(U\psi_k,\psi_j)=0$ for $k\ne j$, $k,j=1,\ldots,M$.

% Source PDF page 213 / original page 206
\hypertarget{origpage-206}{}%
This must hold for every $M$, so $U\psi_k$ is orthogonal to all $\psi_j$ with $k\ne j$, and hence equals $w'_k\psi_k$ for a constant $w'_k$. Thus $\psi_1,\psi_2,\ldots$ are eigenfunctions of $U$, with corresponding eigenvalues $w'_1,w'_2,\ldots$, a permutation of $w_1,w_2,\ldots$. Under these circumstances, however,
\[
U'=\sum_{n=1}^{\infty}(U\psi_n,\psi_n)P_{[\psi_n]}
=\sum_{n=1}^{\infty}w'_nP_{[\psi_n]}=U.
\]
We have therefore found:

Process~\hyperref[eq:V.1.1]{1}, $U\to U'=\sum_{n=1}^{\infty}(U\psi_n,\psi_n)P_{[\psi_n]}$, where $\psi_1,\psi_2,\ldots$ are the eigenfunctions of the operator $R$ of the measured quantity $\mathfrak R$, never decreases entropy. Indeed, it always increases it unless all $\psi_1,\psi_2,\ldots$ are eigenfunctions of $U$, in which case $U=U'$. In that case $U$ also commutes with $R$, and this too is characteristic, since it is equivalent to the existence of the common eigenfunction system $\psi_1,\psi_2,\ldots$; see~\hyperref[sec:II.10]{II.10}.

Thus process~\hyperref[eq:V.1.1]{1} is irreversible in every case in which it produces any change at all.

We now treat reversibility for processes~\hyperref[eq:V.1.1]{1} and~\hyperref[eq:V.1.2]{2} independently of phenomenological thermodynamics, as announced as the second part of our program in~\hyperref[sec:V.2]{V.2}. We already know the mathematical method that makes this possible. If the second law holds, entropy must equal $-N\varkappa\operatorname{Spur}(U\ln U)$, and this must not decrease in any process~\hyperref[eq:V.1.1]{1} or~\hyperref[eq:V.1.2]{2}. We must therefore now treat $-N\varkappa\operatorname{Spur}(U\ln U)$ purely as a mathematical quantity, independently of its interpretation as entropy, and determine its behavior under~\hyperref[eq:V.1.1]{1} and~\hyperref[eq:V.1.2]{2}.\origNote{200}

Under~\hyperref[eq:V.1.2]{2}, $U$ becomes $U_t=e^{-\frac{2\pi i}{h}t\mathsf{H}}Ue^{\frac{2\pi i}{h}t\mathsf{H}}$, or, denoting the unitary operator $e^{-\frac{2\pi i}{h}t\mathsf{H}}$ by $A$, $U\to U_t=AUA^{-1}$. Since $f\to Af$ is, by the unitarity of $A$, an isomorphic mapping of Hilbert space onto itself, taking every operator $P$ into $APA^{-1}$, we have in general $F(APA^{-1})=AF(P)A^{-1}$. Thus $U_t\ln U_t=A\cdot U\ln U\cdot A^{-1}$, and consequently $\operatorname{Spur}(U_t\ln U_t)=\operatorname{Spur}(U\ln U)$. Our quantity $-N\varkappa\operatorname{Spur}(U\ln U)$ is therefore constant under~\hyperref[eq:V.1.2]{2}. We have just determined what happens under~\hyperref[eq:V.1.1]{1}, without invoking the second law: if $U$ changes, $U\ne U'$, the quantity increases; if $U$ remains unchanged, $U=U'$, equivalently if $\psi_1,\psi_2,\ldots$ are eigenfunctions of $U$ or $U,R$ commute, it naturally remains unchanged too. Under an intervention composed of several processes~\hyperref[eq:V.1.1]{1} and~\hyperref[eq:V.1.2]{2}, in any number and order, $-N\varkappa\operatorname{Spur}(U\ln U)$ therefore remains unchanged if every process~\hyperref[eq:V.1.1]{1} is ineffective, producing no change; in all other cases it increases.
% Source PDF page 214 / original page 207
\hypertarget{origpage-207}{}%
Thus, if only interventions~\hyperref[eq:V.1.1]{1} and~\hyperref[eq:V.1.2]{2} are considered, every process~\hyperref[eq:V.1.1]{1} that produces any change at all is irreversible.

There are, incidentally, expressions simpler than $-\operatorname{Spur}(U\ln U)$ that do not increase under~\hyperref[eq:V.1.1]{1} and are constant under~\hyperref[eq:V.1.2]{2}, for example the largest eigenvalue of $U$. Under~\hyperref[eq:V.1.2]{2} it is invariant, like all eigenvalues of $U$. Under~\hyperref[eq:V.1.1]{1}, the eigenvalues $w_1,w_2,\ldots$ of $U$ become the eigenvalues $\sum_{n=1}^{\infty}w_n|x_{1n}|^2$, $\sum_{n=1}^{\infty}w_n|x_{2n}|^2$, \ldots of $U'$; see the earlier considerations of this section. By the unitarity of $\{x_{mn}\}$, $\sum_{n=1}^{\infty}|x_{1n}|^2=1$, $\sum_{n=1}^{\infty}|x_{2n}|^2=1$, and so on. All these numbers are therefore $\leq$ the largest $w_n$. A largest $w_n$ exists because all $w_n\geq0$ and $\sum_{n=1}^{\infty}w_n=1$ implies $w_n\to0$. Now it is readily possible to change $U$ so that $-\operatorname{Spur}(U\ln U)=-\sum_{n=1}^{\infty}w_n\ln w_n$ remains invariant while the largest $w_n$ decreases. We see that there are transitions permitted by phenomenological thermodynamics, and therefore actually performable by our gas processes, that can never be produced by successive applications of~\hyperref[eq:V.1.1]{1} and~\hyperref[eq:V.1.2]{2} alone. This shows that the introduction of the gas methods is unavoidable.

Instead of $-\operatorname{Spur}(U\ln U)$, we can also consider $\operatorname{Spur}(F(U))$ for suitable functions $F(x)$. That it increases under~\hyperref[eq:V.1.1]{1} when $U\ne U'$---it is naturally invariant for $U=U'$ and under~\hyperref[eq:V.1.2]{2}---can be proved as for $F(x)=-x\ln x$, provided $F(x)$ has the only properties of that function used above. These are $F''(x)<0$ and the monotonic decrease of $F'(x)$, the latter following from the former. Thus, for our nonthermodynamic considerations of irreversibility, we can use any $\operatorname{Spur}(F(U))$, provided $F(x)$ is concave, that is, $F''(x)<0$ in $0\leq x\leq1$, where all the eigenvalues of $U$ lie.

Finally, we should also show that mixing two ensembles $U,V$, say in the ratio $\alpha:\beta$, with $\alpha>0$, $\beta>0$, $\alpha+\beta=1$, does not decrease entropy, that is,
\[
-\operatorname{Spur}\bigl((\alpha U+\beta V)\ln(\alpha U+\beta V)\bigr)
\geq-\alpha\operatorname{Spur}(U\ln U)-\beta\operatorname{Spur}(V\ln V).
\]
This also holds for any concave $F(x)$ in place of $-x\ln x$.\tnote{The source calls these functions ``convex upward'' and explicitly requires $F''(x)<0$; this is concavity in current terminology.} The proof is left to the reader.

We wish to find the stationary equilibrium mixture, the mixture of maximum entropy, when the energy is given. The last condition naturally means that the expectation of energy
% Source PDF page 215 / original page 208
\hypertarget{origpage-208}{}%
is prescribed. For the method of thermodynamically investigating statistical ensembles described in note~\hyperlink{note-184}{184}, this is the only possible formulation. Thus we admit only mixtures whose $U$ satisfies $\operatorname{Spur}U=1$, $\operatorname{Spur}(U\mathsf{H})=\mathbf E$, where $\mathsf{H}$ is the energy operator and $\mathbf E$ the prescribed energy expectation. Subject to these constraints, $-N\varkappa\operatorname{Spur}(U\ln U)$ is to be maximized. We make the further simplifying assumption that $\mathsf{H}$ has a pure point spectrum, say eigenvalues $W_1,W_2,\ldots$ and eigenfunctions $\varphi_1,\varphi_2,\ldots$; multiple eigenvalues are allowed.

Let $\mathfrak R$ be a quantity whose operator $R$ has eigenfunctions $\varphi_1,\varphi_2,\ldots$, but all distinct eigenvalues. According to~\hyperref[eq:V.1.2]{2},\tnote{The source refers to process~2 here, although the measurement transformation displayed below is process~\hyperref[eq:V.1.1]{1}.} measuring $\mathfrak R$ transforms $U$ into $U'=\sum_{n=1}^{\infty}(U\varphi_n,\varphi_n)P_{[\varphi_n]}$. The quantity $-N\varkappa\operatorname{Spur}(U\ln U)$ increases unless $U=U'$, while $\operatorname{Spur}(U)$ and $\operatorname{Spur}(U\mathsf{H})$ do not change. The latter follows because the $\varphi_n$ are eigenfunctions of $\mathsf{H}$, so $(\mathsf{H}\varphi_m,\varphi_n)$ vanishes for $m\ne n$:
\begin{align*}
\operatorname{Spur}(U'\mathsf{H})
&=\sum_{n=1}^{\infty}(U\varphi_n,\varphi_n)\operatorname{Spur}(P_{[\varphi_n]}\mathsf{H})\\
&=\sum_{n=1}^{\infty}(U\varphi_n,\varphi_n)(\mathsf{H}\varphi_n,\varphi_n)\\
&=\sum_{m,n=1}^{\infty}(U\varphi_m,\varphi_n)(\mathsf{H}\varphi_n,\varphi_m)
=\operatorname{Spur}(U\mathsf{H}).
\end{align*}
The commutativity of $R,\mathsf{H}$, that is, simultaneous measurability of $\mathfrak R$ and energy, also requires this. Thus the desired maximum is unchanged if we restrict ourselves to the $U'$, statistical operators with eigenfunctions $\varphi_1,\varphi_2,\ldots$, and it is attained only among these.

We therefore have $U=\sum_{n=1}^{\infty}w_nP_{[\varphi_n]}$. Since $U,U\mathsf{H},U\ln U$ all have eigenfunctions $\varphi_n$, with respective eigenvalues $w_n,W_nw_n,w_n\ln w_n$, the problem is to maximize $-N\varkappa\sum_{n=1}^{\infty}w_n\ln w_n$ subject to $\sum_{n=1}^{\infty}w_n=1$, $\sum_{n=1}^{\infty}W_nw_n=\mathbf E$. This is precisely the optimization problem that arises in the corresponding equilibrium calculation in ordinary gas theory,\origNote{201} and it is solved in the same way. By the familiar rules for finding extrema, the maximizing values $w_1,w_2,\ldots$ must satisfy
\[
\frac{\partial}{\partial w_n}\left(\sum_{m=1}^{\infty}w_m\ln w_m\right)
+\alpha\frac{\partial}{\partial w_n}\left(\sum_{m=1}^{\infty}w_m\right)
+\beta\frac{\partial}{\partial w_m}\left(\sum_{m=1}^{\infty}W_mw_m\right)=0,
\]
where $\alpha,\beta$ are suitable constants and $n=1,2,\ldots$. Thus
\[
(\ln w_n+1)+\alpha+\beta W_n=0,\qquad
w_n=e^{-1-\alpha-\beta W_n}=ae^{-\beta W_n},
\]
introducing the constant $a=e^{-1-\alpha}$ in place of $\alpha$.

% Source PDF page 216 / original page 209
\hypertarget{origpage-209}{}%
From $\sum_{n=1}^{\infty}w_n=1$ we obtain $a=1/\sum_{n=1}^{\infty}e^{-\beta W_n}$, hence
\[
w_n=\frac{e^{-\beta W_n}}{\sum_{m=1}^{\infty}e^{-\beta W_m}}.
\]
Since $\sum_{n=1}^{\infty}W_nw_n=\mathbf E$, we must have
\[
\frac{\sum_{n=1}^{\infty}W_ne^{-\beta W_n}}{\sum_{n=1}^{\infty}e^{-\beta W_n}}=\mathbf E,
\]
which determines $\beta$. Introduce, as usual, the ``partition function''
\[
Z(\beta)=\sum_{n=1}^{\infty}e^{-\beta W_n}=\operatorname{Spur}(e^{-\beta \mathsf{H}});
\]
for this and what follows, see note~\hyperlink{note-197}{197}.\tnote{In the following identity for $Z'(\beta)$, the final trace term requires a minus sign. The formula is retained as printed.} Then
\[
Z'(\beta)=-\sum_{n=1}^{\infty}W_ne^{-\beta W_n}=\operatorname{Spur}(\mathsf{H}e^{-\beta \mathsf{H}}),
\]
and the condition for $\beta$ reads
\[
-\frac{Z'(\beta)}{Z(\beta)}=\mathbf E.
\]
We make the further assumption that $\sum_{n=1}^{\infty}e^{-\beta W_n}$ and $\sum_{n=1}^{\infty}W_ne^{-\beta W_n}$ converge for every $\beta>0$: $W_n\to\infty$ sufficiently rapidly as $n\to\infty$. For example, $W_n/\ln n\to\infty$ suffices. For $U$ itself we thus obtain
\[
U=\sum_{n=1}^{\infty}ae^{-\beta W_n}P_{[\varphi_n]}
=ae^{-\beta \mathsf{H}}=\frac{e^{-\beta \mathsf{H}}}{\operatorname{Spur}(e^{-\beta \mathsf{H}})}
=\frac{e^{-\beta \mathsf{H}}}{Z(\beta)}.
\]
The properties of the equilibrium ensemble $U$, determined by specifying $\mathbf E$ or $\beta$, and therefore depending, as they must, on one parameter, can now be found by the method customary in gas theory.

The entropy of our ensemble is
\begin{align*}
\mathbf S&=-N\varkappa\operatorname{Spur}(U\ln U)
=-N\varkappa\operatorname{Spur}\left(\frac{e^{-\beta \mathsf{H}}}{Z(\beta)}\ln\frac{e^{-\beta \mathsf{H}}}{Z(\beta)}\right)\\
&=-\frac{N\varkappa}{Z(\beta)}\operatorname{Spur}\bigl(e^{-\beta \mathsf{H}}(-\beta \mathsf{H}-\ln Z(\beta))\bigr)\\
&=\frac{\beta N\varkappa}{Z(\beta)}\operatorname{Spur}(\mathsf{H}e^{-\beta \mathsf{H}})
+\frac{\ln Z(\beta)N\varkappa}{Z(\beta)}\operatorname{Spur}(e^{-\beta \mathsf{H}})\\
&=N\varkappa\left[-\beta\frac{Z'(\beta)}{Z(\beta)}+\ln Z(\beta)\right],
\end{align*}
% Source PDF page 217 / original page 210
\hypertarget{origpage-210}{}%
and the total energy is
\[
N\mathbf E=-N\frac{Z'(\beta)}{Z(\beta)}.
\]
It is this total energy, not $\mathbf E$ itself, that must be compared with $\mathbf S$. We have thus expressed $U,\mathbf S,N\mathbf E$ in terms of $\beta$. Instead of merely expressing $\beta$ in terms of $\mathbf E$, it is more practical to determine the temperature $\mathbf T$ of the equilibrium mixture and express everything in terms of it. Proceed as follows. Bring our equilibrium mixture into contact with a heat reservoir at temperature $\mathbf T'$, and let it take from the reservoir the energy $N\,d\mathbf E$. By the two laws of thermodynamics, total energy must remain unchanged and entropy must not decrease. The heat reservoir therefore loses energy $N\,d\mathbf E$, so its entropy change is $-N\,d\mathbf E/\mathbf T'$, and at the same time we must have
\[
d\mathbf S-\frac{N\,d\mathbf E}{\mathbf T'}
=\left(\frac{d\mathbf S}{N\,d\mathbf E}-\frac1{\mathbf T'}\right)N\,d\mathbf E\geq0.
\]
On the other hand, $N\,d\mathbf E\gtreqless0$ according as $\mathbf T'\gtreqless\mathbf T$, since the colder body takes energy from the warmer. Thus $\mathbf T'\gtreqless\mathbf T$ implies $d\mathbf S/(N\,d\mathbf E)-1/\mathbf T'\gtreqless0$, or
\[
\mathbf T'\gtreqless\frac{N\,d\mathbf E}{d\mathbf S}
=\frac{N\,d\mathbf E/d\beta}{d\mathbf S/d\beta}.
\]
Consequently
\begin{align*}
\mathbf T&=\frac{N\,d\mathbf E/d\beta}{d\mathbf S/d\beta}
=-\frac1\varkappa\frac{(Z'(\beta)/Z(\beta))'}{(\ln Z(\beta)-\beta Z'(\beta)/Z(\beta))'}\\
&=-\frac1\varkappa\frac{(Z'(\beta)/Z(\beta))'}{-\beta(Z'(\beta)/Z(\beta))'}
=\frac1{\varkappa\beta},
\end{align*}
that is,
\[
\beta=\frac1{\varkappa\mathbf T}.
\]
All of $U,\mathbf S,N\mathbf E$ are now represented as functions of temperature.

The analogy between the expressions obtained above for entropy, the equilibrium ensemble, and so on, and the corresponding results of thermodynamic theory based on classical mechanics is striking. First consider the entropy $-N\varkappa\operatorname{Spur}(U\ln U)$. The operator $U=\sum_{n=1}^{\infty}w_nP_{[\varphi_n]}$ is a mixture of ensembles $P_{[\varphi_1]},P_{[\varphi_2]},\ldots$ in the ratio $w_1:w_2:\cdots$, that is, $Nw_1$ systems in state $\varphi_1$, $Nw_2$ in state $\varphi_2$, and so on. The Boltzmann entropy of this ensemble is obtained from the ``number of complexions''
\[
\frac{N!}{(Nw_1)!(Nw_2)!\cdots}:
\]
it is $\varkappa$ times its logarithm.\textsuperscript{\hyperlink{note-201}{201}} Since $N$ is large, we may approximate the factorial by Stirling's formula $x!\approx\sqrt{2\pi x}\,e^{-x}x^x$. Then $\varkappa\ln[N!/((Nw_1)!(Nw_2)!\cdots)]$ becomes, essentially, $-N\varkappa\sum_{n=1}^{\infty}w_n\ln w_n$, and this is precisely $-\operatorname{Spur}(U\ln U)$.\tnote{The final expression in the source omits the factor $N\varkappa$ present in the preceding entropy formula.}
% Source PDF page 218 / original page 211
\hypertarget{origpage-211}{}%
Furthermore, we obtained the equilibrium ensemble $U=e^{-\mathsf{H}/(\varkappa\mathbf T)}$, omitting the normalization factor $1/Z(\beta)$. This equals $\sum_{n=1}^{\infty}e^{-W_n/(\varkappa\mathbf T)}P_{[\varphi_n]}$: a mixture of the states $P_{[\varphi_1]},P_{[\varphi_2]},\ldots$, the stationary states of energies $W_1,W_2,\ldots$, with the respective relative weights $e^{-W_1/(\varkappa\mathbf T)}:e^{-W_2/(\varkappa\mathbf T)}:\cdots$. If an energy eigenvalue is multiple, say $W_{n_1}=\cdots=W_{n_\nu}=W$, then $P_{[\varphi_{n_1}]}+\cdots+P_{[\varphi_{n_\nu}]}$ occurs in the mixture with weight $e^{-W/(\varkappa\mathbf T)}$, or the properly normalized ensemble $\frac1\nu(P_{[\varphi_{n_1}]}+\cdots+P_{[\varphi_{n_\nu}]})$, see the beginning of~\hyperref[sec:IV.3]{IV.3}, with weight $\nu e^{-W/(\varkappa\mathbf T)}$. This is exactly how the classical ``canonical'' ensemble is defined, apart from the occurrence of the specifically quantum-mechanical construction $\frac1\nu(P_{[\varphi_{n_1}]}+\cdots+P_{[\varphi_{n_\nu}]})$. This is the so-called Boltzmann theorem.\textsuperscript{\hyperlink{note-201}{201}}

As $\mathbf T\to0$, the weights $e^{-W_n/(\varkappa\mathbf T)}$ tend to $1$, and our $U$ therefore tends to $\sum_{n=1}^{\infty}P_{[\varphi_n]}=1$.\tnote{The source prints $\mathbf T\to0$. For fixed $W_n$, the displayed weights tend to $1$ when $\mathbf T\to\infty$. The printed limit is retained.} Thus $U=1$ is the absolute equilibrium state if no energy restrictions apply, a result already obtained in~\hyperref[sec:IV.3]{IV.3}. We see that the ``a priori equal probability of quantum orbits,'' meaning simple, nondegenerate ones, emerges of itself from this theory. In general the multiplicity of the eigenvalue is the ``a priori weight''; see above.

It is worth determining what can be established about the equilibrium ensemble $U$ of given energy without using thermodynamics: solely from the facts that $U$ is stationary, unchanged in the course of time, process~\hyperref[eq:V.1.1]{1}, and that it remains unchanged under every measurement that does not affect energy. These are measurements of quantities simultaneously measurable with energy, process~\hyperref[eq:V.1.2]{2} with commuting $R,\mathsf{H}$, that is, with $\varphi_1,\varphi_2,\ldots$ eigenfunctions of $\mathsf{H}$.\tnote{In this paragraph the source interchanges the process numbers: temporal evolution is process~\hyperref[eq:V.1.2]{2}, while measurement is process~\hyperref[eq:V.1.1]{1}.}

By the differential equation $\partial U/\partial t=(2\pi i/h)(U\mathsf{H}-\mathsf{H}U)$, the first condition says merely that $\mathsf{H},U$ commute. The second says that whenever $\varphi_1,\varphi_2,\ldots$ can serve as a complete eigenfunction system of $\mathsf{H}$, we have $U=U'$: these $\varphi_1,\varphi_2,\ldots$ are also eigenfunctions of $U$. Let their corresponding $\mathsf{H}$ eigenvalues be $W_1,W_2,\ldots$, and those of $U$ be $w_1,w_2,\ldots$. If $W_j=W_k$, we can replace $\varphi_j,\varphi_k$ for $\mathsf{H}$ by $(\varphi_j+\varphi_k)/\sqrt2$, $(\varphi_j-\varphi_k)/\sqrt2$; these must therefore also be eigenfunctions of $U$, whence $w_j=w_k$. A function $F(x)$ with $F(W_n)=w_n$
% Source PDF page 219 / original page 212
\hypertarget{origpage-212}{}%
for $n=1,2,\ldots$ can therefore be constructed, and $F(\mathsf{H})=U$. That this is sufficient is clear, as is the fact that it implies the commutativity of $\mathsf{H},U$.

This method therefore yields $U=F(\mathsf{H})$, but does not determine $F(x)$, which, as we know, is $F(x)=e^{-\beta x}/Z(\beta)$, $\beta=1/(\varkappa\mathbf T)$. From $\operatorname{Spur}U=1$, $\operatorname{Spur}(U\mathsf{H})=\mathbf E$, we further obtain
\[
\sum_{n=1}^{\infty}F(W_n)=1,\qquad
\sum_{n=1}^{\infty}W_nF(W_n)=\mathbf E,
\]
but this exhausts the method.

\section{Macroscopic Measurement}\label{sec:V.4}
Although, as we saw, our entropy expression is entirely analogous to classical entropy, it must nevertheless seem strange that it is invariant during the system's temporal development, process~\hyperref[eq:V.1.2]{2}, and increases only during measurements, process~\hyperref[eq:V.1.1]{1}. In classical theory, where measurements played no role at all, entropy generally increased even during the system's ordinary mechanical evolution in time. This apparently paradoxical behavior therefore needs clarification.

The usual classical thermodynamic argument runs as follows. Take a vessel of volume $\mathcal V$, with $M$ molecules of a gas, ideal for simplicity, at temperature $\mathbf T$ in its right half, volume $\mathcal V/2$, separated from the other half by a partition. If we expanded this gas isothermally and reversibly to volume $\mathcal V$, letting the gas pressure push back the partition, making use of the mechanical work obtained, and keeping the gas temperature constant with a large heat reservoir at temperature $\mathbf T$, entropy outside, in the reservoir, would decrease by $M\varkappa\ln2$; see note~\hyperlink{note-195}{195}. The gas entropy would therefore increase by the same amount. If, instead, we simply withdraw the partition, the gas diffuses into the empty left half, and its volume grows to $\mathcal V$. Its entropy thus increases by $M\varkappa\ln2$ without any compensation. The process is therefore irreversible: entropy has increased in the system's simple mechanical evolution in time, namely during diffusion. Why does our theory yield nothing similar?

The situation is clearest if we put $M=1$. Thermodynamics still applies to such a one-molecule gas, and its entropy does indeed increase by $\varkappa\ln2$ when its volume is doubled. But this difference $\varkappa\ln2$ is actually present only as long as we really know nothing more about the molecule than that it is
% Source PDF page 220 / original page 213
\hypertarget{origpage-213}{}%
in volume $\mathcal V/2$ or $\mathcal V$, respectively. If, for example, the molecule is in volume $\mathcal V$, but we know whether it is to the right or left of the vessel's midpoint, we need only insert a partition at the middle and let the molecule push it, isothermally and reversibly, to the left or right end of the vessel. The mechanical work $\varkappa\mathbf T\ln2$ is performed, that energy being taken from the heat reservoir. At the end the molecule is again in volume $\mathcal V$, but we no longer know whether it is to the right or left of the middle; in compensation, there is an entropy decrease of $\varkappa\ln2$ in the reservoir. We have exchanged our knowledge for the entropy decrease $\varkappa\ln2$.\origNote{202} Thus the entropy in volume $\mathcal V$ is the same as in volume $\mathcal V/2$ if we know which half of the vessel contains the molecule. If we knew the molecule's position and momentum exactly before diffusion, we could calculate for every later instant whether it is in the right or left half; entropy would not have increased at all. Only when we have merely the macroscopic information that the previous volume was $\mathcal V/2$ does entropy actually increase during diffusion.

For a classical observer who knows all coordinates and momenta, entropy is therefore constant, specifically zero, since the Boltzmann ``number of complexions'' is $1$; see the work cited in note~\hyperlink{note-201}{201}. This is exactly as in our theory for states $U=P_{[\varphi]}$, which likewise correspond to the observer's greatest possible knowledge of the system.

The temporal variations of entropy thus arise because the observer does not know everything, or cannot determine, measure, everything that is measurable in principle. The senses permit perception only of the so-called macroscopic quantities. This clarification of the apparent contradiction mentioned at the outset obliges us, however, to seek the exact analogue of classical macroscopic entropy for quantum-mechanical ensembles: entropy as seen by an observer who cannot measure all quantities but only certain distinguished ones, the macroscopic quantities, and perhaps even these only with limited precision.

In~\hyperref[sec:III.3]{III.3} we found that all measurements of limited precision can be replaced by absolutely precise measurements of other quantities that are functions of the former and have pure point spectra. If $\mathfrak R$ is such a quantity, $R$ its operator, and $\lambda^{(1)},\lambda^{(2)},\ldots$ the distinct values among its eigenvalues, measuring $\mathfrak R$ is equivalent to answering the questions ``Is
% Source PDF page 221 / original page 214
\hypertarget{origpage-214}{}%
$\mathfrak R=\lambda^{(1)}$?'', ``Is $\mathfrak R=\lambda^{(2)}$?'', and so on. Indeed, we can say directly that if $\mathfrak S$, with operator $S$, is to be measured with limited precision, say by deciding in which interval $c_{n-1}<\lambda\leq c_n$ it lies, where
\[
\cdots<c_{-2}<c_{-1}<c_0<c_1<c_2<\cdots,
\qquad c_n\to\pm\infty\ \text{as}\ n\to\pm\infty,
\]
then the task is to answer all the questions ``Does $\mathfrak S$ lie in $c_{n-1}<\lambda\leq c_n$?'', $n=0,\pm1,\pm2,\ldots$.

By~\hyperref[sec:III.5]{III.5}, such questions correspond to projection operators $E$, whose quantities $\mathfrak E$, taking only the values $0,1$, are what must actually be measured. In our examples, the $\mathfrak E$ are the functions $F_n(\mathfrak R)$, $n=1,2,\ldots$, where
\[
F(\lambda)=\begin{cases}1,&\lambda=\lambda^{(n)},\\0,&\text{otherwise},\end{cases}
\]
or the functions $G_n(\mathfrak S)$, $n=0,\pm1,\pm2,\ldots$, where
\[
G(\lambda)=\begin{cases}1,&c_{n-1}<\lambda\leq c_n,\\0,&\text{otherwise}.\end{cases}
\]
The corresponding $E$ are $F_n(R)$ or $G_n(S)$. Instead of specifying the macroscopically measurable quantities $\mathfrak S$, including the attainable measurement precision, we can therefore specify the questions $\mathfrak E$ answerable by macroscopic measurements, or their projection operators $E$; see~\hyperref[sec:III.5]{III.5}. This is precisely the characterization of a macroscopic observer: specification of the observer's $E$. Classically one might characterize such an observer as able to measure temperature and pressure in each $\mathrm{cm}^3$ of gas volume, perhaps within certain precision limits, but nothing else.\origNote{203}

Now it is of the essence of macroscopic measurement that everything measurable in this way is also simultaneously measurable. All macroscopically answerable questions $\mathfrak E$ are simultaneously answerable; all the $E$ commute with one another. The fact that some quantum-mechanical quantities cannot be measured simultaneously initially seemed so paradoxical precisely because this concept is foreign to the macroscopic viewpoint. In view of the fundamental importance of this point, a somewhat closer discussion is appropriate.

Consider the method by which two quantities known not to be simultaneously measurable, for example the coordinate $q$ and momentum $p$, see~\hyperref[sec:III.4]{III.4}, are measured at once with limited precision. Let the respective mean measurement errors be $\varepsilon,\eta$, with $\varepsilon\eta\sim h$ by the uncertainty relation. The discussion in~\hyperref[sec:III.4]{III.4} showed that simultaneous measurement is indeed possible with such precision requirements: the $q$, or position, measurement uses light whose wavelength is not too short, and the $p$, or momentum, measurement light wave trains that are not too long. With everything arranged in this way, the actual measurement consists in detecting the two light quanta somehow, for example photographically:
% Source PDF page 222 / original page 215
\hypertarget{origpage-215}{}%
one is the light quantum deflected by the Compton effect in the $q$ measurement; the other is the quantum reflected with a Doppler change of frequency in the $p$ measurement and then deflected by an optical device, a prism or diffraction grating, to determine that frequency. At the end of the experiment there are therefore two light quanta, or two photographic plates, and we must calculate $q,p$ from the directions of the light quanta or the positions of the blackened spots on the plates. It must be emphasized that nothing prevents us from determining these two directions or spots with arbitrary precision. They are evidently simultaneously measurable quantities: momenta or coordinates of two different objects. This does not, however, help much in measuring $q,p$. As shown in~\hyperref[sec:III.4]{III.4}, the coupling of these quantities to $q,p$ is such that precisely the uncertainties $\varepsilon,\eta$ remain for the latter, even with perfectly precise knowledge of the former; one cannot arrange for $\varepsilon\eta\ll h$. Introduce the two directions or spot positions themselves as physical quantities, with operators $Q',P'$. We then see that $Q',P'$ commute exactly, but that the operators $Q,P$ belonging to $q,p$ cannot be expressed by means of them with precision better than $\varepsilon,\eta$, respectively. Let $q',p'$ be the quantities belonging to $Q',P'$. The interpretation that $q',p'$, rather than $q,p$ themselves, are the actual macroscopically measurable quantities is natural---it is indeed $q',p'$ that are measured---and agrees with our postulate of the simultaneous measurability of all macroscopic quantities.

It is useful to interpret this finding generally as a characteristic of the macroscopic viewpoint. On this view, it consists in replacing all possible operators $A,B,C,\ldots$, which generally do not commute, by other operators $A',B',C',\ldots$ of which they are approximately functions, and which do commute. Since we can call these functions of $A',B',C',\ldots$ themselves $A',B',C',\ldots$, we may also say: $A',B',C',\ldots$ approximate $A,B,C,\ldots$, but commute exactly with one another. If $\varepsilon_A,\varepsilon_B,\varepsilon_C,\ldots$ measure the respective magnitudes of $A'-A,B'-B,C'-C,\ldots$, we see that $\varepsilon_A\varepsilon_B$ will be of the order of magnitude of $AB-BA$, generally nonzero, and so on. This gives the limits of attainable approximation. Naturally, it is reasonable, already in setting up $A,B,C,\ldots$, to restrict ourselves to operators whose physical quantities can be grasped macroscopically with at least tolerable approximation.

These wholly qualitative remarks remain an empty program unless we can show that they ask for something mathematically
% Source PDF page 223 / original page 216
\hypertarget{origpage-216}{}%
feasible. For the characteristic case $Q,P$, we therefore continue the discussion of the existence of $Q',P'$ as characterized above in purely mathematical terms. Given two positive numbers $\varepsilon,\eta$ with $\varepsilon\eta=h/(4\pi)$, we seek two commuting $Q',P'$ such that $Q'-Q,P'-P$ have respective orders of magnitude $\varepsilon,\eta$, in a sense still to be defined more precisely.

We do this with quantities $q',p'$ measurable with complete precision: $Q',P'$ are to have pure point spectra. Since they commute, there is a complete orthonormal system $\varphi_1,\varphi_2,\ldots$ of common eigenfunctions; see~\hyperref[sec:II.10]{II.10}. Let the corresponding eigenvalues of $Q',P'$ be $a_1,a_2,\ldots$ and $b_1,b_2,\ldots$, respectively, so that
\[
Q'=\sum_{n=1}^{\infty}a_nP_{[\varphi_n]},\qquad
P'=\sum_{n=1}^{\infty}b_nP_{[\varphi_n]}.
\]
Their measurement is to be performed so that afterward one of the states $\varphi_1,\varphi_2,\ldots$ occurs. One may, for example, measure a quantity $\mathfrak R$ whose operator $R$ has eigenfunctions $\varphi_1,\varphi_2,\ldots$ and distinct eigenvalues $c_1,c_2,\ldots$; $Q',P'$ are then functions of $R$. For such a measurement to amount to an approximate measurement of $Q,P$, the following must evidently hold: in state $\varphi_n$, $Q,P$ are approximately expressed by the respective values of $Q',P'$, namely $a_n,b_n$; their dispersions about these values are small. These dispersions are the expectations of $(q-a_n)^2,(p-b_n)^2$, namely
\begin{align*}
((Q-a_n1)^2\varphi_n,\varphi_n)&=\|(Q-a_n1)\varphi_n\|^2=\|Q\varphi_n-a_n\varphi_n\|^2,\\
((P-b_n1)^2\varphi_n,\varphi_n)&=\|(P-b_n1)\varphi_n\|^2=\|P\varphi_n-b_n\varphi_n\|^2.
\end{align*}
They measure the squared differences between $Q',Q$ and between $P',P$, and must therefore be approximately $\varepsilon^2,\eta^2$. We thus require
\[
\|Q\varphi_n-a_n\varphi_n\|\lesssim\varepsilon,\qquad
\|P\varphi_n-b_n\varphi_n\|\lesssim\eta.
\]
Instead of speaking of $Q',P'$, it is therefore more convenient simply to seek a complete orthonormal system $\varphi_1,\varphi_2,\ldots$ for which these estimates hold with suitable $a_1,a_2,\ldots$ and $b_1,b_2,\ldots$.

From~\hyperref[sec:III.4]{III.4} we know individual $\varphi$, with $\|\varphi\|=1$, for which suitable $a,b$ give
\[
\|Q\varphi-a\varphi\|=\varepsilon,\qquad\|P\varphi-b\varphi\|=\eta:
\]
\[
\varphi_{\varrho,\sigma,\gamma}
=\varphi_{\varrho,\sigma,\gamma}(q)
\equiv\left(\frac{2\gamma}{h}\right)^{1/4}
 e^{-\frac{\pi\gamma}{h}(q-\sigma)^2+\frac{2\pi\varrho}{h}iq}.
\]
Since $\varepsilon\eta=h/(4\pi)$, we have again put $\varepsilon=\sqrt{h\gamma/(4\pi)}$, $\eta=\sqrt{h/(4\pi\gamma)}$, that is, $\gamma=\varepsilon/\eta$; we must choose $a=\sigma$, $b=\varrho$.\tnote{For the Gaussian printed above, the $Q$ and $P$ expectations are $\sigma$ and $\varrho$, with standard deviations $\sqrt{h/(4\pi\gamma)}$ and $\sqrt{h\gamma/(4\pi)}$, respectively. The source reverses these widths here and the two expectations in the following paragraph. Its notation is retained.} The task is now to assemble from these $\varphi_{\varrho,\sigma,\gamma}$ a complete orthonormal system.
% Source PDF page 224 / original page 217
\hypertarget{origpage-217}{}%
Since $\varrho$ is the $Q$ expectation and $\sigma$ the $P$ expectation, it is natural to let $\varrho$ and $\sigma$ range independently over two sets of numbers, with approximate spacings $\varepsilon$ and $\eta$, respectively. It does in fact prove convenient to choose the units $2\sqrt\pi\,\varepsilon=\sqrt{h\gamma}$ and $2\sqrt\pi\,\eta=\sqrt{h/\gamma}$, that is,
\[
\varrho=\sqrt{h\gamma}\,\mu,\qquad\sigma=\sqrt{h/\gamma}\,\nu
\qquad(\mu,\nu=0,\pm1,\pm2,\ldots).
\]
The $\psi_{\mu,\nu}=\varphi_{\sqrt{h\gamma}\,\mu,\sqrt{h/\gamma}\,\nu,\gamma}$, $\mu,\nu=0,\pm1,\pm2,\ldots$, are therefore to correspond to the $\varphi_n$, $n=1,2,\ldots$. Having two indices $\mu,\nu$ instead of one $n$ is evidently immaterial.

These $\psi_{\mu,\nu}$ are not yet orthogonal. They are normalized, and
\[
\|Q\psi_{\mu,\nu}-\sqrt{h\gamma}\,\mu\|=\varepsilon,\qquad
\|P\psi_{\mu,\nu}-\sqrt{h/\gamma}\,\nu\|=\eta.
\]
If we ``orthogonalize'' them by E.~Schmidt's procedure in any order, see~\hyperref[sec:II.2]{II.2}, proof of \hyperref[thm:II.8]{Theorem~8}, then the completeness of the resulting orthonormal system $\psi'_{\mu,\nu}$ can be proved without particular difficulty, and the estimates
\[
\|Q\psi_{\mu,\nu}-\sqrt{h\gamma}\,\mu\|\leq C\varepsilon,\qquad
\|P\psi_{\mu,\nu}-\sqrt{h/\gamma}\,\nu\|\leq C\eta
\]
can be established, with $C$ of order $60$. This value could presumably be reduced considerably. Proving these facts requires rather lengthy calculations but no new ideas, and we omit them. Factors $C\sim60$ matter little, since $\varepsilon\eta=h/(4\pi)$ is extraordinarily small in macroscopic, CGS, units, approximately $10^{-28}$.

We may therefore conclude that it is justified to assume the commutativity of all macroscopic operators, in particular of the macroscopic projection operators $E$ introduced above.

The $E$ correspond to all macroscopically answerable questions $\mathfrak E$: all distinctions between cases concerning the system under investigation that can be made macroscopically. As we saw, they all commute; by~\hyperref[sec:III.5]{III.5}, whenever $E$ is among them, so is $1-E$, and whenever $E,F$ are among them, so are $EF$, $E+F-EF$, and $E-EF$. It is natural to assume that there are only finitely many, $E_1,\ldots,E_n$. For a moment introduce the notation $E^{(+)}=E$, $E^{(-)}=1-E$, and consider all $2^n$ products $E_1^{(s_1)}\cdots E_n^{(s_n)}$, $s_1,\ldots,s_n=\pm$. Any two distinct products have product zero. For if $E_1^{(s_1)}\cdots E_n^{(s_n)}$ and $E_1^{(t_1)}\cdots E_n^{(t_n)}$ are distinct, say $s_\nu\ne t_\nu$, their product contains the factors $E_\nu^{(s_\nu)},E_\nu^{(t_\nu)}$, that is, $E_\nu^{(+)}=E$
% Source PDF page 225 / original page 218
\hypertarget{origpage-218}{}%
and $E_\nu^{(-)}=1-E$, whose product is zero. Every $E_\nu$ is the sum of some of these products, namely
\[
E_\nu=\sum_{\substack{s_1,\ldots,s_{\nu-1},s_{\nu+1},\ldots,s_n=\pm}}
 E_1^{(s_1)}\cdots E_{\nu-1}^{(s_{\nu-1})}E_\nu^{(+)}E_{\nu+1}^{(s_{\nu+1})}\cdots E_n^{(s_n)}.
\]
Call the nonzero products $E'_1,\ldots,E'_m$. Evidently $m\leq2^n$; indeed $m\leq n-1$, since they must occur among $E_1,\ldots,E_n$ and are nonzero. Then $E'_\mu\ne0$, $E'_\mu E'_\nu=0$ for $\mu\ne\nu$, and every $E_\mu$ is a sum of some $E'_\kappa$. The last fact also gives $n=2^m$. Observe that $E_\mu+E_\nu=E'_\varrho$ is impossible unless $E_\mu=0,E_\nu=E'_\varrho$ or $E_\mu=E'_\varrho,E_\nu=0$. Otherwise $E_\mu,E_\nu$ would each be sums of some $E'_\kappa$, so $E'_\varrho$ would be a sum of at least two $E'_\kappa$, possibly with repetitions. By~\hyperref[sec:II.4]{II.4}, Theorems~\hyperref[thm:II.15]{15}, \hyperref[thm:II.16]{16}, these would all be distinct and, since there are at least two, distinct from $E'_\varrho$ as well. Their products with $E'_\varrho$ would therefore vanish, as would that of their sum, contradicting the fact that this sum was to be $E'_\varrho$.

The properties $\mathfrak E'_1,\ldots,\mathfrak E'_m$ corresponding to $E'_1,\ldots,E'_m$ are therefore macroscopic properties of the following kind. None is absurd. Any two exclude one another. Every macroscopic property is obtained by disjunction from some of them. None can be decomposed by disjunction into two sharper macroscopic properties. Thus $\mathfrak E'_1,\ldots,\mathfrak E'_m$ are the most detailed macroscopic distinctions that can be made at all; they are macroscopically indecomposable.

In what follows we shall not require their number to be finite, but only the existence of the macroscopically indecomposable properties $\mathfrak E'_1,\mathfrak E'_2,\ldots$. Let their projection operators be $E'_1,E'_2,\ldots$, again all nonzero, pairwise orthogonal, and with every macroscopic $E$ a sum of some of them.

Thus $1$ too is a sum of some of them. If an $E'_\nu$ were omitted, it would be orthogonal to all those included and hence to $1$, so $E'_\nu=E'_\nu\cdot1=0$, which is impossible. Consequently $E'_1+E'_2+\cdots=1$. We drop the primes, writing $\mathfrak E_1,\mathfrak E_2,\ldots$ and $E_1,E_2,\ldots$. Let the corresponding closed linear manifolds be $\mathfrak M_1,\mathfrak M_2,\ldots$, with dimensions $s_1,s_2,\ldots$.

If the $s_n$ were $1$, so $\mathfrak M_n$ were one-dimensional, we would have $\mathfrak M_n=[\varphi_n]$, $E_n=P_{[\varphi_n]}$, and, since $E_1+E_2+\cdots=1$, the $\varphi_1,\varphi_2,\ldots$ would form a complete orthonormal system. This would mean that macroscopic measurements alone permit complete determination of the state of the system under consideration. Since this is not normally so, in general $s_n>1$, indeed $s_n\gg1$.

Observe, moreover, that the $E_n$, the elementary building blocks of the macroscopic description of the world, in a certain
% Source PDF page 226 / original page 219
\hypertarget{origpage-219}{}%
sense correspond to the division of phase space into cells customary in classical theory. We have already seen that they can approximately reproduce the behavior of noncommuting operators, particularly that of $Q,P$, so important for phase space.

What entropy does a mixture $U$ have for a macroscopic observer whose indecomposable projection operators are $E_1,E_2,\ldots$? More precisely, what is the maximum entropy such an observer can gain in transforming $U$ into $V$? That is, what entropy decrease in external objects, which may naturally be positive, zero, or negative, can the observer at best produce as compensation for the transition $U\to V$?

First, the observer cannot distinguish at all between two ensembles $U,U'$ for which every $E_1,E_2,\ldots$ has the same expectation, that is, $\operatorname{Spur}(UE_n)=\operatorname{Spur}(U'E_n)$ for $n=1,2,\ldots$. After some time the observer might distinguish them, since $U,U'$ change according to~\hyperref[eq:V.1.2]{2}, and
\[
\operatorname{Spur}(AUA^{-1}E_n)=\operatorname{Spur}(AU'A^{-1}E_n),\qquad
A=e^{-\frac{2\pi i}{h}t\mathsf{H}},
\]
need no longer hold.\origNote{204} But we are considering only measurements performable immediately. Under the stated conditions, we may therefore regard $U,U'$ as indistinguishable. Furthermore, the observer can use only semipermeable walls that transmit the $\varphi$ belonging to some $E_n$ and reflect unchanged those belonging to the others. As the method of~\hyperref[sec:V.2]{V.2} readily shows, this suffices to transform $U'=\sum_{n=1}^{\infty}x_nE_n$ reversibly into $V'=\sum_{n=1}^{\infty}y_nE_n$. The entropy difference is therefore still $\varkappa\operatorname{Spur}(U'\ln U')-\varkappa\operatorname{Spur}(V'\ln V')$, and the entropy of $U'$ is $-\varkappa\operatorname{Spur}(U'\ln U')$. Observe, however, that for such $U'$ with $\operatorname{Spur}U'=1$ to exist at all, the $\operatorname{Spur}E_n$, the numbers $s_n$, must be finite. We therefore assume every $s_n$ finite. The eigenvalue $x_1$ of $U'$ has multiplicity $s_1$, $x_2$ multiplicity $s_2$, and so on. Hence $-U'\ln U'$ has eigenvalue $-x_1\ln x_1$ of multiplicity $s_1$, $-x_2\ln x_2$ of multiplicity $s_2$, and so on. Thus $\operatorname{Spur}U'=1$ means $\sum_{n=1}^{\infty}s_nx_n=1$, and the entropy is $-\varkappa\sum_{n=1}^{\infty}s_nx_n\ln x_n$. Since
\[
U'E_m=\sum_{n=1}^{\infty}x_nE_nE_m=x_mE_m,\qquad
\operatorname{Spur}(U'E_m)=x_m\operatorname{Spur}E_m=s_mx_m,
\]
we have $x_m=\operatorname{Spur}(U'E_m)/s_m$, so the entropy in question is
\[
-\varkappa\sum_{n=1}^{\infty}\operatorname{Spur}(U'E_n)
\ln\frac{\operatorname{Spur}(U'E_n)}{s_n}.
\]
For arbitrary $U$, $\operatorname{Spur}U=1$, entropy must likewise equal
\[
-\varkappa\sum_{n=1}^{\infty}\operatorname{Spur}(UE_n)
\ln\frac{\operatorname{Spur}(UE_n)}{s_n}.
\]
% Source PDF page 227 / original page 220
\hypertarget{origpage-220}{}%
For if we put
\[
x_n=\frac{\operatorname{Spur}(UE_n)}{s_n},\qquad U'=\sum_{n=1}^{\infty}x_nE_n,
\]
then $\operatorname{Spur}(UE_n)=\operatorname{Spur}(U'E_n)$; since $U,U'$ are indistinguishable, they have the same entropy.

We should also mention that this entropy is always at least as great as the ordinary entropy:
\[
-\varkappa\sum_{n=1}^{\infty}\operatorname{Spur}(UE_n)
\ln\frac{\operatorname{Spur}(UE_n)}{s_n}
\geq-\varkappa\operatorname{Spur}(U\ln U),
\]
with equality only for $U=\sum_{n=1}^{\infty}x_nE_n$. By the results of~\hyperref[sec:V.3]{V.3}, this certainly holds if $U'=\sum_{n=1}^{\infty}[\operatorname{Spur}(UE_n)/s_n]E_n$ can be produced from $U$ by some, not necessarily macroscopic, applications of process~\hyperref[eq:V.1.1]{1}. For the left side is $-\varkappa\operatorname{Spur}(U'\ln U')$, and $U=\sum_{n=1}^{\infty}x_nE_n$ means the same as $U=U'$. Take an orthonormal system $\varphi_1^{(n)},\ldots,\varphi_{s_n}^{(n)}$ spanning the closed linear manifold $\mathfrak M_n$ belonging to $E_n$. Since $\sum_{n=1}^{\infty}E_n=1$, all the $\varphi_\nu^{(n)}$, $n=1,2,\ldots$, $\nu=1,\ldots,s_n$, together form a complete orthonormal system. Let $R$ be an operator having them as eigenfunctions with all distinct eigenvalues, and $\mathfrak R$ its physical quantity. Measuring $\mathfrak R$ transforms $U$, by~\hyperref[eq:V.1.1]{1}, into
\[
U''=\sum_{n=1}^{\infty}\sum_{\nu=1}^{s_n}
(U\varphi_\nu^{(n)},\varphi_\nu^{(n)})P_{[\varphi_\nu^{(n)}]}.
\]
Then put
\[
\psi_\mu^{(n)}=\frac1{\sqrt{s_n}}\sum_{\nu=1}^{s_n}
 e^{\frac{2\pi i}{s_n}\mu\nu}\varphi_\nu^{(n)}
\qquad(\mu=1,\ldots,s_n).
\]
The $\psi_1^{(n)},\ldots,\psi_{s_n}^{(n)}$ form an orthonormal system spanning the same closed linear manifold as $\varphi_1^{(n)},\ldots,\varphi_{s_n}^{(n)}$, namely $\mathfrak M_n$. Hence all the $\psi_\nu^{(n)}$, $n=1,2,\ldots$, $\nu=1,\ldots,s_n$, also form a complete orthonormal system. Form an operator $S$ with these eigenfunctions and its associated physical quantity $\mathfrak S$. We also record the formulas
\[
(P_{[\varphi_\nu^{(n)}]}\psi_\mu^{(m)},\psi_\mu^{(m)})
=\begin{cases}0,&m\ne n,\\1/s_n,&m=n,\end{cases}
\]
\[
\sum_{\nu=1}^{s_n}P_{[\varphi_\nu^{(n)}]}
=\sum_{\nu=1}^{s_n}P_{[\psi_\nu^{(n)}]}=E_n.
\]
% Source PDF page 228 / original page 221
\hypertarget{origpage-221}{}%
Measuring $\mathfrak S$ therefore transforms $U''$, by~\hyperref[eq:V.1.1]{1}, into
\begin{align*}
&\sum_{m=1}^{\infty}\sum_{\mu=1}^{s_m}(U''\psi_\mu^{(m)},\psi_\mu^{(m)})P_{[\psi_\mu^{(m)}]}\\
&\quad=\sum_{m=1}^{\infty}\sum_{\mu=1}^{s_m}
\left[\sum_{n=1}^{\infty}\sum_{\nu=1}^{s_n}(U\varphi_\nu^{(n)},\varphi_\nu^{(n)})
(P_{[\varphi_\nu^{(n)}]}\psi_\mu^{(m)},\psi_\mu^{(m)})\right]P_{[\psi_\mu^{(m)}]}\\
&\quad=\sum_{m=1}^{\infty}\sum_{\mu=1}^{s_m}
\left[\sum_{\nu=1}^{s_m}\frac{(U\varphi_\nu^{(m)},\varphi_\nu^{(m)})}{s_m}\right]P_{[\psi_\mu^{(m)}]}\\
&\quad=\sum_{m=1}^{\infty}\sum_{\mu=1}^{s_m}
\frac{\operatorname{Spur}(UE_m)}{s_m}P_{[\psi_\mu^{(m)}]}
=\sum_{m=1}^{\infty}\frac{\operatorname{Spur}(UE_m)}{s_m}E_m=U'.
\end{align*}
Thus two processes~\hyperref[eq:V.1.1]{1} suffice to transform $U$ into $U'$, and this is all we needed for the proof.

For states, $U=P_{[\varphi]}$, $\operatorname{Spur}(UE_n)=(E_n\varphi,\varphi)=\|E_n\varphi\|^2$, this entropy is
\[
-\varkappa\sum_{n=1}^{\infty}\|E_n\varphi\|^2\ln\frac{\|E_n\varphi\|^2}{s_n}.
\]
It no longer suffers from the defects of ``microscopic'' entropy. In general it is not constant in time, under process~\hyperref[eq:V.1.2]{2}, nor is it zero for every state $U=P_{[\varphi]}$. Indeed, the fact that the $\operatorname{Spur}(UE_n)$ from which our entropy is built are in general not constant in time was discussed in note~\hyperlink{note-202}{202}.\tnote{The source cites note~202; the time dependence of these expectations is discussed in \hyperlink{note-204}{note~204}.} It is easy to state when the state $U=P_{[\varphi]}$ has entropy zero. Since $0\leq\|E_n\varphi\|^2/s_n\leq1$, all the summands $\|E_n\varphi\|^2\ln(\|E_n\varphi\|^2/s_n)$ in the entropy expression are $\leq0$, so they must all be zero. Thus $\|E_n\varphi\|^2/s_n=0$ or $1$. The former means $E_n\varphi=0$; the latter, $\|E_n\varphi\|=\sqrt{s_n}$. But since
\[
\|E_n\varphi\|\leq1,\qquad s_n\geq1,
\]
this gives $s_n=1$, $\|E_n\varphi\|=\|\varphi\|$, hence $E_n\varphi=\varphi$: $s_n=1$ and $\varphi$ lies in $\mathfrak M_n$. The latter certainly cannot hold for two different $n$, but it also cannot hold for none, since then $E_n\varphi=0$ for all $n$ and $\sum_{n=1}^{\infty}E_n=1$ would imply $\varphi=0$. Thus $\varphi$ would lie in $\mathfrak M_n$ for exactly one $n$, with $s_n=1$. Since we found that in general all $s_n\gg1$, this is impossible. Our entropy is therefore always positive.

Since macroscopic entropy varies with time, the next question is whether it behaves like the entropy of phenomenological thermodynamics in the real world: does it predominantly increase? In classical mechanical theory, the so-called Boltzmann H-theorem answers this question affirmatively,
% Source PDF page 229 / original page 222
\hypertarget{origpage-222}{}%
though certain statistical assumptions, the so-called ``disorder assumptions,'' must be made.\origNote{205} In quantum mechanics, the author succeeded in proving the corresponding theorem without such assumptions.\origNote{206} A closer examination of this subject and of the closely connected ergodic theorem, also proved in the work cited in note~\hyperlink{note-206}{206}, would take us too far. We must therefore refrain from reproducing these investigations. The interested reader is referred to the paper cited.


% Consolidated source: chapter-VI.tex
% Source PDF page 229 / original page 222: chapter begins mid-page.
% The original-page anchor is supplied by the end of Chapter V.
\chapter{The Measurement Process}\label{chap:VI}
\section{Formulation of the Problem}\label{sec:VI.1}
Our discussion so far has examined the relation of quantum mechanics to the various causal and statistical descriptions of nature. In doing so, we encountered a peculiar duality in its procedure that we could not adequately explain. On the one hand, we found that a state $\varphi$, under the influence of an energy operator $\mathsf H$, changes during the time interval $0\leq\tau\leq t$ into a state $\varphi'$ as follows:
\begin{align*}
 \frac{\partial}{\partial\tau}\varphi_\tau
 &=-\frac{2\pi i}{h}\mathsf H\varphi_\tau\qquad(0\leq\tau\leq t),\\
 \varphi_0&=\varphi,\qquad\varphi_t=\varphi',
\end{align*}
that is,
\[
 \varphi'=e^{-\frac{2\pi i}{h}t\mathsf H}\varphi,
\]
and hence in a purely causal manner. A mixture $U$ changes correspondingly into
\[
 U'=e^{-\frac{2\pi i}{h}t\mathsf H}Ue^{\frac{2\pi i}{h}t\mathsf H};
\]
in accordance with the causal change of $\varphi$ into $\varphi'$, states $U=P_{[\varphi]}$ thereby pass into states $U'=P_{[\varphi']}$ (process~\hyperref[eq:V.1.2]{\textbf{2}} in \hyperref[sec:V.1]{V.1}). On the other hand, during a measurement---say, of a quantity having only simple eigenvalues and eigenfunctions $\varphi_1,\varphi_2,\ldots$---the state $\varphi$ undergoes an acausal change: any of the states $\varphi_1,\varphi_2,\ldots$ may arise, with respective probabilities $|(\varphi,\varphi_1)|^2,|(\varphi,\varphi_2)|^2,\ldots$. Thus the mixture
\[
 U'=\sum_{n=1}^{\infty}|(\varphi,\varphi_n)|^2P_{[\varphi_n]}
\]
arises. In general, the mixture $U$ passes into
\[
 U'=\sum_{n=1}^{\infty}(U\varphi_n,\varphi_n)P_{[\varphi_n]}
\]
(process~\hyperref[eq:V.1.1]{\textbf{1}} in \hyperref[sec:V.1]{V.1}). Since states here pass into mixtures, this process is not causal.

The difference between these two processes $U\longrightarrow U'$ is entirely fundamental: apart from their different behavior with respect to the principle of causality, they also differ in that the first is thermodynamically reversible whereas the second is not (cf.\ \hyperref[sec:V.3]{V.3}).
% Source PDF page 230 / original page 223.
\hypertarget{origpage-223}{}%

Let us now compare this account with what actually occurs in nature and in our observation of it. To begin with, it is entirely correct that measurement, or the associated process of subjective apperception, introduces something new that cannot be reduced to the physical environment. It takes us beyond that environment---or, more precisely, inward, into the individual's mental life. This inner life cannot itself be subjected to verification, since every attempt at verification already presupposes it (see the discussion below). Nevertheless, the scientific conception of the world rests on a fundamental requirement, the so-called principle of psychophysical parallelism: we must be able to describe the process of subjective apperception, which in fact lies outside the physical world, as if it occurred within that world. That is, we must be able to associate its parts with physical processes in the objective environment, in ordinary space. (This association repeatedly requires us to locate physical processes within the region of space occupied by our own body. They nevertheless remain part of the environment.)

Consider how this conception applies to a simple example: measuring a temperature. We can carry our calculations as far as determining the temperature around the thermometer's mercury bulb, and then say: the thermometer measures this temperature. But we can also continue. From the properties of mercury as explained by molecular kinetic theory, we can calculate its warming and expansion and the resulting length of the mercury column; we then say: the observer sees this length. We could go further still, taking his light source into account and determining how light quanta are reflected by the opaque mercury column, how the remaining quanta enter his eye, how the lens refracts them, and how an image forms on the retina. Only then would we say: the observer's retina registers this image. If our physiological knowledge were more precise than it is today, we could continue by tracing the chemical reactions caused by the image in the retina, the nerve pathways, and the brain. We would say only at the end: the observer apperceives these chemical changes in his brain cells. Yet however far we carry the calculation---to the mercury bulb, the thermometer's scale, the retina, or the brain---at some point we must say: the observer perceives this. We must therefore always divide the world into two parts: the observed system and the observer.
% Source PDF page 231 / original page 224.
\hypertarget{origpage-224}{}%
Within the observed system, we can trace all physical processes with arbitrary precision, at least in principle; within the observer, such a description has no meaning. The boundary between the two is largely arbitrary: in the example above we saw four different possibilities for it. In particular, the observer in this sense need not coincide with the actual observer's body. In one version of the example above, it included even the thermometer; in another, it did not even include his eyes and nerve pathways. The principle of psychophysical parallelism asserts that this boundary can be moved arbitrarily far inside the actual observer's body. But this does not alter the fact that, in every mode of description, the boundary must be drawn somewhere if the description is not to be empty, that is, if comparison with experience is to be possible. Experience supports statements only of the form ``an observer has had a certain subjective perception,'' never of the form ``a physical quantity has a certain value.''

Quantum mechanics now describes the events occurring in the observed part of the world, so long as they do not interact with the observing part, by means of process~\hyperref[eq:V.1.2]{\textbf{2}} (\hyperref[sec:V.1]{V.1}); as soon as such an interaction occurs, however, that is, a measurement, it prescribes the use of process~\hyperref[eq:V.1.1]{\textbf{1}}. The duality is therefore justified.\origNote{207} Nevertheless, there remains a danger of violating the principle of psychophysical parallelism until it is shown that the boundary between the observed system and the observer can be shifted arbitrarily in the sense indicated above.

To discuss this, we divide the world into three parts, $I,II,III$: $I$ is to be the system actually observed, $II$ the measuring instrument, and $III$ the actual observer.\origNote{208} We must show that the boundary can be drawn both between $I$ and $II+III$ and between $I+II$ and $III$. (In our example above, when comparing the first and second cases, $I$ was the system to be investigated, $II$ the thermometer, and $III$ the light plus the observer; when comparing the second and third cases, $I$ was the system to be investigated plus the thermometer, $II$ the light plus the observer's eye, and $III$ the observer from the retina onward; when comparing the third and fourth cases, $I$ was everything up to the observer's retina, $II$ his retina, nerve pathways, and brain, and $III$ his abstract ``ego.'') Thus, in the one case process~\hyperref[eq:V.1.2]{\textbf{2}} is to be applied to $I$, and process~\hyperref[eq:V.1.1]{\textbf{1}} to the interactions between $I$ and $II+III$; in the other, process~\hyperref[eq:V.1.2]{\textbf{2}} is to be applied to $I+II$, and process~\hyperref[eq:V.1.1]{\textbf{1}} to the interactions between $I+II$ and $III$. ($III$ itself remains outside the calculation in both cases.) Our task is to prove that these two procedures yield the same statements about $I$, the only part that belongs to the observed world in both descriptions.
% Source PDF page 232 / original page 225.
\hypertarget{origpage-225}{}%

We must first examine more closely how two physical systems are combined: how $I$ and $II$ give rise to $I+II$.

\section{Composite Systems}\label{sec:VI.2}
As announced at the end of the preceding section, we consider two physical systems $I,II$ (which need not represent the same systems as $I,II$ in that discussion) and their combination $I+II$. In the classical mechanical description, let $I$ have $k$ degrees of freedom, hence coordinates $q_1,\ldots,q_k$, for which we shall always substitute the single symbol $q$ as an abbreviation. Correspondingly, let $II$ have $l$ degrees of freedom and coordinates $r_1,\ldots,r_l$, to be abbreviated as $r$. Then $I+II$ has $k+l$ degrees of freedom and coordinates $q_1,\ldots,q_k,r_1,\ldots,r_l$, or, briefly, $q,r$. In quantum mechanics, therefore, the wave functions of $I$ have the form $\varphi(q)$, those of $II$ the form $\xi(r)$, and those of $I+II$ the form $\Phi(q,r)$. In the corresponding Hilbert spaces $\mathfrak R^I,\mathfrak R^{II},\mathfrak R^{I+II}$ the inner product is defined, as appropriate, by
\[
 \int\varphi(q)\overline{\psi(q)}\,dq,\qquad
 \int\xi(r)\overline{\eta(r)}\,dr,\qquad
 \iint\Phi(q,r)\overline{\Psi(q,r)}\,dq\,dr,
\]
respectively. The physical quantities of $I,II,I+II$ are correspondingly represented by the hypermaximal Hermitian operators $\mathsf A,A,\boldsymbol A$ in $\mathfrak R^I,\mathfrak R^{II},\mathfrak R^{I+II}$, respectively.

Every physical quantity in $I$ is of course also one in $I+II$, and the corresponding operator $\boldsymbol A$ is obtained from $\mathsf A$ as follows: to determine $\boldsymbol A\Phi(q,r)$, regard $r$ as a constant and apply $\mathsf A$ to the function of $q$, $\Phi(q,r)$.\origNote{209} This rule of association is certainly correct for the coordinate and momentum operators ($Q_1,\ldots,Q_k$ and $P_1,\ldots,P_k$, that is, $q_1,\ldots,q_k$, $\frac{h}{2\pi i}\frac{\partial}{\partial q_1},\ldots,\frac{h}{2\pi i}\frac{\partial}{\partial q_k}$; cf.\ \hyperref[sec:I.2]{I.2}), and accords with the principles of association \hyperref[cond:IV.2.I]{\textbf{I}}, \hyperref[cond:IV.2.II]{\textbf{II}} in \hyperref[sec:IV.2]{IV.2};\origNote{210} we therefore postulate it generally. (It is in general use in quantum mechanics.)

Likewise, every physical quantity in $II$ is also one in $I+II$, and the corresponding operator $\boldsymbol A$ is obtained from $A$ by the same rule: $\boldsymbol A\Phi(q,r)$ equals $A\Phi(q,r)$ when, in the latter expression, $q$ is regarded as constant and $\Phi(q,r)$ as a function of $r$.

If $\varphi_m(q)$ ($m=1,2,\ldots$) is a complete orthonormal system in $\mathfrak R^I$, and $\xi_n(r)$ ($n=1,2,\ldots$) one in $\mathfrak R^{II}$, then $\Phi_{mn}(q,r)=\varphi_m(q)\xi_n(r)$ ($m,n=1,2,\ldots$) is evidently one in $\mathfrak R^{I+II}$. The operators $\mathsf A,A,\boldsymbol A$ can therefore be represented by matrices $\{a_{m/m'}\}$, $\{a_{n/n'}\}$, and $\{\alpha_{mn/m'n'}\}$, respectively ($m,m',n,n'=1,2,\ldots$).\origNote{211} We shall make repeated use of this. Matrix representation means
\begin{align*}
 \mathsf A\varphi_m(q)&=\sum_{m'=1}^{\infty}a_{m/m'}\varphi_{m'}(q),\\
 A\xi_n(r)&=\sum_{n'=1}^{\infty}a_{n/n'}\xi_{n'}(r),
\end{align*}
% Source PDF page 233 / original page 226.
\hypertarget{origpage-226}{}%
and
\[
 \boldsymbol A\Phi_{mn}(q,r)
 =\sum_{m',n'=1}^{\infty}\alpha_{mn/m'n'}\Phi_{m'n'}(q,r),
\]
that is,
\[
 \boldsymbol A\varphi_m(q)\xi_n(r)
 =\sum_{m',n'=1}^{\infty}\alpha_{mn/m'n'}\varphi_{m'}(q)\xi_{n'}(r).
\]
In particular, the association $\mathsf A\longrightarrow\boldsymbol A$ means
\[
 \boldsymbol A\varphi_m(q)\xi_n(r)
 =(\mathsf A\varphi_m(q))\xi_n(r)
 =\sum_{m'=1}^{\infty}a_{m/m'}\varphi_{m'}(q)\xi_n(r),
\]
that is,
\[
 \alpha_{mn/m'n'}=a_{m/m'}\delta_{n/n'},\qquad
 \delta_{n/n'}=\begin{cases}1,&n=n',\\0,&n\ne n'.\end{cases}
\]
Analogously, the association $A\longrightarrow\boldsymbol A$ means $\alpha_{mn/m'n'}=a_{n/n'}\delta_{m/m'}$.

A statistical ensemble in $I+II$ is characterized by its statistical operator $\boldsymbol U$, or by its matrix $\{v_{mn/m'n'}\}$. This determines the statistical properties of all quantities in $I+II$, and hence, in particular, those of the quantities in $I$. Thus there corresponds to it a statistical ensemble in $I$ alone: indeed, an observer capable of perceiving only $I$ and not $II$ would regard the ensemble of systems $I+II$ as an ensemble of systems $I$. What, then, is the statistical operator $\mathsf U$, or its matrix $\{u_{m/m'}\}$, belonging to this $I$-ensemble? We determine it as follows. The $I$-quantity with matrix $\{a_{m/m'}\}$ has, as an $I+II$-quantity, the matrix $\{a_{m/m'}\delta_{n/n'}\}$; thus a calculation in $I$ gives it the expectation value $\sum_{m,m'=1}^{\infty}u_{m/m'}a_{m'/m}$, whereas the calculation in $I+II$ gives
\begin{align*}
 \sum_{m,n,m',n'=1}^{\infty}v_{mn/m'n'}a_{m'/m}\delta_{n'/n}
 &=\sum_{m,m',n=1}^{\infty}v_{mn/m'n}a_{m'/m}\\
 &=\sum_{m,m'=1}^{\infty}\left(\sum_{n=1}^{\infty}v_{mn/m'n}\right)a_{m'/m}.
\end{align*}
For these expressions to be equal, we must have
\[
 u_{m/m'}=\sum_{n=1}^{\infty}v_{mn/m'n}.
\]

In exactly the same way, if only $II$ is considered and $I$ ignored, our $I+II$-ensemble determines a $II$-ensemble, with a statistical operator $U$ and its matrix $\{u_{n/n'}\}$. Analogously one obtains
\[
 u_{n/n'}=\sum_{m=1}^{\infty}v_{mn/mn'}.
\]

Thus we have also obtained rules of association for the statistical operators of $I,II,I+II$, namely $\mathsf U,U,\boldsymbol U$; these, however, differ essentially from the rules of association applying to the operators $\mathsf A,A,\boldsymbol A$ of physical quantities.
% Source PDF page 234 / original page 227.
\hypertarget{origpage-227}{}%

It should also be mentioned that our association of $\mathsf U,U,\boldsymbol U$ only appears to depend on the choice of the complete orthonormal systems $\varphi_m(q)$ and $\xi_n(q)$. It was derived from an invariant condition that uniquely determines the association: agreement of the expectation values of $\mathsf A$ and $\boldsymbol A$, and of $A$ and $\boldsymbol A$, respectively.

$\boldsymbol U$ expresses the statistics in $I+II$, while $\mathsf U$ and $U$ express the statistics restricted to $I$ and $II$, respectively. The question arises: do $\mathsf U,U$ determine $\boldsymbol U$ uniquely, or not? In general, we expect the answer to be no: knowing only $\mathsf U$ and $U$, and thus the separate systems $I$ and $II$, tells us nothing about the ``statistical dependences'' that may exist between them. But if both the state of $I$ and that of $II$ are known exactly, ``statistical dependences'' are out of the question, and $I+II$ is also known exactly. Let us now replace these qualitative considerations with a precise mathematical analysis.

The problem is to find, for two given definite matrices $\{u_{m/m'}\}$ and $\{u_{n/n'}\}$, a third definite matrix $\{v_{mn/m'n'}\}$ such that
\[
 \sum_{n=1}^{\infty}v_{mn/m'n}=u_{m/m'},\qquad
 \sum_{m=1}^{\infty}v_{mn/mn'}=u_{n/n'}.
\]
(From $\sum_{m=1}^{\infty}u_{m/m}=1$, $\sum_{n=1}^{\infty}u_{n/n}=1$, it then follows automatically that $\sum_{m,n=1}^{\infty}v_{mn/mn}=1$, so that the proper normalization is preserved.) The problem is always soluble; for example, $v_{mn/m'n'}=u_{m/m'}u_{n/n'}$ is always a solution (it is easy to see that this matrix is definite). The question, however, is when it is the only one.

We show that this is the case if and only if at least one of the two matrices $\{u_{m/m'}\}$, $\{u_{n/n'}\}$ represents a state. We first prove necessity, that is, the existence of several solutions when both matrices correspond to mixtures. In that case (cf.\ \hyperref[sec:IV.2]{IV.2})
\[
 u_{m/m'}=\alpha v_{m/m'}+\beta w_{m/m'},\qquad
 u_{n/n'}=\gamma v_{n/n'}+\delta w_{n/n'},
\]
where $v_{m/m'},w_{m/m'}$ are definite and do not differ merely by a constant factor, and likewise $v_{n/n'},w_{n/n'}$;
\[
 \sum_{m=1}^{\infty}v_{m/m}=\sum_{m=1}^{\infty}w_{m/m}
 =\sum_{n=1}^{\infty}v_{n/n}=\sum_{n=1}^{\infty}w_{n/n}=1,
\]
and $\alpha,\beta,\gamma,\delta>0$, $\alpha+\beta=1$, $\gamma+\delta=1$. It is easily verified that every\tnote{The source omits the prime on the last $n$ in the left-hand side of the following expression. Its subsequent uses, and the right-hand side, require $v_{mn/m'n'}$. The printed subscript is retained.}
\begin{multline*}
 v_{mn/m'n}=\pi v_{m/m'}v_{n/n'}+\varrho w_{m/m'}v_{n/n'}\\
 {}+\sigma v_{m/m'}w_{n/n'}+\tau w_{m/m'}w_{n/n'}
\end{multline*}
% Source PDF page 235 / original page 228.
\hypertarget{origpage-228}{}%
with
\begin{gather*}
 \pi+\sigma=\alpha,\quad\varrho+\tau=\beta,\quad
 \pi+\varrho=\gamma,\quad\sigma+\tau=\delta,\\
 \pi,\varrho,\sigma,\tau>0,
\end{gather*}
is a solution. The $\pi,\varrho,\sigma,\tau$ can be chosen in infinitely many ways (because $\alpha+\beta=\gamma+\delta$, only three of the four equations are independent; for instance, $\varrho=\gamma-\pi$, $\sigma=\alpha-\pi$, $\tau=(\delta-\alpha)+\pi$. For everything to be positive, $\alpha-\delta=\gamma-\beta<\pi<\alpha,\gamma$ is required, and this holds for infinitely many $\pi$). Different $\pi,\varrho,\sigma,\tau$ lead to different $v_{mn/m'n'}$, since the $v_{m/m'}v_{n/n'},\ldots,w_{m/m'}w_{n/n'}$ are linearly independent, because $v_{m/m'},w_{m/m'}$ and also $v_{n/n'},w_{n/n'}$ are linearly independent.

We now prove sufficiency, and may assume that $u_{m/m'}$ corresponds to a state (the other case is dealt with in the same way). Thus $\mathsf U=P_{[\varphi]}$, and since the complete orthonormal system $\varphi_1,\varphi_2,\ldots$ was arbitrary, we may assume $\varphi_1=\varphi$. Clearly $\mathsf U=P_{[\varphi_1]}$ has the matrix
\[
 u_{m/m'}=\begin{cases}1,&m=m'=1,\\0,&\text{otherwise}.\end{cases}
\]
Hence
\[
 \sum_{n=1}^{\infty}v_{mn/m'n}
 =\begin{cases}1,&m=m'=1,\\0,&\text{otherwise}.\end{cases}
\]
In particular, $\sum_{n=1}^{\infty}v_{mn/mn}=0$ for $m\ne1$; but since definiteness of $v_{mn/m'n'}$ implies that all $v_{mn/mn}\geq0$ [we have $v_{mn/mn}=(\boldsymbol U\Phi_{mn},\Phi_{mn})$], it follows in this case that $v_{mn/mn}=0$. That is, $(\boldsymbol U\Phi_{mn},\Phi_{mn})=0$, and so, by definiteness of $\boldsymbol U$, also $(\boldsymbol U\Phi_{mn},\Phi_{m'n'})=0$ (cf.\ \hyperref[sec:II.5]{II.5}, \hyperref[thm:II.19]{Theorem~19}), where $m',n'$ are arbitrary. Thus $m\ne1$ implies $v_{mn/m'n'}=0$, and by the Hermitian property this follows also from $m'\ne1$. For $m=m'=1$, however, we obtain
\[
 v_{1n/1n'}=\sum_{m=1}^{\infty}v_{mn/mn'}=u_{n/n'}.
\]
Thus the solution $v_{mn/m'n'}$ is uniquely determined, as asserted.

To summarize, we can also formulate our result as follows. A statistical ensemble in $I+II$, with operator $\boldsymbol U=\{v_{mn/m'n'}\}$, is uniquely determined by the statistical ensembles which it determines in $I$ alone and $II$ alone, with operators $\mathsf U=\{u_{m/m'}\}$ and $U=\{u_{n/n'}\}$ respectively, if and only if it satisfies the following conditions:

\textbf{1.} We have $v_{mn/m'n'}=v_{m/m'}v_{n/n'}$. (From
\[
 \operatorname{Spur}\boldsymbol U=\sum_{m,n=1}^{\infty}v_{mn/mn}
 =\sum_{m=1}^{\infty}v_{m/m}\sum_{n=1}^{\infty}v_{n/n}=1
\]
it follows that, by multiplying $v_{m/m'}$ and $v_{n/n'}$ by two reciprocal constant factors, we can arrange $\sum_{m=1}^{\infty}v_{m/m}=1$, $\sum_{n=1}^{\infty}v_{n/n}=1$.
% Source PDF page 236 / original page 229.
\hypertarget{origpage-229}{}%
But then we see that $u_{m/m'}=v_{m/m'}$, $u_{n/n'}=v_{n/n'}$.)

\textbf{2.} Either $v_{m/m'}=\overline{x_m}x_{m'}$ or $v_{n/n'}=\overline{x_n}x_{n'}$. (For $\mathsf U=P_{[\varphi]}$ means $\varphi=\sum_{m=1}^{\infty}y_m\varphi_m$, and hence $u_{m/m'}=\overline{y_m}y_{m'}$, and correspondingly for $v_{m/m'}$; the case $U=P_{[\xi]}$ is analogous.)

We shall call $\mathsf U$ and $U$ the projections of $\boldsymbol U$ into $I$ and $II$, respectively.\origNote{212}

We now turn to the states of $I+II$, $\boldsymbol U=P_\Phi$. The corresponding wave functions $\Phi(q,r)$ can be expanded in the complete orthonormal system $\Phi_{mn}(q,r)=\varphi_m(q)\xi_n(r)$:
\[
 \Phi(q,r)=\sum_{m,n=1}^{\infty}f_{mn}\varphi_m(q)\xi_n(r).
\]

Thus we can also replace them by the coefficients $f_{mn}$ ($m,n=1,2,\ldots$), which are subject only to the restriction that $\sum_{m,n=1}^{\infty}|f_{mn}|^2=\|\Phi\|^2$ be finite.

We can define two operators $F,F^*$ by
\begin{align*}
 F\varphi(q)&=\int\overline{\Phi(qr)}\varphi(q)\,dq,\\
 F^*\xi(r)&=\int\Phi(qr)\xi(r)\,dr.
\end{align*}
These operators are linear, but each maps one Hilbert space into the other: $F$ maps $\mathfrak R^I$ into $\mathfrak R^{II}$, and $F^*$ maps $\mathfrak R^{II}$ into $\mathfrak R^I$. They are adjoints of one another, since evidently $(F\varphi,\xi)=(\varphi,F^*\xi)$ (the inner product on the left is to be formed in $\mathfrak R^{II}$, that on the right in $\mathfrak R^I$!). Since the difference between $\mathfrak R^I$ and $\mathfrak R^{II}$ is mathematically immaterial, we can apply the results of \hyperref[sec:II.11]{II.11}: as these are integral operators, $\Sigma(F)$ and $\Sigma(F^*)$ both equal
\[
 \iint|\Phi(qr)|^2\,dq\,dr=\|\Phi\|^2=1
 \qquad(\|\Phi\|\text{ in }\mathfrak R^{I+II}!),
\]
and are therefore finite. Thus $F,F^*$ are continuous, indeed completely continuous operators, and $F^*F$ and $FF^*$ are definite operators, $\operatorname{Spur}(F^*F)=\Sigma(F)=1$, $\operatorname{Spur}(FF^*)=\Sigma(F^*)=1$. Taking account again of the difference between $\mathfrak R^I,\mathfrak R^{II}$, we see that $F^*F$ is defined in $\mathfrak R^I$, and $FF^*$ in $\mathfrak R^{II}$.

Since $F\varphi_m(q)$ comes out as $\sum_{n=1}^{\infty}\overline{f_{mn}}\xi_n(r)$, $F$ has matrix $\{\overline{f_{mn}}\}$ [using the complete orthonormal systems $\varphi_m(q)$ and $\xi_n(r)$, respectively]; likewise $F^*$ has matrix $\{f_{mn}\}$. Consequently, $F^*F,FF^*$ have the matrices
\[
 \left\{\sum_{n=1}^{\infty}\overline{f_{mn}}f_{m'n}\right\},\qquad
 \left\{\sum_{m=1}^{\infty}\overline{f_{mn}}f_{mn'}\right\},
\]
respectively.
% Source PDF page 237 / original page 230.
\hypertarget{origpage-230}{}%

On the other hand, $\boldsymbol U=P_{[\Phi]}$ has, in the system $\Phi_{mn}(q,r)=\varphi_m(q)\xi_n(r)$, the matrix $\{\overline{f_{mn}}f_{m'n'}\}$, so that its projections into $I$ and $II$, $\mathsf U$ and $U$, have the matrices
\[
 \left\{\sum_{n=1}^{\infty}\overline{f_{mn}}f_{m'n}\right\},\qquad
 \left\{\sum_{m=1}^{\infty}\overline{f_{mn}}f_{mn'}\right\},
\]
respectively.\origNote{213} Hence
\[
 \mathsf U=F^*F,\qquad U=FF^*.
\]
This identity is independent of the choice of $\varphi_m,\xi_n$, because $\mathsf U,U,F$ are themselves independent of that choice.

The operators $\mathsf U,U$ are completely continuous and, by \hyperref[sec:II.11]{II.11} and \hyperref[sec:IV.3]{IV.3}, can be written as
\[
 \mathsf U=\sum_{k=1}^{\infty}w'_kP_{[\psi_k]},\qquad
 U=\sum_{k=1}^{\infty}w''_kP_{[\eta_k]},
\]
where the $\psi_k$ form a complete orthonormal system in $\mathfrak R^I$, the $\eta_k$ one in $\mathfrak R^{II}$, and all $w'_k,w''_k\geq0$. In each formula we now omit the terms for which $w'_k=0$ or $w''_k=0$, respectively, and renumber the remaining terms consecutively by $k=1,2,\ldots$. Then the $\psi_k$ and $\eta_k$ form orthonormal systems, but no longer necessarily complete ones; the sums $\sum_{k=1}^{\infty}$ are replaced by sums $\sum_{k=1}^{M'}$, $\sum_{k=1}^{M''}$, where $M',M''$ may be either finite or infinite; but all the $w'_k,w''_k$ are now positive.

Consider a $\psi_k$. We have $\mathsf U\psi_k=w'_k\psi_k$, hence $F^*F\psi_k=w'_k\psi_k$, $FF^*F\psi_k=w'_kF\psi_k$, $UF\psi_k=w'_kF\psi_k$. Furthermore,
\begin{align*}
 (F\psi_k,F\psi_l)&=(F^*F\psi_k,\psi_l)=(\mathsf U\psi_k,\psi_l)\\
 &=w'_k(\psi_k,\psi_l)
 =\begin{cases}w'_k,&k=l,\\0,&k\ne l.\end{cases}
\end{align*}
In particular, $\|F\psi_k\|^2=w'_k$. The $\frac{1}{\sqrt{w'_k}}F\psi_k$ therefore form an orthonormal system in $\mathfrak R^{II}$, and are eigenfunctions of $U$ with the same eigenvalues as the $\psi_k$ for $\mathsf U$, namely $w'_k$. Thus every eigenvalue of $\mathsf U$ is also an eigenvalue of $U$, with at least the same multiplicity. Interchanging $\mathsf U,U$ shows that they have the same eigenvalues with the same multiplicities. The $w'_k$ and $w''_k$ therefore agree up to their order. Hence $M'=M''=M$, and by renumbering the $w''_k$ we can arrange $w''_k=w'_k=w_k$. Once this is done, we can evidently choose $\eta_k=\frac{1}{\sqrt{w_k}}F\psi_k$ in general. Then
\[
 \frac{1}{\sqrt{w_k}}F^*\eta_k
 =\frac{1}{w_k}F^*F\psi_k=\frac{1}{w_k}\mathsf U\psi_k=\psi_k.
\]
Thus:\origNote{212}
\[
 \eta_k=\frac{1}{\sqrt{w_k}}F\psi_k,\qquad
 \psi_k=\frac{1}{\sqrt{w_k}}F^*\eta_k.
\]
% Source PDF page 238 / original page 231.
\hypertarget{origpage-231}{}%

Now extend the orthonormal system $\psi_1,\allowbreak \psi_2,\allowbreak \ldots$ to a complete one $\psi_1,\allowbreak \psi_2,\allowbreak \ldots,\allowbreak \psi'_1,\allowbreak \psi'_2,\allowbreak \ldots$, and likewise $\eta_1,\allowbreak \eta_2,\allowbreak \ldots$ to $\eta_1,\allowbreak \eta_2,\allowbreak \ldots,\allowbreak \eta'_1,\allowbreak \eta'_2,\allowbreak \ldots$ (each of the two systems $\psi'_1,\allowbreak \psi'_2,\allowbreak \ldots$ and $\eta'_1,\allowbreak \eta'_2,\allowbreak \ldots$ may be empty, finite, or infinite, independently of the other). The preceding considerations were independent of the choice of the complete orthonormal systems $\varphi_1,\allowbreak \varphi_2,\allowbreak \ldots$ and $\xi_1,\allowbreak \xi_2,\allowbreak \ldots$, so we may let them coincide with $\psi_1,\allowbreak \psi_2,\allowbreak \ldots,\allowbreak \psi'_1,\allowbreak \psi'_2,\allowbreak \ldots$ and $\eta_1,\allowbreak \eta_2,\allowbreak \ldots,\allowbreak \eta'_1,\allowbreak \eta'_2,\allowbreak \ldots$, respectively. Let $\psi_k$ correspond to $\varphi_{\mu_k}$ and $\eta_k$ to $\xi_{\nu_k}$ ($\mu_1,\allowbreak \mu_2,\allowbreak \ldots$ distinct, and likewise $\nu_1,\allowbreak \nu_2,\allowbreak \ldots$). Then
\[
 F\varphi_{\mu_k}=\sqrt{w_k}\xi_{\nu_k},\qquad
 F\varphi_m=0\quad\text{for }m\ne\mu_1,\mu_2,\ldots.
\]
Consequently,
\[
 f_{mn}=\begin{cases}\sqrt{w_k},&m=\mu_k,\ n=\nu_k,\ k=1,2,\ldots,\\
 0,&\text{otherwise},\end{cases}
\]
or, equivalently,
\[
 \Phi(q,r)=\sum_{k=1}^{M}\sqrt{w_k}\varphi_{\mu_k}(q)\xi_{\nu_k}(q).
\]

By a suitable choice of the systems $\varphi_m(q)$ and $\xi_n(q)$, we have thus arranged that each row and each column of the matrix $\{f_{mn}\}$ contains at most one nonzero element (the fact that this element is even real and positive, namely $\sqrt{w_k}$, is unimportant for what follows). What is the physical meaning of this formal statement?\tnote{Here the source prints $q$ as the argument of $\xi_{\nu_k}$ and $\xi_n$, although these functions belong to the second system and elsewhere have argument $r$. These printed arguments have been retained.}

Let $A$ be an operator with eigenfunctions $\varphi_1,\varphi_2,\ldots$ and mutually distinct eigenvalues, say $a_1,a_2,\ldots$; likewise let $B$ have $\xi_1,\xi_2,\ldots$ and $b_1,b_2,\ldots$. The operator $A$ corresponds to a physical quantity in $I$, and $B$ to one in $II$, so they are simultaneously measurable. It is easy to see that the statements ``$A$ has value $a_m$ and $B$ has value $b_n$'' together determine the state $\Phi_{mn}(q,r)=\varphi_m(q)\xi_n(r)$, and that, when the system is in the state $\Phi(q,r)$, this joint outcome has probability $(P_{[\Phi_{mn}]}\Phi,\Phi)=|(\Phi,\Phi_{mn})|^2=|f_{mn}|^2$.
Thus our statement means: $A,B$ are simultaneously measurable, and when one of them has been measured in $\Phi$, the value of the other is uniquely determined. (An $a_m$ for which all $f_{mn}=0$ cannot occur, since its total probability $\sum_{n=1}^{\infty}|f_{mn}|^2$ cannot be zero if $a_m$ is found at all; thus $f_{mn}\ne0$ for exactly one $n$. The same holds for $b_n$.) That is, several $A$-values may be possible in the state $\Phi$ (every $a_m$ for which $\sum_{n=1}^{\infty}|f_{mn}|^2>0$, that is, for which an $n$ with $f_{mn}\ne0$ exists; usually all the $a_m$), and just as many $B$-values (every $b_n$ for which $\sum_{n=1}^{\infty}|f_{mn}|^2>0$, that is, for which an $m$ with $f_{mn}\ne0$ exists); but $\Phi$ establishes a one-to-one correspondence between the possible $A$- and $B$-values.\tnote{In the parenthesis concerning the $B$-values, the source repeats the summation index $n$; the corresponding marginal sum is over $m$. The printed index is retained.}
% Source PDF page 239 / original page 232.
\hypertarget{origpage-232}{}%

If we denote the possible $m$ by $\mu_1,\mu_2,\ldots$ and the corresponding possible $n$ by $\nu_1,\nu_2,\ldots$, then
\[
 f_{mn}=\begin{cases}c_k\ne0,&m=\mu_k,\ n=\nu_k,\ k=1,2,\ldots,\\
 0,&\text{otherwise},\end{cases}
\]
and therefore ($M$ finite or infinite)
\[
 \Phi(q,r)=\sum_{k=1}^{M}c_k\varphi_{\mu_k}(q)\xi_{\nu_k}(r).
\]
We then have
\begin{align*}
 u_{m/m'}&=\sum_{n=1}^{\infty}\overline{f_{mn}}f_{m'n}
 =\begin{cases}|c_k|^2,&m=m'=\mu_k,\ k=1,2,\ldots,\\0,&\text{otherwise},\end{cases}\\
 u_{n/n'}&=\sum_{m=1}^{\infty}\overline{f_{mn}}f_{mn'}
 =\begin{cases}|c_k|^2,&n=n'=\nu_k,\ k=1,2,\ldots,\\0,&\text{otherwise},\end{cases}
\end{align*}
and hence
\[
 \mathsf U=\sum_{k=1}^{M}|c_k|^2P_{[\varphi_{\mu_k}]},\qquad
 U=\sum_{k=1}^{M}|c_k|^2P_{[\xi_{\nu_k}]}.
\]
Thus, when projected into $I$ or $II$, $\Phi$ generally becomes a mixture; it is a state only in $I+II$, because it contains information about $I+II$ that cannot be used in either $I$ or $II$ alone: the one-to-one correspondence between the $A$- and $B$-values.

For every $\Phi$, therefore, we can choose $A,B$, that is, the $\varphi_m$ and $\xi_n$, so that this condition is satisfied; for arbitrary $A,B$ it will of course be violated. Every state $\Phi$ thus establishes a particular connection between $I$ and $II$, the correlated quantities $A,B$ depending on $\Phi$. It is not difficult to determine how completely $\Phi$ specifies these quantities, or equivalently their eigenfunctions $\varphi_m$ and $\xi_n$. If all $|c_k|$ are different and nonzero, then $\mathsf U,U$ (which are determined by $\Phi$) determine the $\varphi_m$ and $\xi_n$ uniquely (cf.\ \hyperref[sec:IV.3]{IV.3}); the general discussion is left to the reader.

Finally, we remark that, since $|c_k|^2>0$, neither $\mathsf U$ nor $U$ is a state when $M\ne1$. For $M=1$, both are states: $\mathsf U=P_{[\varphi_{\mu_1}]}$, $U=P_{[\xi_{\nu_1}]}$; then $\Phi(q,r)=c_1\varphi_{\mu_1}(q)\xi_{\nu_1}(r)$. We can absorb $c_1$ into $\varphi_{\mu_1}(q)$. Thus $\mathsf U,U$ are states if and only if $\Phi(q,r)$ has the form $\varphi(q)\xi(r)$; they then equal $P_{[\varphi]}$ and $P_{[\xi]}$, respectively.

The preceding results also establish the following important point: if $I$ is in the state $\varphi(q)$ and $II$ in the state $\xi(r)$, then $I+II$ is in the state $\Phi(q,r)=\varphi(q)\xi(r)$. If, on the other hand, $I+II$ is in a state $\Phi(q,r)$ that is not a product $\varphi(q)\xi(r)$, then $I$ and $II$ are mixtures, but $\Phi$ establishes a one-to-one correspondence between the possible values of certain quantities in $I$ and in $II$.

% Source PDF page 240 / original page 233.
\hypertarget{origpage-233}{}%
\section{Discussion of the Measurement Process}\label{sec:VI.3}
We can now use the formal tools of \hyperref[sec:VI.2]{VI.2} to complete our discussion of measurement along the lines set out in \hyperref[sec:VI.1]{VI.1}. Before doing so, however, we shall use the results of \hyperref[sec:VI.2]{VI.2} to rule out a repeatedly proposed explanation of the statistical character of process~\hyperref[eq:V.1.1]{\textbf{1}} (\hyperref[sec:V.1]{V.1}). It rests on the following idea. Let $I$ be the observed system and $II$ the observer. If before the measurement $I$ is in a state $\mathsf U=P_{[\varphi]}$, but $II$ in a mixture $U=\sum_{n=1}^{\infty}w_nP_{[\xi_n]}$, then $I+II$ is a uniquely determined mixture $\boldsymbol U$; indeed, as one easily calculates from \hyperref[sec:VI.2]{VI.2},
\[
 \boldsymbol U=\sum_{n=1}^{\infty}w_nP_{[\Phi_n]},\qquad
 \Phi_n(q,r)=\varphi(q)\xi_n(r).
\]
If a measurement of the quantity $A$ in $I$ now takes place, it is to be understood as an interaction between $I$ and $II$. This is a process~\hyperref[eq:V.1.2]{\textbf{2}} (\hyperref[sec:V.1]{V.1}) with an energy operator $\boldsymbol H$. If it lasts a time $t$, then $\boldsymbol U$ becomes
\[
 \boldsymbol U'=e^{-\frac{2\pi i}{h}t\boldsymbol H}\boldsymbol Ue^{\frac{2\pi i}{h}t\boldsymbol H},
\]
and evidently
\[
 \boldsymbol U'=\sum_{n=1}^{\infty}w_nP_{[e^{-\frac{2\pi i}{h}t\boldsymbol H}\Phi_n]}.
\]
If every $\Phi_n(q,r)$ were of the form $\psi_n(q)\eta_n(r)$,\tnote{The source writes $\Phi_n(q,r)$ here. In the preceding expression for $\boldsymbol U'$, the states after the interaction are $e^{-\frac{2\pi i}{h}t\boldsymbol H}\Phi_n$; the proposed product form concerns those evolved states. The printed notation is retained.} where the $\psi_n$ are the eigenfunctions of $A$ and the $\eta_n$ some fixed complete orthonormal system, this interaction would constitute a measurement: it would transform every state $\varphi$ of $I$ into a mixture of the eigenfunctions $\psi_n$ of $A$. The statistical character would arise as follows. Before the measurement, $I$ was in a single state but $II$ was a mixture. During the interaction, the mixed character of $II$ would ``infect'' $I+II$, making its projection into $I$, in particular, a mixture. That is, the result of the measurement is indeterminate because the observer's state before the measurement is not known exactly. It is conceivable that such a mechanism might operate, since the observer's knowledge of his own state might have limits imposed by natural law. These limits would be reflected in the values of the $w_n$, which would have to depend solely on the observer (and therefore be independent of $\varphi$!).

It is precisely here that this attempted explanation fails: quantum mechanics requires $w_n=(P_{\psi_n}\varphi,\varphi)=|(\varphi,\psi_n)|^2$, so that $w_n$ depends on $\varphi$! Nor would the possible existence of another decomposition $\boldsymbol U'=\sum_{n=1}^{\infty}w'_nP_{[\Phi'_n]}$ help (the $\Phi'_n(q,r)=\psi_n(q)\eta_n(r)$ are normalized and orthogonal!): for the $w'_n$ are uniquely determined, apart from their order, by $\boldsymbol U'$ (\hyperref[sec:IV.3]{IV.3}), and are therefore equal to the $w_n$.\origNote{214}

The acausality of process~\hyperref[eq:V.1.1]{\textbf{1}} therefore cannot be attributed to incomplete knowledge of the observer's state; in what follows, we shall accordingly always assume that state to be known exactly.
% Source PDF page 241 / original page 234.
\hypertarget{origpage-234}{}%

Let us return to the task formulated at the end of \hyperref[sec:VI.1]{VI.1}. The systems $I,II,III$ are to have the meanings given there; for the quantum-mechanical investigation of $I,II$ we use the notation of \hyperref[sec:VI.2]{VI.2}, while $III$ remains outside the calculations (cf.\ the discussion of this in \hyperref[sec:VI.1]{VI.1}). Let $A$ be the quantity actually to be measured (in $I$), and $\varphi_1(q),\varphi_2(q),\ldots$ its eigenfunctions. Let $I$ be in the state $\varphi(q)$.

If $I$ is the observed system and $II+III$ the observer, we have to apply process~\hyperref[eq:V.1.2]{\textbf{2}},\tnote{The source prints ``process 2'' here; the measurement transformation described in the same sentence is process 1 in the terminology of \hyperref[sec:V.1]{V.1} and \hyperref[sec:VI.1]{VI.1}. The printed numeral is retained.} and find that the measurement takes $I$ from the state $\varphi$ into one of the states $\varphi_n$ ($n=1,2,\ldots$), with respective probabilities $|(\varphi,\varphi_n)|^2$ ($n=1,2,\ldots$). How, then, is the description formulated when $I+II$ is the observed system and only $III$ is the observer?

In this description, $II$ is a measuring instrument whose scale indicates the value of $A$ in $I$. The pointer position on that scale is a physical quantity $B$ in $II$, and it is $B$ that $III$ actually observes. (If $II$ is already inside the observer's body, the scale and pointer position are replaced by their physiological counterparts, such as the retina and the image on it.) The values of $B$ are coupled one-to-one with those of $A$. Let the values of $A$ be $a_1,a_2,\ldots$, and those of $B$ be $b_1,b_2,\ldots$, numbered so that $a_n$ corresponds to $b_n$.

Initially $I$ is in the unknown state $\varphi(q)$ and $II$ in the known state $\xi(r)$, so that $I+II$ is in the state $\Phi(q,r)=\varphi(q)\xi(r)$. The measurement, insofar as it is performed by $II$ on $I$, is effected, as in the previous example, by an energy operator $\boldsymbol H$ (in $I+II$) during the time $t$: this is process~\hyperref[eq:V.1.2]{\textbf{2}}, which transforms $\Phi$ into $\Phi'=e^{-\frac{2\pi i}{h}t\boldsymbol H}\Phi$. From the standpoint of observer $III$, we can speak of a measurement only if the following holds: were $III$ to measure the simultaneously measurable quantities $A,B$ (in $I$ and $II$, respectively, or both in $I+II$) by process~\hyperref[eq:V.1.1]{\textbf{1}}, pairs of values $a_m,b_n$ with $m\ne n$ would have probability zero, while those with $m=n$ would have certain probabilities $w_n$. That is, it suffices to ``look at'' $II$, and $A$ in $I$ has been measured. Quantum mechanics further requires $w_n=|(\varphi,\varphi_n)|^2$.

When this has been achieved, the measurement process, insofar as it occurs in $II$, has been theoretically ``explained''; that is, the boundary discussed in \hyperref[sec:VI.1]{VI.1} has been shifted from $I\mid II+III$ to $I+II\mid III$.

The mathematical task is therefore as follows. A complete orthonormal system $\varphi_1,\varphi_2,\ldots$ in $I$ is given. Find a similar system $\xi_1,\xi_2,\ldots$ in $\mathfrak R^{II}$, a state $\xi$ in $\mathfrak R^{II}$, an energy operator $\boldsymbol H$ in $\mathfrak R^{I+II}$, and a $t$, such that the following holds. If $\varphi$ is an arbitrary state in $\mathfrak R^I$ and we put $\Phi(q,r)=\varphi(q)\xi(r)$, $\Phi'(q,r)=e^{-\frac{2\pi i}{h}t\boldsymbol H}\Phi(q,r)$, then $\Phi'(q,r)$ is to have the form $\sum_{n=1}^{\infty}c_n\varphi_n(q)\xi_n(r)$ (the $c_n$ naturally depending on $\varphi$), with $|c_n|^2=|(\varphi,\varphi_n)|^2$.
% Source PDF page 242 / original page 235.
\hypertarget{origpage-235}{}%
(The equivalence of this last requirement to the physical demand formulated above was discussed in \hyperref[sec:VI.2]{VI.2}.)

In what follows, besides $\varphi_1,\varphi_2,\ldots$, we shall also prescribe $\xi_1,\xi_2,\ldots$ and $\xi$, and seek the unitary operator $\Delta=e^{-\frac{2\pi i}{h}t\boldsymbol H}$ instead of $\boldsymbol H$.

The mathematical problem leads us back to the one solved in \hyperref[sec:VI.2]{VI.2}: there $\Phi'$ was given, and we showed the existence of $c_n,\varphi_n,\xi_n$; now $\varphi_n,\xi_n$ are fixed, and $\Phi,c_n$ are given as functions of $\varphi$, and we must determine $\Delta$ so that the expansion of $\Phi'=\Delta\Phi$ has precisely these coefficients $c_n$ and functions $\varphi_n,\xi_n$.

We shall show that such a $\Delta$ can indeed be found. Our concern is only with the question of principle: whether any such $\Delta$ exists. We shall not, however, concern ourselves with the further question whether the $\Delta=e^{-\frac{2\pi i}{h}t\boldsymbol H}$ of the simplest intuitively comprehensible experimental arrangements (for example, those of \hyperref[sec:III.4]{III.4}) also have this property. We have seen that our conditions are equivalent to the intuitive criterion for a physical intervention to count as a measurement. The experimental arrangements just mentioned certainly meet that criterion; quantum mechanics would therefore have to yield results grossly contradicting intuition if these $\Delta$ did not satisfy our condition, at least approximately.\origNote{215} Thus, in what follows we shall merely specify an abstract $\Delta$ which satisfies our conditions exactly.

Let $\varphi_m$ ($m=0,\pm1,\pm2,\ldots$) and $\xi_n$ ($n=0,\pm1,\pm2,\ldots$) be two given complete orthonormal systems in $\mathfrak R^I$ and $\mathfrak R^{II}$, respectively (we let $m,n$ range over $0,\pm1,\pm2,\ldots$, instead of $1,2,\ldots$, for merely technical reasons of no importance in principle), and for simplicity let the state $\xi$ equal $\xi_0$. Define $\Delta$ by
\[
 \Delta\sum_{m,n=-\infty}^{\infty}x_{mn}\varphi_m(q)\xi_n(r)
 =\sum_{m,n=-\infty}^{\infty}x_{mn}\varphi_m(q)\xi_{m+n}(r).
\]
Since the $\varphi_m(q)\xi_n(r)$, as well as the $\varphi_m(q)\xi_{m+n}(r)$, form a complete orthonormal system in $\mathfrak R^{I+II}$, this $\Delta$ is unitary. Now
\[
 \varphi(q)=\sum_{m=-\infty}^{\infty}(\varphi,\varphi_m)\varphi_m(q),\qquad
 \xi(r)=\xi_0(r),
\]
and hence\tnote{In the first of the following two lines, the source prints $\xi(q)$ instead of $\xi(r)$. This argument is retained as printed.}
\begin{align*}
 \Phi(q,r)&=\varphi(q)\xi(q)
 =\sum_{m=-\infty}^{\infty}(\varphi,\varphi_m)\varphi_m(q)\xi_0(r),\\
 \Phi'(q,r)&=\Delta\Phi(q,r)
 =\sum_{m=-\infty}^{\infty}(\varphi,\varphi_m)\varphi_m(q)\xi_m(r).
\end{align*}
Our aim is thereby achieved: we even have $c_n=(\varphi,\varphi_n)$.
% Source PDF page 243 / original page 236.
\hypertarget{origpage-236}{}%

To see more clearly how this process works, let us give an example using explicit \textsc{Schrödinger} wave functions and specify $\boldsymbol H$ itself rather than $\Delta$.

Let both the observed object and the observer (that is, $I$ and $II$, respectively) be characterized by a single variable $q$ or $r$, ranging continuously from $-\infty$ to $+\infty$. Thus both are to be thought of as points free to move along a line, and their wave functions always have the form $\psi(q)$ and $\eta(r)$, respectively. We assume their masses $m_1,m_2$ to be so large that the kinetic-energy part of the energy operator,
\[
 \frac{1}{2m_1}\left(\frac{h}{2\pi i}\frac{\partial}{\partial q}\right)^2
 +\frac{1}{2m_2}\left(\frac{h}{2\pi i}\frac{\partial}{\partial r}\right)^2,
\]
may be neglected. Only the interaction-energy term of $\boldsymbol H$, which governs the measurement, then remains. We choose it to have the particular form $\frac{h}{2\pi i}q\frac{\partial}{\partial r}$.

The time-dependent \textsc{Schrödinger} differential equation for the $I+II$ wave function $\psi_t=\psi_t(q,r)$ is then
\begin{align*}
 \frac{h}{2\pi i}\frac{\partial}{\partial t}\psi_t(q,r)
 &=-\frac{h}{2\pi i}q\frac{\partial}{\partial r}\psi_t(q,r),\\
 \left(\frac{\partial}{\partial t}+q\frac{\partial}{\partial r}\right)\psi_t(q,r)&=0,
\end{align*}
that is,
\[
 \psi_t(q,r)=f(q,r-tq).
\]
If $\psi_0(q,r)=\Phi(q,r)$ for $t=0$, then $f(q,r)=\Phi(q,r)$, and therefore
\[
 \psi_t(q,r)=\Phi(q,r-tq).
\]
In particular, if the initial states of $I,II$ are represented by $\varphi(q)$ and $\xi(r)$, respectively, then in our scheme of calculation (choosing the time $t$ occurring there equal to $1$)
\begin{align*}
 \Phi(q,r)&=\varphi(q)\xi(r),\\
 \Phi'(q,r)&=\psi_1(q,r)=\varphi(q)\xi(r-q).
\end{align*}
We shall now show that this can be used for a position measurement of $I$ by $II$, that is, that the coordinates $q,r$ are coupled. (Since $q,r$ have continuous spectra and hence can be measured arbitrarily accurately, but not absolutely exactly, this can succeed only approximately.)

For this purpose we assume that $\xi(r)$ differs from zero only in a very narrow interval $-\varepsilon<r<\varepsilon$ (that is, the observer's coordinate $r$ before the measurement is known very accurately). Furthermore, $\xi$ must of course be normalized:
\[
 \|\xi\|=1,\qquad\text{that is,}\qquad\int|\xi(r)|^2\,dr=1.
\]
% Source PDF page 244 / original page 237.
\hypertarget{origpage-237}{}%

The probability that $q$ lies in the interval $q_0-\delta<q<q_0+\delta$, and $r$ in the interval $r_0-\delta'<r<r_0+\delta'$, is
\begin{multline*}
 \int_{q_0-\delta}^{q_0+\delta}\int_{r_0-\delta'}^{r_0+\delta'}
 |\Phi'(qr)|^2\,dq\,dr\\
 =\int_{q_0-\delta}^{q_0+\delta}\int_{r_0-\delta'}^{r_0+\delta'}
 |\varphi(q)|^2|\xi(r-q)|^2\,dq\,dr.
\end{multline*}
If $q_0,r_0$ differ by more than $\delta+\delta'+\varepsilon$, this is zero: thus $q,r$ are so closely coupled that the difference can never exceed $\delta+\delta'+\varepsilon$. And for $r_0=q_0$, if we choose $\delta'\geq\delta+\varepsilon$, the assumptions on $\xi$ make this equal to $\int_{q_0-\delta}^{q_0+\delta}|\varphi(q)|^2\,dq$. Since, however, we can choose $\delta,\delta',\varepsilon$ arbitrarily small (they need only be positive!), this means that $q,r$ are coupled with arbitrary accuracy, and the probability density has the value required by quantum mechanics, $|\varphi(q)|^2$.

We have therefore realized the measurement conditions discussed in \hyperref[sec:VI.1]{VI.1} and in this section.

More complicated examples, such as an analogue of our four-part example in \hyperref[sec:VI.1]{VI.1}, or a check by a second observer $III$ on the correctness of the measurement performed by $II$ on $I$, can be discussed analogously. This is left to the reader.


% Consolidated source: notes.tex

\end{document}
